
\documentclass{book}
%%%%%%%%%%%%%%%%%%%%%%%%%%%%%%%%%%%%%%%%%%%%%%%%%%%%%%%%%%%%%%%%%%%%%%%%%%%%%%%%%%%%%%%%%%%%%%%%%%%%%%%%%%%%%%%%%%%%%%%%%%%%%%%%%%%%%%%%%%%%%%%%%%%%%%%%%%%%%%%%%%%%%%%%%%%%%%%%%%%%%%%%%%%%%%%%%%%%%%%%%%%%%%%%%%%%%%%%%%%%%%%%%%%%%%%%%%%%%%%%%%%%%%%%%%%%
\usepackage{amssymb}
\usepackage{amsfonts}
\usepackage{amsmath}

\setcounter{MaxMatrixCols}{10}
%TCIDATA{OutputFilter=LATEX.DLL}
%TCIDATA{Version=5.50.0.2960}
%TCIDATA{<META NAME="SaveForMode" CONTENT="1">}
%TCIDATA{BibliographyScheme=Manual}
%TCIDATA{Created=Tuesday, April 15, 2025 16:35:49}
%TCIDATA{LastRevised=Monday, August 17, 2026 16:16:01}
%TCIDATA{<META NAME="GraphicsSave" CONTENT="32">}
%TCIDATA{<META NAME="DocumentShell" CONTENT="Standard LaTeX\Standard LaTeX Book">}
%TCIDATA{Language=American English}
%TCIDATA{CSTFile=40 LaTeX Book.cst}
%TCIDATA{ComputeGeneralSettings=0,5,5,0,0,0,1}

\newtheorem{theorem}{Theorem}
\newtheorem{acknowledgement}[theorem]{Acknowledgement}
\newtheorem{algorithm}[theorem]{Algorithm}
\newtheorem{axiom}[theorem]{Axiom}
\newtheorem{case}[theorem]{Case}
\newtheorem{claim}[theorem]{Claim}
\newtheorem{conclusion}[theorem]{Conclusion}
\newtheorem{condition}[theorem]{Condition}
\newtheorem{conjecture}[theorem]{Conjecture}
\newtheorem{corollary}[theorem]{Corollary}
\newtheorem{criterion}[theorem]{Criterion}
\newtheorem{definition}[theorem]{Definition}
\newtheorem{example}[theorem]{Example}
\newtheorem{exercise}[theorem]{Exercise}
\newtheorem{lemma}[theorem]{Lemma}
\newtheorem{notation}[theorem]{Notation}
\newtheorem{problem}[theorem]{Problem}
\newtheorem{proposition}[theorem]{Proposition}
\newtheorem{remark}[theorem]{Remark}
\newtheorem{solution}[theorem]{Solution}
\newtheorem{summary}[theorem]{Summary}
\newenvironment{proof}[1][Proof]{\noindent\textbf{#1.} }{\ \rule{0.5em}{0.5em}}
\input{tcilatex}
\begin{document}

\frontmatter
\title{{\Large Fundamentals of Physics II}\\
{\Large Wave Motion, Optics, and Thermal Physics}}
\author{R. Todd Lines \and (%
%TCIMACRO{\TeXButton{$\copyright$}{$\copyright$}}%
%BeginExpansion
$\copyright$%
%EndExpansion
copyright, R. Todd Lines, All Rights Reserved)}
\date{04/18/2026}
\maketitle
\tableofcontents

\chapter{Preface}

This set of notes is intended to be an aid to the student. It is what I
intend to present in class. There are likely errors and mistakes, so use
these notes, but don't expect perfection. If there are things that are
confusing, please talk to me or ask questions in class.

I consulted many authors and colleagues in creating these notes. Principle
among the authors was Serway and Jewett and Haliday and Resnic. Colleagues
of special note are Kevin Kelly, Brian Pyper, Steve Tercotte, and Ryan
Nielson. I especially appreciate the PH123 students of the Fall of 2006.
They were the first to test run this sets of notes, and I appreciate their
feedback (and patience).

This work is 
%TCIMACRO{\TeXButton{$\copyright$}{$\copyright$} }%
%BeginExpansion
$\copyright$
%EndExpansion
Copyright R. Todd Lines. Permission is given to BYU-Idaho for educational
purposes. All other rights are reserved.

\mainmatter

%TCIMACRO{\TeXButton{\pagenumbering{arabic}}{\pagenumbering{arabic}}}%
%BeginExpansion
\pagenumbering{arabic}%
%EndExpansion

\part{Oscillation and Waves}

\chapter{Introduction to the Course, Simple Harmonic Motion 1.15.1}

%TCIMACRO{%
%\TeXButton{\chaptermark{Simple Harmonic Motion}}{\chaptermark{Simple Harmonic Motion}}}%
%BeginExpansion
\chaptermark{Simple Harmonic Motion}%
%EndExpansion

Last semester, in Principles of Physics I (PH121) you studied how things
move. We identified a moving object (I often refer to this object as the 
\emph{mover object})\emph{\ }and other objects that exerted forces on the
mover (I often refer to these other objects as the \emph{environmental
objects})\emph{.} You learned about forces and torques which get mover
objects moving. You should remember Newton's Second Law and Newton's Second
Law for rotation.

You also learned and practiced a lot of math. We will continue to use the
math you learned in PH121 this semester.

But we will go beyond what we learned in Principles of Physics I to study
new types of motion, and new objects that move.

This semester, we will start with a very special type of motion. It is the
motion that results from oscillation. We call this very special type of
motion, \emph{simple harmonic motion}.

Simple harmonic motion (SHM) means a motion that repeats in the most
special, simple way.

%TCIMACRO{%
%\TeXButton{Fundamental Concepts}{\vspace{0.10in}\hspace{-0.37in}{\Large Fundamental Concepts\vspace{0.2in}}}}%
%BeginExpansion
\vspace{0.10in}\hspace{-0.37in}{\Large Fundamental Concepts\vspace{0.2in}}%
%EndExpansion

\begin{itemize}
\item Simple Harmonic Motion

\item Frequency and Period

\item Mathematical description of Simple Harmonic Motion
\end{itemize}

\section{1 Oscillation and Simple Harmonic Motion}

Simple harmonic motion is a very specific kind of motion. Some
characteristics of this type of motion are as follows:

1.\qquad The motion repeats in a regular way, like a grandfather clock
pendulum swings back and forth in a set amount of time. This set amount of
time is called a period.

2.\qquad The mover object moves about a center position and is symmetric
about that position. This center position is the bottom of the swing for our
grandfather clock. The idea of symmetry just means that the pendulum reaches
the same height on both sides as it swings

3.\qquad The mover object's motion can be traced out in a position vs. time
graph. For SHM, the shape of the position vs. time graph is a sinusoid.

4.\qquad The velocity of the mover object constantly changes. It is zero at
the extreme points (the points farthest from the center, or the largest
positive and the largest negative displacements). That is where it turns
around and goes back the other way. It stops there, but just for a split
second. The velocity is largest at the equilibrium position (the center
position). If we plotted a velocity vs. time graph, the velocity would also
be a sinusoid or snake-like shape.

5.\qquad The acceleration also constantly changes. And it is also a sinusoid.

That is a lot of criteria for a motion that is supposed to be simple! Let's
take an example system to see how simple harmonic motion works. I have drawn
a mass (marked $m$) attached to a spring. And lets assume that the mass is
on a frictionless surface and there is no air drag force. If we pull the
mass so the string stretches, we would get simple harmonic motion.

\FRAME{dtbpF}{2.5261in}{1.9484in}{0pt}{}{}{Figure}{\special{language
"Scientific Word";type "GRAPHIC";maintain-aspect-ratio TRUE;display
"USEDEF";valid_file "T";width 2.5261in;height 1.9484in;depth
0pt;original-width 1.2834in;original-height 0.9833in;cropleft "0";croptop
"1";cropright "1";cropbottom "0";tempfilename
'SUS73U00.wmf';tempfile-properties "XPR";}}The center position for our mass
is called the Equilibrium position, and for convenience we often define it
as the origin of our coordinate system.

\begin{center}
\rule{300pt}{1pt}

{\small Equilibrium Position: The position of the mover mass (not the
spring, which is an environmental object acting on the mover mass) when the
spring is neither stretched nor compressed.}

\rule{300pt}{1pt}
\end{center}

Let's draw a picture of simple harmonic motion. First we need a device to do
this. Suppose you go to the supermarket\footnote{%
Like Albertsons, where they still have checkers.}. But instead of putting
your ramen and granola bars on the belt at the checkout counter, you strap
on a device like this\footnote{%
Please don't really try this at the local supermarkets. They don't seem to
have any sense scientific inquiry or even a sense of humor about such things
at all.}\FRAME{dhF}{4.3474in}{1.5904in}{0pt}{}{}{Figure}{\special{language
"Scientific Word";type "GRAPHIC";maintain-aspect-ratio TRUE;display
"USEDEF";valid_file "T";width 4.3474in;height 1.5904in;depth
0pt;original-width 4.2964in;original-height 1.5549in;cropleft "0";croptop
"1";cropright "1";cropbottom "0";tempfilename
'SUS73U01.wmf';tempfile-properties "XPR";}}You can see the mover object on
the spring. But we have placed a pen on the object and that pen is tracing
out a pattern as the belt moves. You will recognize this pattern as a trig
function. The result might look something like this if you removed the belt.%
\FRAME{dhF}{3.9418in}{1.4226in}{0pt}{}{}{Figure}{\special{language
"Scientific Word";type "GRAPHIC";maintain-aspect-ratio TRUE;display
"USEDEF";valid_file "T";width 3.9418in;height 1.4226in;depth
0pt;original-width 3.8934in;original-height 1.388in;cropleft "0";croptop
"1";cropright "1";cropbottom "0";tempfilename
'SUS73U02.wmf';tempfile-properties "XPR";}}Of course, what we have made is a
position vs. time graph. We remember these from Principles of Physics I!
This gives us a record of the past motion of the object. \FRAME{dtbpFX}{%
3.9583in}{1.8542in}{0pt}{}{}{Plot}{\special{language "Scientific Word";type
"MAPLEPLOT";width 3.9583in;height 1.8542in;depth 0pt;display
"USEDEF";plot_snapshots TRUE;mustRecompute FALSE;lastEngine "MuPAD";xmin
"0";xmax "4.7";xviewmin "0";xviewmax "4.7";yviewmin "-10";yviewmax
"10";viewset"XY";rangeset"X";plottype 4;labeloverrides 3;x-label "t
(s)";y-label "x (cm)";axesFont "Times New
Roman,12,0000000000,useDefault,normal";numpoints 100;plotstyle
"patch";axesstyle "normal";axestips FALSE;xis \TEXUX{t};var1name
\TEXUX{$t$};function \TEXUX{$5\cos \left( 2t+0\right) $};linecolor
"blue";linestyle 1;pointstyle "point";linethickness 1;lineAttributes
"Solid";var1range "0,4.7";num-x-gridlines 100;curveColor
"[flat::RGB:0x000000ff]";curveStyle "Line";VCamFile
'SUS73U4H.xvz';valid_file "T";tempfilename
'SUS73U03.wmf';tempfile-properties "XPR";}}where in this graph, $x_{\max }=5%
\unit{cm}.$ Having the power of mathematics, we know we can write an
equation that would describe this curve. From your Trigonometry experience,
we can guess that an equation for our mover object motion might look
something like%
\begin{equation*}
x\left( t\right) =x_{\max }\cos \left( \theta \right)
\end{equation*}%
Notice that the instantaneous position, $x\left( t\right) $ has to be a
function of time. Back in trigonometry we would have said a cosine function
was a function of an angle, $\theta .$ But we know this is a position vs.
time graph, so our angle must be different for different times. Let's write 
\begin{equation*}
\theta \left( t\right) =\omega t
\end{equation*}%
as strange sort of an angle. This \textquotedblleft angle\textquotedblright\
is a function of time, and let's say how fast our \textquotedblleft
angle\textquotedblright\ is changing is given by 
\begin{equation*}
\frac{d\theta }{dt}=\omega
\end{equation*}%
This is a sort of speed for how fast our angle is changing.

Using our rate of angle change, we can write our equation for simple
harmonic motion as 
\begin{equation*}
x\left( t\right) =x_{\max }\cos \left( \omega t\right)
\end{equation*}%
In the next figure $x\left( t\right) =x_{\max }\cos \left( \omega t\right) $
is plotted with two different values for $\omega .$ \FRAME{dtbpFX}{3.9583in}{%
1.3578in}{0pt}{}{}{Plot}{\special{language "Scientific Word";type
"MAPLEPLOT";width 3.9583in;height 1.3578in;depth 0pt;display
"USEDEF";plot_snapshots TRUE;mustRecompute FALSE;lastEngine "MuPAD";xmin
"0";xmax "4.7";xviewmin "0";xviewmax "4.7";yviewmin "-10";yviewmax
"10";viewset"XY";rangeset"X";plottype 4;labeloverrides 3;x-label "t
(s)";y-label "x (cm)";axesFont "Times New
Roman,12,0000000000,useDefault,normal";numpoints 100;plotstyle
"patch";axesstyle "normal";axestips FALSE;xis \TEXUX{t};var1name
\TEXUX{$t$};function \TEXUX{$5\cos \left( 2t+0\right) $};linecolor
"blue";linestyle 1;pointstyle "point";linethickness 1;lineAttributes
"Solid";var1range "0,4.7";num-x-gridlines 100;curveColor
"[flat::RGB:0x000000ff]";curveStyle "Line";function \TEXUX{$5\cos \left(
4t+0\right) $};linecolor "maroon";linestyle 2;pointstyle
"point";linethickness 1;lineAttributes "Dash";var1range
"0,4.7";num-x-gridlines 100;curveColor "[flat::RGB:0x00800000]";curveStyle
"Line";VCamFile 'SUS75467.xvz';valid_file "T";tempfilename
'SUS73U04.wmf';tempfile-properties "XPR";}}We can see what $\omega $ does
for us. It stretches out or compresses our curve. In the blue (dashed) curve
our angle changes quickly. In the red (solid) curve the angle changes
slowly. Note that we are not plotting position vs. angle. Both plots would
be the same if we did! We are plotting position vs. time. And for the blue
curve our mass is oscillating much more quickly than it is for the red
curve. The quantity $\omega $ tells us something about how fast our mover
object is oscillating.

The quantity $x_{\max }$ is the maximum displacement of our mover object
from the equilibrium position. This maximum displacement from equilibrium
has a name. It is called the \emph{amplitude}. In our equation I\ have used $%
x_{\max }$ where most books use the letter $A$ for amplitude\emph{\ }to
emphasize that the amplitude is the maximum displacement of the mover mass
from the mover mass equilibrium position. In our coordinate system, the mass
is going back and forth in the $x$ direction. Since the amplitude means the
maximum displacement, $x_{\max }$ is a good way to write amplitude. Here is
the same graph but with $x_{\max }=9\unit{cm}$\FRAME{dtbpFX}{3.9583in}{%
1.3578in}{0pt}{}{}{Plot}{\special{language "Scientific Word";type
"MAPLEPLOT";width 3.9583in;height 1.3578in;depth 0pt;display
"USEDEF";plot_snapshots TRUE;mustRecompute FALSE;lastEngine "MuPAD";xmin
"0";xmax "4.7";xviewmin "0";xviewmax "4.7";yviewmin "-10";yviewmax
"10";viewset"XY";rangeset"X";plottype 4;labeloverrides 3;x-label "t
(s)";y-label "x (cm)";axesFont "Times New
Roman,12,0000000000,useDefault,normal";numpoints 100;plotstyle
"patch";axesstyle "normal";axestips FALSE;xis \TEXUX{t};var1name
\TEXUX{$t$};function \TEXUX{$9\cos \left( 2t+0\right) $};linecolor
"blue";linestyle 1;pointstyle "point";linethickness 1;lineAttributes
"Solid";var1range "0,4.7";num-x-gridlines 100;curveColor
"[flat::RGB:0x000000ff]";curveStyle "Line";function \TEXUX{$9\cos \left(
4t+0\right) $};linecolor "maroon";linestyle 2;pointstyle
"point";linethickness 1;lineAttributes "Dash";var1range
"0,4.7";num-x-gridlines 100;curveColor "[flat::RGB:0x00800000]";curveStyle
"Line";VCamFile 'SUS76168.xvz';valid_file "T";tempfilename
'SUS73U05.wmf';tempfile-properties "XPR";}}The two graphs for the two
different $\omega $ values are not more stretched out in time, but now they
are taller along the position axis. This means that the mass is moving
farther from the center as it oscillates.

You will remember from your trigonometry class that the \emph{period,} $T,$
tells us the time it takes for the oscillation to go through a complete
cycle. A complete cycle is when the object, say, goes from $x_{\max }$ to $%
-x_{\max }$ and then back to $x_{\max }.$ You can probably guess that how
long it takes to oscillate and how often it oscillates would be related. How
often the oscillator completes a cycle is called the frequency. The longer
the period, the lower the frequency. 
\begin{equation*}
f=\frac{1}{T}
\end{equation*}%
Think of cars passing you on your way to class. Period is like how long you
wait in between cars. Frequency is like how often cars pass. If you wait
less time between cars, the cars pass more frequently. And that is just what
our equation says! We can see from our graphs that our stretching quantity $%
\omega ,$ must be related to the frequency of our oscillation. If the
frequency is high, then $\omega $ must be large so that we reach different
\textquotedblleft angles\textquotedblright\ faster. But for the cosine
function to work we need angle units. \emph{We will choose radians} for our
units and we will write our stretching quantity as 
\begin{equation*}
\omega =2\pi f
\end{equation*}%
This works! If $\omega $ is bigger then our oscillation happens more
frequently. The $2\pi $ has units of radians. So $\omega $ has units of $%
\unit{rad}\unit{Hz}$ or more commonly $\unit{rad}/\unit{s}.$ That matches
our derivative above. Let's give a name to the quantity $\omega .$ Since it
has radians in it we might guess that it has something to do with circular
motion (more on this later) and it has frequency in it. So we will call $%
\omega $ the \emph{angular frequency}.

\section{Velocity and Acceleration}

So far we have guessed the descriptive equation for SHM. 
\begin{equation}
x\left( t\right) =x_{\max }\cos \left( \omega t\right)
\end{equation}%
This gives us the position of the particle at any given time. Since knowing
the position as a function of time is a good description of the motion of
our mover mass, we can call this equation the \emph{equation of motion} for
our mass. We will see soon that this equation is correct. But let's pretend
that we are early researchers and we just know the equation makes the right
shape. Still, we can learn much from knowing this equation. Since we know
how to take derivatives, and we know the derivative of position with respect
to time is the velocity, we can see that the velocity of the mass at any
given time is given by

\begin{equation}
v\left( t\right) =\frac{dx\left( t\right) }{dt}=-\omega x_{\max }\sin \left(
\omega t\right)
\end{equation}%
From our fond memories of our trigonometry class we know the maximum of a
sine function is always $1.\strut $\FRAME{dtbpFX}{4.4996in}{1.0421in}{0pt}{}{%
}{Plot}{\special{language "Scientific Word";type "MAPLEPLOT";width
4.4996in;height 1.0421in;depth 0pt;display "USEDEF";plot_snapshots
TRUE;mustRecompute FALSE;lastEngine "MuPAD";xmin "-10.001000";xmax
"10.001000";xviewmin "-10.001000";xviewmax "10.001000";yviewmin
"-1.000187";yviewmax "1.000187";viewset"XY";rangeset"X";plottype
4;labeloverrides 3;x-label "x";y-label "sin(x)";axesFont "Times New
Roman,12,0000000000,useDefault,normal";numpoints 100;plotstyle
"patch";axesstyle "normal";axestips FALSE;xis \TEXUX{x};var1name
\TEXUX{$x$};function \TEXUX{$\sin x$};linecolor "blue";linestyle
1;pointstyle "point";linethickness 1;lineAttributes "Solid";var1range
"-10.001000,10.001000";num-x-gridlines 100;curveColor
"[flat::RGB:0x000000ff]";curveStyle "Line";function \TEXUX{$1$};linecolor
"black";linestyle 2;pointstyle "point";linethickness 1;lineAttributes
"Dash";var1range "-10.001000,10.001000";num-x-gridlines 100;curveColor
"[flat::RGB:0000000000]";curveStyle "Line";function \TEXUX{$-1$};linecolor
"black";linestyle 2;pointstyle "point";linethickness 1;lineAttributes
"Dash";var1range "-10.001000,10.001000";num-x-gridlines 100;curveColor
"[flat::RGB:0000000000]";curveStyle "Line";VCamFile
'SUS76269.xvz';valid_file "T";tempfilename
'SUS73U06.wmf';tempfile-properties "XPR";}}Notice that the $\omega $ and the 
$x_{\max }$ aren't changing for our oscillation mover mass. They are
constant. Then if we want the maximum speed we can simply set the $\sin
\left( \omega t\right) =1.$ Then the maximum speed will be 
\begin{equation}
v_{\max }=\omega x_{\max }\left( 1\right)
\end{equation}%
Notice that this does not tell us when the speed is maximum, just what the
maximum speed is. We then have an equation for the speed of our object as a
function of time. 
\begin{equation*}
v\left( t\right) =-v_{\max }\sin \left( \omega t\right)
\end{equation*}%
From this equation we could find when the speed is maximum. We will often
use this trick of knowing the maximum of sine is one.

We can also find the acceleration. We just take another derivative.%
\begin{eqnarray*}
a\left( t\right) &=&\frac{dv}{dt} \\
&=&\frac{d^{2}x\left( t\right) }{dt^{2}} \\
&=&-\omega ^{2}x_{\max }\cos \left( \omega t\right)
\end{eqnarray*}%
This is the acceleration of the object attached to the spring. It's also
true that the maximum for a cosine function is $1,$ so the maximum
acceleration would be

\begin{equation}
a_{\max }=\omega ^{2}x_{\max }\left( 1\right)
\end{equation}%
so we could write our instantaneous acceleration as 
\begin{equation*}
a\left( t\right) =-a_{\max }\cos \left( \omega t\right)
\end{equation*}

These are significant results, so let's summarize. For a simple harmonic
oscillator, the instantaneous position, speed, and acceleration are given by%
\begin{eqnarray}
x\left( t\right) &=&x_{\max }\cos \left( \omega t\right) \\
v\left( t\right) &=&-\omega x_{\max }\sin \left( \omega t\right)  \notag \\
a\left( t\right) &=&-\omega ^{2}x_{\max }\cos \left( \omega t\right)  \notag
\end{eqnarray}

Let's plot $x\left( t\right) ,$ $v\left( t\right) ,$ and $a\left( t\right) $
for a specific case\FRAME{dtbpFX}{3.6357in}{2.4163in}{0pt}{}{}{Plot}{\special%
{language "Scientific Word";type "MAPLEPLOT";width 3.6357in;height
2.4163in;depth 0pt;display "USEDEF";plot_snapshots TRUE;mustRecompute
FALSE;lastEngine "MuPAD";xmin "-5";xmax "5";xviewmin "-5";xviewmax
"5";yviewmin "-20";yviewmax "20";viewset"XY";rangeset"X";plottype
4;labeloverrides 3;x-label "t";y-label "x";axesFont "Times New
Roman,12,0000000000,useDefault,normal";numpoints 100;plotstyle
"patch";axesstyle "normal";axestips FALSE;xis \TEXUX{t};var1name
\TEXUX{$t$};function \TEXUX{$5\cos \left( 2t+0\right) $};linecolor
"blue";linestyle 1;pointstyle "point";linethickness 3;lineAttributes
"Solid";var1range "-5,5";num-x-gridlines 100;curveColor
"[flat::RGB:0x000000ff]";curveStyle "Line";function \TEXUX{$-2\ast 5\sin
\left( 5t+0\right) $};linecolor "green";linestyle 2;pointstyle
"point";linethickness 3;lineAttributes "Dash";var1range
"-5,5";num-x-gridlines 100;curveColor "[flat::RGB:0x00008000]";curveStyle
"Line";function \TEXUX{$-\left( 2\right) ^{2}\left( 5\right) \cos \left(
2t+0\right) $};linecolor "red";linestyle 4;pointstyle "point";linethickness
3;lineAttributes "DotDash";var1range "-5,5";num-x-gridlines 100;curveColor
"[flat::RGB:0x00ff0000]";curveStyle "Line";VCamFile
'SUS7626A.xvz';valid_file "T";tempfilename
'SUS73U07.wmf';tempfile-properties "XPR";}}Red (solid) is the displacement,
green (dashed) is the velocity, and blue (dot-dashed) is the acceleration.
Note that each as a different maximum amplitude. Also note that they don't
rise and fall at the same time. We will describe this as being \emph{not in
phase}. We can see this again in the next graph. The red graph (top graph)
is still the displacement. The bottom graph is the velocity. And note that
the peaks are not at the same time.\FRAME{dtbpF}{2.9888in}{2.8409in}{0pt}{}{%
}{Figure}{\special{language "Scientific Word";type
"GRAPHIC";maintain-aspect-ratio TRUE;display "USEDEF";valid_file "T";width
2.9888in;height 2.8409in;depth 0pt;original-width 6.5639in;original-height
6.2396in;cropleft "0";croptop "1";cropright "1";cropbottom "0";tempfilename
'SUS73U08.wmf';tempfile-properties "XPR";}}The acceleration is $90\unit{%
%TCIMACRO{\U{b0}}%
%BeginExpansion
{{}^\circ}%
%EndExpansion
}$ \emph{out of phase} from the velocity. Let's think about why this would
be. Suppose we attach a mass to a spring and allow the mass to slide on a
frictionless surface.\FRAME{dtbpF}{2.4111in}{1.7461in}{0pt}{}{}{Figure}{%
\special{language "Scientific Word";type "GRAPHIC";maintain-aspect-ratio
TRUE;display "USEDEF";valid_file "T";width 2.4111in;height 1.7461in;depth
0pt;original-width 7.043in;original-height 5.0903in;cropleft "0";croptop
"1";cropright "1";cropbottom "0";tempfilename
'SUS73U09.wmf';tempfile-properties "XPR";}}Let's start by stretching the
spring by pulling the mass to the right and releasing. This is the situation
in the left hand part of the last figure.

\FRAME{dtbpF}{1.4114in}{1.1692in}{0pt}{}{}{Figure}{\special{language
"Scientific Word";type "GRAPHIC";maintain-aspect-ratio TRUE;display
"USEDEF";valid_file "T";width 1.4114in;height 1.1692in;depth
0pt;original-width 2.1084in;original-height 1.7417in;cropleft "0";croptop
"1";cropright "1";cropbottom "0";tempfilename
'SUS73U0A.wmf';tempfile-properties "XPR";}}The spring is pulling strongly on
the mass. We could draw a free body diagram for this situation. We will need
the spring force \FRAME{dtbpF}{1.5592in}{1.0957in}{0pt}{}{}{Figure}{\special%
{language "Scientific Word";type "GRAPHIC";maintain-aspect-ratio
TRUE;display "USEDEF";valid_file "T";width 1.5592in;height 1.0957in;depth
0pt;original-width 2.1577in;original-height 1.5091in;cropleft "0";croptop
"1";cropright "1";cropbottom "0";tempfilename
'SUS73U0B.wmf';tempfile-properties "XPR";}}Back in Principles of Physics I
(PH121) we said that Hooke's Law is not something that is always true, but
by \textquotedblleft law\textquotedblright\ we mean a mathematical
representation (an equation) that comes from our mental model of how the
universe works. In this case, Hooke's law is an equation that comes from
Hooke's model of how springs work. It is a good model for most springs as
long as we don't stretch them too far. You remember Hooke's law tells us
that the spring force is proportional to how far we stretch or compress the
spring from equilibrium $\left( \Delta x_{e}\right) $ and how stiff the
spring is $\left( k\right) .$

\begin{eqnarray}
F_{s} &=&S=-k\Delta x_{e} \\
&=&-k\left( x-x_{e}\right)  \notag
\end{eqnarray}%
where $x_{e}$ is the \emph{equilibrium position of the mass.} Some books
write $x_{e}$ as $x_{o}$ but the \textquotedblleft e\textquotedblright\ is
better for reminding us that it is the equilibrium position.

If we assume the equilibrium position is at $x=0$ then we can write our
spring force as%
\begin{equation}
S=-kx
\end{equation}%
And if we write out Newton's second law we get 
\begin{eqnarray*}
F_{net_{x}} &=&-ma_{x}=-S_{ms} \\
F_{net_{y}} &=&ma_{y}=N_{mT}-W_{mE}
\end{eqnarray*}%
We can expect that $a_{y}$ should be zero because the mass won't lift off
the table in the $y$-direction. But in the $x$-direction 
\begin{equation*}
-a_{x}=-\frac{S_{ms}}{m}
\end{equation*}%
and knowing that 
\begin{equation*}
S_{ms}=-kx
\end{equation*}%
from Hooke's law, we can see that 
\begin{equation*}
-a_{x}=+\frac{kx}{m}
\end{equation*}%
or 
\begin{equation*}
a_{x}=-\frac{kx}{m}
\end{equation*}%
Since at our starting point $x$ is big our acceleration is big. Notice that
there is a minus sign. The minus sign tells us that $a_{x}$ must be to the
left.

But suppose the mass was to the left of $x_{e}=0$ \FRAME{dtbpF}{2.0643in}{%
0.9582in}{0pt}{}{}{Figure}{\special{language "Scientific Word";type
"GRAPHIC";maintain-aspect-ratio TRUE;display "USEDEF";valid_file "T";width
2.0643in;height 0.9582in;depth 0pt;original-width 2.0254in;original-height
0.9254in;cropleft "0";croptop "1";cropright "1";cropbottom "0";tempfilename
'SUS73U0C.wmf';tempfile-properties "XPR";}}Now $x<0$ so it is negative. That
makes 
\begin{equation*}
a_{x}=-\frac{kx}{m}
\end{equation*}%
positive. The acceleration always seems to point toward $x_{e}=0,$ the
equilibrium position for the mass. This is an important part of the
definition of simple harmonic motion, having an acceleration that always
points toward the equilibrium position. We call a force that makes the
acceleration point toward the equilibrium position \emph{restoring force}.

\begin{center}
{\small \rule{300pt}{1pt}}

{\small Restoring force:\ A force that is always directed toward the
equilibrium position}

{\small \rule{300pt}{1pt}}
\end{center}

Now lets look at our object a short time later \FRAME{dtbpF}{1.4096in}{%
1.1087in}{0pt}{}{}{Figure}{\special{language "Scientific Word";type
"GRAPHIC";maintain-aspect-ratio TRUE;display "USEDEF";valid_file "T";width
1.4096in;height 1.1087in;depth 0pt;original-width 1.375in;original-height
1.075in;cropleft "0";croptop "1";cropright "1";cropbottom "0";tempfilename
'SUS73U0D.wmf';tempfile-properties "XPR";}}We can see that in this position, 
$x=0$ so 
\begin{equation*}
a_{x}=\frac{k\left( 0\right) }{m}=0
\end{equation*}%
so at this position the acceleration is zero. That is because right at $x=0$
the spring is neither pushing nor pulling. There is no net force so no
acceleration when $x=0.$ This matches our graphs of $x\left( t\right) ,$ $%
v\left( t,\right) ,$ and $a\left( t\right) .$ The reason for these three
being out of phase is just because the physics says it must be so.

\section{The Idea of Phase}

We said before that $x\left( t\right) ,$ $v\left( t\right) ,$ and $a\left(
t\right) $ are \textquotedblleft out of phase.\textquotedblright\ Let's look
at the idea of \textquotedblleft phase\textquotedblright\ more carefully.
Suppose you return to the grocery store and start your SHM\ device.\FRAME{dhF%
}{3.2292in}{1.1813in}{0pt}{}{}{Figure}{\special{language "Scientific
Word";type "GRAPHIC";maintain-aspect-ratio TRUE;display "USEDEF";valid_file
"T";width 3.2292in;height 1.1813in;depth 0pt;original-width
4.2964in;original-height 1.5549in;cropleft "0";croptop "1";cropright
"1";cropbottom "0";tempfilename 'SUS73U0E.wmf';tempfile-properties "XPR";}}%
But this time you work with a lab partner, and the lab partner tries to
start a stopwatch when you let go of the mass. But, due to having a slow
reaction time, your partner starts the clock too late. The resulting graph
looks like this\FRAME{dtbpFX}{3.9583in}{1.8542in}{0pt}{}{}{Plot}{\special%
{language "Scientific Word";type "MAPLEPLOT";width 3.9583in;height
1.8542in;depth 0pt;display "USEDEF";plot_snapshots TRUE;mustRecompute
FALSE;lastEngine "MuPAD";xmin "0";xmax "4.7";xviewmin "0";xviewmax
"4.7";yviewmin "-10";yviewmax "10";viewset"XY";rangeset"X";plottype
4;labeloverrides 3;x-label "t";y-label "x";axesFont "Times New
Roman,12,0000000000,useDefault,normal";numpoints 100;plotstyle
"patch";axesstyle "normal";axestips FALSE;xis \TEXUX{t};var1name
\TEXUX{$t$};function \TEXUX{$5\cos \left( 2t+\pi /3\right) $};linecolor
"blue";linestyle 1;pointstyle "point";linethickness 1;lineAttributes
"Solid";var1range "0,4.7";num-x-gridlines 100;curveColor
"[flat::RGB:0x000000ff]";curveStyle "Line";VCamFile
'SUS7656B.xvz';valid_file "T";tempfilename
'SUS73U0F.wmf';tempfile-properties "XPR";}}But we know that the only
difference between this graph and the one we had before is that the lab
partner was slow, so the graph is shifted on the time axes. We expect we can
use the same mathematical model for SHM, but we must need to change
something because the graph looks different.

We know from our algebra classes what a shift looks like. Take the
expression 
\begin{equation*}
y=x^{2}
\end{equation*}%
We can plot this to get a parabola\FRAME{dtbpFX}{2.9093in}{1.9407in}{0pt}{}{%
}{Plot}{\special{language "Scientific Word";type "MAPLEPLOT";width
2.9093in;height 1.9407in;depth 0pt;display "USEDEF";plot_snapshots
TRUE;mustRecompute FALSE;lastEngine "MuPAD";xmin "-5.001000";xmax
"10.001000";xviewmin "-5.001000";xviewmax "10.001000";yviewmin "0";yviewmax
"25.00250";viewset"XY";rangeset"X";plottype 4;axesFont "Times New
Roman,12,0000000000,useDefault,normal";numpoints 100;plotstyle
"patch";axesstyle "normal";axestips FALSE;xis \TEXUX{x};var1name
\TEXUX{$x$};function \TEXUX{$x^{2}$};linecolor "blue";linestyle 1;pointstyle
"point";linethickness 3;lineAttributes "Solid";var1range
"-5.001000,10.001000";num-x-gridlines 100;curveColor
"[flat::RGB:0x000000ff]";curveStyle "Line";VCamFile
'SUXH4A0D.xvz';valid_file "T";tempfilename
'SUXH4A01.wmf';tempfile-properties "XPR";}}A shifted parabola would be
expressed as 
\begin{equation*}
y=\left( x-\phi _{o}\right) ^{2}
\end{equation*}%
where $\phi _{o}$ is the amount of the shift. Suppose $\phi _{o}=2,$ then 
\begin{equation*}
y=\left( x-2\right) ^{2}
\end{equation*}%
and our shifted parabola looks like this\FRAME{dtbpFX}{2.8534in}{1.9035in}{%
0pt}{}{}{Plot}{\special{language "Scientific Word";type "MAPLEPLOT";width
2.8534in;height 1.9035in;depth 0pt;display "USEDEF";plot_snapshots
TRUE;mustRecompute FALSE;lastEngine "MuPAD";xmin "-5.001000";xmax
"10.00100";xviewmin "-5.001000";xviewmax "10.00100";yviewmin
"-0.003982";yviewmax "25.00490";viewset"XY";rangeset"X";plottype 4;axesFont
"Times New Roman,12,0000000000,useDefault,normal";numpoints 100;plotstyle
"patch";axesstyle "normal";axestips FALSE;xis \TEXUX{x};var1name
\TEXUX{$x$};function \TEXUX{$\left( x-2\right) ^{2}$};linecolor
"blue";linestyle 1;pointstyle "point";linethickness 3;lineAttributes
"Solid";var1range "-5.001000,10.00100";num-x-gridlines 100;curveColor
"[flat::RGB:0x000000ff]";curveStyle "Line";VCamFile
'SUXH570E.xvz';valid_file "T";tempfilename
'SUXH3V00.wmf';tempfile-properties "XPR";}}Recall that a shift like 
\begin{equation*}
y=\left( x-\phi _{o}\right) ^{2}
\end{equation*}%
will move the parabola to the right while 
\begin{equation*}
y=\left( x+\phi _{o}\right) ^{2}
\end{equation*}%
will move the parabola to the left.

We can use this idea to write our SHM expression for our situation with the
slow lab partner. 
\begin{equation*}
x\left( t\right) =x_{\max }\cos \left( \omega \left( t\pm \tau _{o}\right)
\right)
\end{equation*}%
In our case, the lab partner was late by $\tau _{o}=\allowbreak
4.\,\allowbreak 712\,4\unit{s}.$ This is not usually how we express the
shift, however. We usually distribute the $\omega $ so our equation looks
like%
\begin{eqnarray*}
x\left( t\right) &=&x_{\max }\cos \left( \omega t\pm \omega \tau _{o}\right)
\\
&=&x_{\max }\cos \left( \omega t\pm \omega \tau _{o}\right)
\end{eqnarray*}%
We usually use the symbol $\phi _{o}=\omega \tau _{o}$ so we will write our
SHM expression for position as a function of time as%
\begin{equation*}
x\left( t\right) =x_{\max }\cos \left( \omega t\pm \phi _{o}\right)
\end{equation*}

We could call $\phi _{o}$ the \emph{slow lab partner constant, }but that is
long and not very kind. So let's call $\phi _{o}$ by the name \emph{phase
constant}. It is also customary to call the entire expression in
parenthesis, $\left( \omega t\pm \phi _{o}\right) ,$ the phase of the cosine
function. This is especially used in the fields of Optics and
Electrodynamics.

From what we did before, we know that $x\left( t\right) ,$ $v\left( t\right)
,$ and $a\left( t\right) $ are out of phase, so there must be a phase
constant involved somehow. Let's look for it in the lectures that follow.

%TCIMACRO{%
%\TeXButton{Basic Equations}{\vspace{0.5in}\hspace{-1.3in}{\LARGE Basic Equations\vspace{0.25in}}}}%
%BeginExpansion
\vspace{0.5in}\hspace{-1.3in}{\LARGE Basic Equations\vspace{0.25in}}%
%EndExpansion
\begin{equation*}
f=\frac{1}{T}
\end{equation*}%
\begin{equation*}
\omega =2\pi f
\end{equation*}

\begin{equation*}
x\left( t\right) =x_{\max }\cos \left( \omega t\pm \phi _{o}\right)
\end{equation*}

\chapter{2 Energy and Dynamics of SHM 1.15.2 1.15,3}

%TCIMACRO{\TeXButton{\chaptermark{SHM Energy}}{\chaptermark{SHM Energy}}}%
%BeginExpansion
\chaptermark{SHM Energy}%
%EndExpansion

Back in Principles of Physics I we started with position, velocity, and
acceleration to describe motion. We have done that again for a simple
harmonic oscillator. In Principles of Physics I, after basic motion, we
found that we could use the idea of a force and Newton's second law to find
the acceleration for our description of motion. Then we changed view points
and used the idea of energy to find motion. The energy picture was easier in
some ways. Let's try to develop and energy picture for simple harmonic
oscillators. Of course the goal in physics is to describe the motion of the
object, so we will say that we are looking for the \emph{dynamics} of simple
harmonic oscillators when we look for the position, velocity, and
acceleration as a function of time.

%TCIMACRO{%
%\TeXButton{Fundamental Concepts}{\vspace{0.10in}\hspace{-0.37in}{\Large Fundamental Concepts\vspace{0.2in}}}}%
%BeginExpansion
\vspace{0.10in}\hspace{-0.37in}{\Large Fundamental Concepts\vspace{0.2in}}%
%EndExpansion

\begin{itemize}
\item Initial Conditions

\item Energy and SHM

\item Equation of motion

\item Vertical oscillations
\end{itemize}

\section{Initial Conditions}

Usually, we need to know how we start our oscillator to solve a problem.
Let's see how this works.\FRAME{dtbpF}{1.5833in}{1.2187in}{0pt}{}{}{Figure}{%
\special{language "Scientific Word";type "GRAPHIC";maintain-aspect-ratio
TRUE;display "USEDEF";valid_file "T";width 1.5833in;height 1.2187in;depth
0pt;original-width 1.585in;original-height 1.2134in;cropleft "0";croptop
"1";cropright "1";cropbottom "0";tempfilename
'SUS73U0I.wmf';tempfile-properties "XPR";}}

Suppose we start the motion of a mass attached to a spring (a harmonic
oscillator) by pulling the mass to $x=x_{\max }$ and releasing it at $t=0.$
Let's see if we can find the phase. Our initial conditions require 
\begin{eqnarray*}
x\left( 0\right) &=&x_{\max } \\
v\left( 0\right) &=&0
\end{eqnarray*}

Using our equation for $x\left( t\right) $ 
\begin{eqnarray*}
x\left( 0\right) &=&x_{\max }=x_{\max }\cos \left( 0+\phi _{o}\right) \\
x_{\max } &=&x_{\max }\cos \left( \phi _{o}\right) \\
1 &=&\cos \left( \phi _{o}\right)
\end{eqnarray*}%
then 
\begin{equation*}
\phi _{o}=\cos ^{-1}\left( 1\right)
\end{equation*}%
and our calculator gives us an answer of 
\begin{equation*}
\phi _{o}=0
\end{equation*}%
but of course we know that this is not the only place $x\left( t\right)
=x_{\max }.$ It's a cosine graph, so $\phi _{o}$ could be $0$ but it could
also be $2\pi $ or $4\pi $ or $6\pi $ and so forth. \FRAME{dtbpFX}{4.4997in}{%
1.1735in}{0pt}{}{}{Plot}{\special{language "Scientific Word";type
"MAPLEPLOT";width 4.4997in;height 1.1735in;depth 0pt;display
"USEDEF";plot_snapshots TRUE;mustRecompute FALSE;lastEngine "MuPAD";xmin
"-15.001000";xmax "15.001000";xviewmin "-15.001000";xviewmax
"15.001000";yviewmin "-1";yviewmax "1";viewset"XY";rangeset"X";plottype
4;labeloverrides 3;x-label "phi_o";y-label "x";axesFont "Times New
Roman,12,0000000000,useDefault,normal";numpoints 100;plotstyle
"patch";axesstyle "normal";axestips FALSE;xis \TEXUX{x};var1name
\TEXUX{$x$};function \TEXUX{$\cos \left( x\right) $};linecolor
"blue";linestyle 1;pointstyle "point";linethickness 1;lineAttributes
"Solid";var1range "-15.001000,15.001000";num-x-gridlines 100;curveColor
"[flat::RGB:0x000000ff]";curveStyle "Line";function \TEXUX{$\left[
\MATRIX{2,3}{c}\VR{,,c,,,}{,,c,,,}{,,,,,}\HR{,,,}\CELL{-4\pi
}\CELL{1}\CELL{-2\pi }\CELL{1}\CELL{0}\CELL{1}\right] $};linecolor
"blue";linestyle 1;pointplot TRUE;pointstyle "point";linethickness
1;lineAttributes "Solid";curveColor "[flat::RGB:0x000000ff]";curveStyle
"Point";function \TEXUX{$\left[
\MATRIX{2,3}{c}\VR{,,c,,,}{,,c,,,}{,,,,,}\HR{,,,}\CELL{4\pi
}\CELL{1}\CELL{2\pi }\CELL{1}\CELL{0}\CELL{1}\right] $};linecolor
"blue";linestyle 1;pointplot TRUE;pointstyle "point";linethickness
1;lineAttributes "Solid";curveColor "[flat::RGB:0x000000ff]";curveStyle
"Point";VCamFile 'TDT89C11.xvz';valid_file "T";tempfilename
'TDT7Y004.wmf';tempfile-properties "XPR";}}How do we know which $\phi _{o}$
to choose? The answer is that we need another equation. And we have another
equation. We know $v\left( t\right) $ and we know that at $t=0$ the speed $%
v\left( 0\right) =0$ because we have just let go. It hasn't moved yet. We
can use the $v\left( t\right) $ and our initial condition for $v$%
\begin{eqnarray*}
v\left( 0\right) &=&0=-v_{\max }\sin \left( 0+\phi _{o}\right) \\
0 &=&-v_{\max }\sin \left( \phi _{o}\right) \\
0 &=&\sin \left( \phi _{o}\right)
\end{eqnarray*}%
Again our calculator gives $\phi _{o}=0,$ but we know that $\sin \left( \phi
_{o}\right) =0$ at $\pi $ and $2\pi $ and $3\pi $ as well. \FRAME{dtbpFX}{%
4.4997in}{1.1735in}{0pt}{}{}{Plot}{\special{language "Scientific Word";type
"MAPLEPLOT";width 4.4997in;height 1.1735in;depth 0pt;display
"USEDEF";plot_snapshots TRUE;mustRecompute FALSE;lastEngine "MuPAD";xmin
"-15.001000";xmax "15.001000";xviewmin "-15.001000";xviewmax
"15.001000";yviewmin "-1";yviewmax "1";viewset"XY";rangeset"X";plottype
4;labeloverrides 3;x-label "phi_o";y-label "v";axesFont "Times New
Roman,12,0000000000,useDefault,normal";numpoints 100;plotstyle
"patch";axesstyle "normal";axestips FALSE;xis \TEXUX{x};var1name
\TEXUX{$x$};function \TEXUX{$\sin \left( x\right) $};linecolor
"maroon";linestyle 1;pointstyle "point";linethickness 1;lineAttributes
"Solid";var1range "-15.001000,15.001000";num-x-gridlines 100;curveColor
"[flat::RGB:0x00800000]";curveStyle "Line";function \TEXUX{$\left[
\MATRIX{2,5}{c}\VR{,,c,,,}{,,c,,,}{,,,,,}\HR{,,,,,}\CELL{0}\CELL{0}\CELL{-%
\pi }\CELL{0}\CELL{-2\pi }\CELL{0}\CELL{-3\pi }\CELL{0}\CELL{-4\pi
}\CELL{0}\right] $};linecolor "red";linestyle 1;pointplot TRUE;pointstyle
"point";linethickness 1;lineAttributes "Solid";curveColor
"[flat::RGB:0x00ff0000]";curveStyle "Point";function \TEXUX{$\left[
\MATRIX{2,5}{c}\VR{,,c,,,}{,,c,,,}{,,,,,}\HR{,,,,,}\CELL{0}\CELL{0}\CELL{\pi
}\CELL{0}\CELL{2\pi }\CELL{0}\CELL{3\pi }\CELL{0}\CELL{4\pi }\CELL{0}\right]
$};linecolor "red";linestyle 1;pointplot TRUE;pointstyle
"point";linethickness 1;lineAttributes "Solid";curveColor
"[flat::RGB:0x00ff0000]";curveStyle "Point";VCamFile
'TDT7ZE0O.xvz';valid_file "T";tempfilename
'TDT76Y01.wmf';tempfile-properties "XPR";}}But notice that $\phi _{o}=0$
satisfies both our $x\left( 0\right) $ and our $v\left( 0\right) $
equations. So $\phi _{o}=0$ is a good answer.

\FRAME{dtbpFX}{4.4997in}{1.1735in}{0pt}{}{}{Plot}{\special{language
"Scientific Word";type "MAPLEPLOT";width 4.4997in;height 1.1735in;depth
0pt;display "USEDEF";plot_snapshots TRUE;mustRecompute FALSE;lastEngine
"MuPAD";xmin "-15.001000";xmax "15.001000";xviewmin "-15.001000";xviewmax
"15.001000";yviewmin "-1";yviewmax "1";viewset"XY";rangeset"X";plottype
4;labeloverrides 3;x-label "phi_o";y-label "x";axesFont "Times New
Roman,12,0000000000,useDefault,normal";numpoints 100;plotstyle
"patch";axesstyle "normal";axestips FALSE;xis \TEXUX{x};var1name
\TEXUX{$x$};function \TEXUX{$\cos \left( x\right) $};linecolor
"blue";linestyle 1;pointstyle "point";linethickness 1;lineAttributes
"Solid";var1range "-15.001000,15.001000";num-x-gridlines 100;curveColor
"[flat::RGB:0x000000ff]";curveStyle "Line";function \TEXUX{$\left[
\MATRIX{2,3}{c}\VR{,,c,,,}{,,c,,,}{,,,,,}\HR{,,,}\CELL{-4\pi
}\CELL{1}\CELL{-2\pi }\CELL{1}\CELL{0}\CELL{1}\right] $};linecolor
"blue";linestyle 1;pointplot TRUE;pointstyle "point";linethickness
1;lineAttributes "Solid";curveColor "[flat::RGB:0x000000ff]";curveStyle
"Point";function \TEXUX{$\left[
\MATRIX{2,3}{c}\VR{,,c,,,}{,,c,,,}{,,,,,}\HR{,,,}\CELL{4\pi
}\CELL{1}\CELL{2\pi }\CELL{1}\CELL{0}\CELL{1}\right] $};linecolor
"blue";linestyle 1;pointplot TRUE;pointstyle "point";linethickness
1;lineAttributes "Solid";curveColor "[flat::RGB:0x000000ff]";curveStyle
"Point";function \TEXUX{$\sin \left( x\right) $};linecolor "red";linestyle
1;pointstyle "point";linethickness 1;lineAttributes "Solid";var1range
"-15.001000,15.001000";num-x-gridlines 100;curveColor
"[flat::RGB:0x00ff0000]";curveStyle "Line";function \TEXUX{$\left[
\MATRIX{2,5}{c}\VR{,,c,,,}{,,c,,,}{,,,,,}\HR{,,,,,}\CELL{0}\CELL{0}\CELL{-%
\pi }\CELL{0}\CELL{-2\pi }\CELL{0}\CELL{-3\pi }\CELL{0}\CELL{-4\pi
}\CELL{0}\right] $};linecolor "red";linestyle 1;pointplot TRUE;pointstyle
"point";linethickness 1;lineAttributes "Solid";curveColor
"[flat::RGB:0x00ff0000]";curveStyle "Point";function \TEXUX{$\left[
\MATRIX{2,5}{c}\VR{,,c,,,}{,,c,,,}{,,,,,}\HR{,,,,,}\CELL{0}\CELL{0}\CELL{\pi
}\CELL{0}\CELL{2\pi }\CELL{0}\CELL{3\pi }\CELL{0}\CELL{4\pi }\CELL{0}\right]
$};linecolor "red";linestyle 1;pointplot TRUE;pointstyle
"point";linethickness 1;lineAttributes "Solid";curveColor
"[flat::RGB:0x00ff0000]";curveStyle "Point";VCamFile
'TDT8A812.xvz';valid_file "T";tempfilename
'TDT7G102.wmf';tempfile-properties "XPR";}}It's still not the only answer
though. Notice that at $2\pi $ and $4\pi $ we have both 
\begin{eqnarray*}
x_{\max }\cos \left( 0+2\pi \right) &=&x_{\max } \\
-v_{\max }\sin \left( 0+2\pi \right) &=&0
\end{eqnarray*}%
and%
\begin{eqnarray*}
x_{\max }\cos \left( 0+4\pi \right) &=&x_{\max } \\
-v_{\max }\sin \left( 0+4\pi \right) &=&0
\end{eqnarray*}

We have another equation! we know 
\begin{equation*}
a\left( t\right) =-a_{\max }\cos \left( \omega t+\phi _{o}\right)
\end{equation*}%
and as we start out the acceleration has to be maximum. So 
\begin{eqnarray*}
a\left( 0\right) &=&a_{\max }=-a_{\max }\cos \left( 0+\phi _{o}\right) \\
-1 &=&\cos \left( \phi _{o}\right)
\end{eqnarray*}%
\FRAME{dtbpFX}{4.4997in}{1.1735in}{0pt}{}{}{Plot}{\special{language
"Scientific Word";type "MAPLEPLOT";width 4.4997in;height 1.1735in;depth
0pt;display "USEDEF";plot_snapshots TRUE;mustRecompute FALSE;lastEngine
"MuPAD";xmin "-15.001000";xmax "15.001000";xviewmin "-15.001000";xviewmax
"15.001000";yviewmin "-1";yviewmax "1";viewset"XY";rangeset"X";plottype
4;labeloverrides 3;x-label "phi_o";y-label "a";axesFont "Times New
Roman,12,0000000000,useDefault,normal";numpoints 100;plotstyle
"patch";axesstyle "normal";axestips FALSE;xis \TEXUX{x};var1name
\TEXUX{$x$};function \TEXUX{$-\cos \left( x\right) $};linecolor
"green";linestyle 1;pointstyle "point";linethickness 1;lineAttributes
"Solid";var1range "-15.001000,15.001000";num-x-gridlines 100;curveColor
"[flat::RGB:0x00008000]";curveStyle "Line";function \TEXUX{$\left[
\MATRIX{2,3}{c}\VR{,,c,,,}{,,c,,,}{,,,,,}\HR{,,,}\CELL{-4\pi
}\CELL{-1}\CELL{-2\pi }\CELL{-1}\CELL{0}\CELL{-1}\right] $};linecolor
"green";linestyle 1;pointplot TRUE;pointstyle "point";linethickness
1;lineAttributes "Solid";curveColor "[flat::RGB:0x00008000]";curveStyle
"Point";function \TEXUX{$\left[
\MATRIX{2,3}{c}\VR{,,c,,,}{,,c,,,}{,,,,,}\HR{,,,}\CELL{0}\CELL{-1}\CELL{2\pi
}\CELL{-1}\CELL{4\pi }\CELL{-1}\right] $};linecolor "green";linestyle
1;pointplot TRUE;pointstyle "point";linethickness 1;lineAttributes
"Solid";curveColor "[flat::RGB:0x00008000]";curveStyle "Point";VCamFile
'TDT8470V.xvz';valid_file "T";tempfilename
'TDT7WM03.wmf';tempfile-properties "XPR";}}If we plot all of these on top of
one another we can see that we still have a red, blue, and green dot at $%
\phi _{o}=0,$ $2\pi ,$ $4\pi $ and so forth. \FRAME{dtbpFX}{4.4997in}{%
1.1735in}{0pt}{}{}{Plot}{\special{language "Scientific Word";type
"MAPLEPLOT";width 4.4997in;height 1.1735in;depth 0pt;display
"USEDEF";plot_snapshots TRUE;mustRecompute FALSE;lastEngine "MuPAD";xmin
"-15.001000";xmax "15.001000";xviewmin "-15.001000";xviewmax
"15.001000";yviewmin "-1";yviewmax "1";viewset"XY";rangeset"X";plottype
4;labeloverrides 3;x-label "phi_o";y-label "x";axesFont "Times New
Roman,12,0000000000,useDefault,normal";numpoints 100;plotstyle
"patch";axesstyle "normal";axestips FALSE;xis \TEXUX{x};var1name
\TEXUX{$x$};function \TEXUX{$\cos \left( x\right) $};linecolor
"blue";linestyle 1;pointstyle "point";linethickness 1;lineAttributes
"Solid";var1range "-15.001000,15.001000";num-x-gridlines 100;curveColor
"[flat::RGB:0x000000ff]";curveStyle "Line";function \TEXUX{$\left[
\MATRIX{2,3}{c}\VR{,,c,,,}{,,c,,,}{,,,,,}\HR{,,,}\CELL{-4\pi
}\CELL{1}\CELL{-2\pi }\CELL{1}\CELL{0}\CELL{1}\right] $};linecolor
"blue";linestyle 1;pointplot TRUE;pointstyle "point";linethickness
1;lineAttributes "Solid";curveColor "[flat::RGB:0x000000ff]";curveStyle
"Point";function \TEXUX{$\left[
\MATRIX{2,3}{c}\VR{,,c,,,}{,,c,,,}{,,,,,}\HR{,,,}\CELL{4\pi
}\CELL{1}\CELL{2\pi }\CELL{1}\CELL{0}\CELL{1}\right] $};linecolor
"blue";linestyle 1;pointplot TRUE;pointstyle "point";linethickness
1;lineAttributes "Solid";curveColor "[flat::RGB:0x000000ff]";curveStyle
"Point";function \TEXUX{$\sin \left( x\right) $};linecolor "red";linestyle
1;pointstyle "point";linethickness 1;lineAttributes "Solid";var1range
"-15.001000,15.001000";num-x-gridlines 100;curveColor
"[flat::RGB:0x00ff0000]";curveStyle "Line";function \TEXUX{$\left[
\MATRIX{2,5}{c}\VR{,,c,,,}{,,c,,,}{,,,,,}\HR{,,,,,}\CELL{0}\CELL{0}\CELL{-%
\pi }\CELL{0}\CELL{-2\pi }\CELL{0}\CELL{-3\pi }\CELL{0}\CELL{-4\pi
}\CELL{0}\right] $};linecolor "red";linestyle 1;pointplot TRUE;pointstyle
"point";linethickness 1;lineAttributes "Solid";curveColor
"[flat::RGB:0x00ff0000]";curveStyle "Point";function \TEXUX{$\left[
\MATRIX{2,5}{c}\VR{,,c,,,}{,,c,,,}{,,,,,}\HR{,,,,,}\CELL{0}\CELL{0}\CELL{\pi
}\CELL{0}\CELL{2\pi }\CELL{0}\CELL{3\pi }\CELL{0}\CELL{4\pi }\CELL{0}\right]
$};linecolor "red";linestyle 1;pointplot TRUE;pointstyle
"point";linethickness 1;lineAttributes "Solid";curveColor
"[flat::RGB:0x00ff0000]";curveStyle "Point";function \TEXUX{$-\cos \left(
x\right) $};linecolor "green";linestyle 1;pointstyle "point";linethickness
1;lineAttributes "Solid";var1range "-15.001000,15.001000";num-x-gridlines
100;curveColor "[flat::RGB:0x00008000]";curveStyle "Line";function
\TEXUX{$\left[ \MATRIX{2,3}{c}\VR{,,c,,,}{,,c,,,}{,,,,,}\HR{,,,}\CELL{-4\pi
}\CELL{-1}\CELL{-2\pi }\CELL{-1}\CELL{0}\CELL{-1}\right] $};linecolor
"green";linestyle 1;pointplot TRUE;pointstyle "point";linethickness
1;lineAttributes "Solid";curveColor "[flat::RGB:0x00008000]";curveStyle
"Point";function \TEXUX{$\left[
\MATRIX{2,3}{c}\VR{,,c,,,}{,,c,,,}{,,,,,}\HR{,,,}\CELL{0}\CELL{-1}\CELL{2\pi
}\CELL{-1}\CELL{4\pi }\CELL{-1}\right] $};linecolor "green";linestyle
1;pointplot TRUE;pointstyle "point";linethickness 1;lineAttributes
"Solid";curveColor "[flat::RGB:0x00008000]";curveStyle "Point";VCamFile
'TDT8AN13.xvz';valid_file "T";tempfilename
'TDT86805.wmf';tempfile-properties "XPR";}}so we haven't resolved our $\phi
_{o}$ to just one number. We should pause to see why before we come up with
a solution. In the next figure we can start at the bottom. That is the
situation for the problem we are doing. The spring is stretched to $x_{\max
} $ and then we let go. Going clockwise from the bottom we see what will
happen to our mover mass. The spring will pull it to the left. At some
point, it is at the equilibrium position ($x_{e}=0$ in this case). This is
the left most drawing in the figure. But the mass has momentum so it keeps
going. This compresses the spring and increases the spring force until the
mass stops and turns around (top of the figure). Then the spring pushes the
mass to the right. The mass will pass the equilibrium position again and
keep going due to momentum. Eventually we find our mass at $x_{\max }$ again
and momentarily with $v=0$ (bottom of the figure again). \FRAME{dtbpF}{%
4.7649in}{3.5799in}{0in}{}{}{Figure}{\special{language "Scientific
Word";type "GRAPHIC";maintain-aspect-ratio TRUE;display "USEDEF";valid_file
"T";width 4.7649in;height 3.5799in;depth 0in;original-width
4.827in;original-height 3.6189in;cropleft "0";croptop "1";cropright
"1";cropbottom "0";tempfilename 'TDUM3N00.wmf';tempfile-properties "XPR";}}%
But this is exactly the same situation as our starting point! at some time $%
\omega t=2\pi $ later we are right back where we started. We should expect
that the equation of motion at this point would act just like the equation
for $\omega t=0.$ Another cycle later $\left( \omega t=4\pi \right) $ we are
right back at the bottom of the figure again. This looks just like our start
again! We really do have more than one value for $\phi _{o}$ that gives the
same physics.

Let's make an agreement that when this happens we choose to put in our
equation the smallest value of $\phi _{o}$ that works. We could pick any of
the valid values ford $\phi _{o},$ but let's agree to pick the smallest
value that works. In this case, that is our original value $\phi _{o}=0.$

Notice that we needed to know the starting time and the position and the
velocity at that time to find $\phi _{o}$. These are what we call \emph{%
initial conditions}. It is still true that $x\left( t\right) $ and $v\left(
t\right) $ are out of phase. But we found $\phi _{o}=0.$ There is another
phase term hiding in our expression for $v\left( t\right) $ and to find it
we need a small trig identity.%
\begin{eqnarray*}
\sin \left( \alpha +\phi _{o}\right) &=&\cos \left( \alpha -\left( \frac{\pi 
}{2}-\phi _{o}\right) \right) \\
&=&\cos \left( \alpha -\frac{\pi }{2}+\phi _{o}\right)
\end{eqnarray*}%
so we could write 
\begin{eqnarray*}
v\left( t\right) &=&-\omega x_{\max }\sin \left( \omega t+\phi _{o}\right) \\
&=&-\omega x_{\max }\cos \left( \alpha -\frac{\pi }{2}+\phi _{o}\right)
\end{eqnarray*}%
and we can see that there was a phase shift of $-\pi /2$ hiding in our sine
function. So $v\left( t\right) $ really must be out of phase with $x\left(
t\right) .$

\subsection{A second example}

\FRAME{dtbpF}{1.7039in}{1.4422in}{0pt}{}{}{Figure}{\special{language
"Scientific Word";type "GRAPHIC";maintain-aspect-ratio TRUE;display
"USEDEF";valid_file "T";width 1.7039in;height 1.4422in;depth
0pt;original-width 1.7074in;original-height 1.4413in;cropleft "0";croptop
"1";cropright "1";cropbottom "0";tempfilename
'SUS73U0J.wmf';tempfile-properties "XPR";}}

Using the same equipment, let's start with 
\begin{eqnarray*}
x\left( 0\right) &=&0 \\
v\left( 0\right) &=&v_{i}
\end{eqnarray*}%
that is, the mover object is already moving when we start our experiment,
and we start watching just as it passes the equilibrium point.%
\begin{eqnarray*}
x\left( 0\right) &=&0=x_{\max }\cos \left( 0+\phi _{o}\right) \\
v\left( 0\right) &=&v_{i}=-v_{\max }\sin \left( 0+\phi _{o}\right)
\end{eqnarray*}

From the first equation we have 
\begin{equation*}
0=x_{\max }\cos \left( \phi _{o}\right)
\end{equation*}%
and that gives us 
\begin{equation*}
\phi _{o}=\cos ^{-1}\left( 0\right)
\end{equation*}%
so 
\begin{equation*}
\phi _{o}=\pm \frac{\pi }{2}
\end{equation*}%
Really this is 
\begin{equation*}
\phi _{o}=\pm \frac{\pi }{2}\pm n\pi \qquad n=0,1,2,\cdots
\end{equation*}%
This gives the red dots in the next plot. \FRAME{dtbpFX}{4.4996in}{1.1736in}{%
0pt}{}{}{Plot}{\special{language "Scientific Word";type "MAPLEPLOT";width
4.4996in;height 1.1736in;depth 0pt;display "USEDEF";plot_snapshots
TRUE;mustRecompute FALSE;lastEngine "MuPAD";xmin "-15.001000";xmax
"15.001000";xviewmin "-15.001000";xviewmax "15.001000";yviewmin
"-1";yviewmax "1";viewset"XY";rangeset"X";plottype 4;axesFont "Times New
Roman,12,0000000000,useDefault,normal";numpoints 100;plotstyle
"patch";axesstyle "normal";axestips FALSE;xis \TEXUX{x};var1name
\TEXUX{$x$};function \TEXUX{$\cos \left( x\right) $};linecolor
"black";linestyle 1;pointstyle "point";linethickness 1;lineAttributes
"Solid";var1range "-15.001000,15.001000";num-x-gridlines 100;curveColor
"[flat::RGB:0000000000]";curveStyle "Line";function \TEXUX{$\left[
\MATRIX{2,6}{c}\VR{,,c,,,}{,,c,,,}{,,,,,}\HR{,,,,,,}\CELL{\frac{\pi
}{2}+0}\CELL{0}\CELL{\frac{\pi }{2}+1\pi }\CELL{0}\CELL{\frac{\pi }{2}+2\pi
}\CELL{0}\CELL{\frac{\pi }{2}+3\pi }\CELL{0}\CELL{\frac{\pi }{2}+4\pi
}\CELL{0}\CELL{\frac{\pi }{2}+5\pi }\CELL{0}\right] $};linecolor
"blue";linestyle 1;pointplot TRUE;pointstyle "point";linethickness
1;lineAttributes "Solid";curveColor "[flat::RGB:0x000000ff]";curveStyle
"Point";function \TEXUX{$\left[
\MATRIX{2,6}{c}\VR{,,c,,,}{,,c,,,}{,,,,,}\HR{,,,,,,}\CELL{-\frac{\pi
}{2}+0}\CELL{0}\CELL{-\frac{\pi }{2}-1\pi }\CELL{0}\CELL{\frac{\pi }{2}+2\pi
}\CELL{0}\CELL{-\frac{\pi }{2}-3\pi }\CELL{0}\CELL{-\frac{\pi }{2}-4\pi
}\CELL{0}\CELL{-\frac{\pi }{2}-5\pi }\CELL{0}\right] $};linecolor
"blue";linestyle 1;pointplot TRUE;pointstyle "point";linethickness
1;lineAttributes "Solid";curveColor "[flat::RGB:0x000000ff]";curveStyle
"Point";VCamFile 'SUW6VZ08.xvz';valid_file "T";tempfilename
'SUS73U0K.wmf';tempfile-properties "XPR";}}Notice that at each of these
locations $\cos \left( \phi \right) $ is zero. But we made an agreement that
we will choose the smallest value for $\phi _{o}$ that makes $\cos \left(
\phi _{o}\right) =0.$ That is our%
\begin{equation*}
\phi _{o}=\pm \frac{\pi }{2}
\end{equation*}%
But positive and negative $\pi /2$ are the same \textquotedblleft
smallness.\textquotedblright\ We don't know the sign. Using our initial
velocity condition will let us determine the sign. Let's try it 
\begin{eqnarray*}
v_{i} &=&-v_{\max }\sin \left( \pm \frac{\pi }{2}\right) \\
v_{i} &=&\mp v_{\max } \\
v_{i} &=&\mp \omega x_{\max }
\end{eqnarray*}

From this we can see 
\begin{equation*}
x_{\max }=\pm \frac{v_{i}}{\omega }
\end{equation*}

From our initial conditions we know the initial velocity is positive, and we
insist on having positive amplitudes, so then we choose the minus sign so\ 
\begin{equation*}
\sin \left( -\frac{\pi }{2}\right) =\allowbreak -1
\end{equation*}%
and then this minus sign will cancel the one in our initial velocity
equation that to make our initial velocity positive.%
\begin{eqnarray*}
v_{i} &=&-v_{\max }\sin \left( -\frac{\pi }{2}\right) \\
&=&-v_{\max }\left( -1\right) \\
&=&v_{\max }
\end{eqnarray*}%
Our solutions are%
\begin{eqnarray*}
x\left( t\right) &=&\frac{v_{i}}{\omega }\cos \left( \omega t-\frac{\pi }{2}%
\right) \\
v\left( t\right) &=&v_{i}\sin \left( \omega t-\frac{\pi }{2}\right)
\end{eqnarray*}

Note that our solution is a set of equations! It isn't a set of numbers. We
saw this sometimes in Principles of Physics I. But we will see it more in
our study of oscillations and waves. Sometimes the solutions are equations.
But we did fill in our equations with the number for the phase constant. it
would be even better if we had numbers for $\omega $ and \ $v_{i}.$

\begin{center}
{\small \rule{300pt}{1pt}}

{\small Generally to have a complete solution, you must find all the
constants based on the initial conditions. This would mean we need }$x_{\max
},${\small \ }$\omega ,${\small \ and }$\phi _{o}${\small \ to have a
complete solution.}

{\small \rule{300pt}{1pt}}
\end{center}

Let's try a complete example problem.

\subsection{Example}

A particle moving along the $x$ axis in simple harmonic motion starts from
its equilibrium position, the origin, at $t=0$ and moves to the right. The
amplitude of its motion is $2.00\unit{cm},$ and the frequency is $1.50\unit{%
Hz}.$

a) show that the position of the particle is given by%
\begin{equation*}
x=\left( 2.00\unit{cm}\right) \sin \left( 3.00\pi t\right)
\end{equation*}%
determine

b) the maximum speed and the earliest time $(t>0)$ at which the particle has
this speed,

c) the maximum acceleration and the earliest time $(t>0)$ at which the
particle has this acceleration, and

d) the total distance traveled between $t=0$ and $t=1.00\unit{s}$

\FRAME{dtbpF}{1.8775in}{0.8484in}{0pt}{}{}{Figure}{\special{language
"Scientific Word";type "GRAPHIC";maintain-aspect-ratio TRUE;display
"USEDEF";valid_file "T";width 1.8775in;height 0.8484in;depth
0pt;original-width 1.8403in;original-height 0.8164in;cropleft "0";croptop
"1";cropright "1";cropbottom "0";tempfilename
'SUS73U0L.wmf';tempfile-properties "XPR";}}

Basic Equations

\begin{eqnarray*}
x\left( t\right) &=&x_{\max }\cos \left( \omega t+\phi _{o}\right) \\
v\left( t\right) &=&-\omega x_{\max }\sin \left( \omega t+\phi _{o}\right) \\
a\left( t\right) &=&-\omega ^{2}x_{\max }\cos \left( \omega t+\phi
_{o}\right)
\end{eqnarray*}

\begin{equation*}
\omega =2\pi f
\end{equation*}

\begin{equation*}
v_{m}=\omega x_{m}
\end{equation*}

\begin{equation*}
a_{m}=\omega ^{2}x_{m}
\end{equation*}

\begin{equation*}
T=\frac{1}{f}
\end{equation*}

Variables

\begin{equation*}
\begin{tabular}{lll}
$t$ & time, initial time $=0$ & $t_{i}=0$ \\ 
$x$ & Position, Initial position =0 & $x\left( 0\right) =0$ \\ 
$v$ &  &  \\ 
$a$ &  &  \\ 
$x_{\max }$ & x amplitude & $x_{\max }=2.00\unit{cm}$ \\ 
$v_{m}$ & v amplitude &  \\ 
$a_{m}$ & a amplitude &  \\ 
$\omega $ & angular frequency &  \\ 
$\phi _{o}$ & phase constant &  \\ 
$f$ & frequency & $f=1.50\unit{Hz}$%
\end{tabular}%
\end{equation*}

Symbolic Solution

Part (a)

We can start by recognizing that we know $\omega $ because we know the
frequency. 
\begin{eqnarray*}
\omega &=&2\pi f \\
&=&2\pi \left( 1.50\unit{Hz}\right) \\
&=&9.\,\allowbreak 424\,8\unit{rad}\unit{Hz}
\end{eqnarray*}%
We also know the amplitude $x_{\max }$ which is given. 
\begin{equation*}
x_{\max }=2.00\unit{cm}
\end{equation*}%
Knowing that at $t=0$%
\begin{equation*}
x\left( 0\right) =0=x_{\max }\cos \left( 0+\phi _{o}\right)
\end{equation*}%
which we have seen before! We can guess that 
\begin{equation*}
\phi _{o}=\pm \frac{\pi }{2}
\end{equation*}

Using%
\begin{equation*}
v\left( 0\right) =-\omega x_{\max }\sin \left( 0\pm \frac{\pi }{2}\right)
\end{equation*}%
and demanding that amplitudes be positive values, and noting that at $t=0$
the velocity is positive from the initial conditions:%
\begin{equation*}
\phi _{o}=-\frac{\pi }{2}
\end{equation*}%
We also note from our trig identity that we used above%
\begin{equation*}
\cos \left( \theta -\frac{\pi }{2}\right) =\sin \left( \theta \right)
\end{equation*}%
then we have%
\begin{eqnarray*}
x\left( t\right) &=&x_{\max }\cos \left( 2\pi ft-\frac{\pi }{2}\right) \\
&=&x_{\max }\sin \left( 2\pi ft\right)
\end{eqnarray*}%
which is in the form we want. All we have to do is put in numbers. 
\begin{eqnarray*}
x\left( t\right) &=&x_{\max }\sin \left( 2\pi ft\right) \\
&=&\left( 2.00\unit{cm}\right) \sin \left( 3.00\pi t\right)
\end{eqnarray*}%
And once again our solution for a) is an equation.

Part (b)

We have a formula from our previous work for $v_{\max }$ 
\begin{eqnarray*}
v_{\max } &=&\omega x_{\max } \\
&=&2\pi fx_{\max }
\end{eqnarray*}%
so let's use it. 
\begin{eqnarray*}
v_{\max } &=&2\pi \left( 1.50\unit{Hz}\right) \left( 2.00\unit{cm}\right) \\
&=&\allowbreak 6.0\pi \unit{Hz}\unit{cm} \\
&=&\allowbreak 0.188\,50\frac{\unit{m}}{\unit{s}}
\end{eqnarray*}%
To find when $v_{\max }$ happens, take a derivative of $x\left( t\right) $%
\begin{equation*}
v\left( t\right) =v_{\max }=-\omega x_{\max }\sin \left( 2\pi ft-\frac{\pi }{%
2}\right)
\end{equation*}%
and recognize that $\sin \left( \theta \right) =1$ is at a maximum and this
happens when $\theta =\pi /2$. So we take the stuff that is in the
parenthesis for the sine function and set it equal to our angle that makes
the sine function equal to $1$.%
\begin{equation*}
\frac{\pi }{2}=2\pi ft-\frac{\pi }{2}
\end{equation*}%
We can simplify this and solve for $t$%
\begin{equation*}
\pi =2\pi ft
\end{equation*}%
\begin{equation*}
\frac{1}{2f}=t
\end{equation*}%
\begin{equation*}
t=\frac{1}{2\left( 1.50\unit{Hz}\right) }=0.333\,33\unit{s}
\end{equation*}

Part (c) Like with the velocity we must use our equations. We want
acceleration but we know it is just taking the derivative of $v\left(
t\right) $

\begin{equation*}
a\left( t\right) =-\omega ^{2}x_{\max }\cos \left( \omega t+\phi _{o}\right)
\end{equation*}%
And we recognize that the maximum is achieved when $\cos \left( \omega
t+\phi _{o}\right) =1$ or when $\omega t+\phi _{o}=0.$ We can solve this for 
$t.$

\begin{eqnarray*}
t &=&\frac{\phi _{o}}{\omega } \\
&=&\frac{-\frac{\pi }{2}}{2\pi f} \\
&=&\frac{-1}{4f} \\
&=&-0.166\,67\unit{s}
\end{eqnarray*}%
The formula for $a_{\max }$ is also already in our set of equations 
\begin{eqnarray*}
a_{\max } &=&-\omega ^{2}x_{\max } \\
&=&-(2\pi f)^{2}x_{m}
\end{eqnarray*}%
Putting in numbers gives%
\begin{eqnarray*}
a_{\max } &=&\left( 2\pi 1.5\unit{Hz}\right) ^{2}(2.00\unit{cm}) \\
&=&\allowbreak 1.\,\allowbreak 776\,5\frac{\unit{m}}{\unit{s}^{2}}
\end{eqnarray*}

Part (d)

We know the period is 
\begin{equation*}
T=\frac{1}{f}
\end{equation*}%
We should find the number of periods in $t=1.00\unit{s}$ 
\begin{equation*}
N=\frac{t}{T}
\end{equation*}%
and find the distance traveled in one period, and multiply them together. In
one period the distance traveled is 
\begin{equation*}
d=4x_{m}
\end{equation*}

\begin{equation*}
d_{tot}=d\ast N=d\ast \frac{t}{T}=4fx_{m}t
\end{equation*}%
Putting in numbers gives%
\begin{eqnarray*}
d_{tot} &=&4fx_{m}t \\
&=&4\ast 1.50\unit{Hz}\ast 2.00\unit{cm}\ast 1.00\unit{s} \\
&=&12\unit{cm} \\
&=&\allowbreak 0.12\unit{m}
\end{eqnarray*}

We have come far in only two lectures! We have a set of new equations and a
new problem type. In what we did in this lecture we used Newton's Second
Law. But in PH121 we learned that often using the idea of energy made
problems easier. Can we use energy with simple harmonic motion problems?

\section{Oscillators and Energy}

You might have noticed that we are calling mover objects that experience
simple harmonic motion (SHM) by the name \emph{simple harmonic oscillators
(SHO). }Let's consider such a SHO. Because our SHO is moving, we know there
must be energy associated with it. To understand the energy involved, let's
start with kinetic energy. Recall from PH121 that 
\begin{equation}
K=\frac{1}{2}mv^{2}
\end{equation}%
and we recall that for a spring, we have the spring potential energy given by%
\begin{equation}
U_{s}=\frac{1}{2}kx^{2}
\end{equation}%
Let's try a problem.

Using energy techniques, find the maximum velocity of the mass in terms of
the Amplitude and the angular frequency.

We want to use our problem solving steps:

\textbf{Type of problem:} This is an energy and a SHO problem:

\textbf{Drawing:} \FRAME{dtbpF}{1.8775in}{0.8484in}{0pt}{}{}{Figure}{\special%
{language "Scientific Word";type "GRAPHIC";maintain-aspect-ratio
TRUE;display "USEDEF";valid_file "T";width 1.8775in;height 0.8484in;depth
0pt;original-width 1.8403in;original-height 0.8164in;cropleft "0";croptop
"1";cropright "1";cropbottom "0";tempfilename
'SUS73U0T.wmf';tempfile-properties "XPR";}}

Variables: 
\begin{equation*}
\begin{tabular}{ll}
$E$ & energy \\ 
$K$ & Kinetic energy \\ 
$x$ & Position \\ 
$v$ & velocity or speed \\ 
$m$ & mass of the object \\ 
$k$ & spring stiffness constant%
\end{tabular}%
\end{equation*}

\textbf{Basic Equations:}

\begin{eqnarray*}
K &=&\frac{1}{2}mv^{2} \\
U_{s} &=&\frac{1}{2}kx^{2}
\end{eqnarray*}

Symbolic Answer

You might say, this is an easy problem, we know from last time that 
\begin{equation*}
v_{\max }=x_{\max }\omega
\end{equation*}%
but let's find this again using the ideas of energy. In PH121 we often found
that using energy made problems easier, so this might be worth a little more
work now. The total mechanical energy is 
\begin{equation*}
E=K+U_{s}+U_{g}
\end{equation*}%
Let's say that our oscillator is moving horizontally just like last time,
and that we define the center of mass of our object so that it's $y$%
-position is right at $y=0$ so our object has zero gravitational potential
energy. 
\begin{equation*}
E=K+U_{s}+0
\end{equation*}%
And let's say we have no friction, so that we can say that energy is
conserved. We set the initial and final energies equal to each other. 
\begin{eqnarray*}
E_{i} &=&E_{f} \\
K_{i}+U_{si} &=&K_{f}+U_{sf} \\
\frac{1}{2}mv_{i}^{2}+\frac{1}{2}kx_{i}^{2} &=&\frac{1}{2}mv_{f}^{2}+\frac{1%
}{2}kx_{f}^{2}
\end{eqnarray*}%
And let's start with our initial condition being where the spring is
stretched to $x_{\max }.$ Then at that moment, $v_{i}=0.$ And let's take our
final position when the mass is moving as fast as possible (because that is
what we are looking for, $v_{\max }$). \FRAME{dtbpF}{3.8956in}{2.6131in}{0in%
}{}{}{Figure}{\special{language "Scientific Word";type
"GRAPHIC";maintain-aspect-ratio TRUE;display "USEDEF";valid_file "T";width
3.8956in;height 2.6131in;depth 0in;original-width 3.9417in;original-height
2.6343in;cropleft "0";croptop "1";cropright "1";cropbottom "0";tempfilename
'SV31CS01.wmf';tempfile-properties "XPR";}}We know from our PH121 experience
that this will be right when $U_{sf}=0$ (so all the energy is kinetic). Then 
\begin{equation*}
0+\frac{1}{2}kx_{\max }^{2}=\frac{1}{2}mv_{\max }^{2}+0
\end{equation*}%
or%
\begin{equation*}
\frac{1}{2}mv_{\max }^{2}=\frac{1}{2}kx_{\max }^{2}
\end{equation*}%
We can solve for $v_{\max }$ 
\begin{equation*}
v_{\max }^{2}=\frac{kx_{\max }^{2}}{m}
\end{equation*}%
\begin{equation*}
v_{\max }=x_{\max }\sqrt{\frac{k}{m}}
\end{equation*}%
But is this what we wanted? We expected that 
\begin{equation*}
v_{\max }=x_{\max }\omega
\end{equation*}%
And this is what we got so long as 
\begin{equation*}
\omega =\sqrt{\frac{k}{m}}
\end{equation*}%
And this is always true for a mass on a spring. We don't need a numeric
answer, This is reasonable (just what we expect) and the units check.

Let's look at kinetic and potential energy as a function of time. For our
Simple Harmonic Oscillator (SHO) we know the velocity as a function of time,%
\begin{equation*}
v\left( t\right) =-\omega x_{\max }\sin \left( \omega t+\phi \right)
\end{equation*}%
so the kinetic energy as a function of time must be%
\begin{eqnarray*}
K &=&\frac{1}{2}m\left( -\omega x_{\max }\sin \left( \omega t+\phi \right)
\right) ^{2} \\
&=&\frac{1}{2}m\omega ^{2}x_{\max }^{2}\sin ^{2}\left( \omega t+\phi \right)
\end{eqnarray*}%
and now we know that $\omega =\sqrt{k/m},$ so we can write the kinetic
energy as 
\begin{eqnarray*}
K &=& \\
&=&\frac{1}{2}m\left( \sqrt{\frac{k}{m}}\right) ^{2}x_{\max }^{2}\sin
^{2}\left( \omega t+\phi \right) \\
&=&\frac{1}{2}m\frac{k}{m}x_{\max }^{2}\sin ^{2}\left( \omega t+\phi \right)
\\
&=&\frac{1}{2}kx_{\max }^{2}\sin ^{2}\left( \omega t+\phi \right)
\end{eqnarray*}%
As the spring is stretched or compressed we store energy as spring potential
energy. The potential energy due to a spring acting on our SHO (mover mass)
is given by 
\begin{equation}
U_{s}=\frac{1}{2}kx^{2}
\end{equation}%
For our SHO we also know the position as a function of time 
\begin{equation*}
x\left( t\right) =x_{\max }\cos \left( \omega t+\phi _{o}\right)
\end{equation*}%
So, the potential energy as a function of time must be%
\begin{equation*}
U_{s}=\frac{1}{2}kx_{\max }^{2}\cos ^{2}\left( \omega t+\phi \right)
\end{equation*}%
\FRAME{dtbpF}{2.0906in}{0.981in}{0pt}{}{}{Figure}{\special{language
"Scientific Word";type "GRAPHIC";maintain-aspect-ratio TRUE;display
"USEDEF";valid_file "T";width 2.0906in;height 0.981in;depth
0pt;original-width 2.1013in;original-height 0.9704in;cropleft "0";croptop
"1";cropright "1";cropbottom "0";tempfilename
'T8IB2E01.wmf';tempfile-properties "XPR";}}Let's say, again, that our
oscillator is moving horizontally on a frictionless surface, and that we
define the center of mass of our object so that it's $y$-position is right
at $y=0$ so our object has zero gravitational potential energy%
\begin{equation*}
U_{g}=mgy=0
\end{equation*}%
so the total mechanical energy is given by%
\begin{equation*}
E=K+U_{s}
\end{equation*}%
which we can write out as%
\begin{eqnarray*}
E &=&K+U_{s} \\
&=&\frac{1}{2}kx_{\max }^{2}\sin ^{2}\left( \omega t+\phi \right) +\frac{1}{2%
}kx_{\max }^{2}\cos ^{2}\left( \omega t+\phi \right) \\
&=&\frac{1}{2}kx_{\max }^{2}\left( \sin ^{2}\left( \omega t+\phi \right)
+\cos ^{2}\left( \omega t+\phi \right) \right) \\
&=&\frac{1}{2}kx_{\max }^{2}
\end{eqnarray*}%
Where we have used the trig identity that $\sin ^{2}\left( \theta \right)
+\cos ^{2}\left( \theta \right) =1.$

It tells us that if there are no loss mechanisms (e.g. no friction) then the
energy in a harmonic oscillator never changes. And we remember that we would
say that energy is conserved for such a situation, which is not surprising
because we have already used conservation of energy for springs in PH121 and
in the previous example. But let's give a new name to a quantity that does
not change. Let's call it a \emph{constant of motion}. So we can make a
statement about the total energy for our SHO.

\begin{center}
{\small \rule{300pt}{1pt}}

{\small The total mechanical energy of an ideal SHO is a constant of motion}

{\small \rule{300pt}{1pt}}
\end{center}

If we plot the amount of kinetic and potential energy for an oscillator we
might find something like this:\FRAME{dtbpF}{4.7539in}{3.5699in}{0pt}{}{}{%
Figure}{\special{language "Scientific Word";type
"GRAPHIC";maintain-aspect-ratio TRUE;display "USEDEF";valid_file "T";width
4.7539in;height 3.5699in;depth 0pt;original-width 10.0258in;original-height
7.5247in;cropleft "0";croptop "1";cropright "1";cropbottom "0";tempfilename
'SUS73U0V.wmf';tempfile-properties "XPR";}}Note that the kinetic and
potential energy are out of phase with each other. If we plot them on the
same scale ( for the case $\phi _{o}=0$) we have\FRAME{dtbpF}{3.2655in}{%
2.2364in}{0pt}{}{}{Figure}{\special{language "Scientific Word";type
"GRAPHIC";maintain-aspect-ratio TRUE;display "USEDEF";valid_file "T";width
3.2655in;height 2.2364in;depth 0pt;original-width 3.2206in;original-height
2.1958in;cropleft "0";croptop "1";cropright "1";cropbottom "0";tempfilename
'SUS73U0W.wmf';tempfile-properties "XPR";}}

Let's try another problem using energy of a SHO.

\section{Mathematical Representation of Simple Harmonic Motion}

We have a mathematical representation of simple harmonic motion from looking
at our graph of position vs. time. But as a problem, let's use math and what
we know from PH121 to show that our equation 
\begin{equation*}
x\left( t\right) =x_{\max }\cos \left( \omega t\right)
\end{equation*}%
must be right.

Once again recall from Principles of Physics I 
\begin{equation}
a=\frac{dv}{dt}=\frac{d^{2}x}{dt^{2}}
\end{equation}

and let's assume we have SHO moving only in the $x$-direction. \FRAME{dtbpF}{%
2.0649in}{0.9845in}{0in}{}{}{Figure}{\special{language "Scientific
Word";type "GRAPHIC";maintain-aspect-ratio TRUE;display "USEDEF";valid_file
"T";width 2.0649in;height 0.9845in;depth 0in;original-width
2.0755in;original-height 0.9739in;cropleft "0";croptop "1";cropright
"1";cropbottom "0";tempfilename 'T8ICUT03.wmf';tempfile-properties "XPR";}}%
Further assume the surface the object rests on is frictionless. Also let's
say that the equilibrium position $x_{e}=0.$ Then we can write Newton's
second law as%
\begin{eqnarray}
F_{net_{x}} &=&ma_{x}=-kx \\
m\frac{d^{2}x}{dt^{2}} &=&-kx  \notag
\end{eqnarray}%
We have a new kind of equation. If you are taking this freshman class as
a... well... freshman, you may not have seen this kind of equation before.
It is called a differential equation. The solution of this equation is a
function or functions that will describe the motion of our mass-spring
system as a function of time. It says that the way the object moves is an
equation where the second derivative is almost the same as the original
function. The only difference is some constants that are multiplied.

It is this function that we want. Let's rewrite it in terms of $\omega .$
Start by getting all the constants on the same side of the equation. 
\begin{equation*}
\frac{d^{2}x}{dt^{2}}=-\frac{k}{m}x
\end{equation*}%
We found that $\omega =\sqrt{k/m}.$then 
\begin{equation*}
\omega ^{2}=\frac{k}{m}
\end{equation*}%
Then we can write our differential equation as%
\begin{equation}
\frac{d^{2}x}{dt^{2}}=-\omega ^{2}x
\end{equation}%
We need a function who's second derivative is the negative of itself with
just a constant out front. From our calc1 class (M112) we know a few of these%
\begin{eqnarray*}
x\left( t\right) &=&x_{\max }\cos \left( \omega t+\phi _{o}\right) \\
x\left( t\right) &=&x_{\max }\sin \left( \omega t+\phi _{o}\right)
\end{eqnarray*}%
where $x_{\max },$ $\omega ,$ and $\phi _{o}$ are constants that we must
find. Let's choose the cosine function and explicitly take it's derivatives
to see if this function does solve our differential equation%
\begin{eqnarray}
x\left( t\right) &=&x_{\max }\cos \left( \omega t+\phi _{o}\right) \\
\frac{dx\left( t\right) }{dt} &=&-\omega x_{\max }\sin \left( \omega t+\phi
_{o}\right)  \notag \\
\frac{d^{2}x\left( t\right) }{dt^{2}} &=&-\omega ^{2}x_{\max }\cos \left(
\omega t+\phi _{o}\right)  \notag
\end{eqnarray}%
Let's substitute these expressions into our differential equation for the
motion%
\begin{eqnarray*}
\frac{d^{2}x}{dt^{2}} &=&-\omega ^{2}x \\
-\omega ^{2}x_{\max }\cos \left( \omega t+\phi _{o}\right) &=&-\omega
^{2}x_{\max }\cos \left( \omega t+\phi _{o}\right)
\end{eqnarray*}%
As long as the constant $\omega ^{2}$ is our $\omega ^{2}=k/m$ we have a
solution. We could have found that 
\begin{equation*}
\omega ^{2}=\frac{k}{m}
\end{equation*}%
by solving this differential equation, but it might not have been as
meaningful that way. We can identify $\omega $ as the angular frequency.%
\begin{equation*}
\omega =2\pi f
\end{equation*}

Thus%
\begin{equation}
\omega =\sqrt{\frac{k}{m}}=2\pi f
\end{equation}%
This says that if the spring is stiffer, we get a higher frequency, or if
the mass is larger, we get a lower frequency.

We still don't have a complete solution, because we don't know $x_{\max }$
and $\phi _{o}.$ We recognize $\phi _{o}$ as the initial phase angle. We
will have to find this by knowing the initial conditions of the motion. $%
x_{\max }$ is the amplitude. Let's look at a specific case%
\begin{equation}
\begin{tabular}{|l|}
\hline
$x_{\max }=5\unit{m}$ \\ \hline
$\phi _{o}=0$ \\ \hline
$\omega =2\unit{Hz}$ \\ \hline
\end{tabular}%
\end{equation}

\FRAME{dtbpFX}{2.7812in}{1.8542in}{0pt}{}{}{Plot}{\special{language
"Scientific Word";type "MAPLEPLOT";width 2.7812in;height 1.8542in;depth
0pt;display "USEDEF";plot_snapshots TRUE;mustRecompute FALSE;lastEngine
"MuPAD";xmin "-5";xmax "5";xviewmin "-5";xviewmax "5";yviewmin
"-10";yviewmax "10";viewset"XY";rangeset"X";plottype 4;labeloverrides
3;x-label "t";y-label "x";axesFont "Times New
Roman,12,0000000000,useDefault,normal";numpoints 100;plotstyle
"patch";axesstyle "normal";axestips FALSE;xis \TEXUX{t};var1name
\TEXUX{$t$};function \TEXUX{$5\cos \left( 2t+0\right) $};linecolor
"blue";linestyle 1;pointstyle "point";linethickness 1;lineAttributes
"Solid";var1range "-5,5";num-x-gridlines 100;curveColor
"[flat::RGB:0x000000ff]";curveStyle "Line";VCamFile
'SUXKB70N.xvz';valid_file "T";tempfilename
'SUS73U0Y.wmf';tempfile-properties "XPR";}}The amplitude corresponds to the
maximum displacement $x_{\max }.$ We know from trigonometry that a cosine
function has a period $T.$

\FRAME{fhF}{2.9101in}{1.9484in}{0pt}{}{}{Figure}{\special{language
"Scientific Word";type "GRAPHIC";maintain-aspect-ratio TRUE;display
"USEDEF";valid_file "T";width 2.9101in;height 1.9484in;depth
0pt;original-width 4.4616in;original-height 2.9793in;cropleft "0";croptop
"1";cropright "1";cropbottom "0";tempfilename
'SUS73U0Z.wmf';tempfile-properties "XPR";}}The period is related to the
frequency%
\begin{equation}
T=\frac{1}{f}=\frac{2\pi }{\omega }
\end{equation}%
We can write the period and frequency in terms of our mass and spring
constant%
\begin{eqnarray}
T &=&2\pi \sqrt{\frac{m}{k}} \\
f &=&\frac{1}{2\pi }\sqrt{\frac{k}{m}}
\end{eqnarray}

From our graph we were able to complete our problem. But more importantly we
have two new equations for $T$ and $f$ that will help us solve more problems.

\subsection{Hanging Springs}

Let's do a problem. We now know we have forces involved with our SHOs. So
let's do a force problem. In class so far, I have always hung our masses and
springs, but we have used horizontal systems for calculations. Let's find
the equation of motion for a mass hanging from a spring.

\FRAME{dtbpF}{3.429in}{1.9977in}{0in}{}{}{Figure}{\special{language
"Scientific Word";type "GRAPHIC";maintain-aspect-ratio TRUE;display
"USEDEF";valid_file "T";width 3.429in;height 1.9977in;depth
0in;original-width 3.3831in;original-height 1.9579in;cropleft "0";croptop
"1";cropright "1";cropbottom "0";tempfilename
'SUS73U10.wmf';tempfile-properties "XPR";}}

From observation, we would guess that we can just choose the new equilibrium
point to be $y=0$ and use the same equation%
\begin{equation*}
y\left( t\right) =y_{\max }\cos (\omega t-\phi _{o})
\end{equation*}%
let's see if that is true. Start with a free body diagram for the mass (the
hanging mass is our mover, the spring is part of the environment). There are
two forces acting on the mass. The force due to gravity, and the force due
to the spring.\FRAME{dtbpF}{0.683in}{2.1589in}{0pt}{}{}{Figure}{\special%
{language "Scientific Word";type "GRAPHIC";maintain-aspect-ratio
TRUE;display "USEDEF";valid_file "T";width 0.683in;height 2.1589in;depth
0pt;original-width 0.667in;original-height 2.1713in;cropleft "0";croptop
"1";cropright "1";cropbottom "0";tempfilename
'T8IAHP00.wmf';tempfile-properties "XPR";}}From our free body diagram we can
write out a Newton's second law equation%
\begin{equation*}
\Sigma F_{y}=ma_{y}=S_{ms}-W_{mE}
\end{equation*}%
We know the form of these forces%
\begin{eqnarray*}
S_{ms} &=&-k\Delta y \\
W_{mE} &=&-mg
\end{eqnarray*}%
but we need to carefully choose our origin. Let's try the top of the spring
where it attaches to the stand. If the mass just hangs there we would expect
the spring to stretch to an equilibrium length $y_{e}$ and any other motion
would either shorten or lengthen the spring. We can write our force equation
as 
\begin{eqnarray*}
ma_{y} &=&S_{ms}-W_{mE} \\
&=&-k\left( y-y_{e}\right) -mg
\end{eqnarray*}%
where $y$ can be positive or negative. If $y=0$ and the mass is just sitting
there, not oscillating, there is no acceleration. Then the stretched length
is just $y_{e}.$ In this case%
\begin{eqnarray*}
m\left( 0\right) &=&-k\left( 0-y_{e}\right) -mg \\
ky_{e} &=&mg
\end{eqnarray*}%
But now let's let our mass move again. We can substitute our stationary mass
answer into the previous equation for a moving mass%
\begin{eqnarray*}
ma &=&-k\left( y-y_{e}\right) -mg \\
&=&ky_{e}-ky-mg \\
&=&mg-ky-mg \\
&=&-ky
\end{eqnarray*}%
This gives a net force for the system of 
\begin{equation*}
F_{net}=-ky
\end{equation*}%
It is as though the system were horizontal with no gravitational force and
only a spring force. We can see that we are justified in claiming that we
could simply choose the origin at the distance $y_{e}$ from the top of the
spring, and we can use the equation 
\begin{equation*}
y\left( t\right) =y_{\max }\cos (\omega t-\phi _{o})
\end{equation*}%
as our equation of motion.

From what we have done so far, it seems like we totally redefined $\omega $
for this second physics class. But really there is a relationship between
angular velocity $\omega $ and angular frequency $\omega .$ We will explore
this in our next lecture. We will also ask what to do when there is friction.

%TCIMACRO{%
%\TeXButton{Basic Equations}{\vspace{0.5in}\hspace{-1.3in}{\LARGE Basic Equations\vspace{0.25in}}}}%
%BeginExpansion
\vspace{0.5in}\hspace{-1.3in}{\LARGE Basic Equations\vspace{0.25in}}%
%EndExpansion
\begin{eqnarray}
x\left( t\right) &=&x_{\max }\cos \left( \omega t+\phi _{o}\right) \\
\frac{dx\left( t\right) }{dt} &=&-\omega x_{\max }\sin \left( \omega t+\phi
_{o}\right)  \notag \\
\frac{d^{2}x\left( t\right) }{dt^{2}} &=&-\omega ^{2}x_{\max }\cos \left(
\omega t+\phi _{o}\right)  \notag
\end{eqnarray}%
\begin{equation*}
K=\frac{1}{2}kx_{\max }^{2}\sin ^{2}\left( \omega t+\phi \right)
\end{equation*}%
\begin{equation*}
U_{s}=\frac{1}{2}kx_{\max }^{2}\cos ^{2}\left( \omega t+\phi \right)
\end{equation*}

\begin{equation*}
y\left( t\right) =y_{\max }\cos (\omega t-\phi _{o})
\end{equation*}

\chapter{3 More Oscillators, Forces and Friction 1.15.4 1.15.5 1.15.6}

%TCIMACRO{\TeXButton{\chaptermark{Oscillators}}{\chaptermark{Oscillators}}}%
%BeginExpansion
\chaptermark{Oscillators}%
%EndExpansion

You have probably wondered if anyone actually uses mass-spring systems. And
we do. They were more common in the past where springs were used to store
energy ($U_{s\text{ }}$) to run clocks and toys and machines. But simple
harmonic motion can describe other systems as well. You have probably seen
an old fashioned clock with a pendulum. A pendulum almost experiences simple
harmonic motion. Let's see how this works. Then, we should ask the question
of what happens when there really is friction.

%TCIMACRO{%
%\TeXButton{Fundamental Concepts}{\vspace{0.10in}\hspace{-0.37in}{\Large Fundamental Concepts\vspace{0.2in}}}}%
%BeginExpansion
\vspace{0.10in}\hspace{-0.37in}{\Large Fundamental Concepts\vspace{0.2in}}%
%EndExpansion

\begin{itemize}
\item Relationship between Simple Harmonic Motion and circular motion.

\item Pendula

\item Damping

\item Forced Oscillations

\item Resonance
\end{itemize}

\section{Comparing Simple Harmonic Motion with Uniform Circular Motion}

You might have objected to our use of $\omega $ as angular frequency. Didn't 
$\omega $ mean angular speed in Principles of Physics I?

That circular motion and SHM are related should not be a surprise once we
found the solutions to the equations of motion were trig functions. Recall
that the trig functions are defined on a unit circle

\FRAME{dtbpF}{2.2866in}{2.3186in}{0in}{}{}{Figure}{\special{language
"Scientific Word";type "GRAPHIC";maintain-aspect-ratio TRUE;display
"USEDEF";valid_file "T";width 2.2866in;height 2.3186in;depth
0in;original-width 2.2468in;original-height 2.2788in;cropleft "0";croptop
"1";cropright "1";cropbottom "0";tempfilename
'SUXK6802.wmf';tempfile-properties "XPR";}}%
\begin{eqnarray}
\tan \theta &=&\frac{x}{y} \\
\cos \theta &=&\frac{x}{h} \\
\sin \theta &=&\frac{y}{h}
\end{eqnarray}

Let's relate this to our SHM equations of motion.\FRAME{dtbpF}{1.7876in}{%
1.7487in}{0pt}{}{}{Figure}{\special{language "Scientific Word";type
"GRAPHIC";maintain-aspect-ratio TRUE;display "USEDEF";valid_file "T";width
1.7876in;height 1.7487in;depth 0pt;original-width 1.7504in;original-height
1.7115in;cropleft "0";croptop "1";cropright "1";cropbottom "0";tempfilename
'SUXK6803.wmf';tempfile-properties "XPR";}}Look at the projection $x$ of the
point $P$ on the $x$ axis. This is just the $x$-component of the position!
Let's follow this projection as $P$ travels around the circle. We find the
projection ranges from $-x_{\max }$ to $x_{\max }.$ If we watch closely we
find the projection's velocity is zero at the extreme points and is a
maximum in the middle. This projection is given as the cosine of the vector
from the origin to $P.$ It is just taking the $x$-component. This
projection, indeed fits our SHO solution.

Now lets define a projection of $P$ onto the $y$ axis. Again we have SHM,
but this time the projection is a sine function. Because 
\begin{equation}
\cos \left( \theta -\frac{\pi }{2}\right) =\sin \left( \theta \right)
\end{equation}%
we can see that this is just a SHO that is $90\unit{%
%TCIMACRO{\U{b0}}%
%BeginExpansion
{{}^\circ}%
%EndExpansion
}=\frac{\pi }{2}\unit{rad}$ out of phase. It is probably worth recalling
that a projection of one vector on another can be expressed by a dot
product. We could express our length $x$ as 
\begin{equation*}
x=\overrightarrow{\mathbf{r}}\cdot \mathbf{\hat{\imath}}=r\cos \theta
\end{equation*}%
where $r$ is the radius of the circle. \FRAME{dtbpF}{2.3687in}{2.3186in}{0in%
}{}{}{Figure}{\special{language "Scientific Word";type
"GRAPHIC";maintain-aspect-ratio TRUE;display "USEDEF";valid_file "T";width
2.3687in;height 2.3186in;depth 0in;original-width 2.3281in;original-height
2.2788in;cropleft "0";croptop "1";cropright "1";cropbottom "0";tempfilename
'SUXK6804.wmf';tempfile-properties "XPR";}}We can see that when $\theta =0$
we have $x=r$ and this will be the largest $x$ value, so $r=x_{\max }.$

By projecting circular motion onto the $x$-axis 
\begin{equation*}
x\left( t\right) =x_{\max }\cos \left( \theta \right)
\end{equation*}%
But $\theta $ changes in time. We can recall from our PH121 or Dynamics
experience that the angular speed 
\begin{equation*}
\omega =\frac{\Delta \theta }{\Delta t}
\end{equation*}%
or, if we agree to start from $\theta _{i}=0$ and $t_{i}=0,$ 
\begin{equation*}
\omega =\frac{\theta }{t}
\end{equation*}%
so 
\begin{equation*}
\theta =\omega t
\end{equation*}%
then we have 
\begin{equation*}
x\left( t\right) =x_{\max }\cos \left( \omega t\right)
\end{equation*}%
Which is just our equation for SHM. Now we can see why we used
\textquotedblleft $\omega $\textquotedblright\ for both angular speed and
angular frequency. Really they are very related. Both tell us something
about how fast a cyclic event happens. For motion around a circle, one is
how fast the point $P$ goes around the circle, and the other is how often
the projection goes back and forth. It makes sense that these have to be the
same rate.\FRAME{dtbpFU}{2.5694in}{2.5201in}{0pt}{\Qcb{The projection of
circular motion onto the $x$-axis gives simple harmonic motion.}}{}{Figure}{%
\special{language "Scientific Word";type "GRAPHIC";maintain-aspect-ratio
TRUE;display "USEDEF";valid_file "T";width 2.5694in;height 2.5201in;depth
0pt;original-width 2.5287in;original-height 2.4785in;cropleft "0";croptop
"1";cropright "1";cropbottom "0";tempfilename
'SUXK6805.wmf';tempfile-properties "XPR";}}

Let's go back to our projection on the $y$-axis.\FRAME{dtbpF}{2.2753in}{%
2.2269in}{0pt}{}{}{Figure}{\special{language "Scientific Word";type
"GRAPHIC";maintain-aspect-ratio TRUE;display "USEDEF";valid_file "T";width
2.2753in;height 2.2269in;depth 0pt;original-width 2.2355in;original-height
2.1862in;cropleft "0";croptop "1";cropright "1";cropbottom "0";tempfilename
'SUXK6806.wmf';tempfile-properties "XPR";}}We found that we can describe
this projection as 
\begin{equation*}
y\left( t\right) =y_{\max }\sin \left( \omega t\right)
\end{equation*}

We will choose the cosine function, but from our trig experience it should
be clear that these projections are equivalent, just $90\unit{%
%TCIMACRO{\U{b0}}%
%BeginExpansion
{{}^\circ}%
%EndExpansion
}$ out of phase

\begin{equation*}
\cos \left( \theta \pm \frac{\pi }{2}\right) =\mp \sin \theta
\end{equation*}

So%
\begin{equation*}
x\left( t\right) =x_{\max }\cos \left( \omega t+\phi _{o}\right)
\end{equation*}%
could be a sine function if $\phi _{o}=\pm \pi /2.$ In this way we have
incorporated both possibilities into one function.

Note that this is just the function we guessed from our observation!

\begin{center}
{\small \rule{300pt}{1pt}}

{\small We see that uniform circular motion can be thought of as the
combination of two SHOs, with a phase difference of }$\pi /2\unit{rad}.$

{\small The angular velocity is given by}%
\begin{equation}
\omega =\frac{v_{t}}{r}
\end{equation}

{\small \rule{300pt}{1pt}}
\end{center}

\FRAME{dtbpF}{2.5997in}{2.6184in}{0pt}{}{}{Figure}{\special{language
"Scientific Word";type "GRAPHIC";maintain-aspect-ratio TRUE;display
"USEDEF";valid_file "T";width 2.5997in;height 2.6184in;depth
0pt;original-width 2.6201in;original-height 2.6388in;cropleft "0";croptop
"1";cropright "1";cropbottom "0";tempfilename
'SUXK6807.wmf';tempfile-properties "XPR";}}A particle traveling on the $x$%
-axis in SHM will travel from $x_{\max }\ $to $-x_{\max }$ and from $%
-x_{\max }\ $to $x_{\max }$ (one complete period, $T$) while the particle
traveling with $P$ makes one complete revolution. Thus, the angular
frequency $\omega $ of the SHO and the angular velocity of the particle at $%
P $ are the same. The magnitude of the tangential velocity is then%
\begin{equation}
v_{t}=\omega r=\omega x_{\max }
\end{equation}%
and the projection of this velocity onto the $x$-axis is 
\begin{equation}
v_{tx}=-\omega x_{\max }\sin \left( \omega t+\phi _{o}\right)
\end{equation}%
Which is just what we expected from our earlier observation!\FRAME{dtbpF}{%
2.4794in}{2.5079in}{0pt}{}{}{Figure}{\special{language "Scientific
Word";type "GRAPHIC";maintain-aspect-ratio TRUE;display "USEDEF";valid_file
"T";width 2.4794in;height 2.5079in;depth 0pt;original-width
2.4388in;original-height 2.4664in;cropleft "0";croptop "1";cropright
"1";cropbottom "0";tempfilename 'SUXK6808.wmf';tempfile-properties "XPR";}}

The centripetal acceleration of a particle at $P$ is given by 
\begin{equation}
a_{c}=\frac{v_{t}^{2}}{r}=\frac{v_{t}^{2}}{x_{\max }}=\frac{\omega
^{2}x_{\max }^{2}}{x_{\max }}=\omega ^{2}x_{\max }
\end{equation}%
The direction of the acceleration is inward toward the origin. Of course, we
just want the $x$-component of this, so again we make a projection. If we
project this onto the $x$-axis we have%
\begin{equation}
a_{x}=-\omega ^{2}x_{\max }\cos \left( \omega t+\phi \right)
\end{equation}%
Again this is just want we expected from our observation.

Then, we can think of our set of SHO equations 
\begin{eqnarray}
x\left( t\right) &=&x_{\max }\cos \left( \omega t\right) \\
v\left( t\right) &=&-\omega x_{\max }\sin \left( \omega t\right)  \notag \\
a\left( t\right) &=&-\omega ^{2}x_{\max }\cos \left( \omega t\right)  \notag
\end{eqnarray}%
as a projection of circular motion.

\section{Our first problem type: Simple Harmonic Motion}

You have probably thought by now that we have a new problem type. We can
call it the \emph{simple harmonic motion }problem type or SHM. The equations
we have so far for this problem type are 
\begin{eqnarray}
x\left( t\right) &=&x_{\max }\cos \left( \omega t+\phi _{o}\right) \\
v\left( t\right) &=&-\omega x_{\max }\sin \left( \omega t+\phi _{o}\right) 
\notag \\
a\left( t\right) &=&-\omega ^{2}x_{\max }\cos \left( \omega t+\phi
_{o}\right)  \notag
\end{eqnarray}%
\begin{equation*}
\omega =2\pi f
\end{equation*}

\begin{equation*}
v_{m}=\omega x_{m}
\end{equation*}

\begin{equation*}
a_{m}=\omega ^{2}x_{m}
\end{equation*}%
\begin{equation*}
T=\frac{1}{f}
\end{equation*}

\begin{equation*}
U_{s}=\frac{1}{2}kx_{\max }^{2}\cos ^{2}\left( \omega t+\phi \right)
\end{equation*}%
\begin{equation*}
K=\frac{1}{2}m\omega ^{2}x_{\max }^{2}\sin ^{2}\left( \omega t+\phi \right)
\end{equation*}

\begin{eqnarray}
T &=&2\pi \sqrt{\frac{m}{k}} \\
f &=&\frac{1}{2\pi }\sqrt{\frac{k}{m}}
\end{eqnarray}

We have used energy to describe simple harmonic motion, and we have found
our equation of motion using a differential equation for mass-spring
systems. And we have simple harmonic motion in the $y$-direction including
the weight force due to gravity for mass-spring systems. But there are other
systems that exhibit simple harmonic motion and many others that nearly
exhibit simple harmonic motion. Let's study a few of these systems.

\section{The Simple Pendulum}

\FRAME{dtbpF}{1.7305in}{1.6737in}{0pt}{}{}{Figure}{\special{language
"Scientific Word";type "GRAPHIC";maintain-aspect-ratio TRUE;display
"USEDEF";valid_file "T";width 1.7305in;height 1.6737in;depth
0pt;original-width 1.7349in;original-height 1.6773in;cropleft "0";croptop
"1";cropright "1";cropbottom "0";tempfilename
'SUS73U11.wmf';tempfile-properties "XPR";}}A simple pendulum is a mass on a
string. The mass is called a \textquotedblleft bob.\textquotedblright\
Usually we study the motion of the pendulum bob, so let's consider the
pendulum the mover object. A simple pendulum bob exhibits periodic motion,
but not exactly simple harmonic motion. Let's set up our pendulum and
measure an angle $\theta $ starting with $\theta =0$ when the pendulum is
just hanging. Then as the pendulum moves up $\theta $ will increase.

The forces on the bob are gravity $\overrightarrow{\mathbf{W}}$, and $%
\overrightarrow{\mathbf{T}},$ the tension on the string. The tangential
component of $W$ is always directed toward $\theta =0$, the just hanging
equilibrium position. This is a restoring force! Let's call the path the bob
takes \textquotedblleft $s$\textquotedblright\ because it is an arclength.\
Then from Jr. High geometry we recall\footnote{%
Really I did not recall this, I\ have to look it up every time, but $s$ is
called the \emph{arclength}.}%
\begin{equation}
s=L\theta
\end{equation}%
We will use a the cylindrical or $rtz$ coordinate system. The radial axis is
directed along the string. The tangential direction is along the circular
path the bob takes and is always tangent to the path. In this coordinate
system, we can solve for the part of the force directed along the path. This
is the restoring part of the net force. Remember from Newton's second law
the radial and tangential components of the force are 
\begin{eqnarray*}
F_{r} &=&ma_{r} \\
F_{t} &=&ma_{t}
\end{eqnarray*}%
and 
\begin{equation*}
a_{t}=\frac{d^{2}s}{dt^{2}}
\end{equation*}%
Let's write out Newton's second law equations with the sum of the forces
part.%
\begin{eqnarray*}
-ma_{r} &=&-T+W\sin \left( 90\unit{%
%TCIMACRO{\U{b0}}%
%BeginExpansion
{{}^\circ}%
%EndExpansion
}-\theta \right) \\
ma_{t} &=&-W\cos \left( 90\unit{%
%TCIMACRO{\U{b0}}%
%BeginExpansion
{{}^\circ}%
%EndExpansion
}-\theta \right)
\end{eqnarray*}%
The second, tangential equation is the important one. It tells us how the
bob will move. Let's use another trig identity 
\begin{equation*}
\cos \left( \frac{\pi }{2}-\theta \right) =\sin \left( \theta \right)
\end{equation*}%
Then our tangential Newton's second law equation becomes%
\begin{equation*}
ma_{t}=-W\sin \left( \theta \right)
\end{equation*}%
and we can solve for the tangential acceleration, $a_{t}$ 
\begin{eqnarray*}
a_{t} &=&-\frac{W}{m}\sin \left( \theta \right) \\
&=&-g\sin \theta
\end{eqnarray*}

We have two expressions for $a_{t.}$ We can set them equal 
\begin{equation}
\frac{d^{2}s}{dt^{2}}=-g\sin \theta  \label{pendulum1}
\end{equation}%
Remember that $s=L\theta $ so we could write the left hand side as%
\begin{equation*}
\frac{d^{2}s}{dt^{2}}=\frac{d^{2}}{dt^{2}}\left( L\theta \right) =L\frac{%
d^{2}}{dt^{2}}\left( \theta \right)
\end{equation*}%
then equation (\ref{pendulum1}) becomes%
\begin{equation*}
L\frac{d^{2}}{dt^{2}}\left( \theta \right) =-g\sin \theta
\end{equation*}%
or 
\begin{equation*}
\frac{d^{2}\theta }{dt^{2}}=-\frac{g}{L}\sin \theta
\end{equation*}%
This is a differential equation much like our differential equation for a
harmonic oscillator, 
\begin{equation*}
\frac{d^{2}x}{dt^{2}}=-\omega ^{2}x
\end{equation*}%
except it has a sine function in it.

Let's look at our sine function for a minute. Suppose we measure the angle
in radians (because we do!) and we plot the sine function $y=\sin \left(
\theta \right) $ (red line) \FRAME{dtbpF}{4.0083in}{1.9904in}{0in}{}{}{Figure%
}{\special{language "Scientific Word";type "GRAPHIC";maintain-aspect-ratio
TRUE;display "USEDEF";valid_file "T";width 4.0083in;height 1.9904in;depth
0in;original-width 4.0562in;original-height 1.9993in;cropleft "0";croptop
"1";cropright "1";cropbottom "0";tempfilename
'T8KJUR1B.wmf';tempfile-properties "XPR";}}and we add to the graph $y=\theta 
$ (dark blue line). Notice that below $\theta =1$ the two curves look a lot
a like! If we go even smaller, say $\theta <0.7\unit{rad}$ the curves are
essentially the same. Suppose we only let our pendulum swing with small
angles, then 
\begin{equation}
\sin \left( \theta \right) \approx \theta  \label{small angle approximation}
\end{equation}

and 
\begin{equation*}
\frac{d^{2}\theta }{dt^{2}}=-\frac{g}{L}\sin \theta
\end{equation*}%
becomes 
\begin{equation*}
\frac{d^{2}\theta }{dt^{2}}\approx -\frac{g}{L}\theta
\end{equation*}%
This is a differential equation much like our differential equation for a
harmonic oscillator! Recall that for a harmonic oscillator we expect 
\begin{equation*}
\frac{d^{2}x}{dt^{2}}=-\omega ^{2}x
\end{equation*}%
we see that or pendulum equation is a match if $\theta =x$ and if 
\begin{equation}
\omega ^{2}=\frac{g}{L}
\end{equation}%
or 
\begin{equation*}
\omega =\sqrt{\frac{g}{L}}
\end{equation*}%
And since the differential equations look the same they have to have the
same solutions. We have a solution that is an equation only instead of 
\begin{equation*}
x\left( t\right) =x_{\max }\cos \left( \omega t+\phi _{o}\right)
\end{equation*}%
we have an equation for $\theta \left( t\right) $ 
\begin{equation*}
\theta \left( t\right) =\theta _{\max }\cos \left( \omega t+\phi _{o}\right)
\end{equation*}%
Since $\omega $ is different for the pendulum, the frequency and period will
also be different. Instead of being functions of $k$ and $m$ it is now in
terms of $g$ and $L.$%
\begin{equation}
T=\frac{2\pi }{\omega }=2\pi \sqrt{\frac{L}{g}}
\end{equation}

\begin{center}
{\small \rule{300pt}{1pt}}

{\small For a pendulum that oscillates only over small angles, the period
and frequency depend only on }$L${\small \ and }$g${\small !}

{\small \rule{300pt}{1pt}}
\end{center}

The analysis we just did for the pendulum we can do for other simple
harmonic oscillators (or near simple harmonic oscillators). Let's try a few
more.

\subsection{Physical Pendulum}

Usually when we build a pendulum, we assume the string is so small compared
to the bob, that we can ignore it's mass. What if this is not true? Suppose
we build a pendulum by making a large solid object swing from one point. Can
we describe it's motion?

\FRAME{dtbpF}{0.9509in}{0.9509in}{0pt}{}{}{Figure}{\special{language
"Scientific Word";type "GRAPHIC";maintain-aspect-ratio TRUE;display
"USEDEF";valid_file "T";width 0.9509in;height 0.9509in;depth
0pt;original-width 0.9402in;original-height 0.9402in;cropleft "0";croptop
"1";cropright "1";cropbottom "0";tempfilename
'T8KJ7Q14.wmf';tempfile-properties "XPR";}}Let's pull it to the right\FRAME{%
dtbpF}{2.1341in}{2.6281in}{0pt}{}{}{Figure}{\special{language "Scientific
Word";type "GRAPHIC";maintain-aspect-ratio TRUE;display "USEDEF";valid_file
"T";width 2.1341in;height 2.6281in;depth 0pt;original-width
2.1456in;original-height 2.6494in;cropleft "0";croptop "1";cropright
"1";cropbottom "0";tempfilename 'T8KJ9116.wmf';tempfile-properties "XPR";}}%
then we consider that because of $W$ we will have a torque about an axis
through $O.$ Last pendulum we used Newton's Second Law. This time let's use
Newton's second law for rotation.%
\begin{equation*}
\overrightarrow{\tau }=\overrightarrow{r}\times \overrightarrow{F}
\end{equation*}%
In our case this is%
\begin{eqnarray*}
\overrightarrow{\tau } &=&\overrightarrow{d}\times \overrightarrow{W} \\
\tau &=&-Wd\sin \theta \\
&=&-mgd\sin \theta
\end{eqnarray*}%
Remember that an extended body has a moment of inertia, $\mathbb{I}.$
Remember also that angular acceleration is given by

\begin{equation*}
\alpha =\frac{d\omega }{dt}=\frac{d^{2}\theta }{dt^{2}}
\end{equation*}

From the rotational form of Newton's Second Law for rotation.%
\begin{eqnarray*}
\Sigma \overrightarrow{\tau } &=&\mathbb{I}\overrightarrow{\alpha } \\
\tau &=&\mathbb{I}\alpha
\end{eqnarray*}

we can write%
\begin{equation*}
-mgd\sin \theta =\mathbb{I}\frac{d^{2}\theta }{dt^{2}}
\end{equation*}%
Getting all the constants together gives%
\begin{equation*}
\frac{d^{2}\theta }{dt^{2}}=-\frac{mgd}{\mathbb{I}}\sin \left( \theta \right)
\end{equation*}%
and again if we let $\theta $ be small so that $\sin \left( \theta \right)
\approx \theta $%
\begin{equation*}
\frac{d^{2}\theta }{dt^{2}}\approx -\frac{mgd}{\mathbb{I}}\theta
\end{equation*}%
which we can compare to 
\begin{equation*}
\frac{d^{2}x}{dt^{2}}=-\omega ^{2}x
\end{equation*}%
and we can see that we have the same differential equation if 
\begin{equation}
\omega ^{2}=\frac{mgd}{\mathbb{I}}
\end{equation}

In this case%
\begin{equation*}
\theta \left( t\right) =\theta _{\max }\cos \left( \omega t+\phi _{o}\right)
\end{equation*}%
\begin{equation}
T=2\pi \sqrt{\frac{\mathbb{I}}{mgd}}
\end{equation}

\subsection{Torsional Pendulum}

Physics majors and those taking PH220 in the future, you will see a
torsional pendulum. Mechanical engineering majors might use a torsional
pendulum. A torsional pendulum is made by suspending a rigid object from a
wire.

\FRAME{dhF}{1.1597in}{1.7443in}{0pt}{}{}{Figure}{\special{language
"Scientific Word";type "GRAPHIC";maintain-aspect-ratio TRUE;display
"USEDEF";valid_file "T";width 1.1597in;height 1.7443in;depth
0pt;original-width 1.126in;original-height 1.7071in;cropleft "0";croptop
"1";cropright "1";cropbottom "0";tempfilename
'SUS73U14.wmf';tempfile-properties "XPR";}}In this case, a rigid disk. The
object can rotate. But wires don't like to twist. The spring-like molecular
bonds resist the twisting.\FRAME{dhF}{4.5844in}{1.8273in}{0pt}{}{}{Figure}{%
\special{language "Scientific Word";type "GRAPHIC";maintain-aspect-ratio
TRUE;display "USEDEF";valid_file "T";width 4.5844in;height 1.8273in;depth
0pt;original-width 4.5325in;original-height 1.7902in;cropleft "0";croptop
"1";cropright "1";cropbottom "0";tempfilename
'SUS73U15.wmf';tempfile-properties "XPR";}}The twisted wire exerts a
restoring torque on the body that is proportional to the angular position
(sound familiar?).%
\begin{equation*}
\tau =-\kappa \theta
\end{equation*}%
This looks like a greek version of $F=-kx$! Again, let's use Newton's Second
Law for rotation.%
\begin{eqnarray*}
\Sigma \tau &=&\mathbb{I}\alpha \\
&=&\mathbb{I}\frac{d^{2}\theta }{dt^{2}}
\end{eqnarray*}%
so%
\begin{eqnarray*}
\mathbb{I}\frac{d^{2}\theta }{dt^{2}} &=&-\kappa \theta \\
\frac{d^{2}\theta }{dt^{2}} &=&-\frac{\kappa }{\mathbb{I}}\theta
\end{eqnarray*}%
Once again we have our favorite differential equation so long as 
\begin{equation}
\omega ^{2}=\frac{\kappa }{\mathbb{I}}
\end{equation}%
which gives us a solution equation 
\begin{equation*}
\theta \left( t\right) =\theta _{\max }\cos \left( \omega t+\phi _{o}\right)
\end{equation*}%
and makes the period of the oscillation%
\begin{equation}
T=2\pi \sqrt{\frac{\mathbb{I}}{\kappa }}
\end{equation}

So far we have used only idea springs or wires or frictionless pendula with
no air drag. This made the equations easier. But what if there is friction
of some sort? Let's deal with this next.

\section{Damped Oscillations}

Back in Principles of Physics I we studied drag forces which are a
dissipative force like friction. Let's add drag to our mass-spring
oscillator by putting the mass in a beaker of water. But in Principles of
Physics I we only had air drag. this time we will need water drag and the
equation is a little bit different.%
\begin{equation}
\overrightarrow{D}=-b\overrightarrow{v}  \label{Damping force}
\end{equation}%
This force is proportional to the velocity. This a dissipative
(friction-like) force typical of what we find when we move objects through
viscus fluids like water or oil or even honey. This is a drag force, but a
more extreme drag force than air drag. \FRAME{dtbpF}{2.127in}{1.6383in}{0pt}{%
}{}{Figure}{\special{language "Scientific Word";type
"GRAPHIC";maintain-aspect-ratio TRUE;display "USEDEF";valid_file "T";width
2.127in;height 1.6383in;depth 0pt;original-width 2.1385in;original-height
1.6409in;cropleft "0";croptop "1";cropright "1";cropbottom "0";tempfilename
'SUS73U16.wmf';tempfile-properties "XPR";}}We call $b$ the damping
coefficient and it depends on how much viscosity (friction) the fluid can
supply. Now, from Newton's second law,\FRAME{dtbpF}{1.1078in}{1.9239in}{0pt}{%
}{}{Figure}{\special{language "Scientific Word";type
"GRAPHIC";maintain-aspect-ratio TRUE;display "USEDEF";valid_file "T";width
1.1078in;height 1.9239in;depth 0pt;original-width 1.0999in;original-height
1.9318in;cropleft "0";croptop "1";cropright "1";cropbottom "0";tempfilename
'T8KKEE1C.wmf';tempfile-properties "XPR";}}%
\begin{equation*}
ma_{y}=S_{ms}-D_{mf}-W_{mE}
\end{equation*}%
but we can simplify this equation like we did before for hanging masses on
springs. We can find the equilibrium for the mass as it just hangs from the
spring, $y_{e}$ but redoing Newton's second law with out the fluid and with
no oscillation 
\begin{eqnarray*}
ma_{y} &=&S_{ms}-W_{mE} \\
&=&-k\left( y-y_{e}\right) -mg
\end{eqnarray*}%
and while the mass just sits there the acceleration would be zero so $%
a_{y}=0 $ and let's define our coordinate system so where the mass hangs is $%
y=0$%
\begin{equation*}
m\left( 0\right) =-k\left( 0-y_{e}\right) -mg
\end{equation*}%
then 
\begin{equation*}
ky_{e}=mg
\end{equation*}%
and we can put this back in our Newton's second law for just hanging and no
fluid%
\begin{eqnarray*}
ma_{y} &=&-k\left( y-y_{e}\right) -mg \\
&=&-ky-ky_{e}-mg \\
&=&-ky-mg-mg \\
&=&-ky
\end{eqnarray*}%
So for just hanging, we can say 
\begin{equation*}
ma_{y}=-ky
\end{equation*}%
which includes the effect of gravity. Now let's add back in our fluid so we
have a damping force. We can write the Newton's second law with damping
equation as%
\begin{equation*}
ma_{y}=-ky-bv
\end{equation*}%
We can write the acceleration and velocity as derivatives of the position
just like we have done before 
\begin{equation*}
m\frac{d^{2}y}{dt^{2}}=-ky-b\frac{dy}{dt}
\end{equation*}%
This is another differential equation! But it is harder to guess its
solution, and finding that solution is a subject for a junior level
differential equations class (like M316 or PH332), so we won't learn how to
find the solution here, but we can use the results from our trusted
colleagues in the math department\footnote{%
That is, until you finish M316, then you will know how to do this kind of
problem yourself!}. Here is the solution:%
\begin{equation}
y\left( t\right) =y_{\max }e^{-\frac{b}{2m}t}\cos \left( \omega t+\phi
_{o}\right)  \label{dampped solution}
\end{equation}%
which looks simple enough, but now we $\ $have the added complication that $%
\omega $ is more complex 
\begin{equation}
\omega =\sqrt{\frac{k}{m}-\left( \frac{b}{2m}\right) ^{2}}
\end{equation}%
so we get quite a mess if we write equation (\ref{dampped solution}) with
this new $\omega .$%
\begin{equation}
y\left( t\right) =y_{\max }e^{-\frac{b}{2m}t}\cos \left( \left( \sqrt{\frac{k%
}{m}-\left( \frac{b}{2m}\right) ^{2}}\right) t+\phi _{o}\right)
\end{equation}

To see what this solution means, we should study three cases:

\begin{enumerate}
\item the damping force is small: ($bv_{\max }<ky_{\max })$ The system
oscillates, but the amplitude is smaller as as time goes on. We call this
\textquotedblleft underdamped.\textquotedblright

\item the damping force is large: ($bv_{\max }>ky_{\max })$The system does
not oscillate. we call this \textquotedblleft overdamped.\textquotedblright\
We can also say that $\frac{b}{2m}>\omega _{o}$ (after we define $\omega
_{o} $ below)

\item The system is \textquotedblleft critically damped\textquotedblright\
(see below).
\end{enumerate}

Let's look at an example. Suppose we have an oscillator with the following
characteristics:

\begin{enumerate}
\item 
\begin{equation*}
\begin{tabular}{|l|}
\hline
$y_{\max }=5\unit{cm}$ \\ \hline
$b=0.005\frac{\unit{kg}}{\unit{s}}$ \\ \hline
$k=0.2\frac{\unit{N}}{\unit{m}}$ \\ \hline
$m=.5\unit{kg}$ \\ \hline
$\phi _{o}=0$ \\ \hline
\end{tabular}%
\end{equation*}

Graphing the equation of motion $y\left( t\right) ,$ we we get a graph that
looks like this\FRAME{dtbpFX}{4.4997in}{3.0007in}{0pt}{}{}{Plot}{\special%
{language "Scientific Word";type "MAPLEPLOT";width 4.4997in;height
3.0007in;depth 0pt;display "USEDEF";plot_snapshots TRUE;mustRecompute
FALSE;lastEngine "MuPAD";xmin "0";xmax "500";xviewmin "0";xviewmax
"500";yviewmin "-0.050771";yviewmax "0.05";viewset"XY";rangeset"X";plottype
4;labeloverrides 3;x-label "t(sec)";y-label "y(m)";axesFont "Times New
Roman,12,0000000000,useDefault,normal";numpoints 100;plotstyle
"patch";axesstyle "normal";axestips FALSE;xis \TEXUX{t};var1name
\TEXUX{$t$};function \TEXUX{$\left( 0.05\right) e^{-\frac{0.005}{2\left(
0.5\right) }t}$};linecolor "gray";linestyle 1;pointstyle
"point";linethickness 1;lineAttributes "Solid";var1range
"0,500";num-x-gridlines 100;curveColor "[flat::RGB:0x00c0c0c0]";curveStyle
"Line";rangeset"X";function \TEXUX{$-\left( 0.05\right)
e^{-\frac{0.005}{2\left( 0.5\right) }t}$};linecolor "yellow";linestyle
1;pointstyle "point";linethickness 1;lineAttributes "Solid";var1range
"0,500";num-x-gridlines 100;curveColor "[flat::RGB:0x00808000]";curveStyle
"Line";function \TEXUX{$\left( 0.05\right) e^{-\frac{0.005}{2\left(
0.5\right) }t}\cos \left( \left( \left( \frac{0.2}{0.5}-\left(
\frac{0.005}{2\left( 0.05\right) }\right) ^{2}\right) ^{\frac{1}{2}}\right)
t\right) $};linecolor "blue";linestyle 1;pointstyle "point";linethickness
1;lineAttributes "Solid";var1range "0,500";num-x-gridlines 900;curveColor
"[flat::RGB:0x000000ff]";curveStyle "Line";rangeset"X";VCamFile
'TJXPJY3Y.xvz';valid_file "T";tempfilename
'TJXPJN0Q.wmf';tempfile-properties "XPR";}}The gray lines are given by%
\begin{equation*}
\pm y_{\max }e^{-\frac{b}{2m}t}
\end{equation*}%
Notice that this quantity is as a collection of terms multiplying the cosine
part. It is marked in curly braces blow 
\begin{equation*}
y\left( t\right) =\left\{ y_{\max }e^{-\frac{b}{2m}t}\right\} \cos \left(
\left( \sqrt{\frac{k}{m}-\left( \frac{b}{2m}\right) ^{2}}\right) t+\phi
_{o}\right) 
\end{equation*}%
We know that the part of the equation that multiplies the cosine part is the
amplitude. But now the amplitude is more than just $x_{\max }.$ And notice
that the amplitude $\left\{ y_{\max }e^{-\frac{b}{2m}t}\right\} $ changes
with time. The gray lines in the figure show how the amplitude changes. We
call this the \emph{envelope} of the curve. The oscillation fits within the
gray lines (like an old fashioned letter fits inside a paper envelope).
\end{enumerate}

Now let's change $b$ to a larger value%
\begin{equation*}
\begin{tabular}{|l|}
\hline
$y_{\max }=5\unit{cm}$ \\ \hline
$b=0.05\frac{\unit{kg}}{\unit{s}}$ \\ \hline
$k=0.2\frac{\unit{N}}{\unit{m}}$ \\ \hline
$m=.5\unit{kg}$ \\ \hline
$\phi _{o}=0$ \\ \hline
\end{tabular}%
\end{equation*}%
\FRAME{dtbpFX}{4.4996in}{3in}{0pt}{}{}{Plot}{\special{language "Scientific
Word";type "MAPLEPLOT";width 4.4996in;height 3in;depth 0pt;display
"USEDEF";plot_snapshots TRUE;mustRecompute FALSE;lastEngine "MuPAD";xmin
"0";xmax "500";xviewmin "0";xviewmax "500";yviewmin "-0.050771";yviewmax
"0.05";viewset"XY";rangeset"X";plottype 4;labeloverrides 1;x-label
"t(sec)";axesFont "Times New
Roman,12,0000000000,useDefault,normal";numpoints 100;plotstyle
"patch";axesstyle "normal";axestips FALSE;xis \TEXUX{t};var1name
\TEXUX{$t$};function \TEXUX{$\left( 0.05\right) e^{-\frac{0.05}{2\left(
0.5\right) }t}$};linecolor "gray";linestyle 1;pointstyle
"point";linethickness 1;lineAttributes "Solid";var1range
"0,500";num-x-gridlines 100;curveColor "[flat::RGB:0x00c0c0c0]";curveStyle
"Line";rangeset"X";function \TEXUX{$-\left( 0.05\right)
e^{-\frac{0.05}{2\left( 0.5\right) }t}$};linecolor "yellow";linestyle
1;pointstyle "point";linethickness 1;lineAttributes "Solid";var1range
"0,500";num-x-gridlines 100;curveColor "[flat::RGB:0x00808000]";curveStyle
"Line";function \TEXUX{$\left( 0.05\right) e^{-\frac{0.05}{2\left(
0.5\right) }t}\cos \left( \left( \left( \frac{0.2}{0.5}-\left(
\frac{0.05}{2\left( 0.05\right) }\right) ^{2}\right) ^{\frac{1}{2}}\right)
t\right) $};linecolor "blue";linestyle 1;pointstyle "point";linethickness
1;lineAttributes "Solid";var1range "0,500";num-x-gridlines 900;curveColor
"[flat::RGB:0x000000ff]";curveStyle "Line";rangeset"X";VCamFile
'SUXMGF0T.xvz';valid_file "T";tempfilename
'SUS73U18.wmf';tempfile-properties "XPR";}}

we see we have less oscillation. The envelope has become more restrictive,
making the oscillation die out more quickly. This is a little bit like going
over a bump in your car. The car may go up and down a few times, but not
many. Now let's increase $b$ even more.%
\begin{equation}
\begin{tabular}{|l|}
\hline
$y_{\max }=5\unit{cm}$ \\ \hline
$b=0.5\frac{\unit{kg}}{\unit{s}}$ \\ \hline
$k=0.2\frac{\unit{N}}{\unit{m}}$ \\ \hline
$m=.5\unit{kg}$ \\ \hline
$\phi _{o}=0$ \\ \hline
\end{tabular}%
\end{equation}%
\FRAME{dtbpFX}{4.4996in}{3in}{0pt}{}{}{Plot}{\special{language "Scientific
Word";type "MAPLEPLOT";width 4.4996in;height 3in;depth 0pt;display
"USEDEF";plot_snapshots TRUE;mustRecompute FALSE;lastEngine "MuPAD";xmin
"0";xmax "50";xviewmin "0";xviewmax "50";yviewmin "-0.050005";yviewmax
"0.05";viewset"XY";rangeset"X";plottype 4;labeloverrides 3;x-label
"t(sec)";y-label "y(m)";axesFont "Times New
Roman,12,0000000000,useDefault,normal";numpoints 100;plotstyle
"patch";axesstyle "normal";axestips FALSE;xis \TEXUX{t};var1name
\TEXUX{$t$};function \TEXUX{$\left( 0.05\right) e^{-\frac{1}{2\left(
0.5\right) }t}$};linecolor "gray";linestyle 1;pointstyle
"point";linethickness 1;lineAttributes "Solid";var1range
"0,50";num-x-gridlines 100;curveColor "[flat::RGB:0x00c0c0c0]";curveStyle
"Line";rangeset"X";function \TEXUX{$-\left( 0.05\right) e^{-\frac{1}{2\left(
0.5\right) }t}$};linecolor "yellow";linestyle 1;pointstyle
"point";linethickness 1;lineAttributes "Solid";var1range
"0,50";num-x-gridlines 100;curveColor "[flat::RGB:0x00808000]";curveStyle
"Line";VCamFile 'SUXMGU0U.xvz';valid_file "T";tempfilename
'SUS73U19.wmf';tempfile-properties "XPR";}}What happened?

When the damping force gets bigger, the oscillation eventually stops. Only
the exponential decay is observed. This happens when%
\begin{equation*}
\frac{b}{2m}=\sqrt{\frac{k}{m}}
\end{equation*}%
When that is true, 
\begin{equation*}
\omega =\sqrt{\frac{k}{m}-\left( \frac{b}{2m}\right) ^{2}}=0
\end{equation*}%
We call this situation \emph{critically damped}. We are just on the edge of
oscillation. We define%
\begin{equation*}
\omega _{o}=\omega _{N}=\sqrt{\frac{k}{m}}
\end{equation*}%
as the \emph{natural frequency} of the system. Then the value of $b$ that
gives us critically damped behavior is%
\begin{equation*}
b_{c}=2m\omega _{N}
\end{equation*}%
When $\frac{b}{2\pi }\geq \omega _{o}$ the solution in equation (\ref%
{dampped solution}) is not valid! If you are a physicist or a mechanical
engineer you will find out more about this situation in your advanced
mechanics classes.

\section{Forced Oscillations}

%TCIMACRO{%
%\TeXButton{\chaptermark{Forces and Waves}}{\chaptermark{Forces and Waves}}}%
%BeginExpansion
\chaptermark{Forces and Waves}%
%EndExpansion

We found in the last section that if we added a force like

\begin{equation*}
\mathbf{D}=-b\mathbf{v}
\end{equation*}%
our oscillation died out. An example would be a small child on a swing. You
give them a push, but eventually they stop swinging.

Suppose we want to keep the child going? You know what to do, you push! But
you have to push at just the right time. Suppose we make a machine that can
push the child on the swing. We could apply a periodic force like%
\begin{equation*}
F\left( t\right) =F_{\max }\sin \left( \omega _{f}t\right)
\end{equation*}%
where $\omega _{f}$ is the angular frequency of this new \emph{driving force}
and where $F_{\max }$ is a constant that tells us the maximum force.%
\begin{equation*}
\Sigma F=F_{\max }\sin \left( \omega _{f}t\right) -kx-bv_{x}=ma
\end{equation*}

When this system starts out, the solution is very messy. It is so messy that
we will not give it in this class!\footnote{%
So maybe this isn't really a good way to drive a child on a swing, be
careful to not knock the child out of the swing as the system starts up!}
But after a while, the system calms down and a steady-state is reached. In
this state, the energy added by our driving force $F_{\max }\sin \left(
\omega _{f}t\right) $ is equal to the energy lost by the drag force, and we
have and equation of motion for our forced oscillator that is just 
\begin{equation*}
x\left( t\right) =A\cos \left( \omega _{f}t+\phi \right)
\end{equation*}%
our old friend! BUT NOW the amplitude is given by 
\begin{equation*}
A=\frac{\frac{F_{\max }}{m}}{\sqrt{\left( \omega _{f}^{2}-\omega
_{N}^{2}\right) ^{2}+\left( \frac{b\omega _{f}}{m}\right) ^{2}}}
\end{equation*}%
where%
\begin{equation}
\omega _{N}=\sqrt{\frac{k}{m}}
\end{equation}%
is the natural frequency. This is the frequency we had before when we
weren't pushing the oscillator. And recall that $\omega _{f}$ is the
frequency of our driving force. So now our solution looks more like our
original SHM solution except for the wild formula for the amplitude. In
fact, it operates very like our SHM solution. But what does the new version
of the amplitude mean?

Lets look at $A$ for some values of $\omega _{f}.$ I will pick some nice
numbers for the other values.%
\begin{equation*}
\begin{tabular}{|l|}
\hline
$F_{\max }=2\unit{N}$ \\ \hline
$b=0.5\frac{\unit{kg}}{\unit{s}}$ \\ \hline
$k=0.5\frac{\unit{N}}{\unit{m}}$ \\ \hline
$m=0.5\unit{kg}$ \\ \hline
$\phi _{o}=0$ \\ \hline
\end{tabular}%
\end{equation*}%
The graph looks like this:\FRAME{dtbpF}{3.6428in}{2.5252in}{0pt}{}{}{Figure}{%
\special{language "Scientific Word";type "GRAPHIC";maintain-aspect-ratio
TRUE;display "USEDEF";valid_file "T";width 3.6428in;height 2.5252in;depth
0pt;original-width 3.6845in;original-height 2.5439in;cropleft "0";croptop
"1";cropright "1";cropbottom "0";tempfilename
'T8LYFW02.wmf';tempfile-properties "XPR";}}Now let's calculate $\omega _{N}$ 
\begin{eqnarray*}
\omega _{N} &=&\sqrt{\frac{0.5\frac{\unit{N}}{\unit{m}}}{0.5\unit{kg}}} \\
&=&\allowbreak \frac{1.0}{\unit{s}}
\end{eqnarray*}%
Notice that right at $\omega _{f}=\omega _{N}$ our solution gets very big.
This is called \emph{resonance}. To see why this happens, think of the
velocity for a simple harmonic oscillator%
\begin{equation*}
\frac{dx\left( t\right) }{dt}=-\omega A\sin \left( \omega t+\phi _{o}\right)
\end{equation*}%
Our driving force is still%
\begin{equation*}
F\left( t\right) =F_{\max }\sin \left( \omega t\right)
\end{equation*}%
And remember that work is given by%
\begin{equation*}
w=\int \overrightarrow{\mathbf{F}}\cdot d\overrightarrow{\mathbf{x}}
\end{equation*}%
or just 
\begin{equation*}
w=\overrightarrow{\mathbf{F}}\cdot \Delta \overrightarrow{\mathbf{x}}
\end{equation*}%
if our force is constant. The rate at which work is done (power) is 
\begin{eqnarray}
\mathcal{P} &=&\frac{\overrightarrow{\mathbf{F}}\cdot \Delta \overrightarrow{%
\mathbf{x}}}{\Delta t}=\overrightarrow{\mathbf{F}}\cdot \overrightarrow{%
\mathbf{v}} \\
&=&-F_{\max }\omega A\sin \left( \omega t\right) \sin \left( \omega t+\phi
_{o}\right)
\end{eqnarray}%
if $F$ and $v$ are in phase $(\phi _{o}=0)$, the power will be at a maximum!
Think of pushing a child on a swing. You push in phase with the oscillation
of the child. When you do this the amplitude of the swing get's bigger.

We can plot $A$ for several values of $b,$ our damping (friction)
coefficient.\FRAME{dtbpFU}{3.9763in}{2.7328in}{0pt}{\Qcb{{\protect\small %
Green: b=0.005kg/s; Blue: b=0.05kg/s; Red b=0.01 kg/s}}}{}{Figure}{\special%
{language "Scientific Word";type "GRAPHIC";maintain-aspect-ratio
TRUE;display "USEDEF";valid_file "T";width 3.9763in;height 2.7328in;depth
0pt;original-width 4.0242in;original-height 2.7559in;cropleft "0";croptop
"1";cropright "1";cropbottom "0";tempfilename
'T8LYK304.wmf';tempfile-properties "XPR";}}As $b\rightarrow 0$ we see that
our resonance peak gets larger. In real systems $b$ can never be zero, but
sometimes it can get small. As $b$ gets large, the resonance dies down and
our $A$ gets small.

Resonance can be great, it can make a musical instrument sound louder (more
on this later). But it can also be bad. Here is a picture of the Tacoma
Narrows Bridge.

\FRAME{dtbpFU}{2.8614in}{2.1722in}{0pt}{\Qcb{{\protect\small Fall of the
Tacoma Narrows Bridge (Image in the Public Domain)}}}{}{Figure}{\special%
{language "Scientific Word";type "GRAPHIC";maintain-aspect-ratio
TRUE;display "USEDEF";valid_file "T";width 2.8614in;height 2.1722in;depth
0pt;original-width 2.8871in;original-height 2.1846in;cropleft "0";croptop
"1";cropright "1";cropbottom "0";tempfilename
'Tacoma_Narrows.wmf';tempfile-properties "XNPR";}}The wind formed a periodic
driving force that allowed a driving harmonic oscillation to form. Since the
bridge was resonant with the wind frequency, the amplitude grew until the
bridge materials broke. Resonance can be a bad thing for structures.

We have learned a lot about oscillation. Now let's turn our attention to
something that results from oscillation, waves.

%TCIMACRO{%
%\TeXButton{Basic Equations}{\vspace{0.5in}\hspace{-1.3in}{\LARGE Basic Equations\vspace{0.25in}}}}%
%BeginExpansion
\vspace{0.5in}\hspace{-1.3in}{\LARGE Basic Equations\vspace{0.25in}}%
%EndExpansion

Pendulum%
\begin{equation*}
\omega =\sqrt{\frac{g}{L}}
\end{equation*}%
\begin{equation*}
\theta \left( t\right) =\theta _{\max }\cos \left( \omega t+\phi _{o}\right)
\end{equation*}%
\begin{equation}
T=\frac{2\pi }{\omega }=2\pi \sqrt{\frac{L}{g}}
\end{equation}

Physical Pendulum

\begin{equation}
\omega =\sqrt{\frac{mgd}{\mathbb{I}}}
\end{equation}

\begin{equation*}
\theta \left( t\right) =\theta _{\max }\cos \left( \omega t+\phi _{o}\right)
\end{equation*}%
\begin{equation}
T=2\pi \sqrt{\frac{\mathbb{I}}{mgd}}
\end{equation}

Torsional Pendulum%
\begin{equation}
\omega =\sqrt{\frac{\kappa }{\mathbb{I}}}
\end{equation}%
\begin{equation*}
\theta \left( t\right) =\theta _{\max }\cos \left( \omega t+\phi _{o}\right)
\end{equation*}%
\begin{equation}
T=2\pi \sqrt{\frac{\mathbb{I}}{\kappa }}
\end{equation}

Damped Oscillation%
\begin{equation*}
y\left( t\right) =\left\{ y_{\max }e^{-\frac{b}{2m}t}\right\} \cos \left(
\left( \sqrt{\frac{k}{m}-\left( \frac{b}{2m}\right) ^{2}}\right) t+\phi
_{o}\right)
\end{equation*}

\begin{equation*}
\omega =\sqrt{\frac{k}{m}-\left( \frac{b}{2m}\right) ^{2}}=0
\end{equation*}%
\begin{equation*}
\omega _{N}=\sqrt{\frac{k}{m}}
\end{equation*}%
\begin{equation*}
b_{c}=2m\omega _{N}
\end{equation*}

Forced Oscillation%
\begin{equation*}
x\left( t\right) =A\cos \left( \omega _{f}t+\phi \right) 
\end{equation*}%
\begin{equation*}
A=\frac{\frac{F_{o}}{m}}{\sqrt{\left( \omega _{f}^{2}-\omega _{N}^{2}\right)
^{2}+\left( \frac{b\omega _{f}}{m}\right) ^{2}}}
\end{equation*}%
\begin{equation}
\omega _{N}=\sqrt{\frac{k}{m}}
\end{equation}%
\begin{equation}
\mathcal{P}=-F_{o}\omega A\sin \left( \omega t\right) \sin \left( \omega
t+\phi _{o}\right) 
\end{equation}

\chapter{4 Sinusoidal Waves in one and Two Dimensions 1.16.1, 1.16.2 1,16.3}

%TCIMACRO{%
%\TeXButton{\chaptermark{Sinusoidal Waves}}{\chaptermark{Sinusoidal Waves}}}%
%BeginExpansion
\chaptermark{Sinusoidal Waves}%
%EndExpansion

We now know what a wave is. And we have a way to categorize types of waves.
But we need a mathematical description of waves so we can calculate
predictions of how waves will behave. We will start building a mathematical
model for waves in this lecture.

%TCIMACRO{%
%\TeXButton{Fundamental Concepts}{\vspace{0.10in}\hspace{-0.37in}{\Large Fundamental Concepts\vspace{0.2in}}}}%
%BeginExpansion
\vspace{0.10in}\hspace{-0.37in}{\Large Fundamental Concepts\vspace{0.2in}}%
%EndExpansion

\begin{itemize}
\item A wave requires a disturbance, and a medium that can transfer energy

\item Waves are categorized as longitudinal or transverse (or a combination
of the two).

\item We use snapshot and history graphs to show the $x$ and $t$ dependence
of waves.

\item It takes two variables to describe a wave, position and time

\item Wave graphs have many named parts, amplitude, period, wavelength
crest, trough, etc.

\item There is a spatial analog to temporal frequency called spacial
frequency and it is represented by the factor $k$ in our equations

\item Waves from point sources are spherical
\end{itemize}

\section{What is a Wave?}

Waves are organized motions in a group of objects. We will give a name to
the group of objects, we will call a \emph{medium}.

Waves involve energy transfer, but in the case of waves the energy is
transferred through space without transfer of matter. Winds transfers both
energy and matter. So waves are different than wind. In a wave the objects
don't move far from where they start. Think of an oscillation. The mass does
move, but never gets too far away from its equilibrium position. Waves are
very like this. This is a very specific kind of motion. To be a wave, the
motion must have the following characteristics"

\begin{enumerate}
\item some source of disturbance

\item a medium (group of objects) that can be disturbed

\item some physical mechanism by which the objects of the medium can
influence each other
\end{enumerate}

A wave made by a rock thrown into a pond will go out in all directions away
from the place where the rock started the wave. This is a normal way that
waves are formed. A disturbance starts the wave (the rock disturbs the
water) and the energy from the disturbance moves away from the disturbance
as a wave. But if we have a wave in a rope or string, the wave can't go in
all directions because the string does not go in all directions.

Let's take on this one-dimensional case of a wave on a rope or string first.
In the limit that the string mass is negligible we represent a
one-dimensional wave mathematically as a function of two variables, position
and time, $y\left( x,t\right) .$

We divide the various kinds of waves that occur into two basic types:

\begin{center}
\rule{300pt}{1pt}

\textbf{Transverse wave:} a traveling wave or pulse that causes the elements
of the disturbed medium to move perpendicular to the direction of propagation

\rule{300pt}{1pt}

\textbf{Longitudinal wave}: a traveling wave or pulse that causes the
elements of the medium to move parallel to the direction of propagation

\rule{300pt}{1pt}
\end{center}

\FRAME{dhF}{2.6705in}{1.7218in}{0pt}{}{}{Figure}{\special{language
"Scientific Word";type "GRAPHIC";maintain-aspect-ratio TRUE;display
"USEDEF";valid_file "T";width 2.6705in;height 1.7218in;depth
0pt;original-width 3.781in;original-height 2.4284in;cropleft "0";croptop
"1";cropright "1";cropbottom "0";tempfilename
'SUS73U1D.wmf';tempfile-properties "XPR";}}

Let's look at some specific cases of wave motion.

\FRAME{dhF}{1.3837in}{1.8273in}{0pt}{}{}{Figure}{\special{language
"Scientific Word";type "GRAPHIC";maintain-aspect-ratio TRUE;display
"USEDEF";valid_file "T";width 1.3837in;height 1.8273in;depth
0pt;original-width 1.3482in;original-height 1.7902in;cropleft "0";croptop
"1";cropright "1";cropbottom "0";tempfilename
'SUS73U1E.wmf';tempfile-properties "XPR";}}

In the picture above, you see wave that has just one peak traveling to the
right. We call such a wave a \emph{pulse}. The rope material is the medium
for this wave and notice how the piece of the rope marked $P$ moves up and
down, but the wave is moving to the right. This pulse is a transverse wave
because the parts of the medium (observe point $P$) move perpendicular to
the direction the wave is moving.

Of course, some waves are a combination of these two basic types. You may
have noticed that in Physics we tend to define basic types of things, and
then use these basic types to define more complex objects. Water waves, for
example, are transverse at the surface of the water, but are longitudinal
throughout the water.

\FRAME{dtbpF}{1.6622in}{0.7835in}{0pt}{}{}{Figure}{\special{language
"Scientific Word";type "GRAPHIC";maintain-aspect-ratio TRUE;display
"USEDEF";valid_file "T";width 1.6622in;height 0.7835in;depth
0pt;original-width 5.5901in;original-height 2.6212in;cropleft "0";croptop
"1";cropright "1";cropbottom "0";tempfilename
'SUS73U1F.wmf';tempfile-properties "XPR";}}

Earthquakes produce both transverse and longitudinal waves. The two types of
waves even travel at different speeds! $P$ waves are longitudinal and travel
faster, $S$ waves are transverse and slower.

\section{Wave speed}

We can perform an experiment with a rope or a long spring. Make a wave on
the rope or spring. Then pull the rope or spring tighter and make another
wave. We see that the wave on the tighter spring travels faster.

It is harder to do, but we can also experiment with two different ropes, one
light and one heavy. We would find that the heaver the rope, the slower the
wave. We can express this as 
\begin{equation*}
v=\sqrt{\frac{T_{s}}{\mu }}
\end{equation*}%
where $T_{s}$ is the tension in the rope, and $\mu $ is the linear mass
density%
\begin{equation*}
\mu =\frac{m}{L}
\end{equation*}%
where $m$ is the mass of the rope, and $L$ is the length.

The term $\mu $ might need an analogy to make it seem helpful. So suppose I
have an iron bar that has a mass of $200\unit{kg}$ and is $2\unit{m}$ long.
Further suppose I want to know how much mass there would be in a $20\unit{cm}
$ section cut of the end of the rod. How would I find out?

This is not very hard, We could say that there are $200\unit{kg}$ spread out
evenly over $2\unit{m},$ so each meter of rod has $100\unit{kg}$ of mass,
that is, there is $100\unit{kg}/\unit{m}$ of mass per unit length. Then to
find how much mass there is in a $0.20\unit{m}$ section of the rod I\ take 
\begin{equation*}
m=100\frac{\unit{kg}}{\unit{m}}\times 0.20\unit{m}=20.0\unit{kg}
\end{equation*}%
The $100\unit{kg}/\unit{m}$ is $\mu .$ It is how much mass there is in a
unit length segment of something In this example, it is a unit length of
iron bar, but for waves on a rope, we want the mass per unit length of rope.

If you are buying stock steel bar, you might be able to buy it with a mass
per unit length. If the mass per unit length is higher then the bar is more
massive. The same is true with string. The larger $\mu ,$ the more massive
equal rope segments will be.

We should note that in forming this relationship, we have used our standard
introductory physics assumption that the mass of the rope can mostly be
neglected. Let's consider what would happen if this were not true. Say we
make a wave in a heavy cable that is suspended. The mass at the lower end of
the cable pulls down on the upper part of the cable. The tension will
actually change along the length of the cable, and so will the wave speed.
Such a situation can't be represented by a single wave speed. But for our
class, we will assume that any such changes are small enough to be ignored.

\section{One dimensional waves}

To mathematically describe a wave we will define a function of both time and
position.%
\begin{equation}
y\left( x,t\right)
\end{equation}%
Note, that this is new in our physics experience. Before we usually were
concerned about only one variable at a time. For oscillation we had just $%
y\left( t\right) $ for example. But now we will be concerned about two
variables, $x,$ and $t.$

let's take a specific example\footnote{%
This is not an important wave function, just one I picked because it makes a
nice graphic example.}%
\begin{equation}
y=\frac{2}{\left( x-3.0t\right) ^{2}+1}
\end{equation}%
Let's plot this for $t=0$

\FRAME{dtbpF}{2.5749in}{1.7261in}{0pt}{}{}{Figure}{\special{language
"Scientific Word";type "GRAPHIC";maintain-aspect-ratio TRUE;display
"USEDEF";valid_file "T";width 2.5749in;height 1.7261in;depth
0pt;original-width 2.8241in;original-height 1.8839in;cropleft "0";croptop
"1";cropright "1";cropbottom "0";tempfilename
'T8LZ4I05.wmf';tempfile-properties "XPR";}}What will this look like for $%
t=10 $?\FRAME{dtbpF}{2.6503in}{1.7731in}{0pt}{}{}{Figure}{\special{language
"Scientific Word";type "GRAPHIC";maintain-aspect-ratio TRUE;display
"USEDEF";valid_file "T";width 2.6503in;height 1.7731in;depth
0pt;original-width 3.194in;original-height 2.1288in;cropleft "0";croptop
"1";cropright "1";cropbottom "0";tempfilename
'T8LZ6A09.wmf';tempfile-properties "XPR";}}

The pulse travels along the $x$-axis as a function of time. Note that there
is a value for $y$ for every $x$ position and that these $y$ values change
for different times. That is what we mean by saying we have a function of
both $x$ and $t.$

We denote the speed of the pulse as $v,$ then we can define a function 
\begin{equation}
y\left( x,t\right) =y\left( x-vt,0\right)
\end{equation}%
that describes a pulse as it travels. An element of the medium (rope,
string, etc.) at position $x$ at some time $t,$ will have the displacement
that an element had earlier at $x-vt$ when $t=0.$

We will give $y\left( x-vt,0\right) $ a special name, the \emph{wave function%
}. It represents the $y$ position, the transverse position in our example,
of any element located at a position $x$ at any time $t$.

Notice that wave functions depend on two variables, $x,$ and $t.$ It is hard
to draw a wave so that this dual dependence is clear. Often we draw two
different graphs of the same wave so we can see independently the position
and time dependence. So far we have used one of these graphs. A graph of our
wave at a specific time, $t_{o}.$ This gives $y\left( x,t_{o}\right) $. This
representation of a wave is very like a photograph of the wave taken with a
digital camera. It gives a picture of the entire wave, but only for one
time, the time at which the photograph was taken. Of course we could take a
series of photographs, but still each would be a picture of the wave at just
one time. Here is such a series of graphs at times $t_{1}$ through $t_{4}.$%
\FRAME{dtbpF}{2.4993in}{3.7472in}{0in}{}{}{Figure}{\special{language
"Scientific Word";type "GRAPHIC";maintain-aspect-ratio TRUE;display
"USEDEF";valid_file "T";width 2.4993in;height 3.7472in;depth
0in;original-width 2.4587in;original-height 3.6997in;cropleft "0";croptop
"1";cropright "1";cropbottom "0";tempfilename
'SUS73U1P.wmf';tempfile-properties "XPR";}}

The second representation is to observe the wave at just one point in the
medium, but for many times. This is very like taking a video camera and
using it to record the displacement of just one part of the medium for many
times. You could envision marking just one part of a rope, and then using
the video recorder to make a movie of the motion of that single part of the
rope. We could then go frame by frame through the video, and plot the
displacement of our marked part of the rope as a function of time. Such a
graph is sometimes called a history graph of the wave. Here is such a graph
for the position $x_{1}.$ \FRAME{dtbpF}{2.5918in}{1.7279in}{0pt}{}{}{Figure}{%
\special{language "Scientific Word";type "GRAPHIC";maintain-aspect-ratio
TRUE;display "USEDEF";valid_file "T";width 2.5918in;height 1.7279in;depth
0pt;original-width 2.5503in;original-height 1.6916in;cropleft "0";croptop
"1";cropright "1";cropbottom "0";tempfilename
'SUS73U1Q.wmf';tempfile-properties "XPR";}}To go from the time graphs to the
history graph you observe what happens at the location $x_{1}$ for each of
the times. Then plot those $y$ positions on the $y$ vs. $t$ graph. Then you
must connect the points. This takes some thinking to make sure you connect
them right (or a whole lot of points). But with practice, this is not hard
and both view points are valuable ways to look at a wave.

\section{Sinusoidal Waves}

\FRAME{fhFU}{2.5873in}{2.1128in}{0pt}{\Qcb{{\protect\small Two snapshot
views of a sinusoidal wave.}}}{\Qlb{Sine_wave}}{Figure}{\special{language
"Scientific Word";type "GRAPHIC";maintain-aspect-ratio TRUE;display
"USEDEF";valid_file "T";width 2.5873in;height 2.1128in;depth
0pt;original-width 2.2667in;original-height 1.8455in;cropleft "0";croptop
"1";cropright "1";cropbottom "0";tempfilename
'SUS73U1R.wmf';tempfile-properties "XPR";}}A sinusoidal graph should be
familiar from our study of oscillation. Simple harmonic oscillators are
described by the function%
\begin{equation}
y\left( t\right) =y_{\max }\cos \left( \omega t+\phi \right) \qquad \text{SHM%
}
\end{equation}%
but this only gave us a vertical displacement. Now our sinusoidal function
must also be a function of position along the wave.%
\begin{equation}
y\left( x,t\right) =y_{\max }\sin \left( kx-\omega t+\phi \right) \qquad 
\text{waves}
\end{equation}%
but before we study the nature of this function, lets see what we can learn
from the graph of a sinusoidal wave. We will need both of our two views, the
snapshot graph and the history graph. Look at figure (\ref{Sine_wave}). This
is two camera snap shots superimposed. The red curve shows the wave ($y$
position for each value of $x)$ at $t=0.$ At some some later time $t,$ the
wave pattern has moved to the right as shown by the blue curve. The shift is
by an amount $\Delta x=v\Delta t.$

\section{Parts of a wave in snapshot and history graphs}

The peak of a wave is called the crest. For a sine wave we have a series of
crests. We define a special length called \emph{the wavelength} as the
distance between any to identical points (e.g. crests) on adjacent waves in
a snapshot view.

Notice that this is very similar to the definition of the period, $T,$ when
we graphed SHM on a position vs. time graph. In fact, this similarity is
even more apparent if we plot a sinusoidal wave using our two wave pictures.
In the next figure, the snap-shot graph comes first. We can see that there
will be crests. The distance between the crests is given the name
\textquotedblleft wavelength.\textquotedblright\ This is not the entire
length of the whole wave. But it is a characteristic length of part of the
wave that is easy to identify. It is defined as the crest to crest distance. 
\FRAME{dtbpF}{2.2831in}{4.0136in}{0pt}{}{}{Figure}{\special{language
"Scientific Word";type "GRAPHIC";maintain-aspect-ratio TRUE;display
"USEDEF";valid_file "T";width 2.2831in;height 4.0136in;depth
0pt;original-width 4.34in;original-height 7.6662in;cropleft "0";croptop
"1";cropright "1";cropbottom "0";tempfilename
'SUS73U1S.wmf';tempfile-properties "XPR";}}Note that there are crests in the
history graph view as well. That is because one marked part of the medium is
being displaced as a function of time. Think of a marked piece of the rope
going up and down, or think of floating in the ocean at one point. You
travel up and down (and a little bit back and forth) as the waves go by. But
now the horizontal axis is time. There will be a characteristic time between
crests. That time is called the \emph{period}. This is just like oscillatory
motion--because it is oscillatory motion! Like the wavelength is not the
length of the whole wave, the period is not the time the whole wave exists.
It is just the time it takes the part of the medium we are watching to go
through one complete cycle. Notice that this history picture is exactly the
same as a plot of the motion of a simple harmonic oscillator! For a
sinusoidal wave, each part of the medium experiences simple harmonic motion.

We remember frequency from simple harmonic motion. But now we have a wave,
and the wave is moving. We can extend our view of frequency by defining it
as follows:

\begin{center}
\rule{300pt}{1pt}

\emph{The frequency of a periodic wave is the number of crests (or any other
point of the wave) that pass a given point in a unit time interval. }

\rule{300pt}{1pt}
\end{center}

In the next figure, the blue curve has twice the frequency as the red curve.

\FRAME{dtbpF}{2.9049in}{1.9389in}{0in}{}{}{Figure}{\special{language
"Scientific Word";type "GRAPHIC";maintain-aspect-ratio TRUE;display
"USEDEF";valid_file "T";width 2.9049in;height 1.9389in;depth
0in;original-width 2.9315in;original-height 1.9469in;cropleft "0";croptop
"1";cropright "1";cropbottom "0";tempfilename
'TDYU1B03.wmf';tempfile-properties "XPR";}} Notice how it has two crests for
every red crest. The maximum displacement of the wave is called the \emph{%
amplitude }just as it was for simple harmonic oscillators.

\section{Wavenumber and wave speed}

Consider again a sinusoidal wave 
\begin{equation}
y\left( x,t\right) =y_{\max }\sin \left( kx-\omega t+\phi _{o}\right)
\end{equation}%
with $\phi _{o}=0$ 
\begin{equation}
y\left( x,t\right) =y_{\max }\sin \left( kx-\omega t\right)
\end{equation}%
Let's draw a snapshot graph for this wave. \FRAME{dtbpFX}{2.6109in}{1.7417in%
}{0pt}{}{}{Plot}{\special{language "Scientific Word";type "MAPLEPLOT";width
2.6109in;height 1.7417in;depth 0pt;display "USEDEF";plot_snapshots
TRUE;mustRecompute FALSE;lastEngine "MuPAD";xmin "0";xmax "10";xviewmin
"0";xviewmax "10";yviewmin "-1.000187";yviewmax
"1.000187";viewset"XY";rangeset"X";plottype 4;axesFont "Times New
Roman,12,0000000000,useDefault,normal";numpoints 100;plotstyle
"patch";axesstyle "normal";axestips FALSE;xis \TEXUX{v58130};var1name
\TEXUX{$\theta $};function \TEXUX{$\sin \left( \theta \right) $};linecolor
"blue";linestyle 1;pointstyle "point";linethickness 1;lineAttributes
"Solid";var1range "0,10";num-x-gridlines 100;curveColor
"[flat::RGB:0x000000ff]";curveStyle "Line";VCamFile
'TJXPK440.xvz';valid_file "T";tempfilename
'TDYTRM00.wmf';tempfile-properties "XPR";}}To make this graph, we set $t=0$
and plotted the resulting function%
\begin{equation}
y\left( x,0\right) =y_{\max }\sin \left( kx+0\right)
\end{equation}%
$y_{\max }$ is the amplitude. I want to investigate the meaning of the
constant $k.$ It looks like the spring constant, but it's not! Lets find $k$
like we did for SHM when we found $\omega $. Consider the point $x=0$. We
still have $t=0$ for this graph. And at this position and time 
\begin{equation}
y\left( 0,0\right) =y_{\max }\sin \left( k\left( 0\right) \right) =0
\end{equation}%
The next time $y=0$ is when $x=\frac{\lambda }{2}$ \FRAME{dtbpF}{2.8781in}{%
2.2373in}{0pt}{}{}{Figure}{\special{language "Scientific Word";type
"GRAPHIC";maintain-aspect-ratio TRUE;display "USEDEF";valid_file "T";width
2.8781in;height 2.2373in;depth 0pt;original-width 4.4858in;original-height
3.4809in;cropleft "0";croptop "1";cropright "1";cropbottom "0";tempfilename
'SUS73U1V.wmf';tempfile-properties "XPR";}}then%
\begin{equation}
y\left( \frac{\lambda }{2},0\right) =y_{\max }\sin \left( k\frac{\lambda }{2}%
\right) =0
\end{equation}%
From our trigonometry experience, we know that this is true when 
\begin{equation}
k\frac{\lambda }{2}=\pi
\end{equation}%
solving for $k$ gives%
\begin{equation}
k=\frac{2\pi }{\lambda }
\end{equation}%
Then we now have a feeling for what $k$ means. It is $2\pi $ over the
spacing between the crests. The $2\pi $ must have units of radians attached.
Then%
\begin{equation}
y\left( x,0\right) =A\sin \left( \frac{2\pi }{\lambda }x+0\right)
\end{equation}

We have a special name for the quantity $k.$ It is called the \emph{wave
number}$.$ 
\begin{equation}
k\equiv \frac{2\pi }{\lambda }
\end{equation}

Both the name and the symbol are somewhat unfortunate. Neither gives much
insight into the meaning of this quantity. But from what we have done in
studying oscillation, we can understand this new quantity. For a harmonic
oscillator, we know that 
\begin{equation*}
y\left( t\right) =y_{\max }\sin \left( \omega t\right)
\end{equation*}%
where 
\begin{equation*}
\omega =2\pi f=\frac{2\pi }{T}
\end{equation*}%
$T$ is how far, \emph{in time}, the crests are apart, and the inverse of
this, $\frac{1}{T}$ is the frequency. The frequency tells us how often we
encounter a crest as we march along in time. So $\frac{1}{T}$ must be how
many crests we have in a unit amount of time.

Now think of the relationship between the snapshot and the video
representation for a sinusoidal wave. We have a new quantity 
\begin{equation*}
k=\frac{2\pi }{\lambda }
\end{equation*}%
where $\lambda $ is how far, \emph{in distance}, the crests are apart. This
implies that $\frac{1}{\lambda }$ plays the same role in the snap shot graph
that $f$ plays in the video graph. It must tell us how many crests we have,
but this time it is how many crests in a given amount of distance. We found
above that $k$ told us something about how often the zeros (well, every
other zero) will occur. But the crests must occur at the same rate. So $k$
tells us how often we encounter a crest in our snapshot graph. This is not
too strange. It would be meaningful, for example, to say that a farmer had
plowed his or her field to have two furrows per meter.\FRAME{dtbpF}{2.9853in%
}{1.6198in}{0in}{}{}{Figure}{\special{language "Scientific Word";type
"GRAPHIC";maintain-aspect-ratio TRUE;display "USEDEF";valid_file "T";width
2.9853in;height 1.6198in;depth 0in;original-width 2.9421in;original-height
1.5835in;cropleft "0";croptop "1";cropright "1";cropbottom "0";tempfilename
'SUS73U1W.wmf';tempfile-properties "XPR";}}

The frequency, $\frac{1}{T},$ in the history graph is how often we encounter
a crest, $\frac{1}{\lambda }$ is how often we encounter a crest in the snap
shot graph. Thus $\frac{1}{\lambda }$ is playing the same role for a snap
shot graph as frequency plays for a history graph. We could call $\frac{1}{%
\lambda }$a \emph{spatial frequency. }It is how often we encounter a crest
as we march along in position, or how many crests we have in a unit amount
of distance. And $\lambda $ could be called a \emph{spatial period}. Both $%
1/T$ and $1/\lambda $ answer the question \textquotedblleft how often
something happens in a unit of something\textquotedblright\ but one asks the
question in time and the other in position along the wave.

My mental image for this is the set of groves on the edge of a highway.
There is a distance between them, like a wavelength, and how often I
encounter one as I move a distance along the road is the spatial frequency.
You could say that there are, maybe, $3/\unit{m}$. That is a spatial
frequency! It is how many of something happens in a unit distance. We use
this concept in optics to test how well an optical system resolves details
in a photograph. The next figure is a test image. A good camera will resolve
all spatial frequencies equally well. Notice the test image has sets of bars
with different spatial frequencies. By forming an image of this pattern, you
can see which spacial frequencies are faithfully represented by the optical
system.\FRAME{dtbpFU}{2.6974in}{2.1819in}{0pt}{\Qcb{Resolution test target
based on the USAF\ 1951 Resolution Test Pattern (not drawn to exact
specifications).}}{}{Figure}{\special{language "Scientific Word";type
"GRAPHIC";maintain-aspect-ratio TRUE;display "USEDEF";valid_file "T";width
2.6974in;height 2.1819in;depth 0pt;original-width 2.655in;original-height
2.143in;cropleft "0";croptop "1";cropright "1";cropbottom "0";tempfilename
'SUS73V1X.wmf';tempfile-properties "XPR";}}In class you will see that our
projector does not represent all spatial frequencies equally well! You can
also see this now in the copy you are reading. If you are reading on-line or
an electronic copy, your screen resolution will limit the representation of
some spatial frequencies. Look for the smallest set of three bars where you
can still tell for sure that there are three bars without zooming. A printed
version that has been printed on a laser printer will usually allow you to
see even smaller sets of three bars clearly.

Let's place $k$ in the full equation for the sine wave for any time, $t$.%
\begin{equation}
y\left( x,t\right) =y_{\max }\sin \left( kx-\omega t+\phi _{o}\right)
\end{equation}%
Sometimes you will see this equation written in a slightly different way. We
could change our variable from $x$ to $x-vt$ so that 
\begin{equation*}
y\left( x,t\right) =y\left( x-vt,0\right)
\end{equation*}%
With a little algebra we can do this%
\begin{eqnarray*}
y\left( t\right) &=&y_{\max }\sin \left( kx-\omega t+\phi _{o}\right) \\
&=&y_{\max }\sin \left( \frac{2\pi }{\lambda }x-\frac{2\pi }{T}t+\phi
_{o}\right) \\
&=&y_{\max }\sin \left( \frac{2\pi }{\lambda }\left( x-\frac{\lambda }{T}%
t\right) +\phi _{o}\right)
\end{eqnarray*}%
This is in the form of a wave function so long as%
\begin{equation}
v=\frac{\lambda }{T}
\end{equation}%
which is just how far the wave travels from crest to crest divided by how
long it takes to travel that distance. The wave travels one wavelength in
one period. Indeed, this is the wave speed! Then 
\begin{equation}
y\left( x,t\right) =y_{\max }\sin \left( \frac{2\pi }{\lambda }\left(
x-vt\right) +\phi _{o}\right)
\end{equation}%
or, using $k$ 
\begin{equation}
y\left( x,t\right) =y_{\max }\sin \left( k\left( x-vt\right) +\phi
_{o}\right)
\end{equation}

\subsubsection{Wave speed forms}

We just found%
\begin{equation}
v=\frac{\lambda }{T}
\end{equation}%
but it is easy to see that 
\begin{equation}
v=\frac{2\pi \lambda }{2\pi T}=\frac{\omega }{k}
\end{equation}%
and 
\begin{equation}
v=\lambda f
\end{equation}%
This last formula is, perhaps, the most common form encountered in our study
of light waves.

\subsubsection{Phase}

You may be wondering about the phase constant we learned about in our study
of SHM. We have ignored it up to now. But of course we can shift our sine
just like we did for our plots of position vs. time for oscillation. Only
now with a wave we have two graphs, a history and snapshot graphs, so we
could shift along the $x$ in a snapshot graph or along the $t$ axes in a
history graph. So the sine wave has the form.%
\begin{equation}
y\left( x,t\right) =A\sin \left( kx-\omega t+\phi _{o}\right)
\end{equation}%
were $\phi _{o}$ will need to be determined by initial conditions just like
in SHM problems and those initial conditions will include initial positions
as well as initial times.

Let's consider that we have two views of a wave, the snapshot and history
view. Each of these looks like sinusoids for a sinusoidal wave. Let's
consider a specific wave, 
\begin{equation*}
y\left( x,t\right) =5\sin \left( 3\pi x-\frac{\pi }{5}t+\frac{\pi }{2}\right)
\end{equation*}%
we look at a snapshot graph at $t=0$ 
\begin{equation*}
y\left( x,0\right) =5\sin \left( 3\pi x-\frac{\pi }{5}\left( 0\right) +\frac{%
\pi }{2}\right)
\end{equation*}%
\FRAME{dtbpFX}{4.4996in}{0.7706in}{0pt}{}{}{Plot}{\special{language
"Scientific Word";type "MAPLEPLOT";width 4.4996in;height 0.7706in;depth
0pt;display "USEDEF";plot_snapshots TRUE;mustRecompute FALSE;lastEngine
"MuPAD";xmin "-3";xmax "3";xviewmin "-3";xviewmax "3";yviewmin
"-5.000998";yviewmax "5";viewset"XY";rangeset"X";plottype 4;axesFont "Times
New Roman,12,0000000000,useDefault,normal";numpoints 100;plotstyle
"patch";axesstyle "normal";axestips FALSE;xis \TEXUX{x};var1name
\TEXUX{$x$};function \TEXUX{$5\sin \left( 3\pi x-0+\frac{\pi }{2}\right)
$};linecolor "blue";linestyle 1;pointstyle "point";linethickness
3;lineAttributes "Solid";var1range "-3,3";num-x-gridlines 100;curveColor
"[flat::RGB:0x000000ff]";curveStyle "Line";VCamFile
'SV4MPV0V.xvz';valid_file "T";tempfilename
'SUS73V1Y.wmf';tempfile-properties "XPR";}}and another at $t=2\unit{s}$%
\begin{equation*}
y\left( x,2\unit{s}\right) =5\sin \left( 3\pi x-\frac{\pi }{5}\left( 2\unit{s%
}\right) +\frac{\pi }{2}\right)
\end{equation*}%
\FRAME{dtbpFX}{4.4996in}{0.7706in}{0pt}{}{}{Plot}{\special{language
"Scientific Word";type "MAPLEPLOT";width 4.4996in;height 0.7706in;depth
0pt;display "USEDEF";plot_snapshots TRUE;mustRecompute FALSE;lastEngine
"MuPAD";xmin "-3";xmax "3";xviewmin "-3";xviewmax "3";yviewmin
"-5.000998";yviewmax "5";viewset"XY";rangeset"X";plottype 4;axesFont "Times
New Roman,12,0000000000,useDefault,normal";numpoints 100;plotstyle
"patch";axesstyle "normal";axestips FALSE;xis \TEXUX{x};var1name
\TEXUX{$x$};function \TEXUX{$5\sin \left( 3\pi x-\frac{\pi }{5}\left(
2\right) +\frac{\pi }{2}\right) $};linecolor "blue";linestyle 1;pointstyle
"point";linethickness 3;lineAttributes "Solid";var1range
"-3,3";num-x-gridlines 100;curveColor "[flat::RGB:0x000000ff]";curveStyle
"Line";VCamFile 'SV4MPZ0W.xvz';valid_file "T";tempfilename
'SUS73V1Z.wmf';tempfile-properties "XPR";}}Comparing the two, we could view
the latter as having a different phase constant that is the sum of what we
have called the phase constant, $\phi _{o}$ plus $\omega \Delta t,$ which
tells us how different the times are between the two graphs. This time
difference is what is further shifting the graph. 
\begin{equation*}
\phi _{total}=\omega \Delta t+\phi _{o}=-\frac{\pi }{5}\left( 2\unit{s}%
\right) +\frac{\pi }{2}
\end{equation*}%
that is, within the snapshot view, the time dependent part of the argument
of the sign acts like an additional phase constant.

Likewise, in the history view, we can plot our wave at $x=0$%
\begin{equation*}
y\left( 0,t\right) =5\sin \left( 3\pi \left( 0\right) -\frac{\pi }{5}t+\frac{%
\pi }{2}\right)
\end{equation*}%
\FRAME{dtbpFX}{4.4996in}{0.7706in}{0pt}{}{}{Plot}{\special{language
"Scientific Word";type "MAPLEPLOT";width 4.4996in;height 0.7706in;depth
0pt;display "USEDEF";plot_snapshots TRUE;mustRecompute FALSE;lastEngine
"MuPAD";xmin "-8";xmax "8";xviewmin "-8";xviewmax "8";yviewmin
"-5.000998";yviewmax "5";viewset"XY";rangeset"X";plottype 4;labeloverrides
1;x-label "t";axesFont "Times New
Roman,12,0000000000,useDefault,normal";numpoints 100;plotstyle
"patch";axesstyle "normal";axestips FALSE;xis \TEXUX{t};var1name
\TEXUX{$t$};function \TEXUX{$5\sin \left( 3\pi \left( 0\right) -\frac{\pi
}{5}t+\frac{\pi }{2}\right) $};linecolor "blue";linestyle 1;pointstyle
"point";linethickness 3;lineAttributes "Solid";var1range
"-8,8";num-x-gridlines 100;curveColor "[flat::RGB:0x000000ff]";curveStyle
"Line";VCamFile 'SV4MQ10X.xvz';valid_file "T";tempfilename
'SUS73V20.wmf';tempfile-properties "XPR";}}and at $x=1.5\unit{m}$%
\begin{equation*}
y\left( 1.5\unit{m},t\right) =5\sin \left( 3\pi \left( 1.5\unit{m}\right) -%
\frac{\pi }{5}t+\frac{\pi }{2}\right)
\end{equation*}%
\FRAME{dtbpFX}{4.4996in}{0.7706in}{0pt}{}{}{Plot}{\special{language
"Scientific Word";type "MAPLEPLOT";width 4.4996in;height 0.7706in;depth
0pt;display "USEDEF";plot_snapshots TRUE;mustRecompute FALSE;lastEngine
"MuPAD";xmin "-8";xmax "8";xviewmin "-8";xviewmax "8";yviewmin
"-5.000998";yviewmax "5";viewset"XY";rangeset"X";plottype 4;labeloverrides
1;x-label "t";axesFont "Times New
Roman,12,0000000000,useDefault,normal";numpoints 100;plotstyle
"patch";axesstyle "normal";axestips FALSE;xis \TEXUX{t};var1name
\TEXUX{$t$};function \TEXUX{$5\sin \left( 3\pi \left( 1.5\right) -\frac{\pi
}{5}t+\frac{\pi }{2}\right) $};linecolor "blue";linestyle 1;pointstyle
"point";linethickness 3;lineAttributes "Solid";var1range
"-8,8";num-x-gridlines 100;curveColor "[flat::RGB:0x000000ff]";curveStyle
"Line";VCamFile 'SV4MQ20Y.xvz';valid_file "T";tempfilename
'SUS73V21.wmf';tempfile-properties "XPR";}}Within the history view, the $kx$
part of the argument acts like a additional phase constant.%
\begin{equation*}
\phi _{total}=k\Delta x+\phi _{o}=3\pi \left( 1.5\unit{m}\right) +\frac{\pi 
}{2}
\end{equation*}

Of course neither $kx$ nor $\omega t$ are constant, But within individual
views of the wave we have set them constant to form our snapshot and history
representations. We can see that any part of the argument of the sine, $%
kx-\omega t+\phi _{o}$ could contribute to a phase shift, depending on the
view we are taking.

Because of this, it is customary to call the entire argument of the sine
function, $\phi =kx-\omega t+\phi _{o}$ the \emph{phase of the wave}. Where $%
\phi _{o}$ is the phase constant, $\phi $ is the phase. Of course then, $%
\phi $ must be a function of $x$ and $t,$ so we have a different value for $%
\phi \left( x,t\right) $ for every point on the wave for every time. This
can be a little confusing, $\phi $ and $\phi _{o}$ look a lot the same, but
they are different.

\section{The Speed of Waves}

We have already considered the speed of waves on strings or ropes%
\begin{equation*}
v=\sqrt{\frac{T}{\mu }}
\end{equation*}%
we should see where this comes from so we know it's limitations. We would
expect that there may be things that change the speed of sound waves in air
as well (like a tension or a massiveness of the medium). Let's start with a
formal derivation of our string speed formula, then take on sound waves in
air.

In the following figure we have a yellow string with a pulse. The top of the
pulse has been marked as darker brown. Let's use this figure to find the
wave speed.\FRAME{dhF}{2.7397in}{1.4364in}{0pt}{}{}{Figure}{\special%
{language "Scientific Word";type "GRAPHIC";maintain-aspect-ratio
TRUE;display "USEDEF";valid_file "T";width 2.7397in;height 1.4364in;depth
0pt;original-width 2.6974in;original-height 1.4019in;cropleft "0";croptop
"1";cropright "1";cropbottom "0";tempfilename
'SUS73V22.wmf';tempfile-properties "XPR";}}

Start by asking what forces are acting on our brown element of string?

\begin{enumerate}
\item Tension on the RHS of the element from the rest of the string on the
right, $T_{r}$

\item Tension on the LHS of the element from the rest of the string on the
left, $T_{l}$

\item $W$ the gravitational force
\end{enumerate}

Lets assume that the element of string, $\Delta s,$ at the crest is
approximately an arc of a circle with radius $R$.

There is a force pulling left on the left end of the element that is tangent
to the arc, there is a force pulling right at the right end of the element
which is also tangent to the arc. These forces produce centripetal
accelerations%
\begin{equation}
a_{r}=\frac{v_{t}^{2}}{R}
\end{equation}%
The horizontal components of the forces cancel $\left( T\cos \theta \right)
. $ The vertical component, $\left( T\sin \left( \theta \right) \right) $ is
directed toward the center of the arc. If the rope element is not moving in
the $x$ direction, then $T_{left}=T_{right}.$

Then, the radial force $F_{net_{r}}$ will have matching components from each
side of the element that together are $2T\sin \left( \theta \right) .$
Notice that to make sure the components are the same we must assume that the
string is uniform and that it is not too massive (almost massless string
approximation!). Since the element is small, the angle $\theta $ will also
be small%
\begin{equation}
F_{net_{r}}=2T\sin \left( \theta \right) \approx 2T\theta
\end{equation}

The element has a mass $m.$ 
\begin{equation}
m=\mu \Delta s
\end{equation}

where $\mu $ is the mas per unit length. Using the arc length formula $%
\left( s=R\theta \right) $ the length of the element is 
\begin{equation}
\Delta s=R\left( 2\theta \right)
\end{equation}%
so%
\begin{equation}
m=\mu \Delta s=2\mu R\theta
\end{equation}%
and finally we use the formula for the radial acceleration%
\begin{equation}
F_{net_{r}}=ma_{r}=\left( 2\mu R\theta \right) \frac{v^{2}}{R}
\end{equation}%
Combining these two expressions for $F_{net_{r}}$%
\begin{equation}
2T\theta =\left( 2\mu R\theta \right) \frac{v^{2}}{R}
\end{equation}%
\begin{equation}
T=\left( \mu R\right) \frac{v^{2}}{R}
\end{equation}%
\begin{equation}
\frac{T}{\mu }=v^{2}
\end{equation}%
and we find that 
\begin{equation}
v=\sqrt{\frac{T}{\mu }}
\end{equation}%
which is just what we stated before, only now we see just how much
approximation there was in building this formula! We might guess that it
would not work so well for large cables supporting bridges or with cables
that change size along their length. But it works fine for our semi-massless
strings.

\subsection{Speed of the medium}

Let's review for a moment. Suppose we take a very long jump rope, and shake
one end up and down while a lab partner keeps his or her end stationary. You
can make a sine wave in the rope. But being good physics students, let's
make our \textquotedblleft shake\textquotedblright\ very symmetric. Let's
start our wave using simple harmonic motion.

Let's call an element of the rope $\Delta x.$ Here the \textquotedblleft $%
\Delta $\textquotedblright\ is being used to mean \textquotedblleft a small
amount of.\textquotedblright\ We are taking a small amount of the rope and
calling it's length $\Delta x.$ And it is the motion of this small amount of
rope that I\ want to think about now.

Each element of the rope $\left( \Delta x_{i}\right) $ will also oscillate
with SHM (think of a driven SHO). Note that the elements of the rope
oscillate in the $y$ direction, but the wave travels in the $x$-direction.
This is a transverse wave.

Let's describe the motion of an element of the string at some point $P.$

At $t=0,$%
\begin{equation}
y=y_{\max }\sin \left( kx-\omega t+\phi _{o}\right)
\end{equation}%
The element does not move in the $x$ direction. So, it has no $y$-component $%
\left( v_{x}=0\right) $. But it does move in the $y$-direction so there will
be a $v_{y}.$ Let's give this $y$-componet a name. We can call it the \emph{%
transverse} \emph{speed}, $v_{y},$ and we can call a change in this
transverse speed in an amount of time the \emph{transverse acceleration, }$%
a_{y}.$ These are just the velocity and acceleration of the element of rope
in the $y$ direction. But in our study of waves it is good to know that
these are in what we will call the transverse direction. These are not the
velocity and acceleration of the wave, just the velocity and acceleration of
the element $\Delta x$ of the medium at a point $P.$

Because we are doing this at one specific $x$ location we need partial
derivatives to find the velocity

\begin{equation}
v_{y}=\left. \frac{dy}{dt}\right] _{x=\text{constant}}=\frac{\partial y}{%
\partial t}
\end{equation}%
That is, we take the derivative of $y$ with respect to $t,$ but we pretend
that $x$ is not a variable because we just want one $x$ position. Then%
\begin{equation}
v_{y}=\frac{\partial y}{\partial t}=-\omega y_{\max }\cos \left( kx-\omega
t+\phi _{o}\right)
\end{equation}%
and%
\begin{equation}
a_{y}=\frac{\partial v_{y}}{\partial t}=-\omega ^{2}y_{\max }\sin \left(
kx-\omega t+\phi _{o}\right)
\end{equation}

These solutions should look very familiar! We expect them to be the same as
a harmonic oscillator except that we now have to specify which
oscillator--which part of the rope--we are looking at. That is what the $kx$
part is doing.

\section{The Linear Wave Equation}

We added in some new math in our last two lectures. Differential equations!
We identified simple harmonic motion using a differential equation%
\begin{equation*}
\frac{d^{2}x}{dt^{2}}=-\omega ^{2}x
\end{equation*}%
Perhaps we could find an equation like this for waves and then identify
waves using that equation. Can we find such an equation that would let us
identify waves?

To try to achieve this we can start the same way we did for oscillation.
Let's write the acceleration for a part of the medium.%
\begin{equation}
a_{y}=\frac{d^{2}y}{dt^{2}}=-\omega ^{2}y_{\max }\sin \left( kx-\omega
t+\phi _{o}\right)
\end{equation}

But for waves we have two variables. Let's guess that we might need to take
derivatives with respect to $x$ as well as for $t.$ 
\begin{equation*}
\frac{dy}{dx}=ky_{\max }\cos (kx-\omega t+\phi _{o})
\end{equation*}%
This isn't all that useful on it's own, but it does tell us how the slope of
the curve changes in a snapshot graph. Let's take another derivative%
\begin{equation*}
\frac{d^{2}y}{dx^{2}}=-k^{2}y_{\max }\sin (kx-\omega t+\phi _{o})
\end{equation*}%
We are back to a sine function. That is not a surprise. But it looks a lot
like our time second derivative. We could take the ratio of the two
equations to see just how different they are%
\begin{equation*}
\frac{\frac{d^{2}y}{dt^{2}}}{\frac{d^{2}y}{dx^{2}}}=\frac{-\omega
^{2}y_{\max }\sin \left( kx-\omega t+\phi _{o}\right) }{-k^{2}y_{\max }\sin
(kx-\omega t+\phi _{o})}
\end{equation*}%
Of course the sine parts cancel, and so does $y_{\max }$. 
\begin{equation*}
\frac{\frac{d^{2}y}{dt^{2}}}{\frac{d^{2}y}{dx^{2}}}=\frac{-\omega ^{2}}{%
-k^{2}}
\end{equation*}%
and we know that 
\begin{equation*}
\frac{\omega }{k}=v
\end{equation*}%
the wave speed (not the speed of a piece of the medium, but the speed of the
wave)

\begin{equation*}
\frac{\frac{d^{2}y}{dt^{2}}}{\frac{d^{2}y}{dx^{2}}}=v^{2}
\end{equation*}%
If a function satisfies this equation, we can identify it as a wave. But
this form is awkward, let's rewrite it as 
\begin{equation*}
\frac{d^{2}y}{dt^{2}}=v^{2}\frac{d^{2}y}{dx^{2}}
\end{equation*}%
This is the \emph{linear wave equation}. There are waves that don't follow
this equation. They are called non-linear waves. But that is a subject for
an upper division course.

\section{Waves in two and three dimensions}

So far we have written expressions for waves, but our experience tells us
that waves don't usually come as one dimensional phenomena. In the next
figure, we see the disturbance (a drop) creating a water wave. \FRAME{dtbpFU%
}{2.3709in}{1.5389in}{0pt}{\Qcb{Picture of a water drop (Jon Paul Johnson,
used by permision)}}{}{Figure}{\special{language "Scientific Word";type
"GRAPHIC";maintain-aspect-ratio TRUE;display "USEDEF";valid_file "T";width
2.3709in;height 1.5389in;depth 0pt;original-width 2.3315in;original-height
1.5031in;cropleft "0";croptop "1";cropright "1";cropbottom "0";tempfilename
'Circular_wave_from_a_drip0.wmf';tempfile-properties "XNPR";}}You can see
the wave forming, but the wave is clearly not one dimensional. It appears
nearly circular. In fact, it is closer to hemispherical, and this limit is
only true because the disturbance is at the air-water boundary. Most waves
in a uniform medium will be roughly spherical. As such a wave travels away
from the source, the energy traveling gets more spread out. This causes the
amplitude to decrease. Think of a sound wave, it gets quieter the farther
you are from the source. We change our equation to account for this by
making the amplitude a function of the distance, $r,$ from the source%
\begin{equation}
y\left( r,t\right) =A\left( r\right) \sin \left( kr-\omega t+\phi _{o}\right)
\end{equation}%
Of course, if we look at a very large wave, but we only look at part of the
wave, we see that our part looks flatter as the wave expands. We re just
drawing the crests, and sometimes in such drawings we call the crests \emph{%
wave fronts} (think of wave crests at the beach as they hit you, the name
wave front doesn't seem to bad of a name). \FRAME{dhF}{2.4881in}{1.6189in}{%
0pt}{}{}{Figure}{\special{language "Scientific Word";type
"GRAPHIC";maintain-aspect-ratio TRUE;display "USEDEF";valid_file "T";width
2.4881in;height 1.6189in;depth 0pt;original-width 2.4474in;original-height
1.5826in;cropleft "0";croptop "1";cropright "1";cropbottom "0";tempfilename
'SUS73V23.wmf';tempfile-properties "XPR";}}Very far from the source, our
wave is flat enough that we can ignore the curvature across it's wave
fronts. We call such a wave a \emph{plane wave}. There are no true plane
waves in nature, but this idealization makes our mathematical solutions
simpler and many waves come close to this approximation.

We will soon see that sound is a longitudinal wave with a medium of air.
Really any solid, liquid, or gas will work as a medium for sound. For our
study, we will take sound to be a longitudinal wave and treat liquids and
gasses. Solids have additional forces involved due to the tight bonding of
the atoms, and therefore are more complicated. Technically in a solid sound
can be a transverse wave as well a longitudinal wave, but we usually call
transverse waves of this nature \emph{shear waves.}

We have said that energy travels in a wave. And from Principles of Physics I
we know the change of energy in an amount of time is called a power. So
waves must provide power. We will take up this topic in our next lecture.

%TCIMACRO{%
%\TeXButton{Basic Equations}{\vspace{0.5in}\hspace{-1.3in}{\LARGE Basic Equations\vspace{0.25in}}}}%
%BeginExpansion
\vspace{0.5in}\hspace{-1.3in}{\LARGE Basic Equations\vspace{0.25in}}%
%EndExpansion
Speed of waves on strings

\begin{equation*}
v=\sqrt{\frac{T_{s}}{\mu }}
\end{equation*}%
\begin{equation*}
\mu =\frac{m}{L}
\end{equation*}

\begin{equation}
y=y_{\max }\sin \left( kx-\omega t+\phi _{o}\right)
\end{equation}%
\begin{equation}
v_{y}=\frac{\partial y}{\partial t}=-\omega y_{\max }\cos \left( kx-\omega
t+\phi _{o}\right)
\end{equation}%
\begin{equation}
a_{y}=\frac{\partial v_{y}}{\partial t}=-\omega ^{2}y_{\max }\sin \left(
kx-\omega t+\phi _{o}\right)
\end{equation}%
\begin{equation*}
k=\frac{2\pi }{\lambda }
\end{equation*}%
\begin{equation*}
\omega =2\pi f
\end{equation*}%
\begin{equation}
v=\lambda f
\end{equation}

\begin{equation*}
\frac{d^{2}y}{dt^{2}}=v^{2}\frac{d^{2}y}{dx^{2}}
\end{equation*}

\chapter{5 Power and Superposition 6 1.16.4 1.16.5}

%TCIMACRO{%
%\TeXButton{\chaptermark{Superposition}}{\chaptermark{Superposition}}}%
%BeginExpansion
\chaptermark{Superposition}%
%EndExpansion

In Principles of Physics I we learned about power. Power is the rate of
energy delivery or use. Waves transfer energy from one place to another.
Often the reason we want a wave is to transfer energy! So we should look at
waves and energy delivery.

%TCIMACRO{%
%\TeXButton{Fundamental Concepts}{\vspace{0.10in}\hspace{-0.37in}{\Large Fundamental Concepts\vspace{0.2in}}}}%
%BeginExpansion
\vspace{0.10in}\hspace{-0.37in}{\Large Fundamental Concepts\vspace{0.2in}}%
%EndExpansion

\begin{itemize}
\item Because waves are three dimensional, describing the power or energy
delivered per time of the wave is not enough, We describe how spatially
spread out that power is. We call this spread power \emph{intensity}

\item Waves invert on reflection from \textquotedblleft fixed
ends\textquotedblright\ and don't invert from \textquotedblleft free
ends\textquotedblright

\item If we make more than one wave in a medium, the waves \textquotedblleft
add up\textquotedblright\ or \emph{superimpose.}

\item When two waves interfere we get a new wave with a more complicated
amplitude
\end{itemize}

\section{Energy and Power in waves}

We have said that waves don't transport mass, they transport energy. We used
the example of sound waves. If the air molecules move to another location ,
we call it wind. But if the air molecules just vibrate around an equilibrium
position, we called that sound waves. The molecules don't end up somewhere
else. But something does move.

Maybe a better example is a water wave in a swimming pool. Suppose someone
jumps into a calm pool. The person is a disturbance and the water is a
medium. We will get waves. But the water doesn't end up bunched together on
the other side of the pool. The water is still everywhere in the pool, even
around the person that jumped in. So what did move away from the person? It
is energy.

For sinusoidal waves, the amount of energy in the wave is related to both
it's amplitude and it's frequency. And the specifics of that relationship
depend on the type of wave. We will study both sound and light waves later
in our course and we will find that relationship between energy, amplitude,
and frequency are different for the two different kinds of waves.

Let's look at a specific case of a rope attached to a mechanical oscillator.
You can see the situation in the next figure. The oscillator (black thing)
has a piece that goes up and down (silver thing sticking out to the right).
The rope is attached to this silver oscillating piece. The oscillator acts
as a disturbance. A wave is formed in the rope. \FRAME{dtbpF}{2.9856in}{%
1.5939in}{0pt}{}{}{Figure}{\special{language "Scientific Word";type
"GRAPHIC";maintain-aspect-ratio TRUE;display "USEDEF";valid_file "T";width
2.9856in;height 1.5939in;depth 0pt;original-width 3.014in;original-height
1.5948in;cropleft "0";croptop "1";cropright "1";cropbottom "0";tempfilename
'TE0LO800.wmf';tempfile-properties "XPR";}}

Once again let's take a small part of our medium. Before we took a pice that
was $\Delta x$ long, but this time let's describe our piece of rope as a
small bit of mass $\Delta m.$ We already know about linear mass densities.%
\begin{equation*}
\mu =\frac{M}{L}
\end{equation*}%
and since our rope is uniform, the linear mass density is the same for each
part of the rope. So for our marked part 
\begin{equation*}
\mu =\frac{\Delta m}{\Delta x}
\end{equation*}%
Then for this part of the rope medium there will be a kinetic energy 
\begin{eqnarray*}
K &=&\frac{1}{2}mv^{2} \\
&=&\frac{1}{2}\Delta mv_{y}^{2}
\end{eqnarray*}%
where we realize that the piece of the medium will just go up and down for a
transverse wave, so we only have $v_{y}.$ Using our linear mass density we
can write 
\begin{equation*}
\Delta m=\mu \Delta x
\end{equation*}%
so the kinetic energy for just this little piece of the rope is 
\begin{equation*}
\Delta K=\frac{1}{2}\left( \mu \Delta x\right) v_{y}^{2}
\end{equation*}%
and from before we know that $v_{y}=-\omega A\cos \left( kx-\omega t+\phi
_{o}\right) $ so our kinetic energy is 
\begin{equation*}
\Delta K=\frac{1}{2}\left( \mu \Delta x\right) \left( -\omega A\cos \left(
kx-\omega t+\phi _{o}\right) \right) ^{2}
\end{equation*}%
we can clean this up a bit%
\begin{equation*}
\Delta K=\frac{1}{2}\mu \Delta x\omega ^{2}A^{2}\cos ^{2}\left( kx-\omega
t+\phi _{o}\right)
\end{equation*}%
Now let's simplify our situation by letting $\phi _{o}=0$ and let's take a
snap shot view with $t=0.$

\begin{equation*}
\Delta K=\frac{1}{2}\mu \Delta x\omega ^{2}A^{2}\cos ^{2}\left( kx\right)
\end{equation*}%
Then, if we take only the part of the energy associated with one very tiny
part of our rope medium, our $\Delta x$ becomes just $dx$. We don't have all
the energy of the wave, just a small part of it. So we can write this as 
\begin{equation*}
dK=\frac{1}{2}\mu \omega ^{2}A^{2}\cos ^{2}\left( kx\right) dx
\end{equation*}%
and we can integrate this up to find the energy of just one wavelength of
the wave.

\begin{equation*}
\int_{0}^{K_{\lambda }}dK=\int_{0}^{\lambda }\frac{1}{2}\mu \omega
^{2}A^{2}\cos ^{2}\left( kx\right) dx
\end{equation*}%
A lot of the parts of our equation are constants, let's take them out front.%
\begin{equation*}
\int_{0}^{K_{\lambda }}dK=\frac{1}{2}\mu \omega ^{2}A^{2}\int_{0}^{\lambda
}\cos ^{2}\left( kx\right) dx
\end{equation*}%
The left hand side of this equation is just $K_{\lambda }$%
\begin{equation*}
K_{\lambda }=\frac{1}{2}\mu \omega ^{2}A^{2}\int_{0}^{\lambda }\cos
^{2}\left( kx\right) dx
\end{equation*}%
but if you are like me, the right hand side is an integral that you have
conveniently forgotten how to do. But we can look up this integral in an
integral table. My table gave the form 
\begin{equation*}
\int \left( \cos ^{2}\left( ax\right) \right) dx=\frac{1}{2}x+\frac{1}{4a}%
\sin \left( 2ax\right)
\end{equation*}%
We have to match this to our equation. We see $a$ in the integral table
version is our $k$ so 
\begin{eqnarray*}
K_{\lambda } &=&\frac{1}{2}\mu \omega ^{2}A^{2}\left[ \frac{1}{2}x+\frac{1}{%
4k}\sin \left( 2kx\right) \right\vert _{0}^{\lambda } \\
&=&\frac{1}{2}\mu \omega ^{2}A^{2}\left[ \left\{ \frac{1}{2}\lambda +\frac{1%
}{4k}\sin \left( 2k\lambda \right) \right\} -\left\{ \frac{1}{2}\left(
0\right) +\frac{1}{4k}\sin \left( 2k\left( 0\right) \right) \right\} \right]
\\
&=&\frac{1}{2}\mu \omega ^{2}A^{2}\left[ \left\{ \frac{1}{2}\lambda +\frac{1%
}{4\frac{2\pi }{\lambda }}\sin \left( 2\frac{2\pi }{\lambda }\lambda \right)
\right\} -\frac{1}{4k}\sin \left( 2k\left( 0\right) \right) \right] \\
&=&\frac{1}{2}\mu \omega ^{2}A^{2}\left[ \left\{ \frac{1}{2}\lambda +\frac{%
\lambda }{8\pi }\sin \left( 4\pi \right) \right\} \right] \\
&=&\frac{1}{4}\mu \omega ^{2}A^{2}\lambda
\end{eqnarray*}

This wavelength sized piece of rope will also have some potential energy
because the fibers that make the rope are stretched to make the wave. They
aren't really quite like little springs, but it's not to far wrong to model
them as though they were. So let's take the spring potential energy 
\begin{equation*}
U=\frac{1}{2}k_{s}\left( y-y_{e}\right) ^{2}
\end{equation*}%
and let $y_{e}=0$ for our piece of rope. We know that our piece of rope will
go up and down in SHM. The natural frequency of our rope piece is 
\begin{equation*}
\omega =\sqrt{\frac{k_{s}}{m}}
\end{equation*}%
We can solve this for the spring constant%
\begin{equation*}
k_{s}=m\omega ^{2}
\end{equation*}%
or for our little piece of the rope 
\begin{equation*}
k_{s}=\Delta m\omega ^{2}
\end{equation*}%
Let's put this into our potential energy equation for just our little piece
of the rope%
\begin{equation*}
\Delta U=\frac{1}{2}\Delta m\omega ^{2}\left( y\right) ^{2}
\end{equation*}%
and substitute $\Delta m=\mu \Delta x$ again%
\begin{equation*}
\Delta U=\frac{1}{2}\mu \Delta x\omega ^{2}\left( y\right) ^{2}
\end{equation*}%
and of course, let our $\Delta x$ get small so we have a small bit of
potential energy from a small bit of rope.%
\begin{equation*}
dU=\frac{1}{2}\mu \omega ^{2}\left( y\right) ^{2}dx
\end{equation*}%
and integrate this to find the potential energy in the stretchy bonds of one
wavelength worth of rope. 
\begin{equation*}
\int_{0}^{U_{\lambda }}dU=\int_{0}^{\lambda }\frac{1}{2}\mu \omega
^{2}\left( y\right) ^{2}dx
\end{equation*}%
We need to put in our equation for the position for our wave.%
\begin{equation*}
y\left( x,t\right) =A\sin \left( kx-\omega t+\phi _{o}\right)
\end{equation*}%
But once again let's choose $\phi _{o}=0$ and $t=0$ so that%
\begin{equation*}
y\left( x,0\right) =A\sin \left( kx\right)
\end{equation*}%
then 
\begin{equation*}
\int_{0}^{U_{\lambda }}dU=\int_{0}^{\lambda }\frac{1}{2}\mu \omega
^{2}\left( A\sin \left( kx\right) \right) ^{2}dx
\end{equation*}%
The left hand side is easy. And there are, once again, constants 
\begin{equation*}
U_{\lambda }=\frac{1}{2}\mu \omega ^{2}A^{2}\int_{0}^{\lambda }\left( \sin
\left( kx\right) \right) ^{2}dx
\end{equation*}%
and we have another integral I\ don't remember how to do. But I still have
my table of integrals and my table tells me%
\begin{equation*}
\int \left( \sin ^{2}\left( ax\right) \right) dx=\frac{1}{2}x-\frac{1}{4a}%
\sin \left( 2ax\right)
\end{equation*}%
Which feels very familiar. Let's let $a=k$ again. Then 
\begin{eqnarray*}
U_{\lambda } &=&\frac{1}{2}\mu \omega ^{2}A^{2}\left[ \frac{1}{2}x-\frac{1}{%
4k}\sin \left( 2kx\right) \right\vert _{0}^{\lambda } \\
&&\frac{1}{2}\mu \omega ^{2}A^{2}\left[ \left\{ \frac{1}{2}\lambda -\frac{1}{%
4k}\sin \left( 2k\lambda \right) \right\} -\left\{ \frac{1}{2}0-\frac{1}{4k}%
\sin \left( 2k\left( 0\right) \right) \right\} \right] \\
&=&\frac{1}{4}\mu \omega ^{2}A^{2}\lambda
\end{eqnarray*}%
The total energy in the $\lambda $-sized rope piece would be 
\begin{eqnarray*}
E_{\lambda } &=&K_{\lambda }+U_{\lambda } \\
&=&\frac{1}{4}\mu \omega ^{2}A^{2}\lambda +\frac{1}{4}\mu \omega
^{2}A^{2}\lambda \\
&=&\frac{1}{2}\mu \omega ^{2}A^{2}\lambda
\end{eqnarray*}%
We know that energy is being transferred by the wave. It turns out that
energy transfers in all waves, whether it is a light or sound wave or any
other mechanical wave (thought the equation might be a little different). We
should wonder, how fast is energy transferring? This can mean the difference
between a warm ray of sun on a cool spring day and being burned by a laser
beam. We will start by considering the rate of energy transfer, \emph{power}%
. The concept of power should be familiar to us from Principles of Physics
I. We can find the power as the rate at which energy is transferred.%
\begin{equation*}
\mathcal{P}=\frac{\Delta E}{\Delta t}
\end{equation*}%
Since we picked the amount of energy in one wavelength, and we know the time
it takes for one wavelength of the wave to pass by is $T.$ Then the power is 
\begin{equation*}
P=\frac{\frac{1}{2}\mu \omega ^{2}A^{2}\lambda }{T}
\end{equation*}%
but remember 
\begin{equation*}
v=\frac{\lambda }{T}
\end{equation*}%
so we can write our power as 
\begin{equation*}
P=\frac{1}{2}\mu \omega ^{2}A^{2}v
\end{equation*}%
Power is important. Most detectors (like our ears and eyes) really detect
the power delivered by a wave. And we can see that the power is proportional
to the amplitude squared and the frequency squared for the wave. We will
learn more on this as we study light and sound.

\section{Power and Intensity}

What we have done works fine for linear waves. But if we consider our
spherical wave from a point source, we can see that this description isn't
good enough. In the next figure we have two professors. \FRAME{dtbpF}{%
2.8426in}{2.3488in}{0pt}{}{}{Figure}{\special{language "Scientific
Word";type "GRAPHIC";maintain-aspect-ratio TRUE;display "USEDEF";valid_file
"T";width 2.8426in;height 2.3488in;depth 0pt;original-width
2.8003in;original-height 2.3082in;cropleft "0";croptop "1";cropright
"1";cropbottom "0";tempfilename 'SV53CN0J.wmf';tempfile-properties "XPR";}}%
Thinking from our experience we would say that the sound will seem louder
for professor A. This is because the energy in the sound wave is being
spread over the surface of the wave, and that surface is getting bigger as
the wave moves outward. The energy is more spread out by the time it gets to
professor B.

%TCIMACRO{%
%\TeXButton{Tuning Fork Demo}{\marginpar {
%\hspace{-0.5in}
%\begin{minipage}[t]{1in}
%\small{Tuning Fork Demo}
%\end{minipage}
%}}}%
%BeginExpansion
\marginpar {
\hspace{-0.5in}
\begin{minipage}[t]{1in}
\small{Tuning Fork Demo}
\end{minipage}
}%
%EndExpansion
To be able to describe how much energy we get from our wave we need define
something new.%
\begin{equation}
\mathcal{I}\equiv \frac{\mathcal{P}}{A}
\end{equation}%
that is, the power divided by an area. But what does it mean?

Consider a point source. This could be a loud speaker, or a buzzer, or a
baby crying, etc.\FRAME{dtbpF}{2.4751in}{2.2485in}{0pt}{}{}{Figure}{\special%
{language "Scientific Word";type "GRAPHIC";maintain-aspect-ratio
TRUE;display "USEDEF";valid_file "T";width 2.4751in;height 2.2485in;depth
0pt;original-width 2.4336in;original-height 2.2087in;cropleft "0";croptop
"1";cropright "1";cropbottom "0";tempfilename
'SV53CN0K.wmf';tempfile-properties "XPR";}}The point source sends out waves
in all directions. The wave crests will define a sphere around the points
source. The figure shows a cross section but remember it is a wave from a
point source, so we are really drawing concentric spheres like balloons
inside of balloons. Then form for $I$ for our point source would be%
\begin{equation}
\mathcal{I}=\frac{\mathcal{P}}{4\pi r^{2}}
\end{equation}%
As the wave travels, its the power per unit area decreases with the square
of the distance because the area is getting larger.\FRAME{dhF}{2.655in}{%
1.8273in}{0pt}{}{}{Figure}{\special{language "Scientific Word";type
"GRAPHIC";maintain-aspect-ratio TRUE;display "USEDEF";valid_file "T";width
2.655in;height 1.8273in;depth 0pt;original-width 2.6135in;original-height
1.7902in;cropleft "0";croptop "1";cropright "1";cropbottom "0";tempfilename
'SV53CN0L.wmf';tempfile-properties "XPR";}}

This quantity that tells us how spread out our power has become is called
the \emph{intensity} of the wave. Professor A would agree with us that the
wave he heard was more intense than the wave heard by Professor B. That is
because the wave was less spread out for Professor A.

Suppose we cup our hand to our ear. What are we doing? We are increasing the
area of our ear. Our ears work by transferring the energy of the sound wave
to a mechanical-electro-chemical device that creates a nerve signal.%
\footnote{%
The inner hair cells in the organ of Corti in the cochlea.} The more energy,
the stronger the signal. If we are a distance $r$ away from the source of
the sound then the intensity is 
\begin{equation*}
\mathcal{I}=\frac{\mathcal{P}_{source}}{A_{wave}}
\end{equation*}%
But we are collecting the sound wave with another area, the area of our
hand. The power received is 
\begin{eqnarray*}
\mathcal{P}_{received} &=&\mathcal{I}A_{hand} \\
&=&\frac{A_{hand}}{A_{wave}}\mathcal{P}_{source}
\end{eqnarray*}%
And we can see that, indeed, the larger the hand, the more power, and
therefore more energy we collect. This is the idea behind a dish antenna for
communications and the idea behind the acoustic dish microphones we see at
sporting events. In next figure, we can see that it would take an
increasingly larger dish to maintain the same power gathering capability as
we get farther from the source.\FRAME{dhF}{3.0312in}{2.2131in}{0pt}{}{}{%
Figure}{\special{language "Scientific Word";type
"GRAPHIC";maintain-aspect-ratio TRUE;display "USEDEF";valid_file "T";width
3.0312in;height 2.2131in;depth 0pt;original-width 4.9355in;original-height
3.595in;cropleft "0";croptop "1";cropright "1";cropbottom "0";tempfilename
'SV53CN0M.wmf';tempfile-properties "XPR";}}

\section{Reflection and Transmission}

In our examples so far, our waves could go on forever. But in our house or
church building or classroom this is not usually the case. The waves will
eventually hit something.

In class we will made pulses on a spring with one end of the rope fixed
(held by a class member). What happened when the pulse reached the end of
the rope?

\FRAME{dhF}{3.6063in}{4.0456in}{0pt}{}{}{Figure}{\special{language
"Scientific Word";type "GRAPHIC";maintain-aspect-ratio TRUE;display
"USEDEF";valid_file "T";width 3.6063in;height 4.0456in;depth
0pt;original-width 3.5587in;original-height 3.9963in;cropleft "0";croptop
"1";cropright "1";cropbottom "0";tempfilename
'SV6XQL00.wmf';tempfile-properties "XPR";}}

\subsection{Case I: Fixed rope end.}

There is a big change in the medium at the end of the rope. The rope ends.
There is a person or (as in the next figure) some thing holding the rope in
place. This change in medium causes a reflection.

In the fixed end case, the pulse is inverted (seems to flip upside down).
Why?

In the next figure I\ have envisioned a rope made of small rope segments.

\FRAME{dtbpF}{3.0519in}{1.7711in}{0pt}{}{}{Figure}{\special{language
"Scientific Word";type "GRAPHIC";maintain-aspect-ratio TRUE;display
"USEDEF";valid_file "T";width 3.0519in;height 1.7711in;depth
0pt;original-width 9.6911in;original-height 5.6109in;cropleft "0";croptop
"1";cropright "1";cropbottom "0";tempfilename
'SV6XQL01.wmf';tempfile-properties "XPR";}}The end of the rope pushes up on
the support (person in class, or nail in the figure). By Newton's third law
there must be a normal force downward on the rope end. Compressing the
molecules in the support structure stores potential energy in those
compressed atoms. They will release that potential energy and create kinetic
energy in the rope end. That kinetic energy will have a rope-end velocity 
\begin{equation*}
K=\frac{1}{2}mv_{y}^{2}
\end{equation*}%
and that velocity will be downward. This will pull the rope down, inverting
the wave.

But what happens if the rope end is not fixed?

The rope end rises, and therefore there is no force exerted. The pulse (or
at least part of the pulse energy) is still reflected, but there is no
inversion because there was no downward force or no stored potential energy
in the support!\FRAME{dtbpF}{2.2771in}{1.3958in}{0pt}{}{}{Figure}{\special%
{language "Scientific Word";type "GRAPHIC";maintain-aspect-ratio
TRUE;display "USEDEF";valid_file "T";width 2.2771in;height 1.3958in;depth
0pt;original-width 9.1687in;original-height 5.6109in;cropleft "0";croptop
"1";cropright "1";cropbottom "0";tempfilename
'SV6XQL02.wmf';tempfile-properties "XPR";}}

\subsection{Case III: Partially attached rope end}

Now lets tie the rope to another rope that is larger, more dense, than the
rope we have been using, what will happen?

The light end of the rope exerts a force on the heavy beginning of the new
rope. The atoms in the heavy rope will be pulled and compressed. The heavy
rope will move, but the heaviness of the rope will prevent it from going
very far. The fibers on the end of the heavy rope will build up potential
energy because their bonds will be stretched. They will push downward on the
light end of the rope opposing the stretch. Once again the potential energy
from the stretched heavy rope atoms will transfer to kinetic energy in the
light rope end. Once again that light rope end will move downward. \FRAME{%
dtbpF}{2.6532in}{1.4857in}{0pt}{}{}{Figure}{\special{language "Scientific
Word";type "GRAPHIC";maintain-aspect-ratio TRUE;display "USEDEF";valid_file
"T";width 2.6532in;height 1.4857in;depth 0pt;original-width
3.8in;original-height 2.1162in;cropleft "0";croptop "1";cropright
"1";cropbottom "0";tempfilename 'SV6XQL03.wmf';tempfile-properties "XPR";}}

We expect part of the energy to reflect back along the light rope, and this
pulse will be inverted. Notice that the heavy rope did move upward a little,
so there will also be a pulse on the heavy rope. Since this pulse formed
from the light rope pulling up the heavy rope, and the light rope atoms were
not stretched much, this pulse will not be inverted. \FRAME{dtbpF}{2.8764in}{%
2.2572in}{0pt}{}{}{Figure}{\special{language "Scientific Word";type
"GRAPHIC";maintain-aspect-ratio TRUE;display "USEDEF";valid_file "T";width
2.8764in;height 2.2572in;depth 0pt;original-width 2.8331in;original-height
2.2174in;cropleft "0";croptop "1";cropright "1";cropbottom "0";tempfilename
'SV6XQL04.wmf';tempfile-properties "XPR";}}

This is all great for pulses. But suppose a sinusoidal wave approaches the
boundary. We can envision the crests like pulses, and we expect the first
crest to slow down when it reaches the boundary, letting the other crests
catch up. This is because 
\begin{equation*}
v=\sqrt{\frac{T}{\mu }}
\end{equation*}%
and $\mu $ got bigger so $v$ got smaller. Once the wave passes the boundary,
the crests will be closer together. The wavelength changes as we move to the
slower medium.

But does the frequency change? We know that 
\begin{equation*}
v=\lambda f
\end{equation*}%
so%
\begin{equation*}
f=\frac{v}{\lambda }
\end{equation*}%
Both the speed and the wavelength have changed, but did they change
proportionately so $f$ is constant? This must be so. Think that the change
in wavelength is due to the relative speed of the wave in the two media. If $%
\Delta v$ is small the change in $\lambda $ will be small because the crests
are not delayed too long. If $\Delta v$ is large, the crests are delayed by
a large amount and so the change in $\lambda $ is large. We won't derive the
fact that $f$ is constant, but we can see that is is very believable that it
is true.\FRAME{dhF}{2.2727in}{1.2644in}{0pt}{}{}{Figure}{\special{language
"Scientific Word";type "GRAPHIC";maintain-aspect-ratio TRUE;display
"USEDEF";valid_file "T";width 2.2727in;height 1.2644in;depth
0pt;original-width 3.4618in;original-height 1.9147in;cropleft "0";croptop
"1";cropright "1";cropbottom "0";tempfilename
'SV8THK0F.wmf';tempfile-properties "XPR";}}

This is true for all waves, even light. When a wave crosses a boundary from
a fast to a slow or a slow to a fast medium, $\lambda $ will change and $f$
will remain constant.

\section{Superposition Principle}

What happens if we have more than one wave propagating in a medium? If you
remember being a little child in a bath tub, you will probably remember
making waves in the water. If you made a wave with each hand, the two waves
seemed to \textquotedblleft pile up\textquotedblright\ in the middle and
make a big splash. We should expect something like this for any kind of
wave. We call the \textquotedblleft piling up\textquotedblright\ of waves 
\emph{superposition.} The word literally means putting one wave on top of
another.

\begin{center}
\rule{300pt}{1pt}

\textbf{Superposition}: If two or more traveling waves are moving through
the same medium, the resultant wave formed at any point is the algebraic sum
of the values of the individual wave forms.

\rule{300pt}{1pt}
\end{center}

So if we have 
\begin{equation}
y_{1}\left( x,t\right)
\end{equation}%
and 
\begin{equation}
y_{2}\left( x,t\right)
\end{equation}%
both propagating on a string, then we would see%
\begin{equation}
y_{r}\left( x,t\right) =y_{1}\left( x,t\right) +y_{2}\left( x,t\right)
\end{equation}

This is a fantastically simple way for the universe to act!

Let's look at an example. let's add the top wave (red) to the middle
wave(green). We get the bottom wave (purple)\FRAME{dhF}{2.77in}{3.3788in}{0pt%
}{}{}{Figure}{\special{language "Scientific Word";type
"GRAPHIC";maintain-aspect-ratio TRUE;display "USEDEF";valid_file "T";width
2.77in;height 3.3788in;depth 0pt;original-width 7.2705in;original-height
8.8816in;cropleft "0";croptop "1";cropright "1";cropbottom "0";tempfilename
'SV53CN0N.wmf';tempfile-properties "XPR";}}Of course we are adding these in
the snapshot view. So this is all done for just one instant of time.

Let's see how to do this. \FRAME{dhF}{2.7121in}{2.1715in}{0pt}{}{}{Figure}{%
\special{language "Scientific Word";type "GRAPHIC";maintain-aspect-ratio
TRUE;display "USEDEF";valid_file "T";width 2.7121in;height 2.1715in;depth
0pt;original-width 6.576in;original-height 5.2598in;cropleft "0";croptop
"1";cropright "1";cropbottom "0";tempfilename
'SV53CN0O.wmf';tempfile-properties "XPR";}}Start at $x=-2.$ In the figure, I
drew a red bar to show the $y$ value at $x=-2$ for the red curve. Likewise,
I have a green bar sowing the value of $y$ at $x=-2$ for the green wave.
Note that this is negative. On the bottom graph, the bars have been
repeated, and we can see that the red bar minus the green bar brings us to
the value for the resulting wave at the point $x=-2.$ We need to do this at
every point along all the waves for this instant of time.

This is tedious by hand, so we won't generally do this calculation by hand.
But a computer can do it easily.

Note that this is really only true for \emph{linear} systems. Let's take the
example of a slinky. If we form two waves in the slinky, they behave
according to the superposition principle most of the time. But suppose we
make the amplitude of the individual waves large. They may travel
individually OK, but when the amplitudes add we may overstretch the slinky.
Then it would never return to it's original shape. The wave form would have
to change. Such a wave is not linear. There is a good rule of thumb for when
waves are linear.

\begin{center}
\rule{300pt}{1pt}

A wave is generally linear when its amplitude is much smaller than its
wavelength.

\rule{300pt}{1pt}
\end{center}

\section{Consequences of superposition}

Linear waves traveling in media can pass through each other without being
destroyed or altered!

\FRAME{dhFU}{1.7633in}{2.0817in}{0pt}{\Qcb{{\protect\small Constructive
Interference (Public Domain image by Inductiveload,
http://commons.wikimedia.org/wiki/File:Constructive\_interference.svg)}}}{}{%
Figure}{\special{language "Scientific Word";type
"GRAPHIC";maintain-aspect-ratio TRUE;display "USEDEF";valid_file "T";width
1.7633in;height 2.0817in;depth 0pt;original-width 1.7678in;original-height
2.0924in;cropleft "0";croptop "1";cropright "1";cropbottom "0";tempfilename
'SV53CN0P.wmf';tempfile-properties "XPR";}}

Our wave on the string makes the string segments move in the $y$ direction.
Both waves do this. So putting the two waves together just makes the string
segments move more! There is a special name for what we observe

\begin{center}
\rule{300pt}{1pt}

\emph{interference}: The combination of separate waves in the same region of
space to produce a resultant wave.

\rule{300pt}{1pt}
\end{center}

What happens if one of the pulses is inverted?\FRAME{dhFU}{1.7828in}{2.2033in%
}{0pt}{\Qcb{{\protect\small Destructive Interference (Public Domain image by
Inductiveload,
http://commons.wikimedia.org/wiki/File:Destructive\_interference1.svg)}}}{}{%
Figure}{\special{language "Scientific Word";type
"GRAPHIC";maintain-aspect-ratio TRUE;display "USEDEF";valid_file "T";width
1.7828in;height 2.2033in;depth 0pt;original-width 1.5833in;original-height
1.9638in;cropleft "0";croptop "1";cropright "1";cropbottom "0";tempfilename
'SV53CN0Q.wmf';tempfile-properties "XPR";}}

When the two pulses meet, they \textquotedblleft cancel each other
out.\textquotedblright\ But do they go away? No! the energy is still there,
the string segment motions have just summed vectorially to zero, the energy
carried by each wave is still there hiding in stretched bonds. We have a few
more definitions. The type of interference we have just seen is the first

\begin{enumerate}
\item \emph{Destructive Interference}:\ interference between waves when the
displacements caused by the two waves are opposite in direction

\item \emph{Constructive Interference}: interference between waves when the
displacements caused by the two waves are in the same direction
\end{enumerate}

The combination of waves is important for both scientists and engineers. In
engineering this is the hart of vibrometry. \FRAME{dhFU}{3.5639in}{2.3851in}{%
0pt}{\Qcb{{\protect\small Marshall and Cal Poly technicians wired the
NanoSail-D spacecraft to accelerometers, instruments which measure vibration
response during simulated launch conditions. Image couracy NASA, image in
the Public Domain.}}}{}{Figure}{\special{language "Scientific Word";type
"GRAPHIC";maintain-aspect-ratio TRUE;display "USEDEF";valid_file "T";width
3.5639in;height 2.3851in;depth 0pt;original-width 3.5172in;original-height
2.3454in;cropleft "0";croptop "1";cropright "1";cropbottom "0";tempfilename
'SV53CN0R.wmf';tempfile-properties "XPR";}}Mechanical systems have moving
parts. These moving parts can be the disturbance that creates a wave. If
more than one wave crest arrives at a location in the device, the amplitude
at that location could become large. The oscillation of this part of the
device could rattle apart welds or bolts, destroying the device. Later, as
we study spectroscopy, we will see how to diagnose such a problem and hint
at how to correct it.

\section{Single Frequency Interference in one dimension}

We have seen interference from superimposing two waves in the same medium.
Whether we get constructive or destructive interference 
%TCIMACRO{%
%\TeXButton{Two speaker demo}{\marginpar {
%\hspace{-0.5in}
%\begin{minipage}[t]{1in}
%\small{Two speaker demo}
%\end{minipage}
%}}}%
%BeginExpansion
\marginpar {
\hspace{-0.5in}
\begin{minipage}[t]{1in}
\small{Two speaker demo}
\end{minipage}
}%
%EndExpansion
really depends on the total phase difference. Let's review this, but let's
make our two waves as general as possible.%
\begin{eqnarray*}
y_{1} &=&y_{\max }\sin \left( k_{1}x_{1}-\omega _{1}t_{1}+\phi _{1}\right) \\
y_{2} &=&y_{\max }\sin \left( k_{2}x_{2}-\omega _{2}t_{2}+\phi _{2}\right)
\end{eqnarray*}%
where we have allowed the wavelengths to be different (so then the wave
numbers, $k_{1}$ and $k_{2}$ could be different), the distances traveled by
the waves could be different ($x_{2}$ might not be equal to $x_{1}$), the
frequencies might be different, the times observed might be different, and
the phase constants might be different. We have not made the amplitudes
different (but it is likely to happen in an upcoming homework problem).

Then the resulting wave would be%
\begin{equation*}
y_{r}=y_{\max }\sin \left( k_{2}x_{2}-\omega _{2}t_{2}+\phi _{2}\right)
+A\sin \left( k_{1}x_{1}-\omega _{1}t_{1}+\phi _{1}\right)
\end{equation*}%
but this doesn't tell us much. We can understand much more about what our
combined waves would look like if we use a trig identity%
\begin{equation*}
\sin a+\sin b=2\cos \left( \frac{a-b}{2}\right) \sin \left( \frac{a+b}{2}%
\right)
\end{equation*}%
Let's try this. We could make $a=$ $k_{2}x_{2}-\omega _{2}t_{2}+\phi _{2}$
and $b=k_{1}x_{1}-\omega _{1}t_{1}+\phi _{1}$ then 
\begin{eqnarray*}
&=&2y_{\max }\cos \left( \frac{\left( k_{2}x_{2}-\omega _{2}t_{2}+\phi
_{2}\right) -\left( k_{1}x_{1}-\omega _{1}t_{1}+\phi _{1}\right) }{2}\right)
\\
&&\times \sin \left( \frac{\left( k_{2}x_{2}-\omega _{2}t_{2}+\phi
_{2}\right) +\left( k_{1}x_{1}-\omega _{1}t_{1}+\phi _{1}\right) }{2}\right)
\end{eqnarray*}%
We can rewire this as 
\begin{eqnarray*}
y_{r} &=&2y_{\max }\cos \left( \frac{1}{2}\left[ \left( k_{2}x_{2}-\omega
_{2}t_{2}+\phi _{2}\right) -\left( k_{1}x_{1}-\omega _{1}t_{1}+\phi
_{1}\right) \right] \right) \\
&&\times \sin \left( \frac{\left( k_{2}x_{2}-\omega _{2}t_{2}+\phi
_{2}\right) +\left( k_{1}x_{1}-\omega _{1}t_{1}+\phi _{1}\right) }{2}\right)
\end{eqnarray*}%
and recognize that the sine part is still a function of position and time
(just a complicated one) so it must be a wave. The cosine part must be part
of the amplitude. Let's write out this amplitude%
\begin{equation*}
A=2y_{\max }\cos \left( \frac{1}{2}\left[ \left( k_{2}x_{2}-\omega
_{2}t_{2}+\phi _{2}\right) -\left( k_{1}x_{1}-\omega _{1}t_{1}+\phi
_{1}\right) \right] \right)
\end{equation*}%
The part in square brackets is what we call the phase difference 
\begin{equation*}
\Delta \phi =\left( k_{2}x_{2}-\omega _{2}t_{2}+\phi _{2}\right) -\left(
k_{1}x_{1}-\omega _{1}t_{1}+\phi _{1}\right)
\end{equation*}%
Notice that to make constructive interference we want the amplitude to be as
big as possible. This happens when the cosine is $\pm 1.$ We can find when
this happens by recalling that cosine is $\pm 1$ when the argument of the
cosine is a multiple of $\pi .$ Notice that the amplitude has the cosine of
half the phase difference%
\begin{equation*}
A=2y_{\max }\cos \left( \frac{1}{2}\Delta \phi \right)
\end{equation*}%
so 
\begin{equation*}
\frac{1}{2}\Delta \phi =\pi m\qquad m=0,\pm 1,\pm 2,\pm 3,\cdots
\end{equation*}%
or 
\begin{equation*}
\Delta \phi =2\pi m\qquad m=0,\pm 1,\pm 2,\pm 3,\cdots \qquad \text{%
Constructive}
\end{equation*}%
For total destructive interference we want the amplitude to be zero. To
achieve this the cosine must be zero and this happens at odd multiples of $%
\pi /2.$ 
\begin{equation*}
\frac{1}{2}\Delta \phi =\frac{\pi }{2}m\qquad m=\pm 1,\pm 3,\pm 5,\cdots
\end{equation*}%
Then we could say that our phase difference 
\begin{equation*}
\Delta \phi =\pi m\qquad m=\pm 1,\pm 3,\pm 5,\cdots
\end{equation*}%
Another way to say this is to allow $m$ to be any integer but to write the
destructive interference condition as 
\begin{equation*}
\Delta \phi =\left( m+\frac{1}{2}\right) 2\pi \qquad m=0,\pm 1,\pm 2,\pm
3,\cdots \qquad \text{Destructive}
\end{equation*}

Let's write these out in detail%
\begin{eqnarray*}
\left[ \left( k_{2}x_{2}-\omega _{2}t_{2}+\phi _{2}\right) -\left(
k_{1}x_{1}-\omega _{1}t_{1}+\phi _{1}\right) \right] &=&2\pi m\qquad m=0,\pm
1,\pm 2,\pm 3,\cdots \\
\left[ \left( k_{2}x_{2}-\omega _{2}t_{2}+\phi _{2}\right) -\left(
k_{1}x_{1}-\omega _{1}t_{1}+\phi _{1}\right) \right] &=&\left( m+\frac{1}{2}%
\right) 2\pi \qquad m=0,\pm 1,\pm 2,\pm 3,\cdots
\end{eqnarray*}

We can see that there are at least four sources of phase difference here. We
can change the phase difference by changing wavelength ($k_{2}\neq k_{1}$),
changing frequency ($\omega _{2}\neq \omega _{1}$), having the waves travel
different distances ($x_{2}\neq x_{1}$), having different phase constants ($%
\phi _{2}\neq \phi _{1}$) or even consider somehow mixing the waves at
different times ($t_{2}\neq t_{1}$). This last one is not as interesting, we
want our waves mixed together at the same time so let's set%
\begin{equation*}
t_{2}=t_{1}=t
\end{equation*}%
but let's let any of the other variables be changeable. This gives us many
ways to produce constructive or destructive interference. In our next
lecture we will use this analysis to solve specific problems. The pattern
will be to use our two criteria%
\begin{eqnarray*}
\Delta \phi &=&2\pi m\qquad m=0,\pm 1,\pm 2,\pm 3,\cdots \qquad \text{%
Constructive} \\
\Delta \phi &=&\left( m+\frac{1}{2}\right) 2\pi \qquad m=0,\pm 1,\pm 2,\pm
3,\cdots \qquad \text{Destructive}
\end{eqnarray*}%
to determine if we have constructive, destructive, or partially constructive
or destructive interference.

If both waves have the same frequency, we can simplify this.%
\begin{eqnarray*}
\Delta \phi &=&\left[ \left( k_{2}x_{2}-\omega t_{2}+\phi _{2}\right)
-\left( k_{1}x_{1}-\omega t_{1}+\phi _{1}\right) \right] =2\pi m\qquad
m=0,\pm 1,\pm 2,\pm 3,\cdots \\
\Delta \phi &=&\left[ \left( k_{2}x_{2}-\omega t_{2}+\phi _{2}\right)
-\left( k_{1}x_{1}-\omega t_{1}+\phi _{1}\right) \right] =\left( m+\frac{1}{2%
}\right) 2\pi \qquad m=0,\pm 1,\pm 2,\pm 3,\cdots
\end{eqnarray*}%
becomes 
\begin{eqnarray*}
\Delta \phi &=&\left( k_{2}x_{2}-k_{1}x_{1}\right) +\left( \phi _{2}-\phi
_{1}\right) =2\pi m\qquad m=0,\pm 1,\pm 2,\pm 3,\cdots \\
\Delta \phi &=&\left( k_{2}x_{2}-k_{1}x_{1}\right) +\left( \phi _{2}-\phi
_{1}\right) =\left( m+\frac{1}{2}\right) 2\pi \qquad m=0,\pm 1,\pm 2,\pm
3,\cdots
\end{eqnarray*}%
and our wave equation becomes

\begin{eqnarray*}
y_{r} &=&2y_{\max }\cos \left( \frac{1}{2}\left[ \left(
k_{2}x_{2}-k_{1}x_{1}\right) +\left( \phi _{2}-\phi _{1}\right) \right]
\right) \\
&&\times \sin \left( \frac{k_{2}x_{2}+k_{1}x_{1}}{2}-\frac{\left( \omega
t_{2}+\omega _{1}t_{1}\right) }{2}+\frac{\left( \phi _{2}+\phi _{1}\right) }{%
2}\right)
\end{eqnarray*}

\section{Single frequency interference in more than one dimension}

But what happens if our waves don't travel along the same line? Suppose you
are at a dance, and there are two speakers. Further suppose that you are
testing the system with a constant tone so $\omega _{2}=\omega _{1}$ (either
that or you have somewhat boring music with long, sustained tones). Suppose
the two speakers make waves in phase ($\phi _{1}=\phi _{2}$). There is no
slower medium, so we expect $k_{2}=k_{1}.$ If you are equal distance from
the two speakers, you would expect constructive interference because both $%
\Delta \phi _{o}=0$ and $\Delta x=0$ for this case.%
\begin{eqnarray*}
\Delta \phi &=&\left( k_{2}x_{2}-\omega _{2}t_{2}+\phi _{2}\right) -\left(
k_{1}x_{1}-\omega _{1}t_{1}+\phi _{1}\right) \\
&=&\left( kx_{2}-\omega t+\phi \right) -\left( kx_{1}-\omega t+\phi \right)
\\
&=&k\left( \Delta x\right)
\end{eqnarray*}%
If we pick a spot where $x_{1}=x_{2}$ we would expect 
\begin{equation*}
\Delta \phi =0
\end{equation*}%
and since $0$ is a multiple of $2\pi $ in our criteria, then we expect this
will give us constructive interference.\FRAME{dhF}{1.4218in}{1.9527in}{0pt}{%
}{}{Figure}{\special{language "Scientific Word";type
"GRAPHIC";maintain-aspect-ratio TRUE;display "USEDEF";valid_file "T";width
1.4218in;height 1.9527in;depth 0pt;original-width 3.0441in;original-height
4.1909in;cropleft "0";croptop "1";cropright "1";cropbottom "0";tempfilename
'SVA69G06.wmf';tempfile-properties "XPR";}}But there are more places where
we expect constructive interference, because we know the sound wave is
really spherical. Any time the path difference, $\Delta x=n\lambda $ where $%
n $ is an integer, then 
\begin{equation*}
\Delta \phi =\frac{2\pi }{\lambda }\left( n\lambda \right) =n2\pi
\end{equation*}%
and we will have constructive interference The next figure shows an example
where the path difference is one wavelength.

\FRAME{dhF}{1.5489in}{2.1318in}{0pt}{}{}{Figure}{\special{language
"Scientific Word";type "GRAPHIC";maintain-aspect-ratio TRUE;display
"USEDEF";valid_file "T";width 1.5489in;height 2.1318in;depth
0pt;original-width 3.8233in;original-height 5.2736in;cropleft "0";croptop
"1";cropright "1";cropbottom "0";tempfilename
'SVA69G07.wmf';tempfile-properties "XPR";}}But any of the spots in the next
figure will experience constructive interference. Note the loud spots are
where there are two crests or two troughs together.\FRAME{dhF}{1.6181in}{%
2.1231in}{0pt}{}{}{Figure}{\special{language "Scientific Word";type
"GRAPHIC";maintain-aspect-ratio TRUE;display "USEDEF";valid_file "T";width
1.6181in;height 2.1231in;depth 0pt;original-width 3.7256in;original-height
4.8983in;cropleft "0";croptop "1";cropright "1";cropbottom "0";tempfilename
'SVA69G08.wmf';tempfile-properties "XPR";}}We also expect to see destructive
interference. This should occur where path differences are multiples of $%
\lambda /2$ so that 
\begin{eqnarray*}
\Delta \phi &=&\frac{2\pi }{\lambda }\Delta x \\
&=&\frac{2\pi }{\lambda }\left( n\frac{\lambda }{2}\right) \\
&=&n\pi
\end{eqnarray*}%
The next figure shows a case where $\Delta x=\lambda /2$\FRAME{dhF}{1.6933in%
}{2.0401in}{0pt}{}{}{Figure}{\special{language "Scientific Word";type
"GRAPHIC";maintain-aspect-ratio TRUE;display "USEDEF";valid_file "T";width
1.6933in;height 2.0401in;depth 0pt;original-width 1.9882in;original-height
2.4007in;cropleft "0";croptop "1";cropright "1";cropbottom "0";tempfilename
'SVA69G09.wmf';tempfile-properties "XPR";}}and the next figure shows many
places where there will be destructive interference because the two waves
are out of phase by half a wavelength.\FRAME{dhF}{1.7141in}{1.9969in}{0pt}{}{%
}{Figure}{\special{language "Scientific Word";type
"GRAPHIC";maintain-aspect-ratio TRUE;display "USEDEF";valid_file "T";width
1.7141in;height 1.9969in;depth 0pt;original-width 2.5581in;original-height
2.9836in;cropleft "0";croptop "1";cropright "1";cropbottom "0";tempfilename
'SVA69G0A.wmf';tempfile-properties "XPR";}}

When you moved from one dimension to two dimensions in Principles of Physics
I, we changed from the variables $x$ and $y$ to the variable $r$ where 
\begin{equation*}
r=\sqrt{x^{2}+y^{2}}
\end{equation*}%
Thus for this case our phase becomes 
\begin{equation*}
\Delta \phi =\left( 2\pi \frac{\Delta r}{\lambda }+\Delta \phi _{o}\right)
\end{equation*}

In our dance example, suppose we have speakers that are $4\unit{m}$ apart
and we are standing $3\unit{m}$ directly in front of one of the speakers.
Further suppose that we play an $A$ just above middle $C$ which has a
frequency of $440\unit{Hz}.$ The speed of sound is $343\unit{m}/\unit{s}$.
Our speakers are connected to the same stereo with equal length wires. What
is the phase difference at this spot?

\FRAME{dhF}{1.8609in}{1.9398in}{0pt}{}{}{Figure}{\special{language
"Scientific Word";type "GRAPHIC";maintain-aspect-ratio TRUE;display
"USEDEF";valid_file "T";width 1.8609in;height 1.9398in;depth
0pt;original-width 3.1142in;original-height 3.2474in;cropleft "0";croptop
"1";cropright "1";cropbottom "0";tempfilename
'SVA69G0B.wmf';tempfile-properties "XPR";}}From the geometry we can tell
that the path from the second speaker must be $5\unit{m}.$ So 
\begin{eqnarray*}
\Delta x &=&5\unit{m}-3\unit{m} \\
&=&2\unit{m}
\end{eqnarray*}%
We can tell that the wavelength is%
\begin{eqnarray*}
\lambda &=&\frac{v}{f} \\
&=&\frac{343\unit{m}/\unit{s}}{440\unit{Hz}} \\
&=&\allowbreak 0.779\,55\unit{m}
\end{eqnarray*}%
Since the speakers are connected to the same stereo with equal length wires, 
$\Delta \phi _{o}=0.$ Then 
\begin{eqnarray*}
\Delta \phi &=&\frac{2\pi }{\lambda }\Delta x+\Delta \phi _{o} \\
&=&\frac{2\pi }{\allowbreak 0.779\,55\unit{m}}\left( 2\unit{m}\right) +0 \\
&=&\allowbreak 5.\,\allowbreak 131\,2\pi \\
&=&2\pi +\allowbreak 3.\,\allowbreak 131\,2\pi
\end{eqnarray*}%
$\allowbreak \allowbreak $We should ask, is this constructive or destructive
interference? Well, it is neither purely constructive interference nor total
destructive interference. Our amplitude would be 
\begin{equation*}
2A\cos \left( \frac{1}{2}\left( \frac{2\pi }{\lambda }\Delta x+\Delta \phi
_{o}\right) \right)
\end{equation*}%
so in this case we get%
\begin{equation*}
2A\cos \left( \frac{1}{2}\left( 2\pi +\allowbreak 3.\,\allowbreak 131\,2\pi
\right) \right) =-0.409\,27A
\end{equation*}%
which is smaller (in magnitude) than $A,$ so it is partial destructive
interference. It would be quieter at this spot than if we just had one
speaker playing.

You might guess that this sort of analysis plays a large part in design of
concert halls. It also is important in mechanical designs. But you should
have seen a deficit in what we have learned so far. Up to this point, we
have only mixed waves that have the same frequency. Can we mix waves that
have different frequencies?

%TCIMACRO{%
%\TeXButton{Basic Equations}{\vspace{0.5in}\hspace{-1.3in}{\LARGE Basic Equations\vspace{0.25in}}}}%
%BeginExpansion
\vspace{0.5in}\hspace{-1.3in}{\LARGE Basic Equations\vspace{0.25in}}%
%EndExpansion
\begin{equation*}
U=\frac{1}{2}\mu \omega ^{2}\left( A\sin \left( kx\right) \right) ^{2}
\end{equation*}

\begin{equation*}
\mu =\frac{M}{L}
\end{equation*}%
\begin{equation*}
K=\frac{1}{2}\mu \Delta x\omega ^{2}A^{2}\cos ^{2}\left( kx\right)
\end{equation*}%
\begin{equation*}
P=\frac{1}{2}\mu \omega ^{2}A^{2}v
\end{equation*}%
\begin{equation*}
\mathcal{I}=\frac{\mathcal{P}_{source}}{A_{wave}}
\end{equation*}
\begin{eqnarray*}
\mathcal{P}_{received} &=&\mathcal{I}A_{hand} \\
&=&\frac{A_{hand}}{A_{wave}}\mathcal{P}_{source}
\end{eqnarray*}

\begin{eqnarray*}
\Delta \phi &=&2\pi m\qquad m=0,\pm 1,\pm 2,\pm 3,\cdots \qquad \text{%
Constructive} \\
\Delta \phi &=&\left( m+\frac{1}{2}\right) 2\pi \qquad m=0,\pm 1,\pm 2,\pm
3,\cdots \qquad \text{Destructive}
\end{eqnarray*}

\chapter{6 Standing Waves and Sound Waves 1.16.6, 1.17.1}

So far we have only hinted about sound waves. In this lecture we will study
them in more detail. But first let's consider what happens if we have two
waves with the same frequency and amplitude that superimpose in the same
medium but go the opposite direction. We expect constructive and destructive
interference! And we get it. And it turns out that this situation is very
important for making musical instruments - and for keeping structures from
being destroyed!

%TCIMACRO{%
%\TeXButton{\chaptermark{Standing Waves}}{\chaptermark{Standing Waves}}}%
%BeginExpansion
\chaptermark{Standing Waves}%
%EndExpansion

%TCIMACRO{%
%\TeXButton{Fundamental Concepts}{\vspace{0.10in}\hspace{-0.37in}{\Large Fundamental Concepts\vspace{0.2in}}}}%
%BeginExpansion
\vspace{0.10in}\hspace{-0.37in}{\Large Fundamental Concepts\vspace{0.2in}}%
%EndExpansion

\begin{itemize}
\item Waves traveling different directions can create static patterns of
constructive and destructive interference.

\item These patterns are called \textquotedblleft standing
waves\textquotedblright

\item The frequencies that make the standing wave patterns are quantized

\item Sound is a wave

\item We can express sound waves in terms displacement of air molecules or
as pressure changes
\end{itemize}

\section{Mathematical Description of Superposition}

We know what superposition is, but we don't really want to add values for
millions of points in a medium to find out what a combination of waves will
look like. At the very least, we want to make a computer do that (and
programs like OpenFoam do something very akin to this!). But where we can,
we would like to combine wave functions algebraically. We did this generally
last lecture. Let's do it again, but with more restrictions.

Lets define two wave functions%
\begin{equation*}
y_{1}=y_{\max }\sin \left( kx-\omega t\right)
\end{equation*}%
and 
\begin{equation*}
y_{2}=y_{\max }\sin \left( kx-\omega t+\phi _{o}\right)
\end{equation*}%
These are two waves with the \textbf{same frequency} and \textbf{same wave
number} traveling the \textbf{same direction} in the medium, but at $t=0$
the $y$ values are not the same because of the phase shift. The snapshot
graph of $y_{2}$ is shifted by an amount $\phi _{o}.$

I will pick some values for the constants

\begin{equation*}
\begin{tabular}{l}
$\lambda =2$ \\ 
$k=\frac{2\pi }{\lambda }$ \\ 
$\omega =1$ \\ 
$\phi _{o}=\frac{\pi }{6}$ \\ 
$t=0$ \\ 
$A=1$%
\end{tabular}%
\end{equation*}

Then for $y_{1}$ we have

\begin{eqnarray*}
y_{1} &=&\left( 1\right) \sin \left( \frac{2\pi }{\lambda }x-\left( 1\right)
t\right) \\
&=&\sin \left( \frac{2\pi }{2}x-\left( 1\right) t\right) \\
&=&\sin \left( \pi x-t\right)
\end{eqnarray*}%
And here is a plot of the wave function, $y_{1}$\FRAME{dtbpFX}{2.9689in}{%
1.2816in}{0pt}{}{}{Plot}{\special{language "Scientific Word";type
"MAPLEPLOT";width 2.9689in;height 1.2816in;depth 0pt;display
"USEDEF";plot_snapshots TRUE;mustRecompute FALSE;lastEngine "MuPAD";xmin
"0";xmax "5.001000";xviewmin "0";xviewmax "5.001000";yviewmin "-2";yviewmax
"2";viewset"XY";rangeset"X";plottype 4;axesFont "Times New
Roman,12,0000000000,useDefault,normal";numpoints 100;plotstyle
"patch";axesstyle "normal";axestips FALSE;xis \TEXUX{x};var1name
\TEXUX{$x$};function \TEXUX{$\allowbreak \sin \pi x$};linecolor
"green";linestyle 1;pointstyle "point";linethickness 3;lineAttributes
"Solid";var1range "0,5.001000";num-x-gridlines 100;curveColor
"[flat::RGB:0x00006000]";curveStyle "Line";VCamFile
'SV53L842.xvz';valid_file "T";tempfilename
'SV53NZ0S.wmf';tempfile-properties "XPR";}}

Now let's consider $y_{2.}$ Using the values we chose, $y_{2}$ can be
written as%
\begin{eqnarray*}
y_{2} &=&y_{\max }\sin \left( kx-\omega t+\phi _{o}\right) \\
&=&\sin \left( \pi x-t+\frac{\pi }{6}\right)
\end{eqnarray*}%
which looks like this\FRAME{dtbpFX}{3.0519in}{1.2427in}{0pt}{}{}{Plot}{%
\special{language "Scientific Word";type "MAPLEPLOT";width 3.0519in;height
1.2427in;depth 0pt;display "USEDEF";plot_snapshots TRUE;mustRecompute
FALSE;lastEngine "MuPAD";xmin "0";xmax "5.001000";xviewmin "0";xviewmax
"5.001000";yviewmin "-2";yviewmax "2";viewset"XY";rangeset"X";plottype
4;axesFont "Times New Roman,12,0000000000,useDefault,normal";numpoints
100;plotstyle "patch";axesstyle "normal";axestips FALSE;xis
\TEXUX{x};var1name \TEXUX{$x$};function \TEXUX{$\sin \left( \frac{1}{6}\pi
+\pi x\right) $};linecolor "red";linestyle 1;pointstyle
"point";linethickness 3;lineAttributes "Solid";var1range
"0,5.001000";num-x-gridlines 100;curveColor
"[flat::RGB:0x00ff0000]";curveStyle "Line";VCamFile
'SV53L841.xvz';valid_file "T";tempfilename
'SV53NZ0T.wmf';tempfile-properties "XPR";}}What does it look like if we add
these waves using superposition? Symbolically we have%
\begin{equation}
y_{r}=y_{\max }\sin \left( kx-\omega t\right) +y_{\max }\sin \left(
kx-\omega t+\phi _{o}\right)  \label{Superposition basic}
\end{equation}

and putting in the numbers gives%
\begin{equation*}
y_{r}=\sin \left( \pi x-t\right) +\sin \left( \pi x-t+\frac{\pi }{6}\right)
\end{equation*}%
which is shown in the next graph.

\FRAME{dtbpFX}{3.0727in}{1.3292in}{0pt}{}{}{Plot}{\special{language
"Scientific Word";type "MAPLEPLOT";width 3.0727in;height 1.3292in;depth
0pt;display "USEDEF";plot_snapshots TRUE;mustRecompute FALSE;lastEngine
"MuPAD";xmin "0";xmax "5.001000";xviewmin "0";xviewmax "5.001000";yviewmin
"-2";yviewmax "2";viewset"XY";rangeset"X";plottype 4;axesFont "Times New
Roman,12,0000000000,useDefault,normal";numpoints 100;plotstyle
"patch";axesstyle "normal";axestips FALSE;xis \TEXUX{x};var1name
\TEXUX{$x$};function \TEXUX{$\sin \left( \frac{1}{6}\pi +\pi x\right) +\sin
\pi x$};linecolor "green";linestyle 1;pointstyle "point";linethickness
3;lineAttributes "Solid";var1range "0,5.001000";num-x-gridlines
100;curveColor "[flat::RGB:0x00008000]";curveStyle "Line";VCamFile
'SV53L840.xvz';valid_file "T";tempfilename
'SV53NZ0U.wmf';tempfile-properties "XPR";}}Notice that the resulting wave
form is taller (larger amplitude). Noticed it is shifted along the $x$ axis.

We can find out by how much by rewriting $y_{r}.$ We need a trig identity%
\begin{equation*}
\sin a+\sin b=2\cos \left( \frac{a-b}{2}\right) \sin \left( \frac{a+b}{2}%
\right)
\end{equation*}%
Then let $a=kx-\omega t$ and $b=kx-\omega t+\phi $%
\begin{eqnarray*}
y_{r} &=&y_{\max }\sin \left( kx-\omega t\right) +A\sin \left( kx-\omega
t+\phi _{o}\right) \\
&=&2y_{\max }\cos \left( \frac{\left( kx-\omega t\right) -\left( kx-\omega
t+\phi _{o}\right) }{2}\right) \sin \left( \frac{\left( kx-\omega t\right)
+\left( kx-\omega t+\phi _{o}\right) }{2}\right) \\
&=&2y_{\max }\cos \left( \frac{-\phi _{o}}{2}\right) \sin \left( \frac{%
2kx-2\omega t+\phi _{o}}{2}\right) \\
&=&2y_{\max }\cos \left( \frac{-\phi _{o}}{2}\right) \sin \left( kx-\omega t+%
\frac{\phi _{o}}{2}\right) \\
&=&2y_{\max }\cos \left( \frac{\phi _{o}}{2}\right) \sin \left( kx-\omega t+%
\frac{\phi _{o}}{2}\right)
\end{eqnarray*}

where we used another trig identity $\cos \left( -\theta \right) =\cos
\left( \theta \right) .$

Let's look at the parts of this expression. First take the sine part.

\begin{equation}
\sin \left( kx-\omega t+\frac{\phi _{o}}{2}\right)
\end{equation}%
This part is a traveling wave with the same $k$ and $\omega $ as our
original waves, but it has a phase of $\phi _{o}/2.$ So our combined wave is
shifted by $\phi _{o}/2$ or half the phase shift of $y_{2}.$

Now let's look at other factor%
\begin{equation}
2y_{\max }\cos \left( \frac{\phi _{o}}{2}\right)
\end{equation}%
The sine part of our wave equation is multiplied by all of this factor. So
all of this part is the new amplitude. And this complicated amplitude has a
maximum value when $\phi _{o}=0$

\FRAME{dtbpFX}{2.5314in}{1.0617in}{0pt}{}{}{Plot}{\special{language
"Scientific Word";type "MAPLEPLOT";width 2.5314in;height 1.0617in;depth
0pt;display "USEDEF";plot_snapshots TRUE;mustRecompute FALSE;lastEngine
"MuPAD";xmin "-5";xmax "5";xviewmin "-5.0010000010002";xviewmax
"5.0010000010002";yviewmin "-1.60264739640757";yviewmax "2";plottype
4;labeloverrides 3;x-label "phi_o";y-label "Amplitude";axesFont "Times New
Roman,12,0000000000,useDefault,normal";numpoints 100;plotstyle
"patch";axesstyle "normal";axestips FALSE;xis \TEXUX{x};var1name
\TEXUX{$x$};function \TEXUX{$2\left( 1\right) \cos \left( \frac{x}{2}\right)
$};linecolor "blue";linestyle 1;pointstyle "point";linethickness
1;lineAttributes "Solid";var1range "-5,5";num-x-gridlines 100;curveColor
"[flat::RGB:0x000000ff]";curveStyle "Line";VCamFile
'T96NQT0C.xvz';valid_file "T";tempfilename
'T96NQT00.wmf';tempfile-properties "XPR";}}When $\phi _{o}=\pi ,$ then 
\begin{equation*}
2y_{\max }\cos \left( \frac{\pi }{2}\right) =0
\end{equation*}%
So, when $\phi _{o}=0$ we have a new maximum amplitude of $2y_{\max }$ and
when $\phi _{o}=\pi $ we have a zero amplitude. And these are our
constructive and destructive interference! Here is our wave for several
choices of $\phi _{o}.$ We can see that along with total constructive and
destructive interference we can also have partial constructive and partial
destructive interference.

\FRAME{dtbpFX}{3.1254in}{1.177in}{0pt}{}{}{Plot}{\special{language
"Scientific Word";type "MAPLEPLOT";width 3.1254in;height 1.177in;depth
0pt;display "USEDEF";plot_snapshots TRUE;mustRecompute FALSE;lastEngine
"MuPAD";xmin "0";xmax "5.001000";xviewmin "0";xviewmax "5.001000";yviewmin
"-2";yviewmax "2";viewset"XY";rangeset"X";plottype 4;axesFont "Times New
Roman,12,0000000000,useDefault,normal";numpoints 100;plotstyle
"patch";axesstyle "normal";axestips FALSE;xis \TEXUX{x};var1name
\TEXUX{$x$};function \TEXUX{$\allowbreak 2\sin \left( \frac{1}{2}\left(
0\right) +\pi x\right) \cos \frac{1}{2}\left( 0\right) $};linecolor
"black";linestyle 1;pointstyle "point";linethickness 2;lineAttributes
"Solid";var1range "0,5.001000";num-x-gridlines 100;curveColor
"[flat::RGB:0000000000]";curveStyle "Line";function \TEXUX{$\allowbreak
2\sin \left( \frac{1}{2}\frac{\pi }{9}+\pi x\right) \cos
\frac{1}{2}\frac{\pi }{9}$};linecolor "blue";linestyle 1;pointstyle
"point";linethickness 2;lineAttributes "Solid";var1range
"0,5.001000";num-x-gridlines 100;curveColor
"[flat::RGB:0x00000080]";curveStyle "Line";function \TEXUX{$\allowbreak
2\sin \left( \frac{1}{2}\frac{2\pi }{9}+\pi x\right) \cos
\frac{1}{2}\frac{2\pi }{9}$};linecolor "blue";linestyle 1;pointstyle
"point";linethickness 2;lineAttributes "Solid";var1range
"0,5.001000";num-x-gridlines 100;curveColor
"[flat::RGB:0x000000ff]";curveStyle "Line";function \TEXUX{$\allowbreak
2\sin \left( \frac{1}{2}\frac{3\pi }{9}+\pi x\right) \cos
\frac{1}{2}\frac{3\pi }{9}$};linecolor "green";linestyle 1;pointstyle
"point";linethickness 2;lineAttributes "Solid";var1range
"0,5.001000";num-x-gridlines 100;curveColor
"[flat::RGB:0x00008000]";curveStyle "Line";function \TEXUX{$2\sin \left(
\frac{1}{2}\frac{4\pi }{9}+\pi x\right) \cos \frac{1}{2}\frac{4\pi
}{9}$};linecolor "cyan";linestyle 1;pointstyle "point";linethickness
2;lineAttributes "Solid";var1range "0,5.001000";num-x-gridlines
100;curveColor "[flat::RGB:0x00008080]";curveStyle "Line";function
\TEXUX{$\allowbreak 2\sin \left( \frac{1}{2}\frac{5\pi }{9}+\pi x\right)
\cos \frac{1}{2}\frac{5\pi }{9}$};linecolor "cyan";linestyle 1;pointstyle
"point";linethickness 2;lineAttributes "Solid";var1range
"0,5.001000";num-x-gridlines 100;curveColor
"[flat::RGB:0x000080c0]";curveStyle "Line";function \TEXUX{$2\sin \left(
\frac{1}{2}\frac{6\pi }{9}+\pi x\right) \cos \frac{1}{2}\frac{6\pi
}{9}$};linecolor "green";linestyle 1;pointstyle "point";linethickness
2;lineAttributes "Solid";var1range "0,5.001000";num-x-gridlines
100;curveColor "[flat::RGB:0x0000ff00]";curveStyle "Line";function
\TEXUX{$2\sin \left( \frac{1}{2}\frac{7\pi }{9}+\pi x\right) \cos
\frac{1}{2}\frac{7\pi }{9}$};linecolor "black";linestyle 1;pointstyle
"point";linethickness 2;lineAttributes "Solid";var1range
"0,5.001000";num-x-gridlines 100;curveColor
"[flat::RGB:0x00400000]";curveStyle "Line";function \TEXUX{$2\sin \left(
\frac{1}{2}\frac{8\pi }{9}+\pi x\right) \cos \frac{1}{2}\frac{8\pi
}{9}$};linecolor "black";linestyle 1;pointstyle "point";linethickness
2;lineAttributes "Solid";var1range "0,5.001000";num-x-gridlines
100;curveColor "[flat::RGB:0x00404040]";curveStyle "Line";function
\TEXUX{$2\sin \left( \frac{1}{2}\frac{9\pi }{9}+\pi x\right) \cos
\frac{1}{2}\frac{9\pi }{9}$};linecolor "gray";linestyle 1;pointstyle
"point";linethickness 2;lineAttributes "Solid";var1range
"0,5.001000";num-x-gridlines 100;curveColor
"[flat::RGB:0x00cd99ff]";curveStyle "Line";VCamFile
'SV53L73Y.xvz';valid_file "T";tempfilename
'SV53NZ0W.wmf';tempfile-properties "XPR";}}

\subsection{Mathematical description of standing waves}

Now that we have a way to make two waves to superimpose, we can study the
special case of standing waves. 
%TCIMACRO{%
%\TeXButton{Standing Wave Demo}{\marginpar {
%\hspace{-0.5in}
%\begin{minipage}[t]{1in}
%\small{Standing Wave Demo}
%\end{minipage}
%}}}%
%BeginExpansion
\marginpar {
\hspace{-0.5in}
\begin{minipage}[t]{1in}
\small{Standing Wave Demo}
\end{minipage}
}%
%EndExpansion

A standing wave pattern is the result of the superposition of two traveling
waves with the same frequency going in opposite directions. Let's take two
specific waves%
\begin{eqnarray*}
y_{1} &=&y_{\max }\sin \left( kx-\omega t\right) \\
y_{2} &=&y_{\max }\sin \left( kx+\omega t\right)
\end{eqnarray*}

The sum is 
\begin{equation*}
y_{r}=y_{1}+y_{2}=y_{\max }\sin \left( kx-\omega t\right) +A\sin \left(
kx+\omega t\right)
\end{equation*}%
Which we could just do on a computer. But to gain insight into what these
two waves produce, let's use another of our favorite trig identities%
\begin{equation*}
\sin \left( a\pm b\right) =\sin \left( a\right) \cos \left( b\right) \pm
\cos \left( a\right) \sin \left( b\right)
\end{equation*}%
to get%
\begin{eqnarray*}
y_{r} &=&y_{\max }\sin \left( kx-\omega t\right) +y_{\max }\sin \left(
kx+\omega t\right) \\
&=&y_{\max }\sin \left( kx\right) \cos \left( \omega t\right) -y_{\max }\cos
\left( kx\right) \sin \left( \omega t\right) +y_{\max }\sin \left( kx\right)
\cos \left( \omega t\right) +y_{\max }\cos \left( kx\right) \sin \left(
\omega t\right) \\
&=&2y_{\max }\sin \left( kx\right) \cos \left( \omega t\right) \\
&=&\left( 2y_{\max }\sin \left( kx\right) \right) \cos \left( \omega t\right)
\end{eqnarray*}%
This looks like the harmonic oscillator equation%
\begin{equation*}
y=y_{\max }\cos \left( \omega t+\phi _{o}\right)
\end{equation*}%
with $\phi _{o}=0$ and with another complicated amplitude%
\begin{equation*}
2y_{\max }\sin \left( kx\right)
\end{equation*}%
That is, we have a set of harmonic oscillators all at different values of $x$
who's amplitude is different for each value of $x.$

\FRAME{dhF}{2.5296in}{2.2883in}{0in}{}{}{Figure}{\special{language
"Scientific Word";type "GRAPHIC";maintain-aspect-ratio TRUE;display
"USEDEF";valid_file "T";width 2.5296in;height 2.2883in;depth
0in;original-width 2.4889in;original-height 2.2485in;cropleft "0";croptop
"1";cropright "1";cropbottom "0";tempfilename
'SV53TJ19.wmf';tempfile-properties "XPR";}}We can identify spots along the $%
x $ axis where the amplitude is always zero! we will call these spots \emph{%
nodes.} These happen when $\sin \left( kx\right) =0$ or when 
\begin{equation*}
kx=n\pi
\end{equation*}

By using 
\begin{equation*}
k=\frac{2\pi }{\lambda }
\end{equation*}%
we have%
\begin{equation*}
\frac{2\pi }{\lambda }x=n\pi
\end{equation*}%
We can cancel things, and solve for $x$%
\begin{eqnarray*}
\frac{2}{\lambda }x &=&n \\
x &=&n\frac{\lambda }{2}
\end{eqnarray*}

We can also find the places along $x$ where the amplitude will be largest.
this occurs when $\sin \left( kx\right) =1$ or when%
\begin{equation*}
kx=n\frac{\pi }{2}
\end{equation*}%
or%
\begin{eqnarray*}
\frac{2\pi }{\lambda }x &=&n\frac{\pi }{2} \\
x &=&n\frac{\lambda }{4}
\end{eqnarray*}%
these are called \emph{antinodes}. But notice that if $n=2$ then $x=\frac{%
\lambda }{2}$ and we know that will be a node. So not all values of $n$ work
here. Let's try $n=3.$ Then $x=\frac{3\lambda }{4}$ and that works. To solve
this issue we can say we will only allow odd values of $n.$ 
\begin{equation*}
n=1,3,5,\cdots
\end{equation*}%
Then we get antinodes when 
\begin{equation*}
x=n\frac{\lambda }{4}\text{\qquad }n=1,3,5,\cdots
\end{equation*}

We can also create standing waves with sound or even light waves!

\subsection{Standing Waves in a String Fixed at Both Ends}

\FRAME{dhF}{3.2846in}{2.4699in}{0in}{}{}{Figure}{\special{language
"Scientific Word";type "GRAPHIC";maintain-aspect-ratio TRUE;display
"USEDEF";valid_file "T";width 3.2846in;height 2.4699in;depth
0in;original-width 3.2396in;original-height 2.4284in;cropleft "0";croptop
"1";cropright "1";cropbottom "0";tempfilename
'SV53TJ1A.wmf';tempfile-properties "XPR";}}If we attach a string to
something on both ends, we find something interesting in the standing wave
pattern. Not all possible standing waves can be realized. Some frequencies
are preferred, and some never show up. The standing wave pattern is \emph{%
quantized}. The patterns that are allowed are called \emph{normal modes}. We
will see this any time a wave confined by boundary conditions (light in a
resonant cavity, radio waves in a wave guide, electrons in an atom, etc.).

The figure shows some normal modes for a string.

We find which modes are allowed by first imposing the boundary condition
that each end must be a node. We start with 
\begin{equation*}
y=2y_{\max }\sin \left( kx\right) \cos \left( \omega t\right)
\end{equation*}%
and realize we have two boundary conditions. $y\left( 0,t\right) =0$ and and 
$y\left( L,t\right) =0$ because at these locations the string is attached to
a piece of wood (called the nut or bridge for a ukulele). The string can't
oscillate at these locations. For the first conditions 
\begin{equation*}
y=0
\end{equation*}%
when 
\begin{equation*}
x=0
\end{equation*}%
and for the second $y=0$ when $x=L.$ That happens when%
\begin{equation*}
kL=n\pi
\end{equation*}%
I will write this as%
\begin{equation*}
k_{n}L=n\pi
\end{equation*}%
to indicate there are many values of $k$ that can work. Solving this for $%
\lambda _{n}$ gives%
\begin{eqnarray*}
\frac{2\pi }{\lambda _{n}}L &=&n\pi \\
\frac{2L}{n} &=&\lambda _{n}
\end{eqnarray*}%
Which just says there are many wavelengths that will work. Let's see how to
use this, the first mode will have 
\begin{equation*}
\lambda _{1}=2L
\end{equation*}%
where $L$ is the length of the string. Looking at the figure, this works.

The second mode has three nodes (one on each end and one in the middle).
This gives%
\begin{equation*}
\lambda _{2}=L
\end{equation*}%
We can keep going, the third mode will have five nodes%
\begin{equation*}
\lambda _{3}=\frac{3L}{2}
\end{equation*}%
and so forth to give%
\begin{equation*}
\lambda _{n}=\frac{2L}{n}
\end{equation*}

We use our old friend, the equation for wave speed,%
\begin{equation*}
v=f\lambda
\end{equation*}%
to find the frequencies of the modes%
\begin{equation*}
f=\frac{v}{\lambda }
\end{equation*}

\begin{equation*}
f_{1}=\frac{v}{\lambda _{1}}=\frac{v}{2L}
\end{equation*}%
or, in general%
\begin{eqnarray*}
f_{n} &=&\frac{v}{\lambda _{n}}=n\frac{v}{2L} \\
&=&\frac{n}{2L}v \\
&=&\frac{n}{2L}\sqrt{\frac{T}{\mu }}
\end{eqnarray*}%
The lowest frequency has a special name, the \emph{fundamental frequency}.
The higher frequencies are integer multiples of the fundamental. When this
happens we say that the frequencies form a \emph{harmonic series, }and the
modes are called \emph{harmonics}. Since only certain frequencies work, we
say that the frequencies of waves on the string that make standing waves are 
\emph{quantized}. Quantization means that some values are allowed. This idea
is the basis behind Quantum mechanics (which views light and even matter as
waves).

\subsection{Musical Strings}

So how do we get different notes on a guitar or Piano?%
\begin{equation}
f_{n}=\frac{n}{2L}\sqrt{\frac{T}{\mu }}
\end{equation}

A guitar uses tension to change the frequency or pitch (tuning) and length
of string (your fingers pressing on the strings) to change notes. A Piano
uses both tension and length of string (and mass per unit length as well!).
What do you expect and organ will do?

It turns out that the musical note we hear is related to the fundamental
frequency. So the musical note for a stringed instrument is given by 
\begin{equation}
f_{1}=\frac{1}{2L}\sqrt{\frac{T}{\mu }}
\end{equation}

\section{Sound waves}

Sound is a wave. The medium is air particles. The transfer of energy is done
by collision. \FRAME{dhF}{1.0369in}{1.1744in}{0pt}{}{}{Figure}{\special%
{language "Scientific Word";type "GRAPHIC";maintain-aspect-ratio
TRUE;display "USEDEF";valid_file "T";width 1.0369in;height 1.1744in;depth
0pt;original-width 2.5581in;original-height 2.9006in;cropleft "0";croptop
"1";cropright "1";cropbottom "0";tempfilename
'SV51AY0C.wmf';tempfile-properties "XPR";}}The wave will be a longitudinal
wave. Let's see how it forms. We can take a tube with a piston in it. \FRAME{%
dhF}{3.2707in}{0.7809in}{0pt}{}{}{Figure}{\special{language "Scientific
Word";type "GRAPHIC";maintain-aspect-ratio TRUE;display "USEDEF";valid_file
"T";width 3.2707in;height 0.7809in;depth 0pt;original-width
3.2258in;original-height 0.7498in;cropleft "0";croptop "1";cropright
"1";cropbottom "0";tempfilename 'SV51AY0D.wmf';tempfile-properties "XPR";}}%
As we exert a force on the piston, the air molecules are compressed into a
group. In the next figure, each dot represents a group of air molecules. In
the top picture, the air molecules are not displaced. But when the piston
moves, the air molecules receive energy by collision. They bunch up. We see
this in the second picture. \FRAME{dtbpF}{3.0193in}{1.5061in}{0pt}{}{}{Figure%
}{\special{language "Scientific Word";type "GRAPHIC";maintain-aspect-ratio
TRUE;display "USEDEF";valid_file "T";width 3.0193in;height 1.5061in;depth
0pt;original-width 3.0486in;original-height 1.5061in;cropleft "0";croptop
"1";cropright "1";cropbottom "0";tempfilename
'T96Q0T01.wmf';tempfile-properties "XPR";}}The colored dots represent a few
of the air molecules. The graph below the two pictures shows how much
displacement each molecule group experiences from its equilibrium position
(where it started).

Suppose we now pull the piston back. This would allow the molecules to
bounce back to the left, but the molecules that they have collided with will
receive some energy and go to the right. This is shown in the next figure.
Color coded dots are displayed above the before and after picture so you can
see where each molecule started.\FRAME{dtbpF}{3.131in}{1.6329in}{0pt}{}{}{%
Figure}{\special{language "Scientific Word";type
"GRAPHIC";maintain-aspect-ratio TRUE;display "USEDEF";valid_file "T";width
3.131in;height 1.6329in;depth 0pt;original-width 3.3093in;original-height
1.711in;cropleft "0";croptop "1";cropright "1";cropbottom "0";tempfilename
'T96Q1H02.wmf';tempfile-properties "XPR";}}If we pull the piston back
further, the molecules can pass their original positions.\FRAME{dtbpF}{%
3.2428in}{1.7411in}{0pt}{}{}{Figure}{\special{language "Scientific
Word";type "GRAPHIC";maintain-aspect-ratio TRUE;display "USEDEF";valid_file
"T";width 3.2428in;height 1.7411in;depth 0pt;original-width
3.0716in;original-height 1.6365in;cropleft "0";croptop "1";cropright
"1";cropbottom "0";tempfilename 'T96Q2403.wmf';tempfile-properties "XPR";}}%
Then we can push inward again and compress the gas.\FRAME{dtbpF}{3.0219in}{%
1.6941in}{0pt}{}{}{Figure}{\special{language "Scientific Word";type
"GRAPHIC";maintain-aspect-ratio TRUE;display "USEDEF";valid_file "T";width
3.0219in;height 1.6941in;depth 0pt;original-width 3.0512in;original-height
1.6977in;cropleft "0";croptop "1";cropright "1";cropbottom "0";tempfilename
'T96Q3204.wmf';tempfile-properties "XPR";}}This may seem like a senseless
thing to do, but it is really what a speaker does to produce sound. In
particular, a speaker is a harmonic oscillator. The simple harmonic motion
of the speaker is the disturbance that makes the sound wave.\FRAME{dhF}{%
2.6135in}{2.1344in}{0pt}{}{}{Figure}{\special{language "Scientific
Word";type "GRAPHIC";maintain-aspect-ratio TRUE;display "USEDEF";valid_file
"T";width 2.6135in;height 2.1344in;depth 0pt;original-width
2.5719in;original-height 2.0954in;cropleft "0";croptop "1";cropright
"1";cropbottom "0";tempfilename 'SV51AY0I.wmf';tempfile-properties "XPR";}}

\subsection{Periodic Sound Waves, Pressure}

%TCIMACRO{%
%\TeXButton{Vollyball Demo}{\marginpar {
%\hspace{-0.5in}
%\begin{minipage}[t]{1in}
%\small{Vollyball Demo}
%\end{minipage}
%}}}%
%BeginExpansion
\marginpar {
\hspace{-0.5in}
\begin{minipage}[t]{1in}
\small{Vollyball Demo}
\end{minipage}
}%
%EndExpansion
Suppose I have a ball and I\ ask six people from the class to come up and
press on the ball from all directions. This is a new force situation, unlike
most we dealt with in Dynamics or PH121. Each person exerts a force on the
ball. The person uses the area of their hand to exert the force. The motion
of the ball, and even it's shape depend on both the force (magnitude and
direction) and the area involved in each push.

From our demo, it seems that the force and area of the ball could be related
to better describe the situation. Let's look at the ratio%
\begin{equation*}
\frac{F}{A}
\end{equation*}%
What does it represent? This ratio tells us how spread out an applied force
may be. If I push on the ball with my hand the ball might sail across a net.
If I push with the same force with a sharp nail, I will need a new ball. The
small area of the nail tip makes the force more likely to penetrate into the
ball. The area is important.

It is coinvent to give this concept a name. We will call it \emph{pressure.}

\begin{center}
\rule{300pt}{1pt}

\textbf{Pressure} is a scalar value that describes how a force acts over an
area%
\begin{equation*}
P=\frac{F}{A}
\end{equation*}

\rule{300pt}{1pt}
\end{center}

But in a fluid like air, what is the force? Lets consider a ball hitting a
wall. Is there a force? During the collision, there is a force of the ball
on the wall, and a Newton's third law force of the wall molecules pushing
back on the ball. The ball will bounce back. 
%TCIMACRO{%
%\TeXButton{Ping Pong Ball Demo}{\marginpar {
%\hspace{-0.5in}
%\begin{minipage}[t]{1in}
%\small{Ping Pong Ball Demo}
%\end{minipage}
%}}}%
%BeginExpansion
\marginpar {
\hspace{-0.5in}
\begin{minipage}[t]{1in}
\small{Ping Pong Ball Demo}
\end{minipage}
}%
%EndExpansion

This force is small and only lasts during the collision. But now suppose we
have many balls, and all the balls impact the wall. Further suppose that
every time a ball bounces back it ends up headed back to the wall and
bounces again. If the balls keep coming, there will be a force on the wall
quite a bit of the time. At least, on average there is a force, anyway. This
is the force that causes air pressure. The air molecules in this room are
like the small balls. we will find that they have an amount of kinetic
energy. They impact the walls (and us) all over the surface area. The result
is air pressure.

The water pressure in a swimming pool is caused by moving water particles.
You should convince yourself that the reason the water stays in the pool is
partly because the air molecules bounce against the water surface exerting a
pressure on the waver so it can't leave the pool!

Let's go back to making sounds. Suppose we push our piston as we did before
in figure \ref{Creation of a Pulse}. When we push in the piston, it creates
a region of higher pressure next to it.

\begin{center}
\rule{300pt}{1pt}

\textbf{Compression:} A local region of higher pressure in a fluid

\rule{300pt}{1pt}

\textbf{Rarefaction:} A local region of lower pressure in a fluid

\rule{300pt}{1pt}
\end{center}

When we pull back the piston the fluid expands to fill the void. We create a
rarefaction next to the piston.

Suppose we drive the piston sinusoidally. Can we describe the motion of the
particles and of the wave? We can identify the distance between two
compressions as $\lambda .$ We define $s\left( x,t\right) $ (like we defined
a wave function, $y\left( x,t\right) )$ as the displacement of a particle of
fluid relative to its equilibrium position.%
\begin{equation}
s\left( x,t\right) =s_{\max }\cos \left( kx-\omega t\right)
\end{equation}%
but what is $s_{\max }?$

We remember that $s_{\max }$ is the maximum displacement of a particle of
fluid from its equilibrium position. We plotted this using a bar graph to
show displacement from the equilibrium position for our molecules. As we
push the piston in and out we will get something like this. \FRAME{dtbpF}{%
3.1559in}{4.3657in}{0pt}{}{}{Figure}{\special{language "Scientific
Word";type "GRAPHIC";maintain-aspect-ratio TRUE;display "USEDEF";valid_file
"T";width 3.1559in;height 4.3657in;depth 0pt;original-width
3.1878in;original-height 4.4207in;cropleft "0";croptop "1";cropright
"1";cropbottom "0";tempfilename 'T96QOB06.wmf';tempfile-properties "XPR";}}%
We found before that we get something that looks like a sine wave, but
remember what the bars represent. They represent the displacement from
original position. \FRAME{dtbpF}{2.5048in}{0.9615in}{0pt}{}{}{Figure}{%
\special{language "Scientific Word";type "GRAPHIC";maintain-aspect-ratio
TRUE;display "USEDEF";valid_file "T";width 2.5048in;height 0.9615in;depth
0pt;original-width 2.5235in;original-height 0.9509in;cropleft "0";croptop
"1";cropright "1";cropbottom "0";tempfilename
'T96Q4305.wmf';tempfile-properties "XPR";}}We don't usually draw bar graphs,
we usually just draw the sine wave.

\FRAME{dtbpF}{2.3315in}{1.9294in}{0in}{}{}{Figure}{\special{language
"Scientific Word";type "GRAPHIC";maintain-aspect-ratio TRUE;display
"USEDEF";valid_file "T";width 2.3315in;height 1.9294in;depth
0in;original-width 2.2917in;original-height 1.8913in;cropleft "0";croptop
"1";cropright "1";cropbottom "0";tempfilename
'SUS73V26.wmf';tempfile-properties "XPR";}}

The variation of the gas pressure $\Delta P$ measured from its equilibrium
is also periodic

\FRAME{dhF}{2.6083in}{2.8072in}{0pt}{}{}{Figure}{\special{language
"Scientific Word";type "GRAPHIC";maintain-aspect-ratio TRUE;display
"USEDEF";valid_file "T";width 2.6083in;height 2.8072in;depth
0pt;original-width 6.02in;original-height 6.4809in;cropleft "0";croptop
"1";cropright "1";cropbottom "0";tempfilename
'SUS73V27.wmf';tempfile-properties "XPR";}}which is why we often refer to a
sound wave as a pressure wave. Think of when the wave gets to your ear. the
wave consists of a group of particles all headed for your ear drum. When
they hit, they exert a force. Pressure is a force spread over an area. This
is the beginning of hearing sound!

%TCIMACRO{%
%\TeXButton{Basic Equations}{\vspace{0.5in}\hspace{-1.3in}{\LARGE Basic Equations\vspace{0.25in}}}}%
%BeginExpansion
\vspace{0.5in}\hspace{-1.3in}{\LARGE Basic Equations\vspace{0.25in}}%
%EndExpansion

Standing waves%
\begin{equation*}
y_{r}=\left( 2y_{\max }\sin \left( kx\right) \right) \cos \left( \omega
t\right)
\end{equation*}

Nodes%
\begin{equation*}
x=n\frac{\lambda }{2}
\end{equation*}%
Antinodes%
\begin{equation*}
x=n\frac{\lambda }{4}\text{\qquad }n=1,3,5,\cdots
\end{equation*}%
Standing wave frequencies (for waves on strings)

\begin{equation}
f_{n}=\frac{n}{2L}\sqrt{\frac{T}{\mu }}
\end{equation}

Sound waves%
\begin{equation}
s\left( x,t\right) =s_{\max }\cos \left( kx-\omega t\right)
\end{equation}

\begin{equation*}
P=\frac{F}{A}
\end{equation*}

\chapter{7 Sound Wave Speed and Intensity 1.17.2, 1.17.3}

We know the basics of how a sound wave is formed. Now let's put this on a
more mathematical footing.

%TCIMACRO{%
%\TeXButton{\chaptermark{Speed and Intensity}}{\chaptermark{Speed and Intensity}}}%
%BeginExpansion
\chaptermark{Speed and Intensity}%
%EndExpansion

%TCIMACRO{%
%\TeXButton{Fundamental Concepts}{\vspace{0.10in}\hspace{-0.37in}{\Large Fundamental Concepts\vspace{0.2in}}}}%
%BeginExpansion
\vspace{0.10in}\hspace{-0.37in}{\Large Fundamental Concepts\vspace{0.2in}}%
%EndExpansion

\begin{itemize}
\item We don't hear all frequencies equally well

\item Waves from point sources are spherical

\item Light waves are waves in the electromagnetic field
\end{itemize}

\section{Speed of Sound Waves}

We found that waves on ropes and strings had a speed that depended on the
stretchy bonds in the medium. Our equation for this is 
\begin{equation*}
v=\sqrt{\frac{T}{\mu }}
\end{equation*}%
where $T$ is the tension in the rope or string and $\mu =m/L$ is the linear
mass density.

Let's think about this equation. The tension is an elastic property of the
string or rope medium. And $\mu $ has the mass in it. It is an inertial
property of the string or rope. It tells us how hard the string or rope is
to move.

We need to have something like this for sound waves. And we know that sound
waves have to do with pressure and the compression of the molecules that
make up the medium. Long ago in Principles of Physics I you should have
studied elastic properties of materials. But let's review them here.

\subsection{Young's Modulus}

Remember that a modulus is a constant that describes how a material can be
deformed. It tells us how elastic that material is. Think of it like the
spring constant $k$ that tells us how hard it is to stretch or compress a
spring.

Suppose we pull on a rod with a force $\overrightarrow{\mathbf{F}}\mathbf{.}$%
\FRAME{dtbpF}{3.3271in}{1.7678in}{0pt}{}{}{Figure}{\special{language
"Scientific Word";type "GRAPHIC";maintain-aspect-ratio TRUE;display
"USEDEF";valid_file "T";width 3.3271in;height 1.7678in;depth
0pt;original-width 5.2155in;original-height 2.7576in;cropleft "0";croptop
"1";cropright "1";cropbottom "0";tempfilename
'SV71LI02.wmf';tempfile-properties "XPR";}}

The rod will be stressed. We call this \emph{tensile stress} which comes
from a pull . This pull will stretch the solid's molecular bonds (like a
tension force). We write the stress as 
\begin{equation}
\frac{F}{A}
\end{equation}%
which looks like a pressure, but this time we mean that we are pulling on
the beam so that the pull is spread over the cross sectional area (marked $A$
in the figure above). You can think of the beam as having many slightly
stretchy bonds between molecules across the whole cross sectional area of
the beam. All of them contribute making the bar hard to stretch. That is not
much like a pressure!

But by pulling on the beam, the beam will stretch a little. That stretching
is a strain. We write the strain as the percent change in shape. In this
case, the percent change in length%
\begin{equation}
\frac{\Delta L}{L_{o}}
\end{equation}%
If the stress is not too great, then we can write a linear equation that
relates the stress to the strain. 
\begin{equation}
\frac{F}{A}=Y\frac{\Delta L}{L_{o}}
\end{equation}%
This linear equation only works in the elastic region of a stress vs. strain
curve for the particular object. The constant of proportionality, $Y,$ is
called \emph{Young's Modulus}. If we write this as%
\begin{eqnarray}
F &=&\frac{YA}{L_{o}}\Delta L \\
F &=&k\Delta L
\end{eqnarray}%
it looks very much like Hook's law. Our restoring force or our
\textquotedblleft pull\textquotedblright\ is the rod pulling back with a
force like $F=-k\left( x-x_{o}\right) .$

Young's modulus depends on the microscopic properties of the solid (which we
will not study in detail, but these are our stretchy/squishy forces between
atoms that make normal and tension forces). We will look up Young's modulus
for each material in tables like the following:\smallskip 
\begin{equation*}
\begin{tabular}{|l|l|}
\hline
Material\QQfnmark{%
*Ledbetter, H. M., Physical Properties Data Compilations Relevant to Energy
Storage, US National Bureaus of Standards, 1982. Different alloys have
different properties, so for any real work see the original tables in the
original publication.}$^{,}$\QQfnmark{%
\dag Average values from numbers given in various text books. These numbers
should be taken as example values and more exact numbers found for any real
work.} & Young's Modulus $\left( \unit{N}/\unit{m}^{2}\right) $ \\ \hline
Aluminum$^{\ast }$ & $7\times 10^{10}$ \\ \hline
Titanium$^{\ast }$ & $11.\times 10^{10}$ \\ \hline
Steel$^{\ast }$ & $20.\times 10^{10}$ \\ \hline
Carbon Steel$^{\ast }$ & $21.\times 10^{10}$ \\ \hline
Lead$^{\dag }$ & $1.6\times 10^{10}$ \\ \hline
Brass$^{\dag }$ & $10.\times 10^{10}$ \\ \hline
Concrete$^{\dag }$ & $2.0\times 10^{10}$ \\ \hline
Nylon$^{\dag }$ & $0.5\times 10^{10}$ \\ \hline
Bone$^{\dag }$ (arm or leg) & $1.5\times 10^{10}$ \\ \hline
\begin{tabular}{l}
Pine Wood$^{\dag }$ (parallel to grain) \\ 
(perpendicular to grain)%
\end{tabular}
& 
\begin{tabular}{l}
$1.0\times 10^{10}$ \\ 
$0.1\times 10^{10}$%
\end{tabular}
\\ \hline
\end{tabular}%
\QQfntext{-1}{
*Ledbetter, H. M., Physical Properties Data Compilations Relevant to Energy
Storage, US National Bureaus of Standards, 1982. Different alloys have
different properties, so for any real work see the original tables in the
original publication.}
\QQfntext{1}{
\dag Average values from numbers given in various text books. These numbers
should be taken as example values and more exact numbers found for any real
work.}
\end{equation*}%
\smallskip \smallskip

Let's take an example. Suppose we have a steel piano wire. One end of the
piano wire is fixed in place, and the other is connected to a peg that can
be rotated so the piano wire is tightened. This is what you do when you tune
a piano, you turn the peg and tighten the wire. Suppose we tighten the wire
so that there will be $980\unit{N}$ of tension on the wire. The wire is $1.6%
\unit{m}$ long and has a diameter of $2\unit{mm}.$ How much will the wire
stretch as we tighten the peg? We will need to know that the Young's modulus
for steel is $20\times 10^{10}\unit{N}/\unit{m}^{3}.$

Here is a summary of what we know%
\begin{eqnarray*}
T &=&980\unit{N} \\
L_{o} &=&1.6\unit{m} \\
d &=&2\unit{mm} \\
Y &=&200\times 10^{9}\unit{N}/\unit{m}^{2}
\end{eqnarray*}

Let's start with our stress/strain basic equation%
\begin{equation*}
\frac{F}{A}=Y\frac{\Delta L}{L_{o}}
\end{equation*}%
and solve for $\Delta L$%
\begin{equation*}
\frac{F}{A}\frac{L_{o}}{Y}=\Delta L
\end{equation*}%
the force is our tension, and the cross sectional area will be 
\begin{equation*}
A=\pi \frac{d^{2}}{4}
\end{equation*}%
And we can substitute this into our equation%
\begin{equation*}
\frac{T}{\pi \frac{d^{2}}{4}}\frac{L_{o}}{Y}=\Delta L
\end{equation*}%
or, after dealing with the fraction in the denominator%
\begin{equation*}
\frac{4T}{\pi d^{2}}\frac{L_{o}}{Y}=\Delta L
\end{equation*}%
then 
\begin{eqnarray*}
\Delta L &=&\frac{4\left( 980\unit{N}\right) }{\pi \left( 2\unit{mm}\right)
^{2}}\frac{\left( 1.6\unit{m}\right) }{\left( 200\times 10^{9}\unit{N}/\unit{%
m}^{2}\right) }= \\
&=&\allowbreak 2.\,\allowbreak 495\,5\times 10^{-3}\unit{m}
\end{eqnarray*}%
so the wire stretched about a quarter of a centimeter.

Now that we remember what a modulus is, let's consider the speed of sound
waves in our bar of metal. This is going to be different than our speed of
sound in air. In air the molecules have to travel and crash into each other
to make a sound wave, but in our bar the molecules are bound together. They
can't move. We might guess that this will speed up the sound because there
is no travel time for the medium pieces. And we might guess that the springy
bonding forces would be what transfers energy in the medium. But what is our
speed of sound?

Say, we hit the bar with a hammer. Sound waves would be created in the bar.
But suppose we make our bar thin and stretch it across a guitar. Now as we
hit the string we expect a wave. And that wave will have a wave speed. And
we know that wave speeds are related to the measure of the elasticity of the
metal and the measure of the inertial of the metal material. The tension in
the (now) wire would be our force, but we now know how to express this force
in terms of the elasticity of the wire material. We can start with 
\begin{equation*}
\frac{F}{A}=Y\frac{\Delta L}{L_{o}}
\end{equation*}%
again and put in $T$ for the force%
\begin{equation*}
\frac{T}{A}=Y\frac{\Delta L}{L_{o}}
\end{equation*}%
then%
\begin{equation*}
T=Y\frac{\Delta L}{L_{o}}A
\end{equation*}%
Now let's think back to Newton's second law. but this time for just a small
piece of the wire that has a mass of $dm$ and a length of $dx.$ \FRAME{dtbpF%
}{5.0398in}{1.334in}{0pt}{}{}{Figure}{\special{language "Scientific
Word";type "GRAPHIC";maintain-aspect-ratio TRUE;display "USEDEF";valid_file
"T";width 5.0398in;height 1.334in;depth 0pt;original-width
5.1081in;original-height 1.3305in;cropleft "0";croptop "1";cropright
"1";cropbottom "0";tempfilename 'TEBZJK05.wmf';tempfile-properties "XPR";}}%
Then the net force on this small piece of the wire would be

\FRAME{dtbpF}{1.6906in}{0.4364in}{0pt}{}{}{Figure}{\special{language
"Scientific Word";type "GRAPHIC";maintain-aspect-ratio TRUE;display
"USEDEF";valid_file "T";width 1.6906in;height 0.4364in;depth
0pt;original-width 1.6941in;original-height 0.416in;cropleft "0";croptop
"1";cropright "1";cropbottom "0";tempfilename
'TEBXPP00.wmf';tempfile-properties "XPR";}}%
\begin{equation*}
F_{net}=T_{r}-T_{\ell }=\Delta T
\end{equation*}%
and we want to look at the tensions on either side of our little piece of
wire. Let's use the letter $ds$ for the stretch of the small segment of
wire. It's a very small $\Delta L.$ Then we don't have 
\begin{equation*}
T=Y\frac{\Delta L}{L_{o}}A
\end{equation*}%
for the whole wire, we just have a tension 
\begin{equation*}
T=Y\frac{ds}{dx}A
\end{equation*}%
for our little piece of the wire. But in our Newton's second law equation we
have $\Delta T,$ the tension difference between the two sides of our bit of
wire. That tension difference is causing the stretch, $ds.$ Because our bit
of wire is so small we could write this as $dT.$ And from our last equation
we could find $dT/dx$ to find $dT$%
\begin{equation*}
\frac{d}{dx}T=\frac{d}{dx}\left( Y\frac{ds}{dx}A\right) =YA\frac{d^{2}s}{%
dx^{2}}
\end{equation*}%
so 
\begin{equation*}
dT=YA\frac{d^{2}s}{dx^{2}}dx
\end{equation*}%
Now our $F_{net}=ma$ for our little bit of wire would be 
\begin{equation*}
F_{net}=\left( dm\right) \left( \frac{d^{2}s}{dt^{2}}\right)
\end{equation*}%
We can change this a little bit if we remember what density is 
\begin{equation*}
\rho =\frac{m}{%
%TCIMACRO{\TeXButton{V}{\ooalign{\hfil$V$\hfil\cr\kern0.1em--\hfil\cr}}}%
%BeginExpansion
\ooalign{\hfil$V$\hfil\cr\kern0.1em--\hfil\cr}%
%EndExpansion
}
\end{equation*}%
where $%
%TCIMACRO{\TeXButton{V}{\ooalign{\hfil$V$\hfil\cr\kern0.1em--\hfil\cr}}}%
%BeginExpansion
\ooalign{\hfil$V$\hfil\cr\kern0.1em--\hfil\cr}%
%EndExpansion
$ is the volume. We want the density of our little bit of wire%
\begin{equation*}
\rho =\frac{dm}{Adx}
\end{equation*}%
because $%
%TCIMACRO{\TeXButton{V}{\ooalign{\hfil$V$\hfil\cr\kern0.1em--\hfil\cr}}}%
%BeginExpansion
\ooalign{\hfil$V$\hfil\cr\kern0.1em--\hfil\cr}%
%EndExpansion
=Adx.$ 
\begin{equation*}
dm=\rho Adx
\end{equation*}%
then the net force on our little bit of wire is 
\begin{equation*}
F_{net}=\rho Adx\frac{d^{2}s}{dt^{2}}
\end{equation*}%
And we can put all of this into our Newton's second law equation. Start with 
\begin{equation*}
F_{net}=T_{r}-T_{\ell }=\Delta T
\end{equation*}%
and substitute to get%
\begin{equation*}
\rho Adx\frac{d^{2}s}{dt^{2}}=YA\frac{d^{2}s}{dx^{2}}dx
\end{equation*}%
This looks a lot like our wave equation!. Some things cancel 
\begin{equation*}
\rho \frac{d^{2}s}{dt^{2}}=Y\frac{d^{2}s}{dx^{2}}
\end{equation*}%
Let's put all of our constants on one side.%
\begin{equation*}
\frac{d^{2}s}{dt^{2}}=\frac{Y}{\rho }\frac{d^{2}s}{dx^{2}}
\end{equation*}%
and now all we have to do is to compare this to our linear wave equation.%
\begin{equation*}
\frac{d^{2}y}{dt^{2}}=v^{2}\frac{d^{2}y}{dx^{2}}
\end{equation*}%
In our new equation there is a $Y/\rho $ where the $v^{2}$ should be. This
means for our wire to have a wave on it, the wave speed must be 
\begin{equation*}
v=\sqrt{\frac{Y}{\rho }}
\end{equation*}%
And we have the wave speed in terms of the an elastic quantity (Young's
modulus) and an inertial property (the metal density). This gives a wave
speed on a wire. But this won't do for sound waves in air. We need a
different modulus for gasses. And fortunately we also studied this modulus
in Principles of Physics I. It is the bulk modulus.

\subsection{Bulk Modulus}

This one is harder. It tells us how much a material can be squashed or,
compressed. Let's define 
\begin{equation}
\Delta P=\frac{\Delta F}{A}
\end{equation}%
where now $A$ is the whole outer surface area of the object. We sometimes
call such a force spread out over an object's surface a \emph{surface
pressure}. Often fluids are very hard to compress. Many solids are easy to
compress. Think of Styrofoam or any foam rubber. For solids $\Delta F/A$
where the force is all around the solid and the area is the surface area is
a stress. The strain is the percentage change in volume (change amount over
the original volume\footnote{%
I am using a $%
%TCIMACRO{\TeXButton{V}{\ooalign{\hfil$V$\hfil\cr\kern0.1em--\hfil\cr}}}%
%BeginExpansion
\ooalign{\hfil$V$\hfil\cr\kern0.1em--\hfil\cr}%
%EndExpansion
$ with a line through it for volume so it doesn't look like velocity, $v,$
or voltage, $V.$})%
\begin{equation}
\frac{\Delta 
%TCIMACRO{\TeXButton{V}{\ooalign{\hfil$V$\hfil\cr\kern0.1em--\hfil\cr}}}%
%BeginExpansion
\ooalign{\hfil$V$\hfil\cr\kern0.1em--\hfil\cr}%
%EndExpansion
}{%
%TCIMACRO{\TeXButton{V}{\ooalign{\hfil$V$\hfil\cr\kern0.1em--\hfil\cr}}}%
%BeginExpansion
\ooalign{\hfil$V$\hfil\cr\kern0.1em--\hfil\cr}%
%EndExpansion
}
\end{equation}%
and the relationship is%
\begin{equation}
\Delta P=-B\frac{\Delta 
%TCIMACRO{\TeXButton{V}{\ooalign{\hfil$V$\hfil\cr\kern0.1em--\hfil\cr}}}%
%BeginExpansion
\ooalign{\hfil$V$\hfil\cr\kern0.1em--\hfil\cr}%
%EndExpansion
}{%
%TCIMACRO{\TeXButton{V}{\ooalign{\hfil$V$\hfil\cr\kern0.1em--\hfil\cr}}}%
%BeginExpansion
\ooalign{\hfil$V$\hfil\cr\kern0.1em--\hfil\cr}%
%EndExpansion
}
\end{equation}

$B$ is the \emph{bulk modulus}, which again, we look up in tables.%
\begin{equation*}
\begin{tabular}{|l|l|}
\hline
Material\QQfnmark{$\ddag $Average values from numbers given in various text
books. These numbers should be taken as example values and more exact
numbers found for any real work.} & Bulk Modulus $\left( \unit{N}/\unit{m}%
^{2}\right) $ \\ \hline
Steel$^{\ddag }$ & $140\times 10^{9}$ \\ \hline
Cast Iron$^{\ddag }$ & $90\times 10^{9}$ \\ \hline
Brass$^{\ddag }$ & $80\times 10^{9}$ \\ \hline
Aluminum$^{\ddag }$ & $70\times 10^{9}$ \\ \hline
Water$^{\ddag }$ & $2\times 10^{9}$ \\ \hline
Ethyl Alcohol$^{\ddag }$ & $1\times 10^{9}$ \\ \hline
Mercury$^{\ddag }$ & $2.5\times 10^{9}$ \\ \hline
Air$^{\ddag }$ & $1.01\times 10^{5}$ \\ \hline
\end{tabular}%
\QQfntext{0}{$\ddag $Average values from numbers given in various text
books. These numbers should be taken as example values and more exact
numbers found for any real work.}
\end{equation*}%
The minus sign may be surprising. But in our definition of bulk modulus, we
are thinking of compression. And under compression $\Delta 
%TCIMACRO{\TeXButton{V}{\ooalign{\hfil$V$\hfil\cr\kern0.1em--\hfil\cr}}}%
%BeginExpansion
\ooalign{\hfil$V$\hfil\cr\kern0.1em--\hfil\cr}%
%EndExpansion
$ is negative. The minus sigh means that for compression this formula gives
positive values. Also notice that there is a $\Delta P$ or $\Delta F$ in our
formula. Usually we start compressing a solid while it is experiencing air
pressure. So we don't start from zero stress. We change the force from the
force due to air pressure to some larger compressive force.

Let's find the speed of the wave again using our linear wave equation. 
\FRAME{dtbpF}{3.9932in}{1.7988in}{0pt}{}{}{Figure}{\special{language
"Scientific Word";type "GRAPHIC";maintain-aspect-ratio TRUE;display
"USEDEF";valid_file "T";width 3.9932in;height 1.7988in;depth
0pt;original-width 4.7915in;original-height 2.1421in;cropleft "0";croptop
"1";cropright "1";cropbottom "0";tempfilename
'TEBZIK04.wmf';tempfile-properties "XPR";}}Let's consider a parcel of air.
What is a \textquotedblleft parcel of air?\textquotedblright\ It is just
part of the air that we will consider. It isn't different than the air
around it. The only difference is that we are thinking about this chunk of
air and not the air around it. The word \textquotedblleft
parcel\textquotedblright\ is an old word for \textquotedblleft
package.\textquotedblright\ In our case we have a cylindrical parcel of air.
Then we could use Newton's second law to find the wave equation and pick off
the speed again just like we did before. Newton's law gives

\begin{equation*}
F_{net}=F_{pr}-F_{p\ell }=\Delta F_{p}
\end{equation*}%
where now our forces are due to pressure. We can use 
\begin{equation*}
P=\frac{F}{A}
\end{equation*}%
to change this Newton's second law equation 
\begin{equation*}
F_{net}=P_{r}A-P_{\ell }A
\end{equation*}%
and we can write $P_{r}$ as just $P$ and $P_{\ell }=P+dP$ where $dP$ is the
small change in pressure from left to right. 
\begin{eqnarray*}
F_{net} &=&PA-\left( P+dP\right) A \\
&=&PA-PA-AdP=-AdP \\
&=&-AdP
\end{eqnarray*}%
Now 
\begin{equation*}
F_{net}=ma=\left( dm\right) a
\end{equation*}%
for our small a parcel of air. Then our Newton's second law equation is 
\begin{equation*}
\left( dm\right) a=PA-\left( P+dP\right) A=-AdP
\end{equation*}%
\begin{equation*}
adm=-AdP
\end{equation*}%
Let's look at our stress-strain equation for pressure stress again 
\begin{equation}
\Delta P=-B\frac{\Delta 
%TCIMACRO{\TeXButton{V}{\ooalign{\hfil$V$\hfil\cr\kern0.1em--\hfil\cr}}}%
%BeginExpansion
\ooalign{\hfil$V$\hfil\cr\kern0.1em--\hfil\cr}%
%EndExpansion
}{%
%TCIMACRO{\TeXButton{V}{\ooalign{\hfil$V$\hfil\cr\kern0.1em--\hfil\cr}}}%
%BeginExpansion
\ooalign{\hfil$V$\hfil\cr\kern0.1em--\hfil\cr}%
%EndExpansion
}
\end{equation}%
We need $dP.$ But from what did we in the last problem it seems like for 
\begin{equation*}
\frac{dP}{dx}=\frac{d}{dx}\left( \Delta P\right)
\end{equation*}%
because with our funny notation $\Delta P$ is our how much the pressure is
different from atmospheric pressure and not the change from left to right.
The notation is awkward, but this gives 
\begin{equation*}
dP=\frac{d}{dx}\left( \Delta P\right) dx
\end{equation*}

We could use $s$ again for the strain variable, only this time instead of
the stretch amount it is the squish amount So 
\begin{equation*}
\Delta 
%TCIMACRO{\TeXButton{V}{\ooalign{\hfil$V$\hfil\cr\kern0.1em--\hfil\cr}}}%
%BeginExpansion
\ooalign{\hfil$V$\hfil\cr\kern0.1em--\hfil\cr}%
%EndExpansion
=Ads
\end{equation*}
and the original volume of our parcel is 
\begin{equation*}
%TCIMACRO{\TeXButton{V}{\ooalign{\hfil$V$\hfil\cr\kern0.1em--\hfil\cr}}}%
%BeginExpansion
\ooalign{\hfil$V$\hfil\cr\kern0.1em--\hfil\cr}%
%EndExpansion
=Adx
\end{equation*}%
so our pressure equation is 
\begin{equation}
\Delta P=-B\frac{Ads}{Adx}=-B\frac{ds}{dx}
\end{equation}%
and we can put this into our equation for $dP$%
\begin{eqnarray*}
dP &=&\frac{d}{dx}\left( \Delta P\right) dx \\
&=&\frac{d}{dx}\left( -B\frac{ds}{dx}\right) dx \\
&=&-B\left( \frac{d^{2}s}{dx^{2}}\right) dx
\end{eqnarray*}%
We can put this into right hand side of our Newton's second law. 
\begin{eqnarray*}
adm &=&-AdP \\
adm &=&-A\left( -B\left( \frac{d^{2}s}{dx^{2}}\right) dx\right) \\
adm &=&AB\left( \frac{d^{2}s}{dx^{2}}\right) dx
\end{eqnarray*}%
But we can do more. Last time for the wire we used density. Let's do that
again four our parcel of air.

\begin{equation*}
dm=\rho Adx
\end{equation*}%
so 
\begin{eqnarray*}
F_{net} &=&adm \\
&=&\frac{d^{2}s}{dt^{2}}dm \\
&=&\frac{d^{2}s}{dt^{2}}\rho Adx
\end{eqnarray*}%
and we can put all this into our Newton's second law. Then%
\begin{equation*}
adm=AB\left( \frac{d^{2}s}{dx^{2}}\right) dx
\end{equation*}%
becomes 
\begin{equation*}
\frac{d^{2}s}{dt^{2}}\rho Adx=AB\left( \frac{d^{2}s}{dx^{2}}\right) dx
\end{equation*}%
and some things cancel and we can put all the constants on one side. 
\begin{equation*}
\frac{d^{2}s}{dt^{2}}=\frac{B}{\rho }\left( \frac{d^{2}s}{dx^{2}}\right)
\end{equation*}%
and this is our linear wave equation again only with $B/\rho $ where $v^{2}$
should be. That tells us that 
\begin{equation*}
v=\sqrt{\frac{B}{\rho }}
\end{equation*}%
Our elastic quantity for our medium is, $B.$ And we can still use the
density for the inertial property of the gas, $\rho .$ So our speed of sound
in a gas is 
\begin{equation*}
v=\sqrt{\frac{B}{\rho }}
\end{equation*}%
But how do we find $B?$ That is something we will do much later in this
course. So for now let's borrow a result from our future. For an ideal gas%
\begin{equation*}
v=\sqrt{\frac{\gamma RT_{K}}{M}}
\end{equation*}%
where $\gamma $ is the adiabatic index (related to the type of gas we have)
and $R$ is the universal gas constant, and $T_{K}$ is the temperature in
kelvins and $M$ is the gas molar mass. Most of these quantities are
familiar, but what is a \textquotedblleft universal gas
constant\textquotedblright\ and what is an \textquotedblleft adiabatic
index?\textquotedblright\ These are quantities that we will study at the end
of this course. So we can't complete our study of what wave speed of sound
right now. But we need an approximate way to calculate speed of sound waves.

For now, what is important is that the speed changes with temperature. The
hotter the gas, the faster the waves. If we assume the air is like an ideal
gas, this reduces to just

\begin{equation}
v=v_{o}\sqrt{1+\frac{T_{c}}{T_{o}}}
\end{equation}%
where $v_{o}=331\frac{\unit{m}}{\unit{s}}$ and $T_{o}=273\unit{K}$ ($0\unit{%
%TCIMACRO{\U{2103}}%
%BeginExpansion
{}^{\circ}{\rm C}%
%EndExpansion
}).$\footnote{$v=v_{o}\sqrt{1+\frac{T_{C}}{T_{o}}}=v_{o}\sqrt{\frac{T_{o}}{%
T_{o}}+\frac{T_{C}}{T_{o}}}=v_{o}\sqrt{\frac{T_{o}+T_{C}}{T_{o}}}=v_{o}\sqrt{%
\frac{T_{K}}{T_{o}}}$}

The speed of sound in air at normal room temperature is around $340\unit{m}/%
\unit{s}.$

Notice that we found the speed of waves using a differential equation. But
also notice that we did not solve the differential equation. We just used it
to identify the speed for each case.

\section{Intensity of Sound Waves}

We have already learned that knowing the power of a sound source isn't
enough to tell us how loud that sound source will appear to us. We know the
sound gets's spread out over a big spherical area. And we use the quantity
called intensity to describe how spread out the power is at a particular
distance.

\begin{equation}
\mathcal{I}=\frac{\mathcal{P}}{4\pi r^{2}}
\end{equation}

Let's take a parcel of air again. Only this time let's make our parcel of
air kind of square package shaped. And we wish to have a wave travel through
our parcel of air. The air molecules might both move and stretch out as the
sound wave passes. Let's use $s\left( x\right) $ to be how far the air
molecules move from equilibrium. So the displacement from equilibrium $%
s\left( x\right) $ for the air molecules might be different as the air moves
to location $x+\Delta x.$ The next figure tries to show this.

\FRAME{dtbpF}{3.7608in}{2.5448in}{0in}{}{}{Figure}{\special{language
"Scientific Word";type "GRAPHIC";maintain-aspect-ratio TRUE;display
"USEDEF";valid_file "T";width 3.7608in;height 2.5448in;depth
0in;original-width 3.8034in;original-height 2.5643in;cropleft "0";croptop
"1";cropright "1";cropbottom "0";tempfilename
'SV8RA302.wmf';tempfile-properties "XPR";}}Let's look at the stretch of the
parcel of air for a moment. Let's take the before size of the parcel of air
and compare it to the stretched parcel size.

\FRAME{dtbpF}{3.0938in}{3.2889in}{0pt}{}{}{Figure}{\special{language
"Scientific Word";type "GRAPHIC";maintain-aspect-ratio TRUE;display
"USEDEF";valid_file "T";width 3.0938in;height 3.2889in;depth
0pt;original-width 3.1248in;original-height 3.3235in;cropleft "0";croptop
"1";cropright "1";cropbottom "0";tempfilename
'SV8R3D01.wmf';tempfile-properties "XPR";}}We can see that the size
difference for the displacement from equilibrium is 
\begin{equation*}
\Delta s=s\left( x+\Delta x\right) -s\left( x\right)
\end{equation*}%
The change in volume would be 
\begin{equation*}
\Delta 
%TCIMACRO{\TeXButton{V}{\ooalign{\hfil$V$\hfil\cr\kern0.1em--\hfil\cr}}}%
%BeginExpansion
\ooalign{\hfil$V$\hfil\cr\kern0.1em--\hfil\cr}%
%EndExpansion
=A\Delta s
\end{equation*}%
The fractional change in volume is defined as 
\begin{equation*}
\frac{\Delta 
%TCIMACRO{\TeXButton{V}{\ooalign{\hfil$V$\hfil\cr\kern0.1em--\hfil\cr}}}%
%BeginExpansion
\ooalign{\hfil$V$\hfil\cr\kern0.1em--\hfil\cr}%
%EndExpansion
}{%
%TCIMACRO{\TeXButton{V}{\ooalign{\hfil$V$\hfil\cr\kern0.1em--\hfil\cr}}}%
%BeginExpansion
\ooalign{\hfil$V$\hfil\cr\kern0.1em--\hfil\cr}%
%EndExpansion
}
\end{equation*}%
and we can see that this will be 
\begin{equation*}
\frac{\Delta 
%TCIMACRO{\TeXButton{V}{\ooalign{\hfil$V$\hfil\cr\kern0.1em--\hfil\cr}}}%
%BeginExpansion
\ooalign{\hfil$V$\hfil\cr\kern0.1em--\hfil\cr}%
%EndExpansion
}{%
%TCIMACRO{\TeXButton{V}{\ooalign{\hfil$V$\hfil\cr\kern0.1em--\hfil\cr}}}%
%BeginExpansion
\ooalign{\hfil$V$\hfil\cr\kern0.1em--\hfil\cr}%
%EndExpansion
}=\frac{A\Delta s}{A\Delta x}=\frac{\Delta s}{\Delta x}
\end{equation*}%
If we take our parcel of air to be small so $\Delta x$ is a very small, $dx,$
this is just 
\begin{equation*}
\frac{d%
%TCIMACRO{\TeXButton{V}{\ooalign{\hfil$V$\hfil\cr\kern0.1em--\hfil\cr}}}%
%BeginExpansion
\ooalign{\hfil$V$\hfil\cr\kern0.1em--\hfil\cr}%
%EndExpansion
}{%
%TCIMACRO{\TeXButton{V}{\ooalign{\hfil$V$\hfil\cr\kern0.1em--\hfil\cr}}}%
%BeginExpansion
\ooalign{\hfil$V$\hfil\cr\kern0.1em--\hfil\cr}%
%EndExpansion
}=\frac{ds}{dx}
\end{equation*}

Now let's go back to our bulk modulus, $B.$ 
\begin{equation}
\Delta P=-B\frac{\Delta 
%TCIMACRO{\TeXButton{V}{\ooalign{\hfil$V$\hfil\cr\kern0.1em--\hfil\cr}}}%
%BeginExpansion
\ooalign{\hfil$V$\hfil\cr\kern0.1em--\hfil\cr}%
%EndExpansion
}{%
%TCIMACRO{\TeXButton{V}{\ooalign{\hfil$V$\hfil\cr\kern0.1em--\hfil\cr}}}%
%BeginExpansion
\ooalign{\hfil$V$\hfil\cr\kern0.1em--\hfil\cr}%
%EndExpansion
}
\end{equation}%
for our vary small change in volume we could write this as%
\begin{equation}
\Delta P=-B\frac{d%
%TCIMACRO{\TeXButton{V}{\ooalign{\hfil$V$\hfil\cr\kern0.1em--\hfil\cr}}}%
%BeginExpansion
\ooalign{\hfil$V$\hfil\cr\kern0.1em--\hfil\cr}%
%EndExpansion
}{%
%TCIMACRO{\TeXButton{V}{\ooalign{\hfil$V$\hfil\cr\kern0.1em--\hfil\cr}}}%
%BeginExpansion
\ooalign{\hfil$V$\hfil\cr\kern0.1em--\hfil\cr}%
%EndExpansion
}
\end{equation}%
and this is just 
\begin{equation*}
\Delta P=-B\frac{d%
%TCIMACRO{\TeXButton{V}{\ooalign{\hfil$V$\hfil\cr\kern0.1em--\hfil\cr}}}%
%BeginExpansion
\ooalign{\hfil$V$\hfil\cr\kern0.1em--\hfil\cr}%
%EndExpansion
}{%
%TCIMACRO{\TeXButton{V}{\ooalign{\hfil$V$\hfil\cr\kern0.1em--\hfil\cr}}}%
%BeginExpansion
\ooalign{\hfil$V$\hfil\cr\kern0.1em--\hfil\cr}%
%EndExpansion
}=-B\frac{ds}{dx}
\end{equation*}%
For a sinusoidal wave, 
\begin{equation}
s\left( x,t\right) =s_{\max }\sin \left( kx-\omega t+\phi _{o}\right)
\end{equation}%
This tells us about the distance an air molecule that was at location $x$
when $t=0$ has traveled away from that equilibrium position at some later
time, $t.$ We could find out how fast the molecule is moving by taking a
derivative.%
\begin{equation*}
v\left( x,t\right) =\frac{ds}{dx}=-ks_{\max }\cos \left( kx-\omega t+\phi
_{o}\right)
\end{equation*}%
then 
\begin{equation*}
\Delta P=-B\frac{ds}{dx}=Bks_{\max }\cos \left( kx-\omega t+\phi _{o}\right)
\end{equation*}%
We want the intensity and the intensity is 
\begin{equation}
\mathcal{I}=\frac{\mathcal{P}}{A}
\end{equation}%
We could write the power as the force multiplied by the velocity (something
we remember from Principles of Physics I).%
\begin{equation}
\mathcal{I}=\frac{Fv}{A}=Pv\left( x,t\right)
\end{equation}%
We know something about the pressure, we need the speed. But this isn't the
wave speed, it is the speed of the molecules. And we have just found this
molecule speed to be $v\left( x,t\right) =-ks_{\max }\cos \left( kx-\omega
t+\phi _{o}\right) .$ We can put this all together to get the intensity%
\begin{eqnarray}
\mathcal{I} &=&\Delta Pv \\
&=&\left( Bks_{\max }\cos \left( kx-\omega t+\phi _{o}\right) \right) \left(
-\omega s_{\max }\cos \left( kx-\omega t+\phi _{o}\right) \right) \\
&=&-Bk\omega s_{\max }^{2}\cos ^{2}\left( kx-\omega t+\phi _{o}\right)
\end{eqnarray}%
And low and behold! The intensity is proportional to the amplitude of the
displacement, $s_{\max }$ squared.

We found this to be true before for waves on strings. And now it is true as
well for sound waves (take a good guess what we will find for the intensity
and amplitude relationship for light waves!).

We often can only measure a time average intensity. To average our intensity
function over one period we would integrate%
\begin{eqnarray*}
\left\langle \mathcal{I}\right\rangle &=&\frac{1}{T}\int_{0}^{T}-Bk\omega
s_{\max }^{2}\cos ^{2}\left( kx-\omega t+\phi _{o}\right) \\
&=&-Bk\omega s_{\max }^{2}\frac{1}{T}\int_{0}^{T}\cos ^{2}\left( kx-\omega
t+\phi _{o}\right) dt \\
&=&-Bk\omega s_{\max }^{2}\frac{1}{2}
\end{eqnarray*}%
This gives the average intensity of the wave as being proportional to the
square of the maximum displacement (amplitude) $s_{\max }.$ But we learned
in our last lecture that we can describe sound waves as pressure waves. And
note that the factors $Bks_{\max }$ have units of pressure. We can call this
combination the change in pressure for the sound wave, $\Delta P_{\max
}=Bks_{\max }$ and write our average intensity in terms of this change in
pressure. 
\begin{eqnarray*}
\left\langle \mathcal{I}\right\rangle &=&-\frac{Bk\omega s_{\max }^{2}}{2}=-%
\frac{\omega }{2}\left( Bks_{\max }\right) \left( s_{\max }\right) \\
&=&-\frac{\omega }{2}\left( Bks_{\max }\right) \left( \frac{Bks_{\max }}{Bk}%
\right) \\
&=&-\frac{\omega }{2}\left( \Delta P_{\max }\right) \left( \frac{\Delta
P_{\max }}{Bk}\right) \\
&=&-\frac{\omega }{2Bk}\left( \Delta P_{\max }\right) ^{2}
\end{eqnarray*}%
and we know for sound the wave speed is given by 
\begin{equation*}
v=\frac{\omega }{k}=\sqrt{\frac{B}{\rho }}
\end{equation*}%
so 
\begin{equation*}
B=\frac{\omega ^{2}}{k^{2}}\rho
\end{equation*}%
\begin{equation*}
\left\langle \mathcal{I}\right\rangle =-\frac{\omega }{2\left( \frac{\omega
^{2}}{k^{2}}\rho \right) k}\left( \Delta P_{\max }\right) ^{2}
\end{equation*}%
\begin{equation*}
\left\langle \mathcal{I}\right\rangle =-\frac{1}{2\left( v\rho \right) }%
\left( \Delta P_{\max }\right) ^{2}
\end{equation*}%
This gives us the average intensity in the pressure picture for sound waves
and once again the intensity is proportional to the amplitude squared.

Whichever way we think about sound waves the intensity is proportional to
the amplitude squared. For those of us that will go on to study quantum
mechanics, this square of the amplitude of the wave function is going to be
very important! So chemists and physicists take note! Electrical engineers,
you will need to know this as well.

%TCIMACRO{%
%\TeXButton{Basic Equations}{\vspace{0.5in}\hspace{-1.3in}{\LARGE Basic Equations\vspace{0.25in}}}}%
%BeginExpansion
\vspace{0.5in}\hspace{-1.3in}{\LARGE Basic Equations\vspace{0.25in}}%
%EndExpansion
\begin{equation*}
v=\sqrt{\frac{B}{\rho }}
\end{equation*}%
\begin{equation}
v=v_{o}\sqrt{1+\frac{T_{c}}{T_{o}}}
\end{equation}%
\begin{equation}
\mathcal{I}=\frac{\mathcal{P}}{4\pi r^{2}}
\end{equation}%
\begin{equation*}
\left\langle \mathcal{I}\right\rangle =-\frac{1}{2}Bk\omega s_{\max }^{2}
\end{equation*}

\begin{equation*}
\left\langle \mathcal{I}\right\rangle =-\frac{1}{2\left( v\rho \right) }%
\left( \Delta P_{\max }\right) ^{2}
\end{equation*}

\chapter{8 Sound Standing Waves \ 1.17.4, 1.17.5}

%TCIMACRO{%
%\TeXButton{\chaptermark{Sound Standing Waves}}{\chaptermark{Sound Standing Waves}}}%
%BeginExpansion
\chaptermark{Sound Standing Waves}%
%EndExpansion

We now have sound waves and know a little bit about how these sound waves
act. And we know these waves deliver energy and we can even describe the
intensity of the sound waves. But it turns out our design engineer used a
clever trick in making our auditory system. He wanted us to hear the gentle
sigh of a breeze, and the thunderous roar of a jet engine and every
intensity in between. To make this happen, He made our hearing logarithmic!
We will see how this works in this lecture. We will also look at standing
waves in sound (and even light) and relate them to musical instruments and
structures.

%TCIMACRO{%
%\TeXButton{Fundamental Concepts}{\vspace{0.10in}\hspace{-0.37in}{\Large Fundamental Concepts\vspace{0.2in}}}}%
%BeginExpansion
\vspace{0.10in}\hspace{-0.37in}{\Large Fundamental Concepts\vspace{0.2in}}%
%EndExpansion

\begin{itemize}
\item Sound Level in Decibels

\item Sound Standing waves (music)

\item Light Standing Waves
\end{itemize}

\section{Sound Levels in Decibels}

The intensity of a sound wave gives us how much power per unit area is being
delivered by the wave. This is a great way to measure the power delivered by
a sound wave because most instruments have a detection area. So power per
unit area is a natural measurement for instruments.

%TCIMACRO{%
%\TeXButton{Frequency Range Demo}{\marginpar {
%\hspace{-0.5in}
%\begin{minipage}[t]{1in}
%\small{Frequency Range Demo}
%\end{minipage}
%}}}%
%BeginExpansion
\marginpar {
\hspace{-0.5in}
\begin{minipage}[t]{1in}
\small{Frequency Range Demo}
\end{minipage}
}%
%EndExpansion
\FRAME{dtbpFU}{2.3328in}{2.3647in}{0pt}{\Qcb{Robinson-Dadson equal loudness
curves (Image in the Public Domain courtesy Lindosland)}}{}{Figure}{\special%
{language "Scientific Word";type "GRAPHIC";maintain-aspect-ratio
TRUE;display "USEDEF";valid_file "T";width 2.3328in;height 2.3647in;depth
0pt;original-width 2.3487in;original-height 2.3815in;cropleft "0";croptop
"1";cropright "1";cropbottom "0";tempfilename
'T9BYZY00.wmf';tempfile-properties "XPR";}}

But our ears are not designed to be quantitative scientific instruments
(though they are truly amazing in their range and ability). Sounds with the
same intensity at different frequencies do not appear to us to have the same
loudness. The frequency response graph above shows how this relationship
works for test subjects\footnote{%
Ones that have not gone to Guns n Roses concerts.}.

Our Design Engineer made an interesting choice in building us. We need to
hear very faint sounds, and very loud sounds too. In order to make us able
to hear the soft sounds without causing extreme discomfort when we hear the
loud, He gave up linearity. That is, we don't hear twice the sound intensity
as twice as loud.

The mathematical expression that matches our perception of loudness to the
intensity is 
\begin{equation}
\beta =10\log _{10}\left( \frac{I}{I_{o}}\right)
\end{equation}

where the quantity $I_{o}$ is a reference intensity.

%TCIMACRO{%
%\TeXButton{Sound Meter Demo}{\marginpar {
%\hspace{-0.5in}
%\begin{minipage}[t]{1in}
%\small{Sound Meter Demo}
%\end{minipage}
%}}}%
%BeginExpansion
\marginpar {
\hspace{-0.5in}
\begin{minipage}[t]{1in}
\small{Sound Meter Demo}
\end{minipage}
}%
%EndExpansion
We call $\beta $ the \emph{sound level}. The quantity $I_{o}$ we choose to
be the \emph{threshold of hearing},

\begin{equation*}
I_{o}=1\times 10^{-12}\frac{\unit{W}}{\unit{m}^{2}}
\end{equation*}%
the intensity that is just barely audible for humans (it will be different
for electronic devices or other animals). Thinking back to our previous
chart you can see that $I_{o}$ is an average threshold for hearing. We
really have a different threshold sensitivity for different frequencies. But
on average $1\times 10^{-12}\frac{\unit{W}}{\unit{m}^{2}}$ is where human
hearing starts. Measured this way, we say that intensity is in units of
decibels (dB). The decibel, is useful because it can describe a non-linear
response in a linear way that is easy to match to our human experience. But
it is tricky because, like radians, it is dimensionless because we are
comparing two intensities $\left( I/I_{o}\right) .$ We expect the same
strange unit behavior with dB that we see with radians.

Suppose we double the intensity by a factor of $2.$%
\begin{eqnarray*}
\beta &=&10\log _{10}\left( \frac{2I_{o}}{I_{o}}\right) \\
&=&10\log _{10}2 \\
&=&\allowbreak 3.\,\allowbreak 010\,3dB
\end{eqnarray*}%
The sound intensity level is not twice as large, but only $3dB$ larger (and
notice that we had to supply the $dB$ ourselves, it didn't appear in the
math). It is a tiny increase. This is what we hear. A good rule to remember
is that $3dB$ corresponds to a doubling of the intensity.

The table that follows give some common sounds in units of dB. Just for
reference, I have measured a Guns n Roses concert at $120$ dB outside the
stadium.

\begin{equation*}
\begin{tabular}{|l|l|}
\hline
Sound & Sound Level (dB) \\ \hline
Jet Airplane at 30m & 140 \\ \hline
Rock Concert & 120 \\ \hline
Siren at 30m & 100 \\ \hline
Car interior when Traveling 60mi/hr & 90 \\ \hline
Street Traffic & 70 \\ \hline
Talk at 30cm & 65 \\ \hline
Whisper & 20 \\ \hline
Rustle of Leaves & 10 \\ \hline
Quietest thing we can hear (Io) & 0 \\ \hline
\end{tabular}%
\end{equation*}

In our last lecture we made standing waves with strings and ropes. We should
ask, can we make standing waves with sound waves?

\section{Sound Standing waves.}

We have spent some time studying standing waves on strings. But we also have
studied sound waves. Sound waves can experience constructive and destructive
interference. We should be able to make sound standing waves.

Let's take a pipe as shown in the next figures. The pipe will control the
direction of the sound wave. The sound wave will (mostly) go down the pipe.
But will there be a reflection at the end of the pipe like there was for the
string? It turns out that there will be a reflection! Let's consider our
sound wave as a change in pressure to see why.

For the string wave, we had a reflected pulse because the person (or guitar
bridge, or whatever is holding the string end) exerted a force on the
string. If a sound wave travels down our pipe, the pressure will change as
the wave goes down the pipe. Remember that pressure is a force spread over
an area. When the sound wave gets to the end of the pipe, the rest of the
room's air is waiting. And that air mass pushes back on the air molecules
that are trying to leave the pipe. The air mass of the room won't change
it's pressure much. That resistance to pressure change (the force due to the
rest of the room air molecules colliding with our wave molecules, pushing
them back) will send our tube molecules back down the tube. And that starts
a reflection. Think about this for a moment. We need to picture more than we
have before. The room is full of air, and it turns out that the room air
molecules aren't sitting still. So when we make our wave start by hitting
molecules with a piston, the molecules at the end of the tube will hit room
air molecules that are already moving.

\FRAME{dtbpF}{3.5129in}{2.4154in}{0pt}{}{}{Figure}{\special{language
"Scientific Word";type "GRAPHIC";maintain-aspect-ratio TRUE;display
"USEDEF";valid_file "T";width 3.5129in;height 2.4154in;depth
0pt;original-width 3.4662in;original-height 2.3748in;cropleft "0";croptop
"1";cropright "1";cropbottom "0";tempfilename
'SUS73V37.wmf';tempfile-properties "XPR";}}Think of pool balls. If the cue
ball and another ball, say, the 6 ball, collide, usually the 6 ball is
stationary at the beginning. But if it is not, that will change the outcome
of our conservation of momentum problem. If the cue ball is initially moving
toward the 6 ball, but the 6 ball is also moving toward the cue ball and
there is a collision, the cue ball will bounce back the way it came. That
will happen with some of the molecules at the end of the pipe. But some will
keep going. \FRAME{dtbpF}{3.7144in}{1.7538in}{0pt}{}{}{Figure}{\special%
{language "Scientific Word";type "GRAPHIC";maintain-aspect-ratio
TRUE;display "USEDEF";valid_file "T";width 3.7144in;height 1.7538in;depth
0pt;original-width 3.6668in;original-height 1.7158in;cropleft "0";croptop
"1";cropright "1";cropbottom "0";tempfilename
'SUS73V38.wmf';tempfile-properties "XPR";}}We will split the energy from the
wave into to two waves, much like we did with waves on strings when we came
to an interface. \FRAME{dtbpF}{2.8764in}{2.2572in}{0pt}{}{}{Figure}{\special%
{language "Scientific Word";type "GRAPHIC";maintain-aspect-ratio
TRUE;display "USEDEF";valid_file "T";width 2.8764in;height 2.2572in;depth
0pt;original-width 2.8331in;original-height 2.2174in;cropleft "0";croptop
"1";cropright "1";cropbottom "0";tempfilename
'SUS73V39.wmf';tempfile-properties "XPR";}}

Now that we have a reflected wave in the tube, we can end up with two waves
so long as the piston keeps working at making new waves. And with two waves
in the same medium, we have the possibility of having standing waves!

If we have a pipe open at both ends, we can see that air molecules are free
to move in and out of the ends of the pipes into the room (to collide with
the room molecules). If the air molecules can move, the ends must not be
nodes. In our string case, the part of the string that experienced
destructive interference and did not move was called a node and we always
had a node on the end of the string. But the pipe is different than the
string case! We can see that air molecules will move out and then bounce
back in at the end of the pipe. Still, we expect that there must be a node
somewhere. We can reasonably guess that there will be a node in the middle
of the pipe. Of course, the pressure on both ends must be atmospheric
pressure. So, remembering that pressure and displacement are $90\unit{%
%TCIMACRO{\U{b0}}%
%BeginExpansion
{{}^\circ}%
%EndExpansion
}$ out of phase for sound waves, we can guess that there are pressure nodes
on both ends. But there is a displacement anti-node at the ends of the pipe. 
\FRAME{dhF}{3.7403in}{2.2892in}{0pt}{}{}{Figure}{\special{language
"Scientific Word";type "GRAPHIC";maintain-aspect-ratio TRUE;display
"USEDEF";valid_file "T";width 3.7403in;height 2.2892in;depth
0pt;original-width 5.5054in;original-height 3.3589in;cropleft "0";croptop
"1";cropright "1";cropbottom "0";tempfilename
'SUS73V3A.wmf';tempfile-properties "XPR";}}Then for the first harmonic we
can draw a displacement node in the middle and we see that 
\begin{equation}
\lambda _{1}=2L
\end{equation}%
\FRAME{dtbpF}{3.9009in}{1.7074in}{0pt}{}{}{Figure}{\special{language
"Scientific Word";type "GRAPHIC";maintain-aspect-ratio TRUE;display
"USEDEF";valid_file "T";width 3.9009in;height 1.7074in;depth
0pt;original-width 3.9471in;original-height 1.711in;cropleft "0";croptop
"1";cropright "1";cropbottom "0";tempfilename
'T9C38Q01.wmf';tempfile-properties "XPR";}}and since $v=\lambda f$ we can
find%
\begin{equation*}
f_{1}=\frac{v}{\lambda _{1}}
\end{equation*}%
or 
\begin{equation}
f_{1}=\frac{v}{2L}
\end{equation}%
The next mode fits a whole wavelength\FRAME{dtbpF}{4.0633in}{1.9522in}{0pt}{%
}{}{Figure}{\special{language "Scientific Word";type
"GRAPHIC";maintain-aspect-ratio TRUE;display "USEDEF";valid_file "T";width
4.0633in;height 1.9522in;depth 0pt;original-width 4.112in;original-height
1.9602in;cropleft "0";croptop "1";cropright "1";cropbottom "0";tempfilename
'T9C4D504.wmf';tempfile-properties "XPR";}}%
\begin{eqnarray}
\lambda _{2} &=&L \\
f_{2} &=&\frac{v}{L}
\end{eqnarray}%
And the next mode fits a wavelength and a half\FRAME{dtbpF}{4.0916in}{%
1.9655in}{0pt}{}{}{Figure}{\special{language "Scientific Word";type
"GRAPHIC";maintain-aspect-ratio TRUE;display "USEDEF";valid_file "T";width
4.0916in;height 1.9655in;depth 0pt;original-width 4.1413in;original-height
1.9744in;cropleft "0";croptop "1";cropright "1";cropbottom "0";tempfilename
'T9C3AT03.wmf';tempfile-properties "XPR";}}%
\begin{eqnarray}
\lambda _{3} &=&\frac{3}{2}L \\
f_{3} &=&\frac{2v}{3L}
\end{eqnarray}%
If we keep going 
\begin{eqnarray}
\lambda _{n} &=&\frac{2}{n}L \\
f_{n} &=&n\frac{v}{2L}\qquad n=1,2,3,4\ldots
\end{eqnarray}%
This is the same mathematical form that we achieved for a standing wave on a
string! Pipes \emph{can} make a harmonic series!

\section{Pipes closed on one end}

%TCIMACRO{%
%\TeXButton{Boom Whacker Demo}{\marginpar {
%\hspace{-0.5in}
%\begin{minipage}[t]{1in}
%\small{Boom Whacker Demo}
%\end{minipage}
%}}}%
%BeginExpansion
\marginpar {
\hspace{-0.5in}
\begin{minipage}[t]{1in}
\small{Boom Whacker Demo}
\end{minipage}
}%
%EndExpansion
But what happens if we put a cap on one end of the pipe? The air molecules
cannot move longitudinally once they hit the cap. And where the molecules
don't move, that is a node. So our capped end must be a displacement node.

%TCIMACRO{%
%\TeXButton{Question 123.24.2}{\marginpar {
%\hspace{-0.5in}
%\begin{minipage}[t]{1in}
%\small{Question 123.24.2}
%\end{minipage}
%}} }%
%BeginExpansion
\marginpar {
\hspace{-0.5in}
\begin{minipage}[t]{1in}
\small{Question 123.24.2}
\end{minipage}
}
%EndExpansion
The open end is still a displacement antinode. Our standing wave will have a
node on the capped side, and an anti-node on the open side. It might look
like this \FRAME{dhF}{4.3474in}{2.8876in}{0pt}{}{}{Figure}{\special{language
"Scientific Word";type "GRAPHIC";maintain-aspect-ratio TRUE;display
"USEDEF";valid_file "T";width 4.3474in;height 2.8876in;depth
0pt;original-width 4.2964in;original-height 2.8444in;cropleft "0";croptop
"1";cropright "1";cropbottom "0";tempfilename
'SUS73V3C.wmf';tempfile-properties "XPR";}}In the next figure we draw the
first few harmonics for the \textquotedblleft closed on one
end\textquotedblright\ case.\FRAME{dhF}{4.1796in}{2.0358in}{0pt}{}{}{Figure}{%
\special{language "Scientific Word";type "GRAPHIC";maintain-aspect-ratio
TRUE;display "USEDEF";valid_file "T";width 4.1796in;height 2.0358in;depth
0pt;original-width 6.3252in;original-height 3.0666in;cropleft "0";croptop
"1";cropright "1";cropbottom "0";tempfilename
'SUS73V3D.wmf';tempfile-properties "XPR";}}

Once again the first harmonic for the closed pipe is found by using 
\begin{equation}
v=\lambda f
\end{equation}%
\begin{equation}
f=\frac{v}{\lambda }
\end{equation}%
We know the speed of sound, so we have $v.$ Knowing that the first harmonic
has a node at one end and an anti node at the other end gives just a quarter
of a wavelength in the pipe. If the pipe is $L$ in length, then $L$ must be 
\begin{equation}
L=\frac{1}{4}\lambda _{1}
\end{equation}%
or%
\begin{equation}
\lambda _{1}=4L
\end{equation}%
then%
\begin{equation}
f_{1}=\frac{v}{\lambda _{1}}=\frac{v}{4L}
\end{equation}

The next configuration that will have a node on one end and an antinode on
the other will have 
\begin{equation}
L=\frac{3}{4}\lambda _{2}
\end{equation}%
which gives%
\begin{equation}
\lambda _{2}=\frac{4}{3L}
\end{equation}%
and%
\begin{equation}
f_{2}=\frac{v}{\lambda _{2}}=\frac{3v}{4L}
\end{equation}

If we continued, we would find%
\begin{equation}
\lambda _{n}=\frac{4}{n}L
\end{equation}%
and 
\begin{equation}
f_{n}=n\frac{v}{4L}\qquad n=1,3,5\ldots
\end{equation}%
Notice that we got $1/4$ and $3/4$ and $5/4,$ but not $2/4$ or $6/4$ in our
equations. The even values of $n$ don't work because our boundary conditions
(one end a node and the other an anti-node) prevent a standing wave from
forming with that wavelength. It doesn't fit the pipe. So no standing wave.

Example: organ pipe

\FRAME{dhF}{1.6907in}{1.4088in}{0in}{}{}{Figure}{\special{language
"Scientific Word";type "GRAPHIC";maintain-aspect-ratio TRUE;display
"USEDEF";valid_file "T";width 1.6907in;height 1.4088in;depth
0in;original-width 1.6544in;original-height 1.3742in;cropleft "0";croptop
"1";cropright "1";cropbottom "0";tempfilename
'SUS73V3E.wmf';tempfile-properties "XPR";}}

%TCIMACRO{%
%\TeXButton{Organ Pipe Demo}{\marginpar {
%\hspace{-0.5in}
%\begin{minipage}[t]{1in}
%\small{Organ Pipe Demo}
%\end{minipage}
%}}}%
%BeginExpansion
\marginpar {
\hspace{-0.5in}
\begin{minipage}[t]{1in}
\small{Organ Pipe Demo}
\end{minipage}
}%
%EndExpansion
The organ pipe shown is closed at one end so we expect%
\begin{equation}
f_{n}=n\frac{v}{4L}\qquad n=1,3,5\ldots
\end{equation}%
Measuring the pipe, and assuming about $20\unit{%
%TCIMACRO{\U{2103}}%
%BeginExpansion
{}^{\circ}{\rm C}%
%EndExpansion
}$ for the room temperature we have 
\begin{equation}
\begin{tabular}{l}
$L=0.41\unit{m}$ \\ 
$R=0.06\unit{m}$ \\ 
$v=343\frac{\unit{m}}{\unit{s}}$%
\end{tabular}%
\end{equation}%
There is a detail we have ignored in our analysis, the width of the pipe
matters a little. I will include a fudge factor to account for this. With
the fudge factor, the wavelength is 
\begin{eqnarray*}
\lambda _{1} &=&4\left( L+0.6R\right) \\
&=&178.\,\allowbreak 4\unit{cm}
\end{eqnarray*}%
then our fundamental frequency is 
\begin{eqnarray}
f_{1} &=&\frac{v}{\lambda _{1}} \\
&=&192.\,\allowbreak 26\unit{Hz}
\end{eqnarray}

We can identify this note, and compare to a standard, like a tuning fork or
a piano to verify our prediction.\FRAME{dhF}{4.6345in}{1.8844in}{0pt}{}{}{%
Figure}{\special{language "Scientific Word";type
"GRAPHIC";maintain-aspect-ratio TRUE;display "USEDEF";valid_file "T";width
4.6345in;height 1.8844in;depth 0pt;original-width 6.3529in;original-height
2.5676in;cropleft "0";croptop "1";cropright "1";cropbottom "0";tempfilename
'SUS73V3F.wmf';tempfile-properties "XPR";}}It looks like it is just about a $%
G.$

\section{Lasers and standing waves}

We made standing waves with waves on strings and with sound waves in air.
Light is also a wave (thought we haven't shown that yet). Could we make
standing waves with light waves?

The answer is a big yes! Lasers use light standing waves. A laser is built
with two parallel mirrors that create the reflections. Remember we need two
waves in the light wave medium at the same time going different directions
to make a standing wave.

\FRAME{dhF}{2.6135in}{1.6604in}{0pt}{}{}{Figure}{\special{language
"Scientific Word";type "GRAPHIC";maintain-aspect-ratio TRUE;display
"USEDEF";valid_file "T";width 2.6135in;height 1.6604in;depth
0pt;original-width 2.5719in;original-height 1.6233in;cropleft "0";croptop
"1";cropright "1";cropbottom "0";tempfilename
'SUS73V3G.wmf';tempfile-properties "XPR";}}A flash of light from a flash
lamp starts the wave and one of the mirrors sends the wave back creating the
second wave. These two waves can mix to make a standing wave.\FRAME{dhF}{%
3.1868in}{1.6466in}{0pt}{}{}{Figure}{\special{language "Scientific
Word";type "GRAPHIC";maintain-aspect-ratio TRUE;display "USEDEF";valid_file
"T";width 3.1868in;height 1.6466in;depth 0pt;original-width
3.1427in;original-height 1.6103in;cropleft "0";croptop "1";cropright
"1";cropbottom "0";tempfilename 'SUS73V3H.wmf';tempfile-properties "XPR";}}%
Lasers add in what is called a \textquotedblleft gain
medium\textquotedblright\ which is a special crystal that generates more
light every time the wave passes through it. And that makes lasers very
bright. But you have probably noticed that lasers produce single colors of
light. That single color is related to a single frequency, the frequency
that makes the standing wave. Usually lasers produce a series of light
standing waves just like sound and strings. But our eyes can't see the other
standing wave frequencies, so we just see one.

%TCIMACRO{%
%\TeXButton{Question 123.24.3}{\marginpar {
%\hspace{-0.5in}
%\begin{minipage}[t]{1in}
%\small{Question 123.24.3}
%\end{minipage}
%}}}%
%BeginExpansion
\marginpar {
\hspace{-0.5in}
\begin{minipage}[t]{1in}
\small{Question 123.24.3}
\end{minipage}
}%
%EndExpansion
%TCIMACRO{%
%\TeXButton{Question 123.24.4}{\marginpar {
%\hspace{-0.5in}
%\begin{minipage}[t]{1in}
%\small{Question 123.24.4}
%\end{minipage}
%}}}%
%BeginExpansion
\marginpar {
\hspace{-0.5in}
\begin{minipage}[t]{1in}
\small{Question 123.24.4}
\end{minipage}
}%
%EndExpansion

\section{Standing Waves in Rods and Membranes}

We have hinted all chapter that the analysis techniques we were building
apply to structures. We need more math and computational tools to analyze
complex structures like bridges and buildings, but we can tackle a simple
structure like a rod that is clamped. The atoms in the rod can vibrate
longitudinally \FRAME{dtbpF}{4.318in}{1.2419in}{0in}{}{}{Figure}{\special%
{language "Scientific Word";type "GRAPHIC";maintain-aspect-ratio
TRUE;display "USEDEF";valid_file "T";width 4.318in;height 1.2419in;depth
0in;original-width 4.2661in;original-height 1.2081in;cropleft "0";croptop
"1";cropright "1";cropbottom "0";tempfilename
'SUS73V3I.wmf';tempfile-properties "XPR";}}Since we have motion possible on
both ends and not in the middle, we surmise that this system will have
similar solutions as did the open ended pipe.

\begin{equation*}
f_{n}=n\frac{v}{2L}
\end{equation*}%
The fundamental looks like

\FRAME{dtbpFX}{2.3443in}{1.5638in}{0pt}{}{}{Plot}{\special{language
"Scientific Word";type "MAPLEPLOT";width 2.3443in;height 1.5638in;depth
0pt;display "USEDEF";plot_snapshots TRUE;mustRecompute FALSE;lastEngine
"MuPAD";xmin "-1.5708";xmax "1.5708";xviewmin "-1.5708";xviewmax
"1.5708";yviewmin "-1.2";yviewmax "1.2";viewset"XY";rangeset"X";plottype
4;plottickdisable TRUE;axesFont "Times New
Roman,12,0000000000,useDefault,normal";numpoints 100;plotstyle
"patch";axesstyle "normal";axestips FALSE;xis \TEXUX{x};var1name
\TEXUX{$x$};function \TEXUX{$\sin \left( x\right) $};linecolor
"blue";linestyle 1;pointstyle "point";linethickness 2;lineAttributes
"Solid";var1range "-1.5708,1.5708";num-x-gridlines 100;curveColor
"[flat::RGB:0x000000ff]";curveStyle "Line";function \TEXUX{$-\sin \left(
x\right) $};linecolor "blue";linestyle 1;pointstyle "point";linethickness
2;lineAttributes "Solid";var1range "-1.5708,1.5708";num-x-gridlines
100;curveColor "[flat::RGB:0x000000ff]";curveStyle "Line";VCamFile
'SVA6OA11.xvz';valid_file "T";tempfilename
'SUS73V3J.wmf';tempfile-properties "XPR";}}

But suppose we move the clamp. The clamp forces a node where it is placed.
If we place the clap at $L/4$

\FRAME{dtbpF}{4.6025in}{1.3266in}{0in}{}{}{Figure}{\special{language
"Scientific Word";type "GRAPHIC";maintain-aspect-ratio TRUE;display
"USEDEF";valid_file "T";width 4.6025in;height 1.3266in;depth
0in;original-width 4.5506in;original-height 1.292in;cropleft "0";croptop
"1";cropright "1";cropbottom "0";tempfilename
'SUS73V3K.wmf';tempfile-properties "XPR";}}

\FRAME{dtbpFX}{2.3443in}{1.5638in}{0pt}{}{}{Plot}{\special{language
"Scientific Word";type "MAPLEPLOT";width 2.3443in;height 1.5638in;depth
0pt;display "USEDEF";plot_snapshots TRUE;mustRecompute FALSE;lastEngine
"MuPAD";xmin "-3.1459";xmax "3.1459";xviewmin "-3.1459";xviewmax
"3.1459";yviewmin "-1.2";yviewmax "1.2";viewset"XY";rangeset"X";plottype
4;plottickdisable TRUE;axesFont "Times New
Roman,12,0000000000,useDefault,normal";numpoints 100;plotstyle
"patch";axesstyle "normal";axestips FALSE;xis \TEXUX{x};var1name
\TEXUX{$x$};function \TEXUX{$\cos \left( x\right) $};linecolor
"blue";linestyle 1;pointstyle "point";linethickness 2;lineAttributes
"Solid";var1range "-3.1459,3.1459";num-x-gridlines 100;curveColor
"[flat::RGB:0x000000ff]";curveStyle "Line";function \TEXUX{$-\cos \left(
x\right) $};linecolor "blue";linestyle 1;pointstyle "point";linethickness
2;lineAttributes "Solid";var1range "-3.1459,3.1459";num-x-gridlines
100;curveColor "[flat::RGB:0x000000ff]";curveStyle "Line";VCamFile
'SVA63E01.xvz';valid_file "T";tempfilename
'SUS73V3L.wmf';tempfile-properties "XPR";}}

We can perform a similar analysis for a drum head, but it is much more
complicated. The modes are not points, but lines or curves, and the
frequencies of oscillation are not integer multiples of each other. Of
course structures can also waggle on the ends. the ends can rotate counter
to each other, etc. These are more complex modes than the longitudinal modes
we have considered.

So far, we have combined waves, but we have only combined waves with the
same frequency. But it can't be that every wave has the same frequency. So
what happens if we superimpose waves that have different frequencies? That
is what we will study in our next lecture.

%TCIMACRO{%
%\TeXButton{Basic Equations}{\vspace{0.5in}\hspace{-1.3in}{\LARGE Basic Equations\vspace{0.25in}}}}%
%BeginExpansion
\vspace{0.5in}\hspace{-1.3in}{\LARGE Basic Equations\vspace{0.25in}}%
%EndExpansion

Decibels%
\begin{equation*}
I_{o}=1\times 10^{-12}\frac{\unit{W}}{\unit{m}^{2}}
\end{equation*}

\begin{equation}
\beta =10\log _{10}\left( \frac{I}{I_{o}}\right)
\end{equation}

Standing waves, pipe open on both ends%
\begin{equation*}
f_{n}=n\frac{v}{2L}\qquad n=1,2,3,4\ldots
\end{equation*}%
Standing waves, pipe closed on one end%
\begin{equation*}
f_{n}=n\frac{v}{4L}\qquad n=1,3,5\ldots
\end{equation*}

\chapter{9 Beats and Doppler Shift 1.17.6, 1.17.7}

%TCIMACRO{%
%\TeXButton{\chaptermark{Beats and Doppler}}{\chaptermark{Beats and Doppler}}}%
%BeginExpansion
\chaptermark{Beats and Doppler}%
%EndExpansion

Up till now we only superimposed waves that had the same frequency. But what
happens if we take waves with different frequencies? Let's answer this first
in this lecture. Then, we should worry about something from Principles of
Physics I. We know speeds are relative. But our equations for things like
standing waves involve the speed of the wave. Shouldn't these be relative
speeds viewed from a reference frame? How do we deal with different relative
speeds from different reference frames? We will take on this question in
this lecture.

%TCIMACRO{%
%\TeXButton{Fundamental Concepts}{\vspace{0.10in}\hspace{-0.37in}{\Large Fundamental Concepts\vspace{0.2in}}}}%
%BeginExpansion
\vspace{0.10in}\hspace{-0.37in}{\Large Fundamental Concepts\vspace{0.2in}}%
%EndExpansion

\begin{itemize}
\item If we mix waves of different frequencies we get a periodic change in
loudness

\item If we move the wave source or detector the relative motion causes a
change in detected frequency
\end{itemize}

\section{Beats}

%TCIMACRO{%
%\TeXButton{Beat Demo}{\marginpar {
%\hspace{-0.5in}
%\begin{minipage}[t]{1in}
%\small{Beat Demo}
%\end{minipage}
%}} }%
%BeginExpansion
\marginpar {
\hspace{-0.5in}
\begin{minipage}[t]{1in}
\small{Beat Demo}
\end{minipage}
}
%EndExpansion
Let's mix two waves again, but this time take the case where $\phi _{2}=\phi
_{1}=\phi _{o},$ $k_{2}=k_{1}=k$ and let's let the waves mix at the same
location $x_{2}=x_{1}$ so 
\begin{equation*}
y_{1}=y_{\max }\sin \left( kx-\omega _{1}t+\phi _{o}\right)
\end{equation*}%
\begin{equation*}
y_{2}=y_{\max }\sin \left( kx-\omega _{2}t+\phi _{o}\right)
\end{equation*}%
that is, let's let the frequencies be different so that $\omega _{1}\neq
\omega _{2}.$ Our (total) phase difference would be 
\begin{eqnarray*}
\Delta \phi &=&\left( kx-\omega _{1}t+\phi _{o}\right) -\left( kx-\omega
_{1}t+\phi _{o}\right) \\
&=&\left( \omega _{1}-\omega _{2}\right) t
\end{eqnarray*}%
We can use our criteria for constructive and destructive interference, but
before going on let's put this back into the total amplitude function for
two wave mixing. 
\begin{eqnarray*}
A &=&2y_{\max }\cos \left( \frac{1}{2}\Delta \phi \right) \\
&=&2y_{\max }\cos \left( \frac{1}{2}\left( \omega _{1}-\omega _{2}\right)
t\right)
\end{eqnarray*}%
and we can see that our amplitude will change in time. For a given position
sometimes it will be constructive interference and sometimes destructive
interference and often in between.

We can plot both waves on the same graph and see that this will happen.%
\FRAME{dtbpFX}{4.9018in}{2.29in}{0pt}{}{}{Plot}{\special{language
"Scientific Word";type "MAPLEPLOT";width 4.9018in;height 2.29in;depth
0pt;display "USEDEF";plot_snapshots TRUE;mustRecompute FALSE;lastEngine
"MuPAD";xmin "-15";xmax "15";xviewmin "-8";xviewmax "8";yviewmin
"-2.5";yviewmax "2.5";viewset"XY";rangeset"X";plottype 4;labeloverrides
1;x-label "t";axesFont "Times New
Roman,12,0000000000,useDefault,normal";numpoints 100;plotstyle
"patch";axesstyle "normal";axestips FALSE;xis \TEXUX{x};var1name
\TEXUX{$x$};function \TEXUX{$\cos \left( 10x\right) $};linecolor
"maroon";linestyle 1;pointstyle "point";linethickness 3;lineAttributes
"Solid";var1range "-15,15";num-x-gridlines 900;curveColor
"[flat::RGB:0x00800000]";curveStyle "Line";rangeset"X";function \TEXUX{$\cos
\left( 11.x\right) $};linecolor "green";linestyle 1;pointstyle
"point";linethickness 3;lineAttributes "Solid";var1range
"-10,10";num-x-gridlines 900;curveColor "[flat::RGB:0x00008000]";curveStyle
"Line";rangeset"X";VCamFile 'SV8TFY0G.xvz';valid_file "T";tempfilename
'SUS73V3Y.wmf';tempfile-properties "XPR";}}Notice that there are places
where the waves are in phase, and places where they are not. The
superposition looks like this\FRAME{dtbpFX}{5.0375in}{2.3531in}{0pt}{}{}{Plot%
}{\special{language "Scientific Word";type "MAPLEPLOT";width 5.0375in;height
2.3531in;depth 0pt;display "USEDEF";plot_snapshots TRUE;mustRecompute
FALSE;lastEngine "MuPAD";xmin "-10";xmax "10";xviewmin "-8";xviewmax
"8";yviewmin "-5";yviewmax "5";viewset"XY";rangeset"X";plottype
4;labeloverrides 1;x-label "t";axesFont "Times New
Roman,12,0000000000,useDefault,normal";numpoints 100;plotstyle
"patch";axesstyle "normal";axestips FALSE;xis \TEXUX{x};var1name
\TEXUX{$x$};function \TEXUX{$\cos \left( 10.0x\right) +\cos \left(
11.x\right) $};linecolor "cyan";linestyle 1;pointstyle "point";linethickness
3;lineAttributes "Solid";var1range "-10,10";num-x-gridlines 900;curveColor
"[flat::RGB:0x00008080]";curveStyle "Line";rangeset"X";VCamFile
'SV8TFY0F.xvz';valid_file "T";tempfilename
'SUS73V3Z.wmf';tempfile-properties "XPR";}}

Where there is constructive interference, the resulting wave amplitude is
large, where there is destructive interference, the resulting amplitude is
zero. We get a traveling wave who's amplitude varies. We can find the
amplitude function algebraically.

We can write out the entire resultant wave in our usually way. Our two waves
are%
\begin{equation*}
y_{1}=y_{\max }\sin \left( kx-\omega _{1}t+\phi _{o}\right)
\end{equation*}%
\begin{equation*}
y_{2}=y_{\max }\sin \left( kx-\omega _{2}t+\phi _{o}\right)
\end{equation*}

and the resultant 
\begin{eqnarray*}
y_{r} &=&2y_{\max }\cos \left( \frac{\left( kx-\omega _{2}t+\phi _{o}\right)
-\left( kx-\omega _{1}t+\phi _{o}\right) }{2}\right) \sin \left( \frac{%
kx-\omega _{2}t+\phi _{o}+kx-\omega _{1}t+\phi _{o}}{2}\right) \\
&=&2y_{\max }\cos \left( \frac{kx-2\pi f_{2}t-\left( kx-2\pi f_{1}t\right) }{%
2}\right) \sin \left( \frac{kx-2\pi f_{2}t+kx-2\pi f_{1}t+2\phi _{o}}{2}%
\right) \\
&=&2y_{\max }\cos \left( 2\pi \frac{f_{1}-f_{2}}{2}t\right) \sin \left(
kx-2\pi \frac{f_{1}+f_{2}}{2}t+\phi _{o}\right) \\
&=&\left[ 2y_{\max }\cos \left( 2\pi \frac{f_{1}-f_{2}}{2}t\right) \right]
\sin \left( kx-2\pi \frac{f_{1}+f_{2}}{2}t+\phi _{o}\right)
\end{eqnarray*}

The sine part is a wave, it is a function of position and time. We see that
it has a frequency that is the average of $f_{1}$ and $f_{2}.$ This is the
frequency we hear. But we have another complicated amplitude term, and this
time it is a function of time just as we suspected. The amplitude has its
own frequency that is half the difference of $f_{1}$ and $f_{2}.$ 
\begin{equation*}
A_{\text{resultant}}=2A\cos \left( 2\pi \frac{f_{1}-f_{2}}{2}t\right)
\end{equation*}%
So the sound amplitude will vary in time for a given position in the medium.

The situation is odder still. We have a cosine function, but it is really an
envelope for the higher frequency motion of the air particles. The air
molecules move back and forth for both the crest and the trough of the
envelope function. \FRAME{dtbpFX}{5.0375in}{2.3531in}{0pt}{}{}{Plot}{\special%
{language "Scientific Word";type "MAPLEPLOT";width 5.0375in;height
2.3531in;depth 0pt;display "USEDEF";plot_snapshots TRUE;mustRecompute
FALSE;lastEngine "MuPAD";xmin "-10";xmax "10";xviewmin "-10";xviewmax
"10";yviewmin "-5";yviewmax "5";viewset"XY";rangeset"X";plottype
4;labeloverrides 1;x-label "t";axesFont "Times New
Roman,12,0000000000,useDefault,normal";numpoints 100;plotstyle
"patch";axesstyle "normal";axestips FALSE;xis \TEXUX{x};var1name
\TEXUX{$x$};function \TEXUX{$\cos \left( 10.0x\right) +\cos \left(
11.x\right) $};linecolor "green";linestyle 1;pointstyle
"point";linethickness 1;lineAttributes "Solid";var1range
"-10,10";num-x-gridlines 900;curveColor "[flat::RGB:0x00008000]";curveStyle
"Line";rangeset"X";function \TEXUX{$2\cos (\frac{10-11}{2}x)$};linecolor
"blue";linestyle 1;pointstyle "point";linethickness 3;lineAttributes
"Solid";var1range "-10,10";num-x-gridlines 100;curveColor
"[flat::RGB:0x000000ff]";curveStyle "Line";VCamFile
'SV8TFX0E.xvz';valid_file "T";tempfilename
'SUS73V40.wmf';tempfile-properties "XPR";}}So we will hear two maxima for
every period! This frequency with which we hear the sound get loud at a
given location as the wave goes by is called the \emph{beat frequency}. The
red envelope (solid heavy line in the last figure) has a frequency of%
\begin{equation*}
f_{A}=\frac{f_{1}-f_{2}}{2}
\end{equation*}%
but the cosine part is just the envelope. We can see that the green (thin
line) wave will push and pull air molecules, and therefore our ear drums,
with maximum loudness at twice this frequency. So our beat frequency is 
\begin{equation*}
f_{beat}=\left\vert f_{1}-f_{2}\right\vert
\end{equation*}

Note the absolute value sign. It doesn't matter if $f_{1}$ or if $f_{2}$
bigger. Any time we mix waves of different frequencies we get beating. Often
the beat frequency is very fast, and our hearing system can't track the
beats, so we don't hear them. And if we mixed more than two waves the beats
might not come at easily recognizable intervals. The mixing of waves can
become quite complicated. Yet even a barbershop quartet is a mixing of at
least four waves. So complicated superpositions are common. In the next
lecture we will try to see how we could take on these complicated combined
waves.

\section{Doppler Effect}

Let's start by considering an inertial reference frame (remember this from
Principles of Physics I?)

Suppose we pick two reference frames, one traveling with a velocity $v_{r}$
with respect to the other. Let's also place them far away from any other
object.

\FRAME{dtbpF}{3.2721in}{1.2675in}{0in}{}{}{Figure}{\special{language
"Scientific Word";type "GRAPHIC";maintain-aspect-ratio TRUE;display
"USEDEF";valid_file "T";width 3.2721in;height 1.2675in;depth
0in;original-width 3.3058in;original-height 1.2631in;cropleft "0";croptop
"1";cropright "1";cropbottom "0";tempfilename
'TEKOD500.wmf';tempfile-properties "XPR";}}

Person $A$ sees himself as stationary and sees person $B$ traveling with
velocity $v_{BA}.$ Person $B$ sees himself as stationary, and person $A$
traveling with velocity $-v_{BA}.$ In looking at this situation it is the 
\emph{relative} speed $v_{BA}$ that we must consider. We recall that we
could write the speed of guy $b$ as seen from platform $A$ as 
\begin{equation*}
v_{bA}=v_{bB}+v_{BA}
\end{equation*}%
where $v_{bB}$ is guy $b$'s speed seen from platform $B$ and $v_{BA}$ is the
speed of platform $B$ as seen from platform $A.$ Now suppose we have a wave
generator (a point source) creating spherical waves. Let the point source be
at rest. We will call this point source an emitter and use a subscript $e$
for it. \FRAME{dtbpF}{3.0467in}{1.7322in}{0pt}{}{}{Figure}{\special{language
"Scientific Word";type "GRAPHIC";maintain-aspect-ratio TRUE;display
"USEDEF";valid_file "T";width 3.0467in;height 1.7322in;depth
0pt;original-width 4.0672in;original-height 2.3004in;cropleft "0";croptop
"1";cropright "1";cropbottom "0";tempfilename
'SV8THK05.wmf';tempfile-properties "XPR";}}Let's also assume a detector. If
the detector is stationary with respect to the emitter so the detector is in
the emitter's reference frame, it sees a frequency of the wave in the
emitter frame, $f_{we}.$ But lets have the detector move relative to the
emitter.

\FRAME{dtbpF}{2.5495in}{2.45in}{0in}{}{}{Figure}{\special{language
"Scientific Word";type "GRAPHIC";maintain-aspect-ratio TRUE;display
"USEDEF";valid_file "T";width 2.5495in;height 2.45in;depth
0in;original-width 2.5088in;original-height 2.4085in;cropleft "0";croptop
"1";cropright "1";cropbottom "0";tempfilename
'SV8THK06.wmf';tempfile-properties "XPR";}}Remember, that the frequency is
the number of crests that pass by a given point in a unit time. Does the
moving detector see the same number of crests per unit time as when it was
stationary?

No, the frequency appears to be higher. Every time the detector hits a wave
crest it moves toward the next wave crest so that it hits the wave crests
faster. That is a higher frequency. How about if we let the detector move
the other way?

\FRAME{dtbpF}{2.1223in}{2.0384in}{0pt}{}{}{Figure}{\special{language
"Scientific Word";type "GRAPHIC";maintain-aspect-ratio TRUE;display
"USEDEF";valid_file "T";width 2.1223in;height 2.0384in;depth
0pt;original-width 2.0833in;original-height 2.0003in;cropleft "0";croptop
"1";cropright "1";cropbottom "0";tempfilename
'SV8THK07.wmf';tempfile-properties "XPR";}}Again the frequency seen by the
detector is different, but this time lower. Every time the detector hits a
wave crest the detector moves away from the next wave crest so it takes
longer for that wave crest to catch up to it. That makes a longer period and
a lower frequency.

We can quantify these changes. Take our usual variables $f_{we},$ as the
frequency of the wave in the emitter frame (where the emitter is
stationary). Then $\lambda _{we}$ would be the wavelength in the emitter
frame and in the emitter frame the velocity of sound would be $%
v_{sound}=v_{we}.$ But now let's move our detector. First let's move the
detector toward the source. The detector sees the wave velocity as 
\begin{equation}
v_{wd}=v_{we}+v_{de}
\end{equation}%
where $v_{wd}$ is the velocity of the wave in the detector frame, and $%
v_{de} $ is the velocity of the detector in the emitter frame. Note this is
just our equation for relative motion! And to be clear, for this case the
detector is moving and we are in the emitter reference frame so for us the
emitter is not moving.

The wavelength will not be changed, and so our wave speed in the emitter
frame 
\begin{equation*}
v_{we}=\lambda _{we}f_{we}
\end{equation*}%
can just be written as 
\begin{equation*}
v_{we}=\lambda f_{we}
\end{equation*}%
and let's calculate the frequency of the wave in the detector frame. It
would still be true that 
\begin{equation*}
v_{wd}=f_{wd}\lambda _{wd}=f_{wd}\lambda
\end{equation*}%
because the wavelength is the same in both frames. So then%
\begin{equation*}
f_{wd}=\frac{v_{wd}}{\lambda }
\end{equation*}%
and we know what $v_{wd}$ is! 
\begin{equation*}
f_{wd}=\frac{v_{we}+v_{de}}{\lambda }
\end{equation*}%
We can eliminate $\lambda $ from this expression for $f_{wd}$ by using $%
v_{we}=\lambda f_{we}$ again, this time solving for $\lambda $ we get%
\begin{equation*}
\lambda =\frac{v_{we}}{f_{we}}
\end{equation*}%
and we substitute this into our $f_{wd}$ equation%
\begin{equation*}
f_{wd}=\frac{v_{we}+v_{de}}{\frac{v_{we}}{f_{we}}}
\end{equation*}%
Or, after rearranging%
\begin{equation}
f_{wd}=\frac{v_{we}+v_{de}}{v_{we}}f_{we}\text{ \qquad }%
\begin{tabular}{l}
detector moving toward the emitter \\ 
emitter not moving%
\end{tabular}%
\end{equation}%
or recognizing that $v_{we}$ is the speed of sound in the stationary frame
of the emitter 
\begin{equation}
f_{wd}=\frac{v_{sound}+v_{de}}{v_{sound}}f_{we}\text{ \qquad }%
\begin{tabular}{l}
detector moving toward the emitter \\ 
emitter not moving%
\end{tabular}%
\end{equation}%
Let's stop to think about this. What we have found is that the frequency
measured by the detector is different than the frequency we measure while
sitting next to the emitter. And what makes the difference is the relative
speed $v_{de}.$

Now if the detector is going the other way the speed of the wave seen by the
detector is 
\begin{equation*}
v_{wd}=v_{we}-v_{de}
\end{equation*}%
and the same reasoning we used above gives%
\begin{equation}
f_{wd}=\frac{v_{we}-v_{de}}{v_{we}}f_{we}\text{ \qquad }%
\begin{tabular}{l}
detector moving toward the emitter \\ 
emitter not moving%
\end{tabular}%
\end{equation}%
or 
\begin{equation}
f_{wd}=\frac{v_{sound}-v_{de}}{v_{sound}}f_{we}\text{ \qquad }%
\begin{tabular}{l}
detector moving toward the emitter \\ 
emitter not moving%
\end{tabular}%
\end{equation}

From our thinking about the motion of two inertial reference frames, we
should expect a similar situation if the detector is stationary and the
source moves. The relative motion should be the same. So we have to get a
frequency change. But think, from the detector frame the detector will see a
different wavelength. The emitter makes a wave crest, then moves toward that
crest while it makes the next crest. So the distance between the crests
get's smaller.

\FRAME{dtbpF}{1.5307in}{1.4754in}{0pt}{}{}{Figure}{\special{language
"Scientific Word";type "GRAPHIC";maintain-aspect-ratio TRUE;display
"USEDEF";valid_file "T";width 1.5307in;height 1.4754in;depth
0pt;original-width 1.4944in;original-height 1.4399in;cropleft "0";croptop
"1";cropright "1";cropbottom "0";tempfilename
'SV8THK08.wmf';tempfile-properties "XPR";}}In fact, when we measure the
distance between the crests we must account for the fact that in the emitter
frame the source moved by an amount%
\begin{equation*}
\Delta x_{ed}=v_{ed}T_{we}=\frac{v_{ed}}{f_{we}}
\end{equation*}%
during one period of oscillation of the emitter before making the next
crest. We can see that in this case the speed of the wave in the detector
frame, $v_{wd}=v_{sound},$ because this time the air is not moving with
respect to the detector, so they are in the same reference frame. Then the
wavelength in the detector frame, $\lambda _{wd}$, will be shorter by $%
\Delta x_{ed}$. That is, if we take the wavelength that we would get if the
emitter were at rest, and subtract $\Delta x_{ed}$ we should have our new
wavelength at the detector. Let's start by finding the wavelength we expect
if the emitter were at rest in the detector frame $(v_{ed}=0)$. Then 
\begin{equation*}
\left. \lambda _{we}\right\vert _{v_{ed}=0}=\frac{v_{we}}{f_{we}}=\frac{%
v_{sound}}{f_{we}}
\end{equation*}%
because if the emitter were at rest, then $v_{we}=v_{sound}.$ But in our
case the emitter is moving, so we must subtract $\Delta x_{ed}$ from the
rest wavelength.%
\begin{equation*}
\lambda _{wd}=\left. \lambda _{we}\right\vert _{v_{ed}=0}-\Delta x_{ed}
\end{equation*}%
Then the wavelength we would see at the detector would be%
\begin{equation*}
\lambda _{wd}=\left. \lambda _{we}\right\vert _{v_{ed}=0}-\frac{v_{ed}}{%
f_{we}}
\end{equation*}

Using 
\begin{equation*}
\lambda _{wd}=\frac{v_{wd}}{f_{wd}}
\end{equation*}%
once more, we can write the frequency in the detector frame as 
\begin{equation*}
f_{wd}=\frac{v_{wd}}{\lambda _{wd}}=\frac{v_{wd}}{\left. \lambda
_{we}\right\vert _{v_{ed}=0}-\frac{v_{ed}}{f_{we}}}
\end{equation*}%
and we can substitute in $\left. \lambda _{we}\right\vert _{v_{ed}=0}=\frac{%
v_{sound}}{f_{we}}$ 
\begin{eqnarray*}
f_{wd} &=&\frac{v_{wd}}{\frac{v_{sound}}{f_{we}}-\frac{v_{ed}}{f_{we}}} \\
&=&\frac{v_{wd}}{v_{sound}-v_{ed}}f_{we}
\end{eqnarray*}%
So then we have%
\begin{equation}
f_{wd}=\frac{v_{wd}}{v_{sound}-v_{ed}}f_{we}\qquad 
\begin{tabular}{l}
$\text{emitter moving toward detector}$ \\ 
detector not moving%
\end{tabular}%
\end{equation}%
and sure enough the frequency is different and we see the reason is the
relative speed $v_{ed}$ again.

Now recall that it is the detector that is stationary with respect to the
air this time so $v_{sound}=v_{wd}.$ Then we can write $f_{wd}$ as 
\begin{equation}
f_{wd}=\frac{v_{sound}}{v_{sound}-v_{ed}}f_{we}\qquad 
\begin{tabular}{l}
$\text{emitter moving toward detector}$ \\ 
detector not moving%
\end{tabular}%
\end{equation}%
When the source is moving away from the detector, \FRAME{dtbpF}{1.8507in}{%
1.7841in}{0pt}{}{}{Figure}{\special{language "Scientific Word";type
"GRAPHIC";maintain-aspect-ratio TRUE;display "USEDEF";valid_file "T";width
1.8507in;height 1.7841in;depth 0pt;original-width 1.8135in;original-height
1.7469in;cropleft "0";croptop "1";cropright "1";cropbottom "0";tempfilename
'SV8THK09.wmf';tempfile-properties "XPR";}}we expect the wavelength to be
larger. The same reasoning gives

\begin{equation}
f_{wd}=\frac{v_{wd}}{v_{wd}+v_{ed}}f_{we}\qquad 
\begin{tabular}{l}
$\text{emitter moving away from detector}$ \\ 
detector not moving%
\end{tabular}%
\end{equation}%
or%
\begin{equation}
f_{wd}=\frac{v_{sound}}{v_{sound}+v_{ed}}f_{we}\qquad 
\begin{tabular}{l}
$\text{emitter moving away from detector}$ \\ 
detector not moving%
\end{tabular}%
\end{equation}%
These are great equations to describe Doppler effect. We can see that
relative speed did matter for our waves!

\subsection{Combined Doppler Equation}

We can combine these formulae to make one expression, but to do so we need
to remember what $v_{de}$ and $v_{ed}$ mean. The first was the speed of the
detector as viewed from the emitter frame (where the emitter is not moving).
The second was the speed of the emitter as viewed from the detector frame
(where the detector is not moving). But we are experienced with relative
motion.\ Let's envision a reference frame that is not tied to either the
emitter or detector. In this reference frame $v_{dR}$ is the speed of the
detector, and $v_{eR}$ is the speed of the emitter. In this $R$ reference
frame our first Doppler equation for a moving detector with a stationary
emitter might be written as%
\begin{equation}
f_{wd}=\frac{v_{sound}\pm v_{dR}}{v_{sound}}f_{we}\text{ \qquad detector
moving emitter stationary in R frame}
\end{equation}%
and our second equation for a moving emitter with a stationary detector
might be written as%
\begin{equation}
f_{wd}=\frac{v_{wd}}{v_{wd}\mp v_{eR}}f_{we}\qquad \text{emitter moving
detector stationary in R frame}
\end{equation}%
We could, of course have both the detector and emitter moving in the $R$
frame. This would combine both of our previous scenarios 
\begin{equation}
f_{wd}=\frac{v_{sound}\pm v_{dR}}{v_{sound}\mp v_{eR}}f_{we}
\end{equation}%
where we use the top sign for the speed when the mover (emitter of detector)
is going toward the non-mover.

With this view, that we are neither in the $e$ frame nor the $d$ frame but
in a separate $R$ frame, we could even drop the $R$ subscript without too
much confusion, and we could drop the $w$ for wave as well.%
\begin{equation}
f_{d}=\frac{v_{sound}\pm v_{d}}{v_{sound}\mp v_{e}}f_{e}
\end{equation}%
where now $v_{d}$ is the speed of the detector in the reference frame of the
observer, and $v_{e}$ is the speed of the emitter in the reference frame of
the observer. The quantity $f_{e}$ is the frequency of the wave as seen by
the emitter and $f_{d}$ is the frequency of the wave recorded by the
detector.

\subsection{Doppler effect in light}

Light is also a wave, and so we would expect a Doppler shift in light.
Indeed we do see a Doppler shift when we look at moving objects. Here is an
optical spectrum of the Sun on the top and a spectrum of a similar star
moving away from us in the middle. The final spectrum is for a star moving
toward us.\FRAME{dtbpFU}{1.5757in}{1.3785in}{0pt}{\Qcb{{\protect\small Top:
Normal `dark' spectral line positions at rest. Middle: Source moving away
from observer. Bottom: Source moving towards observer. (Public domain image
courtesy NASA: http://www.jwst.nasa.gov/education/7Page45.pdf)}}}{}{Figure}{%
\special{language "Scientific Word";type "GRAPHIC";maintain-aspect-ratio
TRUE;display "USEDEF";valid_file "T";width 1.5757in;height 1.3785in;depth
0pt;original-width 1.5385in;original-height 1.343in;cropleft "0";croptop
"1";cropright "1";cropbottom "0";tempfilename
'SV8THK0A.wmf';tempfile-properties "XPR";}}Note that the wavelength of the
lines is shifted toward the red part of the spectrum when the glowing object
moves away from us. This is equivalent to lowering of the frequency of a
truck engine noise as it goes away from us. The larger wavelengths indicate
a lower frequency of light because 
\begin{equation*}
f=\frac{c}{\lambda }
\end{equation*}%
This gives us a way to determine if distant stars and galaxies are moving
toward or away from us. We look for the chemical signature pattern of lines,
then see whether they are shifted to the red (moving away from us) or blue
(moving toward us) compared to the position in their spectrum of the Sun.
This photo is of some of the most distant galaxies that are moving very fast
away from us. Their redshift is very large.\FRAME{dhFU}{2.3748in}{2.6671in}{%
0pt}{\Qcb{{\protect\small High Redshift Galaxy Cluster shown here in false
color from the Spitzer Space Telescope. (Public domain image courtesy
NASA/JPL-Caltech/S.A. Stanford (UC Davis/LLNL)}}}{}{Figure}{\special%
{language "Scientific Word";type "GRAPHIC";maintain-aspect-ratio
TRUE;display "USEDEF";valid_file "T";width 2.3748in;height 2.6671in;depth
0pt;original-width 2.3359in;original-height 2.623in;cropleft "0";croptop
"1";cropright "1";cropbottom "0";tempfilename
'SV8THK0B.wmf';tempfile-properties "XPR";}}

Deriving the Doppler equation for light is more tricky because the speed of
light is the same in every reference frame. We really tackle this in our
Principles of Physics IV or \textquotedblleft Modern
Physics\textquotedblright\ class (PH279). So I\ will just quote the result
here.

\begin{eqnarray}
\lambda _{-} &=&\lambda _{o}\sqrt{\frac{1+\frac{v}{c}}{1-\frac{v}{c}}}\text{ 
}\qquad \text{Receding source} \\
\lambda _{-} &=&\lambda _{o}\sqrt{\frac{1-\frac{v}{c}}{1+\frac{v}{c}}}\text{ 
}\qquad \text{Approaching source}
\end{eqnarray}%
For our chemists, think of a group of gas molecules in a sample in a lab.
The gas molecules will be moving randomly due to thermal energy. Some will
be moving closer to us and some away from us. This means the light emitted
from some of the gas molecules will be red shifted and some will be blue
shifted. This affects the spectroscopy we use to identify the molecules!

%TCIMACRO{%
%\TeXButton{Basic Equations}{\vspace{0.5in}\hspace{-1.3in}{\LARGE Basic Equations\vspace{0.25in}}}}%
%BeginExpansion
\vspace{0.5in}\hspace{-1.3in}{\LARGE Basic Equations\vspace{0.25in}}%
%EndExpansion

For beats%
\begin{equation*}
y=\left[ 2y_{\max }\cos \left( 2\pi \frac{f_{1}-f_{2}}{2}t\right) \right]
\sin \left( kx-2\pi \frac{f_{1}+f_{2}}{2}t\right)
\end{equation*}%
\begin{equation*}
A_{\text{resultant}}=2y_{\max }\cos \left( 2\pi \frac{f_{1}-f_{2}}{2}t\right)
\end{equation*}%
\begin{equation*}
f_{beat}=\left\vert f_{1}-f_{2}\right\vert
\end{equation*}

Doppler%
\begin{equation}
f_{wd}=\frac{v_{sound}\pm v_{dR}}{v_{sound}}f_{we}\text{ \qquad detector
moving emitter stationary in R frame}
\end{equation}%
\begin{equation}
f_{wd}=\frac{v_{wd}}{v_{wd}\mp v_{eR}}f_{we}\qquad \text{emitter moving
detector stationary in R frame}
\end{equation}
\begin{equation}
f_{wd}=\frac{v_{sound}\pm v_{dR}}{v_{sound}\mp v_{eR}}f_{we}
\end{equation}

\chapter{10 Shock waves and Non-Sinusoidal Waves 1.17.8}

%TCIMACRO{\TeXButton{\chaptermark{Shock Waves}}{\chaptermark{Shock Waves}}}%
%BeginExpansion
\chaptermark{Shock Waves}%
%EndExpansion

%TCIMACRO{%
%\TeXButton{Fundamental Concepts}{\vspace{0.10in}\hspace{-0.37in}{\Large Fundamental Concepts\vspace{0.2in}}}}%
%BeginExpansion
\vspace{0.10in}\hspace{-0.37in}{\Large Fundamental Concepts\vspace{0.2in}}%
%EndExpansion

\begin{itemize}
\item If the source travels faster than the wave speed in the medium you get
shock waves

\item Actual waves can be viewed as sums of sinusoidal waves
\end{itemize}

\section{Shock Waves}

What happens when the speed of the source is greater than the wave speed?

Remember that the wave speed depends only on the medium. Let's call the
crests of a wave the \emph{wave front}. In the picture below, a point source
is generating a wave and the red lines are the wave fronts. The red dots are
the position of the point source when it made the wave. Each red dot has a
red wave front around it. But the earlier positions were, of course, made
first. So for position $x_{o}$ the wave has had more time to travel. Thus,
the wave front for $x_{o}$ is a bigger sphere The source moved to position $%
x_{1}$ and made another wave. That wave hasn't had as long to travel, so
it's sphere is smaller. The source also made waves at positions $x_{2},$ $%
x_{3,}$ and it has just barley made a wave at position $x_{4.}$

But notice that when the speed of the emitter is equal to or grater than the
speed of sound $\left( v_{ed}\geq v_{sound}\right) $ the waves begin to pile
up. If we allow $v_{ed}>v_{sound}$ then the wave fronts are no longer
generated within each other. \FRAME{dtbpF}{3.1263in}{3.1055in}{0in}{}{}{%
Figure}{\special{language "Scientific Word";type
"GRAPHIC";maintain-aspect-ratio TRUE;display "USEDEF";valid_file "T";width
3.1263in;height 3.1055in;depth 0in;original-width 3.0813in;original-height
3.0614in;cropleft "0";croptop "1";cropright "1";cropbottom "0";tempfilename
'TEOBYS00.wmf';tempfile-properties "XPR";}}The leading edge of the wave
fronts \textquotedblleft build up\textquotedblright\ to form a cone shape.
We recognize this as a superposition of the waves. This is constructive
interference. And constructive interference makes the resulting wave
amplitude bigger! Since these are sound waves, we get a very loud boom when
the superimposed waves reach our ears! This is a sonic boom.

The half angle of this cone is called the \emph{Mach angle} 
\begin{equation}
\sin \theta =\frac{v_{sound}t}{v_{ed}t}=\frac{v_{sound}}{v_{ed}}
\end{equation}%
This ratio $v_{ed}/v_{sound}$ is called the Mach number 
\begin{equation*}
M=\frac{v_{ed}}{v_{sound}}
\end{equation*}%
and the conical wave front is called a shock wave. 
\begin{equation*}
\sin \theta =\frac{1}{M}
\end{equation*}%
We see these conical waves due to constructive interference often in water
(where we call them a wake)

\FRAME{dhFU}{1.9558in}{1.3402in}{0pt}{\Qcb{{\protect\small Boat wakes as a
Doppler cone. Image courtacy US\ Navy. Image is in the Public Domain.}}}{}{%
Figure}{\special{language "Scientific Word";type
"GRAPHIC";maintain-aspect-ratio TRUE;display "USEDEF";valid_file "T";width
1.9558in;height 1.3402in;depth 0pt;original-width 1.9182in;original-height
1.305in;cropleft "0";croptop "1";cropright "1";cropbottom "0";tempfilename
'SV8THK0D.wmf';tempfile-properties "XPR";}}and hear them when jet aircraft
go supersonic. In the next figure we can see a picture of a T-38 aircraft
breaking the sound barrier. You can see the Mach cones, but notice that
there are several! Remember that a disturbance creates a wave. There are
disturbances created by the nose of the plane, the rudder, and the wings,
and perhaps the cockpit in this Schlieren photograph.\FRAME{dhFU}{1.593in}{%
1.2842in}{0pt}{\Qcb{Dr. Leonard Weinstein's Schlieren photograph of a T-38
Talon at Mach 1.1, altitude 13,700 feet, taken at NASA Langley Research
Center, Wallops in 1993. Image Courtacy NASA, image is in the Public Domain.}%
}{}{Figure}{\special{language "Scientific Word";type
"GRAPHIC";maintain-aspect-ratio TRUE;display "USEDEF";valid_file "T";width
1.593in;height 1.2842in;depth 0pt;original-width 1.5575in;original-height
1.2488in;cropleft "0";croptop "1";cropright "1";cropbottom "0";tempfilename
'SV8THK0E.wmf';tempfile-properties "XPR";}}

Superposition is also the basis of making music. We have talked about
resonance making music louder. But there is more to music and superposition.
Let's take this topic on next.

\section{Non Sinusoidal Waves}

We have dealt with pulses, and with sinusoidal waves. But this seems kind of
restrictive. Are all waves either pulses or sinusoidal? Can the universe
really be described by such simple mathematics? The answer is both no, and
yes. There are non-sinusoidal waves, in fact, most waves are not sinusoidal.
But it turns out that we can use a very clever mathematical trick to make
any shape wave out of a superposition of many sinusoidal waves. So our
mathematics for sinusoidal waves turns out to be quite general.

\subsection{Music and Non-sinusoidal waves}

From our example of standing waves on strings, we know that a string can
support a series of standing waves with discrete frequencies--the harmonic
series. We have also hinted that usually we excite more than one standing
wave at a time. The fundamental mode tends to give us the pitch we hear, but
what are the other standing waves for?

To understand, lets take an analogy. Making cookies and cakes.

Here is the beginning of a recipe for cookies.

\FRAME{dtbpF}{3.7896in}{2.1897in}{0in}{}{}{Figure}{\special{language
"Scientific Word";type "GRAPHIC";maintain-aspect-ratio TRUE;display
"USEDEF";valid_file "T";width 3.7896in;height 2.1897in;depth
0in;original-width 3.7421in;original-height 2.1499in;cropleft "0";croptop
"1";cropright "1";cropbottom "0";tempfilename
'SVYFOE00.wmf';tempfile-properties "XPR";}}The recipe is a list of
ingredients, and a symbolic instructions to mix and bake. The product is
chocolate chip cookies. Of course we need more information. We need to know
now much of each ingredient to use.\FRAME{dhF}{3.0052in}{2.3575in}{0pt}{}{}{%
Figure}{\special{language "Scientific Word";type
"GRAPHIC";maintain-aspect-ratio TRUE;display "USEDEF";valid_file "T";width
3.0052in;height 2.3575in;depth 0pt;original-width 2.962in;original-height
2.3177in;cropleft "0";croptop "1";cropright "1";cropbottom "0";tempfilename
'SUS73V42.wmf';tempfile-properties "XPR";}}This graph gives us the amount of
each ingredient by mass.

Now suppose we want chocolate cake.\FRAME{dtbpF}{3.5354in}{2.2001in}{0in}{}{%
}{Figure}{\special{language "Scientific Word";type
"GRAPHIC";maintain-aspect-ratio TRUE;display "USEDEF";valid_file "T";width
3.5354in;height 2.2001in;depth 0in;original-width 3.4886in;original-height
2.1603in;cropleft "0";croptop "1";cropright "1";cropbottom "0";tempfilename
'SVYFPG01.wmf';tempfile-properties "XPR";}}

The predominant taste in each of these foods is chocolate. But chocolate
cake and chocolate chip cookies don't taste exactly the same. We can easily
see that the differences in the other ingredients make the difference
between the \textquotedblleft cookie\ taste\textquotedblright\ and the
\textquotedblleft cake\ taste\textquotedblright\ that goes along with the
chocolate\ taste that predominates.

The sound waves produced by musical instruments work in a similar way. Here
is a recipe for an \textquotedblleft A\textquotedblright\ note from a
clarinet.\FRAME{dhF}{3.7879in}{2.7207in}{0pt}{}{}{Figure}{\special{language
"Scientific Word";type "GRAPHIC";maintain-aspect-ratio TRUE;display
"USEDEF";valid_file "T";width 3.7879in;height 2.7207in;depth
0pt;original-width 3.7395in;original-height 2.6783in;cropleft "0";croptop
"1";cropright "1";cropbottom "0";tempfilename
'SUS73V44.wmf';tempfile-properties "XPR";}}and here is one for a trumpet
playing the same \textquotedblleft A\textquotedblright\ note.\FRAME{dhF}{%
3.2569in}{2.3021in}{0pt}{}{}{Figure}{\special{language "Scientific
Word";type "GRAPHIC";maintain-aspect-ratio TRUE;display "USEDEF";valid_file
"T";width 3.2569in;height 2.3021in;depth 0pt;original-width
3.2119in;original-height 2.2623in;cropleft "0";croptop "1";cropright
"1";cropbottom "0";tempfilename 'SUS73V45.wmf';tempfile-properties "XPR";}}

A trumpet sounds different than a clarinet, and now we see why. There are
more harmonics involved with the trumpet sound than the clarinet sound.
These extra standing waves make up the \textquotedblleft
brassiness\textquotedblright\ of the trumpet sound. As with our baking
example, we need to know how much of each standing wave we have. Each will
have a different amplitude. For our trumpet, we might get amplitudes as
shown.

\FRAME{dhF}{3.0052in}{2.3575in}{0pt}{}{}{Figure}{\special{language
"Scientific Word";type "GRAPHIC";maintain-aspect-ratio TRUE;display
"USEDEF";valid_file "T";width 3.0052in;height 2.3575in;depth
0pt;original-width 2.962in;original-height 2.3177in;cropleft "0";croptop
"1";cropright "1";cropbottom "0";tempfilename
'SUS73V46.wmf';tempfile-properties "XPR";}}Note that the second harmonic has
a larger amplitude, but we still hear the \textquotedblleft
A\textquotedblright\ as at $440\unit{Hz}.$ A fugal horn would still sound
brassy, but would have a different mix of harmonics. For a clarinet we might
have something more like this\FRAME{dtbpF}{3.5379in}{2.7709in}{0in}{}{}{%
Figure}{\special{language "Scientific Word";type
"GRAPHIC";maintain-aspect-ratio TRUE;display "USEDEF";valid_file "T";width
3.5379in;height 2.7709in;depth 0in;original-width 3.4912in;original-height
2.7285in;cropleft "0";croptop "1";cropright "1";cropbottom "0";tempfilename
'SVYFRL02.wmf';tempfile-properties "XPR";}}

We have a tool that you can download to your PC to detect the mix of
harmonics of musical instruments, or mechanical systems. In music, the
different harmonics are called \emph{partials} because they make up part of
the sound. A graph that shows which harmonics are involved is called a \emph{%
spectrum}. The next figure is the spectrum of a six holed bamboo flute. Note
that there are several harmonics involved.

\FRAME{dhF}{3.2145in}{2.3445in}{0in}{}{}{Figure}{\special{language
"Scientific Word";type "GRAPHIC";maintain-aspect-ratio TRUE;display
"USEDEF";valid_file "T";width 3.2145in;height 2.3445in;depth
0in;original-width 3.1695in;original-height 2.3039in;cropleft "0";croptop
"1";cropright "1";cropbottom "0";tempfilename
'SUS73V47.wmf';tempfile-properties "XPR";}}Note that our graph has two
parts. One is the instantaneous spectrum, and one is the spectrum time
history.\FRAME{dhF}{4.6267in}{1.9389in}{0in}{}{}{Figure}{\special{language
"Scientific Word";type "GRAPHIC";maintain-aspect-ratio TRUE;display
"USEDEF";valid_file "T";width 4.6267in;height 1.9389in;depth
0in;original-width 4.574in;original-height 1.9009in;cropleft "0";croptop
"1";cropright "1";cropbottom "0";tempfilename
'SUS73V48.wmf';tempfile-properties "XPR";}}By observing the time history, we
can see changes in the spectrum. We can also see that we don't have pure
harmonics. The graph shows some response off the specific harmonic
frequencies. This six holed flute is very \textquotedblleft
breathy\textquotedblright\ giving a lot of wind noise along with the notes,
and we see this in the spectrum. In the next picture, I played a scale on
the flute.\FRAME{dhF}{2.7674in}{1.8421in}{0pt}{}{}{Figure}{\special{language
"Scientific Word";type "GRAPHIC";maintain-aspect-ratio TRUE;display
"USEDEF";valid_file "T";width 2.7674in;height 1.8421in;depth
0pt;original-width 2.725in;original-height 1.8049in;cropleft "0";croptop
"1";cropright "1";cropbottom "0";tempfilename
'SUS73V49.wmf';tempfile-properties "XPR";}}The instantaneous spectrum is not
active in this figure (since it can't show more than one note at a time) but
in the time history we see that as the fundamental frequency changes by
shorting the length of the flute (uncovering holes), we see that every
partial also goes up in frequency. The flute still has the characteristic
spectrum of a flute, but shifted to new frequencies. We can use this fact to
identify things by their vibration spectrum. In fact, that is how you
recognize voices and musics within your auditory system!

The technique of taking apart a wave into its components is very powerful.
With light waves, the spectrum is an indication of the chemical composition
of the emitter. For example, the spectrum of the sun looks something like
this\FRAME{dhFU}{4.6265in}{1.4644in}{0pt}{\Qcb{{\protect\small Solar coronal
spectrum taken during a solar eclipse. The successive curved lines are each
different wavelengths, and the dark lines are wavelengths that are absorbed.
The pattern of absorbed wavelengths allows a chemical analysis of the
corona. (Image in the Public Domain, orignally published in Bailey, Solon,
L, Popular Science Monthly, Vol 60, Nov. 1919, pp 244)}}}{}{Figure}{\special%
{language "Scientific Word";type "GRAPHIC";maintain-aspect-ratio
TRUE;display "USEDEF";valid_file "T";width 4.6265in;height 1.4644in;depth
0pt;original-width 4.574in;original-height 1.4295in;cropleft "0";croptop
"1";cropright "1";cropbottom "0";tempfilename
'SUS73V4A.wmf';tempfile-properties "XPR";}}The lines in this graph show the
amplitude of each harmonic component of the light. Darker lines have larger
amplitudes. The harmonics come from the excitation of electrons in their
orbitals. Each orbital is a different energy state, and when the electrons
jump from orbital to orbital, they produce specific wave frequencies. By
observing the mix of dark lines in pervious figure, and comparing to
laboratory measurements from each element (see next figure) we can find the
composition of the source. This figure shows the emission spectrum for
Calcium. because it is an emission spectrum the lines are bright instead of
dark. We can even see the color of each line!\FRAME{dhFU}{4.7693in}{0.9774in%
}{0pt}{\Qcb{{\protect\small Emission spectrum of Calcium (Image in the
Public Domain, courtesy NASA)}}}{}{Figure}{\special{language "Scientific
Word";type "GRAPHIC";maintain-aspect-ratio TRUE;display "USEDEF";valid_file
"T";width 4.7693in;height 0.9774in;depth 0pt;original-width
6.2699in;original-height 1.2626in;cropleft "0";croptop "1";cropright
"1";cropbottom "0";tempfilename 'SUS73V4B.wmf';tempfile-properties "XPR";}}

\subsection{Vibrometry}

Just like each atom has a specific spectrum, and each instrument, each
engine, machine, or anything that vibrates has a spectrum. We can use this
to monitor the health of machinery, or even to identify a piece of
equipment. Laser or acoustic vibrometers are available commercially.\FRAME{%
dhFU}{3.7572in}{1.64in}{0pt}{\Qcb{{\protect\small Laser Vibrometer Schematic
(Public Domain Image from Laderaranch:
http://commons.wikimedia.org/wiki/File:LDV\_Schematic.png)}}}{}{Figure}{%
\special{language "Scientific Word";type "GRAPHIC";maintain-aspect-ratio
TRUE;display "USEDEF";valid_file "T";width 3.7572in;height 1.64in;depth
0pt;original-width 9.7293in;original-height 4.2309in;cropleft "0";croptop
"1";cropright "1";cropbottom "0";tempfilename
'SUS73V4C.bmp';tempfile-properties "XPR";}}They provide a way to monitor
equipment in places where it would be dangerous or even impossible to send a
person. The equipment also does not need to be shut down, a great benefit
for factories that are never shut down, or for a satellite system that
cannot be reached by anyone.

\subsubsection{Fourier Series: Mathematics of Non-Sinusoidal Waves}

We should take a quick look at the mathematics of non-sinusoidal waves. We
said earlier that we could make any wave shape out of a superposition of
sinusoidal waves. Let' start with such a superposition of many sinusoidal
waves. The math looks like this%
\begin{equation*}
y\left( t\right) =\dsum\limits_{n}\left( A_{n}\sin \left( k_{n}x-\omega
_{n}t+\phi _{n}\right) +B_{n}\cos \left( k_{n}x-\omega _{n}t+\phi
_{n}\right) \right)
\end{equation*}%
where $A_{n}$ and $B_{n}$ are a series of coefficients and $f_{n}$ are the
harmonic series of frequencies.

\subsubsection{Example: Fourier representation of a square wave.}

Suppose we have a wave that isn't a sinusoid. And suppose it has a
complicated function $f\left( x,t\right) $ that we just don't know. For now,
let's consider a snapshot view of this complicated wave function so we have
just $f\left( x,t_{o}\right) =f\left( x\right) $ where our snapshot graph is
taken at time $t_{o}.$ That makes the wave in our snapshot graph just a
function of $x.$ We could represent a function $f\left( x\right) $ with our
combination of sine or cosine functions (or both) like in the last equation.
It might look something like this:%
\begin{eqnarray}
f\left( x\right) &=&C_{o}+C_{1}\cos \left( \frac{2\pi }{\lambda }%
x+\varepsilon _{1}\right) \\
&&+C_{2}\cos \left( \frac{2\pi }{\frac{\lambda }{2}}x+\varepsilon _{2}\right)
\\
&&+C_{3}\cos \left( \frac{2\pi }{\frac{\lambda }{3}}x+\varepsilon _{3}\right)
\\
&&+\ldots \\
&&+C_{n}\cos \left( \frac{2\pi }{\frac{\lambda }{n}}x+\varepsilon _{n}\right)
\\
&&+\ldots
\end{eqnarray}

where we will let $\varepsilon _{n}=\omega _{n}t_{o}+\phi _{n}$ just to make
the equations less messy.

The $C^{\prime }s$ are just coefficients that tell us the amplitude of the
individual cosine waves. The more terms in the series we take, the better
the approximation we will have, with the series exactly matching $f\left(
x\right) $ when the number of terms, $N\rightarrow \infty .$

Usually we rewrite the terms of the series as 
\begin{equation}
C_{m}\cos \left( mkx+\varepsilon _{m}\right) =A_{m}\cos \left( mkx\right)
+B_{m}\sin \left( mkx\right)
\end{equation}%
where $k$ is the wavenumber%
\begin{equation}
k=\frac{2\pi }{\lambda }
\end{equation}%
and $\lambda $ is the wavelength of the complicated but still periodic
function $f\left( x\right) .$ Then we identify (from our second favorite
trig identity) 
\begin{eqnarray}
A_{m} &=&C_{m}\cos \left( \varepsilon _{m}\right) \\
B_{m} &=&-C_{m}\sin \left( \varepsilon _{m}\right)
\end{eqnarray}%
then%
\begin{equation}
f\left( x\right) =\frac{A_{o}}{2}+\dsum\limits_{m-1}^{\infty }A_{m}\cos
\left( mkx\right) +\dsum\limits_{m-1}^{\infty }B_{m}\sin \left( mkx\right)
\label{Fourier Series}
\end{equation}%
where we separated out the $A_{o}/2$ term because it makes things nicer
later.

\subsubsection{Fourier Analysis}

The process of finding the coefficients of the series is called \emph{%
Fourier analysis}. We start by integrating equation (\ref{Fourier Series})%
\begin{equation}
\int_{0}^{\lambda }f\left( x\right) dx=\int_{0}^{\lambda }\frac{A_{o}}{2}%
dx+\int_{0}^{\lambda }\dsum\limits_{m-1}^{\infty }A_{m}\cos \left(
mkx\right) dx+\int_{0}^{\lambda }\dsum\limits_{m-1}^{\infty }B_{m}\sin
\left( mkx\right) dx
\end{equation}%
We can see immediately that all the sine and cosine terms integrate to zero
(we integrated over a wavelength) so%
\begin{equation}
\int_{0}^{\lambda }f\left( x\right) dx=\int_{0}^{\lambda }\frac{A_{o}}{2}dx=%
\frac{A_{o}}{2}\lambda
\end{equation}%
We solve this for $A_{o}$%
\begin{equation}
A_{o}=\frac{2}{\lambda }\int_{0}^{\lambda }f\left( x\right) dx
\end{equation}

To find the rest of the coefficients we need to remind ourselves of the
orthogonality of sinusoidal functions%
\begin{eqnarray}
\int_{0}^{\lambda }\sin \left( akx\right) \cos \left( bkx\right) dx &=&0 \\
\int_{0}^{\lambda }\cos \left( akx\right) \cos \left( bkx\right) dx &=&\frac{%
\lambda }{2}\delta _{ab} \\
\int_{0}^{\lambda }\sin \left( akx\right) \sin \left( bkx\right) dx &=&\frac{%
\lambda }{2}\delta _{ab}
\end{eqnarray}%
where $\delta _{ab}$ is the Kronecker delta.

To find the coefficients, then, we multiply both sides of equation (\ref%
{Fourier Series}) by $\cos \left( lkx\right) $ where $l$ is a positive
integer. Then we integrate over one wavelength.%
\begin{eqnarray}
\int_{0}^{\lambda }f\left( x\right) \cos \left( lkx\right) dx
&=&\int_{0}^{\lambda }\frac{A_{o}}{2}\cos \left( lkx\right) dx \\
&&+\int_{0}^{\lambda }\dsum\limits_{m-1}^{\infty }A_{m}\cos \left(
mkx\right) \cos \left( lkx\right) dx \\
&&+\int_{0}^{\lambda }\dsum\limits_{m-1}^{\infty }B_{m}\sin \left(
mkx\right) \cos \left( lkx\right) dx
\end{eqnarray}%
which gives 
\begin{equation}
\int_{0}^{\lambda }f\left( x\right) \cos \left( mkx\right)
dx=\int_{0}^{\lambda }A_{m}\cos \left( mkx\right) \cos \left( mkx\right) dx
\end{equation}%
that is, only the term with two cosine functions where $l=m$ will be non
zero. So%
\begin{equation}
\int_{0}^{\lambda }f\left( x\right) \cos \left( mkx\right) dx=\frac{\lambda 
}{2}A_{m}
\end{equation}%
solving for $A_{m}$ we have%
\begin{equation}
A_{m}=\frac{2}{\lambda }\int_{0}^{\lambda }f\left( x\right) \cos \left(
mkx\right) dx
\end{equation}

We can perform the same steps to find $B_{m}$ only we use $\sin \left(
lkx\right) $ as the multiplier. Then we find%
\begin{equation}
B_{m}=\frac{2}{\lambda }\int_{0}^{\lambda }f\left( x\right) \sin \left(
mkx\right) dx
\end{equation}%
Once we have the $A_{m}$ and $B_{m}$ coefficients we can put together our
series and have an approximation for $f\left( x\right) .$ Let's try this for
an example.

\subsection{Square wave}

Let's find the series for a square wave using our Fourier analysis technique.

Let's take 
\begin{equation}
\lambda =2
\end{equation}%
\begin{equation}
f(x)=\left\{ 
\begin{array}{ccl}
1 & \text{if} & 0<x<\frac{\lambda }{2} \\ 
-1 & \text{if} & \frac{\lambda }{2}<x<\lambda%
\end{array}%
\right.
\end{equation}

\FRAME{dtbpF}{2.9253in}{1.9558in}{0in}{}{}{Figure}{\special{language
"Scientific Word";type "GRAPHIC";maintain-aspect-ratio TRUE;display
"USEDEF";valid_file "T";width 2.9253in;height 1.9558in;depth
0in;original-width 2.9519in;original-height 1.9638in;cropleft "0";croptop
"1";cropright "1";cropbottom "0";tempfilename
'T9NI4V01.wmf';tempfile-properties "XPR";}}

Since $f\left( x\right) is$ odd, $A_{m}=0$ for all $m$. We have%
\begin{equation}
B_{m}=\frac{2}{\lambda }\int_{0}^{\frac{\lambda }{2}}\left( 1\right) \sin
\left( mkx\right) dx+\frac{2}{\lambda }\int_{\frac{\lambda }{2}}^{\lambda
}\left( -1\right) \sin \left( mkx\right) dx
\end{equation}%
so%
\begin{equation}
B_{m}=\frac{1}{m\pi }\left( -\cos \left( mkx\right) \mathstrut \right\vert
_{0}^{\frac{\lambda }{2}}+\frac{1}{m\pi }\left( \cos \left( mkx\right)
\mathstrut \right\vert _{\frac{\lambda }{2}}^{\lambda }
\end{equation}%
Which is%
\begin{equation}
B_{m}=\frac{1}{m\pi }\left( 1\cos \left( m\frac{2\pi }{\lambda }x\right)
\mathstrut \right\vert _{0}^{\frac{\lambda }{2}}+\frac{1}{m\pi }\left( \cos
\left( m\frac{2\pi }{\lambda }x\right) \right\vert _{\frac{\lambda }{2}%
}^{\lambda }
\end{equation}%
so%
\begin{eqnarray}
B_{m} &=&\frac{1}{m\pi }\left( \left( -\cos \left( m\frac{2\pi }{\lambda }%
\frac{\lambda }{2}\right) \right) +\cos \left( m\frac{2\pi }{\lambda }\left(
0\right) \right) \right) \\
&&+\frac{1}{m\pi }\left( \left( \cos \left( m\frac{2\pi }{\lambda }\lambda
\right) -\cos \left( m\frac{2\pi }{\lambda }\frac{\lambda }{2}\right)
\right) \right)
\end{eqnarray}%
which is%
\begin{equation}
B_{m}=\frac{2}{m\pi }\left( 1-\cos \left( m\pi \right) \right)
\end{equation}

Our series is then just 
\begin{equation}
f\left( x\right) =\dsum\limits_{m-1}^{\infty }\frac{2}{m\pi }\left( 1-\cos
\left( m\pi \right) \right) \sin \left( mkx\right)
\end{equation}%
and we can write a few terms%
\begin{equation}
\begin{tabular}{|l|l|}
\hline
$Term$ &  \\ \hline
1 & $\frac{4}{\pi }\sin \left( kx\right) $ \\ \hline
2 & $0$ \\ \hline
3 & $\frac{4}{3\pi }\sin \left( 3kx\right) $ \\ \hline
4 & $0$ \\ \hline
5 & $\frac{4}{5\pi }\sin \left( 5kx\right) $ \\ \hline
\end{tabular}%
\end{equation}%
Then the partial sum up to $m=5$ looks like%
\begin{equation}
f\left( x\right) =\frac{4}{\pi }\sin \left( kx\right) +\frac{4}{3\pi }\sin
\left( 3kx\right) +\frac{4}{5\pi }\sin \left( 5kx\right)
\end{equation}%
\FRAME{dtbpF}{2.7461in}{1.8369in}{0pt}{}{}{Figure}{\special{language
"Scientific Word";type "GRAPHIC";maintain-aspect-ratio TRUE;display
"USEDEF";valid_file "T";width 2.7461in;height 1.8369in;depth
0pt;original-width 2.7701in;original-height 1.8432in;cropleft "0";croptop
"1";cropright "1";cropbottom "0";tempfilename
'T9NI5X03.wmf';tempfile-properties "XPR";}}and we can see it kind of works.
We don't expect a great result until we have many many terms. If we take
more terms,%
\begin{eqnarray}
f\left( x\right) &=&\frac{4}{\pi }\sin \left( kx\right) +\frac{4}{3\pi }\sin
\left( 3kx\right) +\frac{4}{5\pi }\sin \left( 5kx\right) +\frac{4}{7\pi }%
\sin \left( 7kx\right) +\frac{4}{9\pi }\sin \left( 95kx\right) \\
&&+\frac{4}{11\pi }\sin \left( 11kx\right) +\frac{4}{13\pi }\sin \left(
13kx\right) +\frac{4}{15\pi }\sin \left( 15kx\right) +\frac{4}{17\pi }\sin
\left( 17kx\right) +\frac{4}{19\pi }\sin \left( 19kx\right)  \notag
\end{eqnarray}%
We see the function get closer and closer to a square wave.\FRAME{dtbpF}{%
4.5484in}{0.9677in}{0pt}{}{}{Figure}{\special{language "Scientific
Word";type "GRAPHIC";maintain-aspect-ratio TRUE;display "USEDEF";valid_file
"T";width 4.5484in;height 0.9677in;depth 0pt;original-width
4.6061in;original-height 0.957in;cropleft "0";croptop "1";cropright
"1";cropbottom "0";tempfilename 'T9NI6T06.wmf';tempfile-properties "XPR";}}

In the limit of infinitely many waves, the match would be perfect. But we
don't usually need an infinite number of terms. we can pick the part of the
spectrum that best represents the phenomena we desire to observe. For
example, oil based compounds all have specific spectral signatures in the
wavelength range between $3-5$ micrometers. If you wish to tell the
difference between gasoline and crude oil, you can restrict your study to
these wavelengths alone.

\section{Frequency Uncertainty for Signals and Particles}

Up to this lecture, when we thought of a wave we have mostly thought of
something like this 
\begin{equation*}
y=y_{\max }\sin \left( kx-\omega t-\phi _{o}\right)
\end{equation*}%
which in practice might look like this\FRAME{dtbpF}{4.4996in}{1.2211in}{0pt}{%
}{}{Plot}{\special{language "Scientific Word";type "MAPLEPLOT";width
4.4996in;height 1.2211in;depth 0pt;display "USEDEF";plot_snapshots
TRUE;mustRecompute FALSE;lastEngine "MuPAD";xmin "-5";xmax "5";xviewmin
"-5.0010000010002";xviewmax "5.0010000010002";yviewmin "-10";yviewmax
"10";plottype 4;axesFont "Times New
Roman,12,0000000000,useDefault,normal";numpoints 100;plotstyle
"patch";axesstyle "normal";axestips FALSE;xis \TEXUX{x};var1name
\TEXUX{$x$};function \TEXUX{$10\sin \left( 5x-0\right) $};linecolor
"blue";linestyle 1;pointstyle "point";linethickness 3;lineAttributes
"Solid";var1range "-5,5";num-x-gridlines 100;curveColor
"[flat::RGB:0x000000ff]";curveStyle "Line";VCamFile
'T9OUDQ01.xvz';valid_file "T";tempfilename
'T9OUE000.wmf';tempfile-properties "XPR";}}And we have noticed that there is
no start or stop to this kind of wave. Our figure starts at $x=-5$ and ends
at $x=5,$ but the equation does not! There is a value of $y$ for every $x$
from $-\infty $ to +$\infty .$ But many signals are not such waves. They may
be very limited in size.\FRAME{dtbpF}{2.0833in}{0.8172in}{0pt}{}{}{Figure}{%
\special{language "Scientific Word";type "GRAPHIC";maintain-aspect-ratio
TRUE;display "USEDEF";valid_file "T";width 2.0833in;height 0.8172in;depth
0pt;original-width 2.0942in;original-height 0.8036in;cropleft "0";croptop
"1";cropright "1";cropbottom "0";tempfilename
'../../../../PH279/Repository/IntroductoryPhysicsIV/Textbook/FigureLimetedWave2.wmf';tempfile-properties "XNPR";}%
}We should investigate what happens when you have a limited wave. I did this
investigation using python. Suppose we have a sine wave with $f=200\unit{Hz}%
, $ but I\ limit this wave's existence by making it start at $t_{i}=0$ and
then make it end at $t_{f}=10\unit{s}.$ I could do this in practice by
turning on a radio transmission or even an acoustic speaker, and then
turning the device off ten seconds later. Our screen resolution is terrible
for plotting such a function, but in the figure below you can see that our
signal only exists from $t=0$ to $t=10\unit{s}.$

\FRAME{dtbpF}{4.2756in}{0.7014in}{0pt}{}{}{Figure}{\special{language
"Scientific Word";type "GRAPHIC";maintain-aspect-ratio TRUE;display
"USEDEF";valid_file "T";width 4.2756in;height 0.7014in;depth
0pt;original-width 4.2246in;original-height 0.6702in;cropleft "0";croptop
"1";cropright "1";cropbottom "0";tempfilename
'T9OUU803.wmf';tempfile-properties "XPR";}} I plotted this using a symbolic
math package. If I zoom in on a part of the graph we can see that it is
really a sine wave.\FRAME{dtbpF}{4.6951in}{0.7688in}{0pt}{}{}{Figure}{%
\special{language "Scientific Word";type "GRAPHIC";maintain-aspect-ratio
TRUE;display "USEDEF";valid_file "T";width 4.6951in;height 0.7688in;depth
0pt;original-width 4.6423in;original-height 0.7368in;cropleft "0";croptop
"1";cropright "1";cropbottom "0";tempfilename
'T9OUU804.wmf';tempfile-properties "XPR";}}

Python did equally bad at plotting this. All we see is a blue band. \FRAME{%
dtbpF}{4.3388in}{1.8317in}{0pt}{}{}{Figure}{\special{language "Scientific
Word";type "GRAPHIC";maintain-aspect-ratio TRUE;display "USEDEF";valid_file
"T";width 4.3388in;height 1.8317in;depth 0pt;original-width
4.3932in;original-height 1.8378in;cropleft "0";croptop "1";cropright
"1";cropbottom "0";tempfilename
'../../../../PH279/Repository/IntroductoryPhysicsIV/Textbook/Figure10s200Hz.wmf';tempfile-properties "XNPR";}%
}But in the second graph, notice that we have plotted frequency. Python and
most scientific programing languages have functions to find a digital
version of a Fourier transform to produce a graph of which frequencies are
in a function. For our EE students, replace the word \textquotedblleft
function\textquotedblright\ with the word \textquotedblleft
signal.\textquotedblright\ We call this kind of a graph a \emph{spectrum}.
It takes in a signal and finds which frequencies are in a signal. It
performs the job of a spectrometer, so we would call the figure to the right
a spectrogram (or just a spectrum).

We expect only one frequency, $200\unit{Hz},$ and that is mostly what we
get. Since our period for our wave is 
\begin{equation*}
T=\frac{1}{200\unit{Hz}}=0.005\,\unit{s}
\end{equation*}%
and we have $10\unit{s}$ of data, that is four orders of magnitude more
signal than one period. The whole signal seems very long compared to a
period. We expect this to look kind of like an infinite signal. But suppose
we take the same wave, but for less time. We limit the wave more.\FRAME{dtbpF%
}{3.8104in}{1.5549in}{0pt}{}{}{Figure}{\special{language "Scientific
Word";type "GRAPHIC";maintain-aspect-ratio TRUE;display "USEDEF";valid_file
"T";width 3.8104in;height 1.5549in;depth 0pt;original-width
3.8548in;original-height 1.5558in;cropleft "0";croptop "1";cropright
"1";cropbottom "0";tempfilename
'../../../../PH279/Repository/IntroductoryPhysicsIV/Textbook/Figure1s200Hz.wmf';tempfile-properties "XNPR";}%
}So we still get a blue blur for our wave picture in python, but now the
wave only exists for one instead of ten seconds. If you look closely at the
frequency graph, you will notice that the $200\unit{Hz}$ peak representing
our wave is a bit wider right at the bottom.

We could limit our wave more, say, so it only lasts $t_{f}=0.1\unit{s}.$ We
would get a set of graphs that look like this. \FRAME{dtbpF}{3.7351in}{%
1.4788in}{0pt}{}{}{Figure}{\special{language "Scientific Word";type
"GRAPHIC";maintain-aspect-ratio TRUE;display "USEDEF";valid_file "T";width
3.7351in;height 1.4788in;depth 0pt;original-width 3.7776in;original-height
1.4786in;cropleft "0";croptop "1";cropright "1";cropbottom "0";tempfilename
'../../../../PH279/Repository/IntroductoryPhysicsIV/Textbook/Figure0.1s200Hz.wmf';tempfile-properties "XNPR";}%
}

Notice that not only can Python render the wave now, but more importantly
the $200\unit{Hz}$ frequency peak is noticeable wider! This is profound! It
means that by limiting the wave, we no longer have just one frequency! The
graph tells us we have mostly $200\unit{Hz}$ but we also have some $199\unit{%
Hz}$ and some $201\unit{Hz}$ and some $190\unit{Hz}$ and some $210\unit{Hz},$
etc. The very fact that the wave does not go on so long requires that we
have more than one frequency in the wave. We could say that as $\Delta t$
gets smaller, our $\Delta f$ is getting bigger. Here are two more examples
with smaller $\Delta t$ values. \FRAME{dtbpF}{3.9288in}{1.4866in}{0pt}{}{}{%
Figure}{\special{language "Scientific Word";type
"GRAPHIC";maintain-aspect-ratio TRUE;display "USEDEF";valid_file "T";width
3.9288in;height 1.4866in;depth 0pt;original-width 3.9746in;original-height
1.4866in;cropleft "0";croptop "1";cropright "1";cropbottom "0";tempfilename
'../../../../PH279/Repository/IntroductoryPhysicsIV/Textbook/Figure0.05s200Hz.wmf';tempfile-properties "XNPR";}%
}\FRAME{dtbpF}{4.6285in}{1.8152in}{0pt}{}{}{Figure}{\special{language
"Scientific Word";type "GRAPHIC";maintain-aspect-ratio TRUE;display
"USEDEF";valid_file "T";width 4.6285in;height 1.8152in;depth
0pt;original-width 3.3164in;original-height 1.2817in;cropleft "0";croptop
"1";cropright "1";cropbottom "0";tempfilename
'../../../../PH279/Repository/IntroductoryPhysicsIV/Textbook/Figure0.005s200Hz.wmf';tempfile-properties "XNPR";}%
}

The cost of limiting our waves is that we can't have a single frequency for
the wave. For an engineer, this means that if you only measure a short
segment of the signal, you have an increased uncertainty in the frequency
you will find from that signal. For chemists, it means that when we look at
quantum wave functions we can expect uncertainty in the frequency (or
wavelength) because the quantum wave particle (like an electron) is limited.
It is important to know what we mean by uncertainty in this case. In our
example above, when we say that $\Delta f$ increased we really mean that we
have more than one frequency. We don't just mean that we don't know the
frequency well. We really are mixing more than one frequency. This is worth
saying again. It's not that we don't know what the frequency is, it is that
there is more than one frequency, so giving one value for the frequency
isn't a good description of our wave. This idea of limiting waves creating
uncertainty shows up in physical chemistry as Heisenberg's uncertainty
principle.

We have talked about waves in general and used sound and light waves as
examples. But light waves have played such an important role in physics that
the study of light is considered a subdiscipline on it's own. We call this
subdiscipline \emph{optics}. So to continue our study of waves, we are going
to start a study of optics and we will begin our study of optics in the next
lecture.

%TCIMACRO{%
%\TeXButton{Basic Equations}{\vspace{0.5in}\hspace{-1.3in}{\LARGE Basic Equations\vspace{0.25in}}}}%
%BeginExpansion
\vspace{0.5in}\hspace{-1.3in}{\LARGE Basic Equations\vspace{0.25in}}%
%EndExpansion
\begin{equation}
\sin \theta =\frac{v_{sound}t}{v_{ed}t}=\frac{v_{sound}}{v_{ed}}
\end{equation}
\begin{equation*}
M=\frac{v_{ed}}{v_{sound}}
\end{equation*}
\begin{equation*}
\sin \theta =\frac{1}{M}
\end{equation*}

\part{Optics}

\chapter{11 Nature of Light and Reflection 3.1.1, 3.1.2}

%TCIMACRO{%
%\TeXButton{\chaptermark{Light and Reflection}}{\chaptermark{Light and Reflection}}}%
%BeginExpansion
\chaptermark{Light and Reflection}%
%EndExpansion

%TCIMACRO{%
%\TeXButton{Question 223.13.4}{\marginpar {
%\hspace{-0.5in}
%\begin{minipage}[t]{1in}
%\small{Question 223.13.4}
%\end{minipage}
%}}}%
%BeginExpansion
\marginpar {
\hspace{-0.5in}
\begin{minipage}[t]{1in}
\small{Question 223.13.4}
\end{minipage}
}%
%EndExpansion
In the science fiction movies inter-galactic star ships blast each other
with laser cannons. The laser beams streak across the screen. This is
dramatic, but not realistic. For us to see the light path, some of the light
leave the path and must get to our eyes. The light must either travel
directly to our eyes from the source, or it must bounce off of something
along its way. We can make a laser beam visible by providing dust for the
light to bounce off of so it travels to our eyes. But unless that galaxy far
far away is really dusty, Hollywood doesn't understand light waves very
well. Let's start our study of geometric optics with a few comments on what
light is. Then let's study reflection, the bouncing of light off of an
object (like dust, or a mirror).

%TCIMACRO{%
%\TeXButton{Fundamental Concepts}{\vspace{0.10in}\hspace{-0.37in}{\Large Fundamental Concepts\vspace{0.2in}}}}%
%BeginExpansion
\vspace{0.10in}\hspace{-0.37in}{\Large Fundamental Concepts\vspace{0.2in}}%
%EndExpansion

\begin{itemize}
\item Light is a wave in the electromagnetic field

\item Light travels fast, with a speed of $c=3\times 10^{8}\unit{m}/\unit{s}$

\item Law of Reflection
\end{itemize}

\section{What is Light?}

Before the 19th century (1800's) light was assumed to be a stream of
particles. Newton was one of the chief proponents of this theory. The theory
was able to explain reflection of light from mirrors and other objects and
therefore explain vision. In 1678 Huygens showed that wave theory could also
explain reflection and vision.

In 1801 Thomas Young demonstrated that light had attributes that were best
explained by wave theory. We will study Young's experiment later. The crux
of his experiment was to show that light displayed constructive and
destructive interference--clearly a wave phenomena! The theory of the nature
of light took a dramatic shift.

In 1805 Joseph Smith was born in Sharon, Vermont.

In September of 1832 Joseph Smith received a revelation that said in part :

\begin{quotation}
For the word of the Lord is truth, and whatsoever is truth is light, and
whatsoever is light is Spirit, even the Spirit of Jesus Christ. And the
Spirit giveth light to every man that cometh into the world; and the Spirit
enlighteneth every man through the world, that hearkeneth to the voice of
the Spirit. (D\&C 84:45-46)
\end{quotation}

In December of 1832 Joseph Smith received another revelation that says in
part:

\begin{quotation}
This Comforter is the promise which I give unto you of eternal life, even
the glory of the celestial kingdom; which glory is that of the church of the
Firstborn, even of God, the holiest of all, through Jesus Christ his
Son---He that ascended up on high, as also he descended below all things, in
that he comprehended all things, that he might be in all and through all
things, the light of truth; which truth shineth. This is the light of
Christ. As also he is in the sun, and the light of the sun, and the power
thereof by which it was made. As also he is in the moon, and is the light of
the moon, and the power thereof by which it was made; as also the light of
the stars, and the power thereof by which they were made; and the earth
also, and the power thereof, even the earth upon which you stand. And the
light which shineth, which giveth you light, is through him who enlighteneth
your eyes, which is the same light that quickeneth your understandings;
which light proceedeth forth from the presence of God to fill the immensity
of space---the light which is in all things, which giveth life to all
things, which is the law by which all things are governed, even the power of
God who sitteth upon his throne, who is in the bosom of eternity, who is in
the midst of all things. (D\&C 88:5-12)
\end{quotation}

Light, even non-metaphorical physical light, seems to be of interest to
members of the Restored Church.

In 1847 the saints entered the Salt Lake Valley.

In 1873 James Clark Maxwell published his findings that light is an
electromagnetic wave (something we will try to show in Principles of Physics
III!). That is, the \textquotedblleft group of objects\textquotedblright\
that form a \textquotedblleft medium\textquotedblright\ for light are not
made out of matter. But are part of what we call the electromagnetic field.
This field is something, but not made out of matter.

Planck's work in quantization theory (1900) was used by Einstein In 1905 to
give an explantation of the photoelectric effect that again made light look
like a particle.

Current theory allows light to exhibit the characteristics of a wave in some
situations and like a particle in others. We will study both before the end
of our optics unit.

The results of Einstein's work give us the concept of a \emph{photon} or a
quantized unit of radiant energy. Each \textquotedblleft piece of
light\textquotedblright\ or photon has energy 
\begin{equation}
E=hf
\end{equation}%
where $f$ is the frequency of the light and $h$ is a constant 
\begin{equation}
h=6.63\times 10^{-34}\unit{J}\unit{s}
\end{equation}

The nature of light is fascinating and useful both in physical and religious
areas of thought.

\section{Measurements of the Speed of Light}

One of the great foundations of modern physical theory is that the speed of
light is constant in a vacuum. Galileo first tried to measure the speed of
light. He used two towers in town and placed a lantern and an assistant on
each tower. The lanterns had shades. The plan was for one assistant to
remove his shade, and then for the assistant on the other tower to remove
his shade as soon as he saw the light from the first lantern. Back at the
first tower, the first assistant would use a clock to determine the time
difference between when the first lantern was un-shaded, and when they saw
the light from the second tower. The light would have traveled twice the
inter-tower distance. Dividing that distance by the time would give the
speed of light. You can probably guess that this did not work. Light travels
very quickly. The clocks of Galileo's day could not measure such a small
time difference. Ole R\o mer was the first to succeed in measuring the speed
of light.

\subsection{R\o mer's Measurement of the speed of light}

\FRAME{dhFU}{1.1761in}{2.0626in}{0pt}{\Qcb{A diagram illustrating R\o mer's
determination of the speed of light. Point A is the Sun, piont B is Jupiter.
Point C is the immersion of Io into Jupiter's shadow at the start of an
eclipse}}{}{Figure}{\special{language "Scientific Word";type
"GRAPHIC";maintain-aspect-ratio TRUE;display "USEDEF";valid_file "T";width
1.1761in;height 2.0626in;depth 0pt;original-width 1.7252in;original-height
3.0468in;cropleft "0";croptop "1";cropright "1";cropbottom "0";tempfilename
'SVAJW80G.wmf';tempfile-properties "XPR";}}R\o mer performed his measurement
in 1675, 269 years before digital devices existed!. He used the orbital
period of Io, a moon of Jupiter, as Jupiter orbited around the sun. He first
measured the period of Io's rotation about Jupiter, then he predicted an
eclipse of Io three months later. But he found his calculation was off by $%
600\unit{s}.$ After careful thought, he realized that the Earth had moved in
its orbit, and that the light had to travel the extra distance due to the
Earth's new position. Given R\o mer's best estimate for the orbital radius
of the earth and his time difference, R\o mer arrived at a estimate of $%
c=2.3\times 10^{8}\frac{\unit{m}}{\unit{s}}.$ Our current value is $%
c=2.99792458\times 10^{8}\frac{\unit{m}}{\unit{s}}.$ Amazing for 1675!

\subsection{Fizeau's Measurement of the speed of light}

\FRAME{dtbpF}{2.5659in}{1.6881in}{0pt}{}{}{Figure}{\special{language
"Scientific Word";type "GRAPHIC";maintain-aspect-ratio TRUE;display
"USEDEF";valid_file "T";width 2.5659in;height 1.6881in;depth
0pt;original-width 2.5253in;original-height 1.6518in;cropleft "0";croptop
"1";cropright "1";cropbottom "0";tempfilename
'SVAJW80H.wmf';tempfile-properties "XPR";}}Hippolyte Fizeau measured the
speed of light in 1849 using the apparatus drawn in the figure above. He
used a toothed wheel and a mirror and a beam of light. The light passed
through the open space in the wheel's teeth as the wheel rotated. Then was
reflected by the mirror. The speed would be%
\begin{equation*}
v=\frac{\Delta x}{\Delta t}
\end{equation*}%
We just need $\Delta x$ and $\Delta t.$

It is easy to see that 
\begin{equation*}
\Delta x=2d
\end{equation*}%
because the light travels twice the distance to the mirror ($d$) and back.
If the wheel turned just at the right angular speed, then the reflected
light would hit the next tooth and be blocked. Think of angular speed $%
\omega $ (not angular frequency this time)%
\begin{equation*}
\omega =\frac{\Delta \theta }{\Delta t}
\end{equation*}%
So, the time difference for a tooth to move would be 
\begin{equation*}
\Delta t=\frac{\Delta \theta }{\omega }
\end{equation*}%
We can find $\Delta \theta $ by taking $2\pi $ and dividing by the number of
teeth on the wheel.%
\begin{equation*}
\Delta \theta =\frac{2\pi }{N_{teeth}}
\end{equation*}

Then the speed of light must be 
\begin{eqnarray*}
c &=&v=\frac{2d}{\frac{\Delta \theta }{\omega }} \\
&=&\frac{2d\omega }{\Delta \theta } \\
&=&\frac{2d\omega N_{teeth}}{2\pi } \\
&=&\frac{d\omega N_{teeth}}{\pi }
\end{eqnarray*}

Then if we have $720$ teeth and $\omega $ is measured to be $d=7500\unit{m}$%
\begin{eqnarray*}
c &=&\frac{\left( 7500\unit{m}\right) \left( 172.\,\allowbreak 79\unit{Hz}%
\right) \left( 720\right) }{\pi } \\
&=&2.\,\allowbreak 97\times 10^{8}\frac{\unit{m}}{\unit{s}}
\end{eqnarray*}%
which is Fizeau's number and it is pretty good!

Modern measurements are performed in very much the same way that Fizeau did
his calculation. The current value is 
\begin{equation}
c=2.9979\times 10^{8}\frac{\unit{m}}{\unit{s}}
\end{equation}

\subsection{Faster than light}

The speed of light in a vacuum is constant, but in matter the speed of light
changes.

We will study this in detail when we look at refraction. But for now, a
dramatic example is Cherenkov radiation. It is an eerie blue glow around the
core of nuclear reactors. It occurs when electrons are accelerated past the
speed of light in the water surrounding the core. The electrons emit light
and the light waves form a Doppler cone or a light-sonic boom! The result is
the blue glow.\FRAME{dtbpFUw}{2.9457in}{2.2024in}{0pt}{\Qcb{{\protect\small %
\ Cherenkov radiation (United States Department of Energy, image in the
public domain)}}}{}{Figure}{\special{language "Scientific Word";type
"GRAPHIC";maintain-aspect-ratio TRUE;display "USEDEF";valid_file "T";width
2.9457in;height 2.2024in;depth 0pt;original-width 6.2595in;original-height
4.6717in;cropleft "0";croptop "1";cropright "1";cropbottom "0";tempfilename
'SVAJW80I.wmf';tempfile-properties "XPR";}}This does bring up a problem in
terminology. What does the word \textquotedblleft medium\textquotedblright\
mean? We have used it to mean the substance through which a wave travels.
This substance must have the property of transferring energy between it's
parts, like the coils of a spring can transfer energy to each other, or like
air molecules can transfer energy by collision. For light the wave medium is
the electromagnetic field. This field can store and transfer energy but it
isn't made of atoms or parts of mater. This may be hard to imagine, but
picture space filled with this non-material substance that can stretch and
hold energy. This is our electromagnetic field. We will study the
electromagnetic field in Principles of Physics III (PH220). So for now, we
just need a mental picture of the field having waves.

But many books on physics call materials like glass a \textquotedblleft
medium\textquotedblright\ through which light travels. The water in our last
example is a \textquotedblleft medium\textquotedblright\ in this sense. Are
glass and water wave mediums for light? The answer is an emphatic NO! Light
does not need any matter to form it's wave. The wave medium is the
electromagnetic field. So we will have to keep this in mind as we allow
light to travel through matter. I will try to say that the light enters a
new \textquotedblleft material\textquotedblright\ to describe something like
light entering a piece of glass. But some books may call the matter a
\textquotedblleft medium.\textquotedblright\ You must remember that by this
they don't mean the wave medium.

Light can't be a wave in some type of matter. We know this because light
travels through empty space. But in Principles of Physics III we will find
that space isn't exactly empty. It has an electromagnetic field in it. And
it turns out that light seems to be a wave in this electromagnetic field. It
will take us a while to fully understand this concept but don't worry.
Physicists knew that light was a wave for almost 80 years before the
electric field was shown to be the medium. We can do a lot just knowing
light behaves like a wave.

This makes the light we see just one small part of a whole class of waves
that are possible in this electromagnetic field medium. Radio waves, and
microwaves, and x-rays are all just different types of electromagnetic
waves. The next figure shows where all of these electromagnetic waves fit
ordered by wavelength (and frequency).\FRAME{dtbpFU}{5.3048in}{0.4013in}{0pt%
}{\Qcb{{\protect\small Image Courtesy NASA}}}{}{Figure}{\special{language
"Scientific Word";type "GRAPHIC";maintain-aspect-ratio TRUE;display
"USEDEF";valid_file "T";width 5.3048in;height 0.4013in;depth
0pt;original-width 7.2032in;original-height 0.518in;cropleft "0";croptop
"1";cropright "1";cropbottom "0";tempfilename
'SVAJW80J.wmf';tempfile-properties "XPR";}}

There is something very unique about this electromagnetic field medium. The
waves in this medium travel at a constant speed no matter what frame of
reference we are in. This fact lead to the formation of the Special Theory
of Relativity and the famous equation 
\begin{equation*}
E=mc^{2}
\end{equation*}%
where $c$ is this speed of light%
\begin{equation*}
c=299792458\frac{\unit{m}}{\unit{s}}
\end{equation*}

Light does slow down when it enters a material, like glass, or even air.
This is not really because it moves slower. what happens is that light is
absorbed by the electrons in the atoms of the material substance. The
electron temporarily takes up all the energy from a bit of the light wave.
But only temporarily. It eventually has to give up the energy and a light
wave is reformed coming from the atoms. There is really a superposition of
what is left of the original light wave, and the new light wave generated by
the atoms. And the resultant wave moves slower. Think of our last lecture on
non-sinusoidal waves. The resultant wave is made of two sinusoidal waves
both traveling at the speed of light. But the combination of the two waves
will also be a wave, and this resultant wave will appear to travel at a
slower speed. How much slower depends on how long the electrons in the atoms
can hang on to the light and the geometry of the situation. Each substance
is different.

We can devise a way to express how much slower light will appear to go in a
substance using the ratio%
\begin{equation*}
\frac{c}{v}
\end{equation*}%
the ratio of the speed of light, $c,$ to the average speed in the substance, 
$v.$ This ratio is so useful that we give it a name, the \emph{index of
refraction}.%
\begin{equation*}
n=\frac{c}{v}
\end{equation*}

Let's look at some a simple interaction of light waves with a material, but
let's not look at the wave and atomic re-emition level. Let's treat our
combined wave as just one resultant wave that travels in a specific
direction. We can call this wave a \emph{ray} (as in a ray of light!).

\section{Reflection}

Using the ray approximation we wish to find what happens when a bundle of
rays reaches a boundary between materials. If the material boundary is very
smooth, then the rays are reflected (bounce off) in a uniform way. This is
called \emph{specular }reflection\FRAME{dhF}{2.3895in}{2.1344in}{0pt}{}{}{%
Figure}{\special{language "Scientific Word";type
"GRAPHIC";maintain-aspect-ratio TRUE;display "USEDEF";valid_file "T";width
2.3895in;height 2.1344in;depth 0pt;original-width 2.3497in;original-height
2.0954in;cropleft "0";croptop "1";cropright "1";cropbottom "0";tempfilename
'SVAJW80K.wmf';tempfile-properties "XPR";}}If the surface is not smooth,
then something different happens. The rays are reflected, but they are
reflected randomly

\FRAME{dtbpF}{2.3886in}{1.8498in}{0pt}{}{}{Figure}{\special{language
"Scientific Word";type "GRAPHIC";maintain-aspect-ratio TRUE;display
"USEDEF";valid_file "T";width 2.3886in;height 1.8498in;depth
0pt;original-width 3.1211in;original-height 2.4102in;cropleft "0";croptop
"1";cropright "1";cropbottom "0";tempfilename
'SW0DNF00.wmf';tempfile-properties "XPR";}}This is called \emph{diffuse}
reflection

This difference can be seen in real life. In the next figure, the surface on
the left is a specular reflector and you can see the reflected light beam.
But the surface on the right photo is diffuse, no reflected beam is seen.%
\FRAME{dtbpF}{4.403in}{1.3553in}{0in}{}{}{Figure}{\special{language
"Scientific Word";type "GRAPHIC";maintain-aspect-ratio TRUE;display
"USEDEF";valid_file "T";width 4.403in;height 1.3553in;depth
0in;original-width 4.4589in;original-height 1.3526in;cropleft "0";croptop
"1";cropright "1";cropbottom "0";tempfilename
'T9OXEH05.wmf';tempfile-properties "XPR";}}

We said the surface must be smooth for there to be specular reflection. What
does smooth mean? Generally the size of the rough spots must be much smaller
than a wavelength to be considered smooth. So suppose we have a red laser.
How small do the surface variations have to be for the surface to be
considered smooth? The wavelength of a $HeNe$ laser is 
\begin{equation*}
\lambda _{HeNe}=633\unit{nm}
\end{equation*}%
This is very small. Modern optics for remote sensing are often manufactured
to $1/10$ of a wavelength, which would be $63\unit{nm}.$ Mirrors are smooth
for visible light.

How about a microwave beam of light like your cell phone uses?

\begin{eqnarray*}
c &=&\lambda f \\
\lambda &=&\frac{c}{f}=\frac{3\times 10^{8}\frac{\unit{m}}{\unit{s}}}{1\unit{%
GHz}} \\
&=&\allowbreak 0.3\unit{m}
\end{eqnarray*}%
Rough brick walls are smooth for cell phone light!

We can see that we must be careful in our definition of \textquotedblleft
smooth.\textquotedblright

\section{Law of reflection}

Experience shows that if we do have a smooth surface, that light bounces
much like a ball. This is why Newton though light was a particle. Suppose we
take a flat surface and we shine a light on it. We have a ray that
approaches at an angle $\theta _{i}$ measured from a line drawn
perpendicular to the flat surface. Then the reflected ray will leave the
surface with an angle $\theta _{r}$ also measured from the same line drawn
perpendicular to the flat surface. IN optics, we often measure angles from
lines drawn perpendicular to a surface.

Way back in Principles of Physics I we had forces that were perpendicular to
surfaces. We called them \emph{normal forces} where the word \emph{normal}
meant \textquotedblleft perpendicular to the surface.\textquotedblright\ We
can use the same word, \emph{normal,} to describe our perpendicular lines.
We can say that we measure $\theta _{i}$ and $\theta $ $_{r}$ from a \emph{%
normal line} (or sometimes, we just say that we measure from a \emph{normal}
and leave off the word \textquotedblleft line\textquotedblright ). And we
can write a very simple equation to describe reflection. 
\begin{equation*}
\theta _{r}=\theta _{i}
\end{equation*}

\FRAME{dhF}{1.8585in}{1.1857in}{0pt}{}{}{Figure}{\special{language
"Scientific Word";type "GRAPHIC";maintain-aspect-ratio TRUE;display
"USEDEF";valid_file "T";width 1.8585in;height 1.1857in;depth
0pt;original-width 1.8213in;original-height 1.1519in;cropleft "0";croptop
"1";cropright "1";cropbottom "0";tempfilename
'SVAJW80N.wmf';tempfile-properties "XPR";}}This is called the \emph{law of
reflection}. We will look at how superposition makes this happen later. But
from experiment we can see that this law of reflection works.

\subsection{Retroreflection}

Let's take and example\FRAME{dtbpF}{1.881in}{1.7625in}{0in}{}{}{Figure}{%
\special{language "Scientific Word";type "GRAPHIC";maintain-aspect-ratio
TRUE;display "USEDEF";valid_file "T";width 1.881in;height 1.7625in;depth
0in;original-width 1.8438in;original-height 1.7244in;cropleft "0";croptop
"1";cropright "1";cropbottom "0";tempfilename
'SW0DXR01.wmf';tempfile-properties "XPR";}}

Let's take our system to be two mirrors set at a right angle. We have a beam
of light incident at angle $\theta _{1}$. By the law of reflection, it must
leave the mirror at $\theta _{2}=\theta _{1}.$ We can see that $\alpha $
must be $90\unit{%
%TCIMACRO{\U{b0}}%
%BeginExpansion
{{}^\circ}%
%EndExpansion
}-\theta _{2}$ and it is clear that $\theta _{3}=\alpha .$ By the law of
reflection, $\theta _{3}=\theta _{4}.$ Then, since 
\begin{equation*}
90\unit{%
%TCIMACRO{\U{b0}}%
%BeginExpansion
{{}^\circ}%
%EndExpansion
}=\theta _{2}+\alpha
\end{equation*}%
But now look at the triangle in the corner (that has angles $\alpha $ and $%
\beta $). It's interior angles have to sum to $180,$ and one of the angles
is $90\unit{%
%TCIMACRO{\U{b0}}%
%BeginExpansion
{{}^\circ}%
%EndExpansion
}$ right in the corner. So 
\begin{equation*}
180\unit{%
%TCIMACRO{\U{b0}}%
%BeginExpansion
{{}^\circ}%
%EndExpansion
}=90\unit{%
%TCIMACRO{\U{b0}}%
%BeginExpansion
{{}^\circ}%
%EndExpansion
}+\alpha +\beta
\end{equation*}%
Then 
\begin{equation*}
90\unit{%
%TCIMACRO{\U{b0}}%
%BeginExpansion
{{}^\circ}%
%EndExpansion
}=\alpha +\beta
\end{equation*}%
so we can see that $\beta $ has to be equal to $\theta _{2}$%
\begin{equation*}
\beta =\theta _{2}=\theta _{1}
\end{equation*}%
and%
\begin{eqnarray*}
90\unit{%
%TCIMACRO{\U{b0}}%
%BeginExpansion
{{}^\circ}%
%EndExpansion
} &=&\beta +\theta _{3} \\
&=&\theta _{1}+\theta _{3}
\end{eqnarray*}%
and of course by the law of reflection $\theta _{3}=\theta _{4}$then the
total angular change is%
\begin{equation*}
\theta _{1}+\theta _{2}+\theta _{3}+\theta _{4}
\end{equation*}%
but if we rearrange this we have%
\begin{eqnarray*}
&&\theta _{1}+\theta _{3}+\theta _{2}+\theta _{4} \\
&=&90\unit{%
%TCIMACRO{\U{b0}}%
%BeginExpansion
{{}^\circ}%
%EndExpansion
}+90\unit{%
%TCIMACRO{\U{b0}}%
%BeginExpansion
{{}^\circ}%
%EndExpansion
}=180\unit{%
%TCIMACRO{\U{b0}}%
%BeginExpansion
{{}^\circ}%
%EndExpansion
}
\end{eqnarray*}%
or the outgoing ray is sent back toward the source! If we do this in three
dimensions we have a corner cube. Here is a picture of a radar corner cube
array.

\FRAME{dhFU}{2.1926in}{2.7913in}{0pt}{\Qcb{{\protect\small Radar
retroreflector tower located in the center of Yucca Flat dry lake bed. Used
as a radar target by maneuvering aircraft during "inert" contact fusing bomb
drops at Yucca Flat. Sandia National Laboratories conducted the tests on the
lake bed from 1954 to 1956. (Image in the Public Domain in the United States)%
}}}{}{Figure}{\special{language "Scientific Word";type
"GRAPHIC";maintain-aspect-ratio TRUE;display "USEDEF";valid_file "T";width
2.1926in;height 2.7913in;depth 0pt;original-width 2.1543in;original-height
2.7484in;cropleft "0";croptop "1";cropright "1";cropbottom "0";tempfilename
'SVAJW80P.wmf';tempfile-properties "XPR";}}And we left visible corner cube
arrays on the Moon so we can bounce laser beams off of them and monitor the
Earth-Moon distance.

\FRAME{dhFU}{2.2219in}{2.7638in}{0pt}{\Qcb{{\protect\small Apollo
Retroreflector (Images in the Public Domain courtesy NASA)}}}{}{Figure}{%
\special{language "Scientific Word";type "GRAPHIC";maintain-aspect-ratio
TRUE;display "USEDEF";valid_file "T";width 2.2219in;height 2.7638in;depth
0pt;original-width 2.1819in;original-height 2.7207in;cropleft "0";croptop
"1";cropright "1";cropbottom "0";tempfilename
'SVAJW80Q.wmf';tempfile-properties "XPR";}}%
%TCIMACRO{%
%\TeXButton{Question 223.13.9}{\marginpar {
%\hspace{-0.5in}
%\begin{minipage}[t]{1in}
%\small{Question 223.13.9}
%\end{minipage}
%}}}%
%BeginExpansion
\marginpar {
\hspace{-0.5in}
\begin{minipage}[t]{1in}
\small{Question 223.13.9}
\end{minipage}
}%
%EndExpansion

\section{Reflections, Objects, and seeing}

Armed with the law of reflection, we can start to understand how we see
things. Using the ray concept, we can say that a ray of light must leave the
light source. That ray then reflects from something. Suppose you look at the
person sitting next to you in class. We should wonder, how is it that we can
see them? We can only detect (see) light that gets to our eyes. Let's trace
the light from it's source (the light fixture in our room) to our eyes to
see how we see our neighbors.

Light comes from the ceiling lights and goes in all directions. Some of that
light hits our neighbor. And some of that light that hits our neighbor will
reflect. But is the person a specular or diffuse reflector? Once again, we
can only give an answer relative to the wavelength of light. For visible
light, your neighbors do not look like mirrors. They are diffuse reflectors.
Light bounces off of them in every direction. Some of that light that
reflects from your neighbor reflects in the direction of your eyes. That
light can be detected. Your eye is designed to take this diverging set of
rays and condense it into a picture of the person that your brain can
interpret.\FRAME{dhF}{3.173in}{1.3534in}{0pt}{}{}{Figure}{\special{language
"Scientific Word";type "GRAPHIC";maintain-aspect-ratio TRUE;display
"USEDEF";valid_file "T";width 3.173in;height 1.3534in;depth
0pt;original-width 3.128in;original-height 1.3188in;cropleft "0";croptop
"1";cropright "1";cropbottom "0";tempfilename
'SVAJW80R.wmf';tempfile-properties "XPR";}}We tend to not draw the rays that
bounce off the diffuse reflector but that don't get to our eyes, because we
don't see them. So a ray diagram is usually much simpler.\FRAME{dhF}{2.2926in%
}{1.9951in}{0pt}{}{}{Figure}{\special{language "Scientific Word";type
"GRAPHIC";maintain-aspect-ratio TRUE;display "USEDEF";valid_file "T";width
2.2926in;height 1.9951in;depth 0pt;original-width 2.2528in;original-height
1.9562in;cropleft "0";croptop "1";cropright "1";cropbottom "0";tempfilename
'SVAJW80S.wmf';tempfile-properties "XPR";}}This is easy to understand, but
we must keep in mind the wildly fluctuating superposition of waves that is
masked by our macroscopic view.

We can use the idea of a ray diagram to solve problems. Suppose you hold a
mirror half a meter in front of you and look at your reflection. Where would
the reflection appear to be?\FRAME{dhF}{1.9847in}{1.9527in}{0pt}{}{}{Figure}{%
\special{language "Scientific Word";type "GRAPHIC";maintain-aspect-ratio
TRUE;display "USEDEF";valid_file "T";width 1.9847in;height 1.9527in;depth
0pt;original-width 1.9458in;original-height 1.9147in;cropleft "0";croptop
"1";cropright "1";cropbottom "0";tempfilename
'SVAJW80T.wmf';tempfile-properties "XPR";}}Our mind assumes that the light
only travels in straight lines. This is not always true, as we shall see.
But in the light processing center in our brain the algorithm assumes the
light traveled in a straight line. Then, we can use the idea of rays to see
where the light appears to be from. We follow the light backwards from our
eyes to the mirror, but because our mind processing doesn't know about the
reflection we keep tracing the light ray through the mirror. Of course the
light doesn't come from behind the mirror. But our brain processing thinks
it does. The light originally came from us, so it is in a way shaped like
us. Our brain processing tells us there is another one of us about a half a
meter behind the mirror. Our frontal lobe knows there is not a second one of
us, so we interpret the light signal from our brain processing as an image
of us in the mirror.

The image is half a meter behind the mirror. Now suppose we look at an image
of that image in a mirror behind us. \FRAME{dhF}{2.6273in}{1.004in}{0pt}{}{}{%
Figure}{\special{language "Scientific Word";type
"GRAPHIC";maintain-aspect-ratio TRUE;display "USEDEF";valid_file "T";width
2.6273in;height 1.004in;depth 0pt;original-width 2.5858in;original-height
0.9712in;cropleft "0";croptop "1";cropright "1";cropbottom "0";tempfilename
'SVAJW80U.wmf';tempfile-properties "XPR";}}The ray diagram makes it easy to
see that the image will appear to be $2\unit{m}$ behind the big mirror. You
might not feel that this was so easy, but you might find it is not so bad in
a problem in your near future.

%TCIMACRO{%
%\TeXButton{Basic Equations}{\vspace{0.5in}\hspace{-1.3in}{\LARGE Basic Equations\vspace{0.25in}}}}%
%BeginExpansion
\vspace{0.5in}\hspace{-1.3in}{\LARGE Basic Equations\vspace{0.25in}}%
%EndExpansion
\begin{equation}
E=hf
\end{equation}
\begin{equation}
h=6.63\times 10^{-34}\unit{J}\unit{s}
\end{equation}

\begin{equation*}
c=299792458\frac{\unit{m}}{\unit{s}}
\end{equation*}%
\begin{equation*}
\theta _{r}=\theta _{i}
\end{equation*}

\chapter{12 Refraction 3.1.3, 3.1.4}

%TCIMACRO{\TeXButton{\chaptermark{Refraction}}{\chaptermark{Refraction}}}%
%BeginExpansion
\chaptermark{Refraction}%
%EndExpansion

\FRAME{dhF}{3.83in}{2.8871in}{0pt}{}{}{Figure}{\special{language "Scientific
Word";type "GRAPHIC";maintain-aspect-ratio TRUE;display "USEDEF";valid_file
"T";width 3.83in;height 2.8871in;depth 0pt;original-width
3.781in;original-height 2.8444in;cropleft "0";croptop "1";cropright
"1";cropbottom "0";tempfilename 'SVJXCR06.wmf';tempfile-properties "XPR";}}%
We have studied reflection of light. Now let's consider what happens when
light enters a material like glass or water.

%TCIMACRO{%
%\TeXButton{Fundamental Concepts}{\vspace{0.10in}\hspace{-0.37in}{\Large Fundamental Concepts\vspace{0.2in}}}}%
%BeginExpansion
\vspace{0.10in}\hspace{-0.37in}{\Large Fundamental Concepts\vspace{0.2in}}%
%EndExpansion

\begin{itemize}
\item Refraction

\item Total internal reflection
\end{itemize}

In the last lecture we split up electromagnetic waves into different kinds
of light. We categorized the different kinds of light by noting they have
different wavelength ranges. We found visible light, of course, but also
infrared and ultraviolet and x-rays and radio waves. Not all surfaces
reflect all the light. Some, like the lenses shown below, reflect some light
at visible wavelengths, but are transparent so most of the light travels
through them. \FRAME{dhF}{3.2846in}{1.6042in}{0pt}{}{}{Figure}{\special%
{language "Scientific Word";type "GRAPHIC";maintain-aspect-ratio
TRUE;display "USEDEF";valid_file "T";width 3.2846in;height 1.6042in;depth
0pt;original-width 3.2396in;original-height 1.5679in;cropleft "0";croptop
"1";cropright "1";cropbottom "0";tempfilename
'SVAJW80W.wmf';tempfile-properties "XPR";}}We need a way to deal with
transparent materials. This is tricky, because different wavelengths of
light penetrate different materials in different ways (we will see more on
why later). As an example, this is also a lens\FRAME{dhFU}{1.0732in}{1.6285in%
}{0pt}{\Qcb{{\protect\small Infrared lens, but visible mirror (Image
courtesy US Navy, image in the public domain)}}}{}{Figure}{\special{language
"Scientific Word";type "GRAPHIC";maintain-aspect-ratio TRUE;display
"USEDEF";valid_file "T";width 1.0732in;height 1.6285in;depth
0pt;original-width 5.4587in;original-height 8.3195in;cropleft "0";croptop
"1";cropright "1";cropbottom "0";tempfilename
'SVAJW80X.bmp';tempfile-properties "XPR";}}but it clearly is not transparent
at visible wavelengths. It is transparent in the infrared. So what might be
transparent at one wavelength might not be at another.

When light travels into a material, we say it is transmitted. This is
nothing new. Light is a wave, so the part that goes through is called the
transmitted part just like when we had waves go from one rope to another.
The situation is shown schematically below.

\FRAME{dtbpF}{2.7207in}{2.2719in}{0in}{}{}{Figure}{\special{language
"Scientific Word";type "GRAPHIC";maintain-aspect-ratio TRUE;display
"USEDEF";valid_file "T";width 2.7207in;height 2.2719in;depth
0in;original-width 2.6783in;original-height 2.2321in;cropleft "0";croptop
"1";cropright "1";cropbottom "0";tempfilename
'SW7RGP00.wmf';tempfile-properties "XPR";}}

In the figure we see a ray incident on an air-glass boundary. Some of the
light is reflected just as we saw in our last lecture. But some passes into
the glass. This isn't surprising, light is a wave and at boundaries we
expect reflected and transmitted waves. Notice that the angle between the
normal and the new transmitted ray is \emph{not} equal to the incident angle
for the incoming ray. We say the ray has been bent or \emph{refracted} by
the change in material. Many experiments were performed to find a
relationship between the incident and the refracted angles. It was found
that 
\begin{equation}
\frac{\sin \left( \theta _{2}\right) }{\sin \left( \theta _{1}\right) }=%
\frac{v_{2}}{v_{1}}=\text{constant}
\end{equation}%
Many optics books write this as 
\begin{equation}
\frac{\sin \left( \theta _{t}\right) }{\sin \left( \theta _{i}\right) }=%
\frac{v_{t}}{v_{i}}=\text{constant}
\end{equation}%
where the subscript $i$ stands for \textquotedblleft
incident\textquotedblright\ and the subscript $t$ stands for
\textquotedblleft transmitted.\textquotedblright\ \FRAME{dtbpF}{3.4229in}{%
2.8548in}{0pt}{}{}{Figure}{\special{language "Scientific Word";type
"GRAPHIC";maintain-aspect-ratio TRUE;display "USEDEF";valid_file "T";width
3.4229in;height 2.8548in;depth 0pt;original-width 3.3771in;original-height
2.8115in;cropleft "0";croptop "1";cropright "1";cropbottom "0";tempfilename
'SW7RJB01.wmf';tempfile-properties "XPR";}}\ Note that we are using the fact
that the speed of light changes in a material. We should probably recall why
this should occur

\section{Speed of light in a material}

In a vacuum, light travels as a disturbance in the electromagnetic field
with nothing to encounter. In a material (like glass) the light waves
continually hit atoms. We have not studied antennas, but I\ think many of
you know that an antenna works because the electrons in the metal act like
driven harmonic oscillators. The incoming radio waves drive the electron
motion. Here each atom has electrons, and the atoms act like little
antennas, their electrons moving and absorbing the light. But the atom
cannot keep the extra energy\footnote{%
If you get to take Modern Physics (PH279) and Quantum Physics (PH433) you
will learn how this works.}, so it is readmitted. It travels to the next
atom and the process repeats. Quantum mechanics tells us that there is a
time delay in the re-emission of the light. This re-emitted light mixes with
what is left of the original wave. And this causes the propagation of the
combined light wave to slow down. Thus the speed of light is slower in a
material.\footnote{%
This is a bit of an approximation. Some of the incident light misses the
atoms. So we really get a superposition of the new re-edmitted wave and the
original wave. But the end result is that the light appears to go slower.}%
\FRAME{dhF}{4.6267in}{1.8135in}{0pt}{}{}{Figure}{\special{language
"Scientific Word";type "GRAPHIC";maintain-aspect-ratio TRUE;display
"USEDEF";valid_file "T";width 4.6267in;height 1.8135in;depth
0pt;original-width 4.574in;original-height 1.7763in;cropleft "0";croptop
"1";cropright "1";cropbottom "0";tempfilename
'SVAJW80Z.wmf';tempfile-properties "XPR";}}As a mechanical analog, consider
a rolling barrel. Suppose you are helping a neighbor move and they have
barrels of wheat. And as often happens in a neighbor move, one of the
barrels get's dropped. It lands on the driveway and rolls toward the grass.%
\FRAME{dhF}{2.4604in}{1.8005in}{0pt}{}{}{Figure}{\special{language
"Scientific Word";type "GRAPHIC";maintain-aspect-ratio TRUE;display
"USEDEF";valid_file "T";width 2.4604in;height 1.8005in;depth
0pt;original-width 2.4189in;original-height 1.7634in;cropleft "0";croptop
"1";cropright "1";cropbottom "0";tempfilename
'SVAJW810.wmf';tempfile-properties "XPR";}}As the barrel rolls from a flat
low-friction concrete to a higher-friction grass lawn, the friction slows
the barrel. If the barrel hits the lawn parallel to the boundary (so it's
velocity vector is perpendicular to the boundary), then the barrel continues
in the same direction at the slower speed. But if it hits at an angle, the
leading edge is slowed first. \FRAME{dhF}{2.3575in}{1.7236in}{0pt}{}{}{Figure%
}{\special{language "Scientific Word";type "GRAPHIC";maintain-aspect-ratio
TRUE;display "USEDEF";valid_file "T";width 2.3575in;height 1.7236in;depth
0pt;original-width 2.725in;original-height 1.9847in;cropleft "0";croptop
"1";cropright "1";cropbottom "0";tempfilename
'SVAJW811.wmf';tempfile-properties "XPR";}}This makes the trailing edge
travel faster than the leading edge, and the barrel turns slightly.\FRAME{dhF%
}{2.4267in}{1.7711in}{0pt}{}{}{Figure}{\special{language "Scientific
Word";type "GRAPHIC";maintain-aspect-ratio TRUE;display "USEDEF";valid_file
"T";width 2.4267in;height 1.7711in;depth 0pt;original-width
3.7948in;original-height 2.7622in;cropleft "0";croptop "1";cropright
"1";cropbottom "0";tempfilename 'SVAJW812.wmf';tempfile-properties "XPR";}}

We expect the same behavior from light.\FRAME{dhF}{2.8236in}{2.0237in}{0pt}{%
}{}{Figure}{\special{language "Scientific Word";type
"GRAPHIC";maintain-aspect-ratio TRUE;display "USEDEF";valid_file "T";width
2.8236in;height 2.0237in;depth 0pt;original-width 2.7804in;original-height
1.9847in;cropleft "0";croptop "1";cropright "1";cropbottom "0";tempfilename
'SVAJW813.wmf';tempfile-properties "XPR";}}We can see that the left hand
side of the wave hits the slower (green) material first and slows down. The
rest of the wave front moves quicker. The result is the turning of the wave.%
\footnote{%
Once again this is a bit of a simplification, but it will due for now. If
you are lucky enough to take a junior level optics class, you will revisit
this. And we will look at this again in our next lecture.}

\section{Change of wavelength}

We have found that when a wave enters a material, its speed may change. But
we remember from wave theory%
\begin{equation}
v=\lambda f
\end{equation}

And it is time to review: does $\lambda $ change, or does $f$ change? If you
will recall, we found that the change in speed at the boundary changes the
wavelength. Recall that if we go from a fast material to a slow material,
the forward part of the wave slows and the rest of the wave catches up to
it. \FRAME{dhF}{2.2727in}{1.2644in}{0pt}{}{}{Figure}{\special{language
"Scientific Word";type "GRAPHIC";maintain-aspect-ratio TRUE;display
"USEDEF";valid_file "T";width 2.2727in;height 1.2644in;depth
0pt;original-width 3.4618in;original-height 1.9147in;cropleft "0";croptop
"1";cropright "1";cropbottom "0";tempfilename
'SVAJW814.wmf';tempfile-properties "XPR";}}This will compress pulses, and
lower the wavelength. Now that we know more about light we can also argue
that $f$ cannot change because 
\begin{equation*}
E=hf
\end{equation*}%
If $f$ changed, then we would either require an input of energy or we would
store energy at the boundary because%
\begin{equation*}
\Delta f=\frac{\Delta E}{h}
\end{equation*}%
This can't be true. If the wavelength changes, there is no change in energy.

And since 
\begin{equation*}
v_{i}=\lambda _{i}f
\end{equation*}%
and 
\begin{equation*}
v_{t}=\lambda _{t}f
\end{equation*}%
then the ratio%
\begin{equation*}
\frac{v_{i}}{v_{t}}=\frac{\lambda _{i}}{\lambda _{t}}
\end{equation*}%
and we have a solution for the wavelength in the material%
\begin{equation*}
\lambda _{t}=\lambda _{i}\frac{v_{t}}{v_{i}}
\end{equation*}%
which agrees with our previous analysis.

\section{Index of refraction and Snell's Law}

Because the equation%
\begin{equation}
\frac{\sin \left( \theta _{i}\right) }{\sin \left( \theta _{t}\right) }=%
\frac{v_{t}}{v_{i}}=\text{constant}
\end{equation}%
has a constant ratio of velocities, it is convenient to define a term that
represents that ratio. Think, it would be kind of nice to say how much
slower the light is going once it enters the slower glass material. We could
take the current velocity and divide by the speed of light in vacuum to get
the percentage of the speed of light we are going%
\begin{equation*}
\frac{v}{c}=\text{percentage of the speed of light}
\end{equation*}%
and that makes a lot of sense. But that is not what the optics community
did. They defined a term. 
\begin{equation}
n\equiv \frac{c}{v}
\end{equation}%
which is $1/$percentage of the speed of light. Which is not as intuitive as
the percentage, itself, was. But it caries the same information 
\begin{equation*}
n=\frac{1}{\text{percentage of the speed of light}}
\end{equation*}%
So we can use this to get an idea of how much slower our transmitted light
is going. This odd ratio of the speed of light to the average speed in the
material has a name. It is called the \emph{index of refraction.} We can
write Snell's law in terms of the index of refraction of air and glass (or
any other two materials). Start with%
\begin{equation}
\frac{\sin \left( \theta _{i}\right) }{\sin \left( \theta _{t}\right) }=%
\frac{v_{i}}{v_{t}}
\end{equation}%
but let's multiply the right hand side by $1,$ only let's write $1$ as 
\begin{equation*}
1=\frac{c}{c}
\end{equation*}%
which is a very odd way to write $1.$ But then%
\begin{equation}
\frac{\sin \left( \theta _{i}\right) }{\sin \left( \theta _{t}\right) }=%
\frac{v_{i}c}{v_{t}c}
\end{equation}%
which we could write as%
\begin{equation}
\frac{\sin \left( \theta _{i}\right) }{\sin \left( \theta _{t}\right) }%
=\left( \frac{c}{v_{t}}\right) \left( \frac{v_{i}}{c}\right)
\end{equation}%
which is just%
\begin{equation}
\frac{\sin \left( \theta _{i}\right) }{\sin \left( \theta _{t}\right) }%
=\left( n_{t}\right) \left( \frac{1}{n_{i}}\right)
\end{equation}%
and for our case $n_{i}=n_{air}\approx 1$ so 
\begin{equation}
\frac{\sin \left( \theta _{i}\right) }{\sin \left( \theta _{t}\right) }%
=n_{glass}
\end{equation}

Suppose we don't have a vacuum (or air that is close to a vacuum). We can
write our formula as 
\begin{equation}
n_{i}\sin \left( \theta _{i}\right) =n_{t}\sin \left( \theta _{t}\right)
\end{equation}%
where we have 
\begin{equation}
n_{i}\equiv \frac{c}{v_{i}}
\end{equation}%
and 
\begin{equation}
n_{t}\equiv \frac{c}{v_{t}}
\end{equation}

This is called \emph{Snell's law of refraction}.

Using the index of refraction we can write our equation relating the ratio
of velocities and wavelengths as%
\begin{equation}
\frac{v_{i}}{v_{t}}=\frac{\lambda _{i}}{\lambda _{t}}=\frac{\frac{c}{n_{i}}}{%
\frac{c}{n_{t}}}=\frac{n_{t}}{n_{i}}
\end{equation}%
which gives 
\begin{equation}
\lambda _{i}n_{i}=\lambda _{t}n_{t}
\end{equation}%
and if we have vacuum and a single material we can find the index of
refraction from 
\begin{equation}
n=\frac{\lambda }{\lambda _{in}}
\end{equation}%
where $\lambda _{in}=\lambda _{t}$ is the wavelength in the material and $%
\lambda =\lambda _{i}$ is the wavelength outside the material.

\section{Total Internal Reflection}

Up to now we have assumed that light was coming from a region of low index
of refraction (air) into a region of high index of refraction (water or
glass). We should pause to look at what can happen if we go the other way.%
\FRAME{dtbpF}{3in}{2.5339in}{0pt}{}{}{Figure}{\special{language "Scientific
Word";type "GRAPHIC";maintain-aspect-ratio TRUE;display "USEDEF";valid_file
"T";width 3in;height 2.5339in;depth 0pt;original-width
3.4601in;original-height 2.9187in;cropleft "0";croptop "1";cropright
"1";cropbottom "0";tempfilename 'SW7RLJ02.wmf';tempfile-properties "XPR";}}

We start with Snell's law%
\begin{equation}
n_{i}\sin \theta _{i}=n_{t}\sin \theta _{t}
\end{equation}%
but this time $n_{glass}=n_{i}$ and $n_{t}\approx 1$ so 
\begin{equation}
n_{glass}\sin \theta _{i}=\sin \theta _{t}
\end{equation}%
which gives 
\begin{equation}
\theta _{t}=\sin ^{-1}\left( n_{glass}\sin \theta _{i}\right)
\end{equation}%
If we take $n_{glass}=1.33$ (which is really $n$ for water) we can plot this
expression as a function of $\theta _{i}$

\FRAME{dtbpF}{2.8643in}{1.9804in}{0pt}{}{}{Figure}{\special{language
"Scientific Word";type "GRAPHIC";maintain-aspect-ratio TRUE;display
"USEDEF";valid_file "T";width 2.8643in;height 1.9804in;depth
0pt;original-width 3.6642in;original-height 2.5235in;cropleft "0";croptop
"1";cropright "1";cropbottom "0";tempfilename
'SW7RMF03.wmf';tempfile-properties "XPR";}}we see that at $\theta
_{i}=\allowbreak 0.850\,91\unit{rad}$ $(48.\,\allowbreak 754\unit{%
%TCIMACRO{\U{b0}}%
%BeginExpansion
{{}^\circ}%
%EndExpansion
})$ the curve becomes infinitely steep. If we use this value in our equation
this gives%
\begin{eqnarray}
\theta _{t} &=&\sin ^{-1}\left( n\sin \left( \allowbreak 0.850\,91\right)
\right) \\
&=&1.\,\allowbreak 570\,8\unit{rad} \\
&=&90\unit{%
%TCIMACRO{\U{b0}}%
%BeginExpansion
{{}^\circ}%
%EndExpansion
}
\end{eqnarray}%
The light skims along the edge of the water! \FRAME{dhF}{2.3756in}{1.9251in}{%
0pt}{}{}{Figure}{\special{language "Scientific Word";type
"GRAPHIC";maintain-aspect-ratio TRUE;display "USEDEF";valid_file "T";width
2.3756in;height 1.9251in;depth 0pt;original-width 2.3359in;original-height
1.887in;cropleft "0";croptop "1";cropright "1";cropbottom "0";tempfilename
'SVAJW817.wmf';tempfile-properties "XPR";}}We can find the value of $\theta
_{i}$ that makes this happen without graphing. Set $\theta _{t}=90\unit{%
%TCIMACRO{\U{b0}}%
%BeginExpansion
{{}^\circ}%
%EndExpansion
}$ then%
\begin{equation}
\theta _{i}=\theta _{c}\equiv \sin ^{-1}\left( \frac{1}{n}\right)
\end{equation}%
We give this value of $\theta _{i}$ a special name. It is the \emph{critical
angle} for internal reflection. But what happens if we go farther than this (%
$\theta _{i}>\theta _{c}$). We will no longer have a transmitted ray. But
the wave energy has to go somewhere! The ray will be reflected. This is why
when you dive into a pool and look up, you see a region of the roof of the
pool area (or sky) but off to the side of the pool the surface looks
mirrored. It is also why you sometimes see the sides of a fish tank appear
to be mirrored when you look through the front. \FRAME{dhF}{1.7616in}{%
1.3258in}{0pt}{}{}{Figure}{\special{language "Scientific Word";type
"GRAPHIC";maintain-aspect-ratio TRUE;display "USEDEF";valid_file "T";width
1.7616in;height 1.3258in;depth 0pt;original-width 1.7236in;original-height
1.2912in;cropleft "0";croptop "1";cropright "1";cropbottom "0";tempfilename
'SVAJW818.wmf';tempfile-properties "XPR";}}More importantly, it is why cut
gems (like diamonds) sparkle. They capture the light with facets that are
cut at angles that create total internal reflection. The light that enters
the gem comes back out the front.

\section{Fiber Optics}

Beyond pretty pebbles, this effect is very useful! It is the heart and soul
of fiber optics. \FRAME{dtbpF}{3.1993in}{1.7882in}{0in}{}{}{Figure}{\special%
{language "Scientific Word";type "GRAPHIC";maintain-aspect-ratio
TRUE;display "USEDEF";valid_file "T";width 3.1993in;height 1.7882in;depth
0in;original-width 3.233in;original-height 1.7935in;cropleft "0";croptop
"1";cropright "1";cropbottom "0";tempfilename
'SVAJW819.wmf';tempfile-properties "XPR";}}

An interior material with a lower index of refraction is inclosed in a
cladding with a higher index. This creates a light pipe that traps the light
in the fiber. \FRAME{dhF}{3.4662in}{1.535in}{0pt}{}{}{Figure}{\special%
{language "Scientific Word";type "GRAPHIC";maintain-aspect-ratio
TRUE;display "USEDEF";valid_file "T";width 3.4662in;height 1.535in;depth
0pt;original-width 3.4203in;original-height 1.4987in;cropleft "0";croptop
"1";cropright "1";cropbottom "0";tempfilename
'SVAJW81A.wmf';tempfile-properties "XPR";}}

Modern fibers don't always have a hard boundary. The fibers have a gradual
change in index of refraction that changes the direction of the light
gradually. This keeps the light in the fiber but tends to direct along the
fiber so the beam is not crisscrossing as it goes.

The cutting edge of fiber design today uses hollow fibers or fibers filled
with different index material.

\FRAME{dtbpFU}{1.6319in}{1.5973in}{0pt}{\Qcb{{\protect\small Hollow-Core
Fiber (Courtesy Defense Advanced Research Projects Agency (DARPA), image in
the public domain)}}}{}{Figure}{\special{language "Scientific Word";type
"GRAPHIC";maintain-aspect-ratio TRUE;display "USEDEF";valid_file "T";width
1.6319in;height 1.5973in;depth 0pt;original-width 1.6338in;original-height
1.5992in;cropleft "0";croptop "1";cropright "1";cropbottom "0";tempfilename
'SVAJW81B.wmf';tempfile-properties "XPR";}}%
%TCIMACRO{%
%\TeXButton{Basic Equations}{\vspace{0.5in}\hspace{-1.3in}{\LARGE Basic Equations\vspace{0.25in}}}}%
%BeginExpansion
\vspace{0.5in}\hspace{-1.3in}{\LARGE Basic Equations\vspace{0.25in}}%
%EndExpansion

\begin{equation*}
n_{i}\sin \theta _{i}=n_{t}\sin \theta _{t}
\end{equation*}%
\begin{equation*}
\lambda _{i}n_{i}=\lambda _{t}n_{t}
\end{equation*}%
\begin{equation*}
\theta _{c}\equiv \sin ^{-1}\left( \frac{1}{n}\right)
\end{equation*}

\chapter{13 Dispersion and Huygen's Principle 3.1.5 3.1.6}

We have discovered some important things about how light travels through
materials. But there is something we missed. It turns out the index of
refraction depends on wavelength! This has some profound effects on what
light will do in a material. We will look at this wavelength dependence on
refraction in this lecture.

%TCIMACRO{\TeXButton{\chaptermark{Dispersion}}{\chaptermark{Dispersion}}}%
%BeginExpansion
\chaptermark{Dispersion}%
%EndExpansion

%TCIMACRO{%
%\TeXButton{Fundamental Concepts}{\vspace{0.10in}\hspace{-0.37in}{\Large Fundamental Concepts\vspace{0.2in}}}}%
%BeginExpansion
\vspace{0.10in}\hspace{-0.37in}{\Large Fundamental Concepts\vspace{0.2in}}%
%EndExpansion

\begin{itemize}
\item The index of refraction is wavelength dependent

\item That different wavelengths bend different amounts when refracted is
called \emph{dispersion}

\item White light is a superposition of many other frequencies
\end{itemize}

\section{Dispersion}

\FRAME{dhF}{4.3613in}{2.2745in}{0pt}{}{}{Figure}{\special{language
"Scientific Word";type "GRAPHIC";maintain-aspect-ratio TRUE;display
"USEDEF";valid_file "T";width 4.3613in;height 2.2745in;depth
0pt;original-width 4.3102in;original-height 2.2347in;cropleft "0";croptop
"1";cropright "1";cropbottom "0";tempfilename
'SVAJW81H.wmf';tempfile-properties "XPR";}}

Who hasn't played with a prism? We immediately recognize a rainbow. But why
does the prism make a rainbow? The secret lies in the nature of the
refractive index. We know that some wavelengths will pass through a
material, say, a piece of glass. But some (like infrared wavelengths) won't.
But there is more to this. Not all the wavelengths that do pass through the
glass will have the same transmitted speed. The speed in the material
depends on the wavelength! Think, blue wavelengths are about the size of
atoms. This means blue photons will hit atoms more frequently. And that will
slow them down more. Red wavelengths are bigger than the atoms in a material
and because of this red photons are not as likely to be absorbed by the
atoms. So red photos will go faster in the material. And this means the
index of refraction depends on the speed.\FRAME{fhFU}{4.3535in}{2.6195in}{0pt%
}{\Qcb{{\protect\small Index of refraction as a function of wavelength (
Ohara optical glass http://www.oharacorp.com/fused-silica-quartz.html data
and Schott optical glass data
http://www.uqgoptics.com/materials\_glasses\_schott.aspx)}}}{\Qlb{Dispersion
of Glass}}{Figure}{\special{language "Scientific Word";type
"GRAPHIC";maintain-aspect-ratio TRUE;display "USEDEF";valid_file "T";width
4.3535in;height 2.6195in;depth 0pt;original-width 5.1436in;original-height
3.084in;cropleft "0";croptop "1";cropright "1";cropbottom "0";tempfilename
'dispersion0.wmf';tempfile-properties "XNPR";}}

We can see light from about $400\unit{nm}$ to about $700\unit{nm}.$ And
notice for the three kinds of glass in the figure the index of refraction
changes over this region of the electromagnetic spectrum. Think of what this
means for Snell's law. We learned last lecture that 
\begin{equation*}
n_{i}\sin \theta _{i}=n_{t}\sin \theta _{t}
\end{equation*}%
and if the light goes from air $\left( n_{i}\approx 1\right) $ into some
kind of glass, then 
\begin{equation*}
\sin \theta _{i}=n_{t}\sin \theta _{t}
\end{equation*}%
and the transmitted angle just depends on $n_{t}.$ This means that as light
enters a material, different wavelengths will be refracted at different
angles. White light is a superposition of many light waves with many
different wavelengths. Thus white light is pulled apart by refraction into a
rainbow. This process is called \emph{dispersion}.

The graph tells us that blue light bends more than red light. \FRAME{dhF}{%
2.4379in}{1.2816in}{0pt}{}{}{Figure}{\special{language "Scientific
Word";type "GRAPHIC";maintain-aspect-ratio TRUE;display "USEDEF";valid_file
"T";width 2.4379in;height 1.2816in;depth 0pt;original-width
6.5207in;original-height 3.4143in;cropleft "0";croptop "1";cropright
"1";cropbottom "0";tempfilename 'SVAJW81I.wmf';tempfile-properties "XPR";}}

We call the change in direction measured from the original direction of
travel, $\delta ,$ the \emph{angle of deviation}. The colors we can see are
called the visible spectrum. Note that this is a second way to make a
visible light spectrum. The first way used a diffraction grating and the
wave nature of light. Although this dispersive element (prism) method of
making a spectrum uses the ides of geometric optics, it is still the wave
nature of light that makes the waves of different wavelengths refract
differently.

Let's look at a natural rainbow. The dispersion is caused buy small droplets
of water. The white sunlight enters the drop and is dispersed. It bounces
off the back of the drop and then leaves the drop, again being dispersed.
Red light leaves the drop at about $42\unit{%
%TCIMACRO{\U{b0}}%
%BeginExpansion
{{}^\circ}%
%EndExpansion
}$ from its input direction, and blue light leaves at about $40\unit{%
%TCIMACRO{\U{b0}}%
%BeginExpansion
{{}^\circ}%
%EndExpansion
}.$\FRAME{dtbpF}{3.2681in}{2.3921in}{0in}{}{}{Figure}{\special{language
"Scientific Word";type "GRAPHIC";maintain-aspect-ratio TRUE;display
"USEDEF";valid_file "T";width 3.2681in;height 2.3921in;depth
0in;original-width 3.2232in;original-height 2.3514in;cropleft "0";croptop
"1";cropright "1";cropbottom "0";tempfilename
'TF1CXN00.wmf';tempfile-properties "XPR";}}This effectively separates all
the different wavelengths and we see a rainbow at angles between $40\unit{%
%TCIMACRO{\U{b0}}%
%BeginExpansion
{{}^\circ}%
%EndExpansion
}$ and $42\unit{%
%TCIMACRO{\U{b0}}%
%BeginExpansion
{{}^\circ}%
%EndExpansion
}$ from the incoming light direction.

\subsection{Calculation of $n$ using a prism}

Let's do a problem using the idea of refraction in a prism. Let's find the
index of refraction of a the prism material. Suppose we make a prism as
shown. We know the angle $\Phi $ and can measure the exit angle $\delta .$
In terms of these two variables, what is $n?$

\FRAME{dtbpF}{2.4595in}{1.9259in}{0in}{}{}{Figure}{\special{language
"Scientific Word";type "GRAPHIC";maintain-aspect-ratio TRUE;display
"USEDEF";valid_file "T";width 2.4595in;height 1.9259in;depth
0in;original-width 2.418in;original-height 1.8887in;cropleft "0";croptop
"1";cropright "1";cropbottom "0";tempfilename
'SWB81E00.wmf';tempfile-properties "XPR";}}

Using the notation indicated in the figure, we choose $\theta _{1}$ such
that the interior ray is horizontal.\footnote{%
WARNING! in an upcoming homework problem you may not be able make the same
assumptions!} This is a refraction problem, so we will want to use Snell's
law. 
\begin{equation*}
n_{1}\sin \theta _{1}=n_{2}\sin \theta _{2}
\end{equation*}%
And this has $n_{2}$ in it. That is what we want to find. We know $n_{1}$
because the prism is in air so $n_{1}\approx 1.$ And suppose we know $\Phi $
and $\delta $ from measurement. Our process will be to find $\theta _{1}$
and $\theta _{2}.$ The we can use 
\begin{equation*}
n_{2}=\frac{\sin \theta _{1}}{\sin \theta _{2}}
\end{equation*}%
to find $n_{2.}$ But we have to find the angles. Let's see one way to do
this. In geometry it is fair game to extend lines and even make some new
lines of our own. I have done this in the figure. You likely will have to do
your own line extensions and additions in the problems you do. But by using
the added lines we can realize that

\begin{equation}
\theta _{1}=\theta _{2}+\alpha
\end{equation}%
and that 
\begin{equation}
\delta =180-\beta
\end{equation}%
and it is also true that

\begin{equation}
180=\beta +2\alpha
\end{equation}%
or 
\begin{equation}
180-\beta =2\alpha
\end{equation}

Then

\begin{equation}
\delta =2\alpha
\end{equation}%
and%
\begin{equation}
\alpha =\frac{\delta }{2}
\end{equation}

Now also realize that

\begin{equation}
90=\alpha +\theta _{2}+\phi
\end{equation}

and

\begin{equation}
180=\Phi +2\alpha +2\phi
\end{equation}%
or%
\begin{equation}
90=\frac{\Phi }{2}+\alpha +\phi
\end{equation}%
Then%
\begin{eqnarray}
\alpha +\theta _{2}+\phi &=&\frac{\Phi }{2}+\alpha +\phi \\
\theta _{2} &=&\frac{\Phi }{2}
\end{eqnarray}%
We can put these in our equation for $\theta _{1}$%
\begin{eqnarray}
\theta _{1} &=&\theta _{2}+\alpha \\
&=&\frac{\Phi }{2}+\frac{\delta }{2} \\
&=&\frac{\Phi +\delta }{2}
\end{eqnarray}%
Now we can use Snell's Law%
\begin{eqnarray}
n_{2} &=&\frac{\sin \left( \theta _{1}\right) }{\sin \left( \theta
_{2}\right) } \\
&=&\frac{\sin \left( \frac{\Phi +\delta }{2}\right) }{\sin \left( \frac{\Phi 
}{2}\right) }
\end{eqnarray}%
and since we know $\Phi $ and $\delta $, we can find $n_{2}.$

\section{Huygen's Principle}

We have learned the laws of reflection and refraction. And we know the
scattering of the light by atoms in a material cause the bending of light
that we call refraction. But let's look at this again in more detail.

To understand what is going on at a fundamental level we need to remember
that light is a wave. And waves can hit and splash off of objects. Sound
waves could bounce off of things. So can light waves. Only, the light waves
that we can see are very small. So they bounce off of very small things like
atoms. Suppose we send our light into a material. It will hit the atoms and
bounce off. We know this, but there is more. The light hits the atoms and is
absorbed by exciting electrons. When the electrons return to their original
state the light is re-emitted. But this makes a new light wave coming from
the atoms. And that new wave is spherical (think, most waves start
spherical). If we have many atoms, the light will hit many atoms and make
many little spherical \emph{scattered} waves. An early researcher had an
idea of how to find the direction of the light using this underlying wave
nature of light. We could follow the light path by finding the little
spherical scattered waves and look for where their crests line up. This
would be constructive interference, and this is where the ray of light would
end up! \FRAME{dtbpF}{0.5047in}{1.2418in}{0pt}{}{}{Figure}{\special{language
"Scientific Word";type "GRAPHIC";maintain-aspect-ratio TRUE;display
"USEDEF";valid_file "T";width 0.5047in;height 1.2418in;depth
0pt;original-width 0.33in;original-height 0.8551in;cropleft "0";croptop
"1";cropright "1";cropbottom "0";tempfilename
'SVJWGL05.wmf';tempfile-properties "XPR";}}This early researcher was named
Huygen, so this method of finding the new wave is called \emph{Huygen's
principle.} He simplified the process by just drawing a line that connected
the crests of each little scattered wave and said that is where the new
combined wave front would be.

We know that light should reflect off a surface. 
\begin{equation*}
\theta _{i}=\theta _{r}
\end{equation*}

\FRAME{dtbpF}{2.841in}{1.8086in}{0pt}{}{}{Figure}{\special{language
"Scientific Word";type "GRAPHIC";maintain-aspect-ratio TRUE;display
"USEDEF";valid_file "T";width 2.841in;height 1.8086in;depth
0pt;original-width 3.6748in;original-height 2.3283in;cropleft "0";croptop
"1";cropright "1";cropbottom "0";tempfilename
'../../../../PH375/repository/OpticsText/text/SVAL9A3C.wmf';tempfile-properties "XPR";}%
}Let's see how Huygen's principle gives this law of reflection. Think about
the top row (or about $\lambda /2$ thickness) of scatterers. In the next
figure we have drawn such a set of scatters with an incoming plane wave.%
\FRAME{dtbpF}{3.3987in}{1.9104in}{0pt}{}{}{Figure}{\special{language
"Scientific Word";type "GRAPHIC";maintain-aspect-ratio TRUE;display
"USEDEF";valid_file "T";width 3.3987in;height 1.9104in;depth
0pt;original-width 4.6778in;original-height 2.6187in;cropleft "0";croptop
"1";cropright "1";cropbottom "0";tempfilename
'../../../../PH375/repository/OpticsText/text/SVAL9A3D.wmf';tempfile-properties "XPR";}%
}On the left, the plane wave crest is striking a scatter and has started a
spherical secondary wave. One period later the situation might look like
this \FRAME{dtbpF}{3.5898in}{2.0176in}{0pt}{}{}{Figure}{\special{language
"Scientific Word";type "GRAPHIC";maintain-aspect-ratio TRUE;display
"USEDEF";valid_file "T";width 3.5898in;height 2.0176in;depth
0pt;original-width 5.1992in;original-height 2.911in;cropleft "0";croptop
"1";cropright "1";cropbottom "0";tempfilename
'../../../../PH375/repository/OpticsText/text/SVAL9A3E.wmf';tempfile-properties "XPR";}%
}where now the first crest has moved on to the second scatterer and the
second crest is at the first scatter. Notice that the first scatter's
spherical secondary wave has grown. This process continues as the wave moves
onward. In another period, we would have three scattered waves \FRAME{dtbpF}{%
3.6547in}{2.0574in}{0pt}{}{}{Figure}{\special{language "Scientific
Word";type "GRAPHIC";maintain-aspect-ratio TRUE;display "USEDEF";valid_file
"T";width 3.6547in;height 2.0574in;depth 0pt;original-width
3.6071in;original-height 2.0185in;cropleft "0";croptop "1";cropright
"1";cropbottom "0";tempfilename
'../../../../PH375/repository/OpticsText/text/SVAL9A3F.wmf';tempfile-properties "XPR";}%
}and in another period, four.\FRAME{dtbpF}{3.9409in}{2.2182in}{0pt}{}{}{%
Figure}{\special{language "Scientific Word";type
"GRAPHIC";maintain-aspect-ratio TRUE;display "USEDEF";valid_file "T";width
3.9409in;height 2.2182in;depth 0pt;original-width 3.8917in;original-height
2.1776in;cropleft "0";croptop "1";cropright "1";cropbottom "0";tempfilename
'../../../../PH375/repository/OpticsText/text/SVAL9A3G.wmf';tempfile-properties "XPR";}%
}and five.\FRAME{dtbpF}{3.7645in}{2.1292in}{0pt}{}{}{Figure}{\special%
{language "Scientific Word";type "GRAPHIC";maintain-aspect-ratio
TRUE;display "USEDEF";valid_file "T";width 3.7645in;height 2.1292in;depth
0pt;original-width 3.717in;original-height 2.0903in;cropleft "0";croptop
"1";cropright "1";cropbottom "0";tempfilename
'../../../../PH375/repository/OpticsText/text/SVAL9A3H.wmf';tempfile-properties "XPR";}%
}By looking for where the wave crests line up and drawing Huygen's line we
can find the new crest of the reflected wave. By superimposing these
scattered waves we the reflected wave is formed.

We also know \emph{Snell's law of refraction}. But let's see how Huygen's
principle leads can be used to show this is true. In our discussion on
reflection we just used the top layer of scatters, but we know that there
are more scatterers in a solid. \FRAME{dtbpF}{3.5509in}{2.0064in}{0pt}{}{}{%
Figure}{\special{language "Scientific Word";type
"GRAPHIC";maintain-aspect-ratio TRUE;display "USEDEF";valid_file "T";width
3.5509in;height 2.0064in;depth 0pt;original-width 4.7063in;original-height
2.6463in;cropleft "0";croptop "1";cropright "1";cropbottom "0";tempfilename
'../../../../PH375/repository/OpticsText/text/SVAL9A39.wmf';tempfile-properties "XPR";}%
}So, what happens inside the material with all the other layers of
scatterers? Of course our secondary waves are kind of spherical shaped, so
the situation is something like this\FRAME{dtbpF}{4.1286in}{2.335in}{0pt}{}{%
}{Figure}{\special{language "Scientific Word";type
"GRAPHIC";maintain-aspect-ratio TRUE;display "USEDEF";valid_file "T";width
4.1286in;height 2.335in;depth 0pt;original-width 4.0785in;original-height
2.2952in;cropleft "0";croptop "1";cropright "1";cropbottom "0";tempfilename
'../../../../PH375/repository/OpticsText/text/SVAL9A3A.wmf';tempfile-properties "XPR";}%
}and we see by drawing another Huygen's line that we are forming another
wave front \emph{inside the material}. Our scattering model predicts the
transmitted wave as well as the reflected wave. We can see that Snell's law
follows from basic geometry just as the law of reflection did. It depends on
the spacing of the atoms and the size of the wavelengths. This is a more
profound view of the laws of reflection and refraction than our ray model
approach.

Huygen's principle is used many times in the study of optics to help us
understand the geometry of an optics problem and from that form an equation
that makes calculations possible!

Snell and others used measurements to find Snell's law. But now that we
understand the geometry we can use it to find Snell's equation. In the next
figure you can see an incoming wave. We have marked a wavefront $\overline{AB%
}$ that is just starting to hit the interface between the two materials. We
can think of the \emph{wavefront} as a particular crest of the wave. And we
have marked another wavefront, $\overline{CD}$ that is the transmitted wave
front that has just become entirely inside the new material.

\FRAME{dtbpF}{5.2183in}{3.2404in}{0pt}{}{}{Figure}{\special{language
"Scientific Word";type "GRAPHIC";maintain-aspect-ratio TRUE;display
"USEDEF";valid_file "T";width 5.2183in;height 3.2404in;depth
0pt;original-width 5.1612in;original-height 3.1955in;cropleft "0";croptop
"1";cropright "1";cropbottom "0";tempfilename
'../../../../PH375/repository/OpticsText/text/SVAL9A38.wmf';tempfile-properties "XPR";}%
}The time it takes the wave front to go from $B$ to $D$ would be 
\begin{equation*}
\Delta t_{BD}=\frac{\overline{BD}}{v_{i}}
\end{equation*}%
but because the new transmitted wave is slower, the transmitted wavefront
has only gone the distance $\overline{AC}$ in the same time%
\begin{equation*}
\Delta t_{AC}=\Delta t_{BD}=\Delta t=\frac{\overline{AC}}{v_{t}}
\end{equation*}%
We have two triangles. And they share a side $\overline{AD}.$ And they are
right triangles. So the two triangles must be similar triangles. Now let's
look at our incident angle $\theta _{i}$ that of course we measure from a
normal line. We are going to want the $\sin \theta _{i}$ for Snell's law and
we can see that 
\begin{equation*}
\sin \theta _{i}=\frac{\overline{BD}}{\overline{AD}}
\end{equation*}%
and also 
\begin{equation*}
\sin \theta _{t}=\frac{\overline{AC}}{\overline{AD}}
\end{equation*}%
and we can solve both for $1/\overline{AD}$ 
\begin{eqnarray*}
\frac{\sin \theta _{i}}{\overline{BD}} &=&\frac{1}{\overline{AD}} \\
\frac{\sin \theta _{t}}{\overline{AC}} &=&\frac{1}{\overline{AD}}
\end{eqnarray*}%
and set them equal so that 
\begin{equation*}
\frac{\sin \theta _{i}}{\overline{BD}}=\frac{\sin \theta _{t}}{\overline{AC}}
\end{equation*}%
but 
\begin{eqnarray*}
\overline{BD} &=&v_{i}\Delta t \\
\overline{AC} &=&v_{t}\Delta t
\end{eqnarray*}%
so 
\begin{equation*}
\frac{\sin \theta _{i}}{v_{i}\Delta t}=\frac{\sin \theta _{t}}{v_{t}\Delta t}
\end{equation*}%
And the $\Delta t$ cancels from both sides%
\begin{equation*}
\frac{\sin \theta _{i}}{v_{i}}=\frac{\sin \theta _{t}}{v_{t}}
\end{equation*}%
We can multiply both sides by $c,$ the speed of light 
\begin{equation*}
c\frac{\sin \theta _{i}}{v_{i}}=c\frac{\sin \theta _{t}}{v_{t}}
\end{equation*}%
to achieve%
\begin{equation*}
n_{i}\sin \theta _{i}=n_{t}\sin \theta _{t}
\end{equation*}%
which is Snell's law, this time from basic geometry.

\section{Filters and other color phenomena}

Of course, we have assumed without statement, that white light is made up of
all the colors of the rainbow. We should ask why a red shirt is red, or why
passing light through a green film makes the light look green as it leaves.

Both of these phenomena are examples of removing wavelengths from white
light.

In the case of the red shirt, the red dye in the cloth absorbs all of the
visible colors except red. The red is reflected, so the shirt looks red. The
filter is much the same. The green pigment in the film causes nearly all
visible colors to be absorbed except green. So only green light is
transmitted. This is why leaves are green. Chlorophyll absorbs red and blue
wavelengths, so the green is reflected or transmitted.

\FRAME{dhFU}{2.4738in}{2.19in}{0pt}{\Qcb{{\protect\small Chlorophyll
Spectrum (Public Domain image courtesy Kurzon)}}}{}{Figure}{\special%
{language "Scientific Word";type "GRAPHIC";maintain-aspect-ratio
TRUE;display "USEDEF";valid_file "T";width 2.4738in;height 2.19in;depth
0pt;original-width 2.4336in;original-height 2.1508in;cropleft "0";croptop
"1";cropright "1";cropbottom "0";tempfilename
'SVAJW81L.wmf';tempfile-properties "XPR";}}

In our next lecture, we will continue to learn about light and we will think
about how light forms images.

\chapter{14 Polarization and Image formation 3.1.7, 3.2.1}

Light is a transverse wave. And it turns out this matters! The direction the
electric field is oscillating leads to what we call \emph{polarization} and
this is used in making things like polarized sun glasses. Let's talk about
this in this lecture.

%TCIMACRO{%
%\TeXButton{\chaptermark{Polarization and Images}}{\chaptermark{Polarization and Images}}}%
%BeginExpansion
\chaptermark{Polarization and Images}%
%EndExpansion

%TCIMACRO{%
%\TeXButton{Fundamental Concepts}{\vspace{0.10in}\hspace{-0.37in}{\Large Fundamental Concepts\vspace{0.2in}}}}%
%BeginExpansion
\vspace{0.10in}\hspace{-0.37in}{\Large Fundamental Concepts\vspace{0.2in}}%
%EndExpansion

\begin{itemize}
\item The direction of the electric field in a plane wave is called the
polarization direction.

\item Natural light is usually a superposition of many waves with random
polarization directions. This light is called unpolarized light.

\item Some materials allow light with one polarization to pass through,
while stopping other polarizations. The polaroid is one such material
polaroids. will have a final intensity that follows the relationship $%
I=I_{\max }\cos ^{2}\left( \theta \right) $

\item Light reflecting off a surface may be polarized because of the
absorption and re-emission pattern of light interacting with the material
atoms.

\item Scattered light may be polarized because of anisotropies in the
scatterers.

\item Birefringent materials have different wave speeds in different
directions. This affects the polarization of light entering these materials.

\item An image is a diverging set of rays in a recognizable pattern

\item Images can be formed by reflection and refraction
\end{itemize}

\section{Polarization of Light Waves}

We said much earlier in our study of light that it was a transverse wave and
that the medium for this wave is the electromagnetic field. Understanding
electric and magnetic fields is a topic for Principles of Physics III
(PH220) so we can't go to deep into electric and magnetic fields here. Let's
assume that all of space (the whole universe) is filled with this stuff we
call electric and magnetic fields. And let's suppose that these fields can
be a medium for waves. And let's identify these waves as light.

The light waves are transverse waves. So when we look at a light wave in
detail we will see the pieces of the electric and magnetic fields
experiencing simple harmonic motion much like the parts of our spring for
waves on springs. It turns out that the electric field medium parts move
perpendicular to the direction the wave goes, and also perpendicular to the
direction the magnetic field medium parts go. This makes them transverse
waves. We will now show some implications of this fact. In a course in
electromagnetic theory, we often draw light as in the figure below.\FRAME{dhF%
}{3.681in}{3.673in}{0pt}{}{}{Figure}{\special{language "Scientific
Word";type "GRAPHIC";maintain-aspect-ratio TRUE;display "USEDEF";valid_file
"T";width 3.681in;height 3.673in;depth 0pt;original-width
3.7218in;original-height 3.7156in;cropleft "0";croptop "1";cropright
"1";cropbottom "0";tempfilename 'SVAM8I46.wmf';tempfile-properties "XPR";}}%
and that looks a bit like our waves on strings or springs only with the
electric field part going up and down and the magnetic field part going back
and forth.

In what follows we will ignore the magnetic field (marked in the figure as $%
B $) because it's amplitude is very much smaller than the electric field
amplitude. So, we will look at the $E$ field an notice that it goes up and
down in the figure. But we could have light in any orientation. If we look
directly at an approaching beam of light we would \textquotedblleft
see\textquotedblright\ many different orientations as shown in the next
figure.\FRAME{dtbpF}{1.5307in}{1.5895in}{0pt}{}{}{Figure}{\special{language
"Scientific Word";type "GRAPHIC";maintain-aspect-ratio TRUE;display
"USEDEF";valid_file "T";width 1.5307in;height 1.5895in;depth
0pt;original-width 1.4944in;original-height 1.5541in;cropleft "0";croptop
"1";cropright "1";cropbottom "0";tempfilename
'SVAM8I47.wmf';tempfile-properties "XPR";}}

When light beams have waves with many orientations, we say they are \emph{%
unpolarized}. But suppose we were able to align all the light so that all
the waves in the beam were transverse waves in the same orientation. Say,
the one in the next figure.\FRAME{dhF}{0.3226in}{1.3508in}{0pt}{}{}{Figure}{%
\special{language "Scientific Word";type "GRAPHIC";maintain-aspect-ratio
TRUE;display "USEDEF";valid_file "T";width 0.3226in;height 1.3508in;depth
0pt;original-width 1.4688in;original-height 6.2923in;cropleft "0";croptop
"1";cropright "1";cropbottom "0";tempfilename
'SVAM8I48.wmf';tempfile-properties "XPR";}}

Then we would describe the light as \emph{linearly polarized}. The plane
that contains the $E$-field is known as the \emph{polarization plane}.

\subsection{Polarization by removing all but one wave orientation}

One way to make polarized light is to remove all but one orientation of an
unpolarized beam. A material that does this at visible wavelengths is called
a \emph{polaroid}. It is made of long-chain hydrocarbons that have been
treated with iodine to make them conductive. The molecules are all oriented
in one direction by stretching during the manufacturing process. The
molecules have electrons that can move when light hits them.

Remember that the reason light slows down when it enters a material is
because the light get's absorbed by the electrons. But in most materials the
electrons have to stay in the atoms. They just go to a higher shell. But in
polaroid material the electrons can move farther in the long direction of
the molecule, And the moving electrons will make light! Along the molecules
act like they are in little radio antennas. The molecules' electrons are
driven into harmonic motion along the length of the molecule. But this means
the energy that used to be our light wave gets turned into kinetic energy of
the electrons. This takes energy (and therefore, light) out of the beam. So
the light that is polarized in this direction stops at the polaroid

But suppose the light is polarized in a direction perpendicular to the long
molecules. Then little electron motion is possible in the short direction of
the molecule. And if the electrons can't move this direction they won't take
up all the light energy. Then the light is passed if it is polarized
perpendicular to the long direction of the molecules. This direction is
called the \emph{transmission axis}.

We can take two pieces of polaroid material to study polarization.\FRAME{dhF%
}{3.7879in}{1.4088in}{0in}{}{}{Figure}{\special{language "Scientific
Word";type "GRAPHIC";maintain-aspect-ratio TRUE;display "USEDEF";valid_file
"T";width 3.7879in;height 1.4088in;depth 0in;original-width
3.7395in;original-height 1.3742in;cropleft "0";croptop "1";cropright
"1";cropbottom "0";tempfilename 'SVAM8I49.wmf';tempfile-properties "XPR";}}

Unpolarized light is initially polarized by the first piece of polaroid
called the \emph{polarizer}. The second piece of polaroid then receives the
light. This piece is called the \emph{analyzer}. If there is an angular
difference in the orientation of the transmission axes of the polarizer and
analyzer, there will be a reduction of light through the system. We expect
that if the transmission axes are separated by $90\unit{%
%TCIMACRO{\U{b0}}%
%BeginExpansion
{{}^\circ}%
%EndExpansion
}$ no light will be seen. If they are separated by $0\unit{%
%TCIMACRO{\U{b0}}%
%BeginExpansion
{{}^\circ}%
%EndExpansion
},$ then there will be a maximum. It is not hard to believe that the
intensity will be given by%
\begin{equation}
I=I_{\max }\cos ^{2}\left( \theta \right)
\end{equation}%
remembering that for waves we have found that the intensity is proportional
to the amplitude squared. This is true for light waves as well (more on this
later, and in Principles of Physics III).

\subsection{Polarization by reflection}

If we look at light reflected off of a desk or table through a piece of
polaroid, we can see that at some angles of orientation, the reflection
diminishes or even disappears! Light is often polarized on reflection. Let's
consider a beam of light made of just two polarizations. We will define a
plane of incidence. This plane is the plane of the paper or computer screen.
This plane is perpendicular to the reflective or refractive surface in the
figure below.\FRAME{dtbpF}{3.3313in}{2.815in}{0pt}{}{}{Figure}{\special%
{language "Scientific Word";type "GRAPHIC";maintain-aspect-ratio
TRUE;display "USEDEF";valid_file "T";width 3.3313in;height 2.815in;depth
0pt;original-width 3.5544in;original-height 3in;cropleft "0";croptop
"1";cropright "1";cropbottom "0";tempfilename
'SVAM8I4A.wmf';tempfile-properties "XPR";}}

One of our polarizations is defined as parallel to this plane. This
direction is represented by orange (lighter grey in black and white) arrows
in the figure. The other polarization is perpendicular to the plane of
incidence (the plane of the paper or your device screen). This is
represented by the black dots in the figure. These dots are supposed to look
like arrows coming out of the paper.

When the light reaches the interface between $n_{1}$ and $n_{2}$ it drives
the electrons in the medium into SHM. The perpendicular polarization finds
electrons that are free to move in the perpendicular direction and
re-radiate in that direction. The electron orbitals change shape and
oscillate with the incoming electromagnetic wave.

The parallel ray is also able to excite SHM, but an electromagnetic analysis
tells us that these little \textquotedblleft antennas\textquotedblright\
will not radiate at an angle $90\unit{%
%TCIMACRO{\U{b0}}%
%BeginExpansion
{{}^\circ}%
%EndExpansion
}$ from their excitation direction. And this is something new! Before when
we used Huygens' principle we thought of the scattered waves as little
spheres. But it turns out that the light doesn't go out in all directions
equally. The disturbance for our wave is a moving charged particle like an
electron. If an electron in a wire experiences simple harmonic motion there
will be a wave in the electric field. But think, that wave will go out away
from the electron perpendicular to the electron's motion \FRAME{dtbpF}{%
3.1816in}{1.7279in}{0pt}{}{}{Figure}{\special{language "Scientific
Word";type "GRAPHIC";maintain-aspect-ratio TRUE;display "USEDEF";valid_file
"T";width 3.1816in;height 1.7279in;depth 0pt;original-width
3.2144in;original-height 1.7332in;cropleft "0";croptop "1";cropright
"1";cropbottom "0";tempfilename 'SVJU7K04.wmf';tempfile-properties "XPR";}}%
In the figure we can see the light wave won't travel up or down. The actual
situation with our light striking a surface is a little more complicated,
The pattern of light coming out from the little atomic oscillators looks
more like a doughnut.

\FRAME{dtbpFU}{2.9616in}{1.9842in}{0pt}{\Qcb{Angular dependence of $S$ for a
dipole scatterer.}}{}{Figure}{\special{language "Scientific Word";type
"GRAPHIC";maintain-aspect-ratio TRUE;display "USEDEF";valid_file "T";width
2.9616in;height 1.9842in;depth 0pt;original-width 2.99in;original-height
1.993in;cropleft "0";croptop "1";cropright "1";cropbottom "0";tempfilename
'SVJTQF02.wmf';tempfile-properties "XPR";}}But the point is that it's not
completely spherical and none of the light wave goes up along the direction
of the oscillation.

In the figure you can see that along the antenna axis, the field amplitude
is zero. This means that the wave really does not go that direction. So in
our case, the amount of polarization in the parallel direction decreases
with the angle between the reflected and refracted rays until at $90\unit{%
%TCIMACRO{\U{b0}}%
%BeginExpansion
{{}^\circ}%
%EndExpansion
}$ there is no reflected ray in the parallel direction. \FRAME{dtbpF}{%
4.1035in}{3.4082in}{0pt}{}{}{Figure}{\special{language "Scientific
Word";type "GRAPHIC";maintain-aspect-ratio TRUE;display "USEDEF";valid_file
"T";width 4.1035in;height 3.4082in;depth 0pt;original-width
4.0542in;original-height 3.3624in;cropleft "0";croptop "1";cropright
"1";cropbottom "0";tempfilename 'SVAM8I3O.wmf';tempfile-properties "XPR";}}

The incidence angle that creates an angular difference between the refracted
and reflected rays of $90\unit{%
%TCIMACRO{\U{b0}}%
%BeginExpansion
{{}^\circ}%
%EndExpansion
}$ is called the \emph{Brewster's\ angle} after its discoverer. At this
angle the reflected beam will be completely linearly polarized.

We can predict this angle. Remember Snell's law.%
\begin{equation*}
n_{1}\sin \theta _{1}=n_{2}\sin \theta _{2}
\end{equation*}%
Let's re-lable the incidence angle $\theta _{1}=$ $\theta _{b}.$ We take $%
n_{1}=1$ and $n_{2}=n$ so%
\begin{equation*}
n=\frac{\sin \theta _{b}}{\sin \theta _{2}}
\end{equation*}%
Now notice that for Brewster's angle, we have 
\begin{equation*}
\theta _{b}+90\unit{%
%TCIMACRO{\U{b0}}%
%BeginExpansion
{{}^\circ}%
%EndExpansion
}+\theta _{2}=180\unit{%
%TCIMACRO{\U{b0}}%
%BeginExpansion
{{}^\circ}%
%EndExpansion
}
\end{equation*}%
so%
\begin{equation*}
\theta _{2}=90\unit{%
%TCIMACRO{\U{b0}}%
%BeginExpansion
{{}^\circ}%
%EndExpansion
}-\theta _{b}
\end{equation*}%
so we have%
\begin{equation*}
n=\frac{\sin \theta _{b}}{\sin \left( 90\unit{%
%TCIMACRO{\U{b0}}%
%BeginExpansion
{{}^\circ}%
%EndExpansion
}-\theta _{b}\right) }
\end{equation*}%
ah, but we remember that $\sin \left( 90\unit{%
%TCIMACRO{\U{b0}}%
%BeginExpansion
{{}^\circ}%
%EndExpansion
}-\theta \right) =\cos \left( \theta \right) $ so%
\begin{equation*}
n=\frac{\sin \theta _{b}}{\cos \theta _{b}}
\end{equation*}%
Again we remember that 
\begin{equation*}
\tan \theta =\frac{\sin \theta }{\cos \theta }
\end{equation*}%
so%
\begin{equation}
n=\tan \theta _{b}
\end{equation}%
which we can solve for $\theta _{b}.$%
\begin{equation*}
\theta _{b}=\tan ^{-1}\left( n\right)
\end{equation*}

This phenomena is why we wear polarizing sunglasses to reduce glare.

\subsection{Birefringence}

Glass is an amorphic solid--that is--it has no crystal structure to speak
of. But some minerals do have definite order. Sometimes the difference in
the crystal structure creates a difference in the speed of propagation of
light in the crystal. This is not to hard to believe. We said before that
the reason light slows down in a substance is because it encounters atoms
which absorb and re-emit the light. If there are more atoms in one direction
than another in a crystal, it makes sense that there could be a different
speed in each direction.

Calcite crystals exhibit this phenomena. We can describe what happens by
defining two polarizations. One parallel to the plane of the figure below,
and one perpendicular.

\bigskip

\FRAME{dhF}{3.1337in}{2.0117in}{0in}{}{}{Figure}{\special{language
"Scientific Word";type "GRAPHIC";maintain-aspect-ratio TRUE;display
"USEDEF";valid_file "T";width 3.1337in;height 2.0117in;depth
0in;original-width 3.1647in;original-height 2.0214in;cropleft "0";croptop
"1";cropright "1";cropbottom "0";tempfilename
'SVAM8I4B.wmf';tempfile-properties "XPR";}}With a careful setup, we can
arrange things so the perpendicular ray is propagated just as we would
expect for glass. We call this the $O$-ray (for \emph{ordinary}). The second
ray is polarized parallel to the incidence plane. It will have a different
speed, and therefore a different index of refraction. We call it the \emph{%
Extraordinary ray} or $E$-ray.

\FRAME{dhF}{2.1137in}{2.1234in}{0in}{}{}{Figure}{\special{language
"Scientific Word";type "GRAPHIC";maintain-aspect-ratio TRUE;display
"USEDEF";valid_file "T";width 2.1137in;height 2.1234in;depth
0in;original-width 2.1252in;original-height 2.135in;cropleft "0";croptop
"1";cropright "1";cropbottom "0";tempfilename
'SVAM8I4C.wmf';tempfile-properties "XPR";}}

If we were to put a light source in a calcite crystal, we would see the $O$%
-ray send out a sphere of light as shown in the figure above. But the $E$%
-ray would send out an ellipse. The speed for the $E$-ray depends on
orientation. There is one direction where the speeds are equal. This
direction is called the \emph{optic axis} of the crystal.\FRAME{dtbpF}{%
3.2206in}{1.8299in}{0in}{}{}{Figure}{\special{language "Scientific
Word";type "GRAPHIC";maintain-aspect-ratio TRUE;display "USEDEF";valid_file
"T";width 3.2206in;height 1.8299in;depth 0in;original-width
3.1756in;original-height 1.7928in;cropleft "0";croptop "1";cropright
"1";cropbottom "0";tempfilename 'SWZRZT01.wmf';tempfile-properties "XPR";}}%
If our light entering our calcite crystal is unpolarized, then we will have
two images leaving the other side that are slightly offset because the $O$%
-rays and $E$-rays both form images.

Some materials (notably plastics) become birefringent under stress. A
plastic or other stress birefringent material is molded in the form planned
for a building or other object (usually made to scale). The model is placed
under a stress, and the system is placed between to polaroids. When
unstressed, no light is seen, but under stress, the model changes the
polarization state of the light, and bands of light are seen.\FRAME{dhF}{%
3.1479in}{2.0676in}{0pt}{}{}{Figure}{\special{language "Scientific
Word";type "GRAPHIC";maintain-aspect-ratio TRUE;display "USEDEF";valid_file
"T";width 3.1479in;height 2.0676in;depth 0pt;original-width
3.1798in;original-height 2.0782in;cropleft "0";croptop "1";cropright
"1";cropbottom "0";tempfilename 'SVAM8I4E.wmf';tempfile-properties "XPR";}}

\subsection{Polarization due to scattering}

It is important to understand that light is also polarized by scattering. It
really takes a bit of electromagnetic theory to describe this. So for a
moment, lets just comment that blue light is scattered more than red light
because the blue wavelength is about the same size as the molecules. In
fact, the relative intensity of scattered light goes like $1/\lambda ^{4}.$
This has nothing to do with polarization, but it is nice to know.

Now suppose we have long pieces of wire in the air, say, a few microns long.
The pieces of wire would have electrons that could be driven into SHM when
light hits them. If the wires were all oriented in a common direction, we
would expect light to be absorbed if it was polarized in the long direction
of the particles and not absorbed in a direction perpendicular to the
orientation of the particles. This is exactly what happens when long ice
particles in the atmosphere orient in the wind (think of the moment of
inertia). We often get impressive halo's around the sun due to scattering
from ice particles. And these halos can be strongly polarized.

Rain drops also have a preferential scattering direction because they are
shaped like oblate spheroids (not \textquotedblleft rain drop
shape\textquotedblright\ like we were told in grade school). This makes
weather radar signals polarized and this lets us measure how much water
there is in a rain storm!

It is also true that small molecules will act like tiny antennas and will
scatter light preferentially in some directions and not in others. This is
called \emph{Rayleigh scattering} and is very like small dipole antennas.

\subsection{Optical Activity}

Some substances will rotate the polarization of a beam of light. This is
called being \emph{optically active}. The polarization state of the light
exiting the material depends on the length of the path through the material.
Your calculator display works this way. An electric field changes the
optical activity of the liquid crystal. There are polarizers over the liquid
crystal, so sometimes light passes through the display and sometimes it is
black.

\subsection{Laser polarization}

One last comment. Lasers are usually polarized. This is because the laser
light is generated in a \emph{cavity} created by two mirrors. The mirror is
tipped so light approaches it at the Brewster angle. Light with the right
polarization (parallel to the plane of the drawing) is reflected back nearly
completely, but light with the opposite polarization is not reflected at
all. This reduces the usual loss in reflection from a mirror, because in one
polarization the light must be reflected completely.

\FRAME{dtbpF}{4.2653in}{1.03in}{0pt}{}{}{Figure}{\special{language
"Scientific Word";type "GRAPHIC";maintain-aspect-ratio TRUE;display
"USEDEF";valid_file "T";width 4.2653in;height 1.03in;depth
0pt;original-width 4.2151in;original-height 0.9963in;cropleft "0";croptop
"1";cropright "1";cropbottom "0";tempfilename
'../../../../PH223/Repository/PH223-Engineering-Physics/Text/SVAM8I3P.wmf';tempfile-properties "XPR";}%
}

\section{Imaging with Flat Mirrors}

All of us have looked in a mirror at some time. We know what to expect. We
see an image of ourselves. And we just thought about fancy mirrors that we
use to make lasers. But we need to apply math to our study of mirrors so we
can calculate how light is reflected by mirrors.

To study mirrors mathematically, we need to establish a sign convention and
some standard notation.

\FRAME{dtbpF}{1.5463in}{1.6267in}{0pt}{}{}{Figure}{\special{language
"Scientific Word";type "GRAPHIC";maintain-aspect-ratio TRUE;display
"USEDEF";valid_file "T";width 1.5463in;height 1.6267in;depth
0pt;original-width 1.5108in;original-height 1.5904in;cropleft "0";croptop
"1";cropright "1";cropbottom "0";tempfilename
'SWDEXI00.wmf';tempfile-properties "XPR";}}In the figure above, we have a
person observing an object $O$ in a mirror. The object is located at a
distance $s$ from the mirror. Just like for lenses, we will call this the 
\emph{object distance, }$d_{o}$. The image appears to be located at a point $%
I$ beyond the mirror a distance $d_{i}.$We will call this the \emph{image
distance}. Be careful, $d_{i}$ looks like an incident or initial distance,
but it's not. It is the distance from the mirror to the \emph{image.} We
know we see an image in the mirror, but to understand this we need to define
what we mean by the word \textquotedblleft image.\textquotedblright

\subsubsection{Images}

Let's think about what an image is. Take a piece of paper and a lens, and
hold up the lens in a darkened room that has some bright object in it. Move
the lens or the paper back and forth, and at just the right distance, a
miniature image of the bright object will appear. We should think about what
the word \textquotedblleft image\textquotedblright\ means in this sense.

\FRAME{dhF}{2.6835in}{2.0237in}{0pt}{}{}{Figure}{\special{language
"Scientific Word";type "GRAPHIC";maintain-aspect-ratio TRUE;display
"USEDEF";valid_file "T";width 2.6835in;height 2.0237in;depth
0pt;original-width 2.6411in;original-height 1.9847in;cropleft "0";croptop
"1";cropright "1";cropbottom "0";tempfilename
'SWDFTN01.wmf';tempfile-properties "XPR";}}We have talked about how we see
objects. Remember the BYU-I guys from a previous lecture.\FRAME{dhF}{3.173in%
}{1.3534in}{0pt}{}{}{Figure}{\special{language "Scientific Word";type
"GRAPHIC";maintain-aspect-ratio TRUE;display "USEDEF";valid_file "T";width
3.173in;height 1.3534in;depth 0pt;original-width 3.128in;original-height
1.3188in;cropleft "0";croptop "1";cropright "1";cropbottom "0";tempfilename
'SWDFTN02.wmf';tempfile-properties "XPR";}}Our eyes gather rays that are
diverging from the object because light has bounced off of the object. Our
eyes intersect a diverging set of rays that form a definite pattern. That
diverging set of rays forming a pattern is the picture that we see of the
object.

So when we say that the lens has formed a miniature image of our object, we
mean that the lens has somehow formed a diverging set of rays that form a
pattern that looks like the pattern formed by the diverging set of rays
coming from the object, itself. In other words, the object forms a diverging
set of rays. And our lens forms a duplicate set of rays in the same pattern,
so we see the same thing. The lens' version is smaller, upside down, and
backwards, but it is still essentially the same pattern.

But how can a mirror form a diverging of rays that form a pattern?

Well, the light that we see in a mirror first bounces off of something, say,
a person. This makes a recognizable pattern in the light. Then the light
reflects off of the mirror. Then the light hits our eyes. So when we see
something in a mirror we do see a diverging set of rays that form a pattern.
But what we see makes more makes sense if we know how our brains processes
the light signals that we see. Our eyes intercept rays of light diverging
from an object. Our brain processes those rays as though they traveled only
in straight lines. Our brain really believes light travels just in lines, no
bouncing or reflecting. This means, if we can create a situation that makes
rays diverge in the same way the object did, we will have an image of the
object. But our brain will be wrong in where it thinks that image is. We
will see an image as though it were behind the mirror.

\subsection{Virtual Images}

We already know that mirrors often create what we call \emph{virtual} images
because the image appears to be created from diverging rays from behind the
mirror. If we look behind the mirror, no rays exist there (if they did
exist, they would not make it through the mirror!).

\FRAME{dtbpF}{2.3869in}{1.4183in}{0in}{}{}{Figure}{\special{language
"Scientific Word";type "GRAPHIC";maintain-aspect-ratio TRUE;display
"USEDEF";valid_file "T";width 2.3869in;height 1.4183in;depth
0in;original-width 2.3471in;original-height 1.3837in;cropleft "0";croptop
"1";cropright "1";cropbottom "0";tempfilename
'SWDFVW03.wmf';tempfile-properties "XPR";}}

Let's look at a simple image as shown in the figure above. The object is an
arrow. We could trace all the rays that diverge from this object and build a
very nice representation of the arrow (like ray tracing-based computer
graphics--the way movies like \emph{Toy Story} are made) but that would take
time and computation power. We only really need to use two rays, and to
remember what the object looked like to figure out where the image will be.

We pick one ray from the top of the arrow that travels straight to the
mirror. This ray will travel a distance $d_{o}$ and bounce back. Using the
law of reflection we know that it will bounce straight back along the path
it came in on. We pick a second ray from point $P$ that travels the path $%
PR. $ This ray bounces off the mirror at an angle $\theta _{r}=\theta _{i}.$
But our brain doesn't believe the light reflected. It thinks the light
traveled straight paths. It mentally traces the rays backwards along a
straight path as though it came from that direction. And this makes it
appear that there is a tip of an arrow is at position $P^{\prime }.$ The
rays from the virtual tip appear to travel the paths $P^{\prime }P$ and $%
P^{\prime }R.$ The light doesn't really come from point $P^{\prime }.$ But
it \emph{appears} to come from point $P^{\prime }$. We could use similar
sets of two rays to find where the light from the middle of the arrow \emph{%
appears} to come from, and for the bottom of the arrow, and all the arrow
parts. The collection of all these apparent locations is image location and
the light that appears to come from this location is the image. But, and
this is important! light doesn't come from $P^{\prime }.$ It only appears to
come from $P^{\prime }.$ When an image is formed, but the light doesn't
really come from the image location, we call that a \emph{virtual image.}

\subsection{Mirror reversal}

Look into a mirror. Raise your right hand. \FRAME{dtbpF}{3.173in}{1.9147in}{%
0pt}{}{}{Figure}{\special{language "Scientific Word";type
"GRAPHIC";maintain-aspect-ratio TRUE;display "USEDEF";valid_file "T";width
3.173in;height 1.9147in;depth 0pt;original-width 3.1289in;original-height
1.8766in;cropleft "0";croptop "1";cropright "1";cropbottom "0";tempfilename
'SWDFXU04.wmf';tempfile-properties "XPR";}}Your image raises what appears to
be a left hand. It looks like a mirror switches the left and right sides of
the image. But it might be better to note that the image raised hand is on
the same side of the room as your real raised hand. Think of lying sideways
on the ground in front of the mirror and raise a hand. Note that your image
feet are on the same side of the room as your real feet. And note that as
you put one hand up, the image hand also goes up. This is not exactly
left-right reversal. What is happening? \FRAME{dtbpF}{4.3716in}{2.6333in}{0pt%
}{}{}{Figure}{\special{language "Scientific Word";type
"GRAPHIC";maintain-aspect-ratio TRUE;display "USEDEF";valid_file "T";width
4.3716in;height 2.6333in;depth 0pt;original-width 4.3206in;original-height
2.5918in;cropleft "0";croptop "1";cropright "1";cropbottom "0";tempfilename
'SWDFYO06.wmf';tempfile-properties "XPR";}}The light from your raised hand
would hit the mirror, reflect, and then be absorbed by your eye. Your brain
tells you the light traveled in a straight line. So it looks like you have a
raised hand in the mirror. But what has happened is that the light has
traveled from your hand, to the mirror and back to your eye. This is not
really left-right reversal. We need a name to describe what is really
happening. The name is odd. To see why we use this name, notice that from
your perspective (moving left to right across the figure) you first have the
back part of your hand, then the front of your hand and then the front of
the image hand, then the back of the image hand. We call this back-front
reversal. It comes from the reversal in the order of parts of the object,
like the front and back of your hand. So we say that a flat mirror performs
a front-back reversal, not a left-right reversal.

It is nice to know about how images are formed by mirrors. We use them all
the time. But you probably own a phone, and your phone probably has a
camera. Cameras also make images. We should work our way up to understanding
how things like cameras work. We will start that in our next lecture.

\chapter{15 Images formed by Curved Mirrors and by Refraction 3.2.2, 3.2.3}

%TCIMACRO{%
%\TeXButton{\chaptermark{Curved Mirrors}}{\chaptermark{Curved Mirrors}}}%
%BeginExpansion
\chaptermark{Curved Mirrors}%
%EndExpansion

Back in Newton's day, there were no electronic computers, or calculators, or
mechanical adding machines. Early optics researchers did math with a pen and
paper. This is one reason they liked small angle approximations. The
approximation allowed them to do harder problems using easy math. And so
long as the things they built worked, the approximations were good enough.
We are going to use another approximation in this lecture. It is called the
thin mirror approximation. It will make the math that describes mirrors for,
say, telescopes, much easier. We will also introduce a way to draw curved
mirrors and light that will tell you how the mirror works. We will use this
drawing approach to do homework and examination problems. Professional
optical designers use computer codes that do this drawing very accurately to
help the optical engineer understand the system they are designing and to
look for mistakes. So this drawing scheme is very useful. Because it makes
optical systems so much easier to understand, let's start with the drawings
and then work toward a mathematical description of thin mirrors.

%TCIMACRO{%
%\TeXButton{Fundamental Concepts}{\vspace{0.10in}\hspace{-0.37in}{\Large Fundamental Concepts\vspace{0.2in}}}}%
%BeginExpansion
\vspace{0.10in}\hspace{-0.37in}{\Large Fundamental Concepts\vspace{0.2in}}%
%EndExpansion

\begin{itemize}
\item Curved Mirrors can form images

\item We can write the magnification in terms of the object and image
distances $m=-\frac{d_{i}}{d_{o}}$

\item For \textquotedblleft thin\textquotedblright\ mirrors $\frac{1}{d_{o}}+%
\frac{1}{d_{i}}=\frac{1}{f}$

\item The focal length of a \textquotedblleft thin\textquotedblright\ mirror
is. $f=\frac{R}{2}$

\item We can describe how a mirror operates with just three easy-to-draw
rays.

\item A semi-infinite bump (single sided lens) can be described by the
equation $\frac{1}{d_{o}}+\frac{1}{d_{i}}=\frac{\left( n_{2}-n_{1}\right) }{R%
}$
\end{itemize}

\subsection{Concave Mirrors}

Concave mirrors can form images. I'm sure you know that many telescopes are
made with mirrors. We should see how this works. We recall the law of
reflection\FRAME{dhF}{1.8585in}{1.1857in}{0pt}{}{}{Figure}{\special{language
"Scientific Word";type "GRAPHIC";maintain-aspect-ratio TRUE;display
"USEDEF";valid_file "T";width 1.8585in;height 1.1857in;depth
0pt;original-width 1.8213in;original-height 1.1519in;cropleft "0";croptop
"1";cropright "1";cropbottom "0";tempfilename
'SVANNN4Z.wmf';tempfile-properties "XPR";}}

\begin{equation*}
\theta _{i}=\theta _{r}
\end{equation*}

Armed with this, we can see what would happen if we curved the mirror
surface. Each ray has a different normal line due to the curvature of the
mirror. \FRAME{dtbpF}{0.8458in}{1.9908in}{0pt}{}{}{Figure}{\special{language
"Scientific Word";type "GRAPHIC";maintain-aspect-ratio TRUE;display
"USEDEF";valid_file "T";width 0.8458in;height 1.9908in;depth
0pt;original-width 0.8138in;original-height 1.9519in;cropleft "0";croptop
"1";cropright "1";cropbottom "0";tempfilename
'SWMB6J00.wmf';tempfile-properties "XPR";}}And the law of reflection must
works at each one of these normal lines, so we get rays that reflect like
this. \FRAME{dtbpF}{3.0649in}{1.9666in}{0pt}{}{}{Figure}{\special{language
"Scientific Word";type "GRAPHIC";maintain-aspect-ratio TRUE;display
"USEDEF";valid_file "T";width 3.0649in;height 1.9666in;depth
0pt;original-width 2.5806in;original-height 1.6466in;cropleft "0";croptop
"1";cropright "1";cropbottom "0";tempfilename
'SWMB8101.wmf';tempfile-properties "XPR";}}The result is that parallel rays
all meet at a spot on the axis.

\FRAME{dhF}{3.243in}{2.079in}{0pt}{}{}{Figure}{\special{language "Scientific
Word";type "GRAPHIC";maintain-aspect-ratio TRUE;display "USEDEF";valid_file
"T";width 3.243in;height 2.079in;depth 0pt;original-width
3.1981in;original-height 2.0401in;cropleft "0";croptop "1";cropright
"1";cropbottom "0";tempfilename 'SVANNN50.wmf';tempfile-properties "XPR";}}%
If we place something at this location (like a pice of paper), we would see
it light up. With enough light we could start a fire! We have just
re-invented the solar cooker. Look at the point where all the rays meet.
Let's give this special place a name. We call it the focal point. It is the
place where the rays come together if they originally came into the mirror
parallel.

\subsection{Paraxial Approximation for Mirrors}

The correct shape of a mirror to make all the parallel rays meet at the
focal point is more like a parabola, but parabolas are hard to machine or
build. Spherical shapes are relatively easy to make. So we often see
spherical mirrors used where we should have parabolas. This will work so
long as we allow only rays that strike the mirror not too far from the
principal axis. We can see why this works if we plot a sphere and a parabola
(and a hyperbola). For small deviations from the center, the shape of the
functions all look alike. \FRAME{dhF}{4.3613in}{1.1628in}{0pt}{}{}{Figure}{%
\special{language "Scientific Word";type "GRAPHIC";display
"USEDEF";valid_file "T";width 4.3613in;height 1.1628in;depth
0pt;original-width 4.3102in;original-height 1.3604in;cropleft "0";croptop
"1";cropright "1";cropbottom "0";tempfilename
'SVANNN51.wmf';tempfile-properties "XPR";}}We would expect the reflections
to be similar under these circumstances. Using only the part of a mirror
where the spherical and parabolic shapes are essentially the same is called
the \emph{paraxial approximation}. So, if we meet the criteria for the
paraxial approximation, our spherical mirrors should work. Note that when
you need the entire mirror, say, in a communications antenna, you must do
better than a spherical approximation to the correct shape for your mirror.
Your satellite dish is likely not a spherical section. We studied flat
mirrors last lecture, and found they made images (well, virtual images).
Will curved mirrors make images? We know how to find out. We can trace two
rays from an object and see where the rays go. If the rays eventually
diverge after forming a recognizable pattern, then we have an image!\FRAME{%
dtbpF}{3.4433in}{2.1367in}{0pt}{}{}{Figure}{\special{language "Scientific
Word";type "GRAPHIC";maintain-aspect-ratio TRUE;display "USEDEF";valid_file
"T";width 3.4433in;height 2.1367in;depth 0pt;original-width
3.4805in;original-height 2.1492in;cropleft "0";croptop "1";cropright
"1";cropbottom "0";tempfilename 'TA9JKZ00.wmf';tempfile-properties "XPR";}}

Like with our flat mirror, we will measure distances from the mirror surface
(from point $V$ in the figure). We can find the image location, $d_{i},$ by
taking two rays. We could use any of billions of rays. But let's try to pick
rays that are easy to draw. One convenient ray is the ray that passes
through the center of curvature, $C.$ This ray will strike the mirror
surface at right angles and bounce back along the same path. The incidence
angle will be zero, so the reflected angle must be zero by the law of
reflection. That is the yellow ray in the figure.

Another convenient ray is the ray from the tip of the object to point $V.$
This ray will bounce back with angle $\theta $. Right at point $V$ the
mirror surface is perpendicular to the optic axis. This makes the bounce of
the blue ray just like the a bounce from a flat mirror! That is easy to
draw. Where these two reflected rays cross and then diverge, we will find
the image of the tip of our tree. Knowing the shape of the tree and that the
bottom of the tree is on the axis, we can fill in the rest of the image.

We can calculate the magnification for this case. We use the gold triangle
to determine that 
\begin{equation*}
\tan \theta =\frac{h_{o}}{d_{o}}
\end{equation*}%
and the light blue triangle to determine that 
\begin{equation*}
\tan \theta =\frac{-h_{i}}{d_{i}}
\end{equation*}%
And you might object to the negative sigh. But we want to indicate that the
image is inverted by making it's sign negative. So because the image is
upside down $h_{i}$ is defined to be negative. To make our tangent equation
work, we need to insert another minus sign. This gives us $\tan \theta $
instead of $-\tan \theta $ when $h^{\prime }$ is inverted (when $h^{\prime
}, $ itself, is negative). When we decide to make the sine of a quantity
negative or positive just because we want it to be that way we say it is
negative (or positive) \emph{by convention.} And we have done this here. It
is just a convention. But we will use it. Our magnification would be%
\begin{equation*}
m\equiv \frac{\text{Image height}}{\text{Object height}}=\frac{h_{i}}{h_{o}}
\end{equation*}%
But we can write 
\begin{eqnarray*}
h_{o} &=&d_{o}\tan \theta \\
h_{i} &=&-d_{i}\tan \theta
\end{eqnarray*}%
so the magnification could be written as 
\begin{equation*}
m=\frac{h_{i}}{h_{o}}=\frac{-d_{i}\tan \theta }{d_{o}\tan \theta }=-\frac{%
d_{i}}{d_{o}}
\end{equation*}%
which gives us the magnification in terms of how far away the object and
image are.

\section{Mirror Equation}

We can further exploit this geometry to get a relationship between $d_{o},$ $%
d_{i},$ and $R.$ Notice that 
\begin{equation*}
\tan \alpha =\frac{h_{o}}{d_{o}-R}
\end{equation*}%
and that 
\begin{equation*}
\tan \alpha =\frac{-h_{i}}{R-d_{i}}
\end{equation*}

Then%
\begin{equation*}
\frac{h_{o}}{d_{o}-R}=\frac{-h_{i}}{R-d_{i}}
\end{equation*}%
or%
\begin{equation*}
\frac{R-d_{i}}{d_{o}-R}=-\frac{h_{i}}{h_{o}}
\end{equation*}

We can use our magnification definition to replace $h_{i}/h_{o}$%
\begin{equation*}
\frac{R-d_{i}}{d_{o}-R}=\frac{d_{i}}{d_{o}}
\end{equation*}%
And we have a useful equation, but not in the form you usually see it. Let's
perform some algebra. Starting with our last equation, let's try to make is
so there are no fractions.%
\begin{eqnarray*}
\left( R-d_{i}\right) d_{o} &=&d_{i}\left( d_{o}-R\right) \\
Rd_{o}-d_{i}d_{o} &=&d_{i}d_{o}-d_{i}R
\end{eqnarray*}%
And now let's get all the terms with an $R$ in them on the same side of the
equation. 
\begin{eqnarray*}
Rd_{o}+d_{i}R &=&d_{i}d_{o}+d_{i}d_{o} \\
d_{i}R+Rd_{o} &=&2d_{o}d_{i}
\end{eqnarray*}%
Now, for no apparent reason, I want to divide both sides of the equation by $%
Rd_{o}d_{i}.$ That might not seem reasonable, but let's do it anyway. 
\begin{equation*}
\frac{Rd_{i}}{Rd_{o}d_{i}}+\frac{Rd_{o}}{Rd_{o}d_{i}}=\frac{2d_{o}d_{i}}{%
Rd_{o}d_{i}}
\end{equation*}%
Some things cancel, and we are left with 
\begin{equation*}
\frac{1}{d_{o}}+\frac{1}{d_{i}}=\frac{2}{R}
\end{equation*}%
We are close to the usual form for this equation. But let's make one more
change to this equation. And it will involve the focal point.

\subsubsection{Focal Point}

Remember that special place where, if the light comes into the mirror
parallel, the reflected rays all meet? We called that the focal point, $f$.

Now that we know the mirror equation, let's let $d_{o}$ be very large (for
example, let $d_{o}$ be the distance to the Sun). Then the rays would come
very nearly parallel. And then after reflection all the rays would come
together at the focal point. They would diverge from the focal point. We
have a diverging set of rays that are shaped like the sun! We would get an
image of the Sun right at the focal point! Since the sun is approximately
infinitely far away, we could say $d_{o}\approx \infty .$ Then, using our
last equation, 
\begin{eqnarray*}
\frac{1}{\infty }+\frac{1}{d_{i}} &\approx &\frac{2}{R} \\
0+\frac{1}{d_{i}} &\approx &\frac{2}{R} \\
\frac{1}{d_{i}} &\approx &\frac{2}{R}
\end{eqnarray*}%
or 
\begin{equation*}
d_{i}\approx \frac{R}{2}
\end{equation*}

This is a special image point. But we really already know what to call this
special place where parallel rays come together after reflection. We call it
the \emph{focal point} and the distance from the mirror to the point $f\ $is
called the \emph{focal length}. We see that for a curved mirror the length $%
f $ is given by 
\begin{equation}
f=\frac{R}{2}
\end{equation}%
So we can write the mirror equation as%
\begin{equation}
\frac{1}{d_{o}}+\frac{1}{d_{i}}=\frac{1}{f}
\end{equation}

This equation is very useful. Even when you use powerful optical design
software it can form the beginning of your optical system design.

\subsection{Ray Diagrams for Mirrors}

For our mirror problems we will need to draw a diagram of the situation. And
often the secret to solving the problem is in making the diagram! So, let's
learn how to do this. We could draw an infinite number of rays to find where
the image will be. And, indeed, optical design software does draw many rays
to check our designs. But for us we need to draw just a few easy rays. But
what rays are easy to draw?

Let's think about what we already know. \FRAME{dtbpF}{2.9386in}{1.224in}{0pt%
}{}{}{Figure}{\special{language "Scientific Word";type
"GRAPHIC";maintain-aspect-ratio TRUE;display "USEDEF";valid_file "T";width
2.9386in;height 1.224in;depth 0pt;original-width 2.9661in;original-height
1.2187in;cropleft "0";croptop "1";cropright "1";cropbottom "0";tempfilename
'SWIXQ903.wmf';tempfile-properties "XPR";}}If parallel rays come into the
mirror, they meet at the focal point. They don't really stop there unless we
put something there.

\FRAME{dtbpF}{3.2907in}{1.3509in}{0pt}{}{}{Figure}{\special{language
"Scientific Word";type "GRAPHIC";maintain-aspect-ratio TRUE;display
"USEDEF";valid_file "T";width 3.2907in;height 1.3509in;depth
0pt;original-width 3.3244in;original-height 1.3473in;cropleft "0";croptop
"1";cropright "1";cropbottom "0";tempfilename
'SWIXTT04.wmf';tempfile-properties "XPR";}}But often we will just stop our
drawing at the point $f.$ This is easy to draw! Just a parallel line in and
bounce through the focal point, $f$. And we know the focal distance $f=R/2$
so it is even easy to find the point, $f.$ This is one way to make a solar
cooker. \FRAME{dtbpF}{2.866in}{2.1577in}{0pt}{}{}{Figure}{\special{language
"Scientific Word";type "GRAPHIC";maintain-aspect-ratio TRUE;display
"USEDEF";valid_file "T";width 2.866in;height 2.1577in;depth
0pt;original-width 1.8386in;original-height 1.3785in;cropleft "0";croptop
"1";cropright "1";cropbottom "0";tempfilename
'SolarCooker.wmf';tempfile-properties "XNPR";}}so we could call any ray that
comes in parallel and goes through the focal point a \textquotedblleft solar
cooker ray.\textquotedblright\ Note that it is a little weird to have a
point marked $f$ and a distance from a mirror to that point named $f.$ But
we will get used to it.

It turns out we can use this geometry twice. Think, the law of reflection
tells us that every ray that bounced off of the mirror did so with equal
angles $\theta _{i}=\theta _{r}.$ So what would happen if we put a light
source at the focal point? The law of reflection would work the same way,
and the rays would come out parallel! \FRAME{dtbpF}{3.9125in}{1.6356in}{0pt}{%
}{}{Figure}{\special{language "Scientific Word";type
"GRAPHIC";maintain-aspect-ratio TRUE;display "USEDEF";valid_file "T";width
3.9125in;height 1.6356in;depth 0pt;original-width 3.9595in;original-height
1.6383in;cropleft "0";croptop "1";cropright "1";cropbottom "0";tempfilename
'SWIXX405.wmf';tempfile-properties "XPR";}}This would be easy to draw! And
this is also useful. Old fashioned flashlights used this to make the light
from a bulb that came out in all directions go in just one direction. Big
versions of this were used to send light out parallel to find enemy aircraft
in World Ward II. These devices were called \emph{spot lights} so we can
call rays that act like this \textquotedblleft spot light\textquotedblright\
rays.

\FRAME{dtbpF}{2.4871in}{2.4002in}{0pt}{}{}{Figure}{\special{language
"Scientific Word";type "GRAPHIC";maintain-aspect-ratio TRUE;display
"USEDEF";valid_file "T";width 2.4871in;height 2.4002in;depth
0pt;original-width 2.5066in;original-height 2.417in;cropleft "0";croptop
"1";cropright "1";cropbottom "0";tempfilename
'SWIYD807.wmf';tempfile-properties "XPR";}}Now days, you might only see such
a device used in movies about your favorite super hero (Murci\'{e}lago Boy)
used to call on the super hero for help.

There is one more ray that is easy to draw. We have already encountered it.
It is the ray that goes through the center of curvature for the spherical
mirror so it hits with $\theta _{i}=0$ and bounces off with $\theta _{r}=0.$

Let's put all of these together to dray ray diagrams for a spherical mirror.

\subsubsection{Principal rays for a concave mirror:}

\begin{enumerate}
\item Ray 1 is drawn from the top of the object such that its reflected ray
must pass through $f$. \FRAME{dtbpF}{2.8729in}{1.3739in}{0pt}{}{}{Figure}{%
\special{language "Scientific Word";type "GRAPHIC";maintain-aspect-ratio
TRUE;display "USEDEF";valid_file "T";width 2.8729in;height 1.3739in;depth
0pt;original-width 2.8987in;original-height 1.3713in;cropleft "0";croptop
"1";cropright "1";cropbottom "0";tempfilename
'SWIYS008.wmf';tempfile-properties "XPR";}}This is the \textquotedblleft
solar cooker\textquotedblright\ ray.

\item Ray 2 is drawn from the top of the object through the center of
curvature. This ray will be incident on the mirror surface at a right angle
and will be reflected back on itself. \FRAME{dtbpF}{2.9457in}{1.4076in}{0pt}{%
}{}{Figure}{\special{language "Scientific Word";type
"GRAPHIC";maintain-aspect-ratio TRUE;display "USEDEF";valid_file "T";width
2.9457in;height 1.4076in;depth 0pt;original-width 2.9732in;original-height
1.405in;cropleft "0";croptop "1";cropright "1";cropbottom "0";tempfilename
'SWIYU00B.wmf';tempfile-properties "XPR";}}This is the \textquotedblleft
center of curvature\textquotedblright\ ray.

\item Ray 3 is drawn from the top of the object through the focal point to
reflect parallel to the principal axis.\FRAME{dtbpF}{2.9634in}{1.8103in}{0pt%
}{}{}{Figure}{\special{language "Scientific Word";type
"GRAPHIC";maintain-aspect-ratio TRUE;display "USEDEF";valid_file "T";width
2.9634in;height 1.8103in;depth 0pt;original-width 2.9918in;original-height
1.8157in;cropleft "0";croptop "1";cropright "1";cropbottom "0";tempfilename
'SWIYTE09.wmf';tempfile-properties "XPR";}}This is the \textquotedblleft
spot light\textquotedblright\ ray.
\end{enumerate}

We can put these together, and where they meet will be where the image will
form. \FRAME{dtbpF}{3.6765in}{2.2423in}{0in}{}{}{Figure}{\special{language
"Scientific Word";type "GRAPHIC";maintain-aspect-ratio TRUE;display
"USEDEF";valid_file "T";width 3.6765in;height 2.2423in;depth
0in;original-width 3.7182in;original-height 2.2565in;cropleft "0";croptop
"1";cropright "1";cropbottom "0";tempfilename
'SWIYYV0C.wmf';tempfile-properties "XPR";}} Of course, we did this for the
top of the tree. We could do the same for the middle of the tree and we
would get \FRAME{dtbpF}{3.7404in}{2.2866in}{0in}{}{}{Figure}{\special%
{language "Scientific Word";type "GRAPHIC";maintain-aspect-ratio
TRUE;display "USEDEF";valid_file "T";width 3.7404in;height 2.2866in;depth
0in;original-width 3.783in;original-height 2.3008in;cropleft "0";croptop
"1";cropright "1";cropbottom "0";tempfilename
'SWIZL90D.wmf';tempfile-properties "XPR";}}and we can do this for every
point on the tree and find the corresponding image point. But it turns out
that if we do the top of the object we pretty much know where the rest of
the image will be. So with just three rays we can find the image of our
object!

We can do the same for an object closer than a focal length\FRAME{dtbpF}{%
2.847in}{1.8542in}{0pt}{}{}{Figure}{\special{language "Scientific Word";type
"GRAPHIC";maintain-aspect-ratio TRUE;display "USEDEF";valid_file "T";width
2.847in;height 1.8542in;depth 0pt;original-width 3.3425in;original-height
2.1664in;cropleft "0";croptop "1";cropright "1";cropbottom "0";tempfilename
'SWIZRI0E.wmf';tempfile-properties "XPR";}}Only the spotlight ray doesn't
work because light that came from the top of the tree and passed through the
focal point won't hit the mirror. But we could draw a ray that goes along
the line from the focal point to the top of the tree and then hits the
mirror. That ray will hit the mirror and bounce just as though it really
came from the focal point. We usually use dotted lines to show when we are
pretending a ray exists somewhere where it really doesn't. So in the drawing
you see a dotted line come from the focal point, $f,$ and then the line
becomes solid blue where the light really starts at the top of the tree and
goes to the mirror and reflects. In this way we can still use our three ray
system to find the image.

But wait. In this case our three rays never meet! So there seems to be no
image! An it is true, there is no real image. But we learned from flat
mirrors that there can be a virtual image. Virtual images come from our
brain processing thinking that the rays only go in straight lines. So we can
extend the rays that leave the mirror backward as though they didn't bounce
at the mirror. Again we use dotted lines for this. And where the dotted
lines meet we have a virtual image!

\subsection{Principal rays for a convex mirror:}

We also may have a mirror that curves, but curves the other way. We can use
three rays, but this time we need to be careful.

\begin{enumerate}
\item Ray 1 is drawn from the top of the object such that its reflected ray
appears to have come from $f$.

\item Ray 2 is drawn from the top of the object so that it appears to have
come from the center of curvature. This ray will be incident on the mirror
surface at a right angle and will be reflected back on itself.

\item Ray 3 is drawn from the top of the object to reflect parallel to the
principal axis by aiming it at $f$.
\end{enumerate}

\FRAME{dtbpF}{2.7804in}{2.2009in}{0pt}{}{}{Figure}{\special{language
"Scientific Word";type "GRAPHIC";maintain-aspect-ratio TRUE;display
"USEDEF";valid_file "T";width 2.7804in;height 2.2009in;depth
0pt;original-width 2.2719in;original-height 1.7928in;cropleft "0";croptop
"1";cropright "1";cropbottom "0";tempfilename
'SWDHRF03.wmf';tempfile-properties "XPR";}}Note again that these rays won't
meet. So there is no real image. But we can extend the rays that leave the
mirror backward as though they hadn't bounced off the mirror we can see that
they do meet, and that is the location our brain assigns for the virtual
image.

We should tabulate our sign convention for mirrors. So far we know that if $%
h_{i}$ is upside down it will be negative. But we can assign signs for all
the other distances in our ray diagrams. 
\begin{equation*}
\begin{tabular}{|l|l|l|}
\hline
\textbf{Quantity} & \textbf{Positive if} & \textbf{Negative if} \\ \hline
{\small Object location }$\left( d_{o}\right) $ & 
\begin{tabular}{l}
{\small Object is in front of surface} \\ 
{\small (real object)}%
\end{tabular}
& 
\begin{tabular}{l}
{\small Object is in back of surface} \\ 
{\small (virtual object)}%
\end{tabular}%
{\small \ } \\ \hline
{\small Image location }$\left( d_{i}\right) $ & 
\begin{tabular}{l}
{\small Image is in front of surface } \\ 
{\small (real image)}%
\end{tabular}
& 
\begin{tabular}{l}
{\small Image is in back of surface} \\ 
{\small (virtual image)}%
\end{tabular}
\\ \hline
{\small Image height }$\left( h_{i}\right) $ & 
\begin{tabular}{l}
{\small Image is upright}%
\end{tabular}
& 
\begin{tabular}{l}
{\small Image is inverted}%
\end{tabular}
\\ \hline
{\small Radius }$\left( R_{1}\text{ and }R_{2}\right) $ & 
\begin{tabular}{l}
{\small Center of curvature is in} \\ 
{\small front of surface}%
\end{tabular}%
{\small \ } & 
\begin{tabular}{l}
{\small Center of curvature is in} \\ 
{\small back of surface}%
\end{tabular}%
{\small \ } \\ \hline
Focal length $\left( f\right) $ & In {\small front of surface} & In {\small %
back of surface} \\ \hline
\end{tabular}%
\end{equation*}

\section{Aberrations}

It's time to ask what happens when we don't use paraxial rays (when the
paraxial ray approximation is not valid) but our equations are based on
using only small angles or paraxial rays. You might guess that we will have
problems in our imagery. That sounds bad, so let's use a special word for
\textquotedblleft problems\textquotedblright\ in imagery. That word is \emph{%
aberration}.

Our first aberration comes from making the mirror the wrong shape. It should
be a parabola, but instead it is a section of a sphere. We will call the
problem this creates, \emph{spherical aberration.}

Remember that reason the we use spherical shapes is that spherical shapes
are easier to make than parabolas or hyperbole, or other shapes. So optics
manufacturers have been using spherical optics for centuries. Later, when we
do lenses we will find there are ways to correct for spherical aberration.

But what does spherical aberration do? We end up with different images of
our object at different image locations. The center rays (blue dashed lines
in the next figure) make an image at $I_{1}$ but the outside (non-paraxial)
rays focus closer at $I_{2.}$ If we put up a screen to see the image, we
will have a blurry image because some of the light is out of focus at $I_{1}$
but a different bit of light is out of fucus at $I_{2}.$ There is no one
place where all the light is in focus. \FRAME{dtbpF}{3.5976in}{2.1975in}{0pt%
}{}{}{Figure}{\special{language "Scientific Word";type
"GRAPHIC";maintain-aspect-ratio TRUE;display "USEDEF";valid_file "T";width
3.5976in;height 2.1975in;depth 0pt;original-width 3.55in;original-height
2.1577in;cropleft "0";croptop "1";cropright "1";cropbottom "0";tempfilename
'SVAM8I41.wmf';tempfile-properties "XPR";}}Most of the time for this class,
we will design our optical system so the object is near the principal axis,
so we can make the paraxial approximation that fixes this problem. But in
real optical systems we can't do this.

This optical problem was the problem with the Hubble Space Telescope!

\FRAME{dtbpFU}{3.5169in}{1.7722in}{0pt}{\Qcb{{\protect\small This before
(left) WFPC and after (right) WFPC2 image of the core of the galaxy M100
shows the dramatic improvement in Hubble's view of the universe after the
first servicing mission in December 1993. (Images Courtesy NASA)}}}{}{Figure%
}{\special{language "Scientific Word";type "GRAPHIC";maintain-aspect-ratio
TRUE;display "USEDEF";valid_file "T";width 3.5169in;height 1.7722in;depth
0pt;original-width 2.0453in;original-height 1.017in;cropleft "0";croptop
"1";cropright "1";cropbottom "0";tempfilename
'HubbleTrouble.wmf';tempfile-properties "XNPR";}}

And it took some work to correct this problem. We will learn more about this
as we study lenses next.

\section{Semi-Infinite Lenses and the Semi-Infinite Lens Image Equation}

We learned how to find an image location graphically for mirrors. But many
optical systems including contacts and glasses, use lenses. Let's look at
lenses next.

Let's start by thinking of a special case for refraction. A curved surface
on a very large piece of glass. We will assume that the piece of glass is
semi-infinate (goes on forever in one direction), but all it has to be is
very large.

I jokingly call this a semi-infinate bump of glass. But I\ have also heard
this called a single sided lens. It may seem like a far fetched situation,
but it turns out that our eyes are essentially a system like this because
our eyes are filled with a jelly like substance, and the light in our eyes
stays in this substance until it is detected in our retina. So let's look at
this situation. \FRAME{dhF}{3.2491in}{1.1865in}{0pt}{}{}{Figure}{\special%
{language "Scientific Word";type "GRAPHIC";maintain-aspect-ratio
TRUE;display "USEDEF";valid_file "T";width 3.2491in;height 1.1865in;depth
0pt;original-width 5.2555in;original-height 1.9009in;cropleft "0";croptop
"1";cropright "1";cropbottom "0";tempfilename
'SVANNN56.wmf';tempfile-properties "XPR";}}

Take a point object that either glows, or has rays of light reflecting from
it. The rays leave the object and reach the surface of the glass. The rays
will refract at the surface. Each bends toward the normal, but because of
the curvature of the glass, every ray has a different normal. \FRAME{dtbpF}{%
3.351in}{1.2489in}{0in}{}{}{Figure}{\special{language "Scientific Word";type
"GRAPHIC";maintain-aspect-ratio TRUE;display "USEDEF";valid_file "T";width
3.351in;height 1.2489in;depth 0in;original-width 3.3856in;original-height
1.2444in;cropleft "0";croptop "1";cropright "1";cropbottom "0";tempfilename
'TA9KVM02.wmf';tempfile-properties "XPR";}}The result is tat the rays all
converge toward the center. We can identify this convergence point as the
image of the point object. For mirrors we used the law of reflection to know
where the light would go. But for lenses, the light is transmitted. So at
the surface we can find the refracted angles using Snell's law%
\begin{equation*}
n_{1}\sin \theta _{1}=n_{2}\sin \theta _{2}
\end{equation*}

We will again use the small angle approximation. Then $\theta _{1}$ and $%
\theta _{2}$ are small and none of the rays are very far away from the axis.
We know this as the paraxial approximation. Snell's law becomes 
\begin{equation*}
n_{1}\theta _{1}=n_{2}\theta _{2}
\end{equation*}

\FRAME{dtbpF}{5.041in}{1.9441in}{0pt}{}{}{Figure}{\special{language
"Scientific Word";type "GRAPHIC";maintain-aspect-ratio TRUE;display
"USEDEF";valid_file "T";width 5.041in;height 1.9441in;depth
0pt;original-width 4.4469in;original-height 1.6985in;cropleft "0";croptop
"1";cropright "1";cropbottom "0";tempfilename
'SWDHT104.wmf';tempfile-properties "XPR";}}Using the more detailed figure
above with it's exaggerated angles (not so paraxial but easier to see), we
observe triangles $SAC$ and $PAC.$ We can see that for triangle $SAC$ the
top angle labeled $\eta $, plus $\theta _{1}$ must be $180.$ 
\begin{equation*}
\eta +\theta _{1}=180\unit{%
%TCIMACRO{\U{b0}}%
%BeginExpansion
{{}^\circ}%
%EndExpansion
}
\end{equation*}%
We also know that the sum of interior angles of a triangle must equal $180.$
So%
\begin{equation*}
180\unit{%
%TCIMACRO{\U{b0}}%
%BeginExpansion
{{}^\circ}%
%EndExpansion
}=\alpha +\beta +\eta
\end{equation*}%
then 
\begin{eqnarray*}
\eta +\theta _{1} &=&\alpha +\beta +\eta \\
\theta _{1} &=&\alpha +\beta
\end{eqnarray*}%
Likewise, from triangle $PAC$,%
\begin{equation*}
\beta +\psi =180\unit{%
%TCIMACRO{\U{b0}}%
%BeginExpansion
{{}^\circ}%
%EndExpansion
}
\end{equation*}

\begin{equation*}
180\unit{%
%TCIMACRO{\U{b0}}%
%BeginExpansion
{{}^\circ}%
%EndExpansion
}=\psi +\gamma +\theta _{2}
\end{equation*}%
so then 
\begin{equation*}
\beta =\theta _{2}+\gamma
\end{equation*}%
then, 
\begin{equation*}
\theta _{2}=\beta -\gamma
\end{equation*}

and we can write our paraxial Snell's law as%
\begin{eqnarray*}
n_{1}\theta _{1} &=&n_{2}\theta _{2} \\
n_{1}\left( \alpha +\beta \right) &=&n_{2}\left( \beta -\gamma \right) \\
n_{1}\alpha +n_{1}\beta &=&n_{2}\beta -n_{2}\gamma \\
n_{1}\alpha +n_{2}\gamma &=&n_{2}\beta -n_{1}\beta \\
n_{1}\alpha +n_{2}\gamma &=&\beta \left( n_{2}-n_{1}\right)
\end{eqnarray*}

Looking at the figure. We see that $d$ is a leg of three different right
triangles ($SAV,$ $ACV$, and $PAV$). The ray in the figure is clearly not a
paraxial ray. If we use a paraxial ray, then the point $V$ will approach the
air-glass boundary. When this happens, then $SV=d_{o},$ $VC=R$, and $%
VP=d_{i}.$ So we can write (using the small angle approximation)

\begin{eqnarray*}
\tan \alpha &\approx &\alpha \approx \frac{d}{d_{o}} \\
\tan \beta &\approx &\beta \approx \frac{d}{R} \\
\tan \gamma &\approx &\gamma \approx \frac{d}{d_{i}}
\end{eqnarray*}%
So, our Snell's law becomes%
\begin{eqnarray*}
n_{1}\alpha +n_{2}\gamma &=&\beta \left( n_{2}-n_{1}\right) \\
n_{1}\frac{d}{d_{o}}+n_{2}\frac{d}{d_{i}} &=&\frac{d}{R}\left(
n_{2}-n_{1}\right)
\end{eqnarray*}%
We can divide out the $d$'s%
\begin{equation}
\frac{n_{1}}{d_{o}}+\frac{n_{2}}{d_{i}}=\frac{\left( n_{2}-n_{1}\right) }{R}
\label{Bump_Equation}
\end{equation}

We can use this formula to convince ourselves that no matter what the angle
is (providing it is small), the rays will form an image at $P.$

Real images will be in the glass for this special case (we will fix this
with non-infinite lenses soon) so we need to change sign conventions.

\begin{equation*}
\begin{tabular}{|l|l|l|}
\hline
\textbf{Quantity} & \textbf{Positive if} & \textbf{Negative if} \\ \hline
{\small Object location }$\left( d_{o}\right) $ & 
\begin{tabular}{l}
{\small Object is in front of interface} \\ 
{\small \ surface}%
\end{tabular}
& 
\begin{tabular}{l}
{\small Object is in back of interface} \\ 
{\small \ surface (virtual object)}%
\end{tabular}%
{\small \ } \\ \hline
{\small Image location }$\left( d_{i}\right) $ & 
\begin{tabular}{l}
{\small Image is in back of interface} \\ 
{\small \ surface (real image)}%
\end{tabular}
& 
\begin{tabular}{l}
{\small Image is in front of interface} \\ 
{\small \ surface (virtual image)}%
\end{tabular}
\\ \hline
{\small Image height }$\left( h_{i}\right) $ & 
\begin{tabular}{l}
{\small Image is upright}%
\end{tabular}
& 
\begin{tabular}{l}
{\small Image is inverted}%
\end{tabular}
\\ \hline
{\small Radius }$\left( R\right) $ & 
\begin{tabular}{l}
{\small Center of curvature is in} \\ 
{\small back of interface surface}%
\end{tabular}%
{\small \ } & 
\begin{tabular}{l}
{\small Center of curvature is in} \\ 
{\small front of interface surface}%
\end{tabular}%
{\small \ } \\ \hline
\end{tabular}%
\end{equation*}

We could go through the entire derivation and switch the indices of
refraction. It turns out we get the same equation. The results will be
different, but the equation is the same.

\subsection{Flat Refracting surfaces}

\FRAME{dtbpF}{2.6282in}{2.4647in}{0in}{}{}{Figure}{\special{language
"Scientific Word";type "GRAPHIC";maintain-aspect-ratio TRUE;display
"USEDEF";valid_file "T";width 2.6282in;height 2.4647in;depth
0in;original-width 2.5867in;original-height 2.4232in;cropleft "0";croptop
"1";cropright "1";cropbottom "0";tempfilename
'SWDHUE05.wmf';tempfile-properties "XPR";}}

Let's consider a fish tank. The fish tank has an interface, but it is flat.
Can we use our equation (\ref{Bump_Equation}) to describe this?

The answer is yes, if we let $R=\infty .$ This makes sense for a flat
surface. If we have an infinitely large sphere, then our small part of that
spherical surface that makes up the fish tank wall will be very flat.

Then%
\begin{equation*}
\frac{n_{1}}{d_{o}}+\frac{n_{2}}{d_{i}}=\frac{\left( n_{2}-n_{1}\right) }{%
\infty }
\end{equation*}%
or%
\begin{equation*}
\frac{n_{1}}{d_{o}}+\frac{n_{2}}{d_{i}}=0
\end{equation*}%
we see that 
\begin{equation*}
d_{i}=-d_{o}\frac{n_{2}}{n_{1}}
\end{equation*}

We can see from this that, since $n_{2}<n_{1}$ for this case $d_{i}$ is less
than $d_{o}.$ The image is in the water! It is in \emph{front} of the
water/air boundary of the tank. But light is not actually coming from that
location. What has happened is that the light path has bent, but our brain
image processing center thinks light travels only in straight lines. So it
tells us there is a fish there, but the fish appears closer than it is. The
light doesn't come from where we think the fish is. And as we know, this
means the image is virtual, just like for flat mirrors!

\chapter{16 Thin Lenses 3.2.4, 3.2.5}

%TCIMACRO{\TeXButton{\chaptermark{Thin Lenses}}{\chaptermark{Thin Lenses}}}%
%BeginExpansion
\chaptermark{Thin Lenses}%
%EndExpansion

In our last lecture we found an equation for finding the image location from
a single sided lens (or semi-infinate bump). But most lenses have two sides
(think of eye glass lenses). We need to do a little more work to find an
equation for the location of an image with a double sided lens. We will to
this work in this lecture.

%TCIMACRO{%
%\TeXButton{Fundamental Concepts}{\vspace{0.10in}\hspace{-0.37in}{\Large Fundamental Concepts\vspace{0.2in}}}}%
%BeginExpansion
\vspace{0.10in}\hspace{-0.37in}{\Large Fundamental Concepts\vspace{0.2in}}%
%EndExpansion

\begin{itemize}
\item For a thin lens, we can describe where the light goes using $\frac{1}{%
d_{o}}+\frac{1}{d_{i}}=\frac{1}{f}$ where $\frac{1}{f}=\left( n-1\right)
\left( \frac{1}{R_{1}}-\frac{1}{R_{2}}\right) $

\item There is a sign convention for all of these equations.

\item We can think of a complicated lens as a series of interfaces between
materials.

\item We can find the image location of a complicated set of interface
surfaces by going through each surface and finding the image created by that
surface, then using that image as the object of the next surface.

\item The wavelength dependence of the index of refraction gives rise to
another problem, chromatic aberration.
\end{itemize}

\section{Images from lenses}

Let's think once again about what an image is. If you haven't done this
before, take a piece of paper and a lens, and hold up the lens in a darkened
room that has some bright object in it. Move the lens or the paper back and
forth, and at just the right distance, a miniature picture of the bright
object will appear. We should think about what the word \textquotedblleft
picture\textquotedblright\ means in this sense.

\FRAME{dhF}{2.6835in}{2.0237in}{0pt}{}{}{Figure}{\special{language
"Scientific Word";type "GRAPHIC";maintain-aspect-ratio TRUE;display
"USEDEF";valid_file "T";width 2.6835in;height 2.0237in;depth
0pt;original-width 2.6411in;original-height 1.9847in;cropleft "0";croptop
"1";cropright "1";cropbottom "0";tempfilename
'SVAL9A3I.wmf';tempfile-properties "XPR";}}We have talked about how we see
objects. Remember the BYU-I guys from a few lectures ago.\FRAME{dhF}{3.173in%
}{1.3534in}{0pt}{}{}{Figure}{\special{language "Scientific Word";type
"GRAPHIC";maintain-aspect-ratio TRUE;display "USEDEF";valid_file "T";width
3.173in;height 1.3534in;depth 0pt;original-width 3.128in;original-height
1.3188in;cropleft "0";croptop "1";cropright "1";cropbottom "0";tempfilename
'SVAL9A3J.wmf';tempfile-properties "XPR";}}Our eyes gather rays that are
diverging from the object because light has bounced off of the object. Our
eyes intersect a diverging set of rays that form a definite pattern. That
diverging set of rays forming a pattern is the picture of the object.

So when we say that the lens has formed a miniature picture of our object,
we mean that the lens has somehow formed a diverging set of rays that form a
pattern that looks like the pattern formed by the diverging set of rays
coming from the object, itself. In other words, the object forms a diverging
set of rays. And our lens forms a duplicate set of rays in the same pattern,
so we see the same thing. The lens' version may be smaller, upside down, and
backwards, but it is still essentially the same pattern. We saw this with
our semi-infinite lens, but it must be true for less infinite lenses as well.

As a first step to see how this works, consider our fish tank again. It
would be bad on the fish, but think about looking at a fish in air. The room
light would bounce off of the fish, and we would have a diverging set of
rays from every point on the fish (see next figure). \FRAME{dtbpF}{4.1468in}{%
2.0643in}{0in}{}{}{Figure}{\special{language "Scientific Word";type
"GRAPHIC";maintain-aspect-ratio TRUE;display "USEDEF";valid_file "T";width
4.1468in;height 2.0643in;depth 0in;original-width 4.0975in;original-height
2.0254in;cropleft "0";croptop "1";cropright "1";cropbottom "0";tempfilename
'SWMIFO02.wmf';tempfile-properties "XPR";}}We can see that the picture is
made from every point on the fish being \textquotedblleft
imaged\textquotedblright\ to a point on our retina. We collect the rays
leaving every point on the fish, and bring them to corresponding points on
the retina to make the picture.

It will take us a few lectures to see exactly how this is done, but as a
first step, let's consider the fish tank, itself. Put the fish back in the
tank and look at it (next figure).

\FRAME{dtbpF}{3.9115in}{2.4811in}{0pt}{}{}{Figure}{\special{language
"Scientific Word";type "GRAPHIC";maintain-aspect-ratio TRUE;display
"USEDEF";valid_file "T";width 3.9115in;height 2.4811in;depth
0pt;original-width 4.3552in;original-height 2.7518in;cropleft "0";croptop
"1";cropright "1";cropbottom "0";tempfilename
'SWMIHW03.wmf';tempfile-properties "XPR";}}Rays still come from the fish.
But we now know that the change from a slow material to a fast material will
bend the light. These bent rays are collected by our eyes, and the picture
of the fish is formed on the retina just as before. But our eyes interpret
the light as though it went in straight lines with no bends (dotted lines in
the last figure). our visual processing center in our brain is designed to
believe light travels in straight lines, so our mind tells us there is a
fish, but that the fish head (and every other part of the fish) is closer
than it really is. We call this apparent fish at the closer location an
image of the fish, because this is where our brain processing center thinks
the diverging set of rays come from that form the fish pattern.

The next figure shows the details of the rays leaving a dot on the fish head
(with the angles exaggerated so it's easier to see them). \FRAME{dtbpF}{%
2.7363in}{2.565in}{0pt}{}{}{Figure}{\special{language "Scientific Word";type
"GRAPHIC";maintain-aspect-ratio TRUE;display "USEDEF";valid_file "T";width
2.7363in;height 2.565in;depth 0pt;original-width 2.6939in;original-height
2.5235in;cropleft "0";croptop "1";cropright "1";cropbottom "0";tempfilename
'SWQ0MD00.wmf';tempfile-properties "XPR";}}A spot on the fish head is our
object for this set of rays. The distance from the fish-head dot and the
edge of the water/air boundary is our object distance,\emph{\ }$d_{o}.$

The distance from the image of the fish-head dot to the edge of the
water/air boundary is our image distance, $d_{i},$ because it appears to be
where the rays come from. We can find where the image is $\left(
d_{i}\right) $ knowing $d_{o}.$ We can see from the figure that%
\begin{eqnarray*}
\ell &=&d_{o}\tan \theta _{1} \\
&=&d_{i}\tan \theta _{2}
\end{eqnarray*}%
so%
\begin{equation*}
d_{o}\tan \theta _{1}=d_{i}\tan \theta _{2}
\end{equation*}%
or%
\begin{equation*}
d_{o}\frac{\tan \theta _{1}}{\tan \theta _{2}}=d_{i}
\end{equation*}%
from Snell's law, we know that 
\begin{equation*}
\frac{\sin \theta _{1}}{\sin \theta _{2}}=\frac{n_{2}}{n_{1}}
\end{equation*}

Usually we can take the small angle approximation. This would limit our
analysis to rays that are near the central axes. We have done this before
with mirrors. We call this central axis the \emph{optics axis} and the rays
that are not too far away from this axis \emph{paraxial rays.} Then for our
small angles we can write 
\begin{equation*}
\tan \theta _{i}\approx \sin \theta _{i}
\end{equation*}%
so%
\begin{equation*}
d_{o}\frac{\tan \theta _{1}}{\tan \theta _{2}}\approx d_{o}\frac{\sin \theta
_{1}}{\sin \theta _{2}}=d_{o}\frac{n_{2}}{n_{1}}=d_{i}
\end{equation*}%
and we have the image distance

\begin{equation*}
d_{i}=d_{o}\frac{n_{2}}{n_{1}}
\end{equation*}%
Note that once again we have a virtual image. The place where we see the
fish does not have rays of light coming from it! But our brain processing
center believes the rays of light didn't bend, so it extends the rays we see
backward into the fish tank and it interprets what we see as a fish being
where those extended rays meet.

This is not so useful unless you have some burning need to know where fish
are in a tank. But we now have the vocabulary to discuss the larger problem
of how a lens works, which we will take up next time.

Before we do a lot of math to describe how lenses work, lets think about our
early childhood experiences. You may have burned things with a magnifying
glass\footnote{%
If you didn't do this as a child, consider trying it now. Be responsible and
safe, but it is valuable to see how this works.}. Using the idea of a ray
diagram, here is what happens.\FRAME{dhF}{3.8017in}{2.2329in}{0pt}{}{}{Figure%
}{\special{language "Scientific Word";type "GRAPHIC";maintain-aspect-ratio
TRUE;display "USEDEF";valid_file "T";width 3.8017in;height 2.2329in;depth
0pt;original-width 3.7533in;original-height 2.1923in;cropleft "0";croptop
"1";cropright "1";cropbottom "0";tempfilename
'SVAJW81M.wmf';tempfile-properties "XPR";}}We know that the rays from the
Sun come from so far away that they are essentially parallel. We know that
these rays come together to a fine point that can start a fire. The point
where these rays converge is important to us. We found such a point from our
mirrors. We called this point the \emph{focal point}.

Now think about Snell's law. We have 
\begin{equation*}
n_{1}\sin \theta _{1}=n_{2}\sin \theta _{2}
\end{equation*}%
The light will bend at a surface, but snell's law doesn't care if $%
n_{1}>n_{2}$ of if $n_{2}>n_{1}.$ It works both ways. The light will follow
the same path either direction. We would expect that if we put a light
source at the focal distance, the rays should come out parallel. \FRAME{dhF}{%
4.1926in}{2.4561in}{0pt}{}{}{Figure}{\special{language "Scientific
Word";type "GRAPHIC";maintain-aspect-ratio TRUE;display "USEDEF";valid_file
"T";width 4.1926in;height 2.4561in;depth 0pt;original-width
4.1433in;original-height 2.4146in;cropleft "0";croptop "1";cropright
"1";cropbottom "0";tempfilename 'SVAJW81N.wmf';tempfile-properties "XPR";}}%
This is one way manufacturers make LED flashlights.

We need one other bit of information that we already have seen from basic
refraction.\FRAME{dhF}{1.996in}{1.6933in}{0pt}{}{}{Figure}{\special{language
"Scientific Word";type "GRAPHIC";maintain-aspect-ratio TRUE;display
"USEDEF";valid_file "T";width 1.996in;height 1.6933in;depth
0pt;original-width 3.6841in;original-height 3.122in;cropleft "0";croptop
"1";cropright "1";cropbottom "0";tempfilename
'SVAJW81O.wmf';tempfile-properties "XPR";}}A flat block does refract the
light, but when the light leaves the block it is only displaced, it retains
the original direction. With these three ray scenarios, we can describe how
a lens works.

We know that for every point on the object, we get millions of reflected
rays that diverge. In the next figure, rays of light are reflecting off of
the top of the large, black arrow. The lens must collect these rays together
to form the corresponding point on the image.

\FRAME{dhF}{4.1373in}{2.4275in}{0pt}{}{}{Figure}{\special{language
"Scientific Word";type "GRAPHIC";maintain-aspect-ratio TRUE;display
"USEDEF";valid_file "T";width 4.1373in;height 2.4275in;depth
0pt;original-width 4.088in;original-height 2.3869in;cropleft "0";croptop
"1";cropright "1";cropbottom "0";tempfilename
'SVAJW81P.wmf';tempfile-properties "XPR";}}Of course, this happens for every
point on the image. \FRAME{dhF}{4.1926in}{2.4561in}{0pt}{}{}{Figure}{\special%
{language "Scientific Word";type "GRAPHIC";maintain-aspect-ratio
TRUE;display "USEDEF";valid_file "T";width 4.1926in;height 2.4561in;depth
0pt;original-width 4.1433in;original-height 2.4146in;cropleft "0";croptop
"1";cropright "1";cropbottom "0";tempfilename
'SVAJW81Q.wmf';tempfile-properties "XPR";}}But we usually pick the top of
the object. If we place the bottom of the object on the \emph{optic axis}%
\footnote{%
The line drawn on the figure thorugh the middle of the lens.}, the bottom of
the image will also be on the optic axis. So knowing where the bottom of the
image is, and finding the top of the image gives a pretty good idea of where
the rest of the image must be. So we will draw diagrams for the top of the
object to find the top of the image.

But suppose this is not true? For example, when we use a camera, we do not
align the object with it's bottom on the axis before we shoot. We center the
object on our viewfinder. So how do we find the location for the bottom of
the image?\FRAME{dtbpF}{4.1329in}{2.4967in}{0pt}{}{}{Figure}{\special%
{language "Scientific Word";type "GRAPHIC";maintain-aspect-ratio
TRUE;display "USEDEF";valid_file "T";width 4.1329in;height 2.4967in;depth
0pt;original-width 4.9199in;original-height 2.962in;cropleft "0";croptop
"1";cropright "1";cropbottom "0";tempfilename
'SWMIR805.wmf';tempfile-properties "XPR";}}We can, of course, trace two more
rays for the bottom of the image and find the location of the bottom of the
image as well.

We know that light is a superposition of millions of waves and that when
light reflects off of something like your neighbor or a fish or a soccer
player there will be millions of rays. Drawing millions of rays is
impractical, and fortunately, not needed. From our study of mirrors, we can
guess that we can find three simple rays that leave the top of the object
and see where these rays converge to form the top of the image. Let's start
with a ray that travels from the top of the object and travels parallel to
the optic axis. We recognize this ray as being like one of the rays from the
Sun. It comes in parallel, so it will leave the lens and travel through the
focal point.\FRAME{dhF}{4.3889in}{2.5815in}{0pt}{}{}{Figure}{\special%
{language "Scientific Word";type "GRAPHIC";maintain-aspect-ratio
TRUE;display "USEDEF";valid_file "T";width 4.3889in;height 2.5815in;depth
0pt;original-width 4.3379in;original-height 2.5399in;cropleft "0";croptop
"1";cropright "1";cropbottom "0";tempfilename
'SVAJW81S.wmf';tempfile-properties "XPR";}}For a second easy ray, lets take
the case that is like our flat block. Near the center of the lens, the sides
are nearly flat. So we expect that the ray will leave in about the same
direction as it was going before it struck the lens.\FRAME{dhF}{3.1522in}{%
1.8542in}{0pt}{}{}{Figure}{\special{language "Scientific Word";type
"GRAPHIC";maintain-aspect-ratio TRUE;display "USEDEF";valid_file "T";width
3.1522in;height 1.8542in;depth 0pt;original-width 8.9949in;original-height
5.2736in;cropleft "0";croptop "1";cropright "1";cropbottom "0";tempfilename
'SVAJW81T.wmf';tempfile-properties "XPR";}}Technically this ray should be
shifted. This is where our thin lens approximation comes in. If the lens is
thin, then the ray going through the center of the lens won't be shifted
much, and we can ignore the shift. We just dray ray 2 as a straight line.

Two rays are really enough to determine where the top of the image will be,
but there is a third ray that is easy to draw, so let's draw it to give us
more confidence in our answer. That ray is one that leaves the top of the
object and passes through the focal point on the object side of the lens.
This situation we also recognize. Since the ray goes through the focal
point, it is as though the light came from that point. This is like our LED
flashlight case. This ray will leave the lens parallel to the optic axis.

\FRAME{dhF}{4.2627in}{2.4976in}{0pt}{}{}{Figure}{\special{language
"Scientific Word";type "GRAPHIC";maintain-aspect-ratio TRUE;display
"USEDEF";valid_file "T";width 4.2627in;height 2.4976in;depth
0pt;original-width 4.2125in;original-height 2.4569in;cropleft "0";croptop
"1";cropright "1";cropbottom "0";tempfilename
'SVAJW81U.wmf';tempfile-properties "XPR";}}Where all three rays intersect,
we will have the top of the image. \FRAME{dhF}{4.9753in}{2.1067in}{0pt}{}{}{%
Figure}{\special{language "Scientific Word";type
"GRAPHIC";maintain-aspect-ratio TRUE;display "USEDEF";valid_file "T";width
4.9753in;height 2.1067in;depth 0pt;original-width 4.9216in;original-height
2.0678in;cropleft "0";croptop "1";cropright "1";cropbottom "0";tempfilename
'SVAJW81V.wmf';tempfile-properties "XPR";}}Notice that in this case, the
image is upside down. That is normal. Also notice that it is smaller than
the object. We say that the image is magnified, which may still seem a
little bit strange. But in optics, a magnification of greater than one means
that the image is bigger than the object. This is like a movie projector
that makes a large image of a small film segment. The magnification can be
equal to one, meaning the object and image are the same size. And finally
the magnification can be less than one. This means that the image is smaller
than the object. This is a convenient definition, because then we can use
the same equation to describe all three situations.%
\begin{equation*}
m\equiv \frac{\text{Image height}}{\text{Object height}}=\frac{h_{i}}{h_{o}}
\end{equation*}%
where $h_{o}$ is the object height, and $h_{i}$ is the image height. This is
the same equation for magnification that we found for mirrors! It's still
true that by convention (meaning physicists got together and voted on this)
we say that an upside down image has a negative magnification. You just have
to memorize this, there is no obvious reason for this except it is
mathematically convenient.

\section{Virtual Images with Lenses}

Lets take another case and draw a ray diagram. This time let's place the
object closer than the focal distance. This is the case when we use a lens
as a magnifying glass. The rays will look like this.\FRAME{dhF}{3.6201in}{%
2.3713in}{0pt}{}{}{Figure}{\special{language "Scientific Word";type
"GRAPHIC";maintain-aspect-ratio TRUE;display "USEDEF";valid_file "T";width
3.6201in;height 2.3713in;depth 0pt;original-width 3.5725in;original-height
2.3315in;cropleft "0";croptop "1";cropright "1";cropbottom "0";tempfilename
'SVAJW81W.wmf';tempfile-properties "XPR";}}Notice that these rays never
converge! We won't get an image that could project on a paper. But we know
that there is an image, we can look through the lens and see it! And that is
the key. The image does not really exist. There is no place where light
gathers and then diverges into a pattern that we recognize. But we look
through the lens, and our mind interprets the diverging rays coming from the
lens as though they had only traveled in straight lines. If we extend these
rays backwards along straight lines, they appear to come from a common
point. This is the point they would have had to have come from if there were
no lens. We believe we see an image at this location. But no light really
goes there! Because this image is not really made from light diverging from
this position, we call it a \emph{virtual image} just like we did for this
situation for mirrors The image we formed before that could be projected on
a screen is called a \emph{real image}. For a real image, the light is
actually at the image location. This is just like our image terminology for
mirrors!

By convention, we say the distance from the lens to the virtual image is a
negative value.

\subsection{Diverging Lenses}

So far our lenses have only been the sort that work as magnifying glasses.
We call these \emph{converging lenses.} These lenses are thicker in the
middle and thinner on the edges. Because of this they are sometimes called 
\emph{convex lenses}. By convention, we say the focal distance for this type
of lens is positive. For this reason, they are often called \emph{positive
lenses.}

But what if we make a lens that is thinner in the middle and thicker on the
edges. We can call this sort of lens a \emph{concave lens}, and we will give
it a negative focal length by convention, so we can also call it a \emph{%
negative lens}. But what would this lens do? If we think about our three
rays, ray $1$ won't be bent toward the optic axis for this type of lens. In
fact, if we observe an object through this lens, ray number 1 will appear to
come from the focal point. Ray number 2 will still go through the middle of
the lens, and if the lens is thin enough, ray 2 will pass through undeviated.

\FRAME{dhF}{3.2984in}{1.7582in}{0pt}{}{}{Figure}{\special{language
"Scientific Word";type "GRAPHIC";maintain-aspect-ratio TRUE;display
"USEDEF";valid_file "T";width 3.2984in;height 1.7582in;depth
0pt;original-width 3.2534in;original-height 1.7201in;cropleft "0";croptop
"1";cropright "1";cropbottom "0";tempfilename
'SVAJW81X.wmf';tempfile-properties "XPR";}}Finally ray three leaves the
object in the direction of the far focal point. It will hit the lens and
leave parallel to the optic axis. From the figure we see that these three
rays will never converge. This is like our convex mirror! And from our
mirror experience we expect the rays (with help from our brain) will form a
virtual image. If we extend the rays backward as shown, we see that the
extensions all meet at a point. The rays leaving the lens appear to come
from this point. This is the location of the top of the virtual image of the
object.

You might wonder what good such a lens could do, but we will find that this
type of lens is used to correct vision for nearsighted people. It's likely
that many people in our class have this type of lens with them either on
their eye (contacts) or on their nose (eye glasses).

\section{Thin Lenses}

Lets' find an equation for a spherical surface once more. But this time,
let's let it be more practical and not make the \textquotedblleft
lens\textquotedblright\ semi-infinate. We will need to deal with two sides
of the lens because (usually) both will be curved.

We found that for refraction%
\begin{equation*}
\frac{n_{1}}{d_{o}}+\frac{n_{2}}{d_{i}}=\frac{\left( n_{2}-n_{1}\right) }{R}
\end{equation*}

but we did this for a spherical bump on a semi-infinite piece of glass. For
this new problem let's make a few assumptions:

\begin{enumerate}
\item We have two spherical surfaces, with $R_{1}$ and $R_{2}$ as the radii
of curvature

\item We have only paraxial rays

\item The image formed by one refractive surface serves as the object for
the second surface

\item The lens is not very thick (the thickness is much smaller than both $%
R_{1}$ and $R_{2})$
\end{enumerate}

The answer we will get is quite simple

\begin{equation}
\frac{1}{d_{o}}+\frac{1}{d_{i}}=\frac{1}{f}  \label{Thin_lens}
\end{equation}%
where%
\begin{equation}
\frac{1}{f}=\left( n-1\right) \left( \frac{1}{R_{1}}-\frac{1}{R_{2}}\right)
\label{Lens_maker}
\end{equation}%
but to appreciate what it means, lets find out where it comes from.

\subsection{Derivation of the lens equation}

Consider the optical element in the figure below. Notice that our object is
a dot, so our image will also be a dot. This is not as boring as it sounds
if we consider any object can be considered as a collection of dots. And in
our past analysis we used the top of our object as our image, which was
really just the dot at the top of the object. We could use that to figure
out where the rest of the image would be. So let's start with a dot.\FRAME{%
dtbpF}{4.5714in}{2.8098in}{0pt}{}{}{Figure}{\special{language "Scientific
Word";type "GRAPHIC";maintain-aspect-ratio TRUE;display "USEDEF";valid_file
"T";width 4.5714in;height 2.8098in;depth 0pt;original-width
4.5186in;original-height 2.7674in;cropleft "0";croptop "1";cropright
"1";cropbottom "0";tempfilename 'SWMJFK06.wmf';tempfile-properties "XPR";}}%
Light enters at a spherical surface on the left hand side. We use a point
object located at $S$ (for \emph{source}) on the principal axis. We could
put our dot object anywhere, but let's put it on the axis and trace two
rays. The ray along the principal axis crosses each spherical surface at
right angles, and therefore traves straight through the optic (this makes
one ray very easy to trace!). The second ray hits the first spherical
surface at point $A.$ It is refracted and travels to point $B$. It is again
refracted and travels toward the principal axis, crossing at $P.$ The image
location is the intersection of these rays, so we have an image at $P.$

Lets study the surfaces separately

\textbf{Surface 1:}

\FRAME{dtbpF}{4.3448in}{2.6117in}{0pt}{}{}{Figure}{\special{language
"Scientific Word";type "GRAPHIC";maintain-aspect-ratio TRUE;display
"USEDEF";valid_file "T";width 4.3448in;height 2.6117in;depth
0pt;original-width 4.9147in;original-height 2.9438in;cropleft "0";croptop
"1";cropright "1";cropbottom "0";tempfilename
'SWMJGL07.wmf';tempfile-properties "XPR";}}

Let's treat surface $1$ as though surface $2$ did not exist. When it's at
surface $1$ The light doesn't know about surface $2.$ It only interacts with
the surface where it is. So ignoring surface $2$ at this point is reasonable.

The light would bend at point $A$ and head off into the lens material. This
is just our single sided lens (semi-infinite bump) problem. And if we
imagine our head to somehow be in the glass looking at the refracted ray,
our brain would think that the ray traveled in a straight line. So our brain
would think the light came from location $P_{1}.$ That would be a virtual
image of our $S$ point, so the distance from the lens to $P_{1}=d_{i1}.$ We
can find this image distance using our single sided lens (semi-infinte bump)
equation.%
\begin{equation*}
\frac{n_{1}}{d_{o1}}+\frac{n_{2}}{d_{i1}}=\frac{\left( n_{2}-n_{1}\right) }{R%
}
\end{equation*}

We can consider $n_{1}=1$ and $n_{2}=n_{glass}$ for a air-glass interface
and noting that $d_{i1}$ is negative (by convention). Then%
\begin{equation}
\frac{1}{d_{o1}}-\frac{n_{glass}}{d_{i1}}=\frac{\left( n_{glass}-1\right) }{%
R_{1}}  \label{Surface1}
\end{equation}

Note that our rays are \emph{not} converging in the glass this time. We can
find the image formed by this side of our lens by tracing the diverging rays
backward as we did for the fish tank. The image formed from the first side
of the lens is very much virtual.

\textbf{Surface 2}

Now consider the second surface.

\FRAME{dtbpF}{4.4348in}{2.2286in}{0pt}{}{}{Figure}{\special{language
"Scientific Word";type "GRAPHIC";maintain-aspect-ratio TRUE;display
"USEDEF";valid_file "T";width 4.4348in;height 2.2286in;depth
0pt;original-width 4.3837in;original-height 2.188in;cropleft "0";croptop
"1";cropright "1";cropbottom "0";tempfilename
'SWMKFE08.wmf';tempfile-properties "XPR";}}The second surface isn't alive
and doesn't know about how light traves. All the second surface sees is
light diverging as though it came from a semi-infinite piece of glass with
the object at $P_{1}.$ It's not true, the light came from $S.$ But the
second surface can't tell the difference. It only knows light is incident in
a particular pattern, and that pattern appears to come from point $P_{1}.$
So we can treat the situation at surface $2$ as though the object is the
virtual image formed by surface $1.$ This is important, so let's say it
again. The object for surface $2$ is the image from surface $1.$ 
\begin{equation*}
d_{o2}=d_{i1}+t
\end{equation*}

We again use our refractive equation%
\begin{equation*}
\frac{n_{1}}{d_{o2}}+\frac{n_{2}}{d_{i2}}=\frac{\left( n_{2}-n_{1}\right) }{R%
}
\end{equation*}%
but this time we identify $n_{1}=n_{glass}$ and $n_{2}=1$. We have for
surface $2$%
\begin{equation}
\frac{n_{glass}}{d_{o2}}+\frac{1}{d_{i2}}=\frac{\left( 1-n_{glass}\right) }{%
R_{2}}
\end{equation}%
or%
\begin{equation}
\frac{n_{glass}}{d_{i1}+t}+\frac{1}{d_{i2}}=\frac{\left( 1-n_{glass}\right) 
}{R_{2}}  \label{Surface2}
\end{equation}

Now we take our thin lens approximation. Let $t\rightarrow 0.$ Then
equations (\ref{Surface1}) and (\ref{Surface2}) become%
\begin{equation*}
\frac{1}{d_{o1}}-\frac{n_{glass}}{d_{i1}}=\frac{\left( n_{glass}-1\right) }{%
R_{1}}
\end{equation*}

\begin{equation*}
\frac{n_{glass}}{d_{i1}}+\frac{1}{d_{i2}}=\frac{\left( 1-n_{glass}\right) }{%
R_{2}}
\end{equation*}

We would prefer a nice equation that has how far away the object is from the
lens and how far away the image is formed like our mirror equation. We can
get such an equation by adding these two equations for the two surfaces%
\begin{equation*}
\frac{1}{d_{o1}}-\frac{n_{glass}}{d_{i1}}+\frac{n_{glass}}{d_{i1}}+\frac{1}{%
d_{i2}}=\frac{\left( n_{glass}-1\right) }{R_{1}}+\frac{\left(
1-n_{glass}\right) }{R_{2}}
\end{equation*}%
or%
\begin{equation*}
\frac{1}{d_{o1}}+\frac{1}{d_{i2}}=\left( n_{glass}-1\right) \left( \frac{1}{%
R_{1}}-\frac{1}{R_{2}}\right)
\end{equation*}

This equation is very useful. It reminds us of the mirror equation (well, a
little) . If we again let $d_{o1}=\infty $ (put the object at $\infty $ so
the rays enter surface $1$ parallel) we find%
\begin{equation*}
\frac{1}{d_{i2}}=\left( n-1\right) \left( \frac{1}{R_{1}}-\frac{1}{R_{2}}%
\right)
\end{equation*}%
The spot where the rays gather if the object is infinitely far away is the
focal point, $f.$ so for parallel rays we can identify $d_{i2}=f$ as the
focal length of the optic. That means that 
\begin{equation*}
\frac{1}{f}=\left( n-1\right) \left( \frac{1}{R_{1}}-\frac{1}{R_{2}}\right)
\end{equation*}%
which is known as the \emph{lens makers' equation}.

Then we have a relationship between the object distance in front of the
lens, and the final image in back of the lens: 
\begin{eqnarray*}
\frac{1}{d_{o1}}+\frac{1}{d_{i2}} &=&\left( n-1\right) \left( \frac{1}{R_{1}}%
-\frac{1}{R_{2}}\right) \\
&=&\frac{1}{f}
\end{eqnarray*}%
We can drop the subscripts (which we can do now that we let $t=0$ since the
internal distances for the inside points are not important).%
\begin{equation*}
\frac{1}{d_{o}}+\frac{1}{d_{i}}=\frac{1}{f}
\end{equation*}%
This is called the \emph{thin lens equation}. The resulting approximate
geometry is shown below.

\FRAME{dtbpF}{5.0548in}{1.5731in}{0pt}{}{}{Figure}{\special{language
"Scientific Word";type "GRAPHIC";maintain-aspect-ratio TRUE;display
"USEDEF";valid_file "T";width 5.0548in;height 1.5731in;depth
0pt;original-width 5.0004in;original-height 1.5376in;cropleft "0";croptop
"1";cropright "1";cropbottom "0";tempfilename
'SWMKHL0A.wmf';tempfile-properties "XPR";}}And this is fantastic! It is the
same equation as the thin mirror equation! Only we have a different equation
for finding the focal length.

We did our math for just one point of light at $S.$ Of course any real
object is made of lots of points, and not all of the points are on an axis.
But each point will be imaged to a corresponding point on the image. Here is
an example with three points in the object (the $s_{i}$ points) and where
their images are (the $p_{i}$ points).\FRAME{dtbpF}{3.1782in}{2.5685in}{0pt}{%
}{}{Figure}{\special{language "Scientific Word";type
"GRAPHIC";maintain-aspect-ratio TRUE;display "USEDEF";valid_file "T";width
3.1782in;height 2.5685in;depth 0pt;original-width 3.1341in;original-height
2.527in;cropleft "0";croptop "1";cropright "1";cropbottom "0";tempfilename
'SWMKIJ0B.wmf';tempfile-properties "XPR";}}but it would work for millions of
points. Our simple analysis explains the formation of actual images and not
just point images.

\subsection{Lens Aberrations}

We found that for mirrors if we used a spherical shape rather than a
parabola we had a problem with our images. We called the problem spiracle
aberration. The reason we used a sphere is that spheres are easy to
manufacture. The same problem happens with lenses if we make their sides
from parts of spheres. The correct shape for a lens is a hyperbola, and
hyperbole are hard to manufacture. Spheres are so much easier.\FRAME{dtbpFU}{%
3.3304in}{1.8991in}{0pt}{\Qcb{{\protect\small Example of Spherical
Aberration for Lenses. Optical Ray Trace done in the OSLO\ Optical Design
Package. Rays drawn as green lines are paraxial and rays drawn as red lines
are not.}}}{}{Figure}{\special{language "Scientific Word";type
"GRAPHIC";maintain-aspect-ratio TRUE;display "USEDEF";valid_file "T";width
3.3304in;height 1.8991in;depth 0pt;original-width 3.2846in;original-height
1.8619in;cropleft "0";croptop "1";cropright "1";cropbottom "0";tempfilename
'SWMLMC0D.wmf';tempfile-properties "XPR";}}The last figure is a ray trace
done by a professional lens design program. The lens is clearly not thin.
The rays drawn as green lines are the paraxial rays and you can see that
they form a tight focus to the right of the lens. But the rays drawn as red
lines are not paraxial. The non-paraxial and paraxial rays focus at
different spots because we have the wrong shape for our lens. We remember
that spherical aberration was made famous as the main problem with the
Hubble Telescope.

We get another image problem or aberration with lenses. Remember that the
index of refraction is different for different wavelengths of light. If we
send white light into our lens, different colors will bend slightly
differently (think Snell's law). So we will get a different focus point for
different colors of light. In the next figure, blue light (drawn with blue
lines) makes an image at a shorter focal distance than red light (drawn with
red lines). This problem is called \emph{chromatic aberration}. \FRAME{dtbpF%
}{3.2353in}{2.1862in}{0pt}{}{}{Figure}{\special{language "Scientific
Word";type "GRAPHIC";maintain-aspect-ratio TRUE;display "USEDEF";valid_file
"T";width 3.2353in;height 2.1862in;depth 0pt;original-width
3.1903in;original-height 2.1465in;cropleft "0";croptop "1";cropright
"1";cropbottom "0";tempfilename 'SWMLRY0F.wmf';tempfile-properties "XPR";}}%
This will also cause blurry pictures because different parts of the light
will focus at different distances from the lens. Generally mirrors can't
have chromatic aberration because they don't refract the light. But Lenses
always will.

There are many aberrations that come from making lenses that are easy to
manufacture, but that are not perfect in some way. We won't study all these
in this class. If you are curious, we cover more of these aberrations in our
junior level optics class (PH375).

\subsection{Sign Convention}

As we have been working we used our sign convention but we need to add to
our sign convention table a second radius.

\begin{equation*}
\begin{tabular}{|l|l|l|}
\hline
\textbf{Quantity} & \textbf{Positive if} & \textbf{Negative if} \\ \hline
{\small Object location }$\left( d_{o}\right) $ & 
\begin{tabular}{l}
{\small Object is in front of surface} \\ 
\end{tabular}
& 
\begin{tabular}{l}
{\small Object is in back of surface} \\ 
{\small (virtual object)}%
\end{tabular}%
{\small \ } \\ \hline
{\small Image location }$\left( d_{i}\right) $ & 
\begin{tabular}{l}
{\small Image is in back of surface} \\ 
{\small (real image)}%
\end{tabular}
& 
\begin{tabular}{l}
{\small Image is in front of surface } \\ 
{\small (virtual image)}%
\end{tabular}
\\ \hline
{\small Image height }$\left( h^{\prime }\right) $ & 
\begin{tabular}{l}
{\small Image is upright}%
\end{tabular}
& 
\begin{tabular}{l}
{\small Image is inverted}%
\end{tabular}
\\ \hline
{\small Radius }$\left( R_{1}\text{ and }R_{2}\right) $ & 
\begin{tabular}{l}
{\small Center of curvature is in} \\ 
{\small back of surface}%
\end{tabular}%
{\small \ } & 
\begin{tabular}{l}
{\small Center of curvature is in} \\ 
{\small front of surface}%
\end{tabular}%
{\small \ } \\ \hline
Focal length $\left( f\right) $ & Converging lens & Diverging lens \\ \hline
\end{tabular}%
\end{equation*}

\subsection{Magnification for thin lenses}

We defined the magnification earlier as a comparison of the image height to
the object height. 
\begin{equation}
m=\frac{h_{i}}{h_{o}}
\end{equation}%
And we found we can put this into terms of $d_{o}$ and $d_{i}$ for mirrors.
But we will this work for lenses? \FRAME{dtbpF}{3.9773in}{1.9951in}{0pt}{}{}{%
Figure}{\special{language "Scientific Word";type
"GRAPHIC";maintain-aspect-ratio TRUE;display "USEDEF";valid_file "T";width
3.9773in;height 1.9951in;depth 0pt;original-width 5.0764in;original-height
2.5322in;cropleft "0";croptop "1";cropright "1";cropbottom "0";tempfilename
'SWML680C.wmf';tempfile-properties "XPR";}}From the diagram, we can see that 
\begin{eqnarray*}
\tan \theta &=&\frac{h_{o}}{d_{o}} \\
\tan \theta &=&\frac{h_{i}}{d_{i}}
\end{eqnarray*}%
then 
\begin{equation*}
\frac{h_{o}}{d_{o}}=\frac{h_{i}}{d_{i}}
\end{equation*}%
or 
\begin{equation*}
\frac{h_{i}}{h_{o}}=\frac{d_{i}}{d_{o}}
\end{equation*}%
And once again we can write the magnification for a converging lens as minus
the ratio of the image distance over the object distance. 
\begin{equation*}
m=-\frac{d_{i}}{d_{o}}
\end{equation*}%
Just like before the minus sign comes from our sign convention because the
image will be upside down.

We found the thin lens formula using converging lenses, but it works for
diverging lenses as well, so long as the thin lens approximation is valid.

\chapter{17 Eyes and Cameras 3.2.6, 3.2.7}

%TCIMACRO{%
%\TeXButton{\chaptermark{Cameras and Magnifiers}}{\chaptermark{Cameras and Magnifiers}}}%
%BeginExpansion
\chaptermark{Cameras and Magnifiers}%
%EndExpansion

in 1900 George Eastman introduced the Brownie Camera. This event has changed
society dramatically. The camera is very like an artificial eye. We will
study both eyes and cameras in this lecture.

%TCIMACRO{%
%\TeXButton{Fundamental Concepts}{\vspace{0.10in}\hspace{-0.37in}{\Large Fundamental Concepts\vspace{0.2in}}}}%
%BeginExpansion
\vspace{0.10in}\hspace{-0.37in}{\Large Fundamental Concepts\vspace{0.2in}}%
%EndExpansion

\begin{itemize}
\item The eye is like a semi-nfinate double lens with the cornea and the
crystalline lens as the two lenses

\item The crystalline lens does focusing but in biology it is called
accommodation

\item We have three color sensors for red green and blue light

\item Myopia and hyperopia are problems with the shape of the eye and are
corrected with additional lenses

\item Optical lenses are measured in optical power which is $1/f\left(
m\right) $ with the unit if a diopter.

\item Medical people measure the power of a lens in diopters. A diopter is
one over the focal length with the focal length measured in meters.

\item Cameras and other imaging systems use a strange term called and
f-number to tell what the intensity of the image will be%
\begin{equation*}
I\propto \frac{1}{\left( f/\#\right) ^{2}}
\end{equation*}%
where $f/\#$ is the symbol for f-number and the f-number is given by%
\begin{equation*}
f/\#\equiv \frac{f}{D}
\end{equation*}
\end{itemize}

\section{The Eye}

\FRAME{dhF}{3.48in}{1.3673in}{0pt}{}{}{Figure}{\special{language "Scientific
Word";type "GRAPHIC";maintain-aspect-ratio TRUE;display "USEDEF";valid_file
"T";width 3.48in;height 1.3673in;depth 0pt;original-width
3.4342in;original-height 1.3327in;cropleft "0";croptop "1";cropright
"1";cropbottom "0";tempfilename 'SVAJW82O.wmf';tempfile-properties "XPR";}}

The figure above shows the parts of the eye. The eye is like a camera in its
operation, but is much more complex. It is truly a marvel. The parts that
concern us are the cornea, crystalline lens, pupil, and the retina. \FRAME{%
dtbpF}{2.6913in}{1.8213in}{0in}{}{}{Figure}{\special{language "Scientific
Word";type "GRAPHIC";maintain-aspect-ratio TRUE;display "USEDEF";valid_file
"T";width 2.6913in;height 1.8213in;depth 0in;original-width
2.6481in;original-height 1.7841in;cropleft "0";croptop "1";cropright
"1";cropbottom "0";tempfilename 'SWQ1PZ01.wmf';tempfile-properties "XPR";}}%
The Cornea-lens system refracts the light onto the retina, which detects the
light. The lens is focused by a set of mussels that flatten the lens to
change it's focal length. The focusing process is different from a standard
camera. The camera moves the lens to achieve a different image distance. Our
eye can't change the distance between the lens system and the retina. So our
eye changes the shape of the lens, changing it's focal length. \FRAME{dhFU}{%
3.2387in}{1.446in}{0pt}{\Qcb{The crystaline lens becomes thicker, and
therefore more curved when the cilliary musscle flexes. Austin Flint,
\textquotedblleft The Eye as an Optical Instrument,\textquotedblright\ \emph{%
Popular Science Monthly,} Vol. 45, p203, 1894 (Image in the public domain)}}{%
}{Figure}{\special{language "Scientific Word";type
"GRAPHIC";maintain-aspect-ratio TRUE;display "USEDEF";valid_file "T";width
3.2387in;height 1.446in;depth 0pt;original-width 4.1987in;original-height
1.8602in;cropleft "0";croptop "1";cropright "1";cropbottom "0";tempfilename
'SVAJW82Q.wmf';tempfile-properties "XPR";}}

The focusing system is called \emph{accommodation}. This system becomes less
effective at about the time you reach an age of 40 years because the lens
becomes less flexible. The closest point that can be focused by
accommodation is called the near point. It is about $25\unit{cm}$ on
average. There is, of course, no such thing as an average person, all of us
are a little bit different. Young students probably have a much shorter near
point than $25\unit{cm}.$ For those of us that are a little older, $25\unit{%
cm}$ or more is more likely.

The farthest point that can be focused should be a long way away. It is
called the far point. Both the near and far points degrade with years
leading to bifocal glasses (and much irritation because you can't see as
well, but I'm really not complaining because at least I can see).

The iris changes the area of the pupil (the hole or \emph{aperture} of the
eye). The pupil is, on average, about $7\unit{mm}$ in diameter.

\subsection{Nearsightedness}

In some people the cornea-lens system focuses in front of the retina.
Usually this is because the shape of their eyes is not the normal spherical
shape. In this case it is elongated along the optic axis of the eye. This is
called nearsightedness or myopia.

\FRAME{dhF}{3.5475in}{2.4189in}{0pt}{}{}{Figure}{\special{language
"Scientific Word";type "GRAPHIC";maintain-aspect-ratio TRUE;display
"USEDEF";valid_file "T";width 3.5475in;height 2.4189in;depth
0pt;original-width 4.7547in;original-height 3.2335in;cropleft "0";croptop
"1";cropright "1";cropbottom "0";tempfilename
'SVAJW82R.wmf';tempfile-properties "XPR";}}So their perfectly good
cornea-lens system makes a great image in the vitreous humor (the jellylike
stuff that fills the eye) and not on the retina where it would be detected.
From your experience with lenses you know this would produce a blurry image.
And that is what nearsighted people see much of the time. We can help
nearsighted people by effectively changing the cornea-lens system of the eye
by adding another lens. We want a diverging lens that makes the light more
spread apart so that the lens system of the eye can make it focus where the
retina actually is. Alternately we could flatten the cornea, itself, with
laser ablation.

\subsection{Farsightedness}

Sometimes the cornea-lens system focuses in back of the retina. This is
usually because the eye grew too flat or \textquotedblleft
oblate.\textquotedblright\ This is called farsightedness or hyperopia. Once
again we can fix the problem by adding an additional lens. This time a
converging lens.

\FRAME{dhF}{3.7593in}{2.693in}{0pt}{}{}{Figure}{\special{language
"Scientific Word";type "GRAPHIC";maintain-aspect-ratio TRUE;display
"USEDEF";valid_file "T";width 3.7593in;height 2.693in;depth
0pt;original-width 3.7118in;original-height 2.6507in;cropleft "0";croptop
"1";cropright "1";cropbottom "0";tempfilename
'SVAJW82S.wmf';tempfile-properties "XPR";}}The converging lens will make the
effective focal length of the system shorter, so that it can form an image
where the retina actually is.

It would be convenient and very understandable if medical professionals
prescribed eye glasses by telling us what focal length we needed to correct
our vision. But that would not be the medical way! Eye glasses use a
different unit of measure to describe how they bend light. The unit is the 
\emph{diopter} and it is equal one over the focal length, but with the focal
length measured in meters. 
\begin{equation}
\text{diopter}=\frac{1}{f\left( \unit{m}\right) }
\end{equation}%
This measurement is called the \emph{power }of the lens. It is just the same
as giving the focal length, but less clear for science students.

\subsection{Color Perception}

The eye detects different colors. The respecters called cones can detect
red, green, and blue light.\FRAME{dhF}{2.4587in}{1.0326in}{0pt}{}{}{Figure}{%
\special{language "Scientific Word";type "GRAPHIC";maintain-aspect-ratio
TRUE;display "USEDEF";valid_file "T";width 2.4587in;height 1.0326in;depth
0pt;original-width 4.3656in;original-height 1.8178in;cropleft "0";croptop
"1";cropright "1";cropbottom "0";tempfilename
'SVAJW82T.wmf';tempfile-properties "XPR";}}The eye combines the red, green,
and blue response to allow us to perceive many different colors.

Most digital cameras also have red, green, and, blue pixels to provide color
to images. The detectors in digital cameras are often have much narrower
frequency bands than the eye. Likewise, television displays and monitors
have red, green, and blue pixels. By targeting the eye receptors, power need
not be wasted in creating light that is not detected well by the eye. The
difference in band-width can cause problems in color mixing. Yellow school
busses (perceived as different amounts of green and red light) may be
reddish or green if the bandwidths are chosen poorly.

The science of human visual perception of imagery is called \emph{image
science}. There are many applications for this field, from forensics to
intelligence gathering.

\section{The Camera}

The idea behind a camera is very simple.\FRAME{dtbpF}{3.0926in}{1.8542in}{0pt%
}{}{}{Figure}{\special{language "Scientific Word";type
"GRAPHIC";maintain-aspect-ratio TRUE;display "USEDEF";valid_file "T";width
3.0926in;height 1.8542in;depth 0pt;original-width 3.0485in;original-height
1.817in;cropleft "0";croptop "1";cropright "1";cropbottom "0";tempfilename
'SWXOK800.wmf';tempfile-properties "XPR";}}

The camera has a lens, often a compound lens made of several individual
lenses, and a screen for projecting a real image created by the lens. We
haven't learned how to deal with multiple lenses yet, but it's not too bad
of an approximation to treat the whole camera lens system as a single
equivalent lens with an equivalent focal length. Photographers speak as
though this were true and it's good enough for them to do their work, but
not good enough to do optical engineering. But we can gain a lot of insight
into how cameras work with this simple model. So let's stick with this
approximation in this lecture.

Let's take an example camera. Say we are at a family reunion and wish to
take a picture of Aunt Agatha. Aunt Agatha is about $1.\,\allowbreak 5\unit{m%
}$ tall. She is standing about $5\unit{m}$ away. Then to fit the image of
Aunt Agatha on our $35\unit{mm}$ detector, we must have

\begin{equation*}
\begin{tabular}{l}
$h_{o}=1.5\unit{m}$ \\ 
$h_{i}=0.035\unit{m}$ \\ 
$d_{o}=5\unit{m}$ \\ 
$f=0.058\unit{m}$%
\end{tabular}%
\end{equation*}

We wish to find $d_{i}$ and $m.$ Let's do $m$ first. 
\begin{eqnarray*}
m &=&\frac{h_{i}}{h_{o}}=\frac{-0.035\unit{m}}{1.5\unit{m}} \\
&=&-2.\,\allowbreak 333\,3\times 10^{-2}
\end{eqnarray*}%
so our image is small and inverted. The small size we wanted. But now we
know from the negative sign that the image in our cameras is upside down. A
digital camera uses it's built-in computer to turn the image right side up
for us on the display on the back of the camera.

Now let's find $d_{i}.$ 
\begin{eqnarray*}
d_{i} &=&\frac{fd_{o}}{d_{o}-f} \\
&=&5.\,\allowbreak 868\,1\times 10^{-2}\unit{m} \\
&=&58.681\unit{mm}
\end{eqnarray*}%
so our detector must be $58.681\unit{mm}$ from the lens.

Now suppose we want to photograph a $1000\unit{m}$ tower from $2\unit{km}$
away. Then%
\begin{eqnarray*}
m &=&-\frac{0.035\unit{m}}{1000\unit{m}} \\
&=&-3.\,\allowbreak 5\times 10^{-5}
\end{eqnarray*}%
and 
\begin{eqnarray*}
d_{i} &=&\frac{\left( 0.058\unit{m}\right) \left( 2000\unit{m}\right) }{2000%
\unit{m}-\left( 0.058\unit{m}\right) } \\
&=&5.\,\allowbreak 800\,2\times 10^{-2}\unit{m} \\
&=&58.002\unit{mm}
\end{eqnarray*}%
Notice that the image distance changed, but not by very much. This is why
you need a focus adjustment on the lens of a good camera. Objects far away
require a different $d_{i}$ value than objects that are close. Usually you
twist the lens housing to make this adjustment in manual focus mode. The
lens housing has a threaded screw system that increases or decreases $d_{i}$
as you twist. Cell phones and consumer cameras often have an motor that
makes this adjustment for you. In some cameras you may see the lens move
back and forth as someone takes a picture.

There are several things that govern whether a picture will be good. When
you buy a quality manual lens, it will be marked in $f/\#$'s. The
specification of an automatic lens will also be given in terms of $f/\#$. To
help us buy such lenses, we should understand what the terminology means.

Most things we want to take a snapshot of are much farther than $58\unit{mm}$
from the camera. For such objects we can revisit the magnification.%
\begin{equation*}
m=-\frac{d_{i}}{d_{o}}
\end{equation*}%
but from the thin lens formula%
\begin{equation*}
\frac{1}{d_{o}}+\frac{1}{d_{i}}=\frac{1}{f}
\end{equation*}%
If $d_{o}\gg f$ then we can say that $1/d_{o}\approx 0$ and so $d_{i}\approx
f.$ Then%
\begin{equation*}
m\approx -\frac{f}{d_{o}}
\end{equation*}%
and we see that the size of the image 
\begin{equation*}
h_{i}=mh_{o}
\end{equation*}%
is directly proportional to the focal distance. If we change the focal
distance, we can change the size of the image. This is how a zoom lens
works. A zoom lens is a compound lens, and the focal length is changed by
increasing the distance between the component lenses (more on compound
lenses soon). You may have a camera with a telephoto button and this button
zooms in and out by moving one part of the complex lens. Cell phones usually
do not do this to zoom in. They just make a photo out of the center of their
photosensor array and because they just use the middle, it looks like they
zoomed in. This is a \emph{digital zoom}. But let's go back to our standard
non-zoom camera lens.\footnote{%
If you are a photographer, we would call this a prime lens.}

Remember we studied intensity%
\begin{equation*}
I=\frac{P}{A}
\end{equation*}

Photographic film and digital focal plane arrays detect the intensity of
light that hits them. We can see that the area of our image depends on our
magnification, which depends on $d_{i}$ and for our distant objects it is
proportional to $f.$ The image area is proportional to $d_{i}^{2}\approx
f^{2}.$ So we can say that the area is proportional to $f^{2}.$ Then 
\begin{equation*}
I\propto \frac{P}{f^{2}}
\end{equation*}%
The power entering the camera is proportional to the size of the aperture
(hole the light goes through). A bigger aperture lets in more light. A
smaller aperture lets in less light. If the camera has a circular opening,
this area is proportional to the square of the diameter of the opening, $%
D^{2}$ so%
\begin{equation*}
I\propto \frac{D^{2}}{f^{2}}
\end{equation*}%
This ratio is useful because it gives us a relative measure of how much
intensity we get in terms of things we can easily know about our camera.
Good cameras have changeable aperture sizes, and many lenses have changeable
focal lengths. By using the combination of these two terms, we can ensure we
will get enough light (but not too much) when we take the picture.

It would be good to give this ratio a special name. But instead, we named
the ratio 
\begin{equation}
f/\#\equiv \frac{f}{D}
\end{equation}%
It is called the $f/\#$ (pronounced f-number, or f-stop) so our intensity is
proportional to one over the f-number. 
\begin{equation}
I\propto \frac{1}{\left( f/\#\right) ^{2}}
\end{equation}%
Good cameras have adjustable lens systems marked in $f/\#^{\prime }s$.
Typical values are $f/2.8,$ $f/4,$ $f/5.6,$ $f/8,$ $f/11,$ and $f/16.$
Notice that since $I$ is proportional to $\left( 1/f/\#\right) ^{2}$ these
common $f/\#$ values give an increase of intensity of a factor of $2$ each
time you change the $f/\#$ by one marked stop. That is because, say, 
\begin{equation*}
\frac{f/4}{f/2.8}=\frac{4}{2.8}=1.\,\allowbreak 428\,6\allowbreak =\sqrt{2}
\end{equation*}
and the intensity goes as $1/\left( f/\#\right) ^{2}$ so the change in
intensity from $f/2.8$ to $f/4$ would be proportional to $\left( \sqrt{2}%
\right) ^{2}=2.$ It is the same between any two of our standard $f/\#$
values.%
\begin{equation*}
\begin{tabular}{|c|c|}
\hline
$f/\#$ change & Intensity change factor \\ \hline
$f/11$ to $f/8$ & $2$ \\ \hline
$f/8$ to $f/5.6$ & $2$ \\ \hline
$f/5.6$ to $f/4$ & $2$ \\ \hline
$f/4$ to $f/2.8$ & $2$ \\ \hline
\end{tabular}%
\end{equation*}

This terminology is used for telescope design as well. The Hubble telescope
is an $f/24$ Ritchey-Chretien Cassegrainian system with a $2.4\unit{m}$
diameter aperture. The effective focal length is $57.6\unit{m}$.

\FRAME{dhF}{3.7593in}{2.3298in}{0pt}{}{}{Figure}{\special{language
"Scientific Word";type "GRAPHIC";maintain-aspect-ratio TRUE;display
"USEDEF";valid_file "T";width 3.7593in;height 2.3298in;depth
0pt;original-width 3.7118in;original-height 2.29in;cropleft "0";croptop
"1";cropright "1";cropbottom "0";tempfilename
'SWOHOQ00.wmf';tempfile-properties "XPR";}}

It is important to realize that electronic (and biological) sensors don't
react instantly to what we see. The intensity is 
\begin{equation*}
I=\frac{P}{A}=\frac{\Delta E}{\Delta tA}
\end{equation*}%
so there is a time involved. The time it takes to collect enough light to
form an image on the sensor is called the \emph{exposure time}. 
\begin{equation*}
\Delta t=\frac{\Delta E}{IA}
\end{equation*}%
So changing our $f/\#$ changes the needed exposure time by changing the
intensity. How sensitive our camera sensor is also affects the exposure
time. Modern sensors have adjustable sensitivity. The photography world
gives the three letters ISO for the name for this detector sensitivity.
There isn't a standard for exactly what ISO setting gives what exposure.
Different manufacturers use slightly different numbers. But a change in ISO
settings usually are equivalent to one $f/\#$ change in exposure.

This is part of what a good photographer does in taking a picture. The
photographer will adjust the $f/\#$ and the exposure time and ISO to get a
photograph that is not over exposed (too light) or underexposed (to dark).

\chapter{18 Optical Systems 3.2.8, 3.3 1}

%TCIMACRO{%
%\TeXButton{\chaptermark{Optical Systems}}{\chaptermark{Optical Systems}}}%
%BeginExpansion
\chaptermark{Optical Systems}%
%EndExpansion

So far we have looked at single lenses and how they work. But we hinted that
actual camera lenses and even our eyes are have lens systems that are made
of more than one lens. We should look at how lens systems work. First,
though, let's go back to the simple magnifier to define a more meaningful
magnification measure for using optical systems.

%TCIMACRO{%
%\TeXButton{Fundamental Concepts}{\vspace{0.10in}\hspace{-0.37in}{\Large Fundamental Concepts\vspace{0.2in}}}}%
%BeginExpansion
\vspace{0.10in}\hspace{-0.37in}{\Large Fundamental Concepts\vspace{0.2in}}%
%EndExpansion

\begin{itemize}
\item Angular magnification $M=\frac{\theta }{\theta _{o}}$ describes how
much bigger the image looks with an optical system compared to with just
your eye.

\item The magnification of a two-lens system is just the product of the
magnifications of the individual lenses $m_{\text{combined}}=m_{1}m_{2}$

\item If the two lenses in a two-lens system are placed so the distance
between them is essentially zero, then the focal length of the two-lens
system is given by 
\begin{equation*}
\frac{1}{f}=\frac{1}{f_{1}}+\frac{1}{f_{2}}
\end{equation*}%
where $f_{1}$ and $f_{2}$ are the focal lengths of the individual lenses

\item If the two lenses in a two-lens system are not placed so the distance
between them is essentially zero, the previous equation is not valid!

\item Telescopes and Microscopes are double lens systems
\end{itemize}

\section{Angular Magnification}

We already encountered the simple magnifier when we studied ray diagrams.
But by this point in our study of optics you are probably wondering about
our definition of magnification. If you are an Idahoan and are out hunting,
when you look through your binoculars or scope you don't want to know how
big the image is compared the actual size of the dear. No, you want to know
if you can see the deer better with the scope than you could with just your
eyes. The magnification on your scope doesn't compare the image size to the
object size. It is a comparison between what you see with two different
optical systems. One is just your eyes, the other is your eyes and the scope
working together.

Let's use the simple magnifier that we already know about (rather than a
complex scope) to define a new kind of magnification that does this
comparison between two optical systems. We can use what we know about rays
that are easy to draw, and about how our eyes work to describe both optical
systems.\FRAME{dtbpF}{3.3278in}{0.8717in}{0pt}{}{}{Figure}{\special{language
"Scientific Word";type "GRAPHIC";maintain-aspect-ratio TRUE;display
"USEDEF";valid_file "T";width 3.3278in;height 0.8717in;depth
0pt;original-width 5.2174in;original-height 1.3456in;cropleft "0";croptop
"1";cropright "1";cropbottom "0";tempfilename
'TFO01907.wmf';tempfile-properties "XPR";}}

We usually use three principle rays, but for this analysis, let's just use
one for each side of an object. Since the image needs to be real and on the
retina, we can see where these singular rays strike the retina and
understand the size of the image on our eye light detection system. Let's
choose to draw the rays that go straight through the middle of the lens of
the eye because they are the easiest ones, they don't bend, they just go
straight through. If we pick a ray from the top of our object that goes
through the center of the lens, that ray won't seem to change direction at
all. It will hit the retina to form the top of the image of the object.

We can do the same for the bottom of the object. Then we can see from the
next figure that the angle $\theta _{o}$ subtends\footnote{%
The angle that \textquotedblleft subtends\textquotedblright\ and object is
the one who's bounding lines just hit either side of the object.} both the
object and the image of the object. If you think about the previous figure,
you will see that if the angle were to increase, so would the size of the
image on the retina. And that would mean that we would see a bigger version
of the image and would be able to see more detail. We can increase this
angle by moving the object closer to our eyes.\FRAME{dtbpF}{3.6089in}{%
1.9484in}{0pt}{}{}{Figure}{\special{language "Scientific Word";type
"GRAPHIC";maintain-aspect-ratio TRUE;display "USEDEF";valid_file "T";width
3.6089in;height 1.9484in;depth 0pt;original-width 3.5613in;original-height
1.9104in;cropleft "0";croptop "1";cropright "1";cropbottom "0";tempfilename
'SWOI8M04.wmf';tempfile-properties "XPR";}}

If we move our object too close we will reach the limit of the eye for
focusing. \FRAME{dtbpF}{3.8882in}{1.0758in}{0pt}{}{}{Figure}{\special%
{language "Scientific Word";type "GRAPHIC";maintain-aspect-ratio
TRUE;display "USEDEF";valid_file "T";width 3.8882in;height 1.0758in;depth
0pt;original-width 4.8585in;original-height 1.324in;cropleft "0";croptop
"1";cropright "1";cropbottom "0";tempfilename
'SWOICM06.wmf';tempfile-properties "XPR";}}If we move the object any closer,
it will appear blurry. Remember, we call this position, the closest point
where we can place an object and still bring it into focus with our eye, the 
\emph{near point.} Thus the maximum value of $\theta _{o}$ where the image
won't be blurry will be at the near point.

But suppose we want to see this object in more detail. That is, we want to
spread the image of the object over more of our retina so our retinal pixels
do a better job of showing the details. We can use a magnifying glass to
spread the light over more of our retina. From our ray diagrams that we have
drawn for thin lenses we know that this is what a magnifying glass does, it
spreads out the rays. We will see that this spreading will make the rays
cover more of our retina. Let's go back to a normal ray diagram and draw two
principle rays for the top of the bug. \FRAME{dtbpF}{5.4846in}{2.565in}{0in}{%
}{}{Figure}{\special{language "Scientific Word";type
"GRAPHIC";maintain-aspect-ratio TRUE;display "USEDEF";valid_file "T";width
5.4846in;height 2.565in;depth 0in;original-width 5.4267in;original-height
2.5235in;cropleft "0";croptop "1";cropright "1";cropbottom "0";tempfilename
'TFT54R00.wmf';tempfile-properties "XPR";}} Sure enough, the two rays spread
out and never meet to form a real image, so we have a virtual image of the
bug. To achieve this we placed the object closer to the magnifying glass
than the focal distance $(d_{o}<f)$, then we have a virtual image with
magnification 
\begin{equation}
m=\frac{-d_{i}}{d_{o}}
\end{equation}%
which is larger than one and positive (because $d_{i}$ is negative by sign
convention). This is what we have studied so far, but you probably guessed
there was more to our virtual images than just this. After all, we do see
the image with our eye. So, we need to think about how our eye sees this
virtual image.

In the next figure I\ have departed from our three principle rays. We know
they were just representative rays, there are millions of actual rays. I
drew one ray that is one of the principle rays (the one that comes in
parallel to the axis) but I also drew a non-principle ray. Still, we see
that these two rays are spreading after the magnifying glass and we can
trace them backwards to find our virtual image. They never meet to form a
real image on their own. But now we introduce another optical system, our
eye! \FRAME{dtbpF}{4.8153in}{2.4448in}{0pt}{}{}{Figure}{\special{language
"Scientific Word";type "GRAPHIC";maintain-aspect-ratio TRUE;display
"USEDEF";valid_file "T";width 4.8153in;height 2.4448in;depth
0pt;original-width 5.1552in;original-height 2.6031in;cropleft "0";croptop
"1";cropright "1";cropbottom "0";tempfilename
'TFT5BQ01.wmf';tempfile-properties "XPR";}}The lens of our eye is another
positive lens and it can refract the rays back so they meet. Our eye
refracts the diverging rays to form a real image on our retina. So in the
end our virtual image from the lens became a real image on our retina, but
it took an additional lens to make this happen. In the last figure I just
drew rays for the top of the bug, but we would need another set of rays for
the bottom of the bug to see how big the real image at the retina became.

Back to magnification. What we really want to know is how much bigger the
image looks with the lens than it did without the lens. By looking at what
happens to the rays when they enter our eye, we now can see why the image
looks bigger. Let's go back to just a ray for the top of the bug and a ray
for the bottom of the bug and align our eye so these rays go through the
cornea as the undeviated rays. Then our diagram is simpler and we can see
why the magnifying lens makes a bigger image on our retina. \FRAME{dtbpF}{%
3.8337in}{3.6028in}{0in}{}{}{Figure}{\special{language "Scientific
Word";type "GRAPHIC";maintain-aspect-ratio TRUE;display "USEDEF";valid_file
"T";width 3.8337in;height 3.6028in;depth 0in;original-width
3.7853in;original-height 3.5544in;cropleft "0";croptop "1";cropright
"1";cropbottom "0";tempfilename 'TFRA3L00.wmf';tempfile-properties "XPR";}}

The magnifying glass has bent the light, and the bent rays make a bigger
angle, $\theta _{i}>\theta _{o}$ so the image on our retina is bigger. Since
the image fills more of our retina, we perceive the image as being bigger.
We need a mathematical formula that will tell us how much bigger the image
on our retina will be. We need to be careful with two parts of our optical
system. We have two images, one virtual and one real. Let's call the virtual
image $I_{1}$ with height $h_{i1}$ and the real image $I_{2}$ with height $%
h_{i2}.$

Notice from the figure that if we look inside our eye for the case of no
extra lens 
\begin{equation*}
\tan \theta _{o}\approx \frac{h_{i,eye}}{d}=\frac{h_{i}}{d}
\end{equation*}%
and for the case with the lens%
\begin{equation*}
\tan \theta _{i}\approx \frac{h_{i,lens+eye}}{d}=\frac{h_{i2}}{d}
\end{equation*}%
so 
\begin{eqnarray*}
h_{i} &=&h_{i,eye}\approx d\tan \theta _{o} \\
h_{i2} &=&h_{i,lens+eye}\approx d\tan \theta _{i}
\end{eqnarray*}%
Then if we compare the new, larger image on the retina formed with the
lens-eye system to the one formed with just the eye, we get 
\begin{equation*}
\frac{h_{i2}}{h_{i}}\approx \frac{d\tan \theta _{i}}{d\tan \theta _{o}}=%
\frac{\tan \theta _{i}}{\tan \theta _{o}}
\end{equation*}%
and if we once again use the small angle approximation 
\begin{equation*}
\frac{h_{i2}}{h_{i}}\approx \frac{\theta _{i}}{\theta _{o}}
\end{equation*}%
This would tell us how much bigger our object looks when viewed with the
magnifying glass compared to how it looked without the magnifying glass.
This is just what we want! Let's give this a new symbol 
\begin{equation}
M=\frac{\theta _{i}}{\theta _{o}}
\end{equation}

Remember, this is really different than the magnification we have studied
before. The magnification we have been using compared the size of the image
with the size of the object. The magnification $M$ compares the size of an
image made with just our eye to the size of the image made by our eye and
another lens system. We call $M$ the \emph{angular magnification} we can
keep these two types of magnification straight.

So, the angular magnification compares how big the object seems to be with
and without an additional lens or lens system. It is really a comparison
between the size of the area the diverging set of rays make on the retina
formed with just our eye, and the one formed with the magnifier.

We now know what angular magnification is, lets find an equation for the
angular magnification for our simple magnifier. \FRAME{dtbpF}{4.2929in}{%
2.9049in}{0pt}{}{}{Figure}{\special{language "Scientific Word";type
"GRAPHIC";maintain-aspect-ratio TRUE;display "USEDEF";valid_file "T";width
4.2929in;height 2.9049in;depth 0pt;original-width 5.0021in;original-height
3.3754in;cropleft "0";croptop "1";cropright "1";cropbottom "0";tempfilename
'TFPJG300.wmf';tempfile-properties "XPR";}}Let's say the image is not right
at the near point, and even make it harder by saying the lens is a distance $%
\ell $ away from the eye. This is really how we use a simple magnifier. Then%
\begin{equation*}
L=\ell +d_{i1}
\end{equation*}%
is how far the image is from the eye. And we can go back to trigonometry. 
\begin{equation*}
\tan \theta _{i}\approx \frac{h_{i1}}{L}\approx \theta _{i}
\end{equation*}%
and it is still true that without the lens to get the biggest view we need
to put the object at our near point. Let's choose this for our case without
the magnifier. 
\begin{equation*}
\tan \theta _{o}=\frac{h_{o}}{25\unit{cm}}\approx \theta _{o}
\end{equation*}%
Our angular magnification is 
\begin{equation*}
M=\frac{\theta _{i}}{\theta _{o}}\approx \frac{\frac{h_{i1}}{L}}{\frac{h_{o}%
}{25\unit{cm}}}=\frac{h_{i1}}{L}\frac{25\unit{cm}}{h_{o}}
\end{equation*}%
Once again we can use 
\begin{equation*}
m=\frac{h_{i1}}{h_{o}}=-\frac{d_{i1}}{d_{o1}}
\end{equation*}%
to write $M$ as 
\begin{equation*}
M\approx -\frac{d_{i1}}{d_{o1}}\frac{25\unit{cm}}{L}
\end{equation*}%
and we can eliminate $d_{o1}$ by solving for it and then substituting%
\begin{equation*}
\frac{1}{d_{o1}}+\frac{1}{d_{i1}}=\frac{1}{f_{1}}
\end{equation*}%
so 
\begin{equation*}
\frac{1}{d_{o1}}=\frac{1}{f_{1}}-\frac{1}{d_{i1}}
\end{equation*}%
and we actually want $1/d_{o1}$ so let's substitute now%
\begin{equation*}
M\approx -d_{i1}\left( \frac{1}{f_{1}}-\frac{1}{d_{i1}}\right) \frac{25\unit{%
cm}}{L}
\end{equation*}%
We need a common denominator%
\begin{eqnarray*}
M &\approx &\left( \frac{-d_{i1}}{f_{1}}-\frac{-d_{i1}}{d_{i1}}\right) \frac{%
25\unit{cm}}{L} \\
&=&\left( \frac{-d_{i1}}{f_{1}}+1\right) \frac{25\unit{cm}}{L} \\
&=&\left( 1+\frac{d_{i1}}{f_{1}}\right) \frac{25\unit{cm}}{L}
\end{eqnarray*}

Then recall that $-d_{i1}=L-\ell $ (the minus sign is because the image is
virtual so by sign convention $d_{i1}$ is negative) so 
\begin{equation*}
M\approx \left( 1+\frac{L-\ell }{f_{1}}\right) \frac{25\unit{cm}}{L}
\end{equation*}%
From this we can get the maximum value for our angular magnification for a
simple magnifier We just let $L=25\unit{cm}$ and we make $\ell \approx 0$
then the virtual image is as close as we can get it and have it not be
blurry. This gives 
\begin{eqnarray*}
M_{\max } &\approx &\left( 1+\frac{25\unit{cm}-0}{f_{1}}\right) \frac{25%
\unit{cm}}{25\unit{cm}} \\
&\approx &\left( 1+\frac{25\unit{cm}}{f_{1}}\right)
\end{eqnarray*}%
This is our maximum magnification (try it, put the magnifying glass right
next to your eye and see how big the image looks. then try other locations
for the lens).

We could also get a minimum angular magnification for our simple magnifier.
We can think of what we know about ray diagrams. If we put the object right
at the focal distance so $d_{o1}=f_{1}$ then the image would be at infinity
because the rays from the object would be bent into parallel rays by the
lens. And looking at our diagram that would mean $L=\infty .$ Of course that
is too far, we wouldn't actually see a virtual image. But suppose we only
let $L\rightarrow \infty $ (so $d_{o1}\rightarrow f_{1}$ but doesn't quite
equal $f_{1}$). Then we should get the minimum magnification 
\begin{eqnarray*}
M_{\min } &=&\underset{L\rightarrow \infty }{\lim }\left( 1+\frac{L-\ell }{%
f_{1}}\right) \frac{25\unit{cm}}{L} \\
&=&\underset{L\rightarrow \infty }{\lim }\left( 25\unit{cm}\left( \frac{1}{L}%
+\frac{L-\ell }{f_{1}L}\right) \right)
\end{eqnarray*}%
and as $L\rightarrow \infty $ little $\ell $ becomes unimportant and $1/L$
becomes zero so 
\begin{eqnarray*}
M_{\min } &=&\underset{L\rightarrow \infty }{\lim }\left( \left( 1+\frac{%
L-\ell }{f_{1}}\right) \frac{25\unit{cm}}{L}\right) \\
&=&\underset{L\rightarrow \infty }{\lim }\left( 25\unit{cm}\left( \frac{1}{L}%
+\frac{L}{Lf_{1}}\right) \right) \\
&=&25\unit{cm}\left( 0+\frac{1}{f_{1}}\right) \\
&=&\frac{25\unit{cm}}{f_{1}}
\end{eqnarray*}%
There were a lot of approximations, but still, we can see that we can have a
magnification approximately in the range 
\begin{equation*}
\frac{25\unit{cm}}{f_{1}}<M<1+\frac{25\unit{cm}}{f_{1}}
\end{equation*}

We use the idea of a simple magnifier in combination with other lenses to
make the magnification happen in telescopes, microscopes, and other
instruments that magnify. So all of these systems can be described using the
idea of an angular magnification.

Angular magnification is important. Many optical systems that we buy are
really designed to make things seem bigger (or smaller) than they seem with
just our eyes. Things like telescopes and microscopes and binoculars. We
will see how to apply this idea of angular magnification to some such
systems in our next lecture.

\section{Combinations of Lenses}

Our eyes and cameras have more than one lens. in using a magnifying glass
technically we have more than one lens because our eye is involved. We found
that to correct chromatic aberration we used two lenses. Together they are
called an \textquotedblleft achromat\textquotedblright\ and all good cameras
use achromats to fix chromatic aberration. We can do something similar to
correct for spherical aberration. But in each case, we are combining two
lenses (or even more!). Microscopes and telescopes are also systems that
have a minimum of two lenses. How do we predict what the lens system will
do? It turns out that we have all we need to know to combine lenses already.

To combine lenses, we do the same thing we did for the two surfaces of a
thin lens or for our simple magnifier case. We form the image from the first
lens as though the second lens is not there. Then we use the image from the
first lens as the object for the second lens. Suppose we take two lenses of
focal lengths $f_{1}$ and $f_{2}$ and place them a distance $d$ apart. 
\FRAME{dtbpF}{4.286in}{1.7876in}{0pt}{}{}{Figure}{\special{language
"Scientific Word";type "GRAPHIC";maintain-aspect-ratio TRUE;display
"USEDEF";valid_file "T";width 4.286in;height 1.7876in;depth
0pt;original-width 4.235in;original-height 1.7504in;cropleft "0";croptop
"1";cropright "1";cropbottom "0";tempfilename
'SX1IQB00.wmf';tempfile-properties "XPR";}}

Because this system would use a magnified image as the object for lens $2,$
the final magnification is the product of the two lens magnifications

\begin{equation}
m_{\text{combined}}=m_{1}m_{2}
\end{equation}

For the first lens we have%
\begin{equation}
\frac{1}{d_{o1}}+\frac{1}{d_{i1}}=\frac{1}{f_{1}}
\end{equation}%
where $d_{i1}$ is our first lens image distance. We can solve for $d_{i1}$%
\begin{equation}
d_{i1}=\frac{d_{o1}f_{1}}{d_{o1}-f_{1}}  \label{q}
\end{equation}%
This first image from the first lens will be the object for the second lens

We take as the second object distance as 
\begin{equation*}
d_{o2}=d-d_{i1}
\end{equation*}%
and we use the lens formula again.%
\begin{equation*}
\frac{1}{d_{o2}}+\frac{1}{d_{i2}}=\frac{1}{f_{2}}
\end{equation*}%
and again find the image distance%
\begin{equation*}
d_{i2}=\frac{d_{o2}f_{2}}{d_{o2}-f_{2}}
\end{equation*}%
but we can use our value of $d_{o2}$ to find%
\begin{eqnarray*}
d_{i2} &=&\frac{\left( d-d_{i1}\right) f_{2}}{\left( d-d_{i1}\right) -f_{2}}
\\
&=&\frac{\left( d-d_{i1}\right) f_{2}}{d-d_{i1}-f_{2}}
\end{eqnarray*}%
We have and expression relating the image distances, $d,$ $d_{i1,}$ $d_{i2}$
and $f_{2}.$ But we would really like to have an expression that relates $%
d_{o1}$ and $d_{i2}.$ That is, we want a relationship that by knowing how
far away the object is from the lens system we get out how far the final
image is from the lens system. Then we can treat the lens system like a
single lens (most of the time).

Lets use our expression for 
\begin{equation*}
d_{i1}=\frac{d_{o1}f_{1}}{d_{o1}-f_{1}}
\end{equation*}%
and substitute it into our expression for $d_{i2}$%
\begin{equation*}
d_{i2}=\frac{\left( d-\frac{d_{o1}f_{1}}{d_{o1}-f_{1}}\right) f_{2}}{d-\frac{%
d_{o1}f_{1}}{d_{o1}-f_{1}}-f_{2}}
\end{equation*}%
This looks messy, but it does the job. Knowing $d_{o1}$ we can find $d_{i2}$%
!. We can do some simplification%
\begin{equation}
d_{i2}=\frac{df_{2}-\frac{d_{o1}f_{1}f_{2}}{d_{o1}-f_{1}}}{d-f_{2}-\frac{%
d_{o1}f_{1}}{d_{o1}-f_{1}}}  \label{TwoLensesSeparated}
\end{equation}%
Well, it is still messy, but we have achieved our goal! We have $d_{i2}$ in
terms of the focal lengths $f_{1}$ and $f_{2}$, the distances $d,$ and $%
d_{o1}.$

Let's think of our achromat lens system where the lenses are glued together.
They are right next to each other. The distance between our lenses is $d.$
Suppose we let $d\rightarrow 0.$ That would mean the two lenses are right
next to each other, like our achromat. Then%
\begin{equation*}
d_{i2}=\frac{-\frac{d_{o1}f_{1}f_{2}}{d_{o1}-f_{1}}}{-f_{2}-\frac{d_{o1}f_{1}%
}{d_{o1}-f_{1}}}
\end{equation*}%
which hardly looks better. But let's rearrange some 
\begin{eqnarray*}
d_{i2} &=&\frac{\frac{d_{o1}f_{1}f_{2}}{d_{o1}-f_{1}}}{\frac{f_{2}\left(
d_{o1}-f_{1}\right) }{d_{o1}-f_{1}}+\frac{d_{o1}f_{1}}{d_{o1}-f_{1}}} \\
&=&\frac{d_{o1}f_{1}f_{2}}{f_{2}d_{o1}-f_{2}f_{1}+d_{o1}f_{1}} \\
&=&\frac{d_{o1}f_{1}f_{2}}{d_{o1}\left( f_{2}+f_{1}\right) -f_{2}f_{1}}
\end{eqnarray*}%
So%
\begin{equation*}
d_{i2}=\frac{d_{o1}f_{1}f_{2}}{d_{o1}\left( f_{2}+f_{1}\right) -f_{2}f_{1}}
\end{equation*}%
Lets undo the math that brought us $d_{i2}$ in the first place. I want to
make this look like the thin lens formula, but with $d_{i2}$ as the image
distance and $d_{o1}$ as the object distance. 
\begin{eqnarray*}
\frac{1}{d_{i2}} &=&\frac{d_{o1}\left( f_{2}+f_{1}\right) -f_{2}f_{1}}{%
d_{o1}f_{1}f_{2}} \\
&=&\frac{d_{o1}\left( f_{2}+f_{1}\right) }{d_{o1}f_{1}f_{2}}-\frac{f_{2}f_{1}%
}{d_{o1}f_{1}f_{2}} \\
&=&\frac{\left( f_{2}+f_{1}\right) }{f_{1}f_{2}}-\frac{1}{d_{o1}}
\end{eqnarray*}%
or%
\begin{equation*}
\frac{1}{d_{i2}}+\frac{1}{d_{o1}}=\frac{\left( f_{2}+f_{1}\right) }{%
f_{1}f_{2}}
\end{equation*}%
Which looks very like the lens formula with 
\begin{equation*}
\frac{1}{f}=\frac{\left( f_{2}+f_{1}\right) }{f_{1}f_{2}}
\end{equation*}%
If we unwind this expression for $f$, we find%
\begin{equation}
\frac{1}{f}=\frac{f_{2}}{f_{1}f_{2}}+\frac{f_{1}}{f_{1}f_{2}}  \notag
\end{equation}%
\begin{equation}
\frac{1}{f}=\frac{1}{f_{1}}+\frac{1}{f_{2}}
\end{equation}

We see that the two lenses are equivalent to a single lens with focal length 
$f$ as long as they are close together.

Of course, we had to place our lenses right next to each other for this to
work. This is not the case for a telescope or microscope. We should look at
such a case. There is no need for more math. We can go back to equation (\ref%
{TwoLensesSeparated}). 
\begin{equation*}
d_{i2}=\frac{df_{2}-\frac{d_{o1}f_{1}f_{2}}{d_{o1}-f_{1}}}{d-f_{2}-\frac{%
d_{o1}f_{1}}{d_{o1}-f_{1}}}
\end{equation*}

Let's look at a case using ray diagrams. For this case, let's take two
lenses, and let's have the first lens make a real image. Once again, let's
have that image be the object for the second lens. But this time, let's move
the second lens so that the image from the first lens (object for the second
lens) is closer to the second lens than $f_{2}.$ If that is the case, the
second lens works like a magnifier. The final image is enlarged.

\FRAME{dtbpF}{3.9409in}{2.5746in}{0pt}{}{}{Figure}{\special{language
"Scientific Word";type "GRAPHIC";maintain-aspect-ratio TRUE;display
"USEDEF";valid_file "T";width 3.9409in;height 2.5746in;depth
0pt;original-width 3.8917in;original-height 2.5339in;cropleft "0";croptop
"1";cropright "1";cropbottom "0";tempfilename
'SX1IQB01.wmf';tempfile-properties "XPR";}}This is the basis of both
microscopes and telescopes!

\section{The Microscope}

Microscopes are designed to allow us to see very small things. They go
beyond a simple magnifier. But it turns out we don't have to go very far
beyond a magnifying glass. To see things that are very small, we add another
lens to our simple magnifier. We will place this lens near the object and
call this lens the \emph{objective} because it is next to the object. We
will keep a simple magnifier and place it near the eye. This lens will be
called the \emph{eyepiece} because it is near your eye.

The objective will have a very short focal length. The eyepiece will have a
longer focal length (a few centimeters).\FRAME{dtbpF}{3.9409in}{2.8383in}{0in%
}{}{}{Figure}{\special{language "Scientific Word";type
"GRAPHIC";maintain-aspect-ratio TRUE;display "USEDEF";valid_file "T";width
3.9409in;height 2.8383in;depth 0in;original-width 3.8917in;original-height
2.7951in;cropleft "0";croptop "1";cropright "1";cropbottom "0";tempfilename
'SWOK380B.wmf';tempfile-properties "XPR";}}

We separate the lenses by a distance $L$ where 
\begin{eqnarray*}
L &>&f_{o} \\
L &>&f_{e}
\end{eqnarray*}

We place the object \emph{just outside} the focal point of the objective
lens. The image formed by the objective lens is then real and inverted. Then
(and here is the trick we just learned!) we use this real image formed by
the objective lens as the new object for the eyepiece. For a real image,
there is really light there. So a real image can act like light bouncing off
of an object. We can use the real image as an object for the eyepiece.

The image formed by the eyepiece is upright and virtual, but it looks upside
down because the object for the eyepiece (the image made by the objective)
is upside down.

To say it once more, what we are doing is forming a real image using the
objective lens and then using the eyepiece lens as a simple magnifier to
make a larger virtual image of the real image made by the objective lens. We
could find $d_{ie}$ using equation (\ref{TwoLensesSeparated}).

But this is a microscope. We want to know how much bigger our final image
is. And of course, we really want to know the magnification of the system.
It's not just the eyepiece that will magnify. We know what the eyepiece will
do because it is being used as just a simple magnifier. Recall that
magnifications are a factors. Think from our basic equation 
\begin{equation*}
m=\frac{h_{i}}{h_{o}}
\end{equation*}%
tells us that 
\begin{equation*}
h_{i}=mh_{o}
\end{equation*}%
or in words, $m$ is the factor by which $h_{i}$ is bigger (or smaller) than $%
h_{o}.$ So magnifications are factors. For example, $d_{i}$ could be $10$
time bigger than $d_{o}$. And we have two lenses. Each lens will have a
magnification of it's object. But how do we combine the magnifications of
the two lenses? We just multiply the two magnifications together! The
objective lens makes an image that is a factor of $m_{o}$ times bigger (or
smaller) than the object. Then the eyepiece make an image of the first image
that is $m_{e}$ times bigger than the first image. So the final image is $%
m_{o}m_{e}$ times bigger than the object.

We should ask if this works for angular magnification. To find the overall
system angular magnification we start with the magnification of the eyepiece 
$M_{e},$ say, $20$ times bigger than it would look with just our eye. So
whatever the object, the eyepiece gives us a virtual image of that object
that is $M_{e}$ times bigger than it would look without the eyepiece. But
the object in our case is the real image formed by the objective lens. And
this real image may be bigger than the actual object. If you think about it,
it makes sense that if the objective makes the object look $10$ times
bigger, and the eyepiece makes the image look $20$ times bigger than if you
looked at it with your eye, then the whole system makes it look $200$ times
bigger. The combined magnification for the two-lens system is 
\begin{equation}
M_{system}=m_{o}M_{e}
\end{equation}%
The minimum magnification of the eye piece will be roughly 
\begin{equation*}
M_{e_{\min }}\approx \frac{25\unit{cm}}{f_{e}}
\end{equation*}%
and the maximum will be 
\begin{equation*}
M_{e_{\max }}=1+\frac{25\unit{cm}}{f_{e}}
\end{equation*}%
But remember the object for the eyepiece is the image from the first lens.
And that image is larger than the object by an amount \ 
\begin{equation*}
m_{o}=\frac{-d_{io}}{d_{oo}}
\end{equation*}%
where $d_{io}$ is the image distance for the objective lens and $d_{oo}$ is
the object distance for the objective lens. And let's estimate how big the
magnification due to the first lens will be. Because for a microscope $%
d_{oo}\approx f_{o}.$ and $d_{io}\approx L$ (very roughly) 
\begin{equation*}
m_{o}=\frac{-d_{io}}{d_{oo}}\approx -\frac{L}{f_{o}}
\end{equation*}

The combined magnification for the two-lens system is about 
\begin{equation}
M_{system}=m_{o}M_{e}\approx -\frac{L}{f_{o}}\frac{25\unit{cm}}{f_{e}}
\end{equation}%
this is the minimum magnification (because we used the minimum magnification
formula for the eyepiece).

This gives us a rough idea of how a microscope works, let's try telescopes
next.

\section{Telescopes}

There are two main types of telescopes \emph{refracting} and \emph{reflecting%
}. We will study refracting telescopes first.

\textbf{Refracting Telescopes}

Like the microscope, we combine two lenses and call one the objective and
the other the eyepiece. The eyepiece again plays the role of a simple
magnifier, magnifying the image produced by the objective.\FRAME{dtbpF}{%
5.1084in}{2.9516in}{0pt}{}{}{Figure}{\special{language "Scientific
Word";type "GRAPHIC";maintain-aspect-ratio TRUE;display "USEDEF";valid_file
"T";width 5.1084in;height 2.9516in;depth 0pt;original-width
4.4996in;original-height 2.5884in;cropleft "0";croptop "1";cropright
"1";cropbottom "0";tempfilename 'SX3HE400.wmf';tempfile-properties "XPR";}}%
We again form a real, inverted image with the objective. We are now looking
at distant objects. This means $d_{oo}$ is very large. and using%
\begin{equation*}
d_{oi}=\frac{d_{oo}f_{o}}{d_{oo}-f_{o}}\approx \frac{d_{oo}f_{o}}{d_{oo}}%
=f_{o}
\end{equation*}%
we can see that the image distance $d_{io}\approx f_{o}.$ Once again, we use
the image from the objective as the object for the eyepiece. The eye piece
forms an upright virtual image (that looks inverted because the object for
the eyepiece is the image from the objective, and the real image from the
objective is inverted). Once again we could find $d_{ie}$ using equation (%
\ref{TwoLensesSeparated}).

Once again we want the magnification. The largest magnification is when the
rays exit the eyepiece parallel to the principal axis. Then the image from
the eyepiece is formed at near infinity (but it is very big, so it is easy
to see). This gives a lens separation of $f_{o}+f_{e}$ which will be roughly
the length of the telescope tube.

The angular magnification will be 
\begin{equation}
M=\frac{\theta _{i}}{\theta _{o}}
\end{equation}%
where $\theta _{o}$ is the angle subtended by the object at the objective
(see figure above) it's what you see if you stand in front of the telescope
and just use your eye. And $\theta _{i}$ is subtended by the final image at
the viewer's eye. Consider the height of the final image, $h_{o},$ is very
large. And let's call the height of the image made by the objective lens $%
h_{io}.$ We see from the figure that 
\begin{equation}
\tan \theta _{o}\approx -\frac{h_{io}}{f_{o}}
\end{equation}%
and with the distance from the objective lens to the object, $d_{oo},$ very
large we can use small angles. 
\begin{equation}
\theta _{o}\approx -\frac{h_{io}}{f_{o}}
\end{equation}

The angle $\theta _{i}$ will be the angle formed by rays bent by the lens of
the eyepiece. This angle can be found by following a ray that leaves the tip
of the first image and goes through the eyepiece and into the person's eye.
The easiest ray to chose is the one that goes through the middle of the
eyepiece. Since we placed the eyepiece so $d_{eo}\approx f_{e}$ 
\begin{equation}
\tan \theta _{i}\approx \frac{h_{io}}{f_{e}}\approx \theta _{i}
\end{equation}%
so

The magnification is nearly%
\begin{equation}
M=\frac{\theta _{i}}{\theta _{o}}\approx \frac{\frac{h_{io}}{f_{e}}}{-\frac{%
h_{io}}{f_{o}}}\approx -\frac{f_{o}}{f_{e}}
\end{equation}

\textbf{Reflecting Telescopes}

Reflecting telescopes use a series of mirrors to replace the objective lens.
Usually, there is an eyepiece that is refractive (although there need not
be, radio frequency telescopes rarely have refractive pieces).\FRAME{dhF}{%
3.6348in}{2.7354in}{0pt}{}{}{Figure}{\special{language "Scientific
Word";type "GRAPHIC";maintain-aspect-ratio TRUE;display "USEDEF";valid_file
"T";width 3.6348in;height 2.7354in;depth 0pt;original-width
3.5864in;original-height 2.693in;cropleft "0";croptop "1";cropright
"1";cropbottom "0";tempfilename 'SVAJW832.wmf';tempfile-properties "XPR";}}

There are two reasons to build reflective telescopes. The first is that
reflective optics do not suffer from chromatic aberration. The second is
that large mirrors are much easier to make and mount than refractive optics.
The Keck Observatory in Hawaii has a $10\unit{m}$ reflective system. The
largest refractive system is a $1\unit{m}$ system. Mechanical engineering
students will appreciate that putting giant $1\unit{m}$ diameter hunk of
glass on the end of a hollow tube and then turning the tube to point in
different directions would cause the tube and lens mass to oscillate. This
would make the imagery from the telescope very bad unless you wait for the
\textquotedblleft ringing\textquotedblright\ oscillations to die out due to
damping. Putting the heavy mass on the bottom of the tube is much more
stable.

The Hubble telescope has a $2.5\unit{m}$ aperture and James Webb has a $6%
\unit{m}$. Reflective telescopes are just easier to build.

The telescope pictured in the figure is a Newtonian, named after Newton, who
designed this type of telescope. Many other designs exist. Popular designs
for space applications include the Cassegrain telescope. The rough design of
a reflective telescope can be worked out using refractive pieces, then the
rough details of the reflective optics can be formed.

But let's take a break from designing optical systems like cameras and
telescopes and return to the fact that light is a wave. In our study of
mirrors and lenses so far, beyond our explanation for reflection and
refraction, we just didn't have much constructive and destructive
interference. But we know for waves constructive and destructive
interference must happen. Let's go looking for interference for light waves
next.

\chapter{19 Double Slits 3.3.2, 3.3.3 \label{DoubleSlitMath}}

%TCIMACRO{%
%\TeXButton{\chaptermark{Double Slits}}{\chaptermark{Double Slits}}}%
%BeginExpansion
\chaptermark{Double Slits}%
%EndExpansion

We know light is a wave and that waves experience constructive and
destructive interference. But so far in our study of light we have used ray
approximation and that didn't include wave effects. It did make the math
easy. But when you buy a camera system there must be some effect of the wave
nature of light that would affect your purchasing decision. And beyond that,
we got some very cool and practical ideas from using waves for sound. Is the
wave nature of light useful?

Let's look at some specific cases where we know light must show interference
phenomena. We can think of when we had two speakers making sound waves. We
found sound wave constructive and destructive interference. Could we set up
two light sources that would make two light waves and mix those two waves?
If we can, it seems like we should see interference!

%TCIMACRO{%
%\TeXButton{Fundamental Concepts}{\vspace{0.10in}\hspace{-0.37in}{\Large Fundamental Concepts\vspace{0.2in}}}}%
%BeginExpansion
\vspace{0.10in}\hspace{-0.37in}{\Large Fundamental Concepts\vspace{0.2in}}%
%EndExpansion

\begin{itemize}
\item Light waves experience constructive and destructive interference like
other waves

\item Young's double slit experiment

\item Visible light wavelengths are small, we can use the small angle
approximation

\item Multiple slits give areas of constructive and destructive interference.
\end{itemize}

\section{Interference and Young's Experiment}

Waves do some funny things when they encounter barriers. Think of a water
wave. If we pass the wave through a small opening in a barrier, the wave
can't all get through the small hole, but it can cause a disturbance right
at the opening. We know that a small disturbance will cause a wave. But this
wave will be due to a very small--almost point--source. Point sources make
spherical waves. So the waves on the other side of the small opening will be
nearly spherical. The smaller the opening the more pronounced the curving of
the wave, because the source (the hole) is more like a point source.\FRAME{%
dhF}{2.1378in}{2.0375in}{0pt}{}{}{Figure}{\special{language "Scientific
Word";type "GRAPHIC";maintain-aspect-ratio TRUE;display "USEDEF";valid_file
"T";width 2.1378in;height 2.0375in;depth 0pt;original-width
2.0989in;original-height 1.9986in;cropleft "0";croptop "1";cropright
"1";cropbottom "0";tempfilename 'TARYM602.wmf';tempfile-properties "XPR";}}%
Now suppose we have two of these openings. We expect the two sources to make
curved waves and those waves can interfere. \FRAME{dhF}{2.2139in}{2.5901in}{%
0pt}{}{}{Figure}{\special{language "Scientific Word";type
"GRAPHIC";maintain-aspect-ratio TRUE;display "USEDEF";valid_file "T";width
2.2139in;height 2.5901in;depth 0pt;original-width 2.962in;original-height
3.4696in;cropleft "0";croptop "1";cropright "1";cropbottom "0";tempfilename
'TARYM603.wmf';tempfile-properties "XPR";}}None of this is really new. This
is just like our two speakers making sound waves. We know waves generally
have a spherical shape, and we have a disturbance in the wave medium (the
electromagnetic field). The field in the two holes is experiencing simple
harmonic motion. This makes the holes become sources for new spherical waves.

In the figure, we can already see that there will be constructive and
destructive interference were the waves from the two holes meet. An early
researcher named Thomas Young predicted that we should see constructive and
destructive interference in light (he drew figures very like the ones we
have used to explain his idea). \FRAME{dtbpF}{2.3765in}{2.6844in}{0pt}{}{}{%
Figure}{\special{language "Scientific Word";type
"GRAPHIC";maintain-aspect-ratio TRUE;display "USEDEF";valid_file "T";width
2.3765in;height 2.6844in;depth 0pt;original-width 2.3367in;original-height
2.642in;cropleft "0";croptop "1";cropright "1";cropbottom "0";tempfilename
'TARYM604.wmf';tempfile-properties "XPR";}}

Young set up a coherent source of light and placed it in front of this
source a barrier with two very thin slits cut in it to test his idea.. He
set up a screen beyond the barrier and observed the pattern on the screen
formed by the light. This something like what he saw

\bigskip

\FRAME{dhF}{2.3618in}{0.7671in}{0pt}{}{}{Figure}{\special{language
"Scientific Word";type "GRAPHIC";maintain-aspect-ratio TRUE;display
"USEDEF";valid_file "T";width 2.3618in;height 0.7671in;depth
0pt;original-width 2.322in;original-height 0.7351in;cropleft "0";croptop
"1";cropright "1";cropbottom "0";tempfilename
'TARYTF06.wmf';tempfile-properties "XPR";}}We see bright spots (constructive
interference) and dark spots (destructive interference). Only wave phenomena
can interfere, so this is fairly good evidence that light is a wave.

\subsection{Constructive Interference}

We can find the condition for getting a bright or a dark spot if we think
about it a bit. The two waves coming from the two slits are just two
different waves. We can use our two wave mixing analysis with our
constructive and destructive interference criteria. For constructive
interference, the difference in phase must be a multiple of $2\pi .$ Let's
review this, but instead of $y_{\max }$ being the displacement of the waves,
let's write them as $E_{\max }$ because it is the electric field that is
carrying the wave. And the electric field doesn't really go up or down, it
just gets a higher value or lower value. Our waves will be 
\begin{eqnarray*}
E_{1} &=&E_{\max }\sin \left( k_{2}r_{2}-\omega _{2}t_{2}+\phi _{2}\right) \\
E_{2} &=&E_{\max }\sin \left( k_{1}r_{1}-\omega _{1}t_{1}+\phi _{1}\right)
\end{eqnarray*}%
and the resulting wave will be 
\begin{eqnarray*}
E_{r} &=&E_{\max }\sin \left( k_{2}r_{2}-\omega _{2}t_{2}+\phi _{2}\right)
+A\sin \left( k_{1}r_{1}-\omega _{1}t_{1}+\phi _{1}\right) \\
&=&2E_{\max }\cos \left( \frac{\left( k_{2}r_{2}-\omega _{2}t_{2}+\phi
_{2}\right) -\left( k_{1}r_{1}-\omega _{1}t_{1}+\phi _{1}\right) }{2}\right)
\\
&&\times \sin \left( \frac{\left( k_{2}r_{2}-\omega _{2}t_{2}+\phi
_{2}\right) +\left( k_{1}r_{1}-\omega _{1}t_{1}+\phi _{1}\right) }{2}\right)
\\
&=&2E_{\max }\cos \left( \frac{1}{2}\left[ \left( k_{2}r_{2}-\omega
_{2}t_{2}+\phi _{2}\right) -\left( k_{1}r_{1}-\omega _{1}t_{1}+\phi
_{1}\right) \right] \right) \\
&&\times \sin \left( \frac{\left( k_{2}r_{2}-\omega _{2}t_{2}+\phi
_{2}\right) +\left( k_{1}r_{1}-\omega _{1}t_{1}+\phi _{1}\right) }{2}\right)
\end{eqnarray*}%
As we are now well aware, the sine part is a combined wave, and the cosine
part is part of the amplitude. The amplitude can be written as 
\begin{equation*}
A=2E_{\max }\cos \left( \frac{1}{2}\left[ \left( k_{2}r_{2}-\omega
_{2}t_{2}+\phi _{2}\right) -\left( k_{1}r_{1}-\omega _{1}t_{1}+\phi
_{1}\right) \right] \right)
\end{equation*}%
We should pause and think about which of our values will change. The light
was originally part of the same incoming wave. So, the frequency won't
change, $\omega _{2}=\omega _{1}=\omega .$ There is no slower material, so
the wavelength won't change, $k_{2}=k_{1}=k.$ We want the waves to mix at
the same time so $t_{2}=t_{1}=t.$ And the new waves are created from the old
wave hitting the slits. As long as the original wave hits both slits at
once, then $\phi _{2}=\phi _{1}=\phi _{o}.$ We are left with 
\begin{eqnarray*}
A &=&2E_{\max }\cos \left( \frac{1}{2}\left[ \left( kr_{2}-\omega t+\phi
_{o}\right) -\left( kr_{1}-\omega t+\phi _{0}\right) \right] \right) \\
&=&2E_{\max }\cos \left( \frac{1}{2}\left[ \left( kr_{2}\right) -\left(
kr_{1}\right) \right] \right) \\
&=&2E_{\max }\cos \left( \frac{1}{2}k\left( r_{2}-r_{1}\right) \right) \\
&=&2E_{\max }\cos \left( \frac{1}{2}k\Delta r\right)
\end{eqnarray*}

That means that to have constructive interference the path difference
between the two slit-sources must be an even number of wavelengths. We have
been calling the path difference in the total phase $\Delta x,$ or for
spherical waves $\Delta r,$ but in optics it is customary to call this path
difference $\delta $ or even $\Delta \ell .$ But since from Principles of
Physics I we would write $r_{2}-r_{1}$ as $\Delta r,$ let's stick with that.%
\begin{equation*}
\Delta r=r_{2}-r_{1}
\end{equation*}%
Let's write the amplitude function as 
\begin{equation*}
A=2E_{\max }\cos \left( \frac{1}{2}k\Delta r\right)
\end{equation*}%
and it is the amplitude that tells us if we have constructive or destructive
interference. Let's do what we did for geometric optics problems and mark up
our double slit situation so we can use some geometry to find our
constructive and destructive interference locations.\FRAME{dtbpF}{4.8793in}{%
2.3647in}{0pt}{}{}{Figure}{\special{language "Scientific Word";type
"GRAPHIC";maintain-aspect-ratio TRUE;display "USEDEF";valid_file "T";width
4.8793in;height 2.3647in;depth 0pt;original-width 4.944in;original-height
2.3815in;cropleft "0";croptop "1";cropright "1";cropbottom "0";tempfilename
'TARXYM01.wmf';tempfile-properties "XPR";}}We can define a change in the
total phase 
\begin{equation*}
\Delta \phi =\left( kr_{2}-\omega t+\phi _{o}\right) -\left( kr_{1}-\omega
t+\phi _{0}\right)
\end{equation*}%
And in our case since $\omega _{2}=\omega _{1}$ and $k_{2}=k_{1}$ and$\phi
_{2}=\phi _{1}$ 
\begin{equation*}
\Delta \phi =k\Delta r
\end{equation*}%
and we can write our amplitude as 
\begin{equation*}
A=2E_{\max }\cos \left( \frac{1}{2}\Delta \phi \right)
\end{equation*}%
We want to use this in our criteria for constructive or destructive
interference. Think for a minute about cosine functions. For cosine $\cos
\left( \theta \right) =\pm 1$ when $\theta =0,$ $\pi ,$ $2\pi ,\cdots $ and
when the cosine in our amplitude is equal to $\pm 1$ our amplitude will be
as big as it can get 
\begin{equation*}
A_{\max }=2E_{\max }\left( \pm 1\right) =\pm 2E_{\max }
\end{equation*}%
So, this is constructive interference. Our $\theta $ in $\cos \left( \theta
\right) $ for our case is $\frac{1}{2}\Delta \phi $ so we are looking for
when 
\begin{equation*}
\theta =\frac{1}{2}\Delta \phi =0,\pi ,2\pi ,\cdots
\end{equation*}%
We could write this as 
\begin{equation*}
\frac{1}{2}\Delta \phi =m\pi \qquad m=0,1,2,3\cdots
\end{equation*}%
where $m$ is an integer. And we could solve for $\Delta \phi $ 
\begin{equation*}
\Delta \phi =2m\pi \qquad m=0,1,2,3\cdots
\end{equation*}%
and we know that $\Delta \phi =k\Delta r$ so to have constructive
interference we have a condition on our $\Delta r$%
\begin{equation*}
k\Delta r=2m\pi \qquad m=0,1,2,3\cdots
\end{equation*}%
\begin{equation*}
\frac{2\pi }{\lambda }\Delta r=2m\pi \qquad m=0,1,2,3\cdots
\end{equation*}%
\begin{equation*}
\Delta r=m\lambda \qquad m=0,1,2,3\cdots
\end{equation*}%
If we are off by a any whole number of wavelengths we have constructive
interference. This is just the same as it was for sound waves!

But let's express $\Delta r$ in terms of geometry that is easy to measure in
our experiment. In this set up, the screen is much farther away than $d,$
the slit distance. We can say that the blue triangle is almost a right
triangle, and then $\Delta r$ is 
\begin{equation*}
\Delta r=r_{2}-r_{1}\approx d\sin \theta
\end{equation*}%
Our wave still repeats every $2\pi $ radians or every wavelength, $\lambda ,$
then we have constructive interference (a bright spot) when%
\begin{equation*}
\Delta \phi =k\Delta r\approx kd\sin \theta =2\pi m\quad \left( m=0,\pm
1,\pm 2\ldots \right)
\end{equation*}%
where $m$ is still an integer. We can solve again for $\Delta r$%
\begin{eqnarray*}
\Delta r &=&d\sin \theta =\frac{2\pi m}{k}\quad \left( m=0,\pm 1,\pm 2\ldots
\right) \\
&=&\frac{2\pi m}{\frac{2\pi }{\lambda }}\quad \left( m=0,\pm 1,\pm 2\ldots
\right)
\end{eqnarray*}%
or 
\begin{equation*}
\Delta r=d\sin \theta =m\lambda \quad \left( m=0,\pm 1,\pm 2\ldots \right)
\end{equation*}%
And we can give $m$ a name. It is called the \emph{order number}. What we
have found is that if we are off by any number of whole wavelengths then our
total phase due to path difference will be a multiple of $2\pi $ and we will
have constructive interference.

If we assume that $\lambda \ll d$ we can find the distance from the axis for
each fringe more easily. This guarantees that $\theta $ will be small. Using
the yellow triangle we see%
\begin{equation*}
\tan \theta =\frac{y}{L}
\end{equation*}%
but if $\theta $ is small this is just about the same as 
\begin{equation*}
\sin \theta \approx \frac{y}{L}
\end{equation*}%
because for small angles $\tan \theta \approx \sin \theta \approx \theta .$
So if $\theta $ is small then 
\begin{eqnarray*}
\Delta r &=&d\sin \theta \\
&\approx &d\frac{y}{L}
\end{eqnarray*}%
and for a bright spot or \emph{\textquotedblleft fringe\textquotedblright }
we find 
\begin{equation*}
d\frac{y}{L}=m\lambda
\end{equation*}%
Solving for the position of the bright spots gives%
\begin{equation}
y_{bright}\approx \frac{\lambda L}{d}m\quad \left( m=0,\pm 1,\pm 2\ldots
\right)
\end{equation}%
We can measure up from the central spot and predict where each successive
bright spot will be.

\subsection{Destructive Interference}

We can also find a condition for destructive interference. We know that a
path difference of an odd multiple of a half wavelength will give
destructive interference. But it is harder to think about how to write an
odd multiple of half wavelengths. Let's go back to our amplitude equation 
\begin{equation*}
A=2E_{\max }\cos \left( \frac{1}{2}\Delta \phi \right)
\end{equation*}
but this time we want when $\cos \left( \frac{1}{2}\Delta \phi \right) =0.$
That gives destructive interference. And we know that $\cos \left( \theta
\right) =0$ when $\theta =\frac{\pi }{2},\frac{3\pi }{2},$ $\frac{5\pi }{2}%
,\cdots .$ We can need a way to write this series of numbers that make $\cos
\left( \theta \right) =0.$ Let's try $\left( m+\frac{1}{2}\right) \pi $
where once again $m=0,\pm 1,\pm 2\ldots $ We can stick in $m$ values to see
if it works. 
\begin{equation*}
\begin{tabular}{ll}
$m=-2$ & $\left( -2+\frac{1}{2}\right) \pi =\allowbreak -\frac{3}{2}\pi $ \\ 
$m=-1$ & $\left( -1+\frac{1}{2}\right) \pi =\allowbreak -\frac{1}{2}\pi $ \\ 
$m=0$ & $\left( 0+\frac{1}{2}\right) \pi =\allowbreak \frac{1}{2}\pi $ \\ 
$m=1$ & $\left( 1+\frac{1}{2}\right) \pi =\allowbreak \frac{3}{2}\pi $ \\ 
$m=2$ & $\left( 2+\frac{1}{2}\right) \pi =\allowbreak \frac{5}{2}\pi $ \\ 
$m=3$ & $\left( 3+\frac{1}{2}\right) \pi =\allowbreak \frac{7}{2}\pi $%
\end{tabular}%
\end{equation*}%
This seems to work! Then let's put this into our amplitude function. Notice
that the amplitude has $\frac{1}{2}\Delta \phi $ and not just $\Delta \phi .$
I want to find the values of $\Delta \phi $ that will work to make
destructive interference. We have to account for the extra factor of $1/2.$
You should convince yourself that 
\begin{equation*}
\Delta \phi =\left( m+\frac{1}{2}\right) 2\pi \qquad m=0,\pm 1,\pm 2,\pm
3,\cdots
\end{equation*}%
will work. Using this we can solve once again for $\Delta r$ 
\begin{equation*}
\Delta \phi =k\Delta r\approx kd\sin \theta =\left( m+\frac{1}{2}\right)
2\pi \qquad m=0,\pm 1,\pm 2,\pm 3,\cdots
\end{equation*}%
or using $k=2\pi /\lambda $ 
\begin{equation*}
\Delta r=d\sin \theta =\left( m+\frac{1}{2}\right) \lambda \quad \left(
m=0,\pm 1,\pm 2\ldots \right)
\end{equation*}%
This will give destructive interference. Will give a dark fringe for these
values of $\Delta r$. Using the same small angle approximation as we did
before, the location of the dark fringes will be 
\begin{equation}
y_{dark}\approx \frac{\lambda L}{d}\left( m+\frac{1}{2}\right) \quad \left(
m=0,\pm 1,\pm 2\ldots \right)
\end{equation}

Let's try an example problem.

Suppose we have a laser beam that produces pretty red light with $\lambda
=632.8\unit{nm}$. And suppose we shine this laser on an opaque screen with
two slits that are $0.020\unit{mm}$ apart. And suppose we let the light hit
a screen that is $5.00\unit{m}$ away from the slits. How far apart will the
bright interference fringes be on the screen?

\FRAME{dtbpF}{2.9031in}{1.0112in}{0pt}{}{}{Figure}{\special{language
"Scientific Word";type "GRAPHIC";maintain-aspect-ratio TRUE;display
"USEDEF";valid_file "T";width 2.9031in;height 1.0112in;depth
0pt;original-width 2.9297in;original-height 1.0023in;cropleft "0";croptop
"1";cropright "1";cropbottom "0";tempfilename
'SVNOVQ02.wmf';tempfile-properties "XPR";}}PT: This is a two slit problem

BE%
\begin{equation*}
y_{bright}=\frac{\lambda L}{d}m\qquad m=0,\pm 1,\pm 2,\cdots
\end{equation*}%
\begin{equation}
y_{dark}\approx \frac{\lambda L}{d}\left( m+\frac{1}{2}\right) \quad \left(
m=0,\pm 1,\pm 2\ldots \right)
\end{equation}

VAR%
\begin{eqnarray*}
\lambda &=&632.8\unit{nm} \\
d &=&0.020\unit{mm} \\
L &=&5.00\unit{m}
\end{eqnarray*}%
Solution

We want the space between two bright fringes. $\Delta y,$ this is 
\begin{eqnarray*}
\Delta y &=&\frac{\lambda L}{d}\left( m+1\right) -\frac{\lambda L}{d}\left(
m\right) \\
&=&\frac{\lambda L}{d}\left( \left( m+1\right) -\left( m\right) \right) \\
&=&\frac{\lambda L}{d}
\end{eqnarray*}%
or%
\begin{equation*}
\Delta y=\frac{\lambda L}{d}
\end{equation*}%
\begin{eqnarray*}
\Delta y &=&\frac{\left( 632.8\unit{nm}\right) \left( 5.00\unit{m}\right) }{%
0.020\unit{mm}} \\
&&0.158\,2\unit{m} \\
&=&15.82\unit{cm}
\end{eqnarray*}

Units check

Reasonableness check

This may seem like a mostly useless thing to measure, but let's think of
this problem backwards. If we have two slits and measure a pattern of $15.82%
\unit{cm}$ peak separation we can know that the slits were $0.200\unit{mm}$
apart? Light interference is used in industry for measuring very small
distances on the order of tens of microns.

\section{Multiple Slits}

What happens if we have more than two slits? \FRAME{dtbpF}{2.3765in}{2.6844in%
}{0in}{}{}{Figure}{\special{language "Scientific Word";type
"GRAPHIC";maintain-aspect-ratio TRUE;display "USEDEF";valid_file "T";width
2.3765in;height 2.6844in;depth 0in;original-width 2.3367in;original-height
2.642in;cropleft "0";croptop "1";cropright "1";cropbottom "0";tempfilename
'SWZSIH0A.wmf';tempfile-properties "XPR";}}

At some point, two of the slits will have a path difference that is a whole
wavelength, and we would expect a bright spot. But what about the other
slits? If we have a slit spacing such that each of the succeeding slits has
a path difference that is just an additional wavelength, then each of the
slits will contribute to the constructive interference at our point, and the
point will become a bright spot.

\FRAME{dtbpF}{1.1554in}{2.8063in}{0pt}{}{}{Figure}{\special{language
"Scientific Word";type "GRAPHIC";maintain-aspect-ratio TRUE;display
"USEDEF";valid_file "T";width 1.1554in;height 2.8063in;depth
0pt;original-width 0.8536in;original-height 2.1136in;cropleft "0";croptop
"1";cropright "1";cropbottom "0";tempfilename
'SWOM8S0K.wmf';tempfile-properties "XPR";}}Let's look at just two slits
again. The light leaves each slit in phase with the light from the rest of
the slits, but at some distance $L$ away and at some angle $\theta $ we will
have a path difference%
\begin{equation}
\Delta r=d\sin \left( \theta _{bright}\right) =m\lambda \quad m=0,\pm 1,\pm
2,\ldots
\end{equation}%
because the path lengths are not all the same. But we can see that if we
have more than two slits, say, four slits, then this same equation would
work to find a bright spot for the top two slits. And it would also work for
the bottom two slits. And it would work for the middle two slits! Then for
every set of two slits we get constructive interference at the same angle.
That means at that spot at that angle all the constructive interference from
sets of two slits add up to make a super constructive interference. This
will be a very bright spot. We call such a bright spot a \emph{principle
maxima}.

Of course we could find a spot where two of the slits were $\lambda /2$ off
so there will be destructive interference from these two slits. But then a
third slit would still send light at that angle. So there would be some
brightness. So this isn't total destructive interference. There will be a
spot less bright than the principle maxima, but not totally dark. We call
these secondary maxima.

In the next figure, you can see laser light that passes through increasing
numbers of slits.

\FRAME{dtbpF}{4.364in}{3.375in}{0pt}{}{}{Figure}{\special{language
"Scientific Word";type "GRAPHIC";maintain-aspect-ratio TRUE;display
"USEDEF";valid_file "T";width 4.364in;height 3.375in;depth
0pt;original-width 7.5145in;original-height 5.8053in;cropleft "0";croptop
"1";cropright "1";cropbottom "0";tempfilename
'SVNLSS00.wmf';tempfile-properties "XPR";}}Notice that as the number of
slits grows, the principle maxima get brighter until they turn white because
they are oversaturating the camera sensor. These patters of light that come
from passing light through holes and letting the light mix together again
are called \emph{diffraction patterns}.

In our next lecture we see if we can find practical uses for light
constructive and destructive interference.

\chapter{20 Thin Films and Light Interference 3.3.4, 3.3.5}

%TCIMACRO{\TeXButton{\chaptermark{Thin Films}}{\chaptermark{Thin Films}}}%
%BeginExpansion
\chaptermark{Thin Films}%
%EndExpansion

Last lecture we reminded ourselves that light is a wave. In this lecture
let's look at some examples of how the wave nature of light can be used.
Applications for the using light waves stretch from radar to digital
communications to atmospheric chemical analysis. But let's take on the ideas
of interference from thin films and from different path lengths.

%TCIMACRO{%
%\TeXButton{Fundamental Concepts}{\vspace{0.10in}\hspace{-0.37in}{\Large Fundamental Concepts\vspace{0.2in}}}}%
%BeginExpansion
\vspace{0.10in}\hspace{-0.37in}{\Large Fundamental Concepts\vspace{0.2in}}%
%EndExpansion

\begin{itemize}
\item Thin films can cause reflections that create interference

\item There may be a phase shift when a wave reflects (the wave may
invert-this is a review)

\item If two waves are out of phase by half a wavelength we have total
destructive interference (review)

\item If two waves are out of phase by a full wavelength we have total
constructive interference (review)

\item Other phases provide partial constructive or partial destructive
interference
\end{itemize}

\section{Mathematical treatment of single frequency interference}

Let's review our equations for interference of two waves. We start with two
waves in the same medium. Only now our waves are light waves. So the medium
is the electromagnetic field. So we can say our amplitude is $E_{\max }$%
\begin{eqnarray*}
E_{1} &=&E_{\max }\sin \left( kx_{1}-\omega t+\phi _{1}\right) \\
E_{2} &=&E_{\max }\sin \left( kx_{2}-\omega t+\phi _{2}\right)
\end{eqnarray*}%
Each wave has its own phase constant. Each wave starts from a different
position (one at $x_{1}$ and the other at $x_{2}$), The superposition yields.%
\begin{equation*}
E_{r}=E_{\max }\sin \left( kx_{1}-\omega t+\phi _{1}\right) +E_{\max }\sin
\left( kx_{2}-\omega t+\phi _{2}\right)
\end{equation*}

which is graphed in the next figure.\FRAME{dtbpFX}{3.0727in}{1.3292in}{0pt}{%
}{}{Plot}{\special{language "Scientific Word";type "MAPLEPLOT";width
3.0727in;height 1.3292in;depth 0pt;display "USEDEF";plot_snapshots
TRUE;mustRecompute FALSE;lastEngine "MuPAD";xmin "0";xmax
"5.001000";xviewmin "0";xviewmax "5.001000";yviewmin "-2";yviewmax
"2";viewset"XY";rangeset"X";plottype 4;labeloverrides 2;y-label
"Er";axesFont "Times New Roman,12,0000000000,useDefault,normal";numpoints
100;plotstyle "patch";axesstyle "normal";axestips FALSE;xis
\TEXUX{x};var1name \TEXUX{$x$};function \TEXUX{$\sin \left( \frac{1}{6}\pi
+\pi x\right) +\sin \pi x$};linecolor "green";linestyle 1;pointstyle
"point";linethickness 2;lineAttributes "Solid";var1range
"0,5.001000";num-x-gridlines 100;curveColor
"[flat::RGB:0x00008000]";curveStyle "Line";VCamFile
'SWOOS614.xvz';valid_file "T";tempfilename
'SWOOS620.wmf';tempfile-properties "XPR";}}Notice that the wave form is
taller (larger amplitude). Noticed it is shifted along the $x$ axis. This
graph is not surprising to us now, because we have done a case like this
before. We can find the shift in general rewriting $y_{r}.$ We need a trig
identity%
\begin{equation*}
\sin a+\sin b=2\cos \left( \frac{a-b}{2}\right) \sin \left( \frac{a+b}{2}%
\right)
\end{equation*}%
then let $a=kx_{2}-\omega t+\phi _{2}$ and $b=kx_{1}-\omega t+\phi _{1}$%
\begin{eqnarray*}
E_{r} &=&E_{\max }\sin \left( kx_{2}-\omega t+\phi _{2}\right) +E_{\max
}\sin \left( kx_{1}-\omega t+\phi _{1}\right) \\
&=&2E_{\max }\cos \left( \frac{\left( kx_{2}-\omega t+\phi _{2}\right)
-\left( kx_{1}-\omega t+\phi _{1}\right) }{2}\right) \sin \left( \frac{%
\left( kx_{2}-\omega t+\phi _{2}\right) +\left( kx_{1}-\omega t+\phi
_{1}\right) }{2}\right) \\
&=&2E_{\max }\cos \left( \frac{kx_{2}-kx_{1}}{2}+\frac{\phi _{2}-\phi _{1}}{2%
}\right) \sin \left( \frac{kx_{2}+kx_{1}-2\omega t+\phi _{2}+\phi _{1}}{2}%
\right) \\
&=&2E_{\max }\cos \left( k\frac{x_{2}-x_{1}}{2}+\frac{\phi _{2}-\phi _{1}}{2}%
\right) \sin \left( k\frac{x_{2}+x_{1}}{2}-\omega t+\frac{\phi _{2}+\phi _{1}%
}{2}\right) \\
&=&2E_{\max }\cos \left( k\frac{\Delta x}{2}+\frac{\Delta \phi _{o}}{2}%
\right) \sin \left( k\frac{x_{2}+x_{1}}{2}-\omega t+\frac{\phi _{2}+\phi _{1}%
}{2}\right) \\
&=&2E_{\max }\cos \left( \frac{1}{2}\left( \frac{2\pi }{\lambda }\Delta
x+\Delta \phi _{o}\right) \right) \sin \left( k\frac{x_{2}+x_{1}}{2}-\omega
t+\frac{\phi _{2}+\phi _{1}}{2}\right) \\
&=&2E_{\max }\cos \left( \frac{1}{2}\left( \Delta \phi \right) \right) \sin
\left( k\frac{x_{2}+x_{1}}{2}-\omega t+\frac{\phi _{2}+\phi _{1}}{2}\right)
\end{eqnarray*}%
where the last line is just a rearrangement to match the form we got last
time. As usual, let's look at the sine part first. It is still a wave. We
can see that more clearly if we define a variable 
\begin{equation*}
\mathsf{x}=\frac{x_{2}+x_{1}}{2}
\end{equation*}%
and another 
\begin{equation*}
\overline{\phi }=\frac{\phi _{2}+\phi _{1}}{2}
\end{equation*}%
Then the sine part is just%
\begin{equation*}
\sin \left( k\mathsf{x}-\omega t+\overline{\phi }\right)
\end{equation*}%
and we can see it is just our basic wave equation form. The rest of the
equation must be the amplitude and we can clearly see that the amplitude
depends on what we call the phase difference, 
\begin{equation*}
\Delta \phi =\frac{2\pi }{\lambda }\left( \Delta x\right) +\Delta \phi _{o}
\end{equation*}%
And from this we can see that for a single frequency situation there can be
two sources two sources of phase difference. One can be from the two waves
traveling different paths and then combining $\left( \Delta x\right) $ and
the other is from the two waves starting with a different phase to begin
with, $\Delta \phi _{o}.$ If the total phase difference between the two
waves is a multiple of $2\pi ,$ then the two waves will experience
constructive interference%
\begin{equation*}
\Delta \phi =m2\pi \qquad m=0,\pm 1,\pm 2,\pm 3,\cdots
\end{equation*}%
Let's see that this works. Our amplitude is 
\begin{equation*}
A=2E_{\max }\cos \left( \frac{1}{2}\left( \Delta \phi \right) \right)
\end{equation*}%
and if we look at a cosine function we see that $\cos \left( \theta \right) $
is either $1$ or $-1$ at $\theta =n\pi .$\FRAME{dtbpFX}{4.4996in}{0.7913in}{%
0pt}{}{}{Plot}{\special{language "Scientific Word";type "MAPLEPLOT";width
4.4996in;height 0.7913in;depth 0pt;display "USEDEF";plot_snapshots
TRUE;mustRecompute FALSE;lastEngine "MuPAD";xmin "-5";xmax "5";xviewmin
"-5.0010000010002";xviewmax "5.0010000010002";yviewmin
"-1.00019994985428";yviewmax "1";plottype 4;labeloverrides 3;x-label "x
(units of pi)";y-label "cos(x)";axesFont "Times New
Roman,12,0000000000,useDefault,normal";numpoints 100;plotstyle
"patch";axesstyle "normal";axestips FALSE;xis \TEXUX{x};var1name
\TEXUX{$x$};function \TEXUX{$\cos \left( \pi x\right) $};linecolor
"black";linestyle 1;pointstyle "point";linethickness 1;lineAttributes
"Solid";var1range "-5,5";num-x-gridlines 100;curveColor
"[flat::RGB:0000000000]";curveStyle "Line";VCamFile
'SWOOTJ15.xvz';valid_file "T";tempfilename
'SWOOTJ21.wmf';tempfile-properties "XPR";}}So if $\Delta \phi =m2\pi $ then
the amplitude is 
\begin{eqnarray*}
A &=&2E_{\max }\left( \frac{1}{2}\left( m2\pi \right) \right) \\
&=&2E_{\max }\cos \left( m\pi \right) \\
&=&\pm 2E_{\max }
\end{eqnarray*}%
We don't really care if the amplitude function is big positively or
negatively. So we get constructive interference for $\cos \left( m\pi
\right) $ being either $1$ or $-1$. Then, this is our case for constructive
interference.%
\begin{equation*}
\Delta \phi =\left( \frac{2\pi }{\lambda }\Delta x+\Delta \phi _{o}\right)
=m2\pi \qquad m=0,\pm 1,\pm 2,\pm 3,\cdots
\end{equation*}

How about for destructive interference? We start again with our amplitude
function%
\begin{equation*}
A=2E_{\max }\cos \left( \frac{1}{2}\left( \Delta \phi \right) \right)
\end{equation*}%
but now we want when the cosine part needs to be zero. 
\begin{equation*}
\cos \left( \theta \right) =0
\end{equation*}

Looking at our cosine graph again, that happens for cosine when $\theta =%
\frac{\pi }{2}$, $\frac{3\pi }{2},$ $\frac{5\pi }{2},$ $\cdots .$ We could
write this as $\theta =\left( m+\frac{1}{2}\right) \pi $ for $m=0,\pm 1,\pm
2,\pm 3,\cdots .$ But remember that in our amplitude function, we already
have the $1/2$ in the function, so we want $\Delta \phi $ to have just the
odd integer multiple of $\pi $. We could write this as 
\begin{equation*}
\Delta \phi =\left( 2m+1\right) \pi \qquad m=0,\pm 1,\pm 2,\pm 3,\cdots 
\text{ }
\end{equation*}%
So our condition for destructive interference is 
\begin{equation*}
\Delta \phi =\left( \frac{2\pi }{\lambda }\Delta x+\Delta \phi _{o}\right)
=\left( 2m+1\right) \pi \qquad m=0,\pm 1,\pm 2,\pm 3,\cdots \text{ }
\end{equation*}%
We have developed a useful matched set of equations that will tell us if we
mix two waves when we will have constructive and destructive interference:

\begin{eqnarray*}
\Delta \phi &=&\left( \frac{2\pi }{\lambda }\Delta r+\Delta \phi _{o}\right)
=m2\pi \qquad m=0,\pm 1,\pm 2,\pm 3,\cdots \text{\qquad Constructive} \\
\Delta \phi &=&\left( \frac{2\pi }{\lambda }\Delta r+\Delta \phi _{o}\right)
=\left( 2m+1\right) \pi \qquad m=0,\pm 1,\pm 2,\pm 3,\cdots \text{\qquad\
Destructive}
\end{eqnarray*}

Let's take an example to see how this can be used.

\subsection{Example of two wave interference: Stealth Fighter}

The U.S. Air Force F-117 Nighthawk Stealth Fighter is an aircraft that is
designed to be mostly invisible to radar. How could this work? \FRAME{dtbpF}{%
2.3549in}{1.5783in}{0pt}{}{}{Figure}{\special{language "Scientific
Word";type "GRAPHIC";maintain-aspect-ratio TRUE;display "USEDEF";valid_file
"T";width 2.3549in;height 1.5783in;depth 0pt;original-width
3.4333in;original-height 2.2917in;cropleft "0";croptop "1";cropright
"1";cropbottom "0";tempfilename 'SXAOG500.wmf';tempfile-properties "XPR";}}

The answer is that I don't know\footnote{%
I don't really know if my guess about how the F-117 works is true. And if I
did, I wouldn't be able to use this as an example! So we are reverse
engineering the stealth fighter technology. But I think it is a good guess.}
. But we can make some guesses. The Stealth Fighter is kept in temperature
and humidity controlled hangers. \FRAME{dtbpF}{2.546in}{1.6959in}{0pt}{}{}{%
Figure}{\special{language "Scientific Word";type
"GRAPHIC";maintain-aspect-ratio TRUE;display "USEDEF";valid_file "T";width
2.546in;height 1.6959in;depth 0pt;original-width 3.6305in;original-height
2.4085in;cropleft "0";croptop "1";cropright "1";cropbottom "0";tempfilename
'SXAOKW01.wmf';tempfile-properties "XPR";}}So maybe there is a coating on
the aircraft that needs protecting so the sun and dryness don't make the
coating peal off of the plane. Let's see if we could make an airplane
\textquotedblleft stealthy\textquotedblright\ by applying a coating. \FRAME{%
dhF}{4.6034in}{2.2053in}{0pt}{}{}{Figure}{\special{language "Scientific
Word";type "GRAPHIC";maintain-aspect-ratio TRUE;display "USEDEF";valid_file
"T";width 4.6034in;height 2.2053in;depth 0pt;original-width
5.8254in;original-height 2.7752in;cropleft "0";croptop "1";cropright
"1";cropbottom "0";tempfilename
'../../../../PH223/Repository/PH223-Engineering-Physics/Text/SWOO921X.wmf';tempfile-properties "XPR";}%
}

So suppose the stealth fighter is coated with an anti-reflective polymer
that is part of it's mechanism for making the plane invisible to radar.
Suppose we have a radar system with a wavelength of $3.00\unit{cm}.$ Further
suppose that the index of refraction of the anti-reflective polymer is $%
n=1.50$, and that the aircraft index of refraction is very large, how thick
would you make the coating?

We want destructive interference, so let's start with our destructive
interference condition 
\begin{equation*}
\Delta \phi =\left( \frac{2\pi }{\lambda }\Delta r+\Delta \phi _{o}\right)
=\left( 2m+1\right) \pi \qquad m=0,\pm 1,\pm 2,\pm 3,\cdots \text{\qquad
Destructive}
\end{equation*}

The radar waves all hit the plane in phase. From the figure, we see that the
radar wave will reflect off of the coating. Because the index of refraction
of the coating is large, this is like a fixed end of a rope. There will be
an inversion.

But some of the wave will penetrate the polymer. This will reflect off of
the plane body. The plane body has a very large index of refraction, so once
again the wave will experience an inversion. The outgoing waves would then
both be in phase as they leave and create constructive interference (if
there were not a path difference) because 
\begin{equation*}
\Delta \phi _{o}=\pi -\pi =0
\end{equation*}%
at this point. Thus 
\begin{equation*}
\Delta \phi =\left( \frac{2\pi }{\lambda }\Delta r\right) =\left(
2m+1\right) \pi \qquad m=0,\pm 1,\pm 2,\pm 3,\cdots \text{\qquad Destructive}
\end{equation*}%
But $\Delta \phi _{o}$ is not the only term in our equation, we also have to
remember the path difference, $\Delta r$. The part of the wave that entered
the polymer travels farther. If that path difference, $\Delta r,$ is just
right we can get destructive interference. For the $m=0$ case 
\begin{equation*}
\Delta \phi =\left( \frac{2\pi }{\lambda }\Delta r\right) =\left( 2\left(
0\right) +1\right) \pi =\pi
\end{equation*}%
then the amplitude function would be 
\begin{eqnarray*}
A &=&2E_{\max }\cos \left( \frac{2\pi }{\lambda }\Delta r\right) \\
&=&2E_{\max }\cos \left( \frac{1}{2}\left( \pi \right) \right) \\
&=&0
\end{eqnarray*}%
and we have destructive interference. Note that these are electromagnetic
waves, so instead of $y_{\max }$ we have used $E_{\max }$ as the individual
wave amplitude. But the important thing is that the plane cannot be seen by
the radar! Of course, this works for $m=1$ and $m=2,$ etc. as well. Any odd
multiple of $\pi $ will work. We can write this as 
\begin{equation*}
\frac{2\pi }{\lambda }\Delta r=\left( 2m+1\right) \pi \qquad m=0,1,2,\cdots
\end{equation*}%
so that we are guaranteed an odd multiple of $\pi .$ This is our condition
for destructive interference.

But we are interested in the thickness. We realize that $\Delta r$ is about
twice the thickness, since the wave travels though the coating and back out.
So let's let $\Delta r\approx 2t$ 
\begin{equation*}
\frac{2\pi }{\lambda }2t\approx \left( 2m+1\right) \pi
\end{equation*}%
\begin{equation*}
2t\approx \left( 2m+1\right) \frac{\lambda }{2}
\end{equation*}%
\begin{equation*}
2t\approx \left( m+\frac{1}{2}\right) \lambda
\end{equation*}%
\begin{equation*}
t\approx \left( m+\frac{1}{2}\right) \frac{\lambda }{2}
\end{equation*}

But there is a further complication. We should write our thickness equation
as 
\begin{equation*}
t\approx \left( m+\frac{1}{2}\right) \frac{\lambda _{in}}{2}
\end{equation*}%
because $\Delta r$ has to provide an odd integer times the wavelength \emph{%
inside the coating} for the phase to be right. After all, the wave is
traveling inside the coating. We know that the wavelength will change as we
enter the slower material.

To see this, consider two waves traveling to the right in the figure below.
One passes through a slower material. We expect the wavelength to shorten.
We can see that, depending on the thickness $t,$ the wave may be in phase or
out of phase with the other wave as it leaves. In the figure, the thickness
is just right so that we have destructive interference. \FRAME{dtbpF}{%
3.7585in}{1.4226in}{0pt}{}{}{Figure}{\special{language "Scientific
Word";type "GRAPHIC";maintain-aspect-ratio TRUE;display "USEDEF";valid_file
"T";width 3.7585in;height 1.4226in;depth 0pt;original-width
3.2206in;original-height 1.2012in;cropleft "0";croptop "1";cropright
"1";cropbottom "0";tempfilename
'../../../../PH223/Repository/PH223-Engineering-Physics/Text/SWOO921Y.wmf';tempfile-properties "XPR";}%
} We have such a wavelength shift in the stealth fighter coating. But we
don't know the wavelength inside the coating. All we know is the radar
wavelength, $\lambda _{out}$. We can fix it by writing the wavelength inside
in terms of the wavelength outside. Earlier in our studies we found that the
new wavelength will be given by equation (\ref{WavelengthChange})%
\begin{equation*}
\lambda _{f}=\frac{v_{f}}{v_{i}}\lambda _{i}
\end{equation*}%
Let's rewrite this for our case%
\begin{equation*}
\lambda _{in}=\frac{v_{in}}{v_{out}}\lambda _{out}
\end{equation*}%
We can express this in terms of the index of refraction%
\begin{equation*}
n=\frac{c}{v}
\end{equation*}%
by multiplying the left hand side by $c/c$ then 
\begin{equation*}
\lambda _{in}=\frac{cv_{in}}{cv_{out}}\lambda _{out}
\end{equation*}%
or%
\begin{eqnarray*}
\lambda _{in} &=&\frac{\frac{c}{v_{out}}}{\frac{c}{v_{in}}}\lambda _{out} \\
&=&\frac{n_{out}}{n_{in}}\lambda _{out}
\end{eqnarray*}%
in the case of our aircraft coating the outside medium is air so $%
n_{out}\approx 1$ 
\begin{equation*}
\lambda _{in}=\frac{1}{n_{in}}\lambda _{out}
\end{equation*}%
This is this wavelength we need to match as the radar signal enters the
medium.

Using this expression for $\lambda _{in}$ in 
\begin{equation*}
t\approx \left( m+\frac{1}{2}\right) \frac{\lambda _{in}}{2}\qquad
m=0,1,2,\cdots
\end{equation*}%
will give us the condition for destructive interference. Let's rewrite our $%
\lambda _{in}$ equation for our case of a coating and air 
\begin{equation*}
\lambda _{in}=\lambda _{coating}=\frac{1}{n_{in}}\lambda _{out}=\frac{1}{%
n_{coating}}\lambda _{air}
\end{equation*}%
thus 
\begin{equation*}
t\approx \left( m+\frac{1}{2}\right) \frac{1}{2}\left( \frac{\lambda _{air}}{%
n_{coating}}\right) \qquad m=0,1,2,\cdots
\end{equation*}%
is our condition for being stealthy.

Let's assume we want the thinnest coating possible, so we set $m=0.$ Then 
\begin{equation*}
t\approx \left( \frac{1}{4}\right) \left( \frac{\lambda _{air}}{n_{coating}}%
\right)
\end{equation*}%
and our thickness would be 
\begin{equation*}
t\approx \left( \frac{1}{4}\right) \left( \frac{3.00\unit{cm}}{1.50}\right)
=0.5\unit{cm}
\end{equation*}%
This seems doable for an aircraft coating!

Of course we could also make a plane that would be more visible to radar by
choosing the constructive interference case. Suppose we are building a
search and rescue plane. We want to enhance it's ability to be seen by radar
in fog. We start with the condition for constructive interference 
\begin{equation*}
\Delta \phi =\left( \frac{2\pi }{\lambda }\Delta r+\Delta \phi _{o}\right)
=m2\pi \qquad m=0,\pm 1,\pm 2,\pm 3,\cdots \text{\qquad Constructive}
\end{equation*}%
It will still be true that $\Delta \phi _{o}=0.$ 
\begin{equation*}
\Delta \phi =\left( \frac{2\pi }{\lambda }\Delta r\right) =m2\pi
\end{equation*}%
and it is still true that $\Delta r\approx 2t.$ 
\begin{equation*}
\left( \frac{2\pi }{\lambda }2t\right) \approx m2\pi
\end{equation*}%
then 
\begin{equation*}
\left( \frac{1}{\lambda }2t\right) \approx m
\end{equation*}%
\begin{equation*}
t\approx \frac{1}{2}m\lambda
\end{equation*}%
and we still have to adjust for the coating index of refraction 
\begin{equation*}
t\approx \frac{m}{2}\left( \frac{\lambda _{air}}{n_{coating}}\right)
\end{equation*}%
And once again we have several choices for $m$ 
\begin{equation*}
t\approx \frac{m}{2}\left( \frac{\lambda _{air}}{n_{coating}}\right) \qquad
m=0,1,2,\cdots
\end{equation*}%
But now the coating will provide constructive interference, making it easier
to track on radar from the command center. For the thinnest possibility, set 
$m=1$ because the $m=0$ case doesn't give us any thickness.%
\begin{eqnarray*}
t &\approx &m\frac{1}{2}\left( \frac{\lambda _{air}}{n_{coating}}\right) \\
&=&\frac{1}{2}\left( \frac{3.00\unit{cm}}{1.50}\right) \\
&=&\allowbreak 1\unit{cm}
\end{eqnarray*}

Note that we reasoned out these equations for the boundary conditions that
we have in our problem (inversion on reflection from both the coating and
the plane body). If the boundary conditions change, so do the equations.

\subsection{Example of two wave interference: Soap Bubble}

Take a soap bubble for a second example.\FRAME{dtbpFU}{3.5293in}{3.141in}{0pt%
}{\Qcb{{\protect\small Interference from a soap bubble. (Bubble image in the
Public Domain, courtesy Marcin Der\k{e}gowski)}}}{}{Figure}{\special%
{language "Scientific Word";type "GRAPHIC";maintain-aspect-ratio
TRUE;display "USEDEF";valid_file "T";width 3.5293in;height 3.141in;depth
0pt;original-width 3.5674in;original-height 3.1718in;cropleft "0";croptop
"1";cropright "1";cropbottom "0";tempfilename
'../../../../PH223/Repository/PH223-Engineering-Physics/Text/SWOO921Z.wmf';tempfile-properties "XPR";}%
}Now we have a phase shift on the first reflection because going from air to
bubble material is like a fixed end. But we won't have an inversion on the
reflection from the inside surface of the bubble. This is because the bubble
is full of air. The index of refraction of air is less than that for the
bubble material. So as we leave the bubble material it is more like having a
free end of a rope. As the waves leave the surface, they are half a
wavelength out of phase due to $\Delta \phi _{o}$ because of the single
inversion from the bubble outer surface. We would have destructive
interference due to just this, but we also have to account for the bubble
thickness. If this thickness is a multiple of a wavelength, then we are
still have half a wavelength out of phase and we have destructive
interference.

Here are our basic equations%
\begin{eqnarray*}
\Delta \phi &=&\left( \frac{2\pi }{\lambda }\Delta r+\Delta \phi _{o}\right)
=m2\pi \qquad m=0,\pm 1,\pm 2,\pm 3,\cdots \text{ Constructive} \\
\Delta \phi &=&\left( \frac{2\pi }{\lambda }\Delta r+\Delta \phi _{o}\right)
=\left( 2m+1\right) \pi \qquad m=0,\pm 1,\pm 2,\pm 3,\cdots \text{
Destructive}
\end{eqnarray*}%
Suppose we want want constructive interference to get our colors, so we take
the first

\begin{equation*}
\Delta \phi =\left( \frac{2\pi }{\lambda }\Delta r+\Delta \phi _{o}\right)
=m2\pi
\end{equation*}%
and this time we have 
\begin{eqnarray*}
\Delta \phi _{o} &=&\phi _{transmitted}-\phi _{reflected} \\
&=&0-\pi \\
&=&-\pi
\end{eqnarray*}%
It is still true that $\Delta r\approx 2t$ so from our constructive
interference equation%
\begin{equation*}
\frac{2\pi }{\lambda }2t-\pi =m2\pi
\end{equation*}%
\begin{equation*}
\frac{2}{\lambda }t-\frac{1}{2}=m
\end{equation*}%
\begin{equation*}
\frac{2}{\lambda }t=m+\frac{1}{2}
\end{equation*}%
\begin{equation*}
t=\frac{\lambda }{2}\left( m+\frac{1}{2}\right)
\end{equation*}%
We again have the problem that this wavelength must be the wavelength inside
the bubble material $\lambda =\lambda _{in}.$ But we see the outside
wavelength $\lambda _{out}$. We can reuse our conversion from outside to
inside wavelength from our last problem because we are once again in air and 
$n_{air}\approx 1$.%
\begin{equation*}
\lambda _{in}=\frac{1}{n_{in}}\lambda _{out}
\end{equation*}%
then%
\begin{equation*}
t=\frac{\lambda _{out}}{2n_{in}}\left( m+\frac{1}{2}\right) \qquad
m=0,1,2,\cdots
\end{equation*}%
Or writing this with $n_{in}=n_{\text{bubble }}$to make it clear that the
inside material is the bubble solution,%
\begin{equation*}
t=\left( m+\frac{1}{2}\right) \frac{1}{2}\left( \frac{\lambda _{air}}{n_{%
\text{bubble}}}\right) \qquad m=0,1,2,\cdots
\end{equation*}%
But this was the equation for destructive interference for the plane! We can
see that memorizing the thickness equations won't work. We need to start
with our conditions on $\Delta \phi $ for constructive and destructive
interference and solve for the thickness to be safe!

How about the dark parts of the bubble with no color (the parts we can see
through). These would be destructive interference%
\begin{equation*}
\Delta \phi =\left( \frac{2\pi }{\lambda }\Delta x+\Delta \phi _{o}\right)
=\left( 2m+1\right) \pi
\end{equation*}%
We can fill in the pieces to obtain%
\begin{eqnarray*}
\left( \frac{2\pi }{\lambda }2t-\pi \right) &=&\left( 2m+1\right) \pi \\
\frac{2}{\lambda }2t-1 &=&\left( 2m+1\right) \\
\frac{2}{\lambda }2t &=&\left( 2m+1\right) +1 \\
\frac{4}{\lambda }t &=&\left( 2m+1\right) +1 \\
t &=&\frac{\lambda }{4}\left( 2m+2\right) \\
t &=&\frac{\lambda }{2}\left( m+1\right) \\
t &=&\frac{m+1}{2}\left( \frac{\lambda _{out}}{n_{\text{bubble}}}\right)
\end{eqnarray*}%
This is our condition for destructive interference for the bubble. We don't
have to, but we could write $m+1=p$ where $p$ is an integer that starts at $%
1 $ instead of zero. 
\begin{equation*}
t=\frac{p}{2}\left( \frac{\lambda _{out}}{n_{\text{bubble}}}\right) \qquad
p=1,2,\cdots
\end{equation*}%
And with this change we can see that our destructive interference equation
for the bubble is very like the condition for constructive interference for
the plane.

Hopefully, it is apparent that we have to start with our basic equations 
\begin{eqnarray*}
\Delta \phi &=&\left( \frac{2\pi }{\lambda }\Delta x+\Delta \phi _{o}\right)
=m2\pi \qquad m=0,\pm 1,\pm 2,\pm 3,\cdots \text{ Constructive} \\
\Delta \phi &=&\left( \frac{2\pi }{\lambda }\Delta x+\Delta \phi _{o}\right)
=\left( 2m+1\right) \pi \qquad m=0,\pm 1,\pm 2,\pm 3,\cdots \text{
Destructive}
\end{eqnarray*}%
each time we attempt an interference problem because the outcome depends on
both $\Delta x$ and $\Delta \phi _{o}.$ We have to construct the equation
each time for the interference condition we want (constructive or
destructive) finding $\Delta x$ and $\Delta \phi _{o}$ for the boundary
conditions we have.

\section{The Michelson Interferometer}

The Michelson interferometer is another device that uses path differences to
create interference fringes. \FRAME{dhF}{3.2007in}{3.1263in}{0pt}{}{}{Figure%
}{\special{language "Scientific Word";type "GRAPHIC";maintain-aspect-ratio
TRUE;display "USEDEF";valid_file "T";width 3.2007in;height 3.1263in;depth
0pt;original-width 3.1566in;original-height 3.0805in;cropleft "0";croptop
"1";cropright "1";cropbottom "0";tempfilename
'SVANNN5D.wmf';tempfile-properties "XPR";}}The device is shown in the
figure. A coherent light source is used. The light beam is split into two
beams that are usually at $90\unit{%
%TCIMACRO{\U{b0}}%
%BeginExpansion
{{}^\circ}%
%EndExpansion
}$ apart. The beams are reflected off of two mirrors back along the same
path and are mixed at the telescope. The result (with perfect alignment) is
a target fringe pattern like the first two shown below. \FRAME{dhF}{3.5362in%
}{1.6466in}{0pt}{}{}{Figure}{\special{language "Scientific Word";type
"GRAPHIC";maintain-aspect-ratio TRUE;display "USEDEF";valid_file "T";width
3.5362in;height 1.6466in;depth 0pt;original-width 3.4895in;original-height
1.6103in;cropleft "0";croptop "1";cropright "1";cropbottom "0";tempfilename
'SVANNN5E.wmf';tempfile-properties "XPR";}}If the alignment is off, you get
smaller fringes, but the system can still work. This is shown in the last
image in the previous figure.

In the figure, we have constructive interference in the center, but if we
move one of the mirrors half a wavelength, we would have destructive
interference and would see a dark spot in the center. This device gives us
the ability to measure distances on the order of the wavelength of the
light. When the distance is continuously changed, the pattern seems to grow
from the center (or collapse into the center).

Notice that if the mirror is moved $\frac{\lambda }{2},$ the path distance
changes by $\lambda $ because the light travels the distance to the mirror
and then back from the mirror (it travels the path twice!).

\chapter{21 Slit Intensity 3.4.1, 3.4.2}

%TCIMACRO{%
%\TeXButton{\chaptermark{Slit Intensity}}{\chaptermark{Slit Intensity}}}%
%BeginExpansion
\chaptermark{Slit Intensity}%
%EndExpansion

Two lectures ago we looked at light passing through two small openings that
we call slits. We found constructive and destructive interference. We now
know where the bright spots will be on a double slit pattern. And we know
where the dark spots are. But in between, the light fades from light to
dark. It would be good to know the intensity at every spot on the pattern.
We will develop methods for making such a calculation in this lecture.

%TCIMACRO{%
%\TeXButton{Fundamental Concepts}{\vspace{0.10in}\hspace{-0.37in}{\Large Fundamental Concepts\vspace{0.2in}}}}%
%BeginExpansion
\vspace{0.10in}\hspace{-0.37in}{\Large Fundamental Concepts\vspace{0.2in}}%
%EndExpansion

\begin{itemize}
\item Single wide openings also make constructive and destructive
interference patterns

\item Single slit minima

\item Phasors

\item Single slit intensity pattern

\item The sinc function
\end{itemize}

\section{Single Slits}

We have looked at interference from two slits, and for many slits. Here is
the figure again that shows how the constructive and destructive
interference pattern changes as we change the number of slits or openings.
Earlier we gave a name for these patterns of constructive and destructive
interference. We called them diffraction patterns. \FRAME{dtbpF}{3.4056in}{%
2.6333in}{0pt}{}{}{Figure}{\special{language "Scientific Word";type
"GRAPHIC";maintain-aspect-ratio TRUE;display "USEDEF";valid_file "T";width
3.4056in;height 2.6333in;depth 0pt;original-width 7.5145in;original-height
5.8053in;cropleft "0";croptop "1";cropright "1";cropbottom "0";tempfilename
'SXCLAP05.wmf';tempfile-properties "XPR";}}The two slits acted like two
sources. We might expect that a single slit will give only a single bright
spot. But if you look carefully at the last figure, you can see that there
is something odd about the one slit pattern. The spot of light is stretched
out. And if you look very carefully it almost looks like there are faint
splotches of red light to the sides of the brighter fringe of light. let's
consider a single slit very closely. To do this, let's return to the work of
Huygens.\footnote{%
Huygens method is technically not a correct representation of what happens.
The actual wave leaving the single opening is a superposition of the
original wave, and the wave scattered from the sides of the opening. You can
see this scattering by tearing a small hole in a piece of paper and looking
through the hole at a light source. You will see the bright ring around the
hole where the edges of the paper are scattering the light. But the
mathematical result we will get using Huygens method gives a mathematically
identical result for the resulting wave leaving the slit with much less high
power math. So we will stick with Huygens in this class.} His idea for the
nature of light was simple. He suggested that every point on the wave front
of a light wave was the source (the disturbance) for a new set of small
spherical waves. In optics, the crests of the waves are often called
wavefronts. The next wavefront would be formed by the superposition of the
little \textquotedblleft wavelets.\textquotedblright\ Here is an example for
a plane wave and a spherical wave.\FRAME{dhF}{2.2208in}{1.8697in}{0pt}{}{}{%
Figure}{\special{language "Scientific Word";type
"GRAPHIC";maintain-aspect-ratio TRUE;display "USEDEF";valid_file "T";width
2.2208in;height 1.8697in;depth 0pt;original-width 4.0594in;original-height
3.4143in;cropleft "0";croptop "1";cropright "1";cropbottom "0";tempfilename
'SUS73V4V.wmf';tempfile-properties "XPR";}}In each case we have drawn spots
on the wave front and drawn spherical waves around those spots. where the
wavefronts of the little wavelets combine, we have new wave front of our
wave. We have already seen that this is sort of what happens in bulk matter.
Remember that light is absorbed and re-emitted by the atoms of the material.
The scattered and original waves superimpose and the combined wave
effectively goes slower. But the light is not necessarily re-emitted in the
same direction. Sometimes it is, but sometimes it is not. This creates a
small, spherical wave (sometimes called a wavelet) that is emitted by that
atom. So Huygens' idea is not too bad. But in our opening there aren't
really atoms to scatter the light wave. What really happens is that the
light scatters from the edges of the aperture and then combines with the
part of the wave that goes straight through. The math to show this is hard.
And it gives the same result as Huygens' method. So we will use the simpler
math, but know that it is not what physically happens.

We can use this idea for a single slit and look at what happens as the light
goes through. Here is such a slit.

\FRAME{fhFU}{2.0686in}{2.693in}{0pt}{\Qcb{{\protect\small Single Slit
Geometry}}}{\Qlb{SingleSlitGeometry}}{Figure}{\special{language "Scientific
Word";type "GRAPHIC";maintain-aspect-ratio TRUE;display "USEDEF";valid_file
"T";width 2.0686in;height 2.693in;depth 0pt;original-width
2.0297in;original-height 2.6507in;cropleft "0";croptop "1";cropright
"1";cropbottom "0";tempfilename 'SUS73V4W.wmf';tempfile-properties "XPR";}}

In the figure above, I have divided a single slit of width $a$ into two
parts, each of size $a/2.$ According to Huygens' principle, each position of
the slit acts as a source of light rays. We can treat half a slit as two
coherent sources. And think of what we have done in our last two lectures.
We have conditions for constructive and destructive interference. 
\begin{eqnarray*}
\Delta \phi &=&\left( \frac{2\pi }{\lambda }\Delta r+\Delta \phi _{o}\right)
=m2\pi \qquad m=0,\pm 1,\pm 2,\pm 3,\cdots \text{ Constructive} \\
\Delta \phi &=&\left( \frac{2\pi }{\lambda }\Delta r+\Delta \phi _{o}\right)
=\left( 2m+1\right) \pi \qquad m=0,\pm 1,\pm 2,\pm 3,\cdots \text{
Destructive}
\end{eqnarray*}%
The light comes from the same source so $\Delta \phi _{o}=0.$ But there
would be a path difference $\Delta r$ for light traveling from the top of
the aperture and light traveling from the middle of the aperture. These two
sources should interfere. And depending on the value of $\Delta r$ we could
have constructive and destructive interference. So what do we see when we
perform such an experiment? Here is a brighter version of our light pattern
from a single slit.

\FRAME{dhF}{2.3065in}{0.8233in}{0pt}{}{}{Figure}{\special{language
"Scientific Word";type "GRAPHIC";maintain-aspect-ratio TRUE;display
"USEDEF";valid_file "T";width 2.3065in;height 0.8233in;depth
0pt;original-width 2.2667in;original-height 0.7904in;cropleft "0";croptop
"1";cropright "1";cropbottom "0";tempfilename
'SUS73V4X.wmf';tempfile-properties "XPR";}}The figure shows a diffraction
pattern for a single thin slit. There are several terms that are in common
use to describe the pattern

\begin{enumerate}
\item Central Maximum: The broad intense central band.

\item Secondary Maxima:The fainter bright bands to both sides of the central
maxima

\item Minima: The dark bands between the maxima
\end{enumerate}

Notice, if we take two of the rays that hit one of our constructive
interference spots, we get a path difference $\Delta r$ that is 
\begin{equation*}
\Delta r=\frac{a}{2}\sin \theta
\end{equation*}%
and we can put this into our conditions for constructive or destructive
interference 
\begin{eqnarray*}
\Delta \phi &=&\left( \frac{2\pi }{\lambda }\left( \frac{a}{2}\sin \theta
\right) +\Delta \phi _{o}\right) =m2\pi \qquad m=0,\pm 1,\pm 2,\pm 3,\cdots 
\text{ Constructive} \\
\Delta \phi &=&\left( \frac{2\pi }{\lambda }\left( \frac{a}{2}\sin \theta
\right) +\Delta \phi _{o}\right) =\left( 2m+1\right) \pi \qquad m=0,\pm
1,\pm 2,\pm 3,\cdots \text{ Destructive}
\end{eqnarray*}

\section{Narrow Slit Minima}

Let's use figure \ref{SingleSlitGeometry} to find the \emph{dark} minima of
the single slit pattern. That would be destructive interference.

First we should notice that figure \ref{SingleSlitGeometry} could have
another set of rays that contribute to the \emph{dark} spot because they
will also have the same path difference of $\left( a/2\right) \sin \theta .$
Let's fill these in. They are rays $2$ and $4$ of the next figure. \FRAME{dhF%
}{2.1525in}{2.6507in}{0pt}{}{}{Figure}{\special{language "Scientific
Word";type "GRAPHIC";maintain-aspect-ratio TRUE;display "USEDEF";valid_file
"T";width 2.1525in;height 2.6507in;depth 0pt;original-width
2.1136in;original-height 2.6091in;cropleft "0";croptop "1";cropright
"1";cropbottom "0";tempfilename 'SUS73V4Y.wmf';tempfile-properties "XPR";}}

Before we started with what we are now calling rays $1$ and $3.$ Ray $1$
travels a distance 
\begin{equation}
\Delta r=\frac{a}{2}\sin \left( \theta \right)
\end{equation}%
farther than ray $3.$ As we just argued, rays $2$ and $4$ also have the same
path difference, and so do rays $3$ and $5.$ We can use our condition for
destructive interference 
\begin{eqnarray*}
\left( \frac{2\pi }{\lambda }\left( \frac{a}{2}\sin \theta \right) +0\right)
&=&\left( 2m+1\right) \pi \\
\frac{\pi }{\lambda }\left( a\sin \theta \right) &=&\left( 2m+1\right) \pi \\
a\sin \theta &=&\left( 2m+1\right) \lambda \\
\sin \theta &=&\left( 2m+1\right) \frac{\lambda }{a}
\end{eqnarray*}%
and for $m=0$ we have 
\begin{equation}
\sin \left( \theta \right) =\pm \frac{\lambda }{a}
\end{equation}

But look at our figure. We could just as well divide the slit into four
equal parts \FRAME{dtbpF}{2.0245in}{2.4933in}{0pt}{}{}{Figure}{\special%
{language "Scientific Word";type "GRAPHIC";maintain-aspect-ratio
TRUE;display "USEDEF";valid_file "T";width 2.0245in;height 2.4933in;depth
0pt;original-width 1.9856in;original-height 2.4526in;cropleft "0";croptop
"1";cropright "1";cropbottom "0";tempfilename
'SXCHJY00.wmf';tempfile-properties "XPR";}}and imagine light traveling from
each of those parts to our viewing screen. Then we have a path difference of 
\begin{equation}
\Delta r=\frac{a}{4}\sin \left( \theta \right)
\end{equation}%
Let's put this into our condition for destructive interference 
\begin{eqnarray*}
\left( \frac{2\pi }{\lambda }\left( \frac{a}{4}\sin \theta \right) +0\right)
&=&\left( 2m+1\right) \pi \\
\frac{\pi }{2\lambda }\left( a\sin \theta \right) &=&\left( 2m+1\right) \pi
\\
a\sin \theta &=&2\left( 2m+1\right) \lambda \\
\sin \theta &=&\left( 2m+1\right) \frac{2\lambda }{a}
\end{eqnarray*}%
And again for the $m=0$ case we have 
\begin{equation}
\sin \left( \theta \right) =\pm \frac{2\lambda }{a}
\end{equation}%
But there is no reason to stop at four equal parts. Suppose we we divide our
opening into six equal parts. \FRAME{dhF}{1.9285in}{2.5114in}{0pt}{}{}{Figure%
}{\special{language "Scientific Word";type "GRAPHIC";maintain-aspect-ratio
TRUE;display "USEDEF";valid_file "T";width 1.9285in;height 2.5114in;depth
0pt;original-width 1.8905in;original-height 2.4708in;cropleft "0";croptop
"1";cropright "1";cropbottom "0";tempfilename
'SUS73V4Z.wmf';tempfile-properties "XPR";}}Then we would find%
\begin{equation}
\Delta r=\frac{a}{6}\sin \left( \theta \right)
\end{equation}%
and 
\begin{eqnarray*}
\left( \frac{2\pi }{\lambda }\left( \frac{a}{6}\sin \theta \right) +0\right)
&=&\left( 2m+1\right) \pi \\
\frac{\pi }{3\lambda }\left( a\sin \theta \right) &=&\left( 2m+1\right) \pi
\\
a\sin \theta &=&3\left( 2m+1\right) \lambda \\
\sin \theta &=&\left( 2m+1\right) \frac{3\lambda }{a}
\end{eqnarray*}%
or for $m=0$ 
\begin{equation}
\sin \left( \theta \right) =\pm \frac{3\lambda }{a}
\end{equation}%
and in general at the angle $\theta $ that makes the sine of that angle
equal to $\mathfrak{m}\lambda /a$ we get a dark spot.%
\begin{equation}
\sin \left( \theta \right) =\mathfrak{m}\frac{\lambda }{a}\quad \mathfrak{m}%
=\pm 1,\pm 2,\pm 3\ldots  \label{Single Slit minima}
\end{equation}

You might object that we did not find the bright spots, only the dark spots.
That's fine because the bright spots have to be in between the dark spots.
But we can do better. Instead of just finding the bright spots, let's find
the brightness of every spot along our intensity pattern for a single slit.

\section{Intensity of the single-slit pattern}

\FRAME{dtbpF}{4.3267in}{2.1828in}{0pt}{}{}{Figure}{\special{language
"Scientific Word";type "GRAPHIC";maintain-aspect-ratio TRUE;display
"USEDEF";valid_file "T";width 4.3267in;height 2.1828in;depth
0pt;original-width 4.2748in;original-height 2.1439in;cropleft "0";croptop
"1";cropright "1";cropbottom "0";tempfilename
'../../../../PH375/repository/OpticsText/text/SVWZBT04.wmf';tempfile-properties "XPR";}%
}

We have found the dark minima, now we would like to find the shape of the
bright bands. We will use phasors to do this. If you have already taken
Principles of Physics III (PH220) or if you are an electrical engineer you
know what a phasor is, but for the rest of us we need a quick introduction.
We will introduce phasors here and use them extensively for alternating
current in Principles of Physics III (PH220) and in our Optics (PH375)
class. Depending on your major, you many see them in many other classes.

\subsection{Phasors}

Let's take two waves%
\begin{equation}
E_{1}=E_{o1}\sin \left( \omega t+\alpha _{1}\left( x\right) \right)
\end{equation}%
and 
\begin{equation}
E_{2}=E_{o2}\sin \left( \omega t+\alpha _{2}\left( x\right) \right)
\end{equation}%
where I have hidden the position dependence in the $\alpha $'s 
\begin{eqnarray*}
\alpha _{1} &=&kx_{1}+\phi _{1} \\
\alpha _{2} &=&kx_{2}+\phi _{2}
\end{eqnarray*}

We wish to find the resultant wave 
\begin{equation}
E=E_{o1}+E_{o2}
\end{equation}%
We actually know how to do the algebra to combine these waves (we have done
it enough times!), but let's find another way to do the math that can extend
to many many waves being added up.

Think of our equation for a wave. 
\begin{equation}
E_{1}=E_{o1}\sin \left( \omega t+\alpha _{1}\right)
\end{equation}%
Because we have essentially picked just one location, $x_{1}$ or $x_{2},$
for our waves this is like a history graph view. We have made our wave
equation look like the equation for a harmonic oscillator. And so long as we
observe our waves at just one location, so we don't change $x_{1}$ or $x_{2}$
this is fine. The electric field at our position where we are mixing the
waves will oscillate like a harmonic oscillator.

Way back when we studied oscillation we said the equation for a harmonic
oscillator position was like the projection of circular motion onto an axis. 
\FRAME{dtbpF}{2.3687in}{2.3186in}{0pt}{}{}{Figure}{\special{language
"Scientific Word";type "GRAPHIC";maintain-aspect-ratio TRUE;display
"USEDEF";valid_file "T";width 2.3687in;height 2.3186in;depth
0pt;original-width 2.3281in;original-height 2.2788in;cropleft "0";croptop
"1";cropright "1";cropbottom "0";tempfilename
'SXCLAP06.wmf';tempfile-properties "XPR";}}That gave us 
\begin{equation*}
x\left( t\right) =x_{\max }\cos \left( \omega t+\phi _{o}\right)
\end{equation*}%
or 
\begin{equation*}
y\left( t\right) =y_{\max }\sin \left( \omega t+\phi _{o}\right)
\end{equation*}%
as our oscillator equations. And see how much our modified wave equation
looks like the second one of these oscillator equations! \textbf{This gives
us an idea of how to use vector addition to add up waves!}

Let's represent $E_{o1}$ as a rotating vector with angular speed $\omega .$
We get projections of $E_{o1}$ on the $x$ and $y$ axis. In this discussion,
we will consider the vertical $y$-axis but we could choose either axis (if
we use the $x$-axis we would then have cosine waves). \FRAME{dtbpF}{3.1176in%
}{1.8057in}{0pt}{}{}{Figure}{\special{language "Scientific Word";type
"GRAPHIC";maintain-aspect-ratio TRUE;display "USEDEF";valid_file "T";width
3.1176in;height 1.8057in;depth 0pt;original-width 3.0727in;original-height
1.7685in;cropleft "0";croptop "1";cropright "1";cropbottom "0";tempfilename
'SXCIU601.wmf';tempfile-properties "XPR";}}At any given time, $E_{1}$ will
be the projection $E_{o1}\sin \left( \omega t+\alpha _{1}\right) $ on the $y$%
-axis. We can also take a second rotating vector for $E_{o2}$ \FRAME{dtbpF}{%
3.1176in}{1.9156in}{0pt}{}{}{Figure}{\special{language "Scientific
Word";type "GRAPHIC";maintain-aspect-ratio TRUE;display "USEDEF";valid_file
"T";width 3.1176in;height 1.9156in;depth 0pt;original-width
2.5832in;original-height 1.5766in;cropleft "0";croptop "1";cropright
"1";cropbottom "0";tempfilename 'SXCIUX02.wmf';tempfile-properties "XPR";}}%
Note that $E_{o2}$ is projected onto the vertical axis through an angle $%
\omega t+\alpha _{2}.$ So the projection $E_{2}$ is $E_{o2}\sin \left(
\omega t+\alpha _{2}\right) $.

These rotating vectors have a name. We call them \emph{phasors}.

To add the phasors, we draw phasor $2$ on the end of phasor $1$. It is just
like adding vectors back in Principles of Physics I (PH121)! \FRAME{dtbpF}{%
4.12in}{2.3333in}{0pt}{}{}{Figure}{\special{language "Scientific Word";type
"GRAPHIC";maintain-aspect-ratio TRUE;display "USEDEF";valid_file "T";width
4.12in;height 2.3333in;depth 0pt;original-width 4.0707in;original-height
2.2935in;cropleft "0";croptop "1";cropright "1";cropbottom "0";tempfilename
'SXCIVV03.wmf';tempfile-properties "XPR";}}

The sum is the resultant phasor like in vector addition. Only, remember now
that the phasors are rotating. But because they have the same angular
frequency, $\omega ,$ phasor 1 and phasor 2 rotate together as a unit. The 
\emph{projection} of this rotating group of phasors gives the sum we are
looking for. But notice the sum's numerical value would change in time as
the phasors rotate. We can choose any time to find the resulting vector, so
we might as well take $t=0,$ then $\omega t=0$ and the vectors will be
rotated as shown. \FRAME{dtbpF}{3.6063in}{1.8585in}{0pt}{}{}{Figure}{\special%
{language "Scientific Word";type "GRAPHIC";maintain-aspect-ratio
TRUE;display "USEDEF";valid_file "T";width 3.6063in;height 1.8585in;depth
0pt;original-width 3.5587in;original-height 1.8204in;cropleft "0";croptop
"1";cropright "1";cropbottom "0";tempfilename
'TB346T0B.wmf';tempfile-properties "XPR";}}

Using the redrawn phasors in the diagram directly above, we can see from the
law of cosines \FRAME{dtbpF}{2.127in}{1.5575in}{0in}{}{}{Figure}{\special%
{language "Scientific Word";type "GRAPHIC";maintain-aspect-ratio
TRUE;display "USEDEF";valid_file "T";width 2.127in;height 1.5575in;depth
0in;original-width 2.1061in;original-height 1.5349in;cropleft "0";croptop
"1";cropright "1";cropbottom "0";tempfilename
'TG65J800.wmf';tempfile-properties "XPR";}}%
\begin{equation*}
b^{2}=a^{2}+c^{2}-2ac\cos \left( B\right)
\end{equation*}%
that 
\begin{equation}
E_{0}^{2}=E_{o1}^{2}+E_{o2}^{2}-2E_{o1}E_{o2}\cos \left( \eta \right)
\end{equation}%
but from the figure we can see that $\eta =\pi -\left( \alpha _{2}-\alpha
_{1}\right) .$ That gives us%
\begin{equation}
E_{0}^{2}=E_{o1}^{2}+E_{o2}^{2}-2E_{o1}E_{o2}\cos \left( \pi -\left( \alpha
_{2}-\alpha _{1}\right) \right)
\end{equation}%
or using the trig identity that $\cos \left( \pi -\gamma \right) =-\cos
\gamma $%
\begin{equation}
E_{0}^{2}=E_{o1}^{2}+E_{o2}^{2}+2E_{o1}E_{o2}\cos \left( \alpha _{2}-\alpha
_{1}\right)
\end{equation}
This is the magnitude of our resulting vector. For the case of a double slit
we know that $E_{o1}=E_{o2}$ so our resulting vector would be 
\begin{eqnarray}
E_{0}^{2} &=&E_{o1}^{2}+E_{o1}^{2}+2E_{o1}E_{o1}\cos \left( \alpha
_{2}-\alpha _{1}\right) \\
&=&2E_{o1}^{2}+2E_{o1}^{2}\cos \left( \alpha _{2}-\alpha _{1}\right)
\end{eqnarray}%
We can see that when $\cos \left( \alpha _{2}-\alpha _{1}\right) =-1$ we get
zero. This is destructive interference!

We should go on to find the angle that makes this happen. In the figure it
is called just $\alpha $ and it is $\alpha _{1}$ plus a little bit more, $%
\theta $. It's not thrilling, but since we have the sides of the triangle we
could just use the law of cosines again to find the little bit more angle%
\begin{equation*}
E_{o1}^{2}=E_{0}^{2}+E_{o1}^{2}+2E_{0}E_{o2}\cos \left( \theta \right)
\end{equation*}%
so 
\begin{equation*}
\theta =\cos ^{-1}\frac{E_{o2}^{2}-E_{0}^{2}+E_{o1}^{2}}{2E_{0}E_{o2}}
\end{equation*}%
or when $E_{o1}=E_{o2}$%
\begin{equation*}
\theta =\cos ^{-1}\frac{2E_{o1}^{2}-E_{0}^{2}}{2E_{0}E_{o1}}
\end{equation*}%
and 
\begin{equation*}
\alpha =\alpha _{1}+\theta
\end{equation*}

Now we ask what the question, \textquotedblleft what is the projection of
this resultant phasor on the vertical axis?\textquotedblright\ Since we
chose $t=0$ for our picture, it is easy to pick off the fact that the
resultant vector has a phase angle of $\alpha $. thus%
\begin{equation}
E_{R}=E_{0}\sin \left( \omega t+\alpha \right)
\end{equation}

This seems like a lot of work for adding just two waves together when we
could have just used our favorite trig identity, but we can use it to add
many many waves together.

\subsection{Using phasors to find the single slit intensity pattern}

Now that we know what a phasor is, let's try to use them to find the single
slit intensity pattern.

We can consider dividing the slit into small regions $\Delta y$ in length. 
\FRAME{dtbpF}{2.2883in}{2.9776in}{0pt}{}{}{Figure}{\special{language
"Scientific Word";type "GRAPHIC";maintain-aspect-ratio TRUE;display
"USEDEF";valid_file "T";width 2.2883in;height 2.9776in;depth
0pt;original-width 2.2485in;original-height 2.9334in;cropleft "0";croptop
"1";cropright "1";cropbottom "0";tempfilename
'TB34BZ0C.wmf';tempfile-properties "XPR";}}Each piece of the slit
contributes $\Delta E$ to the electric field. We can use superposition to
find the resultant wave for a given point $P.$ The phase diagram would look
something like this. \FRAME{dtbpF}{4.0966in}{3.3849in}{0pt}{}{}{Figure}{%
\special{language "Scientific Word";type "GRAPHIC";maintain-aspect-ratio
TRUE;display "USEDEF";valid_file "T";width 4.0966in;height 3.3849in;depth
0pt;original-width 4.0465in;original-height 3.3382in;cropleft "0";croptop
"1";cropright "1";cropbottom "0";tempfilename
'TB34GF0D.wmf';tempfile-properties "XPR";}}We add up all the little $\Delta
E $ parts, but each has a slightly different distance, so it has a slightly
different phase because $kr$ is part of the phase. We put the phase
difference due to the different paths into the phase constant.

Our set of little waves might be like

\begin{eqnarray}
\Delta E_{1} &=&E_{o}\sin \left( \omega t\right) \\
\Delta E_{2} &=&E_{o}\sin \left( \omega t+\Delta \phi _{2}\right)  \notag \\
\Delta E_{3} &=&E_{o}\sin \left( \omega t+\Delta \phi _{3}\right)  \notag \\
&&\vdots  \notag
\end{eqnarray}
Note that these are still waves. We have just put the $kr_{i}$ part into the
phase $\Delta \phi _{i}$ like we did in the $\alpha $'s in the last section.

Note that each wave has an additional phase shift, and each has the
amplitude $E_{o.}$ But we need to find these phase shifts. Look at the next
figure. \FRAME{dtbpF}{1.9225in}{2.5019in}{0pt}{}{}{Figure}{\special{language
"Scientific Word";type "GRAPHIC";maintain-aspect-ratio TRUE;display
"USEDEF";valid_file "T";width 1.9225in;height 2.5019in;depth
0pt;original-width 1.8853in;original-height 2.4604in;cropleft "0";croptop
"1";cropright "1";cropbottom "0";tempfilename
'../../../../PH375/repository/OpticsText/text/SVWZBT07.wmf';tempfile-properties "XPR";}%
}By thinking about this figure we can see that 
\begin{equation}
\Delta r=\Delta y\sin \theta
\end{equation}%
for each pair of adjacent rays. And that gives a phase difference of 
\begin{equation}
\Delta \phi =k\Delta r=\frac{2\pi }{\lambda }\Delta y\sin \theta
\end{equation}%
Let's suppose that as we cross the slit we have $N$ waves each of magnitude $%
\Delta E$. Then the total phase difference between waves from the top to the
bottom of the slit is 
\begin{eqnarray}
\Delta \phi _{total} &=&N\Delta \phi =N\frac{2\pi }{\lambda }\Delta y\sin
\theta \\
&=&\frac{2\pi }{\lambda }a\sin \theta
\end{eqnarray}%
where $a$ is our slit width and 
\begin{equation}
a=N\Delta y
\end{equation}

We are finding dark spots in our diffraction pattern, and the first dark
spot will be when the projection of our phasor sum, $E_{R}=0.$ And think
about this, when you add vectors you put the vectors tip-to-tail, and when
you are done you then take your pencil and start at the tail of the first
and draw a new vector to the tip of the last vector you added. This new
vector is the sum. And to make that new vector have a length $E_{R}=0$ all
we have to do is to have the vectors we add make a circle. Then the distance
from the tail of the first vector to the tip of the last vector is exactly
zero! We are doing this with phasors, so we place our phasors tip-to-tail
and find the first minimum or first dark band when $\Delta \phi
_{total}=2\pi $ or when the phasor diagram circles back on itself.\FRAME{%
dtbpF}{1.7409in}{1.6561in}{0pt}{}{}{Figure}{\special{language "Scientific
Word";type "GRAPHIC";maintain-aspect-ratio TRUE;display "USEDEF";valid_file
"T";width 1.7409in;height 1.6561in;depth 0pt;original-width
1.7028in;original-height 1.6198in;cropleft "0";croptop "1";cropright
"1";cropbottom "0";tempfilename
'../../../../PH375/repository/OpticsText/text/SVWZBT00.wmf';tempfile-properties "XPR";}%
}For our complete circle $\Delta \phi _{total}=2\pi $ so then%
\begin{equation}
2\pi =\frac{2\pi }{\lambda }a\sin \theta
\end{equation}%
We can cancel and rearrange this to give%
\begin{equation}
\frac{\lambda }{a}=\sin \theta
\end{equation}%
which is what we expected for the first minimum! This seems to be working.
Let's now find the whole intensity pattern.

To find the intensity pattern on any point of the screen (not just the
minima), we let $\Delta y$ get very small. We add up $N\rightarrow \infty $
terms. The curve becomes smooth (not shown as smooth in the next figure,
though). \FRAME{dtbpF}{3.4298in}{2.4111in}{0in}{}{}{Figure}{\special%
{language "Scientific Word";type "GRAPHIC";maintain-aspect-ratio
TRUE;display "USEDEF";valid_file "T";width 3.4298in;height 2.4111in;depth
0in;original-width 3.384in;original-height 2.3705in;cropleft "0";croptop
"1";cropright "1";cropbottom "0";tempfilename
'TB35FC0I.wmf';tempfile-properties "XPR";}}

The quantity $E_{R}$ is the resultant from adding together all of our
phasors. We can see the phasor addition in the figure. The angle at the top
of the yellow triangle is $\zeta /2.$ One way to see this is to note that
the lower triangle has angles $\beta ,$ $\gamma ,$ and $\beta $ \FRAME{dtbpF%
}{2.0433in}{1.9604in}{0in}{}{}{Figure}{\special{language "Scientific
Word";type "GRAPHIC";maintain-aspect-ratio TRUE;display "USEDEF";valid_file
"T";width 2.0433in;height 1.9604in;depth 0in;original-width
2.0224in;original-height 1.9395in;cropleft "0";croptop "1";cropright
"1";cropbottom "0";tempfilename 'TG66GF01.wmf';tempfile-properties "XPR";}}

We see that 
\begin{equation}
\gamma =180-\zeta
\end{equation}%
and%
\begin{equation}
\zeta =2\beta
\end{equation}

Note also that 
\begin{equation}
\eta =90-\beta
\end{equation}

For the upper triangle, 
\begin{eqnarray}
180 &=&\psi +\eta +90 \\
90 &=&\psi +\eta
\end{eqnarray}%
we know $\eta ,$ so 
\begin{equation}
90=\psi +90-\beta
\end{equation}%
or%
\begin{equation}
\psi =\beta
\end{equation}%
but%
\begin{equation}
\beta =\frac{\zeta }{2}
\end{equation}%
so%
\begin{equation}
\psi =\frac{\zeta }{2}
\end{equation}

Then%
\begin{equation}
\sin \left( \frac{\zeta }{2}\right) =\frac{\frac{E_{R}}{2}}{R}=\frac{E_{R}}{%
2R}
\end{equation}%
or 
\begin{equation*}
E_{R}=2R\sin \left( \frac{\zeta }{2}\right) =2R\sin \left( \beta \right)
\end{equation*}%
where $R$ is the radius of curvature. Let's see if we can find a value for $%
R.$ From the archlength formula%
\begin{equation}
s=r\left( 2\psi \right) =R\zeta
\end{equation}%
And since we are finding the arc length, we know it will be nearly the sum
of each little phasor vector length $s=NE_{o}$ so%
\begin{equation*}
s=R\zeta =NE_{o}
\end{equation*}%
This gives our value for $R$%
\begin{equation*}
R=\frac{NE_{o}}{\zeta }
\end{equation*}%
Then our resultant wave phasor has length%
\begin{eqnarray}
E_{R} &=&2\left( R\right) \sin \left( \beta \right) \\
&=&2\left( \frac{NE}{\zeta }\right) \sin \left( \beta \right)
\end{eqnarray}%
And we know that $\beta =\zeta /2$ so 
\begin{equation*}
E_{R}=NE_{o}\frac{\sin \left( \beta \right) }{\beta }
\end{equation*}

Earlier we learned that the intensity will go as the amplitude of the field
squared%
\begin{equation}
I=I_{\max }\left( \frac{\sin \left( \beta \right) }{\beta }\right) ^{2}
\end{equation}%
And here is a tricky part. We get the angle $\zeta $ from adding up the
angles from each of our little $\Delta E$ pieces. Each of them has an angle%
\begin{equation}
\Delta \phi =k\Delta r=\frac{2\pi }{\lambda }\Delta y\sin \theta
\end{equation}%
and $N$ of these will be 
\begin{equation*}
\zeta =N\frac{2\pi }{\lambda }\Delta y\sin \theta
\end{equation*}%
but we already found that 
\begin{equation}
N\frac{2\pi }{\lambda }\Delta y\sin \theta =\Delta \phi _{total}
\end{equation}%
So we can see that the angle $\zeta $ is the same as our total phase $\Delta
\phi _{total}$. If we have an aperture width of $a$ then $a=N\Delta y$ and 
\begin{equation*}
\Delta \phi _{total}=N\frac{2\pi }{\lambda }\Delta y\sin \theta =\frac{2\pi 
}{\lambda }a\sin \theta
\end{equation*}%
and that makes 
\begin{equation*}
\beta =\frac{\zeta }{2}=\frac{\Delta \phi _{total}}{2}=\frac{\frac{2\pi }{%
\lambda }a\sin \theta }{2}
\end{equation*}%
We can write our irradiance function 
\begin{equation}
I=I_{\max }\left( \frac{\sin \left( \beta \right) }{\left( \beta \right) }%
\right) ^{2}
\end{equation}%
as 
\begin{equation}
I=I_{\max }\left( \frac{\sin \left( \frac{\frac{2\pi }{\lambda }a\sin \theta 
}{2}\right) }{\frac{\frac{2\pi }{\lambda }a\sin \theta }{2}}\right) ^{2}
\end{equation}%
or more simply%
\begin{equation}
I=I_{\max }\left( \frac{\sin \left( \frac{\pi }{\lambda }a\sin \theta
\right) }{\frac{\pi }{\lambda }a\sin \theta }\right) ^{2}
\end{equation}%
The function is plotted in the next figure\FRAME{dtbpFX}{4.4996in}{2.2502in}{%
0pt}{}{}{Plot}{\special{language "Scientific Word";type "MAPLEPLOT";width
4.4996in;height 2.2502in;depth 0pt;display "USEDEF";plot_snapshots
TRUE;mustRecompute FALSE;lastEngine "MuPAD";xmin "-15";xmax "15";xviewmin
"-15";xviewmax "15";yviewmin "-1";yviewmax
"1.002079";viewset"XY";rangeset"X";plottype 4;labeloverrides 3;x-label
"alpha";y-label "I";axesFont "Times New
Roman,12,0000000000,useDefault,normal";numpoints 100;plotstyle
"patch";axesstyle "normal";axestips FALSE;xis \TEXUX{x};var1name
\TEXUX{$x$};function \TEXUX{$\left( \frac{\sin x}{x}\right) ^{2}$};linecolor
"blue";linestyle 1;pointstyle "point";linethickness 3;lineAttributes
"Solid";var1range "-15,15";num-x-gridlines 100;curveColor
"[flat::RGB:0x000000ff]";curveStyle "Line";VCamFile
'SVWYP10O.xvz';valid_file "T";tempfilename
'TG4C1I00.wmf';tempfile-properties "XPR";}}From this we can find our peaks
where the pattern is bright.

Notice that our intensity pattern equation has in it the form%
\begin{equation*}
\frac{\sin x}{x}
\end{equation*}%
which has a distinctive shape.\FRAME{dtbpFX}{4.4996in}{1.3958in}{0pt}{}{}{%
Plot}{\special{language "Scientific Word";type "MAPLEPLOT";width
4.4996in;height 1.3958in;depth 0pt;display "USEDEF";plot_snapshots
TRUE;mustRecompute FALSE;lastEngine "MuPAD";xmin "-30";xmax "30";xviewmin
"-30";xviewmax "30";yviewmin "-0.217351";yviewmax
"1.000122";viewset"XY";rangeset"X";plottype 4;labeloverrides 3;x-label
"x";y-label "sinc(x)";axesFont "Times New
Roman,12,0000000000,useDefault,normal";numpoints 100;plotstyle
"patch";axesstyle "normal";axestips FALSE;xis \TEXUX{x};var1name
\TEXUX{$x$};function \TEXUX{$\frac{\sin x}{x}$};linecolor "blue";linestyle
1;pointstyle "point";linethickness 3;lineAttributes "Solid";var1range
"-30,30";num-x-gridlines 100;curveColor "[flat::RGB:0x000000ff]";curveStyle
"Line";VCamFile 'SVAN6U3W.xvz';valid_file "T";tempfilename
'SUS73V50.wmf';tempfile-properties "XPR";}}This is known as a sinc function
(pronounced like \textquotedblleft sink\textquotedblright ). It has a
central maximum as we would expect. Of course our pattern has a sinc squared 
\FRAME{dtbpFX}{4.4997in}{2.2494in}{0pt}{}{}{Plot}{\special{language
"Scientific Word";type "MAPLEPLOT";width 4.4997in;height 2.2494in;depth
0pt;display "USEDEF";plot_snapshots TRUE;mustRecompute FALSE;lastEngine
"MuPAD";xmin "-15";xmax "15";xviewmin "-15";xviewmax "15";yviewmin
"-1";yviewmax "1.002079";viewset"XY";rangeset"X";plottype 4;labeloverrides
3;x-label "pi*a*sin(theta)/lambda";y-label "I";axesFont "Times New
Roman,12,0000000000,useDefault,normal";numpoints 100;plotstyle
"patch";axesstyle "normal";axestips FALSE;xis \TEXUX{x};var1name
\TEXUX{$x$};function \TEXUX{$\left( \frac{\sin x}{x}\right) ^{2}$};linecolor
"blue";linestyle 1;pointstyle "point";linethickness 3;lineAttributes
"Solid";var1range "-15,15";num-x-gridlines 100;curveColor
"[flat::RGB:0x000000ff]";curveStyle "Line";VCamFile
'TJXLDZ2A.xvz';valid_file "T";tempfilename
'TJXLDR03.wmf';tempfile-properties "XPR";}}You can see the central maximum
and the much weaker minima produced by this function. Indeed, it seems to
match what we saw very well. Putting it all together, our pattern looks like
this.\FRAME{dhF}{3.0329in}{1.535in}{0pt}{}{}{Figure}{\special{language
"Scientific Word";type "GRAPHIC";maintain-aspect-ratio TRUE;display
"USEDEF";valid_file "T";width 3.0329in;height 1.535in;depth
0pt;original-width 2.9888in;original-height 1.4987in;cropleft "0";croptop
"1";cropright "1";cropbottom "0";tempfilename
'SUS73V52.wmf';tempfile-properties "XPR";}}

You might ask, can we use our phasor approach to find the intensity pattern
for double slit diffraction? Of course, the answer is yes, but I have had
enough phasor fun for this class. You will see phasors again in Principles
of Physics III and maybe again if you take Optics (or maybe for a career if
you are an electrical engineer).

\chapter{22 Double Slit and Grating intensity patterns 3.4.3, 3.4.4}

In the last lecture, we found the diffraction pattern for a single slit
opening. We used phasors to do this. In this lecture we want to do two or
more slits. But let's return to algebraic methods.

%TCIMACRO{\TeXButton{\chaptermark{Double Slit}}{\chaptermark{Double Slit}}}%
%BeginExpansion
\chaptermark{Double Slit}%
%EndExpansion

%TCIMACRO{%
%\TeXButton{Fundamental Concepts}{\vspace{0.10in}\hspace{-0.37in}{\Large Fundamental Concepts\vspace{0.2in}}}}%
%BeginExpansion
\vspace{0.10in}\hspace{-0.37in}{\Large Fundamental Concepts\vspace{0.2in}}%
%EndExpansion

\begin{itemize}
\item Double Slit Intensity Pattern

\item Grating Intensity Pattern

\item Integration over many periods to get what detectors (like eyes) see.
\end{itemize}

\section{Double Slit Intensity Pattern}

\FRAME{dtbpF}{4.5947in}{2.2303in}{0in}{}{}{Figure}{\special{language
"Scientific Word";type "GRAPHIC";maintain-aspect-ratio TRUE;display
"USEDEF";valid_file "T";width 4.5947in;height 2.2303in;depth
0in;original-width 4.5429in;original-height 2.1897in;cropleft "0";croptop
"1";cropright "1";cropbottom "0";tempfilename
'SXEBSK00.wmf';tempfile-properties "XPR";}}

Back in Chapter (\ref{DoubleSlitMath}) we found the bright and dark fringe
locations for a double slit diffraction pattern. But the fringes we have
seen are not just points. They are patterns that fade from a maximum
intensity. In the last lecture we calculated such an intensity pattern for
the light going through a single slit. We used phasors to do so. But for
double slits we already have an algebraic method for finding the
superposition of two waves. Let's use that method to find our intensity
pattern.

Let's let a laser beam hit two parallel slits, just like we did before. The
light coming through the two slits will be like two new sources of light but
with the same frequency and wavelength. We can represent an electromagnetic
wave using just the electric field (remember, the magnetic field pattern is
very similar and can be derived from the electric field pattern)%
\begin{eqnarray*}
E_{2} &=&E_{\max }\sin \left( k_{2}r_{2}-\omega _{2}t_{2}+\phi _{2}\right) \\
E_{1} &=&E_{\max }\sin \left( k_{1}r_{1}-\omega _{1}t_{1}+\phi _{1}\right)
\end{eqnarray*}%
and the resulting wave will be 
\begin{eqnarray*}
E_{r} &=&E_{\max }\sin \left( k_{2}r_{2}-\omega _{2}t_{2}+\phi _{2}\right)
+E_{\max }\sin \left( k_{1}r_{1}-\omega _{1}t_{1}+\phi _{1}\right) \\
&=&2E_{\max }\cos \left( \frac{\left( k_{2}r_{2}-\omega _{2}t_{2}+\phi
_{2}\right) -\left( k_{1}r_{1}-\omega _{1}t_{1}+\phi _{1}\right) }{2}\right)
\\
&&\times \sin \left( \frac{\left( k_{2}r_{2}-\omega _{2}t_{2}+\phi
_{2}\right) +\left( k_{1}r_{1}-\omega _{1}t_{1}+\phi _{1}\right) }{2}\right)
\\
&=&2E_{\max }\cos \left( \frac{1}{2}\left[ \left( k_{2}r_{2}-\omega
_{2}t_{2}+\phi _{2}\right) -\left( k_{1}r_{1}-\omega _{1}t_{1}+\phi
_{1}\right) \right] \right) \\
&&\times \sin \left( \frac{\left( k_{2}r_{2}-\omega _{2}t_{2}+\phi
_{2}\right) +\left( k_{1}r_{1}-\omega _{1}t_{1}+\phi _{1}\right) }{2}\right)
\end{eqnarray*}%
but now we know that we can simplify this because $\omega _{2}=\omega
_{1}=\omega ,$ $k_{2}=k_{1}=k,$ $t_{2}=t_{1}=t,$ and so long as the light
hits the slit screen perpendicular to the screen, $\phi _{2}=\phi _{1}=\phi
_{o}.$ Then%
\begin{equation*}
E_{r}=2E_{\max }\cos \left( \frac{1}{2}k\left( r_{2}-r_{1}\right) \right)
\sin \left( \frac{kr_{2}+kr_{1}}{2}-\omega t+\phi _{o}\right)
\end{equation*}%
This is very familiar territory. We have done this calculation many times
now. But let's go a little further this time.

We have a combined wave at point $P$ that is a traveling wave $\left( \sin
\left( k\frac{\left( r_{2}+r_{1}\right) }{2}-\omega t+\phi _{o}\right)
\right) $ but with amplitude $\left( 2E_{o}\cos \left( \frac{1}{2}k\left(
r_{2}-r_{1}\right) \right) \right) .$ And notice that $r_{2}-r_{1}=d\sin
\theta $ and $k=2\pi /\lambda $ so we can write $E_{r}$ as 
\begin{equation*}
E_{r}=2E_{\max }\cos \left( \frac{1}{2}\frac{2\pi }{\lambda }d\sin \theta
\right) \sin \left( k\frac{r_{2}+r_{1}}{2}-\omega t+\phi _{o}\right)
\end{equation*}%
and we have a wave with an amplitude that depends on our total phase $\Delta
\phi =\frac{2\pi }{\lambda }d\sin \theta .$

But the situation is more complicated because of how we detect light. Our
eyes and most light detectors measure the intensity of the light. We know
that 
\begin{equation*}
I=\frac{\mathcal{P}}{A}
\end{equation*}%
Earlier we surmised that the intensity of a light wave must also be
proportional to the amplitude of the electric field squared.\footnote{%
We will wait until Principles of Physics III to show this and to find the
constant of proportionality.}%
\begin{eqnarray*}
I &=&\frac{P}{A}\varpropto E^{2} \\
&\varpropto &4E_{o}^{2}\cos ^{2}\left( \frac{1}{2}\left( \frac{2\pi }{%
\lambda }d\sin \theta \right) \right) \sin ^{2}\left( \frac{k\left(
r_{2}+r_{1}\right) }{2}-\omega t+\phi _{o}\right)
\end{eqnarray*}

From studying lens systems and cameras we know that light detectors collect
power for a set amount of time. So most light detection will be an intensity
value averaged over a set \emph{integration time}. This means that the
detector sums up (or integrates) the amount of power received over the
detector time. Usually the integration time is much longer than a wave
period, so we need to average our intensity over many periods. back in our
calculus classes we learned that we can take an average of a function using
an integral%
\begin{equation*}
\left\langle f\left( t\right) \right\rangle =\frac{1}{\Delta t}%
\int_{0}^{\Delta t}f\left( t\right) dt
\end{equation*}%
and we used this to integrate sound waves over time to get an average
intensity. Let's do this again!

Suppose we choose to integrate over $N$ periods of our light wave. Then the
average is 
\begin{eqnarray*}
\frac{1}{NT}\int_{0}^{NT}Idt &\varpropto &\frac{1}{NT}%
\int_{0}^{NT}4E_{o}^{2}\cos ^{2}\left( \frac{1}{2}\left( \frac{2\pi }{%
\lambda }d\sin \theta \right) \right) \sin ^{2}\left( \frac{k\left(
r_{2}+r_{1}\right) }{2}-\omega t+\phi _{o}\right) dt \\
&=&4E_{o}^{2}\cos ^{2}\left( \frac{1}{2}\left( \frac{2\pi }{\lambda }d\sin
\theta \right) \right) \left[ \frac{1}{NT}\int_{0}^{NT}\sin ^{2}\left( \frac{%
k\left( rx_{2}+r_{1}\right) }{2}-\omega t+\phi _{o}\right) dt\right]
\end{eqnarray*}%
but the term%
\begin{equation}
\frac{1}{NT}\int_{0}^{NT}\sin ^{2}\left( \frac{k\left( r_{2}+r_{1}\right) }{2%
}-\omega t+\phi _{o}\right) dt=\frac{1}{2}
\end{equation}

To convince yourself of this, think of just $\sin ^{2}\left( \omega t\right)
.$ It has a maximum value of $1$ and a minimum of $0.$ Looking at the graph%
\FRAME{dhF}{1.9986in}{1.4927in}{0pt}{}{}{Figure}{\special{language
"Scientific Word";type "GRAPHIC";maintain-aspect-ratio TRUE;display
"USEDEF";valid_file "T";width 1.9986in;height 1.4927in;depth
0pt;original-width 1.9597in;original-height 1.4572in;cropleft "0";croptop
"1";cropright "1";cropbottom "0";tempfilename
'SVANNN5F.wmf';tempfile-properties "XPR";}}it should be believable that if
we integrate over one period we will get $1/2.$ The average over many
periods will still be $1/2.$ If you want to show this mathematically we can
think that the $\omega $ in $\sin ^{2}\left( \omega t\right) $ is really $%
\frac{2\pi }{T}$ and to take the average we want to integrate 
\begin{equation*}
\frac{1}{T-0}\int_{0}^{T}\sin ^{2}\left( \omega t\right) dt
\end{equation*}
over one period. We can use the trig identity 
\begin{equation*}
\sin ^{2}x=\frac{1-\cos \left( 2x\right) }{2}
\end{equation*}%
and then integrate over a period. 
\begin{eqnarray*}
\frac{1}{T}\int_{0}^{T}\sin ^{2}\left( \omega t\right) dt &=&\frac{1}{2T}%
\int_{0}^{T}\left( 1-\cos \left( 2\omega t\right) \right) dt \\
&=&\frac{1}{2T}\int_{0}^{T}dx-\frac{1}{2T}\int_{0}^{T}\cos \left( 2\omega
t\right) dt \\
&=&\frac{1}{2T}\left( T\right) -\frac{1}{2T}\left( 0\right) -\left. \frac{1}{%
4\omega T}\sin \left( 2\omega t\right) \right\vert _{0}^{T} \\
&=&\frac{1}{2}-\frac{1}{4\frac{2\pi }{T}T}\sin \left( 2\frac{2\pi }{T}%
T\right) -\frac{1}{4\frac{2\pi }{T}T}\sin \left( 0\right) \\
&=&\frac{1}{2}-\frac{1}{4\pi }\sin \left( 4\pi \right) \\
&=&\frac{1}{2}
\end{eqnarray*}%
just as we suspected.

We can treat our more complicated case from a history graph view and say
that $\frac{k\left( rx_{2}+r_{1}\right) }{2}+\phi _{o}$ is like a phase
constant $\alpha $ and 
\begin{equation}
\frac{1}{NT}\int_{0}^{NT}\sin ^{2}\left( -\omega t+\alpha \right) dt=\frac{1%
}{2}
\end{equation}%
because the $\alpha $ won't change our integral (it just shifts the curve on
a history graph)

So we have 
\begin{equation}
\bar{I}=\frac{1}{NT}\int_{0}^{NT}Idt\varpropto 2E_{o}^{2}\cos ^{2}\left( 
\frac{1}{2}\left( \frac{2\pi }{\lambda }d\sin \theta \right) \right)
\end{equation}%
where $\bar{I}$ is the time average intensity. The important part is that
the time varying part has averaged out. So, usually in optics, we ignore the
fast fluctuating parts of such calculations because we can't see them.

We know a lot about cosine functions. They have a maximum of $1$ and the
square of $1$ is still $1.$ So we could define a quantity that is the
maximum intensity that we could get from our average by making the cosine
squared part $1$. We don't know exactly what this will be.\footnote{%
We are still saving the fun of finding the constant of proportionality for
Principles of Physics III (PH220).} But we can see that the $2E_{o}^{2}$
will be part of it%
\begin{equation*}
I_{\max }\varpropto 2E_{o}^{2}
\end{equation*}

Using this $I_{\max }$ we can get rid of our \textquotedblleft is
proportional sign.\textquotedblright\ We can write our average intensity as%
\begin{equation*}
I=I_{\max }\cos ^{2}\left( \frac{1}{2}\left( \frac{2\pi }{\lambda }d\sin
\theta \right) \right)
\end{equation*}%
where we have dropped the bar from the $I,$ but it is understood that the
intensity we report is a time average over many periods. This averaging over
many periods is one of the reasons why we could mostly ignore the wave
nature of light when we studied mirrors and lenses!

We should remind ourselves, our intensity pattern 
\begin{equation*}
I=I_{\max }\cos ^{2}\left( \frac{\pi }{\lambda }d\sin \theta \right)
\end{equation*}%
is really 
\begin{equation*}
I=I_{\max }\cos ^{2}\left( \frac{\Delta \phi }{2}\right)
\end{equation*}%
Which is just our amplitude squared for the mixing of two waves. All we have
done to find the intensity pattern is to fill in the details for the total
phase difference $\Delta \phi .$

Our intensity pattern should give the same location for the center of the
bright spots as we got before. Let's check that it works. We used the small
angle approximation before. So let's use it again now. For for small angles%
\begin{equation*}
\frac{\pi }{\lambda }d\sin \theta \approx \frac{\pi d}{\lambda }\theta
\end{equation*}%
and our intensity is 
\begin{equation*}
I=I_{\max }\cos ^{2}\left( \frac{\pi d}{\lambda }\theta \right)
\end{equation*}%
and for small angles 
\begin{equation*}
\tan \theta =\frac{y}{L}\approx \theta
\end{equation*}%
so we can write $I$ as 
\begin{equation*}
I=I_{\max }\cos ^{2}\left( \frac{\pi d}{\lambda }\frac{y}{L}\right)
\end{equation*}%
Then we have constructive interference when $\cos ^{2}\left( \frac{\pi d}{%
\lambda }\frac{y}{L}\right) =1$ or when 
\begin{equation*}
\frac{\pi d}{\lambda }\frac{y}{L}=m\pi
\end{equation*}%
or%
\begin{equation*}
y=m\frac{L\lambda }{d}
\end{equation*}%
which is what we found before! But now we have an intensity value for every $%
y$ value, not just for the bright and dark spots.

The plot of normalized intensity 
\begin{equation*}
\frac{I}{I_{\max }}=\cos ^{2}\left( \frac{\Delta \phi }{2}\right) 
\end{equation*}%
verses $\Delta \phi /2$ is given next, \FRAME{dtbpFX}{2.292in}{1.5274in}{0pt%
}{}{}{Plot}{\special{language "Scientific Word";type "MAPLEPLOT";width
2.292in;height 1.5274in;depth 0pt;display "USEDEF";plot_snapshots
TRUE;mustRecompute FALSE;lastEngine "MuPAD";xmin "-15";xmax "15";xviewmin
"-15";xviewmax "15";yviewmin "0";yviewmax
"1";viewset"XY";rangeset"X";plottype 4;labeloverrides 3;x-label
"delta_phi/2";y-label "I /I_max";axesFont "Times New
Roman,12,0000000000,useDefault,normal";numpoints 100;plotstyle
"patch";axesstyle "normal";axestips FALSE;xis \TEXUX{x};var1name
\TEXUX{$x$};function \TEXUX{$\cos ^{2}\left( \frac{x}{2}\right) $};linecolor
"blue";linestyle 1;pointstyle "point";linethickness 1;lineAttributes
"Solid";var1range "-15,15";num-x-gridlines 100;curveColor
"[flat::RGB:0x000000ff]";curveStyle "Line";VCamFile
'TJXLEU2E.xvz';valid_file "T";tempfilename
'TJXLEM04.wmf';tempfile-properties "XPR";}}

This is really an interesting result. You might wonder why, when we found
the two slit interference pattern, there was no evidence of the single slit
fringing that we discovered in this chapter. After all, a double slit system
is made from single slits. Shouldn't there be some effect due to the fact
that the slits are individually single slits? The answer is that we did see
some hint of the single slit pattern. Look at the figure below. \FRAME{dhF}{%
4.4719in}{2.2606in}{0pt}{}{}{Figure}{\special{language "Scientific
Word";type "GRAPHIC";maintain-aspect-ratio TRUE;display "USEDEF";valid_file
"T";width 4.4719in;height 2.2606in;depth 0pt;original-width
4.4209in;original-height 2.2208in;cropleft "0";croptop "1";cropright
"1";cropbottom "0";tempfilename 'SUS73V53.wmf';tempfile-properties "XPR";}}%
The intensity of the peaks seems to fall off with distance from the center.
We dealt with only the center-most part of the pattern. If we draw the
pattern for larger angles, we see the following.\FRAME{dtbpF}{4.1883in}{%
2.3497in}{0pt}{}{}{Figure}{\special{language "Scientific Word";type
"GRAPHIC";maintain-aspect-ratio TRUE;display "USEDEF";valid_file "T";width
4.1883in;height 2.3497in;depth 0pt;original-width 4.1381in;original-height
2.3099in;cropleft "0";croptop "1";cropright "1";cropbottom "0";tempfilename
'SUS73V54.wmf';tempfile-properties "XPR";}}

It takes a bright laser or dark room to see the secondary groups of fringes
easily, but we can do it. We can also graph the intensity pattern. It is the
combination of the two slit and single slit pattern with the single slit
pattern acting as an envelope.

\begin{equation}
I=I_{\max }\cos ^{2}\left( \frac{\pi d\sin \left( \theta \right) }{\lambda }%
\right) \left( \frac{\sin \left( \frac{\pi a\sin \left( \theta \right) }{%
\lambda }\right) }{\frac{\pi a\sin \left( \theta \right) }{\lambda }}\right)
^{2}
\end{equation}%
Note that one of the double slit maxima is clobbered by a minimum in the
single slit pattern. We can find out the order of the missing maximum.
Recall that 
\begin{equation*}
d\sin \left( \theta \right) =m\lambda
\end{equation*}%
describes the maxima from the double slit. But for the single slit 
\begin{equation*}
a\sin \left( \theta \right) =\lambda
\end{equation*}%
describes the first minimum. Dividing these yields%
\begin{eqnarray*}
\frac{d\sin \left( \theta \right) }{a\sin \left( \theta \right) } &=&\frac{%
m\lambda }{\lambda } \\
\frac{d}{a} &=&m
\end{eqnarray*}%
so the 
\begin{equation}
m=\frac{d}{a}
\end{equation}%
double slit maximum will be missing.

\section{Diffraction Gratings}

We now have found the single slit diffraction pattern equation and the
double slit diffraction pattern. Let's look once again at many slits. if we
have a whole lot of slits we give the collection of slits a new name. We
call it a \emph{diffraction grating.}

A diffraction grating is an optical element with many many parallel slits
spaced very close together. Here is a photo of part of a typical diffraction
grating created by etching lines in a piece of glass. The etchings scatter
the light, but the un-etched part allows the light to pass through. The
un-etched parts are essentially a series of slits.\FRAME{dhFU}{1.9984in}{%
1.625in}{0pt}{\Qcb{Surface of a diffraction grating (600 lines/mm). Image
taken with optical transmission microscope. (Image in the public domain
courtesy Scapha)}}{}{Figure}{\special{language "Scientific Word";type
"GRAPHIC";maintain-aspect-ratio TRUE;display "USEDEF";valid_file "T";width
1.9984in;height 1.625in;depth 0pt;original-width 2.4474in;original-height
1.9847in;cropleft "0";croptop "1";cropright "1";cropbottom "0";tempfilename
'SVANNN5G.wmf';tempfile-properties "XPR";}}A typical grating might have $%
5000 $ slits per unit centimeter. You have probably used a diffraction
grating to see rainbow colors in a beginning science class. Gratings are
usually made by cutting parallel grooves in a flat surface.

If we use $5000\frac{grooves}{\unit{cm}}$ for an example, we see that the
slit spacing is 
\begin{eqnarray}
d &=&\frac{1}{5000}\unit{cm} \\
&=&\allowbreak 2.0\times 10^{-6}\unit{m}
\end{eqnarray}%
Take a section of diffraction grating as shown below\FRAME{dhF}{1.9951in}{%
3.3788in}{0in}{}{}{Figure}{\special{language "Scientific Word";type
"GRAPHIC";maintain-aspect-ratio TRUE;display "USEDEF";valid_file "T";width
1.9951in;height 3.3788in;depth 0in;original-width 5.5893in;original-height
9.5207in;cropleft "0";croptop "1";cropright "1";cropbottom "0";tempfilename
'SVANNN5H.wmf';tempfile-properties "XPR";}}

At some point, two of the slits will have a path difference that is a whole
wavelength, and we would expect a bright spot. But what about the other
slits? If we have a slit spacing such that each of the succeeding slits has
a path difference that is just an additional wavelength, then each of the
slits will contribute to the constructive interference at our point, and the
point will become a bright spot.

\FRAME{dhF}{0.9755in}{2.3679in}{0pt}{}{}{Figure}{\special{language
"Scientific Word";type "GRAPHIC";maintain-aspect-ratio TRUE;display
"USEDEF";valid_file "T";width 0.9755in;height 2.3679in;depth
0pt;original-width 1.6821in;original-height 4.1217in;cropleft "0";croptop
"1";cropright "1";cropbottom "0";tempfilename
'SVANNN5I.wmf';tempfile-properties "XPR";}}Let's look at just two slits. The
light leaves each slit in phase with the light from the rest of the slits,
but at some distance $L$ away and at some angle $\theta $ we will have a
path difference%
\begin{eqnarray*}
\Delta \phi &=&\left( \frac{2\pi }{\lambda }\Delta r+\Delta \phi _{o}\right)
=m2\pi \qquad m=0,\pm 1,\pm 2,\pm 3,\cdots \text{ Constructive} \\
\Delta \phi &=&\left( \frac{2\pi }{\lambda }\Delta r+\Delta \phi _{o}\right)
=\left( 2m+1\right) \pi \qquad m=0,\pm 1,\pm 2,\pm 3,\cdots \text{
Destructive}
\end{eqnarray*}%
\begin{equation*}
\left( \frac{2\pi }{\lambda }\Delta r+\Delta \phi _{o}\right) =m2\pi \qquad
m=0,\pm 1,\pm 2,\pm 3
\end{equation*}%
\begin{equation*}
\left( \frac{2\pi }{\lambda }\left( d\sin \left( \theta _{bright}\right)
\right) +\Delta \phi _{o}\right) =m2\pi \qquad m=0,\pm 1,\pm 2,\pm 3
\end{equation*}%
Once again we start all the waves the same so $\Delta \phi _{o}=0$ 
\begin{equation*}
\left( \left( d\sin \left( \theta _{bright}\right) \right) +0\right)
=m\lambda \qquad m=0,\pm 1,\pm 2,\pm 3
\end{equation*}%
\begin{equation}
d\sin \left( \theta _{bright}\right) =m\lambda \quad m=0,\pm 1,\pm 2,\ldots
\end{equation}%
All because the path lengths are not all the same.

This equation tells us that each wavelength, $\lambda ,$ will experience
constructive interference at a slightly different angle $\theta _{bright}.$
Different frequencies will create bright spots at different angles. We have
found a way to create a spectrum with light waves. We often call a spectrum
made with visible light frequencies a rainbow. Knowing $d$ and $\theta $
allows an accurate calculation of $\lambda .$ This may seem a silly thing to
do, but suppose we add into our system a sample of a chemical to identify 
\FRAME{dhF}{3.5958in}{1.7465in}{0pt}{}{}{Figure}{\special{language
"Scientific Word";type "GRAPHIC";maintain-aspect-ratio TRUE;display
"USEDEF";valid_file "T";width 3.5958in;height 1.7465in;depth
0pt;original-width 3.6357in;original-height 1.7509in;cropleft "0";croptop
"1";cropright "1";cropbottom "0";tempfilename
'SVANNN5J.wmf';tempfile-properties "XPR";}}

We could then record the intensity of the transmitted light as a function of
angle, which is equivalent to $\lambda .$ We can again generate a spectrum.
This is a traditional way to build a spectrometer and many such devices are
available today.

\subsection{Resolving power of diffraction gratings}

We noticed that with two slits, we got a bright spot for a particular
wavelength, but we didn't just get one bright spot. We got several. The same
is true for diffraction gratings. We expect to get a rainbow, but we really
expect to get a series of rainbows. The integer $m$ tells us which rainbow
we have in the series. The integer $m$ is called the order number.

For two nearly equal wavelengths $\lambda _{1}$ and $\lambda _{2},$ we say
that the diffraction grating can resolve the wavelengths if we can
distinguish the two using the grating. The \emph{resolving power} of the
grating is defined as 
\begin{equation}
R=\frac{\left( \lambda _{1}+\lambda _{2}\right) }{2\left( \lambda
_{1}-\lambda _{2}\right) }=\frac{\bar{\lambda}}{\Delta \lambda }
\end{equation}%
It can be shown that for the $m$-th order diffraction, the resolving power is%
\begin{equation}
R=Nm
\end{equation}%
where $N$ is the number of slits. So our ability to distinguish wavelengths
increases with the number of slits and with the order (which is related to
how far off-axis we look).

Note that for $m=0$ we have no ability to resolve wavelengths. The central
peak is a mix of all wavelengths and usually looks white for normal
illumination.

That the resolution depends on the number of slits, $N,$ means that we can
improve our spectrometer by using more lines. Here is a representation of
what happens as we increase $N$\FRAME{dhF}{3.103in}{2.6654in}{0pt}{}{}{Figure%
}{\special{language "Scientific Word";type "GRAPHIC";maintain-aspect-ratio
TRUE;display "USEDEF";valid_file "T";width 3.103in;height 2.6654in;depth
0pt;original-width 3.0588in;original-height 2.623in;cropleft "0";croptop
"1";cropright "1";cropbottom "0";tempfilename
'SVANNN5K.wmf';tempfile-properties "XPR";}}we can see that the peaks get
narrower as $N$ increases. These graphs are for a particular $\lambda .$ If
the peaks for a particular $\lambda $ get narrower, then there will be less
overlap with adjacent $\lambda ^{\prime }s$ which means that each wavelength
can more easily be distinguished from the others. This is what we mean by
\textquotedblleft resolved.\textquotedblright\ We won't derive the equation
for a diffraction grating intensity pattern in this class. If you are
interested, you can take our junior level optics class. From that class the
intensity pattern is 
\begin{equation*}
I\left( \theta \right) =I_{o}\left( \frac{\sin \beta }{\beta }\right)
^{2}\left( \frac{\sin N\alpha }{\sin \alpha }\right) ^{2}
\end{equation*}%
where $\alpha =\frac{dk}{2}\sin \theta ,$ $\beta =\frac{ka}{2}\sin \theta $
and where $a$ is the slit width and $d$ is the slit spacing and $N$ is the
number of slits.

Spectrometers are used in many places. One that has some public interest
today is monitoring the atmosphere. Instruments like the one shown below
detect the amount of special gasses in the atmosphere using IR spectrometers.

\FRAME{dhFU}{4.6544in}{3.0554in}{0pt}{\Qcb{AIRS\ sensor, spectrometer
design, and global CO$_{2}$ data. (Images in the Public Domain courtesy NASA)%
}}{}{Figure}{\special{language "Scientific Word";type
"GRAPHIC";maintain-aspect-ratio TRUE;display "USEDEF";valid_file "T";width
4.6544in;height 3.0554in;depth 0pt;original-width 4.6017in;original-height
3.0113in;cropleft "0";croptop "1";cropright "1";cropbottom "0";tempfilename
'SVAQIE62.wmf';tempfile-properties "XPR";}}The instrument shown is the AIRS
spectrometer. You can see in the diagram that it uses a grating
spectrometer. The picture of the Earth is a composite of AIRS data showing
the northern and southern bands of CO$_{2}.$

We have come along way in understanding the effects of the wave nature of
light and even found some practical uses for it. But we still don't know how
it affects our ability to make images. Let's take that on in our next
lecture.

\chapter{23 Resolution 3.4.5, 3.4.6, 3.4.7}

%TCIMACRO{\TeXButton{\chaptermark{Resolution}}{\chaptermark{Resolution}}}%
%BeginExpansion
\chaptermark{Resolution}%
%EndExpansion

We have talked about how the wave nature of light can be used in practical
systems like spectrometers and interferometers. But it turns out that the
wave nature of light can't be ignored even for our lens systems like
magnifiers, cameras, and telescopes. In this lecture, let's look at how the
wave nature of light affects imagery. And then, let's discuss when and when
not to expect to see the wave nature of light.

%TCIMACRO{%
%\TeXButton{Fundamental Concepts}{\vspace{0.10in}\hspace{-0.37in}{\Large Fundamental Concepts\vspace{0.2in}}}}%
%BeginExpansion
\vspace{0.10in}\hspace{-0.37in}{\Large Fundamental Concepts\vspace{0.2in}}%
%EndExpansion

\begin{itemize}
\item Circular Apertures

\item Resolution is a name we give to the fundamental blurriness in
geometrical optical systems due to the wave nature of light.

\item Interference from circular apertures

\item X-ray Diffraction

\item Holography
\end{itemize}

\section{Circular apertures}

Our analysis of light going through holes has been somewhat limited by
squarish holes or slits. But most optical systems, including our eyes don't
have square holes. So what happens when the hole is round? The situation is
as shown in the next figure.\FRAME{dhF}{2.341in}{1.2842in}{0pt}{}{}{Figure}{%
\special{language "Scientific Word";type "GRAPHIC";maintain-aspect-ratio
TRUE;display "USEDEF";valid_file "T";width 2.341in;height 1.2842in;depth
0pt;original-width 4.4763in;original-height 2.4431in;cropleft "0";croptop
"1";cropright "1";cropbottom "0";tempfilename
'SUS73V55.wmf';tempfile-properties "XPR";}}Before we discuss this situation,
let's think about the width of a single slit pattern. We remember that 
\begin{equation*}
\sin \left( \theta \right) =\left( 1\right) \frac{\lambda }{a}
\end{equation*}%
for the first minima,or that 
\begin{equation*}
\theta \approx \frac{\lambda }{a}
\end{equation*}%
for small angles. And from the figure \FRAME{dtbpF}{3.2353in}{1.8187in}{0pt}{%
}{}{Figure}{\special{language "Scientific Word";type
"GRAPHIC";maintain-aspect-ratio TRUE;display "USEDEF";valid_file "T";width
3.2353in;height 1.8187in;depth 0pt;original-width 3.9565in;original-height
2.2122in;cropleft "0";croptop "1";cropright "1";cropbottom "0";tempfilename
'SXEIGM01.wmf';tempfile-properties "XPR";}}we can see that 
\begin{equation*}
\theta \approx \frac{y}{L}
\end{equation*}%
so long as $\theta $ is small, then we find the position of the first
minimum to be%
\begin{equation*}
y\approx \frac{\lambda }{a}L
\end{equation*}%
This is the distance from the center bright spot peak intensity to the first
dark spot. The width of the bright spot is twice this distance%
\begin{equation*}
w\approx 2\frac{\lambda }{a}L
\end{equation*}%
We expect something like this for our circular aperture. The derivation is
not to hard, but it involves Bessel functions, which are beyond the math
requirement for this course\footnote{%
We find the whole diffraction pattern for a circular aperture in PH375, 
\emph{Optics.}}. So I will give you the answer%
\begin{equation*}
\theta \approx 1.22\frac{\lambda }{D}
\end{equation*}%
Where we have replaced $a$ with $D$, the diameter of the aperture. But they
are both the size of the hole.

Let's think about this last equation. Compare it to our equation for $\theta 
$ for a square opening. That's right, the circular aperture (hole) only adds
a factor of $1.22$. And as with the slit%
\begin{equation*}
\theta \approx \frac{y}{L}
\end{equation*}%
so%
\begin{equation}
y\approx 1.22\frac{\lambda }{D}L
\end{equation}%
and 
\begin{equation}
w\approx 2.44\frac{\lambda }{D}L
\end{equation}%
The circular aperture pattern looks a little like the slit pattern. But the
secondary maxima are actually smaller for the circular aperture case. Here
is a larger version in red (solid curve) with the single slit pattern in
brown (dashed line).\FRAME{dtbpFUX}{4.907in}{2.7415in}{0pt}{\Qcb{{}}}{}{Plot%
}{\special{language "Scientific Word";type "MAPLEPLOT";width 4.907in;height
2.7415in;depth 0pt;display "USEDEF";plot_snapshots TRUE;mustRecompute
FALSE;lastEngine "MuPAD";xmin "-10.00400";xmax "10.00400";xviewmin
"-10.0060008020004";xviewmax "10.0060008020004";yviewmin
"-0.000098808010947";yviewmax "1.00009999004583";rangeset"X";plottype
4;labeloverrides 3;x-label "Beta";y-label "I";axesFont "Times New
Roman,12,0000000000,useDefault,normal";numpoints 100;plotstyle
"patch";axesstyle "normal";axestips FALSE;xis \TEXUX{r};var1name
\TEXUX{$r$};function \TEXUX{$\frac{1}{\allowbreak 0.25}\left( \text{
}\frac{J_{1}\left( r\right) }{r}\right) ^{2}$};linecolor "blue";linestyle
1;pointstyle "point";linethickness 3;lineAttributes "Solid";var1range
"-10.00400,10.00400";num-x-gridlines 100;curveColor
"[flat::RGB:0x000000ff]";curveStyle "Line";function \TEXUX{$\left(
\frac{\sin \left( r\right) }{r}\right) ^{2}$};linecolor "cyan";linestyle
2;pointstyle "point";linethickness 1;lineAttributes "Dash";var1range
"-10.00400,10.00400";num-x-gridlines 100;curveColor
"[flat::RGB:0x000080c0]";curveStyle "Line";VCamFile
'TJXLCV22.xvz';valid_file "T";tempfilename
'SXEJ1N07.wmf';tempfile-properties "XPR";}}

Here is a three dimensional version of the intensity pattern from the
circular aperture.

\FRAME{dtbpFU}{4.1598in}{2.7758in}{0pt}{\Qcb{{\protect\small Graph of the
Airy disk for the case }$a=0.05\unit{mm},${\protect\small \ }$\protect%
\lambda =500\unit{nm},${\protect\small \ and }$R=1\unit{m}.$}}{}{Figure}{%
\special{language "Scientific Word";type "GRAPHIC";maintain-aspect-ratio
TRUE;display "USEDEF";valid_file "T";width 4.1598in;height 2.7758in;depth
0pt;original-width 8.235in;original-height 5.4819in;cropleft "0";croptop
"1";cropright "1";cropbottom "0";tempfilename
'TGDU0X00.wmf';tempfile-properties "XPR";}}With a bright enough laser, the
pattern becomes visible.\FRAME{dhFU}{2.7415in}{2.6757in}{0pt}{\Qcb{{}}}{}{%
Figure}{\special{language "Scientific Word";type
"GRAPHIC";maintain-aspect-ratio TRUE;display "USEDEF";valid_file "T";width
2.7415in;height 2.6757in;depth 0pt;original-width 1.0404in;original-height
1.0136in;cropleft "0";croptop "1";cropright "1";cropbottom "0";tempfilename
'SUS73V58.wmf';tempfile-properties "XPR";}}But notice that the center of the
pattern looks white. That is because the sensor I used is saturating in the
center. There is too much light there so it can't tell the light is red
anymore. But around the central maximum we can see the faint circles that
are the other small maxima. These are usually so faint that we just ignore
them, but it's good to know they are there.

\section{Resolution}

We have emphasized that an extended object can be viewed as a collection of
point objects. Then the image is formed from the collection if images of
those point objects. Each point on the image is a copy of the aperture or
hole the light traveled through. The image is made up of many overlapping
aperture shapes. In our present case, the aperture shape is a circle, if we
ignore the faint secondary maxima.

We thought about the design of cameras and telescopes and other optical
systems using ray optics. It would be great if optical systems could form
images with infinite precision, but it turns out they can't. And it is the
fact that light acts as a wave prevents this from being true! Our wave
nature of light comes back to complicate our simple optical designs.

Because light is really a wave, the images of the point objects won't be
points, themselves. We just saw that a point of light makes a circle of
light with little rings around it. And because our eyes and cameras have
circular apertures, the image of the point objects will be little circular
central maxima with dim concentric ring patterns. \FRAME{dtbpF}{1.8792in}{%
1.7374in}{0pt}{}{}{Figure}{\special{language "Scientific Word";type
"GRAPHIC";maintain-aspect-ratio TRUE;display "USEDEF";valid_file "T";width
1.8792in;height 1.7374in;depth 0pt;original-width 1.8421in;original-height
1.7002in;cropleft "0";croptop "1";cropright "1";cropbottom "0";tempfilename
'SVAQIE63.wmf';tempfile-properties "XPR";}} If we have more than one point
of light, we will get two such patterns. \FRAME{dtbpF}{1.5281in}{1.2851in}{%
0pt}{}{}{Figure}{\special{language "Scientific Word";type
"GRAPHIC";maintain-aspect-ratio TRUE;display "USEDEF";valid_file "T";width
1.5281in;height 1.2851in;depth 0pt;original-width 1.4909in;original-height
1.2496in;cropleft "0";croptop "1";cropright "1";cropbottom "0";tempfilename
'SVAQIE64.wmf';tempfile-properties "XPR";}}and our image of an extended
object (like Aunt Agatha) is made up of many, many points all reflecting
light into our camera lens. We want each of those points of light to create
corresponding points on our image. But instead they are creating small
circles of light, and from the above figures we can see that those circles
will overlap. All this makes our images blurry.

\section{Resolution}

The quality of our image depends on how poorly a point object is imaged. If
each point object makes a large circle of light on the screen or detector
array, we get a very confusing image (it will look blurry to us). Let's
review why this will happen so we can know how to minimize the effect. Since
we are using eyes or cameras, we expect circular apertures. Let's review
what we know but include some ray diagrams.

We already know that if we take light and pass it through a single slit, we
get an intensity pattern that has a central bright region.\FRAME{dhF}{%
2.9914in}{2.2191in}{0pt}{}{}{Figure}{\special{language "Scientific
Word";type "GRAPHIC";maintain-aspect-ratio TRUE;display "USEDEF";valid_file
"T";width 2.9914in;height 2.2191in;depth 0pt;original-width
2.9473in;original-height 2.1785in;cropleft "0";croptop "1";cropright
"1";cropbottom "0";tempfilename 'SVAQIE65.wmf';tempfile-properties "XPR";}}

Remember that normal objects will be made up of many small points of light
(either due to reflection or glowing) and each of these will form such an
intensity pattern on a screen. Here is a bright point source that is not on
the axis, and we see that it too makes a bright spot on the screen (and
smaller bright spots or rings, depending on the shape of the aperture) 
\FRAME{dhF}{2.7674in}{2.0513in}{0pt}{}{}{Figure}{\special{language
"Scientific Word";type "GRAPHIC";maintain-aspect-ratio TRUE;display
"USEDEF";valid_file "T";width 2.7674in;height 2.0513in;depth
0pt;original-width 2.725in;original-height 2.0124in;cropleft "0";croptop
"1";cropright "1";cropbottom "0";tempfilename
'SVAQIE66.wmf';tempfile-properties "XPR";}}So our images will be made up of
many central bright spots, each of which represents a point of light from
the object. These central bright spots may overlap, (and their secondary
maxima certainly will overlap).

Let's take a simple case of two points of light, $S_{1}$ and $S_{2}.$ If we
take a single slit and pass light from two distant point sources through the
slit, we do not get two sharp images of the point sources. Instead, we get
two diffraction patterns.\FRAME{dhF}{3.3537in}{2.6507in}{0pt}{}{}{Figure}{%
\special{language "Scientific Word";type "GRAPHIC";maintain-aspect-ratio
TRUE;display "USEDEF";valid_file "T";width 3.3537in;height 2.6507in;depth
0pt;original-width 3.3088in;original-height 2.6091in;cropleft "0";croptop
"1";cropright "1";cropbottom "0";tempfilename
'SVAQIE67.wmf';tempfile-properties "XPR";}}

If these patterns are formed sufficiently far from each other, it is easy to
tell they were formed from two distinct objects. Each point became a small
circular blur, but that is really not so bad. We can still tell that the two
blurs came from different sources. If our pixel size is about the same size
of the blur, we may not even notice the blurriness in the digital imagery.%
\FRAME{dhF}{3.6487in}{2.8738in}{0pt}{}{}{Figure}{\special{language
"Scientific Word";type "GRAPHIC";maintain-aspect-ratio TRUE;display
"USEDEF";valid_file "T";width 3.6487in;height 2.8738in;depth
0pt;original-width 3.6011in;original-height 2.8305in;cropleft "0";croptop
"1";cropright "1";cropbottom "0";tempfilename
'SVAQIE68.wmf';tempfile-properties "XPR";}} But if the patterns are formed
close to each other, it gets hard to tell whether they were formed from two
objects or one bright object. We now have a problem.

Suppose you are trying to look at a star and see if it has a planet. But all
you can see is a blur. You can't tell if there is one source of light or two.

Long ago an early researcher titled Lord Rayleigh developed a test to
determine if you can distinguish between two diffraction patterns. When the
central maximum of one point's image falls on the first minimum of anther
point's image, the images are said to be just resolved. This test is known
as \emph{Rayleigh's criterion}.

We can find the required separation for a slit. Remember that 
\begin{equation}
\sin \left( \theta \right) =m\frac{\lambda }{a}\quad m=\pm 1,\pm 2,\pm
3\ldots
\end{equation}%
gives the minima (dark spots) for a single slit. We want the first minimum,
so $m=1$%
\begin{equation}
\sin \left( \theta \right) =\frac{\lambda }{a}
\end{equation}%
Remember that 
\begin{equation*}
\tan \theta =\frac{y}{L}
\end{equation*}%
where $L$ is the distance from the slit to the viewing screen, and $y$ is
the vertical location of the dark spot. If we place the second image maximum
so it is just at this location, the two images will be just barely
resolvable. In the small angle approximation, $\sin \left( \theta \right)
\approx \theta $ so%
\begin{equation}
\theta _{\min }=\frac{\lambda }{a}
\end{equation}

Now you may be saying to yourself that you don't often take pictures through
single illuminated slits, so this is nice, but not really very helpful. But
suppose, instead, that we image a circular aperture. We just pick up a
factor of $1.22$ and change $a$ to $D$ for diameter.%
\begin{equation}
\theta _{\min }=1.22\frac{\lambda }{D}
\end{equation}

The Rayleigh criteria tells you, based on your camera aperture size, how a
point source will be imaged on the film or sensor array. If we consider
extended sources (like your favorite car, or Aunt Agatha) to be collections
of many point sources, then we have a way to tell what features will be
clearly resolved on the image and what features will not (like you may not
be able to see the lettering on the car to tell what model it is, or you may
not be able to distinguish between the gem stones in Aunt Agatha's necklace
because the image is too blurry to see these features clearly).

\FRAME{dtbpFU}{3.6893in}{1.4105in}{0pt}{\Qcb{{\protect\small Pattern from
two resolved circular slits.}}}{}{Figure}{\special{language "Scientific
Word";type "GRAPHIC";maintain-aspect-ratio TRUE;display "USEDEF";valid_file
"T";width 3.6893in;height 1.4105in;depth 0pt;original-width
3.6417in;original-height 1.375in;cropleft "0";croptop "1";cropright
"1";cropbottom "0";tempfilename 'SVAQIE69.wmf';tempfile-properties "XPR";}}%
\FRAME{dtbpFU}{3.6898in}{1.3596in}{0pt}{\Qcb{{\protect\small Rayleigh
Criteria: Pattern from two circular soruces where the sources are close
enough that the maximum from one pattern is placed on the minimum of the
other. Lord Rayleigh gave this as the criteria for just barily being
resolved.}}}{}{Figure}{\special{language "Scientific Word";type
"GRAPHIC";maintain-aspect-ratio TRUE;display "USEDEF";valid_file "T";width
3.6898in;height 1.3596in;depth 0pt;original-width 3.6417in;original-height
1.3249in;cropleft "0";croptop "1";cropright "1";cropbottom "0";tempfilename
'SVAQIE6A.wmf';tempfile-properties "XPR";}}Astronomers sometimes use
Sparrow's criteria for two sources being resolved. It is shown below.\FRAME{%
dtbpFU}{3.6893in}{1.3595in}{0pt}{\Qcb{Sparrow Criteria: This is a less
concervative resolution criteria than Rayleigh. When the intenisty pattern
is flat on the top, there must be two sources. This criterial is used in
astronomy.}}{}{Figure}{\special{language "Scientific Word";type
"GRAPHIC";maintain-aspect-ratio TRUE;display "USEDEF";valid_file "T";width
3.6893in;height 1.3595in;depth 0pt;original-width 3.6417in;original-height
1.3249in;cropleft "0";croptop "1";cropright "1";cropbottom "0";tempfilename
'SVAQIE6B.wmf';tempfile-properties "XPR";}}Since astronomers just want to
know if there are two stars or one, it is enough to see that the intensity
pattern went flat at the center. That must mean that there are two stars.
But if the stars are any closer, the flat center becomes a peak and the
stars are unresolved.

\FRAME{dtbpFU}{3.6893in}{1.3595in}{0pt}{\Qcb{Two circular sources unresolved}%
}{}{Figure}{\special{language "Scientific Word";type
"GRAPHIC";maintain-aspect-ratio TRUE;display "USEDEF";valid_file "T";width
3.6893in;height 1.3595in;depth 0pt;original-width 3.6417in;original-height
1.3249in;cropleft "0";croptop "1";cropright "1";cropbottom "0";tempfilename
'SVAQIE6C.wmf';tempfile-properties "XPR";}}Here is what the easily resolved,
Rayleigh resolved, Sparrow resolved, and unresolved cases look like. \FRAME{%
dtbpF}{4.4171in}{1.3753in}{0pt}{}{}{Figure}{\special{language "Scientific
Word";type "GRAPHIC";maintain-aspect-ratio TRUE;display "USEDEF";valid_file
"T";width 4.4171in;height 1.3753in;depth 0pt;original-width
5.7749in;original-height 1.7782in;cropleft "0";croptop "1";cropright
"1";cropbottom "0";tempfilename 'TGFD2L02.wmf';tempfile-properties "XPR";}}

Of course, astronomers have it easy. They really image points of light. But
for photos of more normal things we don't want to not be able to tell part
of Aunt Sally's eye from part of her nose. Our aperture limits how good the
image will look regardless of the lens that we use in the camera. Because
this limit is due to diffraction patterns we call this resolution limit a 
\emph{diffraction limit}. If you are designing a camera system you generally
want your sensor pixel size to match the size of the central diffraction
spot, because smaller pixels won't get you a less blurry image once you
reach the size of the spot created by imaging a point source.

This really concludes our study of waves and optics! But let's look at two
more applications of what we have learned.

\section{Holography}

You may have seen holograms in the past. We have enough understanding of
light to understand how they are generated now.

\FRAME{dtbpF}{2.853in}{1.7772in}{0pt}{}{}{Figure}{\special{language
"Scientific Word";type "GRAPHIC";maintain-aspect-ratio TRUE;display
"USEDEF";valid_file "T";width 2.853in;height 1.7772in;depth
0pt;original-width 8.0834in;original-height 5.0246in;cropleft "0";croptop
"1";cropright "1";cropbottom "0";tempfilename
'SUS73V5B.wmf';tempfile-properties "XPR";}}

A device for generating a hologram is shown in the figure above. Light from
a laser or other coherent source is expanded and split into two beams. One
travels to a photographic plate, the other is directed to an object. At the
object, light is scattered and the scattered light also reaches the
photographic plate. The combination of the direct and scattered beams
generates a complicated interference pattern.

The pattern can be developed (like you develop photographic film). Once
developed, it can be re-illuminated with a direct beam. The emulsion on the
plate creates complicated patterns of light transmission, which combine to
create interference. It is like a very complicated slit pattern or grating
pattern. The result is a three-dimensional virtual image generated by the
interference. The interference pattern generates an image that looks like
the original object.\FRAME{dtbpF}{2.687in}{1.689in}{0pt}{}{}{Figure}{\special%
{language "Scientific Word";type "GRAPHIC";display "USEDEF";valid_file
"T";width 2.687in;height 1.689in;depth 0pt;original-width
8.0834in;original-height 5.0246in;cropleft "0";croptop "1";cropright
"1";cropbottom "0";tempfilename 'SUS73V5C.wmf';tempfile-properties "XPR";}}

\section{Diffraction of X-rays by Crystals}

If we make the wavelength of light very small, then we can deal with very
small diffraction gratings. This concept is used to investigate the
structure of crystals with x-rays. The crystal latus of molecules or atoms
creates the regular pattern we need for a grating. The pattern is three
dimensional, so the patterns are complex.

Let's start with a simple crystal with a square regular latus. Salt, $NaCl,$
has such a structure.

\FRAME{dtbpF}{3.9156in}{1.7407in}{0pt}{}{}{Figure}{\special{language
"Scientific Word";type "GRAPHIC";maintain-aspect-ratio TRUE;display
"USEDEF";valid_file "T";width 3.9156in;height 1.7407in;depth
0pt;original-width 4.7493in;original-height 2.0965in;cropleft "0";croptop
"1";cropright "1";cropbottom "0";tempfilename
'TGFD1J01.wmf';tempfile-properties "XPR";}}For constructive interference we
need 
\begin{equation*}
\Delta \phi =\left( \frac{2\pi }{\lambda }\Delta r+\Delta \phi _{o}\right)
=m2\pi \qquad m=0,\pm 1,\pm 2,\pm 3,\cdots \text{ Constructive}
\end{equation*}%
If we illuminate the crystal with x-rays, the x-rays can reflect off the top
layer of atoms, or off the second layer of atoms (or off any other layer,
but for now let's just consider two layers). At each reflection we should
get a $\phi _{o}=\pi $ so $\Delta \phi _{o}=\pi -\pi =0.$ But the path for
the two rays is different, so we need to find $\Delta r$ such that we have
constructive interference. If the spacing between the layers is $d,$ then
the path difference will be 
\begin{equation}
\Delta r=2\left( d\sin \left( \theta \right) \right)
\end{equation}%
then for constructive interference%
\begin{equation*}
\left( \frac{2\pi }{\lambda }\left( 2\left( d\sin \left( \theta \right)
\right) \right) +0\right) =m2\pi
\end{equation*}%
\begin{equation*}
\frac{1}{\lambda }\left( 2\left( d\sin \left( \theta \right) \right) \right)
=m
\end{equation*}%
\begin{equation*}
2d\sin \left( \theta \right) =m\lambda
\end{equation*}%
So then for constructive interference%
\begin{equation}
2d\sin \left( \theta \right) =m\lambda \qquad m=1,2,3,\ldots
\end{equation}%
This is known as \emph{Bragg's law.} This relationship can be used to
measure the distance between the crystal planes. Here is a schematic drawing
of a device that does this.

\FRAME{dtbpF}{3.5755in}{2.2709in}{0in}{}{}{Figure}{\special{language
"Scientific Word";type "GRAPHIC";maintain-aspect-ratio TRUE;display
"USEDEF";valid_file "T";width 3.5755in;height 2.2709in;depth
0in;original-width 3.5598in;original-height 2.2508in;cropleft "0";croptop
"1";cropright "1";cropbottom "0";tempfilename
'TGDYPM02.wmf';tempfile-properties "XPR";}}As the crystal structure gets
more complicated, the interference pattern gets more complicated as well.
Here is an example

\FRAME{dhFU}{1.6491in}{1.6604in}{0pt}{\Qcb{{\protect\small Diffraction image
of protein crystal. Hen egg lysozyme, X-ray souce Bruker I}$\protect\mu $%
{\protect\small S, }$\protect\lambda =0.154188\unit{nm}${\protect\small , }$%
45\unit{kV}${\protect\small , Exposure }$10\unit{s}.${\protect\small \
(image in the public domain)}}}{}{Figure}{\special{language "Scientific
Word";type "GRAPHIC";maintain-aspect-ratio TRUE;display "USEDEF";valid_file
"T";width 1.6491in;height 1.6604in;depth 0pt;original-width
1.6129in;original-height 1.6233in;cropleft "0";croptop "1";cropright
"1";cropbottom "0";tempfilename 'SUS73V5E.wmf';tempfile-properties "XPR";}}%
DNA\ makes in interesting diffraction pattern.\FRAME{dhFU}{1.5811in}{1.8279in%
}{0pt}{\Qcb{{\protect\small X-ray diffraction pattern of DNA (image courtesy
of the National Institute of Health, image in the public domain)}}}{}{Figure%
}{\special{language "Scientific Word";type "GRAPHIC";maintain-aspect-ratio
TRUE;display "USEDEF";valid_file "T";width 1.5811in;height 1.8279in;depth
0pt;original-width 1.5437in;original-height 1.7902in;cropleft "0";croptop
"1";cropright "1";cropbottom "0";tempfilename
'SUS73V5F.wmf';tempfile-properties "XPR";}}

\section{Ray Model vs. Wave Model}

We have studied all the new information that we will get on waves and
optics. Notice that we started our study of light using rays, but knowing
that light is a wave. We did most of our work with mirrors and lenses
without using wave properties of light. But then we looked at apertures and
found the wave nature of light was sometimes important. When do we need to
use the equations for light as a wave, and when can we get away with
ignoring the wave nature of light?

In the next figure, two waves of different wavelengths go through a single
opening. The wave representing the central maxima is shown in each case, but
not the secondary maxima.\FRAME{dhF}{2.7951in}{2.776in}{0pt}{}{}{Figure}{%
\special{language "Scientific Word";type "GRAPHIC";maintain-aspect-ratio
TRUE;display "USEDEF";valid_file "T";width 2.7951in;height 2.776in;depth
0pt;original-width 2.7527in;original-height 2.7337in;cropleft "0";croptop
"1";cropright "1";cropbottom "0";tempfilename
'SUS73V5G.wmf';tempfile-properties "XPR";}}Notice that the smaller
wavelength has a narrower central maxima as we would expect from%
\begin{equation*}
\sin \left( \theta _{dark}\right) =1.22\frac{\lambda }{D}
\end{equation*}%
or 
\begin{equation*}
\theta _{dark}\approx 1.22\frac{\lambda }{D}
\end{equation*}%
We see that the ratio of the wavelength to the hole size determines the
angular extent of the central maxima. The smaller the ratio, the smaller the
central region. We can use this to explain why the wave nature of light was
so hard to find.

The patch of light on a screen that is created by light passing through the
aperture is created by the central maximum. \FRAME{dhF}{3.0191in}{2.2468in}{%
0pt}{}{}{Figure}{\special{language "Scientific Word";type
"GRAPHIC";maintain-aspect-ratio TRUE;display "USEDEF";valid_file "T";width
3.0191in;height 2.2468in;depth 0pt;original-width 2.975in;original-height
2.207in;cropleft "0";croptop "1";cropright "1";cropbottom "0";tempfilename
'SUS73V5H.wmf';tempfile-properties "XPR";}}For the long wavelength (red) the
central maximum is larger than the screen. The short wavelength spot will be
wholly on the screen as shown. The geometric spot is what we would see if
the light traveled straight through the opening. Notice that the short
wavelength spot is closer to the size of the geometric spot. In the limit
that 
\begin{equation*}
\lambda \ll a
\end{equation*}%
or for circular openings 
\begin{equation*}
\lambda \ll D
\end{equation*}%
then 
\begin{equation*}
\theta \approx \frac{\lambda }{a}\approx 0
\end{equation*}%
or%
\begin{equation*}
\theta \approx \frac{\lambda }{D}\approx 0
\end{equation*}%
and the spot size would be very nearly equal to the geometric spot size.
This is the limit we will called the \emph{ray approximation}.

For most of mankind's time on the Earth, it was very hard to build holes
that were comparable to the size of a wavelength of visible light (around $%
500\unit{nm}$). So it is no wonder that the waviness of light was missed for
so many years.

This ray limit is very useful for apertures the size of windows and
windshields. It was fine for such big openings to use this small $\lambda ,$
large aperture approximation. For these type apertures we just don't notice
the wave nature of light at all. For professional camera lenses, it starts
to matter, and for tiny cell phone cameras we need to worry.

\section{The Ray Approximation in Geometric Optics}

In the last section we said that when the geometric spot size was about the
same size as the spot due to diffraction, we could ignore diffraction and
use the simpler ray model. This is usually true in our personal experiences.
But this may not be true in experiments or devices we design. We should see
where the crossover point is.

Intuitively, if the aperture and the spot are the same size, that ought to
be some sort of critical point. That is when the aperture size is equal to
the spot size. We found that for a circular aperture the spot width is 
\begin{equation*}
w\approx 2.44\frac{\lambda }{D_{aperture}}L
\end{equation*}%
We want the case where $w=D_{aperture}$

\begin{equation*}
D_{aperture}=2.44\frac{\lambda }{D_{aperture}}L
\end{equation*}%
This gives 
\begin{equation*}
D_{aperture}=\sqrt{2.44\lambda L}
\end{equation*}

Of course this is for round apertures, but for square apertures we know we
remove the $2.44.$ This gives about a millimeter for visible wavelengths.%
\begin{eqnarray*}
D_{aperture} &=&\sqrt{2.44\left( 500\unit{nm}\right) \left( 1\unit{m}\right) 
} \\
&=&1.\,\allowbreak 104\,5\times 10^{-3}\allowbreak \unit{m}
\end{eqnarray*}

For apertures much larger than a millimeter, we expect interference effects
due to diffraction through the aperture to be much harder to see. We expect
diffraction effects to be easy to see if the aperture is smaller than a
millimeter. But what about when the aperture is about a millimeter in size?
That is a subject for junior level Optics class (PH375), and so we will
avoid this case in this class. But this is not too restrictive. Most good
optical systems have apertures larger than $1\unit{mm}.$ Cell phone cameras
may be an exception (but they play tricks to beat this by having multiple
apertures). Even our eyes have an aperture that varies from about $2\unit{mm}
$ to about $7\unit{mm},$ so most common experiences in visible wavelengths
will work fine with what we have learned. Note that for microwave or radio
wave systems this may really not be true!

How about the other extreme? Suppose $\lambda \gg D.$ This is really beyond
our class (requires partial differential equations), but in the extreme
case, we can use reason to find out what happens. If the opening is much
smaller than the wavelength, then the wave does not see the opening, and no
wave is produced on the other side. This is the case of a microwave oven
door. \FRAME{dtbpFU}{2.3843in}{1.8209in}{0pt}{\Qcb{{\protect\small Microwave
oven door https://commons.wikimedia.org/wiki/File:MW\_door.jpg}}}{}{Figure}{%
\special{language "Scientific Word";type "GRAPHIC";maintain-aspect-ratio
TRUE;display "USEDEF";valid_file "T";width 2.3843in;height 1.8209in;depth
0pt;original-width 3.2353in;original-height 2.4639in;cropleft "0";croptop
"1";cropright "1";cropbottom "0";tempfilename
'TBCHP401.wmf';tempfile-properties "XPR";}}If the wavelength is much larger
than the spacing of the little dots or lines that span the door, then the
waves will not leave the interior of the microwave oven. Of course as the
wavelength becomes closer to $D$ this is less true, and this case is more
challenging to calculate, and we will save it for a 300 level course.

To summarize%
\begin{equation*}
\begin{tabular}{ll}
$\lambda \ll D$ & Wave nature of light is not visible \\ 
$\lambda \approx D$ & Wave nature of light is apparent \\ 
$\lambda \gg D$ & Little to no penetration of aperture by the wave%
\end{tabular}%
\end{equation*}

\section{The ray model and phase}

There is a further complication that helps to explain why the wave nature of
light was not immediately apparent. Let's consider a light source.\FRAME{dhF%
}{0.4532in}{0.7827in}{0pt}{}{}{Figure}{\special{language "Scientific
Word";type "GRAPHIC";maintain-aspect-ratio TRUE;display "USEDEF";valid_file
"T";width 0.4532in;height 0.7827in;depth 0pt;original-width
1.3214in;original-height 2.3039in;cropleft "0";croptop "1";cropright
"1";cropbottom "0";tempfilename 'SUS73V5I.wmf';tempfile-properties "XPR";}}%
For an old fashioned light bulb, the filament is larger than about a
millimeter, so we should expect that diffraction should be hard to see. But
the filament is made of hot metal. The atoms of the hot metal emit light
because of the extra energy they have. The method of producing this light is
that the atom's excited electrons are in upper shells because of the extra
thermal energy provided by the electricity flowing through the filament. But
the electrons eventually fall to their proper shell, and in doing so they
give off the extra energy as light. It is not too hard to believe that this
process of exciting electrons and having them fall back down is a random
process. Each electron that moves starts a wave. The atoms have different
positions, so there will be a path difference $\Delta r$ between each atom's
waves. There will also be a time difference $\Delta t$ between when the
waves start. We can model this with a $\Delta \phi _{o}.$

It is also true that not all of the electrons fall from the same shell. This
gives us different frequencies, so we expect beating between different waves
from different atoms. It is also true that we have millions of atoms, so we
have millions of waves.

Let's look at just two of these waves

\begin{equation*}
\begin{tabular}{l}
$\lambda =2$ \\ 
$k=\frac{2\pi }{\lambda }$ \\ 
$\omega =1$ \\ 
$\phi _{o}=\frac{\pi }{6}$ \\ 
$t=0$ \\ 
$E_{o}=1\frac{\unit{N}}{\unit{C}}$%
\end{tabular}%
\end{equation*}

\begin{equation*}
E_{1}=E_{o}\sin \left( kx-\omega t\right)
\end{equation*}

\FRAME{dtbpFX}{2.156in}{0.6875in}{0pt}{}{}{Plot}{\special{language
"Scientific Word";type "MAPLEPLOT";width 2.156in;height 0.6875in;depth
0pt;display "USEDEF";plot_snapshots TRUE;mustRecompute FALSE;lastEngine
"MuPAD";xmin "0";xmax "5.001000";xviewmin "0";xviewmax "5.001000";yviewmin
"-2";yviewmax "2";viewset"XY";rangeset"X";plottype 4;labeloverrides
2;y-label "E";axesFont "Times New
Roman,12,0000000000,useDefault,normal";numpoints 100;plotstyle
"patch";axesstyle "normal";axestips FALSE;xis \TEXUX{x};var1name
\TEXUX{$x$};function \TEXUX{$\allowbreak \sin \pi x$};linecolor
"blue";linestyle 1;pointstyle "point";linethickness 3;lineAttributes
"Solid";var1range "0,5.001000";num-x-gridlines 100;curveColor
"[flat::RGB:0x000000ff]";curveStyle "Line";VCamFile
'TJXLF62F.xvz';valid_file "T";tempfilename
'SW60QA05.wmf';tempfile-properties "XPR";}}

\begin{equation*}
E_{2}=E_{o}\sin \left( kx-\omega t+\phi _{o}\right)
\end{equation*}%
\FRAME{dtbpFX}{2.156in}{0.6875in}{0pt}{}{}{Plot}{\special{language
"Scientific Word";type "MAPLEPLOT";width 2.156in;height 0.6875in;depth
0pt;display "USEDEF";plot_snapshots TRUE;mustRecompute FALSE;lastEngine
"MuPAD";xmin "0";xmax "5.001000";xviewmin "0";xviewmax "5.001000";yviewmin
"-2";yviewmax "2";viewset"XY";rangeset"X";plottype 4;labeloverrides
2;y-label "E";axesFont "Times New
Roman,12,0000000000,useDefault,normal";numpoints 100;plotstyle
"patch";axesstyle "normal";axestips FALSE;xis \TEXUX{x};var1name
\TEXUX{$x$};function \TEXUX{$\sin \left( \frac{1}{6}\pi +\pi x\right)
$};linecolor "maroon";linestyle 1;pointstyle "point";linethickness
3;lineAttributes "Solid";var1range "0,5.001000";num-x-gridlines
100;curveColor "[flat::RGB:0x00800000]";curveStyle "Line";VCamFile
'TJXLF62G.xvz';valid_file "T";tempfilename
'SW60QH06.wmf';tempfile-properties "XPR";}}then%
\begin{equation*}
E_{r}=E_{o}\sin \left( kx-\omega t\right) +E_{o}\sin \left( kx-\omega t+\phi
_{o}\right)
\end{equation*}

We found for these two waves%
\begin{equation*}
E_{r}=2E_{o}\cos \left( \frac{\phi _{o}}{2}\right) \sin \left( kx-\omega t+%
\frac{\phi _{o}}{2}\right)
\end{equation*}%
\FRAME{dtbpFX}{2.354in}{0.5846in}{0pt}{}{}{Plot}{\special{language
"Scientific Word";type "MAPLEPLOT";width 2.354in;height 0.5846in;depth
0pt;display "USEDEF";plot_snapshots TRUE;mustRecompute FALSE;lastEngine
"MuPAD";xmin "0";xmax "5.001000";xviewmin "0";xviewmax "5.001000";yviewmin
"-2";yviewmax "2";viewset"XY";rangeset"X";plottype 4;labeloverrides
2;y-label "E";axesFont "Times New
Roman,12,0000000000,useDefault,normal";numpoints 100;plotstyle
"patch";axesstyle "normal";axestips FALSE;xis \TEXUX{x};var1name
\TEXUX{$x$};function \TEXUX{$\sin \left( \frac{1}{6}\pi +\pi x\right) +\sin
\pi x$};linecolor "green";linestyle 1;pointstyle "point";linethickness
3;lineAttributes "Solid";var1range "0,5.001000";num-x-gridlines
100;curveColor "[flat::RGB:0x00008000]";curveStyle "Line";VCamFile
'TJXLF62H.xvz';valid_file "T";tempfilename
'SW60OR01.wmf';tempfile-properties "XPR";}}

But suppose we complicate the situation by sending lots of waves at random
times, each with different amplitudes and wavelengths, down the rope. If we
look at a single point for a specific time, we might be experiencing
interference, but it would be hard to tell. Lets try this mathematically. I
will combine many waves with random phases, some coming from the right and
some coming from the left.

\begin{eqnarray*}
E_{1} &=&E_{o}\sin \left( 5x-\omega t+\frac{\pi }{4}\right) +0.5E_{o}\sin
\left( 0.2x-\omega t-\frac{\pi }{6}\right)  \\
&&+3.6E_{o}\sin \left( .4x-\omega t+\frac{\pi }{10}\right) +4E_{o}\sin
\left( 20x-\omega t-\frac{\pi }{7}\right)  \\
&&+.2E_{o}\sin \left( 15x-\omega t+1\right) +0.7E_{o}\sin \left( .7x-\omega
t-.25\right) 
\end{eqnarray*}%
\FRAME{dtbpFX}{2.354in}{1.1939in}{0pt}{}{}{Plot}{\special{language
"Scientific Word";type "MAPLEPLOT";width 2.354in;height 1.1939in;depth
0pt;display "USEDEF";plot_snapshots TRUE;mustRecompute FALSE;lastEngine
"MuPAD";xmin "0";xmax "5";xviewmin "0";xviewmax "5";yviewmin "-20";yviewmax
"20";viewset"XY";rangeset"X";plottype 4;labeloverrides 2;y-label
"E";axesFont "Times New Roman,12,0000000000,useDefault,normal";numpoints
100;plotstyle "patch";axesstyle "normal";axestips FALSE;xis
\TEXUX{x};var1name \TEXUX{$x$};function \TEXUX{$0.2\sin \left( 15x+1\right)
+0.7\sin \left( 0.7x-0.25\right) +\allowbreak 0.5\sin \left(
0.2x-\frac{1}{6}\pi \right) +3.\,\allowbreak 6\sin \left( \frac{1}{10}\pi
+0.4x\right) +\allowbreak \sin \left( \frac{1}{4}\pi +5x\right) +4\sin
\left( 20x-\frac{1}{7}\pi \right) $};linecolor "blue";linestyle 1;pointstyle
"point";linethickness 1;lineAttributes "Solid";var1range
"0,5";num-x-gridlines 100;curveColor "[flat::RGB:0x000000ff]";curveStyle
"Line";VCamFile 'TJXLFS2J.xvz';valid_file "T";tempfilename
'TJXLFH05.wmf';tempfile-properties "XPR";}}And now let's make another wave
that is a combination of several waves with random phases%
\begin{eqnarray*}
E_{2} &=&E_{o}\sin \left( 0.2x+\omega t+\pi \right) +2E_{o}\sin \left(
5x+\omega t+\frac{\pi }{6}\right)  \\
&&+6E_{o}\sin \left( 0.4x+\omega t+\frac{\pi }{3.5}\right) +0.4E_{o}\sin
\left( 20x+\omega t-0\right)  \\
&&+E_{o}\sin \left( 15x+\omega t+1\right) +0.7E_{o}\sin \left( .7x+\omega
t-4\right) 
\end{eqnarray*}%
\FRAME{dtbpFX}{2.2716in}{1.0342in}{0pt}{}{}{Plot}{\special{language
"Scientific Word";type "MAPLEPLOT";width 2.2716in;height 1.0342in;depth
0pt;display "USEDEF";plot_snapshots TRUE;mustRecompute FALSE;lastEngine
"MuPAD";xmin "0";xmax "5.001000";xviewmin "0";xviewmax "5.001000";yviewmin
"-20";yviewmax "20";viewset"XY";rangeset"X";plottype 4;labeloverrides
2;y-label "E";axesFont "Times New
Roman,12,0000000000,useDefault,normal";numpoints 100;plotstyle
"patch";axesstyle "normal";axestips FALSE;xis \TEXUX{x};var1name
\TEXUX{$x$};function \TEXUX{$0.7\sin \left( 0.7x-4\right) +6\sin \left(
0.285\,71\pi +0.4x\right) +\allowbreak \sin \left( 15x+1\right) -\sin
0.2x+2\sin \left( \frac{1}{6}\pi +5x\right) +0.4\allowbreak \sin
20x$};linecolor "red";linestyle 1;pointstyle "point";linethickness
2;lineAttributes "Solid";var1range "0,5.001000";num-x-gridlines
100;curveColor "[flat::RGB:0x00ff0000]";curveStyle "Line";VCamFile
'TJXLC31W.xvz';valid_file "T";tempfilename
'TJXLBW01.wmf';tempfile-properties "XPR";}}

Then $E_{1}+E_{2}$ looks like\FRAME{dtbpFX}{2.3336in}{1.1185in}{0pt}{}{}{Plot%
}{\special{language "Scientific Word";type "MAPLEPLOT";width 2.3336in;height
1.1185in;depth 0pt;display "USEDEF";plot_snapshots TRUE;mustRecompute
FALSE;lastEngine "MuPAD";xmin "0";xmax "5.001000";xviewmin "0";xviewmax
"5.001000";yviewmin "-20";yviewmax "20";viewset"XY";rangeset"X";plottype
4;labeloverrides 2;y-label "E";axesFont "Times New
Roman,12,0000000000,useDefault,normal";numpoints 100;plotstyle
"patch";axesstyle "normal";axestips FALSE;xis \TEXUX{x};var1name
\TEXUX{$x$};function \TEXUX{$\allowbreak 0.7\sin \left( 0.7x-4\right) +6\sin
\left( 0.285\,71\pi +0.4x\right) +\allowbreak 1.\,\allowbreak 2\sin \left(
15x+1\right) +0.7\sin \left( 0.7x-0.25\right) +\allowbreak 0.5\sin \left(
0.2x-\frac{1}{6}\pi \right) +3.\,\allowbreak 6\sin \left( \frac{1}{10}\pi
+0.4x\right) -\allowbreak \sin 0.2x+\sin \left( \frac{1}{4}\pi +5x\right)
+2\sin \left( \frac{1}{6}\pi +5x\right) +4\sin \left( 20x-\frac{1}{7}\pi
\right) +\allowbreak 0.4\sin 20x$};linecolor "cyan";linestyle 1;pointstyle
"point";linethickness 1;lineAttributes "Solid";var1range
"0,5.001000";num-x-gridlines 100;curveColor
"[flat::RGB:0x000080c0]";curveStyle "Line";VCamFile
'TJXLCA1X.xvz';valid_file "T";tempfilename
'TJXLCA02.wmf';tempfile-properties "XPR";}}

We get constructive and destructive interference in a changing pattern. In
this example, you could think about the superposition of $E_{1}$ and $E_{2}$
and predict the outcome, but if there were millions of waves, each with it's
own wavelength, phase, and amplitude, the situation would be hopeless. Note
that the fluctuations in these waves are much more frequent than our
original waves. With all the added waves, we get a rapid change in
amplitude. And because the frequencies are not exactly the same the
interference is constantly changing because of beating.

Now if these waves are light waves, our eyes and most detectors are not able
to react fast enough to detect the rapid fluctuations. So if there is
constructive or destructive interference that might be simple enough to
distinguish, we will miss it due to our detection systems' integration
times. To describe this rapidly fluctuating interference pattern that we
can't track with our detectors, we just say that light bulbs emit \emph{%
incoherent light}. The ray approximation assumes incoherent light.

But then light bulbs and hot ovens and most things must emit incoherent
light. Does any thing emit coherent light? Sure, today the easiest source of
coherent light is a laser. That is why professors use lasers in the class
demonstrations. Really though, even a laser is not perfectly coherent. One
property of the laser is that it produces light with a long \emph{coherence
length}, or it produces light that can be treated under most circumstances
as being monochromatic and having a single phase across the wave for much of
the beam length. Radar and microwave transmitters emit coherent light (but
at frequencies we can't see) and so do radio stations.

In the past, one could carefully create a monochromatic beam with filters.
Then split the beam into two beams and remix the two beams. This would
generate two mostly coherent sources if the distances traveled were not too
large. This is what Young did. Now days we just use lasers.

\section{Coherency}

Light that easily shows the wave nature of light is said to be \emph{coherent%
}. To be coherent a light beam must satisfy two conditions:

\begin{enumerate}
\item A given part of the wave must maintain a constant phase with respect
to the rest of the wave.

\item The wave must be monochromatic
\end{enumerate}

These are very hard criteria to achieve. Most light, like that from our
light bulb, is not coherent. So most of the time the destructive and
constructive interference in the light we see is just not apparent. But the
wave nature of light it still affects resolution of cameras and telescopes
and even our eyes. And when we have coherent light, interference is apparent
and even useful as we have seen.

You may be left feeling that we have really just gotten started, and you
would be right. If you are a physics major (or curious and have elective
credit) you will study waves more in PH 279 (Modern Physics) and PH295
(Mathematical and Computational Physics) and you may study more optics in
PH375. All of these are fantastic experiences. Even if you don't take a
further class in these topics, you can study them on your own. There are
many wonderful books on Optics.

Our goal this semester was to study the motion of many things and wave
motion was a partial fulfilment of that goal. But not all motions of many
objects are as uniform as wave motion. What if, say, the air molecules in
our room instead of moving in simple harmonic motion like in a wave, had
instead random velocities? This is the type of motion that we hinted about
when we talked about thermal energy in Principles of Physics I (PH121) and
just now when we discussed hot atoms that emit incoherent light. We will
take on this kind of motion next in this course.

\part{Thermodynamics}

\chapter{24 Fluids and Pressure 1.14.1, 1.14.2}

We have spent about 2/3 of a semester studying waves. Waves are very regular
motion of more than one object. But surely we can have less regular motion
of more than one object!

To study such non-regular motion we are changing topics radically. PH121
taught us how individual objects move. We will use some of what we learned
to address the question of how many objects move. To get started, let's
review some basic properties of matter. Matter is made of many atoms. So
matter is an example of many objects that might just move.

%TCIMACRO{\TeXButton{\chaptermark{Fluids}}{\chaptermark{Fluids}}}%
%BeginExpansion
\chaptermark{Fluids}%
%EndExpansion

%TCIMACRO{%
%\TeXButton{Fundamental Concepts}{\vspace{0.10in}\hspace{-0.37in}{\Large Fundamental Concepts\vspace{0.2in}}}}%
%BeginExpansion
\vspace{0.10in}\hspace{-0.37in}{\Large Fundamental Concepts\vspace{0.2in}}%
%EndExpansion

\begin{itemize}
\item Compressibility of fluids

\item Density of fluids

\item Pressure is a force spread over an area

\item Pressure increases with depth in a fluid

\item Barometers

\item Manometers

\item Buoyancy
\end{itemize}

\section{Fluids}

You are probably aware that there are four general states of matter

\begin{equation*}
\begin{tabular}{|l|}
\hline
solid \\ \hline
liquid \\ \hline
gas \\ \hline
plasma \\ \hline
\end{tabular}%
\end{equation*}

Believe it or not, plasma\footnote{%
This is not the kind of plasma that you donate!} is the most common, because
stars are made of plasma. Planets are sometimes made of solids, liquids, or
gases, but the great glowing stars are plasma. (I am ignoring the mysterious
form of \textquotedblleft dark matter\textquotedblright\ because so far we
don't know what (or if) it is). Plasma is a heated gas that is ionized
(lacking an electron or two or more). We will mostly ignore this state,
because unless you are dealing with neon signs, fluorescent lights, or the
like, on Earth you don't encounter plasmas in every day experience. Let's
look at the other more common states of matter.

\smallskip

\textbf{Solids}

We can view solids as having a set of forces that keep the molecules in
place much as though they were attached using springs. Solids can have
definite organization. If so, they are called crystals.\footnote{%
If you are in Rexburg reading this, you should observe the crystals around
the Romney building if you have not already.} \FRAME{dhF}{2.2926in}{1.6328in%
}{0pt}{}{}{Figure}{\special{language "Scientific Word";type
"GRAPHIC";maintain-aspect-ratio TRUE;display "USEDEF";valid_file "T";width
2.2926in;height 1.6328in;depth 0pt;original-width 2.2528in;original-height
1.5965in;cropleft "0";croptop "1";cropright "1";cropbottom "0";tempfilename
'SUS73V8J.wmf';tempfile-properties "XPR";}}If the solid lacks definite order
in its organization, it is called amorphous.\FRAME{dhF}{1.3258in}{1.5428in}{%
0pt}{}{}{Figure}{\special{language "Scientific Word";type
"GRAPHIC";maintain-aspect-ratio TRUE;display "USEDEF";valid_file "T";width
1.3258in;height 1.5428in;depth 0pt;original-width 2.0721in;original-height
2.4146in;cropleft "0";croptop "1";cropright "1";cropbottom "0";tempfilename
'SUS73V8K.wmf';tempfile-properties "XPR";}}\bigskip

\textbf{Liquids}

The molecules in a liquid are less tightly bound than those in a solid. That
is why they can flow. In the next figure, the atoms are bound by one
spring-like force. But the atoms are not tied together in a tight set of
bonds like a solid. \FRAME{dhF}{2.9352in}{2.2329in}{0pt}{}{}{Figure}{\special%
{language "Scientific Word";type "GRAPHIC";maintain-aspect-ratio
TRUE;display "USEDEF";valid_file "T";width 2.9352in;height 2.2329in;depth
0pt;original-width 2.8911in;original-height 2.1923in;cropleft "0";croptop
"1";cropright "1";cropbottom "0";tempfilename
'SUS73V8L.wmf';tempfile-properties "XPR";}}\bigskip

\textbf{Gasses}

The molecules in a gas are not bound to each other -- not at all.

\FRAME{dhF}{2.3756in}{2.2883in}{0pt}{}{}{Figure}{\special{language
"Scientific Word";type "GRAPHIC";maintain-aspect-ratio TRUE;display
"USEDEF";valid_file "T";width 2.3756in;height 2.2883in;depth
0pt;original-width 2.3359in;original-height 2.2485in;cropleft "0";croptop
"1";cropright "1";cropbottom "0";tempfilename
'SUS73V8M.wmf';tempfile-properties "XPR";}}These last two states of matter
can both be called fluids. We have an intuitive feel for what a fluid is.
But let's make a more formal definition.

\subsection{What is a fluid?}

A \emph{Fluid} is a collection of molecules that are randomly arranged and
are at most held together by weak cohesive forces and by forces exerted by
the walls of a container.

Thus, liquids and gasses are definitely fluids. Solids are generally
considered not fluids. But how about Jell-O%
%TCIMACRO{\TeXButton{TM}{\textsuperscript{\texttrademark}}}%
%BeginExpansion
\textsuperscript{\texttrademark}%
%EndExpansion
? or a combination of corn starch and water (sometimes called
\textquotedblleft ooblick\textquotedblright )? These are non-Newtonian
fluids. That is, things that are sort of solid and sort of not. But we will
stick with things that are definitely fluids (at least at first), and
generally fluids with negligible friction. This is a little like in PH121
when we studied frictionless surfaces. How many surfaces are truly
frictionless? Very few! you might guess that there are few fluids that have
no friction, and you would be right. But just like with PH121, the
assumption of frictionless fluids makes the math easier, and that is good
when we are starting a new topic.

\section{Pressure}

%TCIMACRO{%
%\TeXButton{Vollyball Demo}{\marginpar {
%\hspace{-0.5in}
%\begin{minipage}[t]{1in}
%\small{Vollyball Demo}
%\end{minipage}
%}}}%
%BeginExpansion
\marginpar {
\hspace{-0.5in}
\begin{minipage}[t]{1in}
\small{Vollyball Demo}
\end{minipage}
}%
%EndExpansion
We have already seen pressure in this course. But it's been a while, so
let's review. Consider a situation where we ask six of our class members to
come up and press on the ball from all directions. Suppose further that we
ask each person to exert a force on the ball. And suppose we ask the person
to use the area of their hand to exert the force. The motion of the ball,
and even it's shape would depend on both the force (magnitude and direction)
and the area involved in each push.

Noticing that each person exerts a force but that the force is not acting on
one point. It is spread out over an area, we recognize that each person is
exerting a pressure on the ball\ 
\begin{equation}
P\equiv \frac{F}{A}  \label{PressureDefinition}
\end{equation}%
\mathstrut

We dealt with pressure earlier in our course, but let's review. Pressure is
a force spread over an area. In this case it is the normal force of the hand
on the ball spread over the area of the hand.

Now consider a ball sitting in a room surrounded by air. The air is a fluid,
so it's molecules are quite free to move around. Because there is some
thermal energy in the room\footnote{%
Even in Rexburg.} the molecules will have some kinetic energy. So the air
molecules will hit the ball. This will cause a force on the surface of the
ball. And that force will be spread over the entire surface area of the
ball. This is a force spread over an area. This is a pressure. We call this 
\emph{air pressure}. This air force due to the colliding molecules is like
having the hands pushing on the pushing on the ball.

This force due to individual molecules is small and only lasts during the
collision. But in the room we have many molecules, and many molecules impact
the ball. The molecules also impact the walls of the room. Suppose that
every time a molecule bounces back from one wall it ends up headed back to
the opposite wall bounces back again toward the first wall. If the molecules
keep coming, there will be a force on the wall quite a bit of the time. At
least, on average there is a force, anyway. This is the force that causes
air pressure. The molecules impact the walls, and the ball, and us, and
everything in the room all over the surface area of each object. The result
is air pressure on every object in the room.

Likewise, the water pressure in a swimming pool is caused by moving water
molecules. You should convince yourself that the reason the water stays in
the pool is partly because the air molecules bounce against the water
surface exerting a pressure on the water!

You might notice that there are more molecules bunched together in the water
than in the air. We can describe how bunched together the molecules are by
giving the mass of the molecules for a unit of volume%
\begin{equation}
\frac{m}{%
%TCIMACRO{\TeXButton{V}{\ooalign{\hfil$V$\hfil\cr\kern0.1em--\hfil\cr}}}%
%BeginExpansion
\ooalign{\hfil$V$\hfil\cr\kern0.1em--\hfil\cr}%
%EndExpansion
}=\rho  \label{Density}
\end{equation}%
We call this measure of bunched-up-ness the \emph{density} of the material
and we give it (unfortunately) the symbol $\rho $ (the Greek letter
\textquotedblleft rho\textquotedblright\ pronounced like \textquotedblleft
row\textquotedblright\ as in \textquotedblleft row a boat\textquotedblright
). Sometimes it is useful to compare the density of a substance (like air)
to the density of water%
\begin{equation*}
\frac{\rho _{\text{substance}}}{\rho _{water}}=\text{specific gravity}
\end{equation*}%
This comparison is (unfortunately) called the \emph{specific gravity}. But
we need to be careful. We will find that density can depend on temperature.
So let's define specific gravity as the comparison of the substance's
density to the density of water at $4\unit{%
%TCIMACRO{\U{2103}}%
%BeginExpansion
{}^{\circ}{\rm C}%
%EndExpansion
}$ at $1\unit{atm}$ of pressure. It turns out at this temperature and
pressure 
\begin{equation*}
\rho _{water}=1000\frac{\unit{kg}}{\unit{m}^{3}}
\end{equation*}%
which makes it convenient for comparison.

\subsection{Working with the definition of pressure}

We can work with equation \ref{PressureDefinition} to define the force due
to pressure%
\begin{equation*}
F\equiv PA
\end{equation*}%
and can define a force due to the pressure at a small element of area $%
\Delta A,$ let's call this force on a small part of the area $\Delta F$ 
\begin{equation*}
\Delta F\equiv P\Delta A
\end{equation*}%
If we make the area very small we can write this as 
\begin{equation*}
dF=PdA
\end{equation*}%
Where there is a differential, expect that some time in the future we will
integrate!

But before we go on, let's see what the units of pressure would be. We have
a force divided by an area.

\begin{equation*}
1\allowbreak \frac{\unit{N}}{\unit{m}^{2}}\allowbreak =1\unit{Pa}
\end{equation*}%
The symbol is the $\unit{Pa},$ and the unit is called the \emph{pascal}. It
is named after a famous scientist.

\subsection{Pressure Example:}

Let's do a pressure problem together,

Problem statement:

A $50.0\unit{kg}$ person balances on one heel of a pair of high heeled
shoes. If the heel cross section is circular and has a radius of $0.5000%
\unit{cm},$ what pressure does she exert on the floor?

Drawing

\FRAME{dhF}{1.1044in}{0.4877in}{0pt}{}{}{Figure}{\special{language
"Scientific Word";type "GRAPHIC";maintain-aspect-ratio TRUE;display
"USEDEF";valid_file "T";width 1.1044in;height 0.4877in;depth
0pt;original-width 1.0706in;original-height 0.4575in;cropleft "0";croptop
"1";cropright "1";cropbottom "0";tempfilename
'SUS73V8N.wmf';tempfile-properties "XPR";}}

Variables

\begin{equation*}
\begin{tabular}{lll}
& Known &  \\ 
$M$ & Mass of person & $M=50\unit{kg}$ \\ 
$r$ & Radius of person's heal & $r=0.500\unit{cm}$ \\ 
$g$ & Acceleration do to gravity & $g=9.8\frac{\unit{m}}{\unit{s}^{2}}$ \\ 
& Unknown &  \\ 
$F$ & Force &  \\ 
$x$ & Coordinate Axis &  \\ 
$z$ & Coordinate Axis &  \\ 
$A$ & Area of the heal & 
\end{tabular}%
\end{equation*}

Basic Equations

\begin{eqnarray*}
F &=&ma \\
A &=&\pi r^{2} \\
P &=&\frac{F}{A}
\end{eqnarray*}

Symbolic Solution

\begin{equation*}
A=\pi r^{2}
\end{equation*}

\begin{equation*}
F=ma=Mg
\end{equation*}%
\begin{equation*}
P=\frac{F}{A}=\frac{Mg}{\pi r^{2}}
\end{equation*}

Numerical Solution

\begin{eqnarray*}
P &=&\frac{Mg}{\pi r^{2}} \\
&=&\frac{50\unit{kg}\ast 9.8\frac{\unit{m}}{\unit{s}^{2}}}{\pi \left( 0.500%
\unit{cm}\right) ^{2}} \\
&=&\frac{490.0\frac{\unit{m}}{\unit{s}^{2}}\unit{kg}}{\allowbreak 0.785\,40%
\unit{cm}^{2}}\frac{\left( 100\unit{cm}\right) ^{2}}{\unit{m}^{2}} \\
&=&\frac{490.0\frac{\unit{m}}{\unit{s}^{2}}\unit{kg}}{\allowbreak 0.785\,40}%
\frac{100000}{\unit{m}^{2}} \\
&=&6.\,\allowbreak 238\,9\times 10^{6}\allowbreak \frac{\unit{kg}\unit{m}}{%
\unit{s}^{2}\unit{m}^{2}} \\
&=&6.\,\allowbreak 238\,9\times 10^{6}\unit{Pa}
\end{eqnarray*}%
\begin{equation*}
P=6.24\unit{MPa}
\end{equation*}

Units Check

\begin{equation*}
\frac{\unit{kg}\frac{\unit{m}}{\unit{s}^{2}}}{\unit{cm}^{2}}=\frac{\unit{kg}%
\frac{\unit{m}}{\unit{s}^{2}}}{\unit{cm}^{2}}\frac{\left( 100\unit{cm}%
\right) ^{2}}{\unit{m}^{2}}=10000\frac{\unit{kg}\frac{\unit{m}}{\unit{s}^{2}}%
}{\unit{m}^{2}}=10000\unit{Pa}
\end{equation*}

Units

Check

Reasonableness

This seems like a large number, but I have had a high heeled person step on
my toe, so I believe it! We can compare to average atmospheric pressure
which is about 
\begin{equation*}
P_{atmospheric}=1.13\times 10^{5}\unit{Pa}
\end{equation*}%
Compared to this the toe crunch seems pretty reasonable.

\section{Variation of Pressure with Depth}

If you have been in a swimming pool, you probably have noticed that the
pressure due to the water feels different at the surface than it does at the
bottom of the deep end. Intuitively, we can say that the pressure at the
bottom is larger than the pressure at the top. Let's see if we can show that
this is true. We will do this in a few different ways.

\subsection{Pressure variation in liquids}

Take a glass of water or some other liquid. \FRAME{dtbpF}{1.6968in}{2.002in}{%
0pt}{}{\Qlb{PressureVariationInitial}}{Figure}{\special{language "Scientific
Word";type "GRAPHIC";maintain-aspect-ratio TRUE;display "USEDEF";valid_file
"T";width 1.6968in;height 2.002in;depth 0pt;original-width
1.6604in;original-height 1.9631in;cropleft "0";croptop "1";cropright
"1";cropbottom "0";tempfilename 'SXPQX003.wmf';tempfile-properties "XPR";}}%
And let's look at just a part of the liquid (the section inside the
box-shaped section with the dotted line around it in figure \ref%
{PressureVariationInitial}). The box shaped part of the liquid looks like,
well, a box. And an old word for a box-like package (like those that come
from online shopping) is \textquotedblleft\ parcel.\textquotedblright\ Let's
treat this \textquotedblleft parcel of fluid\textquotedblright\ as thought
it were distinct object and look at the forces acting on it.

We need some way of labeling our forces. My way is kind of simple, but will
work for now.\FRAME{dtbpF}{2.405in}{0.8925in}{0pt}{}{}{Figure}{\special%
{language "Scientific Word";type "GRAPHIC";maintain-aspect-ratio
TRUE;display "USEDEF";valid_file "T";width 2.405in;height 0.8925in;depth
0pt;original-width 3.3702in;original-height 1.2324in;cropleft "0";croptop
"1";cropright "1";cropbottom "0";tempfilename
'SUS73V8O.wmf';tempfile-properties "XPR";}}

Consider Newton's second law. The sum of the forces must be equal to zero.
Why?

Think of the acceleration of the parcel of liquid. Unless the water is
swirling around it's not accelerating. Let's agree for this part of the
course to only study fluids that are in equilibrium so we don't have any
swirling. And then the parcel acceleration is zero which makes the net force
zero. \FRAME{dtbpF}{2.5547in}{1.7538in}{0pt}{}{}{Figure}{\special{language
"Scientific Word";type "GRAPHIC";maintain-aspect-ratio TRUE;display
"USEDEF";valid_file "T";width 2.5547in;height 1.7538in;depth
0pt;original-width 2.5131in;original-height 1.7158in;cropleft "0";croptop
"1";cropright "1";cropbottom "0";tempfilename
'SXT8F105.wmf';tempfile-properties "XPR";}}And what we have in our drawing
looks a lot like a free body diagram for our parcel of liquid. We could draw
the free body diagram like this \FRAME{dtbpF}{1.8766in}{1.2289in}{0pt}{}{}{%
Figure}{\special{language "Scientific Word";type
"GRAPHIC";maintain-aspect-ratio TRUE;display "USEDEF";valid_file "T";width
1.8766in;height 1.2289in;depth 0pt;original-width 1.8386in;original-height
1.1952in;cropleft "0";croptop "1";cropright "1";cropbottom "0";tempfilename
'SXT8FQ06.wmf';tempfile-properties "XPR";}}where you should notice that we
are drawing the forces with their points on the part of the parcel area
where the force acts. That is different than what we did in Principles of
Physics I. This is because we are going to turn these forces into pressures
using 
\begin{equation*}
P=\frac{F}{A}
\end{equation*}%
so we want to identify the area that the force acts on. Let's write out
Newton's second law for our parcel of liquid. 
\begin{equation*}
\overrightarrow{F}_{net}=m\overrightarrow{a}=0=\overrightarrow{F}_{pt}+%
\overrightarrow{F}_{pb}+\overrightarrow{F}_{pmr}+\overrightarrow{F}_{pml}+%
\overrightarrow{W}
\end{equation*}%
but just like in Principles of Physics I, we will break this equation into
components.%
\begin{eqnarray*}
F_{net_{x}} &=&0=0+0-F_{pmr}+F_{pml}-0 \\
F_{net_{y}} &=&0=-F_{pt}+F_{pb}-0+0-W
\end{eqnarray*}

The $x$-part tells us that 
\begin{equation*}
F_{pmr}=F_{pml}
\end{equation*}%
which is not too much of a surprise. It seems reasonable that the side
forces must be equal if the parcel doesn't accelerate. The $y$-part gives%
\begin{equation*}
F_{pb}-F_{pt}=W
\end{equation*}

Now let's use our pressure equation. The definition of pressure gives%
\begin{equation*}
F=PA
\end{equation*}%
\FRAME{dtbpF}{2.7432in}{1.8057in}{0in}{}{}{Figure}{\special{language
"Scientific Word";type "GRAPHIC";maintain-aspect-ratio TRUE;display
"USEDEF";valid_file "T";width 2.7432in;height 1.8057in;depth
0in;original-width 2.7008in;original-height 1.7685in;cropleft "0";croptop
"1";cropright "1";cropbottom "0";tempfilename
'SXT8H407.wmf';tempfile-properties "XPR";}}

So, the force on the top must be%
\begin{equation*}
F_{pt}=P_{t}A
\end{equation*}%
and for the bottom the force must be%
\begin{equation*}
F_{pb}=P_{b}A
\end{equation*}

Recalling that 
\begin{equation}
m=\rho 
%TCIMACRO{\TeXButton{V}{\ooalign{\hfil$V$\hfil\cr\kern0.1em--\hfil\cr}}}%
%BeginExpansion
\ooalign{\hfil$V$\hfil\cr\kern0.1em--\hfil\cr}%
%EndExpansion
\end{equation}%
where $\rho $ is the density and $%
%TCIMACRO{\TeXButton{V}{\ooalign{\hfil$V$\hfil\cr\kern0.1em--\hfil\cr}}}%
%BeginExpansion
\ooalign{\hfil$V$\hfil\cr\kern0.1em--\hfil\cr}%
%EndExpansion
$ is the volume, and substituting, we get%
\begin{eqnarray*}
F_{pb}-F_{pt} &=&W \\
P_{b}A-P_{t}A &=&-mg \\
\left( P_{b}-P_{t}\right) A &=&-\rho 
%TCIMACRO{\TeXButton{V}{\ooalign{\hfil$V$\hfil\cr\kern0.1em--\hfil\cr}}}%
%BeginExpansion
\ooalign{\hfil$V$\hfil\cr\kern0.1em--\hfil\cr}%
%EndExpansion
g
\end{eqnarray*}%
Now consider, the volume must be the area of the top or bottom, $A,$
multiplied by the length, $h,$ of our parcel.%
\begin{equation}
%TCIMACRO{\TeXButton{V}{\ooalign{\hfil$V$\hfil\cr\kern0.1em--\hfil\cr}}}%
%BeginExpansion
\ooalign{\hfil$V$\hfil\cr\kern0.1em--\hfil\cr}%
%EndExpansion
=Ah
\end{equation}%
Then we can write our new pressure equation as 
\begin{equation*}
\left( P_{b}-P_{t}\right) A=-\rho hAg
\end{equation*}%
and the areas cancel 
\begin{equation*}
\left( P_{b}-P_{t}\right) =-\rho hg
\end{equation*}%
\FRAME{dtbpF}{1.3136in}{1.5039in}{0pt}{}{}{Figure}{\special{language
"Scientific Word";type "GRAPHIC";maintain-aspect-ratio TRUE;display
"USEDEF";valid_file "T";width 1.3136in;height 1.5039in;depth
0pt;original-width 1.2791in;original-height 1.4684in;cropleft "0";croptop
"1";cropright "1";cropbottom "0";tempfilename
'SXT8HR08.wmf';tempfile-properties "XPR";}}

We can solve this for the pressure at the bottom 
\begin{equation*}
P_{b}=P_{t}+\rho gh
\end{equation*}

Units Check

\begin{equation*}
\unit{Pa}\Bumpeq \unit{Pa}+\frac{\unit{kg}}{\unit{m}^{3}}\frac{\unit{m}}{%
\unit{s}^{2}}\unit{m}
\end{equation*}%
the last term can be written as%
\begin{equation*}
\left( \frac{\unit{kg}\unit{m}}{\unit{s}^{2}}\right) \frac{1}{\unit{m}^{2}}%
\Rightarrow \frac{\unit{N}}{\unit{m}^{2}}\Rightarrow \unit{Pa}
\end{equation*}%
so the units check.

This is a profound statement! (and a new basic equation for us). The
pressure is larger at the bottom of the pool than the top. And it is larger
by the amount $\rho gh$ where $\rho $ is the liquid density, $g$ is the
acceleration due to gravity, and $h$ is how far down we go to get to the
bottom of our parcel.

You might have thought of an objection in how we came up with this equation
for the pressure at the bottom of a parcel of liquid. We said the sideways
forces on the parcel were equal. But now we know the pressure increases with
depth. Won't that affect the side forces so that the force on the bottom of
the side is larger than the force on the top of the side? Of course the
answer is yes. We could solve this problem by making the sides of the box
infinitesimally small. Let's call the parcel height $\Delta y$ \FRAME{dtbpF}{%
1.9666in}{2.322in}{0pt}{}{}{Figure}{\special{language "Scientific Word";type
"GRAPHIC";maintain-aspect-ratio TRUE;display "USEDEF";valid_file "T";width
1.9666in;height 2.322in;depth 0pt;original-width 1.9285in;original-height
2.2822in;cropleft "0";croptop "1";cropright "1";cropbottom "0";tempfilename
'SXR9HK04.wmf';tempfile-properties "XPR";}}With our parcel of liquid so very
thin now the pressure on the side can't have changed much from the top to
the bottom of the parcel. So now we know the side forces cancel. We could
say that the change in pressure from the top to the bottom 
\begin{eqnarray*}
P_{b} &=&P_{t}+\rho g\Delta y \\
\Delta P_{small} &=&\rho g\Delta y
\end{eqnarray*}%
And we could consider making a larger parcel to be made of many small
parcels of thickness $\Delta y.$ \FRAME{dtbpF}{1.9692in}{2.322in}{0pt}{}{}{%
Figure}{\special{language "Scientific Word";type
"GRAPHIC";maintain-aspect-ratio TRUE;display "USEDEF";valid_file "T";width
1.9692in;height 2.322in;depth 0pt;original-width 1.9303in;original-height
2.2822in;cropleft "0";croptop "1";cropright "1";cropbottom "0";tempfilename
'SXSXWL00.wmf';tempfile-properties "XPR";}}For each of the little parcels
the side forces cancel. We just have top and bottom forces. For the topmost
parcel we would have 
\begin{equation*}
P_{b1}=P_{t}+\rho g\Delta y_{1}
\end{equation*}%
And note, the bottom pressure from the topmost parcel is the top pressure
for the second parcel%
\begin{eqnarray*}
P_{b2} &=&P_{t2}+\rho g\Delta y_{2} \\
&=&P_{b1}+\rho g\Delta y_{2}
\end{eqnarray*}%
And the bottom pressure from the second parcel is the top pressure for the
third parcel%
\begin{eqnarray*}
P_{b3} &=&P_{t3}+\rho g\Delta y_{3} \\
&=&P_{b2}+\rho g\Delta y_{3}
\end{eqnarray*}

To get to the bottom of the big parcel we would have 
\begin{eqnarray*}
P_{b(total)} &=&P_{t}+\sum_{i=1}^{N}\rho g\Delta y_{i} \\
&=&P_{t}+\rho g\sum_{i=1}^{N}\Delta y_{i} \\
&=&P_{t}+\rho gh
\end{eqnarray*}%
just as we thought.

Of course we could take a limit to make $\Delta y\rightarrow dy$ and then we
would have 
\begin{eqnarray*}
P_{b} &=&P_{t}+\int_{d}^{d+h}\rho gdy \\
&=&P_{t}+\left. \rho gy\mathstrut \right\vert _{d}^{d+h} \\
&=&P_{t}+\rho g\left( \left( d+h\right) -d\right) \\
&=&P_{t}+\rho gh
\end{eqnarray*}%
just as we said before.

Note that we used a box shaped volume, but the result is general. To see
this, view any arbitrary volume as consisting of little boxes. As we make
the box size small, we approximate the actual volume better and better until
it can be exact. Since it works for each box that makes up the volume
individually, it will work for the whole volume (it takes a minute to think
about this).

\subsection{Pressure variation in gasses}

You probably noticed that in our last section we didn't let $\rho $ change.
Another way to say that the density doesn't change is to call our liquid 
\emph{incompressible}. Liquids are fluids, but they are difficult to
compress. So to a good approximation their density doesn't change. But
gasses are also fluids, and it is easy to compress a gas. We can't make the
approximation of incompressibility for gasses.

\FRAME{dtbpF}{2.4898in}{2.156in}{0in}{}{}{Figure}{\special{language
"Scientific Word";type "GRAPHIC";maintain-aspect-ratio TRUE;display
"USEDEF";valid_file "T";width 2.4898in;height 2.156in;depth
0in;original-width 2.4491in;original-height 2.1171in;cropleft "0";croptop
"1";cropright "1";cropbottom "0";tempfilename
'SXRH2H00.wmf';tempfile-properties "XPR";}}

We don't know enough yet to know how $\rho $ changes with $y.$ But if you
took high school physics or chemistry we can make an approximation using the
Ideal Gas Law. If you have never heard of the Ideal Gas Law, don't despair.
We will study it soon. But for now here is the equation 
\begin{equation*}
P%
%TCIMACRO{\TeXButton{V}{\ooalign{\hfil$V$\hfil\cr\kern0.1em--\hfil\cr}}}%
%BeginExpansion
\ooalign{\hfil$V$\hfil\cr\kern0.1em--\hfil\cr}%
%EndExpansion
=Nk_{B}T
\end{equation*}%
where $P$ is the pressure, $%
%TCIMACRO{\TeXButton{V}{\ooalign{\hfil$V$\hfil\cr\kern0.1em--\hfil\cr}}}%
%BeginExpansion
\ooalign{\hfil$V$\hfil\cr\kern0.1em--\hfil\cr}%
%EndExpansion
$ is the volume of the air, $N$ is a number of air molecules, $%
k_{B}=1.3806568\times 10^{-23}\frac{\unit{J}}{\unit{K}}$ is a constant, and $%
T$ is the temperature. We can solve this for pressure

\begin{equation*}
P=\frac{Nk_{B}T}{%
%TCIMACRO{\TeXButton{V}{\ooalign{\hfil$V$\hfil\cr\kern0.1em--\hfil\cr}}}%
%BeginExpansion
\ooalign{\hfil$V$\hfil\cr\kern0.1em--\hfil\cr}%
%EndExpansion
}
\end{equation*}%
And now we see that our pressure depends on the volume. This makes sense for
compressible fluids like gases.

We can change the form of this a bit by realizing that the number of
molecules is the total mass of our parcel of gas divided by the mass of one
gas molecule. 
\begin{equation*}
N=\frac{m_{total}}{m_{molecule}}
\end{equation*}%
so that 
\begin{equation*}
P=\frac{\frac{m_{total}}{m_{molecule}}k_{B}T}{%
%TCIMACRO{\TeXButton{V}{\ooalign{\hfil$V$\hfil\cr\kern0.1em--\hfil\cr}}}%
%BeginExpansion
\ooalign{\hfil$V$\hfil\cr\kern0.1em--\hfil\cr}%
%EndExpansion
}
\end{equation*}%
Then 
\begin{equation*}
P=\frac{\frac{m_{total}}{%
%TCIMACRO{\TeXButton{V}{\ooalign{\hfil$V$\hfil\cr\kern0.1em--\hfil\cr}}}%
%BeginExpansion
\ooalign{\hfil$V$\hfil\cr\kern0.1em--\hfil\cr}%
%EndExpansion
}k_{B}T}{m_{molecule}}=\rho \frac{k_{B}T}{m_{molecule}}
\end{equation*}%
and we can solve this for the density so we can see how density of the gas
changes with pressure.

\begin{equation*}
\frac{Pm_{molecule}}{k_{B}T}=\rho
\end{equation*}%
Now we can use our equation for a small parcel of fluid (gas this time)

\begin{equation*}
P_{b}=P_{t}+\rho g\Delta y
\end{equation*}%
but this time we want to know $P_{t}$ at the top so let's solve for that. 
\begin{equation*}
P_{t}=P_{b}-\rho g\Delta y
\end{equation*}%
If we bring $P_{b}$ over to the right we have and 
\begin{equation*}
P_{t}-P_{b}=-\rho g\Delta y
\end{equation*}%
which we could write it as $\Delta P.$ Dividing both sides by $\Delta y$
gives 
\begin{equation*}
\frac{\Delta P}{\Delta y}=-\rho g=-\frac{Pm_{molecule}g}{k_{B}T}
\end{equation*}

This seems pretty terrible, but we can take the limit as $\Delta
y\rightarrow dy$%
\begin{equation*}
\frac{dP}{dy}=-P\left( \frac{m_{molecule}g}{k_{B}T}\right)
\end{equation*}%
and we note that so long as the temperature doesn't change, all the things
in the parentheses are constant. So we could give this group of constants a
symbol, say, $\alpha $%
\begin{equation*}
\alpha =\frac{m_{molecule}g}{k_{B}T}
\end{equation*}%
Then our equation for how pressure varies with height is just 
\begin{equation*}
\frac{dP}{dy}=-\alpha P
\end{equation*}%
and this is a nice differential equation. But rather than use our experience
from our differential equations class (that we haven't taken yet) to solve
this let's rearrange it 
\begin{equation*}
\frac{dP}{P}=-\alpha dy
\end{equation*}%
and now let's integrate over the height of our parcel of air from $y=0$ to $%
y=h.$ On the pressure side we need to integrate from $P_{b}$ to $P_{t}$%
\begin{equation*}
\int_{P_{b}}^{P_{t}}\frac{1}{P}dP=-\int_{0}^{h}\alpha dy
\end{equation*}%
We get 
\begin{equation*}
\left. \ln P\mathstrut \right\vert _{P_{b}}^{P_{t}}=-\alpha \left.
y\mathstrut \right\vert _{o}^{h}
\end{equation*}%
\begin{equation*}
\ln \left( P_{t}\right) -\ln \left( P_{b}\right) =-\alpha h
\end{equation*}%
Using one of our favorite log identities $\ln a-\ln b=\ln \frac{a}{b}$ we
can rewrite this as 
\begin{equation*}
\ln \left( \frac{P_{t}}{P_{b}}\right) =-\alpha h
\end{equation*}%
and we can un-log this by exponentiating both sides 
\begin{equation*}
\frac{P_{t}}{P_{b}}=e^{-\alpha h}
\end{equation*}%
which gives finally 
\begin{equation*}
P_{t}=P_{b}e^{-ah}
\end{equation*}%
We could let our parcel of air be arbitrarily tall by letting the height be $%
h=y$ above the surface $(y=0)$ and making $P_{t}=P\left( y\right) $ and
letting the surface pressure $P_{b}=P_{surface}.$ Then 
\begin{equation*}
P\left( y\right) =P_{surface}e^{-ay}
\end{equation*}%
but this can only work for small distances above the Earth's surface because
in reality the temperature also changes with $y.$ But in our classrooms this
approximation is good enough (because the hvac system is keeping the whole
lab at the same temperature). Here is a graph of the pressure and density as
a function of height from the equations given in the 1976 US\ Standard
Atmosphere. \FRAME{dtbpFU}{2.5822in}{3.2755in}{0pt}{\Qcb{{\protect\small %
Total pressure and mass density as a function of geometric altitude. 1976 US
Standard Atmosphere}}}{}{Figure}{\special{language "Scientific Word";type
"GRAPHIC";maintain-aspect-ratio TRUE;display "USEDEF";valid_file "T";width
2.5822in;height 3.2755in;depth 0pt;original-width 2.563in;original-height
3.259in;cropleft "0";croptop "1";cropright "1";cropbottom "0";tempfilename
'TGH7Z500.wmf';tempfile-properties "XPR";}}You can see that it is more
complicated than just and exponential. If you need a more accurate estimate
for larger altitudes, the Standard Atmosphere is a good place to start.

\section{Pressure Measurements}

Now that we understand pressure and how it changes in fluids, we can study
several pressure measurement devices. But there are some quirks in pressure
measurement. Humans have measured pressure for a long time. But for most of
our existence what we cared about was how much larger a pressure was than
the surrounding air pressure. Think of the pressure inside a steam engine.
The steam pressure has to be larger than atmospheric pressure to make the
steam engine work. So humans invented pressure gages to put on things like
steam engines that tell us how much higher than atmospheric pressure our
particular pressure is. We call this a \emph{gauge pressure.} But now days
we understand that most of the universe is a vacuum. And that vacuum of
space has almost no pressure.\footnote{%
You might have heard an old saying which says that nature abhors a vacuum.
But since most of nature is free space, this old saying only applies to
things on or in planets or stars. The universe seems to love vacuum.} We
need to calculate the \emph{absolute}\textbf{\ }pressure, and if your gauge
measured gauge pressure you need to add atmospheric pressure to your reading.%
\begin{equation*}
P_{abs}=P_{gague}+P_{atmospheric}
\end{equation*}%
But with this understanding, let's look at pressure measuring devices.

\subsection{The Barometer}

The barometer is a common pressure measuring device. It consists of \FRAME{%
dtbpF}{1.2799in}{2.0578in}{0pt}{}{}{Figure}{\special{language "Scientific
Word";type "GRAPHIC";maintain-aspect-ratio TRUE;display "USEDEF";valid_file
"T";width 1.2799in;height 2.0578in;depth 0pt;original-width
1.2755in;original-height 2.0684in;cropleft "0";croptop "1";cropright
"1";cropbottom "0";tempfilename 'SUS73V8U.wmf';tempfile-properties "XPR";}}

\begin{enumerate}
\item a long tube closed at one end

\item a dish

\item Mercury or another fluid
\end{enumerate}

The dish and the tube are filled with the fluid. The tube is inverted in the
dish. The pressure at the top of the tube is essentially zero. This is
because the weight of the fluid pulls the mass of fluid downward. If the
fluid is mercury, it is massive enough to leave a vacuum behind.

The pressure at point $A$ and point $B$ must be the same (or fluid would
flow until $P_{A}=P_{B}$). This might be a little hard to see. But remember
the air has a pressure, and just because it is a gas it does not mean that
pressure is lower than the liquid pressure (think of the gas pressure in a
hurricane that blows over trees and buildings).

It would be useful to find out how high up the mercury (or other fluid) will
stay. And we have a way to do this! We know 
\begin{equation*}
P_{b}=P_{t}+\rho gh
\end{equation*}%
At the top we have zero pressure, so $P_{t}=0.$ Note, this is not
atmospheric pressure, \textbf{it is vacuum}. How did we get vacuum at the
top? We got it by filling the tube with mercury and then turning the tube
upside down. The mercury is heavy. So heavy that it falls down the tube a
bit, leaving nothing behind. So there is no air at the top of the tube, just
empty tube because the heavy mercury moved downward. But the mercury can
only fall down so far. The air pressure on the mercury in the dish pushes,
keeping the mercury from all falling out. Think of all the little air
molecules hitting the top of the mercury in the dish pushing it downward and
back into the tube. Then for this special case 
\begin{equation}
P_{b}=0+\rho gh
\end{equation}

And we can solve for the height of the fluid%
\begin{equation}
h=\frac{P_{b}}{\rho g}
\end{equation}%
Let's consider what $P_{b}$ must be. The air outside the barometer pushes
down on the liquid mercury. Remember, when we first construct a barometer,
we fill the tall tube with fluid and then turn it upside down in the dish.
The fluid will fall until the force due to air pressure matches the weight
of the column of fluid. Then the pressure at $B$ must be atmospheric
pressure. Since the pressure at point $A$ equals the pressure at point $B$
(think $P_{b}=P_{t}+\rho gh$), we can use this as a convenient $P_{b}.$ So 
\begin{equation*}
P_{b}=P_{atm}
\end{equation*}%
and 
\begin{equation}
h=\frac{P_{atm}}{\rho g}
\end{equation}%
Thus as the atmospheric air pressure changes, the height of the mercury
changes. We often hear of atmospheric pressure given in $\unit{mm}$ of $Hg$.
This is why.

\subsection{Manometer}

The manometer is often used in industrial settings to find the pressure in
pipes. \FRAME{dtbpF}{1.5629in}{1.962in}{0pt}{}{}{Figure}{\special{language
"Scientific Word";type "GRAPHIC";maintain-aspect-ratio TRUE;display
"USEDEF";valid_file "T";width 1.5629in;height 1.962in;depth
0pt;original-width 1.5646in;original-height 1.9709in;cropleft "0";croptop
"1";cropright "1";cropbottom "0";tempfilename
'SUS73V8V.wmf';tempfile-properties "XPR";}}The atmospheric pressure $P_{atm}$
is applied on one end of the small manometer tube. In our example it is on
the right side, where the tube is open at the top. The pressure to be tested
is applied to the other side, in our case it is on the left side. It is the
pressure of the fluid in the large pipe. The pressure at point $A$ must
equal the pressure at point $B$. Why? (think of their heights above the
u-shaped bend in the manometer pipe and $\rho gh$). The pressure at point $A$
is very close to the pressure to be measured from the big pipe. Again 
\begin{equation*}
P_{b}=P_{t}+\rho gh
\end{equation*}%
where we take $P_{A}$ to be the test pressure, and for this design the tube
above $B$ is open to the atmosphere so the top pressure will be atmospheric
pressure. 
\begin{equation*}
P_{A}=P_{B}=P_{atm}+\rho gh
\end{equation*}%
We can measure the height $h,$ and knowing $P_{atm}$ we can calculate $%
P_{A}. $

These are just two of many types of pressure measuring devices. In the next
lecture we will use what we know of pressure to design some useful devices
like hydraulic jacks and things that float (like boats).

\chapter{25 Pascal and Archimedes 1.14.3, 1.14.4}

So now we know how pressure changes with depth, and we know that liquids are
generally hard to compress, so we can approximate them as incompressible,
let's build some things with our new knowledge.

%TCIMACRO{%
%\TeXButton{\chaptermark{Pascal, Archimedes}}{\chaptermark{Pascal, Archimedes}}}%
%BeginExpansion
\chaptermark{Pascal, Archimedes}%
%EndExpansion

%TCIMACRO{%
%\TeXButton{Fundamental Concepts}{\vspace{0.10in}\hspace{-0.37in}{\Large Fundamental Concepts\vspace{0.2in}}}}%
%BeginExpansion
\vspace{0.10in}\hspace{-0.37in}{\Large Fundamental Concepts\vspace{0.2in}}%
%EndExpansion

\begin{itemize}
\item Pascal's Principle

\item Hydraulic Presses

\item Archimedes Principle

\item Buoyant Forces
\end{itemize}

\section{Hydraulic Press}

We use fluids and their incompressibility to convert pressure into forces to
smash or lift things. A hydraulic press can smash carbon into diamonds. Here
is a picture of the BYU hydraulic press that made the first synthetic
diamonds! \FRAME{dtbpF}{3.2932in}{2.6429in}{0pt}{}{}{Figure}{\special%
{language "Scientific Word";type "GRAPHIC";maintain-aspect-ratio
TRUE;display "USEDEF";valid_file "T";width 3.2932in;height 2.6429in;depth
0pt;original-width 3.2482in;original-height 2.6013in;cropleft "0";croptop
"1";cropright "1";cropbottom "0";tempfilename
'SXRLBU01.wmf';tempfile-properties "XPR";}}A hydraulic jack can lift cars. 
\FRAME{dtbpF}{3.3927in}{1.9934in}{0pt}{}{\Qlb{Hydrolic_jack}}{Figure}{%
\special{language "Scientific Word";type "GRAPHIC";maintain-aspect-ratio
TRUE;display "USEDEF";valid_file "T";width 3.3927in;height 1.9934in;depth
0pt;original-width 3.3468in;original-height 1.9545in;cropleft "0";croptop
"1";cropright "1";cropbottom "0";tempfilename
'SXRLYD02.wmf';tempfile-properties "XPR";}}

Let's take on the hydraulic jack and see if we can mathematically describe
how it works.\footnote{%
Another project in reverse engineering!} In the last figure, you can see a
real Hydraulic jack. The jack is made from two cylinders filled with fluid,
usually an oil. The jack in the picture is a little harder to think about.
Let's draw a simplified jack:

\FRAME{dtbpF}{2.3981in}{2.0306in}{0in}{}{}{Figure}{\special{language
"Scientific Word";type "GRAPHIC";maintain-aspect-ratio TRUE;display
"USEDEF";valid_file "T";width 2.3981in;height 2.0306in;depth
0in;original-width 2.3583in;original-height 1.9917in;cropleft "0";croptop
"1";cropright "1";cropbottom "0";tempfilename
'SVAQIE6E.wmf';tempfile-properties "XPR";}}

We usually have a constant force acting on the jack area. And that force
will be related to the pressure in the liquid.

\begin{equation}
F=PA
\end{equation}%
Let's look at just the left cylinder, the input of our jack. And further
suppose that instead of a sensible liquid like oil, we fill the jack with a
gas that is compressible. \FRAME{dtbpF}{3.659in}{2.5763in}{0pt}{}{}{Figure}{%
\special{language "Scientific Word";type "GRAPHIC";maintain-aspect-ratio
TRUE;display "USEDEF";valid_file "T";width 3.659in;height 2.5763in;depth
0pt;original-width 3.2794in;original-height 2.3013in;cropleft "0";croptop
"1";cropright "1";cropbottom "0";tempfilename
'SXT5DA02.wmf';tempfile-properties "XPR";}}As we press down on the piston,
the force is spread over the piston area. That causes an increase in
pressure, $\Delta P.$ An early researcher named Pascal reasoned that at
first the gas molecules that are near the piston would go flying away from
the piston because of the force of the piston hitting them. But they would
hit other molecules (and the walls) and after a time all the gas molecules
would be effected. When the system comes to equilibrium, we find that not
only has the top pressure by the piston changed by $\Delta P$ so that (using
the figure) 
\begin{equation*}
P_{t2}=P_{t1}+\Delta P
\end{equation*}%
but also the bottom pressure has changed. 
\begin{equation*}
P_{b2}=P_{b1}+\Delta P
\end{equation*}%
This is called Pascal's principle. We can restate it as the change in
pressure in a fluid is transferred to the entire fluid and the magnitude of
the change is the same everywhere.

\begin{center}
\rule{300pt}{1pt}

\textbf{Pascal's Law}: a change in the pressure applied to a fluid is
transmitted undiminished to every point of the fluid and to the walls of the
container.

\rule{300pt}{1pt}
\end{center}

Now let's consider the jack again. In the jack we really have a liquid, and
liquids are not easily compressible. When our jack is just sitting in our
garage, we have just atmospheric pressure pushing on the jack input piston, 
\begin{equation}
F_{atm}=P_{atm}A_{in}
\end{equation}%
so Newton's second law 
\begin{eqnarray*}
F_{net_{y}} &=&ma_{y} \\
&=&\sum_{i}F_{yi}
\end{eqnarray*}%
for the input piston is \FRAME{dtbpF}{0.9968in}{1.6186in}{0pt}{}{}{Figure}{%
\special{language "Scientific Word";type "GRAPHIC";maintain-aspect-ratio
TRUE;display "USEDEF";valid_file "T";width 0.9968in;height 1.6186in;depth
0pt;original-width 1.1075in;original-height 1.8157in;cropleft "0";croptop
"1";cropright "1";cropbottom "0";tempfilename
'TGJ3BT00.wmf';tempfile-properties "XPR";}} 
\begin{equation*}
\Sigma F_{in}=F_{fluid}-F_{atm}-W_{in}=ma_{y}
\end{equation*}%
But $ma_{y}=0$ because our piston isn't yet accelerating so 
\begin{equation}
0=F_{fluid}-P_{atm}A_{in}-W_{in}
\end{equation}%
which tells us that 
\begin{equation*}
F_{fluid}=P_{atm}A_{in}+W_{in}
\end{equation*}%
But we wish to add a force $F_{in}$ to this. Usually this is done using the
lever that you can see in figure (\ref{Hydrolic_jack}) on the small
cylinder. So on the input side, \FRAME{dtbpF}{1.0421in}{1.7293in}{0pt}{}{}{%
Figure}{\special{language "Scientific Word";type
"GRAPHIC";maintain-aspect-ratio TRUE;display "USEDEF";valid_file "T";width
1.0421in;height 1.7293in;depth 0pt;original-width 1.282in;original-height
2.1462in;cropleft "0";croptop "1";cropright "1";cropbottom "0";tempfilename
'TGJ3D701.wmf';tempfile-properties "XPR";}}We have, using Newton's second law%
\begin{equation*}
\Sigma F_{in}=F_{fluid}-F_{atm}-W_{in}-F_{in}=ma_{y}
\end{equation*}%
and think that the force due to the atmosphere is a pressure force so $%
F_{atm}=P_{atm}A_{in}$ 
\begin{equation}
\Sigma F_{in}=F_{fluid}-P_{atm}A_{in}-W_{in}-F_{in}
\end{equation}%
Our $F_{in}$ will cause a change in pressure in the fluid. The piston will
accelerate for a moment, then come to a stop. When it comes to a stop the
pressure will be different because of our force by an amount $\Delta P$ 
\begin{equation}
\Delta P=\frac{F_{in}}{A_{in}}
\end{equation}%
But the first three forces in our Newton's second law haven't changed much,
and before we applied $F_{in}$ they summed to zero $\left(
0=F_{fluid}-P_{atm}A_{in}-W_{in}\right) $. So after we have applied $F_{in}$
we have%
\begin{equation*}
\Sigma F_{in}\approx 0-\Delta PA_{in}
\end{equation*}%
Because we really have a liquid in our jack this time we can't compress the
fluid in the cylinder. We still expect our $\Delta P$ to be transmitted
throughout the entire fluid. By changing the pressure on the input side we
have changed the pressure by $\Delta P$ on the output side (the lifting
side, with the big cylinder). The only way this can happen without
compressing the liquid is for the output piston to move! So the output will
have an amount of force added to it that will make it accelerate briefly.
The output side is shown in the next figure. Before we pushed the lever we
had atmospheric pressure pushing down on the output piston. The output
piston wasn't accelerating. So before we push the lever we have 
\begin{equation}
\Sigma F_{out}=F_{fluid}-F_{atm}-W_{out}=0
\end{equation}%
After we push the lever we will have \FRAME{dtbpF}{3.2153in}{1.8436in}{0pt}{%
}{}{Figure}{\special{language "Scientific Word";type
"GRAPHIC";maintain-aspect-ratio TRUE;display "USEDEF";valid_file "T";width
3.2153in;height 1.8436in;depth 0pt;original-width 3.4508in;original-height
1.9657in;cropleft "0";croptop "1";cropright "1";cropbottom "0";tempfilename
'TGJ3LL02.wmf';tempfile-properties "XPR";}}%
\begin{equation}
\Sigma F_{out}=F_{fluid}-F_{atm}-W_{out}+F_{out}
\end{equation}%
We have a new force $F_{out}=\Delta PA_{out}$ that we have added to our net
force on the output piston. This new force came from the input piston moving
and changing the pressure by $\Delta P.$ Then once the input piston makes
the pressure change there is a net force on the output piston.%
\begin{equation}
\Sigma F_{out}\approx 0+\Delta PA_{out}
\end{equation}%
Let's solve for $\Delta P$ in our two net force equations%
\begin{equation*}
\frac{\Sigma F_{in}}{A_{in}}=\Delta P=\frac{F_{in}}{A_{in}}
\end{equation*}%
and 
\begin{equation}
\frac{\Sigma F_{out}}{A_{out}}=\Delta P=\frac{F_{out}}{A_{out}}
\end{equation}%
Since $\Delta P$ in both equations is the same when the two sides are at the
same elevation, then%
\begin{equation}
\frac{F_{out}}{A_{out}}=\frac{F_{in}}{A_{in}}
\end{equation}

So how does your hydraulic jack work?%
\begin{equation}
F_{out}=F_{in}\frac{A_{out}}{A_{in}}
\end{equation}%
If $A_{out}>A_{in}$ a much smaller $F_{in}$ can produce a large $F_{out}$.%
\footnote{%
You probably already knew that from your experience jacking up cars, but
it's nice to see that math that shows how it works.}

It is important to note that we have assumed our jack fluid is not
compressible. This is mostly true at normal pressures. So the volume of
fluid leaving the cylinder at the input side must be the same as the volume
of fluid entering the output side. We could say 
\begin{equation*}
%TCIMACRO{\TeXButton{V}{\ooalign{\hfil$V$\hfil\cr\kern0.1em--\hfil\cr}}}%
%BeginExpansion
\ooalign{\hfil$V$\hfil\cr\kern0.1em--\hfil\cr}%
%EndExpansion
_{out}=%
%TCIMACRO{\TeXButton{V}{\ooalign{\hfil$V$\hfil\cr\kern0.1em--\hfil\cr}}}%
%BeginExpansion
\ooalign{\hfil$V$\hfil\cr\kern0.1em--\hfil\cr}%
%EndExpansion
_{in}
\end{equation*}%
but 
\begin{eqnarray*}
%TCIMACRO{\TeXButton{V}{\ooalign{\hfil$V$\hfil\cr\kern0.1em--\hfil\cr}}}%
%BeginExpansion
\ooalign{\hfil$V$\hfil\cr\kern0.1em--\hfil\cr}%
%EndExpansion
_{out} &=&A_{out}\Delta y_{out} \\
%TCIMACRO{\TeXButton{V}{\ooalign{\hfil$V$\hfil\cr\kern0.1em--\hfil\cr}}}%
%BeginExpansion
\ooalign{\hfil$V$\hfil\cr\kern0.1em--\hfil\cr}%
%EndExpansion
_{in} &=&A_{in}\Delta y_{in}
\end{eqnarray*}%
Then we can see that 
\begin{equation*}
A_{out}\Delta y_{out}=A_{in}\Delta y_{in}
\end{equation*}%
Since $A_{out}>A_{in,}$ it is clear that the output piston will travel up a
much smaller distance than the input piston moved down. This is why you have
to pump quite a lot on the input side of your jack to move a car a
relatively small distance upward.

Note that our equation 
\begin{equation*}
P_{b}=P_{t}+\rho gh
\end{equation*}%
still applies! The pressure is not the same at the top and bottom of the
cylinders. Pascal's principle tells us that the change in pressure $\Delta P$
was the same in the entire liquid, not that the pressure, itself, is the
same. This is an important distinction! We will need to be careful with
Pascal's principle.

In many cases we can sum up pressure over an area because the pressure will
be the same at every point within the area. Our hydraulic jack is one of
these cases. But let's turn our attention now to pressure variation with
depth again and look at things that float.

\section{Buoyant Forces, Archimedes' Principle}

Remember our parcel of liquid from last lecture. \FRAME{dtbpF}{2.6135in}{%
1.5653in}{0pt}{}{}{Figure}{\special{language "Scientific Word";type
"GRAPHIC";maintain-aspect-ratio TRUE;display "USEDEF";valid_file "T";width
2.6135in;height 1.5653in;depth 0pt;original-width 2.5719in;original-height
1.529in;cropleft "0";croptop "1";cropright "1";cropbottom "0";tempfilename
'SXT7XY04.wmf';tempfile-properties "XPR";}}Let's investigate the net force
on our parcel and let's say our liquid is water this time. The net force
must be zero, or it would be accelerating. Let's let $A_{top}=A_{bottom}=A$
and let's let the sides of the box have area $A_{m}$ ($m$ for
\textquotedblleft middle\textquotedblright ). Then%
\begin{eqnarray*}
\Sigma F_{x} &=&ma_{x}=0=P_{m}A_{m}-P_{m}A_{m} \\
\Sigma F_{y} &=&ma_{y}=0=-P_{top}A+P_{bottom}A-W
\end{eqnarray*}%
Then from the $y$ equation we have 
\begin{equation*}
W=-P_{top}A+P_{bottom}A
\end{equation*}%
But suppose I\ ignore the weight of the parcel of water, $W=mg$. Then there
would be a net force due to all other forces in the $\widehat{\mathbf{y}}$
direction. But our parcel isn't accelerating. That force must be matching
the downward force due to gravity. And that is just what we see 
\begin{equation*}
W=P_{b}A-P_{t}A
\end{equation*}%
The right hand side is the part that is just due to the fluid. The left hand
side is the weight. The right hand side is like a net force that leaves out
weight. It is a net force for just the pressure forces. We could give this
right hand side net force due to all the water pressure forces a symbol, say
\textquotedblleft $B$.\textquotedblright 
\begin{equation*}
B=P_{b}A-P_{t}A
\end{equation*}%
Then $B$ would be equal to the weight of the parcel of water!%
\begin{equation*}
B=W
\end{equation*}

Now take a beach ball sized parcel of water\FRAME{dtbpF}{2.6844in}{2.0306in}{%
0pt}{}{}{Figure}{\special{language "Scientific Word";type
"GRAPHIC";maintain-aspect-ratio TRUE;display "USEDEF";valid_file "T";width
2.6844in;height 2.0306in;depth 0pt;original-width 2.6411in;original-height
1.9917in;cropleft "0";croptop "1";cropright "1";cropbottom "0";tempfilename
'SUS73V8X.wmf';tempfile-properties "XPR";}}The parcel of water is in
equilibrium just as before. We know there is a force on the parcel of water
due to gravity. It is downward. Just as before, the ball-shaped parcel of
water is not sinking so again there is no net force. If we add up all the
forces due to the water pressure only, we must have a net pressure force
that is upward to balance the downward force due to gravity. Again this is
what we called \textquotedblleft $B$.\textquotedblright\ And again $B$ is
equal to the weight of the parcel of water.

The idea of a net force just due to fluid pressure is so useful we give it
its own name We will call the net force due to just pressure from the fluid
the \emph{Buoyant force}.

Note that this is \textbf{not} the net force, it is just the sum of the
forces due to the water pressure. It does not include gravity or any tension
or spring or any other forces. It just contains the pressure forces due to
the fluid.

Note also that the buoyant force is not some new kind of force. It is a name
we give to the sum of the pressure forces acting on an object in the fluid.
Like we call a group of soldiers a battalion or a group of cows a heard, the
sum of the group of pressure forces due to a fluid is given the name
\textquotedblleft buoyant force.\textquotedblright\ But herds and battalions
are not new things, they are groups of cows and soldiers. A Buoyant force is
not a new force, it is the vector sum of a group of pressure forces.

Because the parcel is not sinking, we can determine that in this special
case 
\begin{equation}
B=W_{\text{fluid parcel}}\qquad \text{Floating}
\end{equation}%
for this situation that we call floating (these are magnitudes, are the
directions the same?). The term on the right is the weight of the parcel of
water or fluid.

\FRAME{dtbpF}{1.3413in}{1.8611in}{0in}{}{}{Figure}{\special{language
"Scientific Word";type "GRAPHIC";maintain-aspect-ratio TRUE;display
"USEDEF";valid_file "T";width 1.3413in;height 1.8611in;depth
0in;original-width 1.3067in;original-height 1.8239in;cropleft "0";croptop
"1";cropright "1";cropbottom "0";tempfilename
'TBHUOS00.wmf';tempfile-properties "XPR";}}

Suppose we replace this amount of water with the actual beach ball, and
let's suppose our beach ball is rigid so it won't be compressed.\footnote{%
Yes, I know real beach balls would be compressed by the fluid. But we have a
very special beach ball in our problem!} The weight of the beach ball is
different than the weight of the water parcel. But will any of the pressure
forces be different? The two volumes are exactly the same!

It turns out that the buoyant force will be exactly the same for the beach
ball as it would for the beach ball-shaped parcel of water.

This is great! It means we can calculate the buoyant force by thinking of
replacing the actual object we have with an equal volume parcel of water.
The buoyant force will be equal to the weight of that equal volume of the
fluid.

We could say that to replace the water with the beach ball we have \emph{%
displaced} a beach ball's volume worth of water. The water that was at the
location of the beach ball is now somewhere else, so it is displaced. Then
the buoyant force would be equal to the weight of the water that was
displaced.%
\begin{equation}
B=W_{\text{displaced fluid}}
\end{equation}%
Of course, the weight of the beach ball is far less than the weight of the
water displaced%
\begin{equation*}
W_{ball}<W_{\text{displaced fluid}}
\end{equation*}%
and this is why the beach ball accelerates upward.%
\begin{eqnarray*}
m_{ball}a_{_{^{y}ball}} &=&B-W_{ball} \\
&=&W_{\text{displaced fluid}}-W_{ball}
\end{eqnarray*}

It turns out that this idea of the buoyant force being equal to the weight
of an equal volume of fluid displaced is general! This concept of the
Buoyant force being equal to the weight of the water that fits inside our
volume is called \emph{Archimedes' principle}.

\begin{center}
\rule{300pt}{1pt}

The magnitude of the buoyant force equals the weight of the fluid displaced
by the object.

\rule{300pt}{1pt}
\end{center}

\section{Buoyant Force: Two Cases}

We can gain insight into this concept of a buoyant force by considering more
floating objects. Let's consider an object that is floating submerged in the
fluid (like a fish or submarine) and one that is floating at the surface and
partially extends out of the fluid (like a floating log or boat).

\subsection{Totally submerged object}

\FRAME{dtbpF}{2.0557in}{1.7789in}{0pt}{}{}{Figure}{\special{language
"Scientific Word";type "GRAPHIC";maintain-aspect-ratio TRUE;display
"USEDEF";valid_file "T";width 2.0557in;height 1.7789in;depth
0pt;original-width 2.0176in;original-height 1.7417in;cropleft "0";croptop
"1";cropright "1";cropbottom "0";tempfilename
'SUS73V8Y.wmf';tempfile-properties "XPR";}}Consider a solid object, a block
of some material, placed in water. The weight of the solid object is $W=Mg.$

The magnitude of the buoyant force is equal to the weight of an equal volume
of water. If we imagine replacing the block with a block-sized volume of
water, then 
\begin{eqnarray*}
B &=&m_{_{fluid-displaced}}g \\
&=&\rho _{_{fluid-displaced}}%
%TCIMACRO{\TeXButton{V}{\ooalign{\hfil$V$\hfil\cr\kern0.1em--\hfil\cr}}}%
%BeginExpansion
\ooalign{\hfil$V$\hfil\cr\kern0.1em--\hfil\cr}%
%EndExpansion
_{object}g \\
&=&\rho _{_{fluid-displaced}}g%
%TCIMACRO{\TeXButton{V}{\ooalign{\hfil$V$\hfil\cr\kern0.1em--\hfil\cr}}}%
%BeginExpansion
\ooalign{\hfil$V$\hfil\cr\kern0.1em--\hfil\cr}%
%EndExpansion
_{block}
\end{eqnarray*}%
Think Archimedes: The weight of an equal volume of water was displaced by
the block. So we have the volume of the block as part of our equation for
the buoyant force.

If the object has mass $M$ then we can write the mass of the block as 
\begin{eqnarray*}
M &=&\rho _{object}%
%TCIMACRO{\TeXButton{V}{\ooalign{\hfil$V$\hfil\cr\kern0.1em--\hfil\cr}}}%
%BeginExpansion
\ooalign{\hfil$V$\hfil\cr\kern0.1em--\hfil\cr}%
%EndExpansion
_{object} \\
&=&\rho _{block}%
%TCIMACRO{\TeXButton{V}{\ooalign{\hfil$V$\hfil\cr\kern0.1em--\hfil\cr}}}%
%BeginExpansion
\ooalign{\hfil$V$\hfil\cr\kern0.1em--\hfil\cr}%
%EndExpansion
_{block}
\end{eqnarray*}%
and the weight of the block is then 
\begin{eqnarray*}
W &=&Mg \\
&=&\rho _{object}%
%TCIMACRO{\TeXButton{V}{\ooalign{\hfil$V$\hfil\cr\kern0.1em--\hfil\cr}}}%
%BeginExpansion
\ooalign{\hfil$V$\hfil\cr\kern0.1em--\hfil\cr}%
%EndExpansion
_{object}g \\
&=&\rho _{block}%
%TCIMACRO{\TeXButton{V}{\ooalign{\hfil$V$\hfil\cr\kern0.1em--\hfil\cr}}}%
%BeginExpansion
\ooalign{\hfil$V$\hfil\cr\kern0.1em--\hfil\cr}%
%EndExpansion
_{block}g
\end{eqnarray*}%
The net force is then 
\begin{eqnarray*}
F_{net} &=&B-W \\
&=&\rho _{fluid-displaced}g%
%TCIMACRO{\TeXButton{V}{\ooalign{\hfil$V$\hfil\cr\kern0.1em--\hfil\cr}}}%
%BeginExpansion
\ooalign{\hfil$V$\hfil\cr\kern0.1em--\hfil\cr}%
%EndExpansion
_{object}-\rho _{object}%
%TCIMACRO{\TeXButton{V}{\ooalign{\hfil$V$\hfil\cr\kern0.1em--\hfil\cr}}}%
%BeginExpansion
\ooalign{\hfil$V$\hfil\cr\kern0.1em--\hfil\cr}%
%EndExpansion
_{object}g \\
&=&\rho _{fluid-displaced}g%
%TCIMACRO{\TeXButton{V}{\ooalign{\hfil$V$\hfil\cr\kern0.1em--\hfil\cr}}}%
%BeginExpansion
\ooalign{\hfil$V$\hfil\cr\kern0.1em--\hfil\cr}%
%EndExpansion
_{block}-\rho _{block}%
%TCIMACRO{\TeXButton{V}{\ooalign{\hfil$V$\hfil\cr\kern0.1em--\hfil\cr}}}%
%BeginExpansion
\ooalign{\hfil$V$\hfil\cr\kern0.1em--\hfil\cr}%
%EndExpansion
_{block}g \\
&=&g%
%TCIMACRO{\TeXButton{V}{\ooalign{\hfil$V$\hfil\cr\kern0.1em--\hfil\cr}}}%
%BeginExpansion
\ooalign{\hfil$V$\hfil\cr\kern0.1em--\hfil\cr}%
%EndExpansion
_{block}(\rho _{fluid-displaced}-\rho _{block})
\end{eqnarray*}%
Consider this last equation. If the density of the block is large, larger
than the density of the fluid, then $(\rho _{fluid}-\rho _{block})$ is
negative, and the force will be negative. The block will accelerate
downward. If the density of the block is smaller than the density of the
fluid, then $(\rho _{fluid}-\rho _{block})$ will be positive and the block
will accelerate upward. If the density of the block is just the same as the
density of the fluid, the block will float in place (or at least sink or
float upward at a constant rate--the acceleration is zero when the net force
is zero). Now we can see how a submarine can rise or dive!. It adjusts it's
density to dive or to surface!

We did this calculation for our block, but it could have been any object. So
in general%
\begin{equation}
F_{net}=g%
%TCIMACRO{\TeXButton{V}{\ooalign{\hfil$V$\hfil\cr\kern0.1em--\hfil\cr}}}%
%BeginExpansion
\ooalign{\hfil$V$\hfil\cr\kern0.1em--\hfil\cr}%
%EndExpansion
_{object}(\rho _{fluid-displaced}-\rho _{object})
\end{equation}

If the object has a density that is less than the fluid, the object will be
accelerated upward. And if the density of the object is greater than the
density of the fluid the object will be accelerated downward.%
%TCIMACRO{%
%\TeXButton{bb in a test tube demo}{\marginpar {
%\hspace{-0.5in}
%\begin{minipage}[t]{1in}
%\small{bb in a test tube demo}
%\end{minipage}
%}}}%
%BeginExpansion
\marginpar {
\hspace{-0.5in}
\begin{minipage}[t]{1in}
\small{bb in a test tube demo}
\end{minipage}
}%
%EndExpansion

Note that the motion of a totally submerged object is determined only by the
relative density! In a way, this could be confusing because we did not allow
our object to change volume as it moved. If we had an object that could
change size, then it changes density as it changes size. A balloon full of
air released from the bottom of a swimming pool does change volume as it
rises. So in that case we would need to know $\Delta 
%TCIMACRO{\TeXButton{V}{\ooalign{\hfil$V$\hfil\cr\kern0.1em--\hfil\cr}}}%
%BeginExpansion
\ooalign{\hfil$V$\hfil\cr\kern0.1em--\hfil\cr}%
%EndExpansion
_{object}$ and how the density, $\rho _{object},$ changes as the balloon
moves and expands. But let's not yet allow a change in volume for our
objects.

\subsection{Partially submerged object}

Now let's consider a partially submerged object\FRAME{dtbpF}{1.9545in}{%
1.8204in}{0pt}{}{}{Figure}{\special{language "Scientific Word";type
"GRAPHIC";maintain-aspect-ratio TRUE;display "USEDEF";valid_file "T";width
1.9545in;height 1.8204in;depth 0pt;original-width 1.9164in;original-height
1.7832in;cropleft "0";croptop "1";cropright "1";cropbottom "0";tempfilename
'SUS73V8Z.wmf';tempfile-properties "XPR";}}

Let's assume $\rho _{object}<\rho _{fluid}$ (why? think of $\rho
_{fluid-displaced}-\rho _{object}$). We assume static equilibrium, which
means we are observing this floating object when it is just sitting still
relative to the objects around it, not accelerating.

Since $a_{y}=0$ then 
\begin{equation*}
\Sigma F_{yi}=ma_{y}=0
\end{equation*}%
and%
\begin{equation*}
\Sigma F_{yi}=B-W
\end{equation*}%
So, for an object floating on the surface 
\begin{equation*}
B=W
\end{equation*}

This isn't a surprise for floating objects. Let's call the volume of the
fluid displaced by the object $%
%TCIMACRO{\TeXButton{V}{\ooalign{\hfil$V$\hfil\cr\kern0.1em--\hfil\cr}}}%
%BeginExpansion
\ooalign{\hfil$V$\hfil\cr\kern0.1em--\hfil\cr}%
%EndExpansion
_{fd}$. Now this is not the same as $%
%TCIMACRO{\TeXButton{V}{\ooalign{\hfil$V$\hfil\cr\kern0.1em--\hfil\cr}}}%
%BeginExpansion
\ooalign{\hfil$V$\hfil\cr\kern0.1em--\hfil\cr}%
%EndExpansion
_{object}$ for this case! Part of the object is sticking out of the fluid,
so $%
%TCIMACRO{\TeXButton{V}{\ooalign{\hfil$V$\hfil\cr\kern0.1em--\hfil\cr}}}%
%BeginExpansion
\ooalign{\hfil$V$\hfil\cr\kern0.1em--\hfil\cr}%
%EndExpansion
_{fd}<%
%TCIMACRO{\TeXButton{V}{\ooalign{\hfil$V$\hfil\cr\kern0.1em--\hfil\cr}}}%
%BeginExpansion
\ooalign{\hfil$V$\hfil\cr\kern0.1em--\hfil\cr}%
%EndExpansion
_{object}.$ But, it is still true that the weight of the displaced volume of
fluid gives the buoyant force 
\begin{equation*}
B=\rho _{fd}g%
%TCIMACRO{\TeXButton{V}{\ooalign{\hfil$V$\hfil\cr\kern0.1em--\hfil\cr}}}%
%BeginExpansion
\ooalign{\hfil$V$\hfil\cr\kern0.1em--\hfil\cr}%
%EndExpansion
_{fd}
\end{equation*}

Then returning to our force equation%
\begin{eqnarray*}
B &=&-W \\
\rho _{fd}g%
%TCIMACRO{\TeXButton{V}{\ooalign{\hfil$V$\hfil\cr\kern0.1em--\hfil\cr}}}%
%BeginExpansion
\ooalign{\hfil$V$\hfil\cr\kern0.1em--\hfil\cr}%
%EndExpansion
_{fd} &=&-\rho _{object}%
%TCIMACRO{\TeXButton{V}{\ooalign{\hfil$V$\hfil\cr\kern0.1em--\hfil\cr}}}%
%BeginExpansion
\ooalign{\hfil$V$\hfil\cr\kern0.1em--\hfil\cr}%
%EndExpansion
_{object}g
\end{eqnarray*}%
or%
\begin{equation}
\frac{%
%TCIMACRO{\TeXButton{V}{\ooalign{\hfil$V$\hfil\cr\kern0.1em--\hfil\cr}}}%
%BeginExpansion
\ooalign{\hfil$V$\hfil\cr\kern0.1em--\hfil\cr}%
%EndExpansion
_{fd}}{%
%TCIMACRO{\TeXButton{V}{\ooalign{\hfil$V$\hfil\cr\kern0.1em--\hfil\cr}}}%
%BeginExpansion
\ooalign{\hfil$V$\hfil\cr\kern0.1em--\hfil\cr}%
%EndExpansion
_{object}}=\frac{\rho _{object}}{\rho _{fd}}
\end{equation}%
The fraction of the volume of a floating object that is below the fluid
surface is equal to the ratio of the density of the object to that of the
fluid.

In the case of an iceberg, not all of the ice is visible above the surface.
The density of ice is about $920\frac{\unit{kg}}{\unit{m}^{3}}$ and that of
sea water is about $1000\frac{\unit{kg}}{\unit{m}^{3}}.$ So for ice%
\begin{equation*}
\frac{%
%TCIMACRO{\TeXButton{V}{\ooalign{\hfil$V$\hfil\cr\kern0.1em--\hfil\cr}}}%
%BeginExpansion
\ooalign{\hfil$V$\hfil\cr\kern0.1em--\hfil\cr}%
%EndExpansion
_{fluid}}{%
%TCIMACRO{\TeXButton{V}{\ooalign{\hfil$V$\hfil\cr\kern0.1em--\hfil\cr}}}%
%BeginExpansion
\ooalign{\hfil$V$\hfil\cr\kern0.1em--\hfil\cr}%
%EndExpansion
_{object}}=\frac{920\frac{\unit{kg}}{\unit{m}^{3}}}{1030\frac{\unit{kg}}{%
\unit{m}^{3}}}=0.893\,2
\end{equation*}%
\begin{equation*}
%TCIMACRO{\TeXButton{V}{\ooalign{\hfil$V$\hfil\cr\kern0.1em--\hfil\cr}}}%
%BeginExpansion
\ooalign{\hfil$V$\hfil\cr\kern0.1em--\hfil\cr}%
%EndExpansion
_{fluid}=0.89%
%TCIMACRO{\TeXButton{V}{\ooalign{\hfil$V$\hfil\cr\kern0.1em--\hfil\cr}}}%
%BeginExpansion
\ooalign{\hfil$V$\hfil\cr\kern0.1em--\hfil\cr}%
%EndExpansion
_{object}
\end{equation*}%
that is the fluid displaced is $89\%$ of the ice. Just about $10\%$ of the
ice sticks up out of the water. I \FRAME{dtbpFU}{3.5302in}{2.3931in}{0pt}{%
\Qcb{{\protect\small \ Iceberg with both beautiful blue-green submerged
portion and a reflection of the surface ice and snow. Approximately 90 per
cent of the iceberg is submerged. Antarctica, Palmer Peninsula, Northern
area. (Dr. Mike Goebel, NOAA NMFS SWFSC, NOAA NMFS SWFSC Antarctic Marine
Living Resources (AMLR) Program.)}}}{}{Figure}{\special{language "Scientific
Word";type "GRAPHIC";maintain-aspect-ratio TRUE;display "USEDEF";valid_file
"T";width 3.5302in;height 2.3931in;depth 0pt;original-width
3.5692in;original-height 2.4099in;cropleft "0";croptop "1";cropright
"1";cropbottom "0";tempfilename 'SUS73V90.wmf';tempfile-properties "XPR";}}

We have gained some good insight using pressure. But think, pressure is a
force spread over an area. And we know forces from Principles of Physics I
(PH121). But we didn't just use forces in Principles of Physics I. We also
found we could use energy to solve motion problems. In our next lecture,
let's try to bring the powerful technique of conservation of energy into our
study of fluid flow.

\chapter{26 Conservation of Energy for Fluid Flow 1.14.5, 1.14.6, 1.14.7}

%TCIMACRO{%
%\TeXButton{\chaptermark{Energy in Fluids}}{\chaptermark{Energy in Fluids}}}%
%BeginExpansion
\chaptermark{Energy in Fluids}%
%EndExpansion

So far, we have only dealt with fluids in equilibrium. The topic of fluid
dynamics is a complicated mathematical field that requires differential
equations to do with any exactness. There are whole upper division classes
that study this called \emph{Fluid Dynamics}\textbf{. }But with some
simplifying assumptions, we can get a feel for how fluids flow.

%TCIMACRO{%
%\TeXButton{Fundamental Concepts}{\vspace{0.10in}\hspace{-0.37in}{\Large Fundamental Concepts\vspace{0.2in}}}}%
%BeginExpansion
\vspace{0.10in}\hspace{-0.37in}{\Large Fundamental Concepts\vspace{0.2in}}%
%EndExpansion

\begin{enumerate}
\item Laminar Flow

\item Turbulent Flow

\item Viscosity

\item Continuity

\item Bernoulli's Equation

\item Reynolds number gives a measure of how laminar or turbulen our flow is.
\end{enumerate}

\section{Moving Fluids}

To start our study of moving fluids, let's make some definitions.

New Definitions: \FRAME{dtbpF}{2.7086in}{3.5803in}{0in}{}{}{Figure}{\special%
{language "Scientific Word";type "GRAPHIC";maintain-aspect-ratio
TRUE;display "USEDEF";valid_file "T";width 2.7086in;height 3.5803in;depth
0in;original-width 2.6671in;original-height 3.5336in;cropleft "0";croptop
"1";cropright "1";cropbottom "0";tempfilename
'SUS73V91.wmf';tempfile-properties "XPR";}}

\begin{enumerate}
\item Laminar: Steady flow, if each particle of the fluid follows a smooth
path, such that the paths of the different particles never cross each other.

\item Turbulent: non-laminar flow characterized by whirlpool or eddy regions

\item Viscosity: term describing the internal friction of a fluid (honey is
more viscus than water)

\item Irrotational: flow having no angular momentum

\item Streamline: a line indicating the path taken by a particle under
steady flow
\end{enumerate}

In the figure above we can see streamlines going around airfoils. The reason
we can see them is that small particles are being deposited by evenly spaced
sources in an air tunnel. This makes it so we can see the streamlines. Here
is the same type of imaging for a soccer ball (futbol) being tested by NASA. 
\FRAME{dtbpFU}{3.1412in}{3.1456in}{0pt}{\Qcb{{\protect\small NASA Ames
measuring turbulence from a soccor ball. There is a cool video at this URL!
https://www.nasa.gov/centers-and-facilities/ames/nasa-turns-world-cup-into-lesson-in-aerodynamics/%
}}}{}{Figure}{\special{language "Scientific Word";type
"GRAPHIC";maintain-aspect-ratio TRUE;display "USEDEF";valid_file "T";width
3.1412in;height 3.1456in;depth 0pt;original-width 3.7048in;original-height
3.7101in;cropleft "0";croptop "1";cropright "1";cropbottom "0";tempfilename
'SXT9N809.wmf';tempfile-properties "XPR";}}

The last two figures show a way to tell if we have laminar flow around an
object. Small particles are inserted into the wind or fluid that passes by
the object. In laminar flow the stream of particles just goes around the
object. But with turbulent flow the stream of particles breaks up. We can
see this for an airfoil shape and for a soccer ball (futbol).

Sadly, we will limit our study to laminar, non-viscus, irrotational,
incompressible fluids. This is a little bit like when we said we would study
frictionless surfaces in Principles of Physics I! But to go beyond
non-viscus laminar flow we need the math of partial differential equations,
which is not a prerequisite for this class. So we will have to content
ourselves with this restriction (mostly). For those of you who are
interested in learning more (or are ME\ majors) you can take our junior
level fluid dynamic course taught by the Mechanical Engineering Department.

\section{Equation of continuity}

Let's start our discussion of fluid flow with a pipe that changes diameter
along it's length.

\FRAME{dtbpF}{2.1638in}{2.6083in}{0pt}{}{\Qlb{Continuity}}{Figure}{\special%
{language "Scientific Word";type "GRAPHIC";maintain-aspect-ratio
TRUE;display "USEDEF";valid_file "T";width 2.1638in;height 2.6083in;depth
0pt;original-width 2.1248in;original-height 2.5668in;cropleft "0";croptop
"1";cropright "1";cropbottom "0";tempfilename
'SUS73V92.wmf';tempfile-properties "XPR";}}

And just for good measure we can have the pipe go up in position. We allow
an ideal (non-compressible) fluid to flow through the pipe in laminar flow.
The segment of fluid is the shaded part in the top part of figure \ref%
{Continuity}. This is like our parcel of fluid that we used before, only now
it is moving. The whole pipe has water in it. But we are going to pay
attention to the shaded part of the fluid in the pipe.

In a time $\Delta t$ the fluid at the left hand side moves a distance 
\begin{equation*}
\Delta x_{1}=v_{1}\Delta t
\end{equation*}

If $A_{1}$ is the area of the pipe at the left hand side (LHS) then the mass
contained in the grey shaded region (bottom left part of figure \ref%
{Continuity}) is%
\begin{equation*}
m_{1}=\rho 
%TCIMACRO{\TeXButton{V}{\ooalign{\hfil$V$\hfil\cr\kern0.1em--\hfil\cr}}}%
%BeginExpansion
\ooalign{\hfil$V$\hfil\cr\kern0.1em--\hfil\cr}%
%EndExpansion
_{1}=\rho A_{1}\Delta x_{1}=\rho A_{1}v_{1}\Delta t
\end{equation*}%
where $\rho $ is the density of the fluid. Likewise, for the grey RHS region
in the bottom of figure \ref{Continuity} 
\begin{equation*}
m_{2}=\rho 
%TCIMACRO{\TeXButton{V}{\ooalign{\hfil$V$\hfil\cr\kern0.1em--\hfil\cr}}}%
%BeginExpansion
\ooalign{\hfil$V$\hfil\cr\kern0.1em--\hfil\cr}%
%EndExpansion
_{2}=\rho A_{2}\Delta x_{2}=\rho A_{2}v_{2}\Delta t
\end{equation*}%
Unless we create, destroy, or pool mass, the mass that crosses $A_{2}$ in $%
\Delta t$ must equal the mass that crosses $A_{1}$ in $\Delta t$%
\begin{eqnarray*}
m_{1} &=&m_{2} \\
\rho A_{1}v_{1}\Delta t &=&\rho A_{2}v_{2}\Delta t
\end{eqnarray*}%
or%
\begin{equation}
A_{1}v_{1}=A_{2}v_{2}
\end{equation}%
This is the \emph{equation of continuity for fluids} (remember that we
assumed ideal incompressible fluids).

Since we assumed that no mass was created, destroyed, or pooled as the fluid
flowed, this is what we might call \textquotedblleft conservation of
mass.\textquotedblright\ And as usual, when something is conserved, we can
use it to solve problems. With this equation, knowing the diameter of a pipe
and how it changes tells us about the speed of the fluid in the pipe and how
it changes. Of course, there could be situations where mass is not conserved
(say, a leaky pipe). So conservation of mass is an assumption that we have
to check. But for now let's assume that the equation of continuity holds for
our problems.

The quantity $Av$ has a name. We call it the flow rate and for some reason I
don't know we give it the symbol $\mathbb{Q}$%
\begin{equation*}
\mathbb{Q}=Av
\end{equation*}

\section{Bernoulli's Equation}

It would be great if we could find the pressure as well as the speed of the
fluid inside our non-uniform pipe. We would need an equation that included
both pressure and speed. Let's try to find such an equation!

In PH121 we developed to ways of looking at motion problems. One was the
force picture where we had to deal with the vector nature of forces and
motion. The other was the energy picture where we could solve problems
without considering directions (but we gave up being able to find
directions). Usually the energy picture made solving problems easier. And
except for the flow direction, we don't really want to know the direction of
the motion of each molecule in our fluid, so giving up direction doesn't
seem too bad for fluid flow.

The work energy theorem is something we know from Principles of Physics I.
Let's start there, 
\begin{equation*}
w=\Delta K+\Delta U
\end{equation*}%
We still need to use forces to find the work done on the parcel of water,
but if we choose our axes carefully, we can keep the work equation in one
dimension. \FRAME{dtbpF}{2.1923in}{1.1857in}{0pt}{}{}{Figure}{\special%
{language "Scientific Word";type "GRAPHIC";maintain-aspect-ratio
TRUE;display "USEDEF";valid_file "T";width 2.1923in;height 1.1857in;depth
0pt;original-width 3.058in;original-height 1.6414in;cropleft "0";croptop
"1";cropright "1";cropbottom "0";tempfilename
'SUS73V93.wmf';tempfile-properties "XPR";}}

\subsection{Work done in $\Delta t$}

Let's review a little of what we learned in Principles of Physics I. Suppose
we have two characters, Normal Person and Super Person. Normal Person and
Super Person have the same mass (but it is distributed differently).

\FRAME{dtbpF}{1.7339in}{1.2721in}{0pt}{}{}{Figure}{\special{language
"Scientific Word";type "GRAPHIC";maintain-aspect-ratio TRUE;display
"USEDEF";valid_file "T";width 1.7339in;height 1.2721in;depth
0pt;original-width 2.1949in;original-height 1.6034in;cropleft "0";croptop
"1";cropright "1";cropbottom "0";tempfilename
'SXTBDC0A.wmf';tempfile-properties "XPR";}}And suppose they want to get to
the top of a tall building. Normal Person has to take the stairs, but in
some incredible way Super Person can jump right up to the top. \FRAME{dtbpF}{%
2.3549in}{2.5944in}{0pt}{}{}{Figure}{\special{language "Scientific
Word";type "GRAPHIC";maintain-aspect-ratio TRUE;display "USEDEF";valid_file
"T";width 2.3549in;height 2.5944in;depth 0pt;original-width
2.3709in;original-height 2.6157in;cropleft "0";croptop "1";cropright
"1";cropbottom "0";tempfilename 'SUS73V95.wmf';tempfile-properties "XPR";}}%
Who has done more work?

We know the change in potential energy is the same for both. And we know
that 
\begin{equation*}
w=-\Delta U
\end{equation*}%
so the amount of work is the same for both cases! Now let's return to our
fluid flow. \FRAME{dtbpF}{4.1105in}{1.5679in}{0pt}{}{}{Figure}{\special%
{language "Scientific Word";type "GRAPHIC";maintain-aspect-ratio
TRUE;display "USEDEF";valid_file "T";width 4.1105in;height 1.5679in;depth
0pt;original-width 3.6322in;original-height 1.369in;cropleft "0";croptop
"1";cropright "1";cropbottom "0";tempfilename
'SXTBQF0B.wmf';tempfile-properties "XPR";}}

On the left hand side (LHS) of our pipe we have a force due to pressure from
the fluid in the rest of the LHS\ of the pipe. That fluid has molecules that
will bang in to our small marked parcel of fluid. The force from the LHS
fluid in the rest of the pipe will be 
\begin{equation}
F_{1}=P_{1}A_{1}\widehat{\mathbf{x}}
\end{equation}

On the right hand side (RHS) there is also more fluid in the rest of the RHS
of the pipe. Our fluid will bang into that fluid, and since the fluid is
non-compressible, the rest of the fluid in the RHS\ will push back. The
force will be%
\begin{equation}
F_{2}=-P_{2}A_{2}\widehat{\mathbf{x}}
\end{equation}

From Principles of Physics I we remember that work is 
\begin{equation*}
w=\int \overrightarrow{\mathbf{F}}\cdot \overrightarrow{\mathbf{dx}}
\end{equation*}%
and since our force is not changing this would be simply%
\begin{equation}
w=\overrightarrow{\mathbf{F}}\cdot \overrightarrow{\mathbf{\Delta x}}
\end{equation}%
And for a pressure force moving a small parcel of fluid a distance $\Delta x$
this would be 
\begin{equation*}
w=PA\Delta x
\end{equation*}%
and this is equal to a change in potential energy%
\begin{equation*}
\Delta U=-PA\Delta x
\end{equation*}%
for the fluid.

We need to find the work done in moving our particular parcel of fluid from
the left hand side to the right hand side. As the parcel moves in the $%
\widehat{\mathbf{x}}$ direction the amount of work will slowly change (like
Normal Person going up the stairs) but really the amount of work only
depends on the beginning and ending conditions (like Super Person who does
the same amount of work but skips the stairs). So we could ignore all the
details of moving the parcel and envision the parcel jumping from the left
up to the right side of the pipe (like in the Super Person case).

Then for the LHS 
\begin{equation*}
w_{1}=P_{1}A_{1}\Delta x_{1}
\end{equation*}

and for the RHS 
\begin{equation*}
w_{2}=-P_{2}A_{2}\Delta x_{2}
\end{equation*}%
The total work (total change in potential energy) will be the sum of these
two works. We can add them together to get the total work done on the parcel
of fluid. Thinking about this work like a change in potential energy helps.
We can just add up the changes in work energy due to the work done in moving
the fluid.

But before we add them together, let's make some substitutions that will
make our final formula for the total work more meaningful. Lets define the
volume of our fluid segment again as

\begin{equation}
%TCIMACRO{\TeXButton{V}{\ooalign{\hfil$V$\hfil\cr\kern0.1em--\hfil\cr}}}%
%BeginExpansion
\ooalign{\hfil$V$\hfil\cr\kern0.1em--\hfil\cr}%
%EndExpansion
=A_{1}\Delta x_{1}
\end{equation}%
Note that this volume is NOT the volume of the entire segment. It is just
the part that is $\Delta x_{1}$ long. Then from the same arguments that lead
to the equation of continuity, we see that the volume of the marked part of
the fluid must not change as it flows%
\begin{equation*}
%TCIMACRO{\TeXButton{V}{\ooalign{\hfil$V$\hfil\cr\kern0.1em--\hfil\cr}}}%
%BeginExpansion
\ooalign{\hfil$V$\hfil\cr\kern0.1em--\hfil\cr}%
%EndExpansion
=A_{1}\Delta x_{1}=A_{2}\Delta x_{2}
\end{equation*}%
The we can write the work equations as%
\begin{equation*}
w_{1}=P_{1}%
%TCIMACRO{\TeXButton{V}{\ooalign{\hfil$V$\hfil\cr\kern0.1em--\hfil\cr}}}%
%BeginExpansion
\ooalign{\hfil$V$\hfil\cr\kern0.1em--\hfil\cr}%
%EndExpansion
\end{equation*}%
and 
\begin{equation*}
w_{2}=-P_{2}%
%TCIMACRO{\TeXButton{V}{\ooalign{\hfil$V$\hfil\cr\kern0.1em--\hfil\cr}}}%
%BeginExpansion
\ooalign{\hfil$V$\hfil\cr\kern0.1em--\hfil\cr}%
%EndExpansion
\end{equation*}%
So finally, the total work done is 
\begin{eqnarray}
w &=&w_{1}+w_{2} \\
&=&P_{1}%
%TCIMACRO{\TeXButton{V}{\ooalign{\hfil$V$\hfil\cr\kern0.1em--\hfil\cr}}}%
%BeginExpansion
\ooalign{\hfil$V$\hfil\cr\kern0.1em--\hfil\cr}%
%EndExpansion
-P_{2}%
%TCIMACRO{\TeXButton{V}{\ooalign{\hfil$V$\hfil\cr\kern0.1em--\hfil\cr}} }%
%BeginExpansion
\ooalign{\hfil$V$\hfil\cr\kern0.1em--\hfil\cr}
%EndExpansion
\notag \\
&=&\left( P_{1}-P_{2}\right) 
%TCIMACRO{\TeXButton{V}{\ooalign{\hfil$V$\hfil\cr\kern0.1em--\hfil\cr}} }%
%BeginExpansion
\ooalign{\hfil$V$\hfil\cr\kern0.1em--\hfil\cr}
%EndExpansion
\notag
\end{eqnarray}%
so, from the work-energy theorem%
\begin{eqnarray*}
w &=&\Delta K+\Delta U \\
\left( P_{1}-P_{2}\right) 
%TCIMACRO{\TeXButton{V}{\ooalign{\hfil$V$\hfil\cr\kern0.1em--\hfil\cr}} }%
%BeginExpansion
\ooalign{\hfil$V$\hfil\cr\kern0.1em--\hfil\cr}
%EndExpansion
&=&\Delta K+\Delta U
\end{eqnarray*}%
We have an expression for the left hand side of our work-energy theorem! Now
for the right hand side $\Delta K+\Delta U$

\subsection{Kinetic Energy}

\FRAME{dtbpF}{2.4664in}{1.8049in}{0pt}{}{\Qlb{fluids shaded region}}{Figure}{%
\special{language "Scientific Word";type "GRAPHIC";maintain-aspect-ratio
TRUE;display "USEDEF";valid_file "T";width 2.4664in;height 1.8049in;depth
0pt;original-width 2.4249in;original-height 1.7677in;cropleft "0";croptop
"1";cropright "1";cropbottom "0";tempfilename
'SUS73V97.wmf';tempfile-properties "XPR";}}If we think about the energy of
the parcel, we will realize that after time $\Delta t,$ it is as though most
of the fluid did not move. We can't tell one part of the fluid from another.
The shaded region in figure \ref{fluids shaded region} is occupied before
and after $\Delta t.$ We can treat this problem as if we moved fluid from
region 1 to region 2 (Super Person vs. Normal Person again!) and left the
rest of the segment alone!

So, ignoring the rest of the fluid, we can see that if the velocity of our
parcel of water changed, then the kinetic energy of the parcel must change
as well

\begin{equation}
\Delta K=\frac{1}{2}mv_{2}^{2}-\frac{1}{2}mv_{1}^{2}
\end{equation}%
(the mass is the same, as we found before).

\subsection{Potential Energy}

\FRAME{dtbpF}{2.911in}{1.2739in}{0pt}{}{}{Figure}{\special{language
"Scientific Word";type "GRAPHIC";maintain-aspect-ratio TRUE;display
"USEDEF";valid_file "T";width 2.911in;height 1.2739in;depth
0pt;original-width 3.9583in;original-height 1.7158in;cropleft "0";croptop
"1";cropright "1";cropbottom "0";tempfilename
'SUS73V98.wmf';tempfile-properties "XPR";}}

Let's look at the change in potential energy for the parcel of fluid. Once
again we can consider just the beginning and ending case (think of Super
Person and Normal Person one more time). We work as though we moved a mass $%
m=m_{1}=m_{2}$ from $y_{1}$ to $y_{2}.$

\begin{equation}
\Delta U=mgy_{2}-mgy_{1}
\end{equation}%
Since our pipe went upward as the parcel of water flows, we gained
gravitational potential energy. But if the pipe went downward we would loose
potential energy.

\subsection{Total Work done on the system}

Let's now assemble our work-energy theorem. We had found the work so far, 
\begin{equation*}
\left( P_{1}-P_{2}\right) 
%TCIMACRO{\TeXButton{V}{\ooalign{\hfil$V$\hfil\cr\kern0.1em--\hfil\cr}}}%
%BeginExpansion
\ooalign{\hfil$V$\hfil\cr\kern0.1em--\hfil\cr}%
%EndExpansion
=\Delta K+\Delta U
\end{equation*}%
but now we know $\Delta K$ and $\Delta U$ We have calculated $w$, $\Delta K$%
, and $\Delta U$ for our parcel of water, so let's substitute in what we
have found. 
\begin{equation}
\left( \left( P_{1}-P_{2}\right) 
%TCIMACRO{\TeXButton{V}{\ooalign{\hfil$V$\hfil\cr\kern0.1em--\hfil\cr}}}%
%BeginExpansion
\ooalign{\hfil$V$\hfil\cr\kern0.1em--\hfil\cr}%
%EndExpansion
\right) =\left( \frac{1}{2}mv_{2}^{2}-\frac{1}{2}mv_{1}^{2}\right) +\left(
mgy_{2}-mgy_{1}\right)
\end{equation}%
This equation describes our fluid flow much like the equation of continuity,
but now it has pressure involved! We have achieved our goal. But we can make
this equation a little easier to understand if we play a trick by dividing
by $%
%TCIMACRO{\TeXButton{V}{\ooalign{\hfil$V$\hfil\cr\kern0.1em--\hfil\cr}}}%
%BeginExpansion
\ooalign{\hfil$V$\hfil\cr\kern0.1em--\hfil\cr}%
%EndExpansion
$ 
\begin{equation}
\left( \left( P_{1}-P_{2}\right) \right) =\left( \frac{1}{2}\frac{m}{%
%TCIMACRO{\TeXButton{V}{\ooalign{\hfil$V$\hfil\cr\kern0.1em--\hfil\cr}}}%
%BeginExpansion
\ooalign{\hfil$V$\hfil\cr\kern0.1em--\hfil\cr}%
%EndExpansion
}v_{2}^{2}-\frac{1}{2}\frac{m}{%
%TCIMACRO{\TeXButton{V}{\ooalign{\hfil$V$\hfil\cr\kern0.1em--\hfil\cr}}}%
%BeginExpansion
\ooalign{\hfil$V$\hfil\cr\kern0.1em--\hfil\cr}%
%EndExpansion
}v_{1}^{2}\right) +\left( \frac{m}{%
%TCIMACRO{\TeXButton{V}{\ooalign{\hfil$V$\hfil\cr\kern0.1em--\hfil\cr}}}%
%BeginExpansion
\ooalign{\hfil$V$\hfil\cr\kern0.1em--\hfil\cr}%
%EndExpansion
}gy_{2}-\frac{m}{%
%TCIMACRO{\TeXButton{V}{\ooalign{\hfil$V$\hfil\cr\kern0.1em--\hfil\cr}}}%
%BeginExpansion
\ooalign{\hfil$V$\hfil\cr\kern0.1em--\hfil\cr}%
%EndExpansion
}gy_{1}\right)
\end{equation}%
and recalling that 
\begin{equation*}
\rho =\frac{m}{%
%TCIMACRO{\TeXButton{V}{\ooalign{\hfil$V$\hfil\cr\kern0.1em--\hfil\cr}}}%
%BeginExpansion
\ooalign{\hfil$V$\hfil\cr\kern0.1em--\hfil\cr}%
%EndExpansion
}
\end{equation*}%
Then%
\begin{equation*}
\left( P_{1}-P_{2}\right) =\left( \frac{1}{2}\rho v_{2}^{2}-\frac{1}{2}\rho
v_{1}^{2}\right) +\left( \rho gy_{2}-\rho gy_{1}\right)
\end{equation*}

Notice that we have subscripts $1$ and $2.$ Let's put all the $1$'s on one
side of our equation and all the $2$'s on the other side. 
\begin{equation}
P_{1}+\frac{1}{2}\rho v_{1}^{2}+\rho gy_{1}=P_{2}+\frac{1}{2}\rho
v_{2}^{2}+\rho gy_{2}
\end{equation}%
which completes our goal. We have an equation that contains the velocity 
\emph{and} the pressure of the fluid.

But notice! The left hand side looks just the same as the right hand side
except for the subscripts! We have a before and after picture like we did in
energy and momentum problems back in Principles of Physics I. Something is
being conserved! Sometimes you see this written as 
\begin{equation}
P+\frac{1}{2}\rho v^{2}+\rho gy=\text{constant}
\end{equation}%
This is another conservation equation. The thing that is being conserved is $%
P+\frac{1}{2}\rho v^{2}+\rho gy.$ We have not named this quantity, but
whatever it is, it is not changing as the fluid flows. Our experience with
conservation equations tells us that this new conservation law will help us
solve problems. This equation is so useful we give it a name. It is called 
\emph{Bernoulli's equation} (after the scientist that came up with it). And
really, what is being conserved is energy. We found Bernoulli's equation
from conservation of energy. But this form of conservation of energy is
useful for solving pressure and fluid flow problems.

\section{Viscosity}

Way back in Principles of Physics I we learned about friction. So far we
haven't allowed any friction with our fluids. But we know this can't be
reality! We know there is drag force. So there must be some friction. And we
can think of there being two different cases where friction can matter.
There is the friction between parts of the water, itself, and there must be
something like a drag force between the water and the solid earth or a
container wall. Let's take on the water/earth friction first.

In the past, we used two different equations for drag for fluids, one for air%
\begin{equation*}
D=\frac{1}{2}\rho CAv^{2}
\end{equation*}%
where $\rho $ is the air density, $C$ is the drag coefficient, $A$ is the
cross sectional area, and $v$ is the speed of an object through air. And we
had a different one for a liquid like oil when we had oscillators being
damped by the fluid 
\begin{equation*}
D=bv
\end{equation*}

\FRAME{dtbpF}{2.8106in}{1.7832in}{0pt}{}{}{Figure}{\special{language
"Scientific Word";type "GRAPHIC";maintain-aspect-ratio TRUE;display
"USEDEF";valid_file "T";width 2.8106in;height 1.7832in;depth
0pt;original-width 3.5345in;original-height 2.2321in;cropleft "0";croptop
"1";cropright "1";cropbottom "0";tempfilename
'SW62VF08.wmf';tempfile-properties "XPR";}}where $b$ is the damping
coefficient and $v$ is the speed. These two equations, one for a gas fluid
like air and one for a thicker liquid like oil, show that friction for
fluids can be quite different. And of course it is much easier to get
through air than oil, so that is not a surprise. Recall that we call
friction in liquids \emph{viscosity} and we would say the oil is more viscus
than the air.

There will be a drag force between the river bed and the lowest water. And
that drag force will likely be more of the $D=bv$ type for a liquid. This
slows the water that is right next to the river bed down to a $v$ that is
quite slow.

But we know fluids don't just have friction with solid surfaces. Because the
fluid can move around internally there is internal friction or viscosity.
When a river flows down its bed we know the drag force at the bottom causes
the water to slow down. But this slower water is also viscus and so it slows
the water just above it a little. So the water at the bottom of the river
will be slow and the water just above it will be less slow. The water above
this second layer will be faster still, and this keeps going until we reach
the top of the water where there is some friction with the air, but it is
much less. So the fastest water is on top. This works for air too.

You might have tried flying a kite. \FRAME{dtbpFU}{3.1656in}{3.0784in}{0pt}{%
\Qcb{{\protect\small Kite flown on the June 19th 2025 BUYU-I High Altitude
Research Team flight as part of the high altitude balloon payload retrieval
effort. The kite can carry radios that connect to the payload trackers.}}}{}{%
Figure}{\special{language "Scientific Word";type
"GRAPHIC";maintain-aspect-ratio TRUE;display "USEDEF";valid_file "T";width
3.1656in;height 3.0784in;depth 0pt;original-width 5.3419in;original-height
5.1949in;cropleft "0";croptop "1";cropright "1";cropbottom "0";tempfilename
'SY5VCP00.wmf';tempfile-properties "XPR";}}When you first start, you have to
run with the kite. But once the kite is higher, the winds tend to be higher
and the kite can stay up without you running.\FRAME{dtbpF}{2.7008in}{1.7149in%
}{0pt}{}{}{Figure}{\special{language "Scientific Word";type
"GRAPHIC";maintain-aspect-ratio TRUE;display "USEDEF";valid_file "T";width
2.7008in;height 1.7149in;depth 0pt;original-width 2.6584in;original-height
1.6777in;cropleft "0";croptop "1";cropright "1";cropbottom "0";tempfilename
'SY5WD601.wmf';tempfile-properties "XPR";}}

Both cases work well as long as nothing disrupts the fluid flow. This nice
well behaved flow we described as being \emph{laminar}. In laminar fluid
flow layers may have different speeds, but they flow without mixing.

All this works pretty well if we don't have turbulence and the viscosity
isn't too high.

We can measure the viscosity of a liquid by placing it between two sheets

\FRAME{dtbpF}{3.0856in}{1.2142in}{0pt}{}{}{Figure}{\special{language
"Scientific Word";type "GRAPHIC";maintain-aspect-ratio TRUE;display
"USEDEF";valid_file "T";width 3.0856in;height 1.2142in;depth
0pt;original-width 3.0415in;original-height 1.1796in;cropleft "0";croptop
"1";cropright "1";cropbottom "0";tempfilename
'SW67AO0C.wmf';tempfile-properties "XPR";}}If a force, $F$ is applied to the
top sheet we have a sheer force if the sheet moves horizontally. The
viscosity makes the fluid move with it. The fluid at the bottom won't move
at all, but the fluid at the top moves with speed $v$ and the layers in
between will have speeds in between $v$ and $0$. The viscosity $\eta $ is
given by 
\begin{equation*}
\eta =\frac{F\Delta y}{vA}
\end{equation*}

In our society we put a lot of effort to make fluids flow in pipes. Our
water comes and goes in pipes. Fuel flows in pipes. Air flows in organ pipes
during General Conference, etc. Last lecture we learned that the rate at
which the water flows is called a flow rate and unfortunately it is often
given the symbol $\mathbb{Q}.$ But what makes the fluid flow? Generally is
is pressure. We recall that pressure is a stress, a force spread over an
area.

\begin{equation*}
P=\frac{F}{A}
\end{equation*}%
\FRAME{dtbpF}{3.9522in}{0.64in}{0in}{}{}{Figure}{\special{language
"Scientific Word";type "GRAPHIC";maintain-aspect-ratio TRUE;display
"USEDEF";valid_file "T";width 3.9522in;height 0.64in;depth
0in;original-width 3.9029in;original-height 0.6088in;cropleft "0";croptop
"1";cropright "1";cropbottom "0";tempfilename
'SW7TSU04.wmf';tempfile-properties "XPR";}}

For our fluid in a pipe the area will be the pipe cross sectional area. We
would like to know how the pressure affects the flow rate, $\mathbb{Q}.$ It
turns out that this is fairly hard to solve (there are multivariate
differential equations involved) so for those who are curious or need to be
able to do this professionally I recommend taking a fluid dynamics class%
\footnote{%
At BYU-I this class is ME360 and is called Fluid Mechanics.} but for our
class let's just borrow a result.

\begin{equation*}
\mathbb{Q}=\frac{P_{in}-P_{out}}{R}
\end{equation*}%
This is called Poiseuille's Law. We can see the pressure in, $P_{in}$ and
the pressure out, $P_{out}.$ The denominator, $R$ is the viscosity
resistance. This equation won't work for fluids with no viscosity, because $%
R $ would be zero and the flow rate infinite. But fluids with no viscosity
are rare, so it works for many normal situations. We would not expect the
flow to be uniform across the pipe. Fluid closer to the pipe walls should
move slower. Fluid near the center should move faster. We can model the
resistance in the pipe by borrowing again from a fluid dynamics class 
\begin{equation*}
R=\frac{8\eta L}{\pi r^{4}}
\end{equation*}%
where $\eta $ is the viscosity of the fluid, $L$ is the length of the pipe,
and $r$ is the radius of the pipe. We can put this simple model for $R$ into
Poiseuille's law 
\begin{eqnarray*}
Q &=&\frac{P_{in}-P_{out}}{\frac{8\eta L}{\pi r^{4}}} \\
&=&\frac{\pi r^{4}\left( P_{in}-P_{out}\right) }{8\eta L}
\end{eqnarray*}%
We can see that this gives us a dependence on the radius. The wider the
pipe, the more flow we get. The more viscus the fluid, the less the flow
rate. Notice also that longer pipes have lower flow. Friction will have more
of an affect if the fluid interacts with the pipe longer, so this makes
sense. And of course if we increase the pressure on one end, the flow
changes. \FRAME{dtbpFU}{2.0678in}{2.5469in}{0pt}{\Qcb{{\protect\small Pipe
with no viscosity (top) and with more realistic viscosity (bottom)}}}{}{%
Figure}{\special{language "Scientific Word";type
"GRAPHIC";maintain-aspect-ratio TRUE;display "USEDEF";valid_file "T";width
2.0678in;height 2.5469in;depth 0pt;original-width 2.0289in;original-height
2.5054in;cropleft "0";croptop "1";cropright "1";cropbottom "0";tempfilename
'SXTEGX0C.wmf';tempfile-properties "XPR";}}

The math that lead to this equation assumed laminar flow. But what if we
don't have laminar flow?

\subsection{Turbulence}

We said we would ignore turbulence. But if we are cautious we can at least
get some intuition about what will happen if we have friction or non-laminar
flow.

We can disrupt laminar flow by, say, putting a rock in the way of the fluid.
Here is a rock in a river.

\FRAME{dtbpF}{2.8003in}{1.8161in}{0pt}{}{}{Figure}{\special{language
"Scientific Word";type "GRAPHIC";maintain-aspect-ratio TRUE;display
"USEDEF";valid_file "T";width 2.8003in;height 1.8161in;depth
0pt;original-width 3.9729in;original-height 2.5668in;cropleft "0";croptop
"1";cropright "1";cropbottom "0";tempfilename
'SW67YQ0E.wmf';tempfile-properties "XPR";}}Now the laminar flow is disrupted
by the rock and the layers will mix. We call a situation like this \emph{%
turbulent }flow. As you can imagine, the equations for turbulent flow are
harder. Finding the flow rate for turbulent flow is really a topic for a
fluid dynamics class. But we can give a measure of how turbulent our
situation is. 
\begin{equation*}
R_{e}=\frac{\rho vD}{\eta }
\end{equation*}%
where $\rho $ is the fluid density, $v$ is the flow speed, $D$ is the pipe
diameter, and of course $\eta $ is the viscosity of the fluid. This quantity
is called the Reynolds number and if this number is high enough the flow
will be turbulent. Engineering books give a table like the following for
determining if the flow is still laminar.%
\begin{equation*}
\begin{tabular}{ll}
laminar & $R_{e}<2300$ \\ 
Transient & $2300<R_{e}<4000$ \\ 
Turbulent & $R_{e}>4000$%
\end{tabular}%
\end{equation*}%
although not all books agree on the value for the transition to turbulence.
But generally if $R_{e}<2300$ our equations for this class work and for $%
R_{e}>2300$ we don't have laminar flow and our equations really don't work.
The transient range is a transition to turbulence and for $R_{e}$ above $%
3000 $ to $4000$ most agree that the flow is turbulent. For many engineering
designs we really don't want turbulence. So this is enough to do some good
work. For cool physics problems turbulence is a form of \textquotedblleft
complexity\textquotedblright\ which is a current emphasis in physics
research.

Fluid dynamics is the movement of a bulk of fluid. This is all we can do for
fluid dynamics in this class. But we can still look at the motion of the
molecules and atoms inside of the fluids (and solids). So lets move on to
thinking about what we will call internal motion in the fluid. It turns out
that this internal motion is related to how hot or cold our fluid is. The
motion of the fluid will have kinetic energy, but it's not like a ball being
thrown or a wave on water. The kinetic energy will have a randomness to it.
Still it is energy and if energy is related to hotness or coldness this is a
comfort because we know something about how to use energy to solve problems.
The study of this internal energy is what we will take on next. We will call
this the thermodynamics of the fluid!

\chapter{27 Temperature 2.1.1, 2.1.2, and 2.1.3}

Last lecture we reviewed some basics of how materials are made and how
chemists and physicists and material scientists determine how many moving
parts a sample of a material might have. We also talked about measuring
temperature. But didn't actually talk about what temperature is or even how
temperature measuring devices work. We will start with measurement
techniques in this lecture and build up to a model of what temperature is.

%TCIMACRO{\TeXButton{\chaptermark{Temperature}}{\chaptermark{Temperature}}}%
%BeginExpansion
\chaptermark{Temperature}%
%EndExpansion

%TCIMACRO{%
%\TeXButton{Fundamental Concepts}{\vspace{0.10in}\hspace{-0.37in}{\Large Fundamental Concepts\vspace{0.2in}}}}%
%BeginExpansion
\vspace{0.10in}\hspace{-0.37in}{\Large Fundamental Concepts\vspace{0.2in}}%
%EndExpansion

\begin{itemize}
\item Number Density

\item Atomic Mass

\item Atomic Mass Number

\item The Mole

\item Temperature

\item Kelvin Temperature Scale

\item volume of a liquid

\item Pressure of a gas at constant volume

\item Volume of a gas at constant pressure

\item The color of an object and temperature

\item Phase Changes and Phase Diagrams

\item Thermal Expansion
\end{itemize}

\section{Temperature and the Zeroth Law of Thermodynamics}

We need a basis for talking about thermal systems. We sort of have an
intuition for temperature, so let's use our intuition for now.

\FRAME{dtbpF}{2.2053in}{1.5895in}{0pt}{}{}{Figure}{\special{language
"Scientific Word";type "GRAPHIC";maintain-aspect-ratio TRUE;display
"USEDEF";valid_file "T";width 2.2053in;height 1.5895in;depth
0pt;original-width 2.1646in;original-height 1.5541in;cropleft "0";croptop
"1";cropright "1";cropbottom "0";tempfilename
'TJXLXK06.wmf';tempfile-properties "XPR";}}In the figure we have two
objects. I measure the temperature of both objects. I find that they are the
same. Suppose I\ put the two objects in \textquotedblleft
contact.\textquotedblright

\FRAME{dtbpF}{2.2105in}{1.6665in}{0pt}{}{}{Figure}{\special{language
"Scientific Word";type "GRAPHIC";maintain-aspect-ratio TRUE;display
"USEDEF";valid_file "T";width 2.2105in;height 1.6665in;depth
0pt;original-width 4.5117in;original-height 3.3944in;cropleft "0";croptop
"1";cropright "1";cropbottom "0";tempfilename
'TJXLXK07.wmf';tempfile-properties "XPR";}}

Would you expect there to be energy transferred between the two objects? For
most of us our intuition says \textquotedblleft no.\textquotedblright\ Given
this intuition, lets make some more formal definitions:

\begin{itemize}
\item Thermal Contact: Two objects are in thermal contact if they can
exchange energy due to a temperature difference

\item Thermal Equilibrium : A state where two objects would not exchange
energy by heat or electromagnetic radiation if they were placed in thermal
contact

\item Temperature: The property that determines if an object is in thermal
equilibrium with other objects (two objects in thermal equilibrium have the
same temperature)
\end{itemize}

Given all this, we can state a law of thermal dynamics:

\begin{center}
\rule{300pt}{1pt}

\textbf{If objects A and B are separately in thermal equilibrium with a
third object C, then A and B are in thermal equilibrium with each other.}

\rule{300pt}{1pt}
\end{center}

Again this fits our intuition, if we have two containers of water that are
both at the same temperature as a third container of water, we expect they
all have the same temperature.

In a sense, this seems obvious, so obvious that it was always assumed
without statement. But in retrospect, it is a fundamental concept that
deserves it's place as one of the laws of thermodynamics. By the time this
was acknowledged, the names \textquotedblleft first law\textquotedblright\
and \textquotedblleft second\ law\textquotedblright\ were taken, so this
became the \emph{Zeroth law of Thermodynamics.}

Before we move on, let's review some things we know from high school
chemistry.

\section{Numbers of Atoms and Number Densities}

Atoms are small. They are really incredibly small. A typical atom measures a
few Angstroms across. An angstrom is $1\times 10^{-10}\unit{m}.$ That is
unbelievably small. Let's try to get an understanding of this scale. Picture
a meter stick. Then divide it up into $10000000000$ equal units. That is how
small an atoms is. Or, consider scaling our meter stick to be the size of
the continental US, that is about $4.\,\allowbreak 087\,7\times 10^{6}\unit{m%
}$. Then for our US\ sized meter stick, on this scale an atom would still
only be about half a millimeter across! About the size of a cake sprinkle.

It is no surprise that in a typical sample of a solid we can have something
like $10^{25}$ atoms. We use the capital letter $N$ to mean the number of
atoms in a sample of matter.

Because these numbers are so large, it is sometimes convenient to talk about
how many atoms there are in a small volume of a substance. This is often
less than in the whole sample. We can say, for example, that there are $1000$
atoms per cubic micrometer of material. We call such a ratio of number of
atoms per volume a \emph{number density. }In the SI system, the units for
the number of atoms per unit volume are $1/\unit{m}^{3}.$ So number
densities can still be quite large! Number densities are more useful for
samples where the density is uniform, but we can talk about substances who's
number density changes. Unfortunately you will often see a small $n$ used to
mean number density. This over use of $n$ is confusing, so sometimes we will
write number density as $N/%
%TCIMACRO{\TeXButton{V}{\ooalign{\hfil$V$\hfil\cr\kern0.1em--\hfil\cr}}}%
%BeginExpansion
\ooalign{\hfil$V$\hfil\cr\kern0.1em--\hfil\cr}%
%EndExpansion
$ so it is obvious what we mean.

\subsection{Atomic Mass}

You probably ran into atomic mass in High School. The idea of atomic mass is
that each nucleon (proton or neutron) has about the same mass, and the
electrons don't have much mass at all, so we can sort of count nucleons and
say that an atom has a mass in units of nucleon masses. Carbon has 12
nucleons (six protons and six neutrons) so it would have a mass of $12\unit{u%
}.$ The unit $\unit{u}$ is the atomic mas unit, the mass of one nucleon.
Because of the electrons (and other effects)\ this is not an exact approach,
and it is not good enough if we attempt to do nuclear physics. But it is
close enough for us for now. So we can say that hydrogen $\left(
^{1}H\right) $ has an atomic mass of $1\unit{u}.$ If we have a molecule, we
can just add the atomic masses of each atom in the molecule together to find
the molecule's mass. Note that to be more accurate, we can use a periodic
table and find the atomic mass, or, if we need many digits of accuracy, we
can go to a table of the nucleotides. We can convert these masses to SI
units using%
\begin{equation*}
1\unit{u}=1.66\times 10^{-27}\unit{kg}
\end{equation*}

\subsection{Moles}

Chances are that you are familiar with the concept of a mole as well. But
experience has shown me that the mole will be new to some people. So let's
start with something we know, a dozen.

What is a dozen? Well, it is $12$ of something. We could have a dozen
doughnuts or a dozen eggs. or a dozen large white pickup trucks. It is just
an amount. It is $12$ of whatever we are counting.

A mole is like a dozen, but it is a BIG\ dozen. Is is $6.02\times 10^{23}$
of something. I\ could have $6.02\times 10^{23}$ doughnuts, or $6.02\times
10^{23}$ eggs. I would need quite a few hungry friends to eat it all. But a
mole is just a number of something. The mole is not a very practical amount
for everyday objects like doughnuts or trucks. But think of the large number
of atoms in a sample of matter. For this situation, a mole seems like a
reasonable number to use.

Like we say I want to buy a dozen eggs, we could say we want to buy a mole
of hydrogen atoms. Since atoms are so small, packaging them in big numbers
makes sense.

This number, $6.02\times 10^{23},$ is called \emph{Avogadro's number}%
\begin{equation*}
N_{A}=6.02\times 10^{23}
\end{equation*}%
To find the number of dozens of a something, we take the number of items and
divide by 12.%
\begin{equation*}
n_{dozen}=\frac{N}{12}
\end{equation*}%
To find the number of moles of a something, we take the number of items and
divide by Avogadro's number%
\begin{equation*}
n_{moles}=\frac{N}{N_{A}}
\end{equation*}%
Usually we will use small $n$ to be the number of moles of something. This
could be a problem because $n$ can also mean number density. So we will have
to be careful to state what we mean with our variables.

\subsection{Molar mass}

We could give a mass for a dozen eggs. If each egg has the same mass (more
or less) then the dozenal mass would be the mass of a dozen eggs and it
would be twelve times the mass of one egg. The dozenal mass of a dozen
trucks would be larger than the dozenal mass of eggs because the mass of a
single truck is larger than the mass of single egg, and the mass of a dozen
trucks is $12$ times the mass of one truck.

We can also define a molar mass, which would be the mass of $6.02\times
10^{23}$ of something. Usually, we choose the units for this molar mass to
be the mass \emph{in grams} of a mole of something. And when we do this, the
molar mass is equal to the numerical value of the atomic or molecular mass.
for example, the molar mass of He, with $m=4\unit{u}$ is $M=4\unit{g}/\unit{%
mol}.$ The molar mass of diatomic nitrogen would be $28\unit{g}/\unit{mol}$
because each nitrogen atom has an atomic mass of $m=14\unit{u}$ and we add
the two atomic masses together for the molecular mass. With this definition
of molar mass we can see that for a sample of a material with mass $%
m_{sample}$ 
\begin{equation*}
n_{moles}=\frac{m_{sample}\left( \text{in grams}\right) }{M\left( \text{in
grams per mol}\right) }
\end{equation*}

Now that we have these new terms to use, let's take up measuring temperature
in our next lecture.

\section{How to build a thermometer}

We will find that several material properties change with temperature

\begin{equation*}
\begin{tabular}{|c|}
\hline
\textbf{Things that Vary with Temperature} \\ \hline
Volume of a liquid \\ \hline
The dimensions of a solid \\ \hline
The pressure of a gas at constant volume \\ \hline
The volume of a gas at constant pressure \\ \hline
The electric resistance of a conductor \\ \hline
The color of an object \\ \hline
The size of an object \\ \hline
\end{tabular}%
\end{equation*}

Since these things change with temperature, we can use any of these to make
a temperature measuring device. Think of old fashioned thermometers, how do
they work?

\FRAME{dtbpF}{2.4474in}{0.2421in}{0pt}{}{}{Figure}{\special{language
"Scientific Word";type "GRAPHIC";maintain-aspect-ratio TRUE;display
"USEDEF";valid_file "T";width 2.4474in;height 0.2421in;depth
0pt;original-width 2.4059in;original-height 0.2136in;cropleft "0";croptop
"1";cropright "1";cropbottom "0";tempfilename
'SUS73W9N.wmf';tempfile-properties "XPR";}}Old fashioned mercury
thermometers (and their glycerine replacements) use the first property,
volume of a liquid increases with temperature. The liquid is placed in an
evacuated tube, and the amount of the tube that is filled depends on the
temperature. The liquid fills more of the tube when the temperature goes up
because the liquid volume gets bigger with temperature. By placing a ruled
background behind the tube or even by making rule marks on the outside of
the tube, we can find out how much the liquid volume changed and calculate
the temperature change.

\subsection{Celsius temperature scale}

To make our old fashioned mercury thermometer give meaningful values we need
to have units of temperature to mark on the tube. One of the sets of units
for temperature is called the \emph{Celsius temperature scale.} The most
important thing to say about the Celsius scale is not to use it for doing
thermodynamic problems (The same goes double for the Fahrenheit scale). But
since it is common to see Celsius temperatures, we really have to talk about
them.

\FRAME{dtbpF}{3.0193in}{2.6583in}{0in}{}{}{Figure}{\special{language
"Scientific Word";type "GRAPHIC";maintain-aspect-ratio TRUE;display
"USEDEF";valid_file "T";width 3.0193in;height 2.6583in;depth
0in;original-width 3.0486in;original-height 2.6805in;cropleft "0";croptop
"1";cropright "1";cropbottom "0";tempfilename
'SUS73W9O.wmf';tempfile-properties "XPR";}}

The idea is to make marks on our thermometer tube that show how the change
in liquid volume is related to temperature. To do this for the Celsius scale
we put our thermometer in ice water and mark where the top of the liquid is.
And we assign the value of $0\unit{%
%TCIMACRO{\U{2103}}%
%BeginExpansion
{}^{\circ}{\rm C}%
%EndExpansion
}$ to this mark. Then we take our thermometer and put it in boiling water.
The liquid expands and that makes it move up the tube. Where it stops, we
make a mark and label this mark $100\unit{%
%TCIMACRO{\U{2103}}%
%BeginExpansion
{}^{\circ}{\rm C}%
%EndExpansion
}.$ This scale sets the zero point at the ice point of water and sets $100%
\unit{%
%TCIMACRO{\U{2103}}%
%BeginExpansion
{}^{\circ}{\rm C}%
%EndExpansion
}$ at the steam point of water. Then we make $100$ evenly spaced divisions
between these points. We can also mark the tube with the same division
spacing below the $0\unit{%
%TCIMACRO{\U{2103}}%
%BeginExpansion
{}^{\circ}{\rm C}%
%EndExpansion
}$ mark and above the $100\unit{%
%TCIMACRO{\U{2103}}%
%BeginExpansion
{}^{\circ}{\rm C}%
%EndExpansion
}$ mark. But this calibration does not work well if high accuracy is
required. There are often large problems when temperatures are measured
outside of the $0\unit{%
%TCIMACRO{\U{2103}}%
%BeginExpansion
{}^{\circ}{\rm C}%
%EndExpansion
}$ to $100\unit{%
%TCIMACRO{\U{2103}}%
%BeginExpansion
{}^{\circ}{\rm C}%
%EndExpansion
}$ range.

Worse yet, water only boils at $100\unit{%
%TCIMACRO{\U{2103}}%
%BeginExpansion
{}^{\circ}{\rm C}%
%EndExpansion
}$ at sea level, and does not always freeze at $0\unit{%
%TCIMACRO{\U{2103}}%
%BeginExpansion
{}^{\circ}{\rm C}%
%EndExpansion
}.$

Even worse, there is nothing really very $0$-ish about $0\unit{%
%TCIMACRO{\U{2103}}%
%BeginExpansion
{}^{\circ}{\rm C}%
%EndExpansion
}.$ We picked it because it was a nice value for water and for us humans
water is important. But we can certainly have lower temperatures than $0%
\unit{%
%TCIMACRO{\U{2103}}%
%BeginExpansion
{}^{\circ}{\rm C}%
%EndExpansion
}.$\footnote{%
We often do in Rexburg from about November to March.} This makes $0\unit{%
%TCIMACRO{\U{2103}}%
%BeginExpansion
{}^{\circ}{\rm C}%
%EndExpansion
}$ more like an origin of a coordinate system. We could say that the
intersection of Main St. and Center St. is our origin, but there is noting
really zeroish about that intersection. But, it turns out there \emph{is} a
natural zero point for temperature and it is not the ice point of water. A
better temperature scale would start at the natural zero point. Let's see if
we can make such a scale.

\subsection{The Constant-Volume Gas Thermometer and the Absolute Temperature
Scale}

We wish to define a temperature scale that is based on a physical zero
point, an point where there really is not temperature. To do this we need a
better thermometer. So here is one in the next figure.\FRAME{dtbpF}{3.6832in%
}{1.5229in}{0pt}{}{}{Figure}{\special{language "Scientific Word";type
"GRAPHIC";maintain-aspect-ratio TRUE;display "USEDEF";valid_file "T";width
3.6832in;height 1.5229in;depth 0pt;original-width 3.6357in;original-height
1.4866in;cropleft "0";croptop "1";cropright "1";cropbottom "0";tempfilename
'SY672300.wmf';tempfile-properties "XPR";}}This device uses the rise in
pressure at a constant volume (the number 3 property from our list at the
beginning of the lecture). It works by placing the gas bulb in thermal
contact with the object who's temperature is to be measured (reservoir $B$
in the figure\footnote{%
A reservoir is just a large container of something.}). A scale (like a
ruler) is placed in the apparatus such that the mercury in reservoir $A$ on
the left is at the zero point. We could take water at its ice point, for
example, to find this first part of the scale. Then the bulb is placed in
thermal contact with another object by swapping out what is in reservoir $B$%
, say we use water at its steam point this time. The mercury will move
because the pressure increases in the bulb, and the gas tries to expand. We
add mercury into reservoir $A$ (just when you thought you were done with
U-tubes!) until the level the reservoir is again at the zero point. The
height of the additional mercury we added is proportional to the change in
pressure. The Pressure for both measurements is calculated and plotted
against the temperature and a linear curve is drawn between the two
pressure-temperature points. This curve serves as a calibration for other
temperatures. We could fill the bulb with different gasses. If we do this,
we get different sloped lines. You can see this in the next figure.

But think! it seems like with this new thermometer design we should have
zero temperature if we have zero gas in the bulb. We have different lines
for each gas, but if we then extend these lines back to where $P=0$ (so that
we have no gas), we find they intersect with $P=0$ at $T=-273.15\unit{%
%TCIMACRO{\U{2103}}%
%BeginExpansion
{}^{\circ}{\rm C}%
%EndExpansion
}.$ This is the basis we need for an absolute temperature scale! We have
found a place where the temperature really should be zero because we have no
gas to have a temperature!

\FRAME{dtbpF}{3.3492in}{1.9886in}{0in}{}{}{Figure}{\special{language
"Scientific Word";type "GRAPHIC";maintain-aspect-ratio TRUE;display
"USEDEF";valid_file "T";width 3.3492in;height 1.9886in;depth
0in;original-width 3.3847in;original-height 1.9984in;cropleft "0";croptop
"1";cropright "1";cropbottom "0";tempfilename
'SUS73W9Q.wmf';tempfile-properties "XPR";}}

We will take this common zero pressure point, $T=273.15\unit{%
%TCIMACRO{\U{2103}}%
%BeginExpansion
{}^{\circ}{\rm C}%
%EndExpansion
}$ as our zero point for our improved temperature scale, but we need another
point. We will choose the triple point of water, the one point where water,
ice and steam exist in equilibrium; this happens at $0.01\unit{%
%TCIMACRO{\U{2103}}%
%BeginExpansion
{}^{\circ}{\rm C}%
%EndExpansion
}$ and $4.58\unit{mmHg}=$ $610.\,\allowbreak 61\unit{Pa}.$ We will call this
point $273.15$ on our new scale which seems kind of random. But think, with
this choice each unit of temperature will then be $1/273.15^{th}$ of the
distance between our zero point and the water triple point. This way the
degrees in this scale are the same size as the degrees in the Celsius scale
(which is kind of nice). This new scale that starts at the physical zero
point we will call the Kelvin Temperature scale and the units are Kelvins ($%
\unit{K}$). This is the temperature system we really should use for our work.

\subsection{Temperature scale conversions}

We will need some conversion factors so you can take temperatures from the
Celsius and Fahrenheit scales to the Kelvin scale.%
\begin{eqnarray}
T_{C} &=&T_{K}-273.15 \\
T_{F} &=&\frac{9}{5}T_{C}+32\unit{%
%TCIMACRO{\U{2109}}%
%BeginExpansion
{}^{\circ}{\rm F}%
%EndExpansion
} \\
\Delta T_{C} &=&\Delta T_{K}=\frac{5}{9}\Delta T_{F}
\end{eqnarray}

Since the Kelvin scale is not the scale your doctor or your meteorologist
uses, it is less familiar. We should try to get some feeling for the Kelvin
scale. Here are some interesting temperatures placed on the Kelvin scale.

\FRAME{dtbpF}{2.2237in}{2.3017in}{0pt}{}{}{Figure}{\special{language
"Scientific Word";type "GRAPHIC";maintain-aspect-ratio TRUE;display
"USEDEF";valid_file "T";width 2.2237in;height 2.3017in;depth
0pt;original-width 2.237in;original-height 2.3168in;cropleft "0";croptop
"1";cropright "1";cropbottom "0";tempfilename
'SUS73W9R.wmf';tempfile-properties "XPR";}}

\section{Thermal Expansion}

%TCIMACRO{%
%\TeXButton{Ring and Ball Demo}{\marginpar {
%\hspace{-0.5in}
%\begin{minipage}[t]{1in}
%\small{Ring and Ball Demo}
%\end{minipage}
%}}}%
%BeginExpansion
\marginpar {
\hspace{-0.5in}
\begin{minipage}[t]{1in}
\small{Ring and Ball Demo}
\end{minipage}
}%
%EndExpansion

We still only kind of know what temperature is. But we can now measure
temperature, so suppose we increase the temperature of a substance. We could
take the mercury in a thermometer, for example. When the temperature is
higher, the mercury takes up more space. It expands. This isn't just true
for liquids. Here is the change in length for a laser crystal as a function
of temperature. \FRAME{dtbpF}{3.5976in}{2.7043in}{0pt}{}{}{Figure}{\special%
{language "Scientific Word";type "GRAPHIC";maintain-aspect-ratio
TRUE;display "USEDEF";valid_file "T";width 3.5976in;height 2.7043in;depth
0pt;original-width 3.55in;original-height 2.661in;cropleft "0";croptop
"1";cropright "1";cropbottom "0";tempfilename
'SY689202.wmf';tempfile-properties "XPR";}}Notice that the change in length
makes a curve that bends. We would prefer to make our thermometers out of a
material that has a linear expansion. That is why we used mercury in
thermometers for many years. Mercury has expansion that is linear in
temperature for many situations. Twice as hot gives twice the volume change.
Usually when the expansion length, $\Delta L,$ is much less than the
original dimensions of the object $\left( L_{i}\right) $, we can consider
the expansion to be linear in the change in temperature $\Delta T$. Of
course we could increase the temperature above this nice linear region, and
we could get interesting effects like melting or exploding. But for this
lecture we will limit our discussion to a linear change in length with
temperature for now. Let's see how this linear region of expansion works.

\subsection{Expansion Coefficient}

What does it mean to have a linear change in length with temperature? 
\begin{equation}
L_{f}=L_{i}+\alpha L_{i}\Delta T  \label{Linear Expansion}
\end{equation}%
It means that the final length of the object is equal to the initial length
of the object plus a little bit more that depends on $\Delta T$ and not $%
\Delta T^{2}$ or $\Delta T^{3}$ or some other function of $\Delta T.$ We
will get straight line plots instead of curved plots as $T$ changes. You can
see that this is the case for the equation above for thermal expansion. We
can write it next to an equation for a straight line to show the linear
nature of the equation%
\begin{eqnarray*}
y &=&b+mx \\
L_{f} &=&L_{i}+\alpha L_{i}\Delta T
\end{eqnarray*}%
If we plot $L_{f}$ vs. $\Delta T$ this would make a straight line. The slope
of the line would be $\alpha L_{i}$ so the amount of expansion depends on
how big the object was to begin with. It also depends on what material we
have. That is what the $\alpha $ tells us. Different substances expand at
different rates, even if they have the same initial length. So if you have a
metal cavity of length $L_{i}$ and fill it with a laser crystal of length $%
L_{i}$ and heat them both up, you may find the crystal expands at a
different rate than the metal fasteners when the temperature rises--and may
need to buy a replacement for the shattered crystal.\footnote{%
It was a very sad day when this happened to me personally.}. The coefficient 
$\alpha $ is different for every substance, because every substance expands
a little differently than other substances. We should give $\alpha $ a name.
It is called the \emph{average coefficient of linear expansion} where from
equation (\ref{Linear Expansion}) we can see that 
\begin{eqnarray*}
\alpha &\equiv &\frac{\frac{\Delta L}{L_{i}}}{\Delta T} \\
&=&\frac{\Delta L}{L_{i}\Delta T}
\end{eqnarray*}%
where $\Delta L$ is the change in length due to a temperature change $\Delta
T.$ And of course we already know $L_{i}$ is the original length of the
sample.

As long as $\Delta T$ is \textquotedblleft not too big\textquotedblright\ $%
\alpha $ is constant. We can write this as%
\begin{equation*}
L_{f}-L_{i}=\alpha L_{i}\left( T_{f}-T_{i}\right)
\end{equation*}%
where the subscript $f$ means final and $i$ means initial as usual. The
units of $\alpha $ are inverse temperature (unfortunately it is often in $%
\unit{%
%TCIMACRO{\U{2103}}%
%BeginExpansion
{}^{\circ}{\rm C}%
%EndExpansion
},$ but we would prefer to use $\unit{K}$). Laboratories publish tables of
average coefficients of linear expansion and manufacturers give this as a
specification for the materials they produce. Values for a few materials are
given in the table below%
\begin{equation*}
\begin{tabular}{ll}
Material & Average $\alpha \left( \unit{K}^{-1}\right) $ \\ 
Aluminum & $25\times 10^{-6}$ \\ 
Copper & $17\times 10^{-6}$ \\ 
Gold & $14\times 10^{-6}$ \\ 
Lead & $29\times 10^{-6}$ \\ 
Steel & $11\times 10^{-6}$ \\ 
Brass & $18.7\times 10^{-6}$ \\ 
quartz & $0.4\times 10^{-6}$ \\ 
Glass & $9\times 10^{-6}$ \\ 
Concrete & $12\times 10^{-6}$%
\end{tabular}%
\end{equation*}

Suppose we increase the temperature of something that has a hole in it. Does
the hole increase or decrease in size? A cavity in a piece of material
expands in the same way as if the cavity were filled with the material.
Let's see that this must be true.

\subsection{Volume expansion}

Suppose we have a cube with sides of length $L_{i}.$ What happens when we
increase the cube's temperature? Each side will become larger. So what
happens to the volume?%
\begin{eqnarray*}
%TCIMACRO{\TeXButton{V}{\ooalign{\hfil$V$\hfil\cr\kern0.1em--\hfil\cr}}}%
%BeginExpansion
\ooalign{\hfil$V$\hfil\cr\kern0.1em--\hfil\cr}%
%EndExpansion
_{i}+\Delta 
%TCIMACRO{\TeXButton{V}{\ooalign{\hfil$V$\hfil\cr\kern0.1em--\hfil\cr}} }%
%BeginExpansion
\ooalign{\hfil$V$\hfil\cr\kern0.1em--\hfil\cr}
%EndExpansion
&=&\left( L_{i}+\Delta L\right) \left( L_{i}+\Delta L\right) \left(
L_{i}+\Delta L\right) \\
&=&\left( L_{i}+\alpha L_{i}\Delta T\right) \left( L_{i}+\alpha L_{i}\Delta
T\right) \left( L_{i}+\alpha L_{i}\Delta T\right) \\
&=&L_{i}^{3}\left( 1+\alpha \Delta T\right) \left( 1+\alpha \Delta T\right)
\left( 1+\alpha \Delta T\right) \\
&=&L_{i}^{3}\left( \allowbreak \Delta T^{3}\alpha ^{3}+3\Delta T^{2}\alpha
^{2}+3\Delta T\alpha +1\right) \\
&=&%
%TCIMACRO{\TeXButton{V}{\ooalign{\hfil$V$\hfil\cr\kern0.1em--\hfil\cr}}}%
%BeginExpansion
\ooalign{\hfil$V$\hfil\cr\kern0.1em--\hfil\cr}%
%EndExpansion
_{i}\left( \allowbreak \Delta T^{3}\alpha ^{3}+3\Delta T^{2}\alpha
^{2}+3\Delta T\alpha +1\right)
\end{eqnarray*}%
Now let's divide both sides by $%
%TCIMACRO{\TeXButton{V}{\ooalign{\hfil$V$\hfil\cr\kern0.1em--\hfil\cr}}}%
%BeginExpansion
\ooalign{\hfil$V$\hfil\cr\kern0.1em--\hfil\cr}%
%EndExpansion
_{i}$ 
\begin{eqnarray*}
1+\frac{\Delta 
%TCIMACRO{\TeXButton{V}{\ooalign{\hfil$V$\hfil\cr\kern0.1em--\hfil\cr}}}%
%BeginExpansion
\ooalign{\hfil$V$\hfil\cr\kern0.1em--\hfil\cr}%
%EndExpansion
}{%
%TCIMACRO{\TeXButton{V}{\ooalign{\hfil$V$\hfil\cr\kern0.1em--\hfil\cr}}}%
%BeginExpansion
\ooalign{\hfil$V$\hfil\cr\kern0.1em--\hfil\cr}%
%EndExpansion
_{i}} &=&\allowbreak \Delta T^{3}\alpha ^{3}+3\Delta T^{2}\alpha
^{2}+3\Delta T\alpha +1 \\
\frac{\Delta 
%TCIMACRO{\TeXButton{V}{\ooalign{\hfil$V$\hfil\cr\kern0.1em--\hfil\cr}}}%
%BeginExpansion
\ooalign{\hfil$V$\hfil\cr\kern0.1em--\hfil\cr}%
%EndExpansion
}{%
%TCIMACRO{\TeXButton{V}{\ooalign{\hfil$V$\hfil\cr\kern0.1em--\hfil\cr}}}%
%BeginExpansion
\ooalign{\hfil$V$\hfil\cr\kern0.1em--\hfil\cr}%
%EndExpansion
_{i}} &=&\allowbreak \Delta T^{3}\alpha ^{3}+3\Delta T^{2}\alpha
^{2}+3\Delta T\alpha
\end{eqnarray*}%
and let's make an approximation. For normal conditions, $\alpha \Delta T\ll
1 $ (good for $T<373.\,\allowbreak 15\unit{K}$ ($100\unit{%
%TCIMACRO{\U{2103}}%
%BeginExpansion
{}^{\circ}{\rm C}%
%EndExpansion
}$)). Then $\alpha ^{2}\Delta T^{2}$ and $\alpha ^{3}\Delta T^{3}$ are
really really small. Let's ignore them.%
\begin{equation}
\frac{\Delta 
%TCIMACRO{\TeXButton{V}{\ooalign{\hfil$V$\hfil\cr\kern0.1em--\hfil\cr}}}%
%BeginExpansion
\ooalign{\hfil$V$\hfil\cr\kern0.1em--\hfil\cr}%
%EndExpansion
}{%
%TCIMACRO{\TeXButton{V}{\ooalign{\hfil$V$\hfil\cr\kern0.1em--\hfil\cr}}}%
%BeginExpansion
\ooalign{\hfil$V$\hfil\cr\kern0.1em--\hfil\cr}%
%EndExpansion
_{i}}\approx \allowbreak 3\alpha \Delta T
\end{equation}%
We can write this as 
\begin{eqnarray}
\Delta 
%TCIMACRO{\TeXButton{V}{\ooalign{\hfil$V$\hfil\cr\kern0.1em--\hfil\cr}} }%
%BeginExpansion
\ooalign{\hfil$V$\hfil\cr\kern0.1em--\hfil\cr}
%EndExpansion
&\approx &\allowbreak 3\alpha 
%TCIMACRO{\TeXButton{V}{\ooalign{\hfil$V$\hfil\cr\kern0.1em--\hfil\cr}}}%
%BeginExpansion
\ooalign{\hfil$V$\hfil\cr\kern0.1em--\hfil\cr}%
%EndExpansion
_{i}\Delta T \\
&=&\beta 
%TCIMACRO{\TeXButton{V}{\ooalign{\hfil$V$\hfil\cr\kern0.1em--\hfil\cr}}}%
%BeginExpansion
\ooalign{\hfil$V$\hfil\cr\kern0.1em--\hfil\cr}%
%EndExpansion
_{i}\Delta T  \notag
\end{eqnarray}%
where $\beta $ is called the \emph{average coefficient of volume expansion.}

We could do the same for a two dimensional metal plate (in the approximation
that a plate is 2D!)

\begin{equation}
\Delta A=2\alpha A_{i}\Delta T
\end{equation}

We need to be careful that we don't use these linear expansion equations
outside of their valid range. For example, suppose we build a satellite and
launch it into a low Earth orbit (LEO). The space environment for a LEO
orbit (sun side to shade side) has a very large $\Delta T.$

\begin{equation}
\begin{tabular}{l}
$T_{sun}=120\unit{%
%TCIMACRO{\U{2103}}%
%BeginExpansion
{}^{\circ}{\rm C}%
%EndExpansion
}=393.\,\allowbreak 15\unit{K}$ \\ 
$T_{shade}=-100\unit{%
%TCIMACRO{\U{2103}}%
%BeginExpansion
{}^{\circ}{\rm C}%
%EndExpansion
}=173.\,\allowbreak 15\unit{K}$%
\end{tabular}%
\end{equation}%
This $\Delta T$ range is too large for most materials to still be linear.
Special low $\alpha $ materials are used to allow satellite structures to
survive

But how about on Earth? Even on Earth we can have extremes

\begin{equation}
\begin{tabular}{l}
$T_{\max }=57.6\unit{%
%TCIMACRO{\U{2103}}%
%BeginExpansion
{}^{\circ}{\rm C}%
%EndExpansion
}=330.\,\allowbreak 75\unit{K}$ \\ 
$T_{\min }=-89\unit{%
%TCIMACRO{\U{2103}}%
%BeginExpansion
{}^{\circ}{\rm C}%
%EndExpansion
}=184.\,\allowbreak 15\unit{K}$%
\end{tabular}%
\end{equation}

This $\Delta T$ range would be a challenge to cover with most materials.%
\footnote{%
Expecially metal structures and laser crystals. \TEXTsymbol{>}sigh%
\TEXTsymbol{<}} So if you are going to build a device or instrument that
will travel around the globe, you need to use special low $\alpha $
materials. Often making these materials is very expensive and harder to
obtain because the process used to make such materials is a tightly held
industrial secret. Alternately you can air condition or heat your system.
But you do have to do something to stay in the linear range.\footnote{%
I am told by the BYU-I facilities people that the BYUIC has one huge beam
across the entire auditorium that has to stay heated or it shrinks enough to
let the balcony fall off the walls!}

Looking at the table of average linear expansion coefficients we see
different $\alpha $ values for different materials. Sometimes $\alpha $
values are quite different. Suppose we take two materials with different $%
\alpha $'s and put them together. \FRAME{dtbpF}{3.4229in}{1.3189in}{0pt}{}{}{%
Figure}{\special{language "Scientific Word";type
"GRAPHIC";maintain-aspect-ratio TRUE;display "USEDEF";valid_file "T";width
3.4229in;height 1.3189in;depth 0pt;original-width 3.4601in;original-height
1.3154in;cropleft "0";croptop "1";cropright "1";cropbottom "0";tempfilename
'SVAQIE6G.wmf';tempfile-properties "XPR";}}The figure shows a bimetallic
strip. What would happen if we heated this bimetallic strip? The two metals
will expand at different rates, forcing the bonded pair of metals to bend.
And we could make a temperature measuring device that measures how much this
bimetallic strip bends.

\subsection{Water is weird (and the Lord knows what He is doing!)}

%TCIMACRO{%
%\TeXButton{Question123.9.2}{\marginpar {
%\hspace{-0.5in}
%\begin{minipage}[t]{1in}
%\small{Question 123.9.2}
%\end{minipage}
%}}}%
%BeginExpansion
\marginpar {
\hspace{-0.5in}
\begin{minipage}[t]{1in}
\small{Question 123.9.2}
\end{minipage}
}%
%EndExpansion
\FRAME{dtbpFU}{2.8746in}{2.3082in}{0pt}{\Qcb{{\protect\small Density of
water near }$0\unit{%
%TCIMACRO{\U{2103}}%
%BeginExpansion
{}^{\circ}{\rm C}%
%EndExpansion
}.${\protect\small \ (Public Domain image PD-1923, PD-Australia)}}}{}{Figure%
}{\special{language "Scientific Word";type "GRAPHIC";maintain-aspect-ratio
TRUE;display "USEDEF";valid_file "T";width 2.8746in;height 2.3082in;depth
0pt;original-width 2.8314in;original-height 2.2684in;cropleft "0";croptop
"1";cropright "1";cropbottom "0";tempfilename
'SVAQIE6H.wmf';tempfile-properties "XPR";}}

For most substances, the volume increases with increasing temperature.
Usually liquids increase in volume more than solids. But water is different.
Look at the graph of water density vs. Temperature.

Just as ice melts, it starts to \emph{increase} in density, but then it
becomes less dense with increased temperature after about $277.\,\allowbreak
15\unit{K}$ ($4\unit{%
%TCIMACRO{\U{2103}}%
%BeginExpansion
{}^{\circ}{\rm C}%
%EndExpansion
}$). We expect the behavior we see after $277.\,\allowbreak 15\unit{K}$ ($4%
\unit{%
%TCIMACRO{\U{2103}}%
%BeginExpansion
{}^{\circ}{\rm C}%
%EndExpansion
}$). The density is decreasing because we have the same number of water
molecules, but the temperature change has increased the volume. So the
density goes down.

But from $273.\,\allowbreak 15\unit{K}$ ($0\unit{%
%TCIMACRO{\U{2103}}%
%BeginExpansion
{}^{\circ}{\rm C}%
%EndExpansion
}$) to $277.\,\allowbreak 15\unit{K}$ ($4\unit{%
%TCIMACRO{\U{2103}}%
%BeginExpansion
{}^{\circ}{\rm C}%
%EndExpansion
}$) we see the volume decreases with $\Delta T.$ This is very odd. But very
cool! This is why a lake freezes from the top down. The colder water near $%
273.\,\allowbreak 15\unit{K}$ ($0\unit{%
%TCIMACRO{\U{2103}}%
%BeginExpansion
{}^{\circ}{\rm C}%
%EndExpansion
}$) is less dense, and so before it freezes it floats to the surface. This
is great, because if it became more dense just before it froze, the water
would sink to the bottom and freeze, leaving the liquid water exposed to the
cold atmosphere to freeze more water. The lake would fill up from the bottom
with solid ice, killing fish and water life.

\FRAME{dtbpFU}{3.3555in}{2.5918in}{0pt}{\Qcb{Lake Erie from Space. Notice
the ice covering part of the lake. (Public Domain Image courtesy NASA\
Visible Earth, http://visibleearth.nasa.gov/view\_rec.php?id=1286}}{}{Figure%
}{\special{language "Scientific Word";type "GRAPHIC";maintain-aspect-ratio
TRUE;display "USEDEF";valid_file "T";width 3.3555in;height 2.5918in;depth
0pt;original-width 5.2166in;original-height 4.0231in;cropleft "0";croptop
"1";cropright "1";cropbottom "0";tempfilename
'SVAQIE6I.wmf';tempfile-properties "XPR";}}Because the water becomes less
dense before freezing we don't freeze large lakes and oceans solid. And that
keeps life going on Earth!

We have made progress. But we still don't know what temperature is. We know
some things about what temperature does. And this can give us a clue.
Temperature is making things move by expansion. And we know that there must
be forces behind making things move. Forces do work. Work is energy. Could
it be that temperature has to do with exchanging energy?

%TCIMACRO{%
%\TeXButton{Question123.9.2}{\marginpar {
%\hspace{-0.5in}
%\begin{minipage}[t]{1in}
%\small{Question 123.9.2}
%\end{minipage}
%}}}%
%BeginExpansion
\marginpar {
\hspace{-0.5in}
\begin{minipage}[t]{1in}
\small{Question 123.9.2}
\end{minipage}
}%
%EndExpansion
%TCIMACRO{%
%\TeXButton{Question123.9.3}{\marginpar {
%\hspace{-0.5in}
%\begin{minipage}[t]{1in}
%\small{Question 123.9.3}
%\end{minipage}
%}}}%
%BeginExpansion
\marginpar {
\hspace{-0.5in}
\begin{minipage}[t]{1in}
\small{Question 123.9.3}
\end{minipage}
}%
%EndExpansion

\chapter{28 Thermal Properties of Matter 2.1.4, 2.1.5}

%TCIMACRO{\TeXButton{\chaptermark{Chapter 30}}{\chaptermark{Chapter 30}}}%
%BeginExpansion
\chaptermark{Chapter 30}%
%EndExpansion

In this lecture we want to make our study of temperature more mathematical.
And we want to consider how energy relates to temperature.

We learned how to use the concept of energy in Principles of Physics I. We
studied work and potential energy, and we studied kinetic energy. But these
were always the energy \emph{of} an object. We described these energies in
terms of the motion of the center of mass of an object using the particle
model. Now we know that there is energy \emph{in} an object. This is
different. Particles don't have internal parts, but real things do. Our
particle model was an approximation. So we need to go beyond the particle
model here. We will call the energy associated with the inside parts of an
object \emph{internal energy.} And this internal energy is associated with
thermal energy.

To study this, Let's refine our definition of internal energy by saying that
it is all the energy of a system that is associated with its microscopic
components -- atoms and molecules -- when viewed form a reference frame at
rest with respect to the center of mass of the system. Any motion of the
center of mass is our Principles of Physics I mechanical energy. Now we want
the motion of the parts of the object with respect to the center of mass.

Suppose we let all of the masses involved be tied together with spring-like
bonding forces. These spring-like forces would allow the masses to oscillate
a little. But if the oscillations were random, the actual center of mass
would not change. This is the kind of motion involved in internal energy.

By making all the little molecules move, we have changed the energy state of
the object, but we have not moved it, so the mechanical energy has not
changed. The energy we have described is the thermal internal energy (there
is still the nuclear internal energy, but that is a subject for the class
called \emph{Modern Physics} (PH279). Let's consider forms of internal
energy in more detail in this lecture.

%TCIMACRO{%
%\TeXButton{Fundamental Concepts}{\vspace{0.20in}\hspace{-0.37in}{\Large Fundamental Concepts\vspace{0.2in}}}}%
%BeginExpansion
\vspace{0.20in}\hspace{-0.37in}{\Large Fundamental Concepts\vspace{0.2in}}%
%EndExpansion

\begin{enumerate}
\item Heat capacity relates energy transfer by heat to temperature change in
solids and liquids

\item Specific heat is a heat capacity per unit mass

\item A phase change is changing from solid to liquid, or from liquid to gas
for a sample of a material

\item Calorimetry is the use of thermal conservation of energy to study a
sample of material
\end{enumerate}

\section{Heat}

Temperature, in our intuitive model, seems connected to heat. But what is
heat? If we can figure this out we might be able to give a better definition
of temperature. Let's give heat a symbol\footnote{%
I'm just so sorry. $Q$ is not the same thing as $\mathbb{Q}.$ I wish we
could use more different symbols.}, $Q.$ This is a different Q than our flow
rate $\mathbb{Q}$.

In this lecture we want to give a specific definition of heat:

\begin{center}
\rule{300pt}{1pt}

Heat is defined as the transfer of energy across the boundary of a system
due to a temperature difference between the system and its surroundings.

\rule{300pt}{1pt}
\end{center}

When we \textquotedblleft heat\textquotedblright\ a substance, we add energy
from it's surroundings. When the temperature changes, it is energy that is
making the temperature change. An object could gain or lose energy \emph{by
the process of heat}. This is like back in Principles of Physics I when we
could add energy to an object \emph{by doing work. }We will use the word
\textquotedblleft heat\textquotedblright\ for the transfer of energy no
matter which direction the energy is flowing. This is a little like using
the word \textquotedblleft acceleration\textquotedblright\ no matter whether
we are speeding up or slowing down.

In the dim past, scientists thought heat was a fluid (called caloric). This
model of heat is not correct. Science changes! We now define heat as a \emph{%
transfer} of energy. But because of this history, we have some left-over
names that are not very descriptive of the modern ideas they represent.
Examples are \textquotedblleft latent heat\textquotedblright\ and
\textquotedblleft heat capacity\textquotedblright\ (we will define these
soon).

The thing to remember is that a system \emph{has} a temperature but it can 
\emph{give or receive} energy by heat.

\subsection{Units}

There are several units for heat. The calorie (cal) is the amount of energy
transfer necessary to raise the temperature of $1\unit{g}$ of water from $%
14.5\unit{%
%TCIMACRO{\U{2103}}%
%BeginExpansion
{}^{\circ}{\rm C}%
%EndExpansion
}$ ($287.\,\allowbreak 65\unit{K}$) to $15.5\unit{%
%TCIMACRO{\U{2103}}%
%BeginExpansion
{}^{\circ}{\rm C}%
%EndExpansion
}$ ($288.\,\allowbreak 65\unit{K}$). This is somewhat arbitrary, but it
works for many things, especially biological systems. The British thermal
unit (Btu) is the unit used by refrigeration and heating contractors. One $%
\unit{Btu}$ is the amount of energy transfer necessary to raise the
temperature of $1\unit{lb}$ of water from $63\unit{%
%TCIMACRO{\U{2109}}%
%BeginExpansion
{}^{\circ}{\rm F}%
%EndExpansion
}$ $(290.\,\allowbreak 37\unit{K})$ to $64\unit{%
%TCIMACRO{\U{2109}}%
%BeginExpansion
{}^{\circ}{\rm F}%
%EndExpansion
}$ ($290.\,\allowbreak 93\unit{K}$). This is an enormous amount of energy
being transferred! If you wish to heat an entire building, this is a fine
unit, but for us it is a bit large. In the SI system we already have a unit
for energy, the Joule. We will try to stick to this nice, medium sized,
unit. But let's see where it came from.

\subsection{Mechanical equivalent of heat}

\FRAME{dhFU}{1.2436in}{1.4852in}{0pt}{\Qcb{{\protect\small James Prescott
Joule. (Image in the Public Domain)}}}{}{Figure}{\special{language
"Scientific Word";type "GRAPHIC";maintain-aspect-ratio TRUE;display
"USEDEF";valid_file "T";width 1.2436in;height 1.4852in;depth
0pt;original-width 4.6985in;original-height 5.6204in;cropleft "0";croptop
"1";cropright "1";cropbottom "0";tempfilename
'SVAQIE6K.wmf';tempfile-properties "XPR";}}

The name \textquotedblleft Joule\textquotedblright\ is a person's name,
Joule was an early researcher and is credited with showing that heat is a
transfer of energy. Joule was able to build a device to prove that
mechanical energy could be converted to heat energy. His device is shown in
the picture. \FRAME{dtbpFU}{3.4979in}{1.9971in}{0pt}{\Qcb{{\protect\small %
Joule Apparatus (Image in the Public Domain courtesy Gaius Cornelius)}}}{}{%
Figure}{\special{language "Scientific Word";type
"GRAPHIC";maintain-aspect-ratio TRUE;display "USEDEF";valid_file "T";width
3.4979in;height 1.9971in;depth 0pt;original-width 3.4523in;original-height
1.9579in;cropleft "0";croptop "1";cropright "1";cropbottom "0";tempfilename
'SY6A6R03.wmf';tempfile-properties "XPR";}}

The paddle device was placed in the drum that is sitting behind it. The
crank was operated by strings attached to it that were pulled by masses tied
to the other end of the string (see the right hand side of the figure). When
the masses were released, the paddle wheel stirred a liquid. The friction
from the paddle wheel stirring the liquid raised the temperature of the
liquid. The masses would lose energy equal to 
\begin{equation}
\Delta U=2mgh
\end{equation}%
assuming the masses were equal. By carefully measuring the temperature of
the water, Joule found that 
\begin{equation}
\Delta U=C\Delta T
\end{equation}%
with the proportionality constant approximately 
\begin{equation}
C=4.18\frac{\unit{J}}{\unit{g}}\unit{%
%TCIMACRO{\U{2103}}%
%BeginExpansion
{}^{\circ}{\rm C}%
%EndExpansion
}
\end{equation}%
which means that $4.18\unit{J}$ of mechanical energy raises the temperature
of $1\unit{g}$ of water by $1\unit{%
%TCIMACRO{\U{2103}}%
%BeginExpansion
{}^{\circ}{\rm C}%
%EndExpansion
}.$ (Here $C$ is a generic symbol for a constant). We have more precise
measurements now that set this value at $4.186\frac{\unit{J}}{\unit{g}}\unit{%
%TCIMACRO{\U{2103}}%
%BeginExpansion
{}^{\circ}{\rm C}%
%EndExpansion
}.$ But you can see that Joule did very good for working with what amounts
to a fancy butter churn in the 1800's!

Notice that we are considering raising the temperature of $1\unit{g}$ of
water by $1\unit{%
%TCIMACRO{\U{2103}}%
%BeginExpansion
{}^{\circ}{\rm C}%
%EndExpansion
}.$ This is the basis of the unit called the calorie. Using the same $4.5%
\unit{%
%TCIMACRO{\U{2103}}%
%BeginExpansion
{}^{\circ}{\rm C}%
%EndExpansion
}$ ($287.\,\allowbreak 65\unit{K}$) to $15.5\unit{%
%TCIMACRO{\U{2103}}%
%BeginExpansion
{}^{\circ}{\rm C}%
%EndExpansion
}$ ($288.\,\allowbreak 65\unit{K}$) change in temperature, we realize%
\begin{equation}
1\unit{cal}=4.186\unit{J}
\end{equation}%
which gives us a convenient way to convert from calories to Jules. But more
importantly, it shows that heat is really tied to energy, not to some
substance. Work can raise a temperature in a system just like
\textquotedblleft heating\textquotedblright\ the system with a fire or stove.

Because of this experiment by Joule, this idea that work can be turned into
thermal energy is known as the \emph{mechanical equivalent of heat}.

\section{State Variables}

As you know, we are at the beginning of a new topic, and we will find that
physics did not develop as a cohesive whole. The new topic is called \emph{%
Thermodynamics} and it developed independently of the theories of
oscillation and waves and optics. Thus the language we use will be
different, and sometimes old words will need new meanings.

We learned before about solids, liquids, and gases. We learned that the
degree of bondedness made the difference between these states. The states of
matter are often called \emph{phases} in thermodynamics. This use of the
word \textquotedblleft phase\textquotedblright\ has nothing to do with the
phase of an oscillator or wave. It is just the condition of being solid,
liquid or gas.

\begin{center}
\rule{300pt}{1pt}

To change from solid to liquid, or from liquid to gas, is called a \emph{%
phase change}.

\rule{300pt}{1pt}
\end{center}

To describe the state of a sample of matter, we use a set of variables like
mass, volume, pressure, density, and energy. We have used these throughout
Principles of Physics I. Thermodynamics has it's own set of descriptive
variables to add to these: number of moles, number density, and Temperature.
These variables, along with our original set, are called \emph{state
variables} and together, they describe the conditions of a particular piece
of matter. We have hinted that one of these state variables will be the
amount of internal energy in our sample of material, let's call that $\Delta
E_{int}.$ Note here the \textquotedblleft $\Delta $\textquotedblright\ means
\textquotedblleft a small amount of\textquotedblright\ and not a
\textquotedblleft difference between two amounts of.\textquotedblright\ So $%
\Delta E_{int}$ is an amount of internal energy. We know from our study of
thermal expansion that the work, $w,$ might be such a state variable. And we
also know now that energy can be transferred by heat, $Q,$ so maybe heat is
a state variable.

It would be good to be able to relate these new variables. So let's try to
develop a relationship between $\Delta E_{int},$ $W,$ and $Q$ for our
samples. Suppose you are heating something on your stove, say, a sample of
steak. \FRAME{dtbpF}{1.6134in}{1.9451in}{0pt}{}{}{Figure}{\special{language
"Scientific Word";type "GRAPHIC";maintain-aspect-ratio TRUE;display
"USEDEF";valid_file "T";width 1.6134in;height 1.9451in;depth
0pt;original-width 1.6152in;original-height 1.954in;cropleft "0";croptop
"1";cropright "1";cropbottom "0";tempfilename
'SY6CU905.wmf';tempfile-properties "XPR";}}You would expect that as energy
is transferred from the hot stove to the cold beef the temperature of the
steak would increase. It makes sense that the more energy you transfer, the
more the temperature will rise. Gasses are more complicated, because their
volume will change, but the volume of solids does not change much with
temperature (we know that it changes a little) so we don't have as many
complications as we would with a gas.\FRAME{dtbpF}{2.522in}{2.0084in}{0in}{}{%
}{Figure}{\special{language "Scientific Word";type
"GRAPHIC";maintain-aspect-ratio TRUE;display "USEDEF";valid_file "T";width
2.522in;height 2.0084in;depth 0in;original-width 2.5029in;original-height
1.9875in;cropleft "0";croptop "1";cropright "1";cropbottom "0";tempfilename
'TH3D0V00.wmf';tempfile-properties "XPR";}}

We could guess that the temperature increase might be linear in the amount
of energy transfer we give to the solid.%
\begin{equation*}
\Delta T\propto \Delta E_{int}
\end{equation*}%
where our only form of increasing the energy here is $Q$, since the change
in volume of the stake sample is small enough that any work, $w,$ is
negligible. Since $P$ and $%
%TCIMACRO{\TeXButton{V}{\ooalign{\hfil$V$\hfil\cr\kern0.1em--\hfil\cr}}}%
%BeginExpansion
\ooalign{\hfil$V$\hfil\cr\kern0.1em--\hfil\cr}%
%EndExpansion
$ are not changing much. So to a good approximation 
\begin{equation*}
\Delta T\propto Q
\end{equation*}

Experimentalists did this experiment and found that for many temperature
regions this is true. Of course, we could over heat our sample\FRAME{dtbpF}{%
2.7612in}{2.0019in}{0pt}{}{}{Figure}{\special{language "Scientific
Word";type "GRAPHIC";maintain-aspect-ratio TRUE;display "USEDEF";valid_file
"T";width 2.7612in;height 2.0019in;depth 0pt;original-width
2.7851in;original-height 2.0117in;cropleft "0";croptop "1";cropright
"1";cropbottom "0";tempfilename 'SY6CU904.wmf';tempfile-properties "XPR";}}%
changing it's form (from stake to charcoal, and smoke), but then the sample
is no longer the same substance. In this case $\Delta T$ is no longer
proportional to $Q.$ But as long as we don't change the composition of the
substance (we stay in the linear regime), we can say that $\Delta T$ is
proportional to $Q.$

Physicists don't like proportionality signs, so they invented a constant
with the right units that would make this an equality%
\begin{equation*}
\Delta T=\frac{1}{K}Q
\end{equation*}%
But, they found that the internal structure of the solid mattered. There was
a different constant, $K,$ for every material.

They called the constant, $K,$ the \emph{heat capacity} of the material. The
name is one of our throwbacks from the past. But for us it means the
constant of proportionality between $Q$ and $\Delta T$ that holds all the
specific material properties of that substance that change the slope of the $%
\Delta T$ vs. $Q$ graph.

Of course, the more material we have in our sample, the longer it takes to
heat it up. If we have a roast, it takes longer to cook than a small steak.
So it is customary to define a heat capacity per unit mass, this is called
the \emph{specific heat} of the material.%
\begin{equation*}
\Delta T=\frac{1}{mc}Q
\end{equation*}%
and we often rearrange this equation to give 
\begin{equation*}
Q=mc\Delta T
\end{equation*}%
As an example, it takes $4190\unit{J}$ to raise $1\unit{kg}$ of water by $1%
\unit{K}.$ So 
\begin{equation*}
4190\unit{J}=\left( 1\unit{kg}\right) \left( c\right) \left( 1\unit{K}\right)
\end{equation*}%
solving for $c$ yields 
\begin{equation*}
c_{water}=4190\frac{\unit{J}}{\unit{kg}\unit{K}}
\end{equation*}%
Which is familiar to us from Joule's work.

Here are a few more substances and their specific heat values.%
\begin{equation*}
\begin{tabular}{|l|l|l|}
\hline
\textbf{Substance} & $c\left( \frac{\unit{J}}{\unit{kg}\unit{K}}\right) $ & $%
c\left( \frac{\unit{J}}{\unit{mol}\unit{K}}\right) $ \\ \hline
Aluminum & $900$ & $24.3$ \\ \hline
Copper & $385$ & $24.4$ \\ \hline
Iron & $449$ & $25.1$ \\ \hline
Gold & $129$ & $25.4$ \\ \hline
Lead & $128$ & $26.5$ \\ \hline
Ice & $2090$ & $37.6$ \\ \hline
Mercury & $140$ & $28.1$ \\ \hline
Water & $4190$ & $75.4$ \\ \hline
\end{tabular}%
\end{equation*}%
Notice that things that heat up or cool down quickly have small specific
heats. Things that heat up slowly or cool slowly have large specific heats.
Water in any of it's phases has large values of specific heat. This is why
living next to a lake or the ocean keeps your temperatures moderate. It
takes quite a lot of energy transfer to heat up or cool off the water, so
the temperature of the water does not change much.

This is why Buffalo NY temperatures lower when the lake freezes. This is
also the source of lake or ocean breezes. In the morning, the land changes
temperature more quickly than the water. The warmer air over the land
becomes less dense and rises, leaving a lower pressure area. This causes a
cool breeze to form from the sea to the land. \FRAME{dtbpF}{2.2001in}{%
1.6034in}{0in}{}{}{Figure}{\special{language "Scientific Word";type
"GRAPHIC";maintain-aspect-ratio TRUE;display "USEDEF";valid_file "T";width
2.2001in;height 1.6034in;depth 0in;original-width 2.1594in;original-height
1.567in;cropleft "0";croptop "1";cropright "1";cropbottom "0";tempfilename
'SYDLMU05.wmf';tempfile-properties "XPR";}}

Conversely, in a large city, the concrete and blacktop heat up much quicker,
so we experience increases in daytime temperature as a city grows.

A specific heat is a per-unit-mass property, but we could just as well
define specific heat as a per-unit-mole quantity. Then the specific heat
would be the amount of energy it takes to increase the temperature of a mole
of material. This incarnation of the idea of specific heat is called a \emph{%
molar specific heat,} and we will give it the symbol $C.$%
\begin{equation*}
Q=nC\Delta T
\end{equation*}

It is worth noting that the molar heat capacities given in the last table
are remarkably alike. Except for water, which is a compound, they are all
about $25\unit{J}/\left( \unit{mol}\unit{K}\right) .$ This gives us a clue
towards understanding what is happening microscopically, but we will have to
wait to investigate this clue until a further lecture.

WARNING: Specific heat values are really not constant. They do vary with
temperature a little. We know from our study of thermal expansion that the
volume of a sample of a material will change a small amount as it is heated
or cooled. After all, this is one way we found to measure temperature! There
is a small amount of work done in changing the volume of the material. Not
all of the energy will go into changing the temperature. So our heat
capacity will change a little as the volume changes. But if the change in
temperature is not too big, we can treat heat capacities for solids and
liquids as constants. Oh, but for gasses...we know the volume might change
quite a lot as the temperature changes for a gas, so we will find that this
analysis needs a few changes for gasses. We will study this is a later
lecture.

We should make a sign convention about our energy transfer by heat. Let's
say that if energy is added by heat, $Q$ is positive. If the energy leaves
by heat, we will say that $Q$ is negative.

\FRAME{dtbpF}{1.8983in}{1.9182in}{0in}{}{}{Figure}{\special{language
"Scientific Word";type "GRAPHIC";maintain-aspect-ratio TRUE;display
"USEDEF";valid_file "T";width 1.8983in;height 1.9182in;depth
0in;original-width 1.8602in;original-height 1.8801in;cropleft "0";croptop
"1";cropright "1";cropbottom "0";tempfilename
'SYQNX300.wmf';tempfile-properties "XPR";}}

\section{Phase Changes and Phase Diagrams}

Let's consider heating ice. You take a chunk of ice from outside on a cold
February day in Rexburg, Idaho. The ice has a temperature of about $-25\unit{%
%TCIMACRO{\U{2103}}%
%BeginExpansion
{}^{\circ}{\rm C}%
%EndExpansion
}.$\footnote{%
OK, that might be an exaggeration, it only occasionally gets that cold in
Rexburg.} You place it in a pan on your stove and heat it up. Then, being a
physics student, you plot the temperature as a function of time. Here is
what you get. \FRAME{dhF}{3.2292in}{2.0237in}{0pt}{}{}{Figure}{\special%
{language "Scientific Word";type "GRAPHIC";maintain-aspect-ratio
TRUE;display "USEDEF";valid_file "T";width 3.2292in;height 2.0237in;depth
0pt;original-width 3.1842in;original-height 1.9847in;cropleft "0";croptop
"1";cropright "1";cropbottom "0";tempfilename
'SUS73W9S.wmf';tempfile-properties "XPR";}}Notice the places where you are
still transferring energy by heat, but the temperature does not change.
These are phase changes. A phase change is like melting and boiling. In our
chart, the first is melting, the second is boiling.

One of the things you will notice is that the melted ice does not boil at $%
100\unit{%
%TCIMACRO{\U{2103}}%
%BeginExpansion
{}^{\circ}{\rm C}%
%EndExpansion
}$ here in Rexburg. This is because we are up so high that the air pressure
is less than $1\unit{atm}=1.013\times 10^{5}\unit{Pa}.$ In the next graph, I
have plotted the melting, and boiling point of a substance as a function of
pressure. The melting point and boiling point are not single points, but
lines on this graph. This is because melting and boiling points depend on
pressure. There is also a line for the sublimation point. Along this line, a
solid will change directly into a gas. Dry ice does this at normal room
temperature and pressure. \FRAME{dhF}{2.3359in}{1.8931in}{0pt}{}{}{Figure}{%
\special{language "Scientific Word";type "GRAPHIC";maintain-aspect-ratio
TRUE;display "USEDEF";valid_file "T";width 2.3359in;height 1.8931in;depth
0pt;original-width 3.6703in;original-height 2.9706in;cropleft "0";croptop
"1";cropright "1";cropbottom "0";tempfilename
'SUS73W9T.wmf';tempfile-properties "XPR";}}If we observe the boiling line,
we see that as pressure goes down, the boiling temperature goes down. Look
at $T_{1}$ and $T_{2}$ in the next figure, where $T_{2}<T_{1}.$ Note that $%
P_{2}<P_{1}.$\FRAME{dhF}{1.9683in}{1.8222in}{0pt}{}{}{Figure}{\special%
{language "Scientific Word";type "GRAPHIC";maintain-aspect-ratio
TRUE;display "USEDEF";valid_file "T";width 1.9683in;height 1.8222in;depth
0pt;original-width 10.8577in;original-height 10.0474in;cropleft "0";croptop
"1";cropright "1";cropbottom "0";tempfilename
'SUS73W9U.wmf';tempfile-properties "XPR";}}This agrees with our Rexburg
boiling point experience.\footnote{%
So if we have a water boil alert here, you have to boil your water longer in
Rexburg than you would in NY City in order to make it safe to drink.} We can
see from the graph that melting or boiling depends on pressure, and so does
sublimation point.

The triple point is a special point where all three states are possible at
the same pressure and temperature. Ice, steam, and water will all exist in
equilibrium, for example, at the triple point of water.

The critical point is where the clear distinction between liquid and gas
stops. beyond this boiling the material is in a fluid state, but not clearly
liquid or gas. Along each of the lines in our figure, we can say that the
material is in phase equilibrium, meaning that both phases from either side
of the line are possible and equally probable. Only at the triple point are
all three phases equally probable.

Our diagram is rather schematic, and this is on purpose. An actual phase
diagram for water is given in the next figure.

\FRAME{dtbpFU}{4.203in}{1.9545in}{0pt}{\Qcb{More Complete Phase Diagram for
Water (Public Domain Image courtesy Karlhahn)}}{}{Figure}{\special{language
"Scientific Word";type "GRAPHIC";maintain-aspect-ratio TRUE;display
"USEDEF";valid_file "T";width 4.203in;height 1.9545in;depth
0pt;original-width 3.3589in;original-height 1.5471in;cropleft "0";croptop
"1";cropright "1";cropbottom "0";tempfilename
'SUS73W9V.wmf';tempfile-properties "XPR";}}You can see there are more phases
for water than our simple model of solid, liquid, and gas give.

But why do we care about phase changes? Because it takes energy to make a
phase change happen. Phase changes direct energy into breaking molecular
bonds instead of increasing temperature. So in our graph of temperature vs.
heating time the flat parts are where bonds are being broken and a phase
change is happening.

\section{Phase changes and heat of transformation}

Recall the graph of temperature vs. time for heating a chunk of Rexburg ice. 
\FRAME{dhF}{3.2595in}{2.0202in}{0pt}{}{}{Figure}{\special{language
"Scientific Word";type "GRAPHIC";maintain-aspect-ratio TRUE;display
"USEDEF";valid_file "T";width 3.2595in;height 2.0202in;depth
0pt;original-width 4.8101in;original-height 2.9706in;cropleft "0";croptop
"1";cropright "1";cropbottom "0";tempfilename
'SW9P0O0V.wmf';tempfile-properties "XPR";}}Here is the graph again, but this
time the horizontal axis is labeled \textquotedblleft Energy
added\textquotedblright\ because we realize that we are adding internal
energy as we wait and watch the ice. \FRAME{dtbpF}{3.5267in}{2.6498in}{0pt}{%
}{}{Figure}{\special{language "Scientific Word";type
"GRAPHIC";maintain-aspect-ratio TRUE;display "USEDEF";valid_file "T";width
3.5267in;height 2.6498in;depth 0pt;original-width 5.0004in;original-height
3.7498in;cropleft "0";croptop "1";cropright "1";cropbottom "0";tempfilename
'SW9P0O0W.wmf';tempfile-properties "XP";}}We see regions of the graph where
the ice changes temperature, and we can understand that during these times
the temperature gain will be proportional to the amount of energy
transferred by heat.%
\begin{equation*}
\text{Slope}=\frac{\Delta T}{Q}=\frac{1}{mc}
\end{equation*}

But what about the horizontal parts of the graph? We identified these before
as phase changes. The slope is zero during the melting and boiling times.

We have an amount of energy transferred by heat. You might guess that how
much energy it takes for the phase change to completely change from one
state to the other depends on the strength of the bonds of the material. The
added energy is going into breaking those bonds. So the amount of energy it
takes to, say, melt different substances will be different. From experience
we know that it takes more energy to heat large amounts of material than
small amounts of material. We can express this mathematically as 
\begin{equation*}
Q=\pm Lm
\end{equation*}%
where $L$ is the symbol we give to the constant that expresses the how
easily the material melts or freezes. We need at least two of these, one for
melting/freezing and one for boiling/condensing. Then 
\begin{eqnarray*}
Q &=&\pm L_{f}m\qquad \text{melt/freeze} \\
Q &=&\pm L_{v}m\qquad \text{boil/condense}
\end{eqnarray*}%
where we call these two constants 
\begin{equation*}
\begin{tabular}{ll}
$L_{f}$ & heat of fusion \\ 
$L_{v}$ & heat of vaporization%
\end{tabular}%
\end{equation*}%
Generically we call each of these a \emph{heat of transformation} but
generally they are still referred to by their old title of \textquotedblleft 
\emph{Latent heat}.\textquotedblright

\begin{equation*}
\begin{tabular}{|l|l|l|l|l|}
\hline
\textbf{Substance} & $T_{m}\left( \unit{%
%TCIMACRO{\U{2103}}%
%BeginExpansion
{}^{\circ}{\rm C}%
%EndExpansion
}\right) $ & $L_{f}\left( \frac{\unit{J}}{\unit{kg}}\right) $ & $T_{m}\left( 
\unit{%
%TCIMACRO{\U{2103}}%
%BeginExpansion
{}^{\circ}{\rm C}%
%EndExpansion
}\right) $ & $L_{v}\left( \frac{\unit{J}}{\unit{kg}}\right) $ \\ \hline
Nitrogen $\left( N_{2}\right) $ & $-210$ & $0.26\times 10^{5}$ & $-196$ & $%
1.99\times 10^{5}$ \\ \hline
Ethyl alcohol & $-114$ & $1.09\times 10^{5}$ & $78$ & $8.79\times 10^{5}$ \\ 
\hline
Mercury & $-39$ & $0.11\times 10^{5}$ & $357$ & $2.96\times 10^{5}$ \\ \hline
Water & $0$ & $3.33\times 10^{5}$ & $100$ & $22.6\times 10^{5}$ \\ \hline
Lead & $328$ & $0.25\times 10^{5}$ & $1750$ & $8.58\times 10^{5}$ \\ \hline
\end{tabular}%
\end{equation*}

Notice that our equations each have a \textquotedblleft $\pm $%
\textquotedblright\ sign in them. We must supply the sign by context. We
need a sign convention like back when we were studying optics!

And we already know what this sign convention is. We choose that energy that
leaves is negative and energy that is gained by the system is positive. Also
notice that so long as we warm the ice somewhat slowly, that the ice
completely melts before the temperature of the water starts to rise.

Let's do a problem. Suppose we want to increase the temperature of $1\unit{g}%
=\allowbreak 0.001\,\unit{kg}$ of ice at $-30.0\unit{%
%TCIMACRO{\U{2103}}%
%BeginExpansion
{}^{\circ}{\rm C}%
%EndExpansion
}$ and change it to steam at $120.0\unit{%
%TCIMACRO{\U{2103}}%
%BeginExpansion
{}^{\circ}{\rm C}%
%EndExpansion
}.$ How much energy transfer by heat is required to do this?

A graph of temperature vs. heat energy is shown. Let's take it one piece at
a time.\FRAME{dhF}{3.1176in}{2.1093in}{0pt}{}{}{Figure}{\special{language
"Scientific Word";type "GRAPHIC";maintain-aspect-ratio TRUE;display
"USEDEF";valid_file "T";width 3.1176in;height 2.1093in;depth
0pt;original-width 5.4916in;original-height 3.7057in;cropleft "0";croptop
"1";cropright "1";cropbottom "0";tempfilename
'SW9P0O0X.wmf';tempfile-properties "XPR";}}

\textbf{Part A:} Warming up the ice to $0\unit{%
%TCIMACRO{\U{2103}}%
%BeginExpansion
{}^{\circ}{\rm C}%
%EndExpansion
}.$

This part acts as we learned in the last section. We have 
\begin{eqnarray*}
Q &=&m_{i}c_{i}\Delta T \\
&=&\left( 0.001\,\unit{kg}\right) \left( 2090\frac{\unit{J}}{\unit{kg}\unit{%
%TCIMACRO{\U{2103}}%
%BeginExpansion
{}^{\circ}{\rm C}%
%EndExpansion
}}\right) \left( 0\unit{%
%TCIMACRO{\U{2103}}%
%BeginExpansion
{}^{\circ}{\rm C}%
%EndExpansion
}-\left( -30\unit{%
%TCIMACRO{\U{2103}}%
%BeginExpansion
{}^{\circ}{\rm C}%
%EndExpansion
}\right) \right) \\
&=&62.\,\allowbreak 7\unit{J}
\end{eqnarray*}%
We can get away with using $\unit{%
%TCIMACRO{\U{2103}}%
%BeginExpansion
{}^{\circ}{\rm C}%
%EndExpansion
}$ because we only need $\Delta T$ and the divisions of the Celsius and
Kelvin scales are the same size.

\textbf{Part B:} Melting ice

Now the ice changes to water. We have a heat of transformation. To find out
what happens we use the equation%
\begin{eqnarray*}
Q &=&m_{i}L_{f} \\
&=&\left( 0.001\,\unit{kg}\right) \left( 3.33\times 10^{5}\frac{\unit{J}}{%
\unit{kg}}\right) \\
&=&\allowbreak 333.0\unit{J}
\end{eqnarray*}

The total energy so far is $62.\,\allowbreak 7\unit{J}+330.0\unit{J}%
=\allowbreak 392.\,\allowbreak 7\unit{J}.$ Note that during Part B the
temperature didn't go up!

\textbf{Part C:} Warming the melted ice (water)

Again we have the normal case where we can use the specific heat of water
(now that the ice has melted)

\begin{eqnarray*}
Q &=&m_{w}c_{w}\Delta T \\
&=&\left( 1\unit{g}\right) \left( 4190\frac{\unit{J}}{\unit{kg}\unit{%
%TCIMACRO{\U{2103}}%
%BeginExpansion
{}^{\circ}{\rm C}%
%EndExpansion
}}\right) \left( 100\unit{%
%TCIMACRO{\U{2103}}%
%BeginExpansion
{}^{\circ}{\rm C}%
%EndExpansion
}-0\unit{%
%TCIMACRO{\U{2103}}%
%BeginExpansion
{}^{\circ}{\rm C}%
%EndExpansion
}\right) \\
&=&\allowbreak 419.0\unit{J}
\end{eqnarray*}

The total energy so far is $62.\,\allowbreak 7\unit{J}+333.0\unit{J}%
+419.\,\allowbreak 0\unit{J}=\allowbreak 814.\,\allowbreak 7\unit{J}$

\textbf{Part D:} Boiling the melted ice (water)

Now we convert the water to steam. 
\begin{eqnarray*}
Q &=&m_{w}L_{v} \\
&=&\left( 0.001\,\unit{kg}\right) \left( 2.26\times 10^{6}\frac{\unit{J}}{%
\unit{kg}}\right) \\
&=&\allowbreak 2260.0\unit{J}
\end{eqnarray*}

The total energy so far is $62.\,\allowbreak 7\unit{J}+333.0\unit{J}%
+419.\,\allowbreak 0\unit{J}+2260.0\unit{J}=\allowbreak 3074.\,\allowbreak 7%
\unit{J}$

\textbf{Part E:} Warming the steam

Now we have steam, and can use the specific heat of steam (we don't usually
do this! we usually use the ideal gas law when it becomes a gas!)%
\begin{eqnarray*}
Q &=&m_{s}c_{s}\Delta T \\
&=&\left( 0.001\,\unit{kg}\right) \left( 2010\frac{\unit{J}}{\unit{kg}\unit{%
%TCIMACRO{\U{2103}}%
%BeginExpansion
{}^{\circ}{\rm C}%
%EndExpansion
}}\right) \left( 120\unit{%
%TCIMACRO{\U{2103}}%
%BeginExpansion
{}^{\circ}{\rm C}%
%EndExpansion
}-100\unit{%
%TCIMACRO{\U{2103}}%
%BeginExpansion
{}^{\circ}{\rm C}%
%EndExpansion
}\right) \\
&=&\allowbreak 40.\,\allowbreak 2\unit{J}
\end{eqnarray*}%
The total energy transfer required is $62.\,\allowbreak 7\unit{J}+333.0\unit{%
J}+419\unit{J}+2260.0\unit{J}+40.\,\allowbreak 2\unit{J}=\allowbreak
3114.\,\allowbreak 9\unit{J}$

\section{Calorimetry}

We can now calculate results for many common thermal experiences. We have
all experienced eating something with a large temperature. The solution is
to quickly drink something cold. This lowers the temperature. We would like
to be able to quantify this. But our mouth is not a great device to use to
do quantitative analysis on temperature changes. Let's consider a more
suitable apparatus.

Suppose we have an insulated container with a thermometer. We can call this
apparatus a \emph{calorimeter.}

\FRAME{dtbpF}{3.6659in}{1.5567in}{0pt}{}{}{Figure}{\special{language
"Scientific Word";type "GRAPHIC";maintain-aspect-ratio TRUE;display
"USEDEF";valid_file "T";width 3.6659in;height 1.5567in;depth
0pt;original-width 9.2033in;original-height 3.8925in;cropleft "0";croptop
"1";cropright "1";cropbottom "0";tempfilename
'SW9P0O0Y.wmf';tempfile-properties "XPR";}}Inside of this container let's
place a known mass of water and measure the temperature, $T_{w}$. Then we
can introduce our sample of hot material. Suppose we know that the hot
material has a temperature $T_{x}.$

Let's assume that the container walls are perfectly insulating meaning that
no energy can escape through the walls by heat. Then knowing the sample
temperature $T_{x}$ and the water temperature $T_{w}$ with $T_{w}<T_{x}$ we
can find $c$ or $K.$(or even $C$) for our material. We recognize for this
ideal system no energy can leave, because of the insulation. So what energy
the cold water gains through heat must have come from the hot sample by
heat. We can write this as 
\begin{equation}
Q_{\text{cold}}=-Q_{\text{hot}}
\end{equation}%
This tells us that energy is being conserved. Then for the water%
\begin{equation*}
Q_{w}=m_{w}c_{w}\left( T_{f}-T_{wi}\right)
\end{equation*}%
where $T_{f}$ is the final temperature for the interior of the container and 
$T_{wi}$ is the initial temperature of the water. For the sample 
\begin{equation*}
Q_{x}=m_{x}c_{x}\left( T_{f}-T_{xi}\right)
\end{equation*}%
where $T_{f}$ is the same final temperature for the interior of the
container and $T_{xi}$ is the sample initial temperature. Using energy
conservation%
\begin{eqnarray*}
-Q_{x} &=&Q_{w} \\
-m_{x}c_{x}\left( T_{f}-T_{xi}\right) &=&m_{w}c_{w}\left( T_{f}-T_{wi}\right)
\end{eqnarray*}%
and we can solve for $c_{x}$%
\begin{equation}
c_{x}=-\frac{m_{w}c_{w}\left( T_{f}-T_{wi}\right) }{m_{x}\left(
T_{f}-T_{xi}\right) }
\end{equation}

Note the sign convention! If energy leaves the object, the value of $Q$ is
negative. Our denominator is, indeed, negative because for the sample $T_{f}$
is less than $T_{xi}.$ We could use our value of $c_{x}$ to identify our
material by looking it up in a table of specific heat values.

Let's try a second problem. Suppose we have $2\unit{kg}$ of ice at $0\unit{%
%TCIMACRO{\U{2103}}%
%BeginExpansion
{}^{\circ}{\rm C}%
%EndExpansion
}$ and $1\unit{kg}$ of water at $50\unit{%
%TCIMACRO{\U{2103}}%
%BeginExpansion
{}^{\circ}{\rm C}%
%EndExpansion
}.$ What is the final temperature? This time we could have a phase change,
we will need to be careful. We can look up the heat of transformation and
the specific heat if ice and water:%
\begin{eqnarray*}
L_{f} &=&3.33\times 10^{5}\frac{\unit{J}}{\unit{kg}} \\
c_{w} &=&4190\frac{\unit{J}}{\unit{kg}\unit{K}}
\end{eqnarray*}%
and we know%
\begin{eqnarray*}
m_{i} &=&2\unit{kg} \\
m_{w} &=&1\unit{kg} \\
T_{ii} &=&0\unit{%
%TCIMACRO{\U{2103}}%
%BeginExpansion
{}^{\circ}{\rm C}%
%EndExpansion
} \\
T_{wi} &=&50\unit{%
%TCIMACRO{\U{2103}}%
%BeginExpansion
{}^{\circ}{\rm C}%
%EndExpansion
}
\end{eqnarray*}

We still don't expect to lose any energy, so

\begin{equation*}
Q_{\text{cold}}=-Q_{\text{hot}}
\end{equation*}%
We identify the hot thing as the water this time and the cold thing as the
ice. The water only changes temperature, but the ice experiences a phase
change. We need to account for all of the energy transferred so for this ice 
\begin{equation*}
Q_{\text{cold}}=Q_{ice}=Q_{\text{melt}}+Q_{\text{heat up}}
\end{equation*}%
and 
\begin{equation*}
Q_{\text{hot}}=Q_{\text{water}}
\end{equation*}%
Then we can put these into our conservation of energy equation 
\begin{eqnarray*}
Q_{\text{ice}} &=&-Q_{\text{water}} \\
Q_{\text{melt}}+Q_{\text{heat up}} &=&-Q_{\text{water}} \\
L_{f}m_{i}+m_{i}c_{w}\left( T_{f}-T_{ii}\right) &=&-m_{w}c_{w}\left(
T_{f}-T_{wi}\right)
\end{eqnarray*}%
The final temperature will be the same for both. To find that final
temperature we have a little algebra to do. Let's distribute on both sides
of our equation. Then 
\begin{equation*}
L_{f}m_{i}+m_{i}c_{w}\left( T_{f}-T_{ii}\right) =-m_{w}c_{w}\left(
T_{f}-T_{wi}\right)
\end{equation*}%
becomes%
\begin{equation*}
L_{f}m_{i}+m_{i}c_{w}T_{f}-m_{i}c_{w}T_{ii}=-m_{w}c_{w}T_{f}+m_{w}c_{w}T_{wi}
\end{equation*}%
Now let's rearrange so all the terms with $T_{f}$ are on the same side.%
\begin{equation*}
L_{f}m_{i}-m_{i}c_{w}T_{ii}-m_{w}c_{w}T_{wi}=-m_{w}c_{w}T_{f}-m_{i}c_{w}T_{f}
\end{equation*}

\begin{equation*}
L_{f}m_{i}-m_{i}c_{w}T_{ii}-m_{w}c_{w}T_{iw}=-\left(
m_{w}c_{w}+m_{i}c_{w}\right) T_{f}
\end{equation*}%
We can solve for $T_{f}$%
\begin{equation*}
\frac{L_{f}m_{i}-m_{i}c_{w}T_{ii}-m_{w}c_{w}T_{wi}}{-\left(
m_{w}c_{w}+m_{i}c_{w}\right) }=T_{f}
\end{equation*}%
or

\begin{equation*}
T_{f}=\frac{m_{i}c_{w}T_{ii}+m_{w}c_{w}T_{wi}-L_{f}m_{i}}{\left(
m_{w}c_{w}+m_{i}c_{w}\right) }
\end{equation*}

Really we should do this in kelvin.

\begin{eqnarray*}
T_{f} &=&\frac{\left( 2\unit{kg}\right) \left( 4190\frac{\unit{J}}{\unit{kg}%
\unit{K}}\right) \left( 273\unit{K}\right) +\left( 1\unit{kg}\right) \left(
4190\frac{\unit{J}}{\unit{kg}\unit{K}}\right) \left( 273+50\right) \unit{K}%
-\left( 3.33\times 10^{5}\frac{\unit{J}}{\unit{kg}}\right) \left( 2\unit{kg}%
\right) }{\left( \left( 1\unit{kg}\right) \left( 4190\frac{\unit{J}}{\unit{kg%
}\unit{K}}\right) +\left( 2\unit{kg}\right) \left( 4190\frac{\unit{J}}{\unit{%
kg}\unit{K}}\right) \right) } \\
&=&236.\,\allowbreak 68\unit{K} \\
&=&-36.\,\allowbreak 47\unit{%
%TCIMACRO{\U{2103}}%
%BeginExpansion
{}^{\circ}{\rm C}%
%EndExpansion
}
\end{eqnarray*}

Now let's ask, is this reasonable?

The answer is no! We can't end up with a temperature less than the two
temperatures we started with. Let's see why this went wrong.

We should have checked in advance to see if all the ice melts. It would take 
\begin{equation*}
Q=m_{i}L_{f}=\left( 2\unit{kg}\right) \left( 3.33\times 10^{5}\frac{\unit{J}%
}{\unit{kg}}\right) =\allowbreak 6.\,\allowbreak 66\times 10^{5}\unit{J}
\end{equation*}%
to melt all the ice. Since the water in the calorimeter is losing energy, we
expect that it will have a lower temperature. But it is unlikely that it
will all freeze, so the lowest temperature for the water would be $0\unit{%
%TCIMACRO{\U{2103}}%
%BeginExpansion
{}^{\circ}{\rm C}%
%EndExpansion
}.$ That means the available energy transfer from the water is just 
\begin{eqnarray*}
Q &=&m_{w}C_{w}\left( T_{0}-T_{wi}\right) \\
&=&\left( 1\unit{kg}\right) \left( 4190\frac{\unit{J}}{\unit{kg}\unit{K}}%
\right) \left( 273\unit{K}-323.\,\allowbreak 15\unit{K}\right) \\
&=&\allowbreak -2.095\times 10^{5}\unit{J}
\end{eqnarray*}%
That is, the warmer water can provide $2.095\times 10^{5}\unit{J}.$ But
since $2.095\times 10^{5}\unit{J}<\allowbreak 6.\,\allowbreak 66\times 10^{5}%
\unit{J}$ the warm water does not have enough energy available to melt all
the ice. Once the water is all at $0\unit{%
%TCIMACRO{\U{2103}}%
%BeginExpansion
{}^{\circ}{\rm C}%
%EndExpansion
},$ the ice and water are in thermal equilibrium. No energy will be
transferred. So our final temperature of the mixture is $273\unit{K}$ or $0%
\unit{%
%TCIMACRO{\U{2103}}%
%BeginExpansion
{}^{\circ}{\rm C}%
%EndExpansion
}.$

Of course if our system of ice and water isn't isolated, eventually all the
ice will melt. But that would be due to a transfer of energy by heat from
the outside environment.

This problem is a bit of a trick, but such a situation can really happen. So
it is really important that we stop to consider whether our numerical
answers are reasonable once we have done a calorimetry problem.

We have filled in the missing piece of information for solids and liquids.
We can say that so long as we don't change the structure of the material, 
\begin{equation*}
Q=mc\Delta T
\end{equation*}%
or 
\begin{equation*}
Q=nC\Delta T
\end{equation*}%
so we can find $Q$, $w$, and $\Delta E_{int}$ for solids and liquids But we
said gasses will be more difficult. We will take up the relationship between
energy transfer by heat and the change in temperature for gasses in our next
lecture.

\chapter{29 Heat Transfer and Gas Laws 2.1.6, 2.2.1}

Now that we understand heat as a transfer of energy, we should look at how
we can move energy around. We will study this in this lecture. We will also
develop a relationship between pressure, volume, and temperature for a gas.
You probably know this relationship from your high school science classes.
But we will look at how we got it and when it is valid.

%TCIMACRO{\TeXButton{\chaptermark{Chapter 31}}{\chaptermark{Chapter 31}}}%
%BeginExpansion
\chaptermark{Chapter 31}%
%EndExpansion

%TCIMACRO{%
%\TeXButton{Fundamental Concepts}{\vspace{0.10in}\hspace{-0.37in}{\Large Fundamental Concepts\vspace{0.2in}}}}%
%BeginExpansion
\vspace{0.10in}\hspace{-0.37in}{\Large Fundamental Concepts\vspace{0.2in}}%
%EndExpansion

\begin{itemize}
\item Conduction

\item Convection

\item Radiation and Stefan's law

\item The Ideal Gas Law

\item The Van der Waals Gas Law
\end{itemize}

\section{Heat Transfer Mechanisms}

In Principles of Physics I we found we could change the energy state of a
system by work, $w.$ Joule showed us that this works for internal energy as
well. And we already know about $Q,$ but how does energy transfer by heat
happen?

Let's start by thinking about how to change $E_{int}.$ The change in
internal energy is the sum of all the ways we can make $E_{int}$ become
different. This is just how energy works. Think back to Principles of
Physics I. Our mechanical energy was the sum of all the kinds of energy we
could have. So 
\begin{equation*}
E_{mech}=K+U_{g}+U_{s}+\cdots
\end{equation*}%
We should expect a sum like this for $E_{int}.$ So far we have%
\begin{equation*}
\Delta E_{int}=Q+w
\end{equation*}%
We know a lot about work, but let's become more familiar with $Q.$ For
mechanical energy we had a bunch of different potential energy types, like
gravitational potential energy $\left( U_{b}\right) $ and spring potential
energy $\left( U_{s}\right) .$ For internal energy there are a bunch of
different types of energy transfer by heat. We could say each of these is an
generic energy transfer mechanism, $\Delta E_{\text{Transfer mechanism}}$
which would include our work, $w.$ Then we can sum up all the energy
transfer to find the change in internal energy. 
\begin{equation}
\Delta E_{int}=\sum_{i}\Delta E_{\text{Transfer mechanism}}{}_{i}
\end{equation}
And we are concentrating on types of $Q$ right now, so let's list a few
kinds of $Q.$

\subsection{Thermal Conduction}

This is what we have been calling heat up to this point. We define it a
little more carefully as the exchange of kinetic energy between microscopic
particles (molecules and atoms) through direct contact. Really we mean 
\textbf{collisions} between the atoms and molecules of the two objects. 
\FRAME{dtbpF}{1.8498in}{2.4543in}{0pt}{}{}{Figure}{\special{language
"Scientific Word";type "GRAPHIC";maintain-aspect-ratio TRUE;display
"USEDEF";valid_file "T";width 1.8498in;height 2.4543in;depth
0pt;original-width 1.8118in;original-height 2.4128in;cropleft "0";croptop
"1";cropright "1";cropbottom "0";tempfilename
'SY6E3T07.wmf';tempfile-properties "XPR";}}

When we add energy by heat to part of a solid, the atoms are locally
displaced from their equilibrium locations. They are still bound to their
neighbors. But they stretch these bonds and knock into their neighbors. This
causes the neighboring atoms to be displaced and increases the neighbor's
kinetic energy. This process continues through the whole solid.

The rate of heat transfer will depend on the properties of the atoms that
make up the substance. Some substances more tightly bind their atoms, some
allow for kinetic energy more easily. \FRAME{dtbpF}{1.8248in}{2.2485in}{0pt}{%
}{}{Figure}{\special{language "Scientific Word";type
"GRAPHIC";maintain-aspect-ratio TRUE;display "USEDEF";valid_file "T";width
1.8248in;height 2.2485in;depth 0pt;original-width 1.7876in;original-height
2.2087in;cropleft "0";croptop "1";cropright "1";cropbottom "0";tempfilename
'TC5WQY00.wmf';tempfile-properties "XPR";}}We can envision a slab shaped
piece of solid as in the figure. It has thickness $\Delta L$ and
cross-sectional area $A.$ Suppose one face of the slab is at temperature $%
T_{c}$ and the other at temperature $T_{h}.$ Experience tells us that energy 
$Q$ will transfer in a time interval $\Delta t$ from the hotter face to the
colder one. This is just a burner on your stove heating a pan, but it takes
time for the pan to heat up before you can cook your egg for breakfast. The
rate at which we use energy is what we called power, $\mathcal{P}$. Then the
power is the amount of energy transferred, $Q$ in an amount of time. 
\begin{equation*}
\mathcal{P=}\frac{E}{\Delta t}=\frac{Q}{\Delta t}
\end{equation*}%
The energy transferred, $Q,$ is proportional to $A$ as well as and $\Delta
T. $ Think, if you have a bigger pan, it takes longer to heat up.%
\begin{equation*}
\mathcal{P}=\frac{Q}{\Delta t}\varpropto A\frac{\Delta T}{\Delta L}
\end{equation*}%
We can take $\Delta L$ very small%
\begin{equation}
\mathcal{P}=k_{therm}A\left\vert \frac{dT}{dL}\right\vert
\end{equation}%
where we now have another constant $k$! This $k_{therm}$ is the \emph{%
thermal conductivity} of the material and we define $\left\vert \frac{dT}{dL}%
\right\vert $ as the \emph{temperature gradient.}

A gradient is a rate of change, so this gradient is the rate of change of
the temperature with respect to position. An example of this kind of
gradient might be a sleeping bag\footnote{%
Or, if you don't go camping, the pile of blankets on you use to keep you
warm as you sleep.}. You transfer energy from your body to the sleeping bag.
The temperature of the bag material next you goes up. But think of burrowing
through the sleeping bag material from the inside out. The temperature would
go down as you move away from the energy source (you). We would find that
the outside of the bag has a much lower temperature than the inside. That is
a temperature gradient.\footnote{%
If you are camping in Idaho in the winter, the outside of your bag might be
at $-23.\,\allowbreak 333\unit{%
%TCIMACRO{\U{2103}}%
%BeginExpansion
{}^{\circ}{\rm C}%
%EndExpansion
}$ $\left( -10\unit{%
%TCIMACRO{\U{2109}}%
%BeginExpansion
{}^{\circ}{\rm F}%
%EndExpansion
}\right) ,$ but if you have a really good sleeping bag the inside might be a
more comfortable $21\unit{%
%TCIMACRO{\U{2103}}%
%BeginExpansion
{}^{\circ}{\rm C}%
%EndExpansion
}$ $\left( -70\unit{%
%TCIMACRO{\U{2109}}%
%BeginExpansion
{}^{\circ}{\rm F}%
%EndExpansion
}\right) .$}

We can say that our slab is a one dimensional problem, because the energy
transfer is only in one direction. If we assume that $k_{therm}$ is not
temperature dependent and the slab is uniform, then%
\begin{equation*}
\left\vert \frac{dT}{dL}\right\vert =\frac{T_{h}-T_{c}}{\Delta L}
\end{equation*}%
and our energy transfer rate is 
\begin{eqnarray}
\mathcal{P}_{rod} &\mathcal{=}&\frac{Q}{\Delta t}=k_{therm}A\frac{T_{h}-T_{c}%
}{\Delta L} \\
&=&\frac{A\left( T_{h}-T_{c}\right) }{\frac{\Delta L}{k_{therm}}}  \notag
\end{eqnarray}

If we have a complicated multipart slab with $n$ different substances we
would just find the energy transferred from the first slab part and use that
as the input to the next slab part. To a good approximation we could write
the result as%
\begin{equation}
\mathcal{P}\approx \frac{A\left( T_{h}-T_{c}\right) }{\dsum\limits_{i}^{n}%
\frac{\Delta L_{i}}{k_{therm_{i}}}}
\end{equation}%
The amount of energy transferred by thermal conduction is 
\begin{equation*}
Q=\mathcal{P}\Delta t\approx \frac{A\left( T_{h}-T_{c}\right) \Delta t}{%
\dsum\limits_{i}^{n}\frac{\Delta L_{i}}{k_{therm_{i}}}}
\end{equation*}

All this wasn't too bad mathematically. But notice that our slab looks a lot
like a wall, and the building industry does some further simplifications to
our thermal conduction equation. We define a term%
\begin{equation}
R=\frac{\Delta L}{k_{therm}}
\end{equation}%
as the \textquotedblleft R-value\textquotedblright\ of a material. Then a
wall or multipart solid would have%
\begin{equation}
\mathcal{P=}\frac{Q}{\Delta t}\approx \frac{A\left( T_{h}-T_{c}\right) }{%
\dsum\limits_{i}^{n}R_{i}}
\end{equation}%
so that 
\begin{equation*}
Q_{conduction}\approx \frac{A\left( T_{h}-T_{c}\right) \Delta t}{%
\dsum\limits_{i}^{n}R_{i}}
\end{equation*}%
These R-values are used in home insulation and heating design and are
usually in English units

\subsection{Convection}

This is a fundamentally new type of energy transport in our study of
thermodynamics, although we are all familiar with it from normal life. In
the last section we used the example of an electric burner that touched our
pot of vegetables. But there are other kinds of stoves. \FRAME{dtbpF}{%
3.6677in}{1.836in}{0pt}{}{}{Figure}{\special{language "Scientific Word";type
"GRAPHIC";maintain-aspect-ratio TRUE;display "USEDEF";valid_file "T";width
3.6677in;height 1.836in;depth 0pt;original-width 5.2088in;original-height
2.5927in;cropleft "0";croptop "1";cropright "1";cropbottom "0";tempfilename
'SY6E9E08.wmf';tempfile-properties "XPR";}}Convection is the transfer of
energy by the movement of the atoms or molecules, themselves. Because the
group of moving atoms or molecules has internal energy, moving these
molecules also moves the energy. So convection is an energy transfer. This
is like circulation of air in your home, or the transfer of energy from a
gas flame to your cooking pot. In the case of a gas stove, the burning gas
flame has high internal energy and you move the hot gas to your pan to bring
the energy to your pan.

Your student apartment might have base-board heading. But your parent's home
would likely use forced convection where a fan blows the air. That air
carries energy from the furnace to your room. \FRAME{dtbpF}{2.4365in}{%
2.5173in}{0pt}{}{}{Figure}{\special{language "Scientific Word";type
"GRAPHIC";maintain-aspect-ratio TRUE;display "USEDEF";valid_file "T";width
2.4365in;height 2.5173in;depth 0pt;original-width 2.4543in;original-height
2.5368in;cropleft "0";croptop "1";cropright "1";cropbottom "0";tempfilename
'SVAQIE6O.wmf';tempfile-properties "XPR";}}For several decades houses have
been made with the idea that you can heat or cool a house with passive
convection. Here is an example.\FRAME{dtbpF}{4.9744in}{1.5532in}{0pt}{}{}{%
Figure}{\special{language "Scientific Word";type
"GRAPHIC";maintain-aspect-ratio TRUE;display "USEDEF";valid_file "T";width
4.9744in;height 1.5532in;depth 0pt;original-width 4.9199in;original-height
1.5177in;cropleft "0";croptop "1";cropright "1";cropbottom "0";tempfilename
'SYH9JO01.wmf';tempfile-properties "XPR";}}This house tries to control the
thermal conduction mechanisms to keep the interior of the house at nearly
the same temperature year round.

\subsection{Radiation}

Let's go back to thinking about cooking. Suppose you are cooking over an old
wood stove. You open the stove door and see this\FRAME{dtbpF}{2.7095in}{%
2.7095in}{0pt}{}{}{Figure}{\special{language "Scientific Word";type
"GRAPHIC";maintain-aspect-ratio TRUE;display "USEDEF";valid_file "T";width
2.7095in;height 2.7095in;depth 0pt;original-width 2.6671in;original-height
2.6671in;cropleft "0";croptop "1";cropright "1";cropbottom "0";tempfilename
'SY6EJR09.wmf';tempfile-properties "XPR";}}And feel a blast of energy
transfer by heat. But you don't feel a blast of wind. This isn't convection.
And you are not touching the stove. So it's not conduction.

Think about the light coming from the stove, that orange glow. We know light
is a wave and waves transfer energy. Light is a form of energy transfer! We
call this energy removal or delivery by light \textquotedblleft
radiation.\textquotedblright\ This is different than the use of the word
\textquotedblleft radiation\textquotedblright\ in nuclear physics.

All objects radiate energy according to Stefan's law%
\begin{equation}
\mathcal{P=}\frac{Q}{\Delta t}=\sigma AeT^{4}
\end{equation}%
The parts of this law are Stephen's constant $\sigma =5.6696\times 10^{-8}%
\frac{\unit{W}}{\unit{m}^{2}\unit{K}^{4}}$, the factor, $A,$ is the surface
area, $e$ is the emissivity of the material. This $e$ tells us how like a
perfect absorber (called a black body) our surface is, and $T$ is the
temperature. Note there are other forms for these quantities! But this is
the form we will use.

The emissivity is a value from $0$ to $1.$ If $e=0$ then the material is a
perfect reflector. This would be a perfect mirror. If the emissivity is $1,$
then it is a perfect absorber. This is more like blacktop in sunlight. The
blacktop warms up, even in the winter, because it absorbs the sunlight. Of
course, nothing has an emissivity of exactly $0$ or exactly $1.$ Real
objects are somewhere between these extremes.

It may seem strange that \textquotedblleft black\textquotedblright\ objects
glow, but they do. Think again of that blacktop. The ice on the blacktop
melts because the blacktop not only absorbs the radiation from the sun, it
also emits radiation. In our optics section, materials absorbed light and
re-emitted it at the same wavelength. But this is not always true. The
blacktop is absorbing a lot of visible light from the sun but is radiating
more infrared light that is then absorbed by the snow.

An extreme example is the Sun, itself. Light that hits the Sun from outside
will scatter off the gasses in the Sun. There is a lot of gas in the Sun, so
the light just keeps bouncing in the vast sea of gas. It takes a very long
time for the light to reemerge, if it ever does. Most likely it will be
absorbed by the gasses. So the Sun is considered a \textquotedblleft black
body!\textquotedblright\ But it glows in the visible part of the spectrum,
so it does not look black.

We can explain many common experiences using the idea of emissivity. Take
this leaf for example. Why is it sinking into the snow?

\FRAME{dhF}{1.6527in}{1.9908in}{0pt}{}{}{Figure}{\special{language
"Scientific Word";type "GRAPHIC";maintain-aspect-ratio TRUE;display
"USEDEF";valid_file "T";width 1.6527in;height 1.9908in;depth
0pt;original-width 5.3938in;original-height 6.5086in;cropleft "0";croptop
"1";cropright "1";cropbottom "0";tempfilename
'SVAQIE6Q.wmf';tempfile-properties "XPR";}}The leaf has a different
emissivity than the snow, so it absorbs more radiation from the sun, the
temperature of the leaf rises and it melts the snow around it.

With this understanding, we can see that all objects absorb radiation as
well as radiate. If an object is at temperature $T_{obj}$ and it's
surroundings are at temperature $T_{env}.$ then the net rate of energy
gained or lost by radiation is 
\begin{equation}
\mathcal{P}=\frac{Q}{\Delta t}=\sigma Ae\left( T_{obj}^{4}-T_{env}^{4}\right)
\end{equation}%
and the energy transfer by radiation would be 
\begin{equation*}
Q_{radiation}=\frac{Q}{\Delta t}=\sigma Ae\left(
T_{obj}^{4}-T_{env}^{4}\right) \Delta t
\end{equation*}

Back when we were thinking of building thermometers we said that the color
of an object could vary with temperature. What we meant is that the
temperature determines the color of the light radiating from the object. If
you get to go on in physics, you will study what wavelengths of light come
from objects of different temperature (in Principles of Physics IV, called
Modern Physics). But here is a graph of what wavelengths are produced by an
ideal object who's temperature is $0\unit{%
%TCIMACRO{\U{2103}}%
%BeginExpansion
{}^{\circ}{\rm C}%
%EndExpansion
}.$ \FRAME{dtbpFU}{2.6126in}{1.9882in}{0pt}{\Qcb{{\protect\small Thermal
Emission Spectrum (Planck Curve)}}}{}{Figure}{\special{language "Scientific
Word";type "GRAPHIC";maintain-aspect-ratio TRUE;display "USEDEF";valid_file
"T";width 2.6126in;height 1.9882in;depth 0pt;original-width
3.0199in;original-height 2.2917in;cropleft "0";croptop "1";cropright
"1";cropbottom "0";tempfilename 'SYHCZM08.wmf';tempfile-properties "XPR";}}

You can see that more than one wavelength is produced but there is a peak
wavelength. And if you are paying attention you can see that the peak is not
in the visible spectrum. It is more like a radio wave. The object that made
this glow is fairly cold and this is why we don't see the glow from a cold
object. But if we increase the temperature, say, to $1200\unit{%
%TCIMACRO{\U{2103}}%
%BeginExpansion
{}^{\circ}{\rm C}%
%EndExpansion
}$ we would get an object that glows red. And at $6000\unit{%
%TCIMACRO{\U{2103}}%
%BeginExpansion
{}^{\circ}{\rm C}%
%EndExpansion
}$ we would have our object glow blue. Here are the curves for these two
cases, along with our $0\unit{%
%TCIMACRO{\U{2103}}%
%BeginExpansion
{}^{\circ}{\rm C}%
%EndExpansion
}$ curve for reference. \FRAME{dtbpFU}{2.5131in}{1.913in}{0pt}{\Qcb{%
{\protect\small Normalized emission curves for objects with }$T=0\unit{%
%TCIMACRO{\U{2103}}%
%BeginExpansion
{}^{\circ}{\rm C}%
%EndExpansion
}${\protect\small \ (black curve), }$T=1200\unit{%
%TCIMACRO{\U{2103}}%
%BeginExpansion
{}^{\circ}{\rm C}%
%EndExpansion
}${\protect\small \ (red curve), and }$T=6000\unit{%
%TCIMACRO{\U{2103}}%
%BeginExpansion
{}^{\circ}{\rm C}%
%EndExpansion
}${\protect\small \ (blue curve).}}}{}{Figure}{\special{language "Scientific
Word";type "GRAPHIC";maintain-aspect-ratio TRUE;display "USEDEF";valid_file
"T";width 2.5131in;height 1.913in;depth 0pt;original-width
3.1289in;original-height 2.3748in;cropleft "0";croptop "1";cropright
"1";cropbottom "0";tempfilename 'SYHDDJ0B.wmf';tempfile-properties "XPR";}}

The curves moved down into the visible part of the spectrum. Note that in
this last graph the three temperatures wouldn't give the same amount of
light. Stefan's law tells us this. So the blue $6000\unit{%
%TCIMACRO{\U{2103}}%
%BeginExpansion
{}^{\circ}{\rm C}%
%EndExpansion
}$ curve would be much taller. In order to see the three curves on the same
graph I have normalized the curves by dividing the values in a curve by
their maximum. This makes the curves easier to compare, but it makes them
not tell us about the relative strength of the different glows. If we expand
the visible light part of the graph we see this \FRAME{dtbpF}{2.6066in}{%
1.9847in}{0pt}{}{}{Figure}{\special{language "Scientific Word";type
"GRAPHIC";maintain-aspect-ratio TRUE;display "USEDEF";valid_file "T";width
2.6066in;height 1.9847in;depth 0pt;original-width 3.0415in;original-height
2.3082in;cropleft "0";croptop "1";cropright "1";cropbottom "0";tempfilename
'SYHDFG0C.wmf';tempfile-properties "XPR";}}

The blue curve is for the $6000\unit{%
%TCIMACRO{\U{2103}}%
%BeginExpansion
{}^{\circ}{\rm C}%
%EndExpansion
}$ object and it does peak in the blue part of the spectrum. The $1200\unit{%
%TCIMACRO{\U{2103}}%
%BeginExpansion
{}^{\circ}{\rm C}%
%EndExpansion
}$ curve is the red curve and it peaks off the chart, but the light it makes
is all on the red side of the spectrum, so we would see red. And the black
curve for $0\unit{%
%TCIMACRO{\U{2103}}%
%BeginExpansion
{}^{\circ}{\rm C}%
%EndExpansion
}$ doesn't rise in the visible spectrum as we already know.

But this brings us a question. Can we determine temperature by looking at
the wavelength of the radiation glow coming from an object. Of course, the
answer is yes. We can find the peak of the radiation, and use our curve
backwards to find the temperature. Thermal cameras do exactly that. \FRAME{%
dtbpFU}{2.6171in}{2.6171in}{0pt}{\Qcb{{\protect\small Thermal Image of
Physics Students. The colors you see are not the real colors of the
radiative glow. The glow colors are all in the infrared part of the
spectrum. The camera manufacturer assigned a visible color to all of the
infrared colors so we could see them on the display.}}}{}{Figure}{\special%
{language "Scientific Word";type "GRAPHIC";maintain-aspect-ratio
TRUE;display "USEDEF";valid_file "T";width 2.6171in;height 2.6171in;depth
0pt;original-width 4.3024in;original-height 4.3024in;cropleft "0";croptop
"1";cropright "1";cropbottom "0";tempfilename
'SYHDKS0D.wmf';tempfile-properties "XPR";}}Stefan's law adds up the energy
from all the wavelengths that we see in our graphs to get the total power
emitted.

\subsection{Multiple transfer mechanisms}

As we said before, we might have many different energy transfer mechanisms.
We could have all of these (and more).%
\begin{equation*}
\Delta E_{int}=Q_{conduction}+Q_{convection}+Q_{radiation}+w
\end{equation*}

Let's think about an example that uses radiation, convection, and insulation
to keep something cold. It is called a dewar flask.

\FRAME{fhFU}{3.3702in}{1.8369in}{0pt}{\Qcb{{\protect\small You might see our
liquid nitrogen dewar in demonstrations. The flask is pictured on the left,
and a schematic diagram of how a dewar works is given on the right.}}}{}{%
Figure}{\special{language "Scientific Word";type
"GRAPHIC";maintain-aspect-ratio TRUE;display "USEDEF";valid_file "T";width
3.3702in;height 1.8369in;depth 0pt;original-width 3.4051in;original-height
1.8432in;cropleft "0";croptop "1";cropright "1";cropbottom "0";tempfilename
'SVAQIE6R.wmf';tempfile-properties "XPR";}}This is the same idea as a
thermos bottle, only better made (and more expensive). The flask has doubled
walls with the space between evacuated to reduce convection and eliminate
conduction. The walls are silvered to reduce heat transfer by radiation. The
cap has a long piece of insulation attached to it to prevent conduction
through the flask opening. New designs use multi-layer films which make
their surfaces super Reflective.

\section{Gas Laws}

We are trying to find the change in internal energy as $Q$ and $w$ are
applied to our systems. To go further, we need to know how the pressure,
temperature and volume of a gas are related. You probably studied the ideal
gas law in your pre-college science classes, but before we use it in our
class let's try a though experiment.

Let's consider the following data\FRAME{dhF}{2.6472in}{1.5947in}{0pt}{}{}{%
Figure}{\special{language "Scientific Word";type
"GRAPHIC";maintain-aspect-ratio TRUE;display "USEDEF";valid_file "T";width
2.6472in;height 1.5947in;depth 0pt;original-width 6.7014in;original-height
4.0248in;cropleft "0";croptop "1";cropright "1";cropbottom "0";tempfilename
'SUS73WA2.wmf';tempfile-properties "XPR";}}The black straight line is a
curve fit to the data. What can we say about the curve fit?

It might look useless. But look at the region from about $-6<x<5.$ In this
region, the curve fit is not too bad. Of course, the equation we got for the
curve fit is not the right equation, and we know this from the regions $x<-6$
and $x>4.$ If the equation we used for the fit comes from our model for how
this physical system works, we know that something is wrong with the model.
The data tells us that do not have an exact theory that explains the
behavior of this system well. But could the curve fit be useful? If all you
needed was an estimate of the value of the actual function at $x=-1,$ then
the curve fit might be enough to get your work done! This may not be
terribly satisfying, but sometimes it is a useful way to work.

Another way to think about this is that if most of the time, under usual
conditions, the phenomena we experience fall within the range $-6<x<5$ and
the curve fit equation is much simpler, it might be good enough and
convenient to use the curve fit. For example, we know that Newton's laws are
not exact. We have to use General Relativity to be accurate. But in
calculating our average speed on a trip to Idaho Falls, Newton's laws are
plenty good enough, and \emph{much} more convenient than their General
Relativistic forms.

We will use the \emph{Ideal Gas Law} to study the expansion of gasses as
temperature changes. The Ideal Gas Law is an approximation very like our
curve fit we have been considering. It is an equation that works pretty well
under normal conditions, and it is much easier to use than the exact law. It
allows us to gain great insight, without the mathematical difficulty of the
exact relationship. Given all this, we should realize that there is no such
thing as an ideal gas.

In this approximation, the atoms are so weakly bound that there is no
equilibrium position between atoms. Thus there is no definite volume defined
for an ideal gas. This is generally mostly true around standard pressures
and temperatures for real gasses. We will let $%
%TCIMACRO{\TeXButton{V}{\ooalign{\hfil$V$\hfil\cr\kern0.1em--\hfil\cr}}}%
%BeginExpansion
\ooalign{\hfil$V$\hfil\cr\kern0.1em--\hfil\cr}%
%EndExpansion
$ be a variable for gases. Let's look at where the ideal gas law came from.
It is much like our hypothetical experiment. The experimentalists gave us
some clues on how to relate $%
%TCIMACRO{\TeXButton{V}{\ooalign{\hfil$V$\hfil\cr\kern0.1em--\hfil\cr}}}%
%BeginExpansion
\ooalign{\hfil$V$\hfil\cr\kern0.1em--\hfil\cr}%
%EndExpansion
,$ $P,$ and $T.$

\subsection{What lead to the Ideal Gas Law}

A researcher named Boyle found that if a system is kept at a constant
temperature then 
\begin{equation}
P\propto \frac{1}{%
%TCIMACRO{\TeXButton{V}{\ooalign{\hfil$V$\hfil\cr\kern0.1em--\hfil\cr}}}%
%BeginExpansion
\ooalign{\hfil$V$\hfil\cr\kern0.1em--\hfil\cr}%
%EndExpansion
}\text{ for constant }T
\end{equation}%
\FRAME{dtbpF}{3.2785in}{2.1906in}{0pt}{}{}{Figure}{\special{language
"Scientific Word";type "GRAPHIC";maintain-aspect-ratio TRUE;display
"USEDEF";valid_file "T";width 3.2785in;height 2.1906in;depth
0pt;original-width 3.2335in;original-height 2.1517in;cropleft "0";croptop
"1";cropright "1";cropbottom "0";tempfilename
'TC5YGG09.wmf';tempfile-properties "XPR";}}That is, if you don't let the
temperature of the gas change, then a change in pressures is inversely
related to a change in the volume of the gas. The more pressure, the smaller
the volume. We could envision taking a gas in a cylinder, and putting a
piston in the cylinder that can be used to compress our gas. \FRAME{dtbpF}{%
0.9074in}{1.931in}{0pt}{}{}{Figure}{\special{language "Scientific Word";type
"GRAPHIC";maintain-aspect-ratio TRUE;display "USEDEF";valid_file "T";width
0.9074in;height 1.931in;depth 0pt;original-width 1.3367in;original-height
2.8783in;cropleft "0";croptop "1";cropright "1";cropbottom "0";tempfilename
'SYBUKW02.wmf';tempfile-properties "XPR";}}Boyle tells us that as we exert a
pressure on the gas using the piston the volume of the gas changes. That
seems reasonable.

Two researchers, Charles and Gay-Lussac, found that if a system is kept at
constant pressure then%
\begin{equation}
%TCIMACRO{\TeXButton{V}{\ooalign{\hfil$V$\hfil\cr\kern0.1em--\hfil\cr}}}%
%BeginExpansion
\ooalign{\hfil$V$\hfil\cr\kern0.1em--\hfil\cr}%
%EndExpansion
\propto T\text{ for constant }P
\end{equation}%
\FRAME{dtbpF}{3.365in}{2.2485in}{0in}{}{}{Figure}{\special{language
"Scientific Word";type "GRAPHIC";maintain-aspect-ratio TRUE;display
"USEDEF";valid_file "T";width 3.365in;height 2.2485in;depth
0in;original-width 3.3183in;original-height 2.2087in;cropleft "0";croptop
"1";cropright "1";cropbottom "0";tempfilename
'TC5YFI08.wmf';tempfile-properties "XPR";}}The warmer the gas, the more it
will expand. Note that in the graph temperature is not one of the axes! That
is a little weird. But you can see the volume of the gas getting larger with
the pressure staying the same. Something must be making this happen, and
that something is the increase in temperature.

Combining these two findings gives an equation that describes ideal gasses.

\begin{equation}
P%
%TCIMACRO{\TeXButton{V}{\ooalign{\hfil$V$\hfil\cr\kern0.1em--\hfil\cr}}}%
%BeginExpansion
\ooalign{\hfil$V$\hfil\cr\kern0.1em--\hfil\cr}%
%EndExpansion
=nRT
\end{equation}%
where $n$ is the number of moles of gas, and $R$ is called the \emph{%
universal gas constant}. This is the \emph{ideal gas law}. The law came from
researchers that could not access extreme pressures or temperatures, and so
for normal conditions, it works just fine. But it will not work at all for
extremely low or high pressures and temperatures.

The variables $n,$ $%
%TCIMACRO{\TeXButton{V}{\ooalign{\hfil$V$\hfil\cr\kern0.1em--\hfil\cr}}}%
%BeginExpansion
\ooalign{\hfil$V$\hfil\cr\kern0.1em--\hfil\cr}%
%EndExpansion
,$ $P,$ and $T,$ are state variables for ideal gasses. Depending on the
units for $%
%TCIMACRO{\TeXButton{V}{\ooalign{\hfil$V$\hfil\cr\kern0.1em--\hfil\cr}}}%
%BeginExpansion
\ooalign{\hfil$V$\hfil\cr\kern0.1em--\hfil\cr}%
%EndExpansion
,$ $P,$ and $T,$ the value for $R$ may be different.%
\begin{eqnarray}
R &=&8.314\frac{\unit{J}}{\unit{mol}\unit{K}} \\
&=&0.08214\frac{\unit{atm}\unit{l}}{\unit{mol}\unit{K}}  \notag
\end{eqnarray}

\subsection{Alternate form of the Ideal Gas Law}

We know we can describe the number of moles as 
\begin{equation}
n=\frac{N}{N_{A}}
\end{equation}%
where $N$ is the number of molecules of gas that we have. Then we can write%
\begin{equation}
P%
%TCIMACRO{\TeXButton{V}{\ooalign{\hfil$V$\hfil\cr\kern0.1em--\hfil\cr}}}%
%BeginExpansion
\ooalign{\hfil$V$\hfil\cr\kern0.1em--\hfil\cr}%
%EndExpansion
=\frac{N}{N_{A}}RT
\end{equation}%
if we define a new constant 
\begin{equation}
k_{B}=\frac{R}{N_{A}}
\end{equation}%
Then%
\begin{equation}
P%
%TCIMACRO{\TeXButton{V}{\ooalign{\hfil$V$\hfil\cr\kern0.1em--\hfil\cr}}}%
%BeginExpansion
\ooalign{\hfil$V$\hfil\cr\kern0.1em--\hfil\cr}%
%EndExpansion
=Nk_{B}RT
\end{equation}%
The constant $k_{B}$ is called \emph{Boltzmann's constant}

\begin{equation}
k_{B}=1.38\times 10^{-23}\frac{\unit{J}}{\unit{K}}
\end{equation}

%TCIMACRO{%
%\TeXButton{Fire Syringe Demo}{\marginpar {
%\hspace{-0.5in}
%\begin{minipage}[t]{1in}
%\small{Fire Syringe Demo}
%\end{minipage}
%}}}%
%BeginExpansion
\marginpar {
\hspace{-0.5in}
\begin{minipage}[t]{1in}
\small{Fire Syringe Demo}
\end{minipage}
}%
%EndExpansion

We now have two forms of the ideal gas law!

\section{PV diagrams}

You may have noticed that in the diagrams in the previous sections we put
pressure and volume on the axes. \FRAME{dtbpF}{3.6158in}{2.4154in}{0pt}{}{}{%
Figure}{\special{language "Scientific Word";type
"GRAPHIC";maintain-aspect-ratio TRUE;display "USEDEF";valid_file "T";width
3.6158in;height 2.4154in;depth 0pt;original-width 3.5682in;original-height
2.3748in;cropleft "0";croptop "1";cropright "1";cropbottom "0";tempfilename
'TC5YOZ0D.wmf';tempfile-properties "XPR";}}This was awkward because our gas
laws have three state variables, $P,$ $%
%TCIMACRO{\TeXButton{V}{\ooalign{\hfil$V$\hfil\cr\kern0.1em--\hfil\cr}}}%
%BeginExpansion
\ooalign{\hfil$V$\hfil\cr\kern0.1em--\hfil\cr}%
%EndExpansion
,$ and $T$ 
\begin{equation*}
P=\frac{nRT}{%
%TCIMACRO{\TeXButton{V}{\ooalign{\hfil$V$\hfil\cr\kern0.1em--\hfil\cr}}}%
%BeginExpansion
\ooalign{\hfil$V$\hfil\cr\kern0.1em--\hfil\cr}%
%EndExpansion
}
\end{equation*}%
and there wasn't another axis for $T.$ We commonly do this, so we need to
know what happens if we change $T.$ In the previous picture we graphed $P$
vs. $V$ holding $T$ constant. In the next figure, you can see the red dashed
line is the same \FRAME{dtbpF}{3.4082in}{2.2753in}{0pt}{}{}{Figure}{\special%
{language "Scientific Word";type "GRAPHIC";maintain-aspect-ratio
TRUE;display "USEDEF";valid_file "T";width 3.4082in;height 2.2753in;depth
0pt;original-width 3.9833in;original-height 2.6498in;cropleft "0";croptop
"1";cropright "1";cropbottom "0";tempfilename
'TC5YQ80F.wmf';tempfile-properties "XPR";}}but now there is a green solid
line. This is a line for a higher temperature. We expect a $1/%
%TCIMACRO{\TeXButton{V}{\ooalign{\hfil$V$\hfil\cr\kern0.1em--\hfil\cr}}}%
%BeginExpansion
\ooalign{\hfil$V$\hfil\cr\kern0.1em--\hfil\cr}%
%EndExpansion
$ line for both, but with a higher temperature the line moves up and a
little to the right on the graph. Since the temperature is the same along
the whole red dotted line, we call such a line an \emph{isotherm.} The green
solid line is a different isotherm with a different temperature. By thinking
of which isotherm line a point is on we can get an idea of temperature on a
PV diagram.

\section{Van der Waals Gasses}

But of course our ideal gas is missing pieces of the physics of actual
gases. For our physics students we have an upper division class in
thermodynamics that fills in many of these missing pieces. We also have some
quantum mechanics classes that help. Chemists have a thermodynamics class
and their physical chemistry class. Mechanical engineers have a thermal
class as well. So some of the details will just have to wait. But if you are
in one of the disciplines that won't take these fun classes, we should leave
you with a slight improvement. We can add in two pieces of physics.

\begin{enumerate}
\item As a real gas get's compressed the molecules are more and more
attracted to each other. They are closer to becoming bonded into a liquid or
a solid. This reduces the pressure in the gas

\item At higher pressures the gas molecules, themselves, constrain where
other gas molecules can go (at some point they hit each other!).
\end{enumerate}

A researcher named Van der Waals came up with a new approximation, so the
result of adding in these two additional bits of physics is called the \emph{%
Van der Waals gas law}. Let's see what it looks like.

For the first bit of new physics we can add a term to the pressure that is
proportional to the molar density squared.%
\begin{equation*}
P+a\left( \frac{n}{%
%TCIMACRO{\TeXButton{V}{\ooalign{\hfil$V$\hfil\cr\kern0.1em--\hfil\cr}}}%
%BeginExpansion
\ooalign{\hfil$V$\hfil\cr\kern0.1em--\hfil\cr}%
%EndExpansion
}\right) ^{2}
\end{equation*}%
where the constant $a$ depends on the specific gas. You may not find it
obvious why this is the right form for our new gas law equation to make the
pressure come out right, but it does adjust the pressure under compression
and it turns out it kind of works. So we will replace the solitary $P$ in
our ideal gas law with $P+a\left( \frac{n}{%
%TCIMACRO{\TeXButton{V}{\ooalign{\hfil$V$\hfil\cr\kern0.1em--\hfil\cr}}}%
%BeginExpansion
\ooalign{\hfil$V$\hfil\cr\kern0.1em--\hfil\cr}%
%EndExpansion
}\right) ^{2}.$

For the second bit of new physics we can subtract a term proportional to the
volume of the molecules, themselves, from the volume of the whole gas. We
can call this (sort of) molar volume $b$. 
\begin{equation*}
%TCIMACRO{\TeXButton{V}{\ooalign{\hfil$V$\hfil\cr\kern0.1em--\hfil\cr}}}%
%BeginExpansion
\ooalign{\hfil$V$\hfil\cr\kern0.1em--\hfil\cr}%
%EndExpansion
-nb
\end{equation*}%
and put it where the $%
%TCIMACRO{\TeXButton{V}{\ooalign{\hfil$V$\hfil\cr\kern0.1em--\hfil\cr}}}%
%BeginExpansion
\ooalign{\hfil$V$\hfil\cr\kern0.1em--\hfil\cr}%
%EndExpansion
$ was in our ideal gas law. All together our new gas law becomes 
\begin{equation*}
\left( P+a\left( \frac{n}{%
%TCIMACRO{\TeXButton{V}{\ooalign{\hfil$V$\hfil\cr\kern0.1em--\hfil\cr}}}%
%BeginExpansion
\ooalign{\hfil$V$\hfil\cr\kern0.1em--\hfil\cr}%
%EndExpansion
}\right) ^{2}\right) \left( 
%TCIMACRO{\TeXButton{V}{\ooalign{\hfil$V$\hfil\cr\kern0.1em--\hfil\cr}}}%
%BeginExpansion
\ooalign{\hfil$V$\hfil\cr\kern0.1em--\hfil\cr}%
%EndExpansion
-nb\right) =nRT
\end{equation*}

Although it is an improvement, it still doesn't have all the physics that
might be needed to describe a real gas. For example, look at the graph
below. You see a plot of isotherms for a Van der Waals gas (CO$_{2}$).

\FRAME{dtbpF}{3.0519in}{2.3549in}{0pt}{}{}{Figure}{\special{language
"Scientific Word";type "GRAPHIC";maintain-aspect-ratio TRUE;display
"USEDEF";valid_file "T";width 3.0519in;height 2.3549in;depth
0pt;original-width 3.0087in;original-height 2.3151in;cropleft "0";croptop
"1";cropright "1";cropbottom "0";tempfilename
'SYDIL002.wmf';tempfile-properties "XPR";}}For high temperatures we see
something that looks a lot like our ideal gas graph. But at a certain
temperature for each gas the graph flattens out (middle orange line). That
is really not behavior that we expect from an ideal gas. For this line the
Van der Walls equation is better. The temperature for the orange flattened
line is called a \emph{critical temperature, }$T_{c},$ and for our CO$_{2}$
it is about $304\unit{K}.$ But look at the lower green line) this line dips
down and make a valley. This isn't possible. Our \textquotedblleft
improved\textquotedblright\ gas law is giving us unphysical predictions.
What is missing is that we can have a change from gas to liquid at this
pressure and temperature and we didn't include such changes in our Van der
Waals gas law equation. Above $T_{c}$ our Van der Waals equation isn't too
bad. But below $T_{c}$ it is not as useful. Of course we could improve on
our Van der Waals equation, but that is a subject for those upper division
classes in thermodynamics. For this class, let's concentrate on the ideal
gas equation even though we know it has some limitations.

\chapter{30 Kinetic Theory of Gases 2.2.2, 2.2.3}

%TCIMACRO{\TeXButton{\chaptermark{Chapter 32}}{\chaptermark{Chapter 32}}}%
%BeginExpansion
\chaptermark{Chapter 32}%
%EndExpansion

We now know we have gases with internal energy. But so far, as we have
considered the motion of atoms with thermal energy we have only been able to
treat whole samples of a gas at a time. Usually we have measured the size of
these samples as numbers of moles of atoms or molecules. But you might guess
that to get the internal energy or temperature or pressure of the sample we
are sort of taking averages of the motions of the atoms or molecules that
make up that sample. Let's make this more explicit in this lecture.

%TCIMACRO{%
%\TeXButton{Fundamental Concepts}{\vspace{0.10in}\hspace{-0.37in}{\Large Fundamental Concepts\vspace{0.2in}}}}%
%BeginExpansion
\vspace{0.10in}\hspace{-0.37in}{\Large Fundamental Concepts\vspace{0.2in}}%
%EndExpansion

\begin{itemize}
\item In a gas, the molecules have different speeds

\item The distribution of speeds is given by the Boltzmann distribution law

\item the root mean square speed $\left( v_{rms}=\sqrt{\overline{v^{2}}}%
\right) $is a better measure of the speed of gas molecules than the average
velocity

\item Pressure for an ideal gas is given by $P=\frac{2}{3}\frac{N}{%
%TCIMACRO{\TeXButton{V}{\ooalign{\hfil$V$\hfil\cr\kern0.1em--\hfil\cr}}}%
%BeginExpansion
\ooalign{\hfil$V$\hfil\cr\kern0.1em--\hfil\cr}%
%EndExpansion
}\left( \frac{1}{2}m\overline{v^{2}}\right) $

\item Temperature is related directly to the average kinetic energy of the
molecules of a gas

\item We call a way the molecule can move a \emph{degree of freedom.}

\item For an ideal gas $C_{V}=\frac{3}{2}R$

\item For an ideal gas $E_{int}=\frac{3}{2}k_{B}T$

\item Mean Free Path
\end{itemize}

\section{Molecular Model of an Ideal Gas}

We have hinted that the internal energy of a system must be the energy
associated with the atoms and molecules that make the system. By making this
association in our mental model more concrete, and more mathematical, we
will be able to see how the thermodynamic laws we have studied are generated
by the basic laws of motion. But we will limit our study to ideal gases in
this class. So let's state some assumptions that follow from the ideal gas
approximation.%
\begin{equation*}
\begin{tabular}{|c|}
\hline
{\large Ideal Gas Model} \\ \hline
The number of molecules in the gas is very large, \\ \hline
The average separation between molecules is large compared the their
dimensions \\ \hline
The molecules obey Newton's laws of motion, \\ \hline
On the whole the molecules move randomly \\ \hline
The molecules interact only by short-range forces during elastic collisions
\\ \hline
The molecules make elastic collisions with the walls \\ \hline
The gas under consideration is a pure substance; that is; all molecules are
identical. \\ \hline
\end{tabular}%
\end{equation*}

We will often say \textquotedblleft molecule\textquotedblright\ but for an
ideal gas atoms and molecules act alike.

\section{Microscopic Definition of Pressure}

We want to express macroscopic quantities like pressure and volume in terms
of microscopic effects. The easiest to start with is pressure

\FRAME{dtbpF}{2.2911in}{2.2538in}{0in}{}{}{Figure}{\special{language
"Scientific Word";type "GRAPHIC";maintain-aspect-ratio TRUE;display
"USEDEF";valid_file "T";width 2.2911in;height 2.2538in;depth
0in;original-width 2.3062in;original-height 2.2689in;cropleft "0";croptop
"1";cropright "1";cropbottom "0";tempfilename
'SW9LUS05.wmf';tempfile-properties "XPR";}}Picture a particle in a cubical
box of side length $d.$

The particle could have any direction in the box. And that direction would
have $x,$ $y,$ and $z$ components. And as physicists often do we would
divide our three dimensional problem into three one dimensional problems. So
we can start by considering just the $x$ part of the motion. As the particle
approaches the wall of the box we may say that it has velocity $v_{x}.$ The
particle impacts the wall and bounces off again. Now it travels with
velocity $-v_{x}$.

\FRAME{dtbpF}{1.7011in}{2.2494in}{0pt}{}{}{Figure}{\special{language
"Scientific Word";type "GRAPHIC";maintain-aspect-ratio TRUE;display
"USEDEF";valid_file "T";width 1.7011in;height 2.2494in;depth
0pt;original-width 2.9499in;original-height 3.9081in;cropleft "0";croptop
"1";cropright "1";cropbottom "0";tempfilename
'SW9LUS06.wmf';tempfile-properties "XPR";}}

We can think of why this would be. Think of conservation of momentum. Before
the collision we may take the velocity of the wall to be zero, and we may
imagine the wall to have a very large mass (which it does, compared to the
particle!). We can call the particle mass $m$ and the wall mass $M,$ but the
wall will not move with any speed after the collision so $v_{wall_{f}}=0.$
The initial momentum for the system would be%
\begin{eqnarray*}
P_{i} &=&Mv_{wall_{i}}+mv_{x_{i}} \\
&=&mv_{x_{i}}
\end{eqnarray*}%
and the final momentum of the system would be 
\begin{eqnarray*}
P_{f} &=&Mv_{wall_{f}}+mv_{xf} \\
&=&-mv_{x_{f}}
\end{eqnarray*}%
because the wall still isn't going to move. Then for conservation of
momentum 
\begin{equation*}
P_{i}=P_{f}
\end{equation*}%
We have 
\begin{equation*}
mv_{x_{i}}=-mv_{x_{f}}
\end{equation*}%
Then the magnitude of the velocities must be equal but the direction
changed. And the \emph{change} in momentum for the ball must be 
\begin{eqnarray*}
\Delta p_{ball} &=&p_{xf}-p_{xi} \\
&=&-mv_{x_{f}}-\left( mv_{x_{i}}\right) \\
&=&-2mv_{x_{i}}
\end{eqnarray*}

%TCIMACRO{%
%\TeXButton{Question 123.14.6}{\marginpar {
%\hspace{-0.5in}
%\begin{minipage}[t]{1in}
%\small{Question 123.14.6}
%\end{minipage}
%}}}%
%BeginExpansion
\marginpar {
\hspace{-0.5in}
\begin{minipage}[t]{1in}
\small{Question 123.14.6}
\end{minipage}
}%
%EndExpansion
Recall that the impulse from the collision of the particle on the wall would
be 
\begin{eqnarray*}
\bar{F}_{i,\text{on molecule}}\Delta t_{\text{collision}} &=&\Delta p_{x_{i}}
\\
&=&-2mv_{x_{i}}
\end{eqnarray*}%
This is the impulse momentum theorem from Principles of Physics I. The time
of the collision is very short. This force is the force \emph{of} the wall 
\emph{on} the particle.

Now consider that if a particle bounces off one wall in the $x$-direction,
it must travel to the opposite wall and back before it can bounce off the
same wall again. That is, the time between bounces is roughly%
\begin{equation}
\Delta t=\frac{2d}{v_{x_{i}}}
\end{equation}

Averages are funny things, it is perfectly legal to redefine the time of our
average force to be this travel time $\Delta t.$ What does this mean? In the
first graph of the next figure we have the full definition of impulse 
\begin{equation}
\Delta p=\int_{t_{i}}^{t_{f}}\mathbf{F}dt
\end{equation}%
which says the impulse is the area under the $F$ vs. $t$ graph. That area is
colored in (orange if you are seeing this in color) in the first graph of
the figure. \FRAME{dtbpF}{4.5697in}{3.0856in}{0pt}{}{}{Figure}{\special%
{language "Scientific Word";type "GRAPHIC";maintain-aspect-ratio
TRUE;display "USEDEF";valid_file "T";width 4.5697in;height 3.0856in;depth
0pt;original-width 4.5178in;original-height 3.0415in;cropleft "0";croptop
"1";cropright "1";cropbottom "0";tempfilename
'TC7K5900.wmf';tempfile-properties "XPR";}}The average force $\bar{F}$ is
pictured in the second graph. It has the same orange area as the first
graph, but is flat, representing the average force (red dotted line). The
final graph has a $\bar{F}$ averaged over a much larger $\Delta t.$ The area
is the same, but the magnitude of the average force is much smaller.

This may seem useless, but at any rate, we can do this. If we take a larger
averaging time we get a lower average force. And this is what we want to do.
Let's take our average of $F$ over the time 
\begin{equation*}
\Delta t=\frac{2d}{v_{x_{i}}}
\end{equation*}%
which is the time it takes our particle to travel away from our wall, and
bounce back to the wall. We know that sometime within this $\Delta t$ the
collision with our wall actually occurs. But this time is much longer than
the actual collision time, so if we use this time our average force will be
smaller than the short collision force. The change in momentum is still 
\begin{equation*}
\bar{F}\Delta t=-2mv_{x_{i}}
\end{equation*}%
So, we can write%
\begin{eqnarray*}
\bar{F}\frac{2d}{v_{x_{i}}} &=&-2mv_{x_{i}} \\
\bar{F} &=&-\frac{mv_{x_{i}}^{2}}{d}
\end{eqnarray*}%
So far our force has been the force of the wall on the particle, but we know
by Newton's third law that the particle force must be equal and opposite. 
\begin{equation*}
\bar{F}_{i,\text{on molecule}}=-\bar{F}_{i,\text{ on the wall}}
\end{equation*}%
Then if we shorten $\bar{F}_{i,\text{ on the wall}}$ to just $\bar{F}_{iw}$ 
\begin{equation*}
\bar{F}_{iw}=\frac{mv_{x_{i}}^{2}}{d}
\end{equation*}%
for one particle (molecule or atom).

For many particles%
\begin{eqnarray*}
\bar{F}_{w} &=&\dsum\limits_{i=1}^{N}\frac{mv_{x_{i}}^{2}}{d} \\
&=&\frac{m}{d}\dsum\limits_{i=1}^{N}v_{x_{i}}^{2}
\end{eqnarray*}%
And now that we have the force on the wall our average force seems more
valid. The many particles will hit the wall over and over again and not all
at the same time, so there is always some force on the wall. And the force
the wall will feel is an average over all these little forces. We have found
the average force on the container wall!

%TCIMACRO{%
%\TeXButton{Question 123.14.7}{\marginpar {
%\hspace{-0.5in}
%\begin{minipage}[t]{1in}
%\small{Question 123.14.7}
%\end{minipage}
%}}}%
%BeginExpansion
\marginpar {
\hspace{-0.5in}
\begin{minipage}[t]{1in}
\small{Question 123.14.7}
\end{minipage}
}%
%EndExpansion
Let's rewire this equation for our average force using the mathematical
definition of an average. 
\begin{equation*}
\bar{x}=\frac{1}{N}\dsum\limits_{i=1}^{N}x_{i}
\end{equation*}%
It almost looks like we have an average velocity in our equation. If we were
to calculate the average velocity,\emph{\ but square the velocity} before we
take the average, it would look like this%
\begin{equation*}
\overline{v_{x_{i}}^{2}}=\frac{1}{N}\dsum\limits_{i=1}^{N}v_{x_{i}}^{2}
\end{equation*}%
And this is really close to what we see in our force equation! Let's put in
the missing factor of $1/N,$ and let's use a bar over $v_{x_{i}}^{2}$ to
mean the average, so that 
\begin{eqnarray*}
\bar{F} &=&\frac{m}{d}\frac{N}{N}\dsum\limits_{i=1}^{N}v_{x_{i}}^{2} \\
&=&\frac{m}{d}N\overline{v_{x}^{2}}
\end{eqnarray*}%
From this we can see that the force on the wall is proportional to the
average of the velocity squared!

Going back to a single molecule, we remember that the velocity of that
single molecule has components. So far we have just used the $x$-component.
But we have $y$ and $z$-components. The velocity squared would be (think of
a dot product)%
\begin{equation*}
v_{i}^{2}=v_{x_{i}}^{2}+v_{y_{i}}^{2}+v_{z_{i}}^{2}
\end{equation*}%
where $v_{i}^{2}$ is the magnitude of the velocity squared for the $i$-th
molecule.

Then if we take the averages of each piece like we did with $v_{x_{i}}^{2}$ 
\begin{equation*}
\overline{v^{2}}=\overline{v_{x}^{2}}+\overline{v_{y}^{2}}+\overline{%
v_{z}^{2}}
\end{equation*}

Our atoms or molecules are so small that the force due to gravity on them is
very small. To a good approximation we can mostly ignore it. Then there is
not much reason to prefer the $x$-direction, we could have chosen the $y$%
-direction and had the same value. In this approximation, then 
\begin{equation*}
v_{x}^{2}=v_{y}^{2}=v_{z}^{2}
\end{equation*}%
and%
\begin{equation*}
\overline{v^{2}}=3v_{x}^{2}
\end{equation*}%
So our average force on the wall from the molecules could be written as%
\begin{eqnarray*}
\bar{F} &=&\frac{m}{d}N\overline{v_{x}^{2}} \\
&=&\frac{m}{d}N\frac{1}{3}\overline{v^{2}} \\
&=&\frac{mN}{3d}\overline{v^{2}} \\
&=&\frac{N}{3d}\left( m\overline{v^{2}}\right)
\end{eqnarray*}%
Let's rewire this average force from all the molecules 
\begin{equation*}
\bar{F}=\frac{N}{3d}\left( m\overline{v^{2}}\right)
\end{equation*}
And we know that 
\begin{equation*}
P=\frac{F}{A}
\end{equation*}%
Finally we can find the pressure on the wall. We know%
\begin{eqnarray*}
P &=&\frac{F}{A} \\
&=&\frac{F}{d^{2}} \\
&=&\frac{1}{d^{2}}\left( \frac{N}{3d}\left( m\overline{v^{2}}\right) \right)
\\
&=&\frac{1}{d^{3}}\frac{N}{3}\left( m\overline{v^{2}}\right) \\
&=&\frac{1}{3}\frac{N}{%
%TCIMACRO{\TeXButton{V}{\ooalign{\hfil$V$\hfil\cr\kern0.1em--\hfil\cr}}}%
%BeginExpansion
\ooalign{\hfil$V$\hfil\cr\kern0.1em--\hfil\cr}%
%EndExpansion
}\left( m\overline{v^{2}}\right) \\
&=&\frac{2}{3}\frac{N}{%
%TCIMACRO{\TeXButton{V}{\ooalign{\hfil$V$\hfil\cr\kern0.1em--\hfil\cr}}}%
%BeginExpansion
\ooalign{\hfil$V$\hfil\cr\kern0.1em--\hfil\cr}%
%EndExpansion
}\left( \frac{1}{2}m\overline{v^{2}}\right)
\end{eqnarray*}%
That last step might seem unnecessary, but look how cute the $1/2$ looks
with a $m\overline{v^{2}}$ next to it! We can recognize the part in
parenthesis as the average kinetic energy of the molecule! So

\begin{equation}
P=\frac{2}{3}\frac{N}{%
%TCIMACRO{\TeXButton{V}{\ooalign{\hfil$V$\hfil\cr\kern0.1em--\hfil\cr}}}%
%BeginExpansion
\ooalign{\hfil$V$\hfil\cr\kern0.1em--\hfil\cr}%
%EndExpansion
}\left( \frac{1}{2}m\overline{v^{2}}\right)
\end{equation}%
and we see that the pressure of the gas is proportional to the number of
molecules per unit volume and to the average kinetic energy of the
molecules! This tells us where pressure comes from. It is due to the average
kinetic energy of the molecules. In fact, this equation shows some of the
things we have learned about gas laws. \ For example, $P$ is inversely
proportional to $%
%TCIMACRO{\TeXButton{V}{\ooalign{\hfil$V$\hfil\cr\kern0.1em--\hfil\cr}}}%
%BeginExpansion
\ooalign{\hfil$V$\hfil\cr\kern0.1em--\hfil\cr}%
%EndExpansion
$!

\section{Temperature}

Finally, we are ready to define temperature!

%TCIMACRO{%
%\TeXButton{Question 123.15.1}{\marginpar {
%\hspace{-0.5in}
%\begin{minipage}[t]{1in}
%\small{Question 123.15.1}
%\end{minipage}
%}}}%
%BeginExpansion
\marginpar {
\hspace{-0.5in}
\begin{minipage}[t]{1in}
\small{Question 123.15.1}
\end{minipage}
}%
%EndExpansion
%TCIMACRO{%
%\TeXButton{Question 123.15.2}{\marginpar {
%\hspace{-0.5in}
%\begin{minipage}[t]{1in}
%\small{Question 123.15.2}
%\end{minipage}
%}}}%
%BeginExpansion
\marginpar {
\hspace{-0.5in}
\begin{minipage}[t]{1in}
\small{Question 123.15.2}
\end{minipage}
}%
%EndExpansion
%TCIMACRO{%
%\TeXButton{Equation on Board}{\marginpar {
%\hspace{-0.5in}
%\begin{minipage}[t]{1in}
%\small{Equation on Board}
%\end{minipage}
%}}}%
%BeginExpansion
\marginpar {
\hspace{-0.5in}
\begin{minipage}[t]{1in}
\small{Equation on Board}
\end{minipage}
}%
%EndExpansion
We now know that 
\begin{equation*}
P=\frac{2}{3}\frac{N}{%
%TCIMACRO{\TeXButton{V}{\ooalign{\hfil$V$\hfil\cr\kern0.1em--\hfil\cr}}}%
%BeginExpansion
\ooalign{\hfil$V$\hfil\cr\kern0.1em--\hfil\cr}%
%EndExpansion
}\left( \frac{1}{2}m\overline{v^{2}}\right)
\end{equation*}%
and that this has some of the properties of the ideal gas law. We can
rearrange this to look a little more like the ideal gas law. Let's start
with our pressure, but put the volume on the other side next to the pressure
where it is in the ideal gas law. 
\begin{equation*}
P%
%TCIMACRO{\TeXButton{V}{\ooalign{\hfil$V$\hfil\cr\kern0.1em--\hfil\cr}}}%
%BeginExpansion
\ooalign{\hfil$V$\hfil\cr\kern0.1em--\hfil\cr}%
%EndExpansion
=\frac{2}{3}N\left( \frac{1}{2}m\overline{v^{2}}\right)
\end{equation*}%
But temperature does not appear in this equation. It should be on the right
and side. It must be that something on that right hand side is really our
temperature!

Let's recall that we can write the ideal gas law as 
\begin{equation*}
P%
%TCIMACRO{\TeXButton{V}{\ooalign{\hfil$V$\hfil\cr\kern0.1em--\hfil\cr}}}%
%BeginExpansion
\ooalign{\hfil$V$\hfil\cr\kern0.1em--\hfil\cr}%
%EndExpansion
=Nk_{B}T
\end{equation*}%
If we set these two equations' right hand sides equal to each other%
\begin{equation*}
Nk_{B}T=\frac{2}{3}N\left( \frac{1}{2}m\overline{v^{2}}\right)
\end{equation*}%
then%
\begin{equation}
T=\frac{2}{3}\frac{1}{k_{B}}\left( \frac{1}{2}m\overline{v^{2}}\right)
\end{equation}%
This is fantastic! We have finally defined temperature. Temperature is the
average kinetic energy of the molecules! There is a constant out front to
fix the units. But Temperature is the average kinetic energy of the
molecules.

Solving for the average kinetic energy gives 
\begin{equation*}
\frac{1}{2}m\overline{v^{2}}=\frac{3}{2}k_{B}T
\end{equation*}%
So, the average kinetic energy is 
\begin{equation*}
\bar{K}_{mol}=\frac{1}{2}m\overline{v^{2}}=\frac{3}{2}k_{B}T
\end{equation*}%
Note that in each direction we should have 
\begin{equation*}
\bar{K}_{x}=\frac{1}{2}m\overline{v_{x}^{2}}=\frac{1}{2}k_{B}T
\end{equation*}%
and%
\begin{equation*}
\bar{K}_{y}=\frac{1}{2}k_{B}T
\end{equation*}%
and 
\begin{equation*}
\bar{K}_{z}=\frac{1}{2}k_{B}T
\end{equation*}%
So in each direction the molecules have 
\begin{equation}
\bar{K}_{i}=\frac{1}{2}k_{B}T
\end{equation}%
We call a way the molecule can move a \emph{degree of freedom. }Our
molecules can move in the $x,$ $y,$ and $z,$ direction, so our molecules
have three degrees of freedom. Each degree of freedom contributes $\frac{1}{2%
}k_{B}T$ worth of energy.

Of course, we only allow translational energy in our formulation for ideal
gasses, but for more complex molecules we could have rotational energy
(another way to move) and vibrational energy (yet another way to move) etc.
And each way the molecule can move is another degree of freedom that would
contribute $\frac{1}{2}k_{B}T.$ But remember, these more complex molecules
would not be ideal gasses.

For the collection of $N$ molecules, we have all together 
\begin{eqnarray}
K &=&N\left( \frac{3}{2}k_{B}T\right) \\
&=&\frac{3}{2}Nk_{B}T \\
&=&\frac{3}{2}nRT  \notag
\end{eqnarray}%
and for our ideal gas the only form of internal energy we have is the
kinetic energy, so we see that the internal energy depends only on the
temperature.

This is only true under our assumption that the molecules or atoms are
structureless. So this works well for actual monotonic gases, but not as
well for diatomic gasses, etc.

Let's do a problem. Suppose a colleague tells you that in her Statistical
Mechanics class they defined a root-mean-squared velocity for atoms in a
gas. That is, they took the velocity, squared it, and then took the average.
They got 
\begin{equation*}
v_{rms}=\sqrt{\frac{3k_{B}T}{m}}
\end{equation*}%
Let's show that this works from our kinetic theory of gasses.

We know that each degree of freedom contributes $\frac{1}{2}k_{B}T$ worth of
internal energy for the gas molecules. We just found that 
\begin{equation*}
\bar{K}=\frac{K}{N}=\frac{3}{2}k_{B}T
\end{equation*}%
and we know that 
\begin{equation*}
\bar{K}=\frac{1}{2}m\overline{v^{2}}
\end{equation*}%
Setting these equal gives%
\begin{eqnarray*}
\frac{3}{2}k_{B}T &=&\frac{1}{2}m\overline{v^{2}} \\
3k_{B}T &=&m\overline{v^{2}} \\
\frac{3k_{B}T}{m} &=&\overline{v^{2}}
\end{eqnarray*}%
so

\begin{equation*}
v_{rms}=\sqrt{\overline{v^{2}}}=\sqrt{\frac{3k_{B}T}{m}}
\end{equation*}%
which is just what our colleague said. This shows the power of the simple
idea of the ideal gas approximation.

The $rms$ speed tells us that if the molecule or atom is tiny (has little
mass) it must move very fast. If it is larger, it will move more slowly.
Let's look at how molecules travel as they go through a gas.

\section{Mean Free Path}

We started this chapter by saying how molecules interact with the walls of a
box. But the molecules also bounce off each other. The average distance a
molecule travels before a molecule-molecule collision is called the \emph{%
mean free path}.\FRAME{fhF}{3.109in}{1.5584in}{0pt}{}{}{Figure}{\special%
{language "Scientific Word";type "GRAPHIC";maintain-aspect-ratio
TRUE;display "USEDEF";valid_file "T";width 3.109in;height 1.5584in;depth
0pt;original-width 11.1353in;original-height 5.5659in;cropleft "0";croptop
"1";cropright "1";cropbottom "0";tempfilename
'SW9MF20R.wmf';tempfile-properties "XPR";}}We could give this a symbol,
let's re-use the Greek letter $\lambda $%
\begin{eqnarray*}
\lambda &=&\frac{\text{Length traveled}}{\text{Number of collisions}} \\
&=&\frac{L}{N_{\func{col}}}
\end{eqnarray*}%
This will give the average length before a collision happens.

The average length traveled is easy%
\begin{equation*}
L=\bar{v}\Delta t
\end{equation*}

But to predict the number of collisions is harder. To do this problem, let's
play a geometrical trick. First let's assume the molecules are spheres with
diameter $d.$ \FRAME{fhF}{4.3171in}{2.162in}{0pt}{}{}{Figure}{\special%
{language "Scientific Word";type "GRAPHIC";maintain-aspect-ratio
TRUE;display "USEDEF";valid_file "T";width 4.3171in;height 2.162in;depth
0pt;original-width 12.0252in;original-height 6.0096in;cropleft "0";croptop
"1";cropright "1";cropbottom "0";tempfilename
'SW9P0O0Z.wmf';tempfile-properties "XPR";}}We see that a collision does not
happen unless the distance between the molecules is less than or equal to $%
2d.$ All we want is the number of collisions (we don't care about what
happens after the collisions) and this way of counting collisions is easier
(and easier to code). We can model this interaction as a large particle of
size $2d$ and many point particles. We let the large particle move in a
straight line to create a cylindrical path. \FRAME{dtbpF}{3.1453in}{2.4578in%
}{0pt}{}{}{Figure}{\special{language "Scientific Word";type
"GRAPHIC";maintain-aspect-ratio TRUE;display "USEDEF";valid_file "T";width
3.1453in;height 2.4578in;depth 0pt;original-width 3.0995in;original-height
2.4163in;cropleft "0";croptop "1";cropright "1";cropbottom "0";tempfilename
'SW9P0O10.wmf';tempfile-properties "XPR";}}All point particles in the
cylinder will collide with our large particle. The volume of such a big
cylindrical path is%
\begin{equation*}
%TCIMACRO{\TeXButton{V}{\ooalign{\hfil$V$\hfil\cr\kern0.1em--\hfil\cr}}}%
%BeginExpansion
\ooalign{\hfil$V$\hfil\cr\kern0.1em--\hfil\cr}%
%EndExpansion
=(\pi d^{2})L
\end{equation*}%
But the length, $L,$ of the cylinder will be 
\begin{equation*}
L=\bar{v}\Delta t
\end{equation*}%
so the volume will be%
\begin{equation}
%TCIMACRO{\TeXButton{V}{\ooalign{\hfil$V$\hfil\cr\kern0.1em--\hfil\cr}}}%
%BeginExpansion
\ooalign{\hfil$V$\hfil\cr\kern0.1em--\hfil\cr}%
%EndExpansion
=\pi d^{2}\bar{v}\Delta t
\end{equation}%
Then, if $n_{V}$ is the number of molecules per unit volume, our big
equivalent particle will hit 
\begin{equation}
N_{\func{col}}=\pi d^{2}\bar{v}\Delta tn_{V}
\end{equation}%
particles as it travels. This is the number of point molecules in the
cylinder. And it is the number of collisions.

The mean free path is the average distance traveled in between collisions. A
little thought shows that it is the distance traveled in a time interval $%
\Delta t$ divided by the number of collisions that occur in $\Delta t.$ We
now know both of these, so 
\begin{eqnarray}
\lambda &=&\frac{L}{N_{\func{col}}}=\frac{\bar{v}\Delta t}{\pi d^{2}\bar{v}%
\Delta tn_{V}}  \notag \\
&=&\frac{1}{\pi d^{2}n_{V}}
\end{eqnarray}%
We can also find the the average time between collisions. This is called the
mean free time. we could use 
\begin{equation*}
v=\frac{\Delta x}{\Delta t}
\end{equation*}%
and solve for $\Delta t$%
\begin{equation*}
\Delta t=\frac{\Delta x}{v}
\end{equation*}%
as we did in Principles of Physics I. But let's put in average values. The
average $\Delta x=\lambda $ and we could use $v_{rms}$ as our average speed. 
\begin{eqnarray*}
\tau &=&\frac{\lambda }{v_{rms}} \\
&=&\frac{1}{\pi d^{2}n_{V}v_{rms}}
\end{eqnarray*}%
The frequency of collisions is 
\begin{equation}
f=\frac{1}{\tau }=\pi d^{2}v_{rms}n_{V}
\end{equation}

It turns out that our simple model is off by a factor of $\sqrt{2}$ because
we assumed the molecules are stationary.\footnote{%
This is not obvious. But we won't do the derivation in our class. If you are
taking an upper division thermodynamics class you may get to do the more
rigorous math.} We can fix that easily%
\begin{equation}
\lambda =\frac{1}{\sqrt{2}\pi d^{2}n_{V}}=\frac{%
%TCIMACRO{\TeXButton{V}{\ooalign{\hfil$V$\hfil\cr\kern0.1em--\hfil\cr}}}%
%BeginExpansion
\ooalign{\hfil$V$\hfil\cr\kern0.1em--\hfil\cr}%
%EndExpansion
}{\sqrt{2}\pi d^{2}N}
\end{equation}%
and%
\begin{equation}
f=\sqrt{2}\pi d^{2}\bar{v}n_{V}=\frac{\sqrt{2}\pi d^{2}v_{rms}N}{%
%TCIMACRO{\TeXButton{V}{\ooalign{\hfil$V$\hfil\cr\kern0.1em--\hfil\cr}}}%
%BeginExpansion
\ooalign{\hfil$V$\hfil\cr\kern0.1em--\hfil\cr}%
%EndExpansion
}
\end{equation}%
and finally%
\begin{equation*}
\tau =\frac{1}{\sqrt{2}\pi d^{2}n_{V}\bar{v}}=\frac{%
%TCIMACRO{\TeXButton{V}{\ooalign{\hfil$V$\hfil\cr\kern0.1em--\hfil\cr}}}%
%BeginExpansion
\ooalign{\hfil$V$\hfil\cr\kern0.1em--\hfil\cr}%
%EndExpansion
}{\sqrt{2}\pi d^{2}Nv_{rms}}
\end{equation*}%
We can make a further modification if we assume an ideal gas. We know 
\begin{equation*}
P%
%TCIMACRO{\TeXButton{V}{\ooalign{\hfil$V$\hfil\cr\kern0.1em--\hfil\cr}}}%
%BeginExpansion
\ooalign{\hfil$V$\hfil\cr\kern0.1em--\hfil\cr}%
%EndExpansion
=Nk_{B}T
\end{equation*}%
so 
\begin{equation*}
\frac{%
%TCIMACRO{\TeXButton{V}{\ooalign{\hfil$V$\hfil\cr\kern0.1em--\hfil\cr}}}%
%BeginExpansion
\ooalign{\hfil$V$\hfil\cr\kern0.1em--\hfil\cr}%
%EndExpansion
}{N}=\frac{k_{B}T}{P}
\end{equation*}%
Then where there is a $%
%TCIMACRO{\TeXButton{V}{\ooalign{\hfil$V$\hfil\cr\kern0.1em--\hfil\cr}}}%
%BeginExpansion
\ooalign{\hfil$V$\hfil\cr\kern0.1em--\hfil\cr}%
%EndExpansion
/N$ we can substitute%
\begin{equation}
\lambda =\frac{k_{B}T}{\sqrt{2}\pi d^{2}P}
\end{equation}%
and 
\begin{eqnarray}
f &=&\frac{\sqrt{2}\pi d^{2}v_{rms}P}{k_{B}T} \\
\tau &=&\frac{k_{B}T}{\sqrt{2}\pi d^{2}v_{rms}P}
\end{eqnarray}

We should consider a problem to see how this works. Suppose we take a $%
41.\,\allowbreak 582\unit{mol}$ of helium (which really acts very like an
ideal gas) at $T=293\unit{K}$ and $P=1.013\times 10^{5}\unit{Pa}$ and we put
them in a container with volume $%
%TCIMACRO{\TeXButton{V}{\ooalign{\hfil$V$\hfil\cr\kern0.1em--\hfil\cr}}}%
%BeginExpansion
\ooalign{\hfil$V$\hfil\cr\kern0.1em--\hfil\cr}%
%EndExpansion
=1\unit{m}^{3}.$ The helium atoms have a diameter of about $d=0.031\unit{nm}$
and a mass of $6.64\times 10^{-27}\unit{kg}$. How often do we get collisions
between the molecules?

The mean free path would be 
\begin{equation*}
\lambda =\frac{k_{B}T}{\sqrt{2}\pi d^{2}P}=\frac{\left( 1.3806568\times
10^{-23}\frac{\unit{J}}{\unit{K}}\right) \left( 273\unit{K}\right) }{\sqrt{2}%
\pi \left( 0.031\unit{nm}\right) ^{2}\left( 1.013\times 10^{5}\unit{Pa}%
\right) }=\allowbreak 8.\,\allowbreak 714\,7\times 10^{-6}\unit{m}
\end{equation*}

The mean time between collisions needs to know the average speed and we
chose $v_{rms}$ for this which is a nice choice because we found earlier
that 
\begin{equation*}
v_{rms}=\sqrt{\frac{3k_{B}T}{m}}
\end{equation*}%
so 
\begin{equation*}
v_{rms}=\sqrt{\frac{3\left( 1.3806568\times 10^{-23}\frac{\unit{J}}{\unit{K}}%
\right) \left( 273\unit{K}\right) }{6.64\times 10^{-27}\unit{kg}}}%
=1305.\,\allowbreak 0\frac{\unit{m}}{\unit{s}}\allowbreak
\end{equation*}%
then the average time between collisions 
\begin{eqnarray*}
\tau &=&\frac{k_{B}T}{\sqrt{2}\pi d^{2}v_{rms}P} \\
&=&\frac{\left( 1.3806568\times 10^{-23}\frac{\unit{J}}{\unit{K}}\right)
\left( 273\unit{K}\right) }{\sqrt{2}\pi \left( 0.031\unit{nm}\right)
^{2}\left( 1305.\,\allowbreak 0\frac{\unit{m}}{\unit{s}}\right) \left(
1.013\times 10^{5}\unit{Pa}\right) } \\
&=&\allowbreak 6.\,\allowbreak 677\,9\times 10^{-9}\unit{s}
\end{eqnarray*}%
is and the frequency of collisions is%
\begin{equation*}
f=\frac{1}{6.\,\allowbreak 677\,9\times 10^{-9}\unit{s}}=1.\,\allowbreak
497\,5\times 10^{8}\frac{1}{\unit{s}}
\end{equation*}%
We would get quite a few collisions among our helium atoms.

Now that we know what temperature is, let's go back to thinking about
temperature rise as we transfer energy by heat. We studied heat capacities
and specific heats a few lectures ago, but we avoided gasses as we did so.
We need to find the temperature rise $\Delta T$ in a gas after transferring
energy by heat $Q$ for gasses.

\section{Molar Specific Heat}

We learned that 
\begin{equation*}
Q=mc\Delta T
\end{equation*}%
or that we could use the molar form of this equation 
\begin{equation*}
Q=nC\Delta T
\end{equation*}%
But we need to find the $C$ values for different gasses. To do this,
remember for an ideal gas, we only have translational kinetic energy, so it
must be true that 
\begin{equation*}
E_{int}=K_{tran}=\frac{3}{2}nRT
\end{equation*}%
But let's take a special case. If we take the system through a process at
constant volume,

\begin{equation*}
\Delta E_{int}=Q+w
\end{equation*}%
\begin{equation*}
Q=\Delta E_{int}
\end{equation*}%
because no work is done (we didn't change the volume) and for this spacial
case we can write $C$ as $C_{%
%TCIMACRO{\TeXButton{V}{\ooalign{\hfil$V$\hfil\cr\kern0.1em--\hfil\cr}}}%
%BeginExpansion
\ooalign{\hfil$V$\hfil\cr\kern0.1em--\hfil\cr}%
%EndExpansion
}$ to remind us that the volume didn't change.

\begin{equation*}
Q=nC_{%
%TCIMACRO{\TeXButton{V}{\ooalign{\hfil$V$\hfil\cr\kern0.1em--\hfil\cr}}}%
%BeginExpansion
\ooalign{\hfil$V$\hfil\cr\kern0.1em--\hfil\cr}%
%EndExpansion
}\Delta T
\end{equation*}%
so%
\begin{equation*}
\Delta E_{int}=nC_{%
%TCIMACRO{\TeXButton{V}{\ooalign{\hfil$V$\hfil\cr\kern0.1em--\hfil\cr}}}%
%BeginExpansion
\ooalign{\hfil$V$\hfil\cr\kern0.1em--\hfil\cr}%
%EndExpansion
}\Delta T
\end{equation*}%
Notice this is a change it internal energy. We added energy to our gas 
\begin{eqnarray*}
E_{int_{f}}-E_{int_{i}} &=&nC_{%
%TCIMACRO{\TeXButton{V}{\ooalign{\hfil$V$\hfil\cr\kern0.1em--\hfil\cr}}}%
%BeginExpansion
\ooalign{\hfil$V$\hfil\cr\kern0.1em--\hfil\cr}%
%EndExpansion
}\left( T_{f}-T_{i}\right) \\
&=&nC_{%
%TCIMACRO{\TeXButton{V}{\ooalign{\hfil$V$\hfil\cr\kern0.1em--\hfil\cr}}}%
%BeginExpansion
\ooalign{\hfil$V$\hfil\cr\kern0.1em--\hfil\cr}%
%EndExpansion
}T_{f}-nC_{%
%TCIMACRO{\TeXButton{V}{\ooalign{\hfil$V$\hfil\cr\kern0.1em--\hfil\cr}}}%
%BeginExpansion
\ooalign{\hfil$V$\hfil\cr\kern0.1em--\hfil\cr}%
%EndExpansion
}T_{i}
\end{eqnarray*}%
This strongly suggests that if $C_{%
%TCIMACRO{\TeXButton{V}{\ooalign{\hfil$V$\hfil\cr\kern0.1em--\hfil\cr}}}%
%BeginExpansion
\ooalign{\hfil$V$\hfil\cr\kern0.1em--\hfil\cr}%
%EndExpansion
}$ is constant%
\begin{equation*}
E_{int}=nC_{%
%TCIMACRO{\TeXButton{V}{\ooalign{\hfil$V$\hfil\cr\kern0.1em--\hfil\cr}}}%
%BeginExpansion
\ooalign{\hfil$V$\hfil\cr\kern0.1em--\hfil\cr}%
%EndExpansion
}T
\end{equation*}%
Then we can solve for $C_{%
%TCIMACRO{\TeXButton{V}{\ooalign{\hfil$V$\hfil\cr\kern0.1em--\hfil\cr}}}%
%BeginExpansion
\ooalign{\hfil$V$\hfil\cr\kern0.1em--\hfil\cr}%
%EndExpansion
}$%
\begin{equation*}
C_{%
%TCIMACRO{\TeXButton{V}{\ooalign{\hfil$V$\hfil\cr\kern0.1em--\hfil\cr}}}%
%BeginExpansion
\ooalign{\hfil$V$\hfil\cr\kern0.1em--\hfil\cr}%
%EndExpansion
}=\frac{1}{n}\frac{\Delta E_{int}}{\Delta T}
\end{equation*}%
And for very small changes in temperature%
\begin{equation*}
C_{V}=\frac{1}{n}\frac{dE_{int}}{dT}
\end{equation*}%
But for an ideal gas the internal energy is kinetic energy so we know that 
\begin{equation*}
E_{int}=\frac{3}{2}nRT
\end{equation*}%
So we can easily find $\frac{dE_{int}}{dT}$%
\begin{equation*}
\frac{dE_{int}}{dT}=\frac{3}{2}nR
\end{equation*}%
Then our value for $C_{V}$ is%
\begin{eqnarray*}
C_{%
%TCIMACRO{\TeXButton{V}{\ooalign{\hfil$V$\hfil\cr\kern0.1em--\hfil\cr}}}%
%BeginExpansion
\ooalign{\hfil$V$\hfil\cr\kern0.1em--\hfil\cr}%
%EndExpansion
} &=&\frac{1}{n}\frac{dE_{int}}{dT} \\
&=&\frac{1}{n}\frac{3}{2}nR \\
&=&\frac{3}{2}R
\end{eqnarray*}%
Knowing $R=8.314\frac{\unit{J}}{\unit{mol}\unit{K}},$ the numerical value is%
\begin{eqnarray*}
C_{%
%TCIMACRO{\TeXButton{V}{\ooalign{\hfil$V$\hfil\cr\kern0.1em--\hfil\cr}}}%
%BeginExpansion
\ooalign{\hfil$V$\hfil\cr\kern0.1em--\hfil\cr}%
%EndExpansion
} &=&\frac{3}{2}8.314\frac{\unit{J}}{\unit{mol}\unit{K}} \\
&=&12.\,\allowbreak 471\frac{\unit{J}}{\unit{mol}\unit{K}}
\end{eqnarray*}%
for all monotonic gasses. It turns out that this is a fairly good
approximation. Then 
\begin{eqnarray*}
Q &=&nC_{%
%TCIMACRO{\TeXButton{V}{\ooalign{\hfil$V$\hfil\cr\kern0.1em--\hfil\cr}}}%
%BeginExpansion
\ooalign{\hfil$V$\hfil\cr\kern0.1em--\hfil\cr}%
%EndExpansion
}\Delta T \\
&=&n\frac{3}{2}R\Delta T
\end{eqnarray*}%
for the case of no volume change. This seems useful!

You might find the restriction of heating at a constant volume somewhat
constrictive. We will lift this restriction in a later lecture.

\chapter{31 Speeds and Systems 2.2.4, 2.3.1}

%TCIMACRO{\TeXButton{\chaptermark{Chapter 33}}{\chaptermark{Chapter 33}}}%
%BeginExpansion
\chaptermark{Chapter 33}%
%EndExpansion

We found in our last lecture that 
\begin{equation*}
E_{int}=K_{tran}=\frac{3}{2}nRT
\end{equation*}%
for ideal gasses. But most gasses are not ideal. We can extend our theory to
diatomic gasses by considering more than just kinetic energy.

We also found in our last lecture that $v_{rms}$ is a nice average kind of
speed that tells us in bulk how our molecules are moving. But surely if we
have an average speed, some molecules are moving more slowly and some are
moving more quickly. Can we say how many molecules have a particular speed?
We will take on this problem in this lecture.

We will also consider energy conservation in thermodynamic processes. In
thermodynamics, conservation of energy is known as the \emph{first law of
thermodynamics}.

%TCIMACRO{%
%\TeXButton{Fundamental Concepts}{\vspace{0.10in}\hspace{-0.37in}{\Large Fundamental Concepts\vspace{0.2in}}}}%
%BeginExpansion
\vspace{0.10in}\hspace{-0.37in}{\Large Fundamental Concepts\vspace{0.2in}}%
%EndExpansion

\begin{itemize}
\item Equipartition of Energy

\item Partial Pressure and Vapor Pressure

\item For solids $C_{%
%TCIMACRO{\TeXButton{V}{\ooalign{\hfil$V$\hfil\cr\kern0.1em--\hfil\cr}}}%
%BeginExpansion
\ooalign{\hfil$V$\hfil\cr\kern0.1em--\hfil\cr}%
%EndExpansion
}\approx 3R$

\item The number of molecules that have a specific speed in an ideal gas is
given by the Maxwell-Boltsman distribution $N_{v}=4\pi N\left( \frac{m}{2\pi
k_{B}T}\right) ^{\frac{3}{2}}v^{2}e^{-\frac{mv^{2}}{2k_{B}T}}$

\item We characterize a gas by giving the $v_{rms,}$ the average speed $\bar{%
v}$, and the most probable speed $v_{mp}$
\end{itemize}

\section{The Equipartition of Energy}

For an ideal gas 
\begin{equation*}
E_{int}=K_{tran}=\frac{3}{2}nRT
\end{equation*}%
But there are few gasses that really act like ideal gasses. Nobel gasses
can, because they don't easily form bonds. So a gas of Argon or Helium atoms
is very like an ideal gas. But our atmosphere is mostly made of nitrogen and
oxygen which form diatomic molecules. Diatomic molecules have internal
structure. They are not ideal gasses. It would be good to try to apply what
we have learned to gasses like these that are so important to our life in
this mortal state.

We can extend our ideal gas model a little by using what we know about
degrees of freedom. Remember that we found that for each degree of freedom
the internal energy was 
\begin{equation}
E_{i}=\frac{1}{2}k_{B}T
\end{equation}%
We found for an ideal monotonic gas that the internal energy was 
\begin{equation}
E_{int}=3E_{i}=\frac{3}{2}k_{B}T
\end{equation}

where each $E_{i}$ came from moving in the $x,$ $y,$ or $z$-directions. We
call moving in the $x,$ $y,$ or $z$-directions \emph{translational degrees
of freedom}. \FRAME{dtbpF}{1.9519in}{1.9476in}{0pt}{}{}{Figure}{\special%
{language "Scientific Word";type "GRAPHIC";maintain-aspect-ratio
TRUE;display "USEDEF";valid_file "T";width 1.9519in;height 1.9476in;depth
0pt;original-width 2.4673in;original-height 2.4621in;cropleft "0";croptop
"1";cropright "1";cropbottom "0";tempfilename
'THGWPX00.wmf';tempfile-properties "XPR";}}But in a gas made of diatomic
molecules, the molecules have more degrees of freedom. They can rotate about
any of the axes. And from Principles of Physics I we know it takes energy to
rotate. This is another way we can put energy into our gas. Rotation is an
additional degree of freedom. But there are several ways we can rotate our
diatomic molecule.\FRAME{dtbpF}{2.3307in}{2.3307in}{0pt}{}{}{Figure}{\special%
{language "Scientific Word";type "GRAPHIC";maintain-aspect-ratio
TRUE;display "USEDEF";valid_file "T";width 2.3307in;height 2.3307in;depth
0pt;original-width 2.29in;original-height 2.29in;cropleft "0";croptop
"1";cropright "1";cropbottom "0";tempfilename
'SYOEUF01.wmf';tempfile-properties "XPR";}}

From the figure, rotation about the $y$ axis does not contribute
significantly because the moment of inertia of a sphere about it's axis
is(we will take the atoms to be roughly spherical) 
\begin{equation}
\mathbb{I}=\frac{2}{5}mr^{2}
\end{equation}%
where $m$ is the mass of the atom. Most of the mass is centered in the
nucleus (proton mass $=1.67\times 10^{-27}\unit{kg},$ electron mass $%
=9.11\times 10^{-31}\unit{kg})$, which has a radius of about $r=1.7\times
10^{-5}\unit{%
%TCIMACRO{\U{212b}}%
%BeginExpansion
\text{\AA}%
%EndExpansion
}.$ This is the effective radius of the atom as far as a moment of inertial
is concerned. For nitrogen $\left( 2.33\times 10^{-26}\unit{kg}\right) $ the
moment of inertial for rotation about the $y$-axis would be about 
\begin{eqnarray}
\mathbb{I} &=&\frac{2}{5}\left( 2.33\times 10^{-26}\unit{kg}\right) \left(
1.7\times 10^{-5}\unit{%
%TCIMACRO{\U{212b}}%
%BeginExpansion
\text{\AA}%
%EndExpansion
}\right) ^{2} \\
&=&\allowbreak 2.\,\allowbreak 693\,5\times 10^{-56}\unit{m}^{2}\unit{kg}
\end{eqnarray}%
which is very small. We would have twice this, because we have two nitrogen
atoms. Suppose we have an rotational speed of $\omega =2.\,\allowbreak
740\,4\times 10^{17}.$ This may seem high, but it's not hard to rotate
little diatomic molecules. The rotational energy would be 
\begin{eqnarray*}
K_{rot} &=&\frac{1}{2}\mathbb{I\omega }^{2} \\
&=&\frac{1}{2}\left( 2\times 2.\,\allowbreak 693\,5\times 10^{-56}\unit{m}%
^{2}\unit{kg}\right) \left( 2.\,\allowbreak 740\,4\times 10^{17}\right) ^{2}
\\
&=&2.\,\allowbreak 022\,8\times 10^{-21}\unit{J}
\end{eqnarray*}

We can also rotate the diatomic molecule about the $x$ or $z$-axes. The
moment of inertia for rotation about the center of the two mass system
(about the $x$ or $z$-axis) is about 
\begin{equation}
\mathbb{I}=\sum_{i}m_{i}R_{i}^{2}
\end{equation}%
where $R$ is the distance from the center of mass. For diatomic nitrogen, $R=%
\frac{1}{2}1.098\unit{%
%TCIMACRO{\U{212b}}%
%BeginExpansion
\text{\AA}%
%EndExpansion
}=\allowbreak 0.549\,\unit{%
%TCIMACRO{\U{212b}}%
%BeginExpansion
\text{\AA}%
%EndExpansion
}$ 
\begin{eqnarray}
\mathbb{I} &=&2\left( 4.65\times 10^{-26}\unit{kg}\right) \left( \allowbreak
0.549\,\unit{%
%TCIMACRO{\U{212b}}%
%BeginExpansion
\text{\AA}%
%EndExpansion
}\right) ^{2} \\
&=&2.\,\allowbreak 803\times 10^{-46}\unit{m}^{2}\unit{kg}
\end{eqnarray}%
This is ten $\left( 10\right) $ orders of magnitude larger than rotation
about the center $\left( y\right) $ axis of the molecule. We could have
energies like 
\begin{eqnarray*}
K_{rot} &=&\frac{1}{2}\mathbb{I\omega }^{2} \\
&=&\frac{1}{2}\left( 2.\,\allowbreak 803\times 10^{-46}\unit{m}^{2}\unit{kg}%
\right) \left( 2.\,\allowbreak 740\,4\times 10^{17}\frac{\unit{rad}}{\unit{s}%
}\right) ^{2} \\
&=&\allowbreak 1.\,\allowbreak 052\,5\times 10^{-11}\unit{J}
\end{eqnarray*}%
Again, we are about $10$ orders of magnitude larger than for rotation about
the $y$-axis. This seems like a place where we could put some energy. Let's
count this as a degree of freedom. And we would get this much energy for
both the $x$ and $z$ rotation so each would be a degree of freedom. But
let's ignore rotation about the $y$-axis.

We are left with three translational and two rotational degrees of freedom.
This gives%
\begin{equation}
E_{int}=\left( \frac{3}{2}k_{B}T\right) _{trans}+\left( \frac{2}{2}%
k_{B}T\right) _{rot}=\frac{5}{2}k_{B}T
\end{equation}%
Diatomic hydrogen really isn't an ideal gas. But we can get away with using
the ideal gas law for hydrogen sometimes by adding in these additional
degrees of freedom. Writing this in molar terms%
\begin{equation}
E_{int}=\frac{5}{2}nRT
\end{equation}%
then 
\begin{equation}
C_{%
%TCIMACRO{\TeXButton{V}{\ooalign{\hfil$V$\hfil\cr\kern0.1em--\hfil\cr}}}%
%BeginExpansion
\ooalign{\hfil$V$\hfil\cr\kern0.1em--\hfil\cr}%
%EndExpansion
}=\frac{1}{n}\frac{dE_{int}}{dT}=\frac{5}{2}R
\end{equation}

BUT WAIT, we did not include vibration! The atoms are bond together with an
electrical attraction that acts quite like a spring force. So vibration
along the axis is possible and we need to add in one more degree of freedom.
We also have potential energy involved for a spring force. That is an
additional place to put energy! So, we expect an additional degree of
freedom for vibration for the potential energy in vibration. \FRAME{dtbpF}{%
1.9908in}{1.9925in}{0pt}{}{}{Figure}{\special{language "Scientific
Word";type "GRAPHIC";maintain-aspect-ratio TRUE;display "USEDEF";valid_file
"T";width 1.9908in;height 1.9925in;depth 0pt;original-width
1.9519in;original-height 1.9527in;cropleft "0";croptop "1";cropright
"1";cropbottom "0";tempfilename 'SYOEVU02.wmf';tempfile-properties "XPR";}}%
When we add all these up, we get 
\begin{equation}
E_{int}=\left( \frac{3}{2}k_{B}T\right) _{trans}+\left( \frac{2}{2}%
k_{B}T\right) _{rot}+\left( \frac{2}{2}k_{B}T\right) _{vib}=\frac{7}{2}k_{B}T
\end{equation}%
which gives 
\begin{equation}
C_{%
%TCIMACRO{\TeXButton{V}{\ooalign{\hfil$V$\hfil\cr\kern0.1em--\hfil\cr}}}%
%BeginExpansion
\ooalign{\hfil$V$\hfil\cr\kern0.1em--\hfil\cr}%
%EndExpansion
}=\frac{7}{2}R
\end{equation}%
So if we include all of these degrees of freedom we might be able to use the
ideal gas approximation for diatomic hydrogen which would be convenient.

We should pause to ask, what values do we use? for diatomic gasses, is $%
C_{V}=\frac{7}{2}R$ all the time?

It turns out that when energy is added to a collection of molecules, it does
not pick randomly from the degrees of freedom. \FRAME{dhF}{3.2569in}{2.4422in%
}{0in}{}{}{Figure}{\special{language "Scientific Word";type
"GRAPHIC";maintain-aspect-ratio TRUE;display "USEDEF";valid_file "T";width
3.2569in;height 2.4422in;depth 0in;original-width 3.2119in;original-height
2.4007in;cropleft "0";croptop "1";cropright "1";cropbottom "0";tempfilename
'SW9P0O14.wmf';tempfile-properties "XPR";}}

At low temperature, the translational degrees of freedom are preferred.
Then, as the temperature rises, the rotational degrees of freedom are
filled. Finally the vibrational degrees of freedom are used. The figure
shows this relationship for diatomic Hydrogen. Note that there are plateaus
at each of our values of $C_{%
%TCIMACRO{\TeXButton{V}{\ooalign{\hfil$V$\hfil\cr\kern0.1em--\hfil\cr}}}%
%BeginExpansion
\ooalign{\hfil$V$\hfil\cr\kern0.1em--\hfil\cr}%
%EndExpansion
}$ that we found ($\frac{3}{2}R,\frac{5}{2}R,\frac{7}{2}R$).

We have talked about the simplest of molecules. If we had a more complex
molecule, there would be more complex rotational and vibrational degrees of
freedom. That is why for larger molecules the values of $C_{%
%TCIMACRO{\TeXButton{V}{\ooalign{\hfil$V$\hfil\cr\kern0.1em--\hfil\cr}}}%
%BeginExpansion
\ooalign{\hfil$V$\hfil\cr\kern0.1em--\hfil\cr}%
%EndExpansion
}$ in table in the book don't follow a simple fraction of $R.$

\section{Quantization}

We have come to the end of what classical theory can explain. We know that
different degrees of freedom are filled in groups, but we don't know why. We
will only hint at the reason, but it has to do with wave motion that we
studied!\FRAME{dhF}{3.7178in}{3.8787in}{0pt}{}{}{Figure}{\special{language
"Scientific Word";type "GRAPHIC";maintain-aspect-ratio TRUE;display
"USEDEF";valid_file "T";width 3.7178in;height 3.8787in;depth
0pt;original-width 3.6703in;original-height 3.8311in;cropleft "0";croptop
"1";cropright "1";cropbottom "0";tempfilename
'SW9P0O15.wmf';tempfile-properties "XPR";}}

Quantum theory tells us that atoms can be described as waves under boundary
conditions. Like a string with fixed ends, these atomic waves have quantized
frequencies. The energy of the molecule is proportional to the frequency in
quantum theory, so quantum mechanics tells us that the energy of the
molecule will be quantized. The figure shows an energy level diagram for a
molecule. Higher energy is in the $y$ direction. The vibrational and
rotational states are marked. (the maroon arrows show transitions from
higher to lower states. These transitions create radiation).

The lowest allowed state is called the \emph{ground state}. It is analogous
to the fundamental frequency. At low temperatures the molecule only has
enough energy to populate the lower states, but when the temperature rises,
there is enough energy transfer due to collisions to push molecules into the
rotational states. With higher temperatures, you will see the collisions
transfer molecules to the higher vibrational states. This gives rise to the
order of our $C_{V}$ values.

\section{Molar specific heat of solids}

Let's do a problem. Let's find $C_{%
%TCIMACRO{\TeXButton{V}{\ooalign{\hfil$V$\hfil\cr\kern0.1em--\hfil\cr}}}%
%BeginExpansion
\ooalign{\hfil$V$\hfil\cr\kern0.1em--\hfil\cr}%
%EndExpansion
}$ for a solid, and compare it to the specific heat of elemental solids.

We can view the atoms of solids as having a structure of springs and atoms.

\FRAME{dtbpF}{2.5719in}{2.4543in}{0in}{}{}{Figure}{\special{language
"Scientific Word";type "GRAPHIC";maintain-aspect-ratio TRUE;display
"USEDEF";valid_file "T";width 2.5719in;height 2.4543in;depth
0in;original-width 2.5304in;original-height 2.4128in;cropleft "0";croptop
"1";cropright "1";cropbottom "0";tempfilename
'SYOF3N05.wmf';tempfile-properties "XPR";}}We can see that we should have
three translational degrees of freedom for each atom. We also have three
degrees of freedom associated with the potential energy from the spring-like
bonding forces. Recall that%
\begin{equation}
E_{x}=\frac{1}{2}mv_{x}^{2}-\frac{1}{2}kx^{2}
\end{equation}%
and we have and $E_{y}$ and $E_{z}$ as well.

From equipartition of energy we have 
\begin{equation}
E_{i}=6\left( \frac{1}{2}k_{B}T\right)
\end{equation}%
per atom! So we have%
\begin{eqnarray}
E_{int} &=&3Nk_{B}T \\
&=&3nRT
\end{eqnarray}%
and then for solids 
\begin{equation}
C_{%
%TCIMACRO{\TeXButton{V}{\ooalign{\hfil$V$\hfil\cr\kern0.1em--\hfil\cr}}}%
%BeginExpansion
\ooalign{\hfil$V$\hfil\cr\kern0.1em--\hfil\cr}%
%EndExpansion
}=\frac{1}{n}\frac{dE_{int}}{dT}=3R
\end{equation}%
numerically this is 
\begin{eqnarray*}
C_{%
%TCIMACRO{\TeXButton{V}{\ooalign{\hfil$V$\hfil\cr\kern0.1em--\hfil\cr}}}%
%BeginExpansion
\ooalign{\hfil$V$\hfil\cr\kern0.1em--\hfil\cr}%
%EndExpansion
} &=&3\times 8.314\frac{\unit{J}}{\unit{mol}\unit{K}} \\
&=&24.\,\allowbreak 942\frac{\unit{J}}{\unit{mol}\unit{K}}
\end{eqnarray*}%
This is the \emph{DuLong-Petit law}. Here is our table of specific heat
values again.%
\begin{equation*}
\begin{tabular}{|l|l|l|}
\hline
\textbf{Substance} & $c\left( \frac{\unit{J}}{\unit{kg}\unit{K}}\right) $ & $%
C_{V}\left( \frac{\unit{J}}{\unit{mol}\unit{K}}\right) $ \\ \hline
Aluminum & 900 & 24.3 \\ \hline
Copper & 385 & 24.4 \\ \hline
Iron & 449 & 25.1 \\ \hline
Gold & 129 & 25.4 \\ \hline
Lead & 128 & 26.5 \\ \hline
Ice & 2090 & 37.6 \\ \hline
Mercury & 140 & 28.1 \\ \hline
Water & 4190 & 75.4 \\ \hline
\end{tabular}%
\end{equation*}%
we see that our simple analysis did not do too bad! For elemental solids we
get about the right number.

This law also breaks down for very cold temperatures where we need even more
quantum mechanics to explain what happens. Einstein and Debye both
contributed to explaining this departure, but their work is beyond this
freshman course. This is a topic for a solid state physics course which is
an elective that you might consider if you are interested is seeing a more
complete solution.\FRAME{dhFU}{2.5832in}{2.2174in}{0pt}{\Qcb{{\protect\small %
Debye and Einstein models. (Image in the Public Domain courtesy Fr\'{e}d\'{e}%
ric Perez)}}}{}{Figure}{\special{language "Scientific Word";type
"GRAPHIC";maintain-aspect-ratio TRUE;display "USEDEF";valid_file "T";width
2.5832in;height 2.2174in;depth 0pt;original-width 9.7222in;original-height
8.3333in;cropleft "0";croptop "1";cropright "1";cropbottom "0";tempfilename
'SW9P0O17.bmp';tempfile-properties "XPR";}}

\section{Partial Pressure and Vapor Pressure}

So far we have assumed we only have one type of molecule in our gas. But
what if we had several kinds of molecules, like in our air that we breath
that is mostly N$_{2}$ but also has some very important O$_{2}$ that we need
to live (and a little bit of CO$_{2}$ that we worry about)?

We could define a pressure that the gas would make if it occupied the entire
volume on it's own. We call this a \emph{partial pressure}. But think of our
kinetic theory of gasses that we have just developed. If the gasses are in
thermal equilibrium, then they have the same temperature, and therefore the
same average kinetic energy. They would each make a pressure on the walls of
our container. 
\begin{equation*}
P=\frac{2}{3}\frac{N}{%
%TCIMACRO{\TeXButton{V}{\ooalign{\hfil$V$\hfil\cr\kern0.1em--\hfil\cr}}}%
%BeginExpansion
\ooalign{\hfil$V$\hfil\cr\kern0.1em--\hfil\cr}%
%EndExpansion
}\left( \frac{1}{2}m\overline{v^{2}}\right)
\end{equation*}%
But each gas is doing this at the same time. So the total pressure on the
container walls would be the sum of the pressures from the individual gasses
in the mix%
\begin{equation*}
P_{total}=\sum_{i}P_{i}
\end{equation*}%
This is called Dalton's law. A result of this law is that for any two gasses
in equilibrium 
\begin{equation*}
\frac{P_{1}}{n_{1}}=\frac{P_{2}}{n_{2}}
\end{equation*}

There is a safety issue hidden in what we have just said. Suppose we have a
large container of liquid nitrogen in our lab (because we do). And suppose
our container has a leak. The pressure in the room might not change much at
all, but the portion of that pressure due to O$_{2}$ might change. And we
need that O$_{2}$ pressure to stay alive. But since the overall pressure
didn't change noticeably, we wouldn't have any warning that there was a
problem until we felt trouble moving or passed out.

Suppose we have a vapor in a container but also have some of the same
material in it's liquid state in the same container. And suppose that the
vapor and liquid are in thermal equilibrium. The partial pressure of the
vapor under these conditions is called a \emph{vapor pressure}. This is
important in many fields including weather prediction. The partial pressure
of the water vapor in our atmosphere can't exceed the vapor pressure of
water at the same temperature or the water vapor will condense out of the
air. This causes dew and the temperature at which the condensation starts to
happen is called the \emph{dew point temperature.}

If you listen to a weather report your might hear about the \emph{relative
humidity}.%
\begin{equation*}
R.H.=\frac{\text{Partial pressure of water vapor at }T}{\text{Vapor pressure
of water at }T}\times 100
\end{equation*}%
Think about what this means. If at a given temperature the partial pressure
of water vapor is the same as the vapor pressure of water, then the air has
as much water as it can hold at that temperature. We say the air is
saturated and it feels very muggy.

This almost completes our kinetic theory of ideal gasses. But before we move
on we need to look at the speeds of our molecules. We know average speed is
related to temperature, but not every molecule is moving at that average
speed. What distribution of speeds like?

\section{Speed of molecules}

In this course we have found that internal energy is mostly due to the
motion of the molecules in the system. That motion has kinetic energy
associated with it. But not a kinetic energy of the whole system moving
together. It is a kinetic energy of the parts of the system, the molecules.
Think of driving a car. We could describe the kinetic energy of the whole
car. That is the kinetic energy we learned about in PH121. But we could also
talk about the kinetic energy of one of the pistons in the engine. These two
kinetic energies would be different! Just like this, the kinetic energy of
the molecules will be different than the kinetic energy of the system as a
whole. \FRAME{dtbpF}{3.5453in}{2.3629in}{0pt}{}{}{Figure}{\special{language
"Scientific Word";type "GRAPHIC";maintain-aspect-ratio TRUE;display
"USEDEF";valid_file "T";width 3.5453in;height 2.3629in;depth
0pt;original-width 3.5843in;original-height 2.3789in;cropleft "0";croptop
"1";cropright "1";cropbottom "0";tempfilename
'SW9LUS08.wmf';tempfile-properties "XPR";}}To find this internal kinetic
energy of the molecules, we need to know how fast the molecules are going in
an ideal gas at a particular temperate and pressure in a particular volume.
Let's review a bit from our last lecture.

The internal kinetic energy depends on the speed of the molecules. But not
all molecules have the same velocity. So when we say \textquotedblleft speed
of the molecules\textquotedblright\ we have to be specific, what speed? The
average? The speed the most particles have?

The average velocity won't do. \FRAME{dtbpF}{3.0725in}{2.1988in}{0in}{}{}{%
Figure}{\special{language "Scientific Word";type
"GRAPHIC";maintain-aspect-ratio TRUE;display "USEDEF";valid_file "T";width
3.0725in;height 2.1988in;depth 0in;original-width 3.1027in;original-height
2.2121in;cropleft "0";croptop "1";cropright "1";cropbottom "0";tempfilename
'SW9LUS09.wmf';tempfile-properties "XPR";}}

The molecule motions have random directions. So the average velocity is
zero. 
\begin{equation*}
\overline{\overrightarrow{v}}=0.
\end{equation*}%
Some velocities are in negative directions and some are in positive
directions, so the velocities cancel out. But if we square the velocities
they are all positive. Then we could find the average of the squared
velocities and that would be better. And we recognize this! The quantity $%
\overline{v^{2}}$ showed up in our kinetic model for pressure and
temperature. It is part of the expression for the average kinetic energy of
the molecules. But a velocity squared is not a velocity. So to get back to
velocity units, let's take a square root.

\begin{equation*}
v_{rms}=\sqrt{\overline{v^{2}}}
\end{equation*}%
This quantity is like an average, and gives a representative value that is
near the most probable speeds. We have seen this quantity before, it is the $%
rms$ speed.

Using this, we can get an expression that relates the temperature of our gas
to the speed of the molecules. But it we need to know something about how
many molecules have what speed. To find the \emph{distribution} of speeds,
let's remember the idea of number density%
\begin{equation*}
n_{V}=\frac{\#\text{ of molecules }}{%
%TCIMACRO{\TeXButton{V}{\ooalign{\hfil$V$\hfil\cr\kern0.1em--\hfil\cr}}}%
%BeginExpansion
\ooalign{\hfil$V$\hfil\cr\kern0.1em--\hfil\cr}%
%EndExpansion
}
\end{equation*}%
We can express the number of molecules this way assuming we know the volume.
But we want the speed of the molecules, and we have learned that the concept
of energy usually can get us a speed if we don't care to know the direction.
Knowing the direction for every air molecule in a room would not be helpful
in most cases, so we can be content with just knowing the speed. Lets find a
sort of \textquotedblleft energy density,\textquotedblright\ that is, the
density of molecules that have energy between two amounts of energy, say, $%
E_{1}$ and $E_{2}$. We want $E_{1}$ and $E_{2}$ to be quite close together.
So let's let $E_{1}=$ $E$ and $E_{2}=E+\Delta E$ where $\Delta E$ is a small
amount of energy. If you take quantum mechanics, or statistical mechanics,
you will likely see this quantity quite a lot. Then the number of molecules
with a particular energy between $E_{1}$ and $E_{2}$ could be written as%
\begin{equation}
n_{V}\left( E\right) dE=\frac{\#\text{ of molecules with energy between }E%
\text{ and }E+\Delta E}{%
%TCIMACRO{\TeXButton{V}{\ooalign{\hfil$V$\hfil\cr\kern0.1em--\hfil\cr}}}%
%BeginExpansion
\ooalign{\hfil$V$\hfil\cr\kern0.1em--\hfil\cr}%
%EndExpansion
}
\end{equation}

This is called a \emph{distribution function}. We find distribution
functions in statistics. They are associated with probabilities. The
standard \textquotedblleft bell curve\textquotedblright\ used sometimes in
grading is a distribution function. It tells the total number of students
that got a particular number of points in a class. Think of grades on a
test. If we plot the grades, the height of the bar for a \textquotedblleft
B\textquotedblright\ grade is the number of students that scored between
83\% and 87\%. \FRAME{dtbpF}{3.7161in}{2.1612in}{0pt}{}{}{Figure}{\special%
{language "Scientific Word";type "GRAPHIC";maintain-aspect-ratio
TRUE;display "USEDEF";valid_file "T";width 3.7161in;height 2.1612in;depth
0pt;original-width 3.6685in;original-height 2.1223in;cropleft "0";croptop
"1";cropright "1";cropbottom "0";tempfilename
'SYQJ6G00.wmf';tempfile-properties "XPR";}}This graph is a distribution of
grades and it tells us for each score segment, say $S$ and $S+\Delta S$
(think $83\%$ and $83\%+4\%=87\%$) how many students had scores in that
range. We are doing the same for our molecules and their energy. We will
make a graph that shows for each energy segment from $E$ to $E+\Delta E$ how
many molecules have that energy. This gives us the probability that the
molecules will have a particular energy (or speed, since this is kinetic
energy). A function that gives the amount of molecules that have a
particular amount of energy also is called a \emph{distribution function}
just like it was for grades. The distribution function that we will use is
written symbolically in this cryptic fashion, $n_{V}\left( E\right) dE.$ It
is the number of molecules with a particular energy (because our $\Delta E$
is so small it is a $dE$) divided by the total number of molecules. That
makes it like a percent.

We won't derive an equation for this quantity, instead we will borrow a
result from our junior level thermal physics class.%
\begin{equation}
n_{V}\left( E\right) =n_{E}e^{-\frac{E}{k_{B}T}}
\end{equation}%
where $n_{E}dE$ is the number of molecules per unit volume having energy
between $E=0$ and $E=dE.$ This distribution function is called the \emph{%
Boltzmann distribution law.} It tells us that the probability of finding the
molecules in a particular energy state varies exponentially as the negative
of the energy divided by $k_{B}T.$

\section{Distribution of Molecular Speeds}

Since there is a distribution of energies, we expect our gas molecules to
have a distribution of velocities. That is the same as saying the molecules
in the gas do not all go the same speed. Again we won't derive this (at
least not in this class!), but the distribution should depend on
temperature, $T,$ since we know the internal energy is tied to temperature.
The distribution is as follows:%
\begin{equation}
N_{v}\left( v\right) =4\pi N\left( \frac{m}{2\pi k_{B}T}\right) ^{\frac{3}{2}%
}v^{2}e^{-\frac{mv^{2}}{2k_{B}T}}
\end{equation}%
where $m$ is the mass of the molecule. The kinetic energy of the molecules
is hiding in the exponent, so we could write this as 
\begin{equation}
N_{v}=4\pi N\left( \frac{m}{2\pi k_{B}T}\right) ^{\frac{3}{2}}v^{2}e^{-\frac{%
K}{k_{B}T}}
\end{equation}%
Note that this is of the form for our distribution function because for an
ideal gas the only internal energy we have is kinetic energy, $E=K$ so we
have $E/(k_{B}T)$ as the exponent. We used this to turn our energy
distribution function into a speed distribution function.

If there are $dN$ molecules with speeds between $v$ and $v+dv$ then%
\begin{equation}
dN=N_{v}dv
\end{equation}%
so there should be 
\begin{equation}
N=\int_{0}^{\infty }N_{v}dv
\end{equation}%
total molecules. This tells us that the $4\pi N\left( \frac{m}{2\pi k_{B}T}%
\right) ^{\frac{3}{2}}$ is a \emph{normalization constant} that makes sure
our integral of $N_{v}dv$ gives us $N$, the total number of particles.

If we plot $N_{v}$ vs. $v$ we get the figure below. We get a plot that is a
little like our grade distribution plot, only for molecular speeds. The
number of molecules with speeds between $v$ and $v+dv$ is the area under the
blue curve at $v.$ It would be the area of the thin red rectangle right at $%
v $ with width $dv$ shown in the next figure. \FRAME{dtbpF}{2.1058in}{%
2.0885in}{0pt}{}{}{Figure}{\special{language "Scientific Word";type
"GRAPHIC";maintain-aspect-ratio TRUE;display "USEDEF";valid_file "T";width
2.1058in;height 2.0885in;depth 0pt;original-width 2.0669in;original-height
2.0505in;cropleft "0";croptop "1";cropright "1";cropbottom "0";tempfilename
'SYQK9K02.wmf';tempfile-properties "XPR";}}The peak of the curve tells us
the most probable speed, that is, the speed the most molecules have, $v_{mp}$
(like in our grade example, the most students got an \textquotedblleft
A\textquotedblright\ because that bar was the tallest). The curve is not
symmetric, so the most probable speed is not the average speed, $\bar{v}$.
There is also our new speed estimate marked $v_{rms}$.\FRAME{dtbpF}{3.3998in%
}{3.3546in}{0pt}{}{}{Figure}{\special{language "Scientific Word";type
"GRAPHIC";maintain-aspect-ratio TRUE;display "USEDEF";valid_file "T";width
3.3998in;height 3.3546in;depth 0pt;original-width 3.4353in;original-height
3.3892in;cropleft "0";croptop "1";cropright "1";cropbottom "0";tempfilename
'SW9LUS0A.wmf';tempfile-properties "XPR";}}

If we plot $N_{v}$ for different temperatures, we observe that the peak
shifts, and the curve broadens\FRAME{dtbpFU}{3.3243in}{2.4491in}{0pt}{\Qcb{%
{\protect\small Temperature dependence of the Maxwell-Boltsman distribution
(Image in the Public Domain couracy Fred Stober)}}}{}{Figure}{\special%
{language "Scientific Word";type "GRAPHIC";maintain-aspect-ratio
TRUE;display "USEDEF";valid_file "T";width 3.3243in;height 2.4491in;depth
0pt;original-width 3.359in;original-height 2.4667in;cropleft "0";croptop
"1";cropright "1";cropbottom "0";tempfilename
'SW9LUS0B.wmf';tempfile-properties "XPR";}}A motivated student could now
find the most probably speed by finding the maximum of $N_{v}.$ We know from
our calculus experience how to find a maximum. We take a derivative

\begin{eqnarray*}
\frac{dN_{v}}{dv} &=&\frac{d}{dv}\left( 4\pi N\left( \frac{m}{2\pi k_{B}T}%
\right) ^{\frac{3}{2}}v^{2}e^{-\frac{mv^{2}}{2k_{B}T}}\right) \\
&=&4\pi N\left( \frac{m}{2\pi k_{B}T}\right) ^{\frac{3}{2}}\frac{d}{dv}%
\left( v^{2}e^{-\frac{m}{2k_{B}T}v^{2}}\right) \\
&=&4\pi N\left( \frac{m}{2\pi k_{B}T}\right) ^{\frac{3}{2}}\frac{d}{dv}%
\left( v^{2}\frac{d}{dv}e^{-\frac{m}{2k_{B}T}v^{2}}+e^{-\frac{m}{2k_{B}T}%
v^{2}}\frac{d}{dv}v^{2}\right) \\
&=&4\pi N\left( \frac{m}{2\pi k_{B}T}\right) ^{\frac{3}{2}}\left( -2\frac{m}{%
2k_{B}T}v^{3}e^{-\frac{m}{2k_{B}T}v^{2}}+2ve^{-\frac{m}{2k_{B}T}v^{2}}\right)
\\
&=&-2ve^{-\frac{mv^{2}}{2k_{B}T}}\left( 4\pi N\left( \frac{m}{2\pi k_{B}T}%
\right) ^{\frac{3}{2}}\right) \left( \frac{mv^{2}}{2k_{B}T}-1\right)
\end{eqnarray*}

and set it equal to zero. I didn't spend a lot of time simplifying, because
all the factors out front will go away when we set it equal to zero: $\frac{d%
}{dt}e^{at}=\allowbreak ae^{at}$%
\begin{eqnarray*}
0 &=&-2ve^{-\frac{mv^{2}}{2k_{B}T}}\left( 4\pi N\left( \frac{m}{2\pi k_{B}T}%
\right) ^{\frac{3}{2}}\right) \left( \frac{mv^{2}}{2k_{B}T}-1\right) \\
\frac{mv^{2}}{2k_{B}T} &=&1 \\
v_{mp} &=&\sqrt{\frac{2k_{B}T}{m}}
\end{eqnarray*}%
\begin{equation}
v_{mp}=\sqrt{\frac{2k_{B}T}{m}}
\end{equation}%
and this is great! We have related the temperature, $T$, to the most
probable speed of the molecules. But it is more convenient to use $v_{rms,}$
so let's see if we can modify this expression to be in terms of $v_{rms}.$
And it gives us some insight into using a distribution function. To do this
let's once again review averages, but this time using our distribution
function.

The average value of $v^{n}$ is given by%
\begin{equation*}
\overline{v^{n}}=\frac{1}{N}\int_{0}^{\infty }v^{n}N_{v}dv
\end{equation*}%
Think of this as summing up $v^{n}$ but using the fact that we already know
that for each $v$ there are $N_{v}dv$ molecules with speed $v$. So our
distribution function makes the integral easier. That motivated student
could also use this to find the average speed (not the average velocity,
which is zero). He or she will want the average value of $v^{1}$ 
\begin{eqnarray*}
\overline{v^{1}} &=&\frac{1}{N}\int_{0}^{\infty }v^{1}4\pi N\left( \frac{m}{%
2\pi k_{B}T}\right) ^{\frac{3}{2}}v^{2}e^{-\frac{mv^{2}}{2k_{B}T}}dv \\
&=&4\pi \left( \frac{m}{2\pi k_{B}T}\right) ^{\frac{3}{2}}\int_{0}^{\infty
}v^{3}e^{-\frac{mv^{2}}{2k_{B}T}}dv
\end{eqnarray*}%
You might be thinking that this integral looks a little hard to do. You
would be right. You could use a symbolic math program to do this. Or you can
dust off your second semester calculus knowledge. I prefer the symbolic math
program, but let's do it the calc II way here. Looking just at the integral

\begin{equation*}
\int_{0}^{\infty }v^{3}e^{-\frac{mv^{2}}{2k_{B}T}}dv=\int_{0}^{\infty
}v^{2}e^{-\frac{mv^{2}}{2k_{B}T}}vdv
\end{equation*}

Let's do a substitution. let $x=v^{2},$ $dx=2vdv$ and $a=-\frac{m}{2k_{B}T}$
then we have 
\begin{equation*}
\int_{0}^{\infty }v^{3}e^{-\frac{mv^{2}}{2k_{B}T}}dv=\int_{0}^{\infty
}xe^{ax}dx
\end{equation*}%
and now we have one of those many integration identities that you memorized
(and now have forgotten) from Calc II. So we look in an integral table and
find

\begin{equation*}
\int_{0}^{\infty }xe^{ax}dx=\frac{e^{ax}}{a^{2}}\left( ax-1\right)
\end{equation*}%
This matches our integral that we need. We now know the answer. We just have
to undo our substitution.

\begin{eqnarray*}
\int_{0}^{\infty }v^{3}e^{-\frac{mv^{2}}{2k_{B}T}}dv &=&\left. \frac{e^{%
\frac{-m}{2k_{B}T}v^{2}}}{\left( \frac{m}{2k_{B}T}\right) ^{2}}\left( \left(
-\frac{m}{2k_{B}T}\right) v^{2}-1\right) \right\vert _{0}^{\infty } \\
&=&\left[ \frac{e^{\frac{-m}{2k_{B}T}\left( \infty \right) ^{2}}}{\left( 
\frac{m}{2k_{B}T}\right) ^{2}}\left( \left( -\frac{m}{2k_{B}T}\right) \left(
\infty \right) ^{2}-1\right) \right] -\left[ \frac{e^{\frac{-m}{2k_{B}T}%
\left( 0\right) ^{2}}}{\left( \frac{m}{2k_{B}T}\right) ^{2}}\left( \left( -%
\frac{m}{2k_{B}T}\right) \left( 0\right) ^{2}-1\right) \right] \\
&=&0-\left[ \frac{1}{\left( \frac{m}{2k_{B}T}\right) ^{2}}\left( \left( -%
\frac{m}{2k_{B}T}\right) \left( 0\right) ^{2}-1\right) \right] \\
&=&\frac{1}{\left( \frac{m}{2k_{B}T}\right) ^{2}}
\end{eqnarray*}

And we can put our result back in our equation for $\overline{v^{1}}$

\begin{eqnarray*}
\overline{v^{1}} &=&4\pi \left( \frac{m}{2\pi k_{B}T}\right) ^{\frac{3}{2}%
}\int_{0}^{\infty }v^{3}e^{-\frac{mv^{2}}{2k_{B}T}}dv \\
&=&2\pi \left( \frac{m}{2\pi k_{B}T}\right) ^{\frac{3}{2}}\left( \frac{1}{%
\left( \frac{\pi m}{2\pi k_{B}T}\right) ^{2}}\right) \\
&=&2\pi \left( \frac{m}{2\pi k_{B}T}\right) ^{\frac{3}{2}}\left( \left( 
\frac{m}{2\pi k_{B}T}\right) ^{-\frac{4}{2}}\pi ^{-2}\right) \\
&=&2\pi \left( \frac{m}{2\pi k_{B}T}\right) ^{-\frac{1}{2}}\left( \pi
^{-2}\right) \\
&=&2\left( \frac{2\pi k_{B}T}{m}\right) ^{\frac{1}{2}}\left( \pi ^{-1}\right)
\\
&=&\sqrt{\frac{8k_{B}T}{\pi m}}
\end{eqnarray*}%
so

\begin{equation*}
\bar{v}=\sqrt{\frac{8k_{B}T}{\pi m}}
\end{equation*}%
which is also great, but not what we wanted. But it is close. This time
let's find the average value of $v^{2}.$ We hinted that the root mean
squared value would be useful. that is, $\overline{v^{2}}=$ $v_{rms}$. We
can find this like we found the average velocity

\begin{eqnarray*}
\overline{v^{2}} &=&\frac{1}{N}\int_{0}^{\infty }v^{2}4\pi N\left( \frac{m}{%
2\pi k_{B}T}\right) ^{\frac{3}{2}}v^{2}e^{-\frac{mv^{2}}{2k_{B}T}}dv \\
&=&4\pi \left( \frac{m}{2\pi k_{B}T}\right) ^{\frac{3}{2}}\int_{0}^{\infty
}v^{4}e^{-\frac{mv^{2}}{2k_{B}T}}dv
\end{eqnarray*}%
But the math is a bit more cumbersome. Let%
\begin{equation*}
a=\frac{m}{2k_{B}T}
\end{equation*}%
Then we have 
\begin{equation*}
\overline{v^{2}}=4\pi \left( \frac{m}{2\pi k_{B}T}\right) ^{\frac{3}{2}%
}\int_{0}^{\infty }v^{4}e^{-av^{2}}dv
\end{equation*}%
It is time to look up our integral in a table of integrals. My table tells
me that 
\begin{eqnarray*}
\int x^{4}e^{-ax^{2}}dx &=&-\left. \frac{1}{8a^{\frac{5}{2}}}\left( 6\sqrt{a}%
xe^{-ax^{2}}-3\sqrt{\pi }\func{erf}\left( \sqrt{a}x\right) +4a^{\frac{3}{2}%
}x^{3}e^{-ax^{2}}\right) \allowbreak \right\vert _{0}^{\infty } \\
&=&-\frac{1}{8a^{\frac{5}{2}}}\left( -3\sqrt{\pi }\right)
\end{eqnarray*}

The quantity $\func{erf}\left( \sqrt{a}x\right) $ is called the
\textquotedblleft error function.\textquotedblright\ If you study this
function in a good mathematical book on integration you will find that 
\begin{equation*}
\func{erf}\left( \infty \right) =\frac{2}{\sqrt{\pi }}\int_{0}^{\infty
}e^{-t^{2}}dt=\allowbreak 1
\end{equation*}

Let's apply all this to our $\overline{v^{2}}$ integral. Let's let $x=v$ and 
$a=\frac{m}{2k_{B}T}$ just as we said before. Then%
\begin{eqnarray*}
\overline{v^{2}} &=&4\pi \left( \frac{m}{2\pi k_{B}T}\right) ^{\frac{3}{2}%
}\int_{0}^{\infty }x^{4}e^{-ax^{2}}dx \\
&=&4\pi \left( \frac{m}{2\pi k_{B}T}\right) ^{\frac{3}{2}}\left( -\frac{1}{%
8a^{\frac{5}{2}}}\left( -3\sqrt{\pi }\right) \right)
\end{eqnarray*}%
and substituting back in for $x$ and $a$ 
\begin{equation}
\overline{v^{2}}=4\pi \left( \frac{m}{2\pi k_{B}T}\right) ^{\frac{3}{2}%
}\left( -\frac{1}{8\left( \frac{m}{2k_{B}T}\right) ^{\frac{5}{2}}}\left( -3%
\sqrt{\pi }\right) \right)  \notag
\end{equation}%
Some things cancel%
\begin{eqnarray}
\overline{v^{2}} &=&\frac{4\pi }{\left( 2\pi \right) ^{\frac{3}{2}}}\left( 
\frac{m}{k_{B}T}\right) ^{\frac{3}{2}}\left( \frac{1}{8\frac{1}{2^{\frac{5}{2%
}}}\left( \frac{m}{k_{B}T}\right) ^{\frac{5}{2}}}\left( 3\sqrt{\pi }\right)
\right)  \notag \\
&=&\frac{4\pi }{\left( 2\pi \right) ^{\frac{3}{2}}}\left( \frac{1}{8\frac{1}{%
2^{\frac{5}{2}}}\left( \frac{m}{k_{B}T}\right) ^{\frac{2}{2}}}\left( 3\left(
\pi \right) ^{\frac{1}{2}}\right) \right) \\
&=&\frac{4\pi ^{\frac{2}{2}}}{2^{\frac{3}{2}}\pi ^{\frac{3}{2}}}\left( \frac{%
1}{8\frac{1}{2^{\frac{5}{2}}}\left( \frac{m}{k_{B}T}\right) ^{\frac{2}{2}}}%
\left( 3\left( \pi \right) ^{\frac{1}{2}}\right) \right) \\
&=&4\left( \frac{3}{8\frac{2^{\frac{3}{2}}}{2^{\frac{5}{2}}}\left( \frac{m}{%
k_{B}T}\right) }\right) \\
&=&\frac{4}{8}\frac{3}{\frac{1}{2}\left( \frac{m}{k_{B}T}\right) } \\
&=&\frac{3k_{B}T}{m}
\end{eqnarray}%
This is $\overline{v^{2}}.$ But we want $v_{rms}=\sqrt{\overline{v^{2}}}$ so%
\begin{equation}
v_{rms}=\sqrt{\frac{3k_{B}T}{m}}
\end{equation}%
We have an expression that relates the temperature of the gas to the $rms$
speed of the molecules of the gas. And you might object that we have already
seen this form for $v_{rms}$ and it was much easier to find last time! But
it is worth seeing that our distribution of molecular speeds produces this
result. It gives us confidence in our distribution function. Notice that,
indeed, $v_{rms}$ is close to the average and most probable speeds.

That we have some speeds higher than others explains why things do not
evaporate or all boil away at once.%
\begin{equation*}
\begin{tabular}{|c|c|c|}
\hline
& \textbf{A few rms Speeds for Gasses} &  \\ \hline
Gas & Molar Mass ($\unit{kg}/\unit{mol}$) & $v_{rms}$ at $20\unit{%
%TCIMACRO{\U{2103}}%
%BeginExpansion
{}^{\circ}{\rm C}%
%EndExpansion
}$ $\left( \unit{m}/\unit{s}\right) $ \\ \hline
H$_{2}$ & $2.02\times 10^{-3}$ & $1902$ \\ \hline
He & $4.0\times 10^{-3}$ & $1352$ \\ \hline
H$_{2}$O & $18\times 10^{-3}$ & $637$ \\ \hline
N$_{2}$ & $28\times 10^{-3}$ & $511$ \\ \hline
O$_{2}$ & $28\times 10^{-3}$ & $478$ \\ \hline
CO$_{2}$ & $44\times 10^{-3}$ & $408$ \\ \hline
\end{tabular}%
\end{equation*}

Armed with a molecular model for gasses, let's return to our state variables
for a system and see how to find the state variables for different
thermodynamic processes. Let's start this by looking at some special ways to
change thermodynamic systems.

\chapter{32 First Law and Work 2.3.2, 2.3.3}

%TCIMACRO{\TeXButton{\chaptermark{Chapter 34}}{\chaptermark{Chapter 34}}}%
%BeginExpansion
\chaptermark{Chapter 34}%
%EndExpansion

We have studied the zeroth law of thermodynamics. Zero is an odd place to
start numbering (unless you are programing a computer in python or C++).
That no energy is transferred between to objects in thermodynamic
equilibrium is so obvious that for a while it went unstated. But now we
recognize from our kinetic theory of gasses why this must be the case and we
can label this fact a law of thermodynamics. But there is more than one law
of thermodynamics (other wise, why would we number them at all?). Let's
learn another law of thermodynamics, which has to do with conservation of
energy in thermodynamic systems.

%TCIMACRO{%
%\TeXButton{Fundamental Concepts}{\vspace{0.10in}\hspace{-0.37in}{\Large Fundamental Concepts\vspace{0.2in}}}}%
%BeginExpansion
\vspace{0.10in}\hspace{-0.37in}{\Large Fundamental Concepts\vspace{0.2in}}%
%EndExpansion

\begin{itemize}
\item Thermodynamic systems are just the part where we want to know state
it's variables $P,$ $%
%TCIMACRO{\TeXButton{V}{\ooalign{\hfil$V$\hfil\cr\kern0.1em--\hfil\cr}}}%
%BeginExpansion
\ooalign{\hfil$V$\hfil\cr\kern0.1em--\hfil\cr}%
%EndExpansion
,$ $n,$ and $T.$ All the material and space around the system is called the
surroundings.

\item We think of quasi-static processes in thermodynamics and would prefer
reversible processes.

\item The first law of Thermodynamics is just conservation of energy

\item $\Delta E_{int}$ is an invariant quantity in thermodynamics

\item Work in thermodynamic processes can be described as the integral $%
w_{int}=-\int_{%
%TCIMACRO{\TeXButton{V}{\ooalign{\hfil$V$\hfil\cr\kern0.1em--\hfil\cr}}}%
%BeginExpansion
\ooalign{\hfil$V$\hfil\cr\kern0.1em--\hfil\cr}%
%EndExpansion
_{i}}^{%
%TCIMACRO{\TeXButton{V}{\ooalign{\hfil$V$\hfil\cr\kern0.1em--\hfil\cr}}}%
%BeginExpansion
\ooalign{\hfil$V$\hfil\cr\kern0.1em--\hfil\cr}%
%EndExpansion
_{f}}Pd%
%TCIMACRO{\TeXButton{V}{\ooalign{\hfil$V$\hfil\cr\kern0.1em--\hfil\cr}}}%
%BeginExpansion
\ooalign{\hfil$V$\hfil\cr\kern0.1em--\hfil\cr}%
%EndExpansion
$

\item $\Delta E_{int}=Q+w_{int}$ is the first law of thermodynamics
\end{itemize}

\section{Thermodynamic systems}

In the next few lectures we are going to talk a lot about thermodynamic
systems. But what is a thermodynamic system? It is just the group of
material that we care about, that we want to know the state variables of in
our problem or experiment. In our calorimetry problems we had an insolated
container. The stuff in the container was our thermodynamic system and the
container and air and table and room and everything else around the
container is what we call the surroundings.\FRAME{dtbpF}{3.7481in}{2.3955in}{%
0pt}{}{}{Figure}{\special{language "Scientific Word";type
"GRAPHIC";maintain-aspect-ratio TRUE;display "USEDEF";valid_file "T";width
3.7481in;height 2.3955in;depth 0pt;original-width 3.7005in;original-height
2.3557in;cropleft "0";croptop "1";cropright "1";cropbottom "0";tempfilename
'SYQLXC03.wmf';tempfile-properties "XPR";}}We could draw a more general
picture of a thermodynamic system \FRAME{dtbpF}{3.1531in}{2.0557in}{0pt}{}{}{%
Figure}{\special{language "Scientific Word";type
"GRAPHIC";maintain-aspect-ratio TRUE;display "USEDEF";valid_file "T";width
3.1531in;height 2.0557in;depth 0pt;original-width 3.1081in;original-height
2.0167in;cropleft "0";croptop "1";cropright "1";cropbottom "0";tempfilename
'SYQM5K04.wmf';tempfile-properties "XPR";}}

In our first system example we had an insulated container. That insulation
prevents energy transfer by heat, and if it is ridged it prevents energy
transfer by work. But this is a special case. We call such a system a \emph{%
closed system}. But most thermodynamic systems do exchange energy with their
surroundings. We call them \emph{open systems}.

In our problems, we generally think of the material inside of our system as
being in thermodynamic equilibrium. This is not always the case.

Consider the situation shown below\FRAME{dhF}{1.6215in}{2.2745in}{0pt}{}{}{%
Figure}{\special{language "Scientific Word";type
"GRAPHIC";maintain-aspect-ratio TRUE;display "USEDEF";valid_file "T";width
1.6215in;height 2.2745in;depth 0pt;original-width 1.5852in;original-height
2.2347in;cropleft "0";croptop "1";cropright "1";cropbottom "0";tempfilename
'SW9P0O18.wmf';tempfile-properties "XPR";}}We have a gas in one side of a
sealed container and a vacuum on the other side. A membrane separates the
two. We have a very definite temperature in the gas, and also a definite
temperature in the vacuum. But if I puncture the membrane, \FRAME{dhF}{%
1.7755in}{2.4976in}{0pt}{}{}{Figure}{\special{language "Scientific
Word";type "GRAPHIC";maintain-aspect-ratio TRUE;display "USEDEF";valid_file
"T";width 1.7755in;height 2.4976in;depth 0pt;original-width
1.7383in;original-height 2.4569in;cropleft "0";croptop "1";cropright
"1";cropbottom "0";tempfilename 'SW9P0O19.wmf';tempfile-properties "XPR";}}%
gas goes swirling into the vacuum. The density will not be constant as the
gas swirls. Where there is gas, we would say we have a temperature closer to 
$T_{i}.$ But where there is no gas, we have a temperature nearer $0\unit{K}.$
It takes a while for the gas to come to equilibrium. During this process we
can't say that we have a definite temperature, or definite pressure for the
system. Higher physics classes (and engineering classes) deal with the
difficult situation that exists before equilibrium is achieved. For now, we
will only look at situations that have reached an equilibrium, or are not
too far from equilibrium. We call states that change slowly so they are
never far from equilibrium, \emph{quazi-static}. The situation we have just
described is called a \emph{free expansion} because the gas is allowed to
expand with no outside influences. A free expansion is \emph{not} a
quasi-static process. Let's think of how we might make a quasi-static
process happen.

\section{Reversible and Irreversible Processes}

A reversible process is one in which every point along some path is an
equilibrium state and one for which the system can be returned to its
initial state along the same path.

All natural processes are known to be irreversible. Reversible processes are
an idealization, like frictionless surfaces, or massless strings. But some
real processes are good approximations to reversible processes.

Our free expansion, of course, is not a reversible process. Imagine trying
to get all the gas atoms back through the hole in the membrane!

We could try to put the gas back into a smaller volume by putting a piston
in the system and applying a force. We would be transferring energy by work.
This would increase the temperature, so we would have to let energy transfer
out of the system by heat to keep the temperature the same. This would allow
the system to return to it's original state, but the surroundings will have
changed. We will have heated up the surroundings. This is not a reversible
process!

But suppose we apply our force on our piston in a quasi-static way. \FRAME{%
dtbpF}{1.3569in}{1.3742in}{0pt}{}{}{Figure}{\special{language "Scientific
Word";type "GRAPHIC";maintain-aspect-ratio TRUE;display "USEDEF";valid_file
"T";width 1.3569in;height 1.3742in;depth 0pt;original-width
4.3059in;original-height 4.3595in;cropleft "0";croptop "1";cropright
"1";cropbottom "0";tempfilename 'SW9P0O1A.wmf';tempfile-properties "XPR";}}%
Imagine dropping sand on a frictionless piston to compress a gas one grain
at a time. Then slowly removing the sand one grain at a time. The system
stays almost in equilibrium as it changes, so that it returns to a new
equilibrium right away. After all, you only moved one grain of sand at a
time. This process is almost reversible. Of course, you still have to do
work in moving the sand, so the outside environment is changed. Even this
process is not truly reversible. But it is a good mental model for the sort
of processes we will study in this class; slow, almost in equilibrium all
the time.

For our quasi-static process we can say we stay in thermodynamic
equilibrium. Recall that we have one \textquotedblleft
law\textquotedblright\ of thermodynamics. The zeroth law. It tells us that
if two objects are in thermodynamic equilibrium they have the same
temperature. From what we have done in this and the last lecture we now see
this zeroth law is telling us about forces and collisions between molecules.
When the molecules of the two objects have the same average kinetic energy,
no net energy is transferred between them and this is thermodynamic
equilibrium.

\section{Conservation of Energy: First Law of Thermodynamics}

Remember when we studied energy in Principles of Physics I? We used the work
energy theorem%
\begin{equation*}
\Delta K=w
\end{equation*}%
and expanded the work side to show the different types of work explicitly.
We have work done by conservative forces $w_{c}$ and work done by
non-conservative forces (dissipative forces like friction) $w_{nc}$.%
\begin{equation*}
\Delta K=w_{c}+w_{nc}
\end{equation*}%
and we could further divide up the work by categorizing the source of the
work, say, work done by the force of gravity, $w_{g},$ or work done by a
spring force $w_{s}.$%
\begin{equation*}
\Delta K=w_{g}+w_{s}+w_{nc}
\end{equation*}%
and we could go on in more complicated systems listing all the different
kinds of work.

We relabeled these different types of work as potential energies. For
example,%
\begin{eqnarray*}
\Delta U_{g} &=&-w_{g} \\
\Delta U_{s} &=&-w_{s}
\end{eqnarray*}%
and rewrote the work energy theorem%
\begin{equation*}
\Delta K+\Delta U_{g}+\Delta U_{s}=w_{nc}
\end{equation*}%
This gave us a way to think of mechanical energy in a system. But now we
realize that we have not covered all types of energy yet. Suppose we
consider your car. You put gas in the car. This is a source of energy, \emph{%
a chemical potential energy}, $U_{chem}$. So we should really include it.%
\begin{equation*}
\Delta K+\Delta U_{g}+\Delta U_{s}+\Delta U_{chem}=w_{nc}
\end{equation*}%
And if your car is a DeLorean that has been modified to use nuclear energy
to run, then you might have a nuclear potential energy\footnote{%
This is a reference to the car in the \emph{Back to the Future} movie
series. Sadly, it isn't something that actually exists.}%
\begin{equation*}
\Delta K+\Delta U_{g}+\Delta U_{s}+\Delta U_{chem}+\Delta U_{nuclear}=w_{nc}
\end{equation*}%
and back in PH121 we called the non-conservative work $w_{nc}=\Delta E_{th}$
because friction-like forces dissipate mechanical energy by turning them
into thermal energy. Then 
\begin{equation*}
\Delta K+\Delta U_{g}+\Delta U_{s}+\Delta U_{chem}+\Delta U_{nuclear}=\Delta
E_{th}
\end{equation*}

The quantities $\Delta U_{chem},$ $\Delta U_{nuclear}$, and $\Delta E_{th}$
are different than the mechanical energies we studied before. They have to
do with the internal workings of the material, itself. We will call an
amount of energy that comes from the material, itself \emph{internal energy.}
And from what we have done so far, we might expect that temperature has
something to do with an internal energy of the object. And we would be right.

\subsection{Work and Internal energy}

Our potential energy terms are all derived from work. Long ago we found that
the work done on an object may depend on the path the object takes. We found
work using 
\begin{equation}
w=\int_{x_{i}}^{x_{f}}\overrightarrow{F}\cdot \overrightarrow{dr}
\end{equation}%
and the work was the area under the force vs. displacement curve. But this
was work done moving an object as a whole. Now we can do work on an object
by compressing it, for example. This is different from moving the whole
object around. Let's call this new kind of work \emph{internal work. }This
new kind of work should be similar in form to regular mechanical work. We
expect it to depend on the details of how the internal work is done.

For mechanical work we defined the potential energy, $U,$ as 
\begin{equation}
w=-\Delta U
\end{equation}%
Using this concept, we showed that when $\Delta E_{th}=0$ mechanical energy, 
$E_{mech}$ . 
\begin{equation}
E_{mech}=K+U
\end{equation}%
was often conserved.%
\begin{eqnarray*}
E_{mech,i} &=&E_{mech,f} \\
K_{i}+U_{i} &=&K_{f}+U_{f}
\end{eqnarray*}

Remember we solved a lot of problems back in Principles of Physics I by
knowing that for a closed system, $E_{mech}$ did not change. In other words, 
$E_{mech}$ was independent of the path we take on our $F$ vs. $x$ diagram.
This makes $E_{mech}$ much easier to work with. A path independent quantity
makes problems easier.

We want another variable to work with that removes the difficulty of
worrying about the details of how the internal motion happens (the path).
But we can't use the kinetic energy for this. The kinetic energy is part of
the mechanical work of the object as a whole, and does not affect the
internal workings of the object.

Frictional forces do affect the object. Think of rubbing your hands
together. There is mechanical work done by your muscles. There is kinetic
energy while your hands move. But your hands heat up. This higher
temperature is still there for a while after your hands stop moving.
Eventually the temperature of your hands returns to normal. Back in
Principles of Physics I we called this $\Delta E_{th}$ but now we now know
to call this transfer of energy heat, $Q$ and to think about the internal
work $w_{int}.$

Now that we remember all of this, we are ready to find our path independent
quantity for internal energy. We define the change in internal energy as%
\begin{equation*}
\Delta E_{int}=Q\pm w_{int}
\end{equation*}

As with the mechanical energy, the change in internal energy does not depend
on the details of how the change happens (what we are calling the
\textquotedblleft path\textquotedblright ). The change in energy is used
because $Q$ and $w_{int}$ are transfer variables, so they \emph{do} depend
on the details of the transfer. For small changes we can write%
\begin{equation*}
\Delta E_{int}=\Delta Q\pm \Delta w_{int}
\end{equation*}%
or for very small changes, we could replace the $\Delta $'s with small $d$'s

\begin{equation*}
dE_{int}=dQ\pm dw_{int}
\end{equation*}%
but be careful! these are not true differentials because $dw$ depends on the
path! Integration of $dw$ is tricky and you would have to think about it
carefully. The relationship 
\begin{equation}
\Delta E_{int}=Q\pm w_{int}
\end{equation}%
is known as the \emph{first law of thermodynamics.}

Think for a moment about what the first law of thermodynamics tells us. It
says the internal energy of a sample of a material goes up if you add energy
to it. That is not at all surprising. It also tells us we can add energy to
the sample by doing heat $\left( Q\right) $ or by doing work $\left(
w_{int}\right) $ on the sample. This is really just conservation of energy!

Let's look at some systems and think about their internal energies.

\subsection{An Isolated System}

Last lecture we defined isolated systems. An isolated system is a system
separated from the rest of the universe. Nothing can act on it. When a
system is isolated, we have%
\begin{equation}
w_{int}=0
\end{equation}%
because no work can be done on an isolated system and 
\begin{equation}
Q=0
\end{equation}%
because no energy can flow by heat to or from an isolated system. Then 
\begin{equation}
\Delta E_{int}=\pm w_{int}+Q=0
\end{equation}

\subsection{A Cyclic Process}

Suppose we have a process, something that is done to a sample of matter, say
a gas, over and over again the say way. We call such a process a \emph{cyclic%
} process. This is like the pistons in your car going through the same
motion over and over again to compress the gas in your engine cylinders.
Let's practice with the idea of internal energy, work, heat and PV\ diagrams
for such a process. If a process repeats, it starts and ends in the same
state so then 
\begin{equation}
\Delta E_{int}=0
\end{equation}%
and so 
\begin{equation}
\Delta E_{int}=\pm w_{int}+Q=0
\end{equation}%
but this time $w_{int}$ and $Q$ are not zero. In this case. 
\begin{equation}
Q=-w_{int}
\end{equation}

But these examples were too easy. Let's take on a harder problem where we
actually find the amount of work done.

\section{Work}

Consider the work done by a piston as it compresses a gas. It might be good
if the amount of work was equal to the amount of energy transfer by heat
like in our last example. We often want this to be true! Otherwise you could
melt your engine as your car operates. We want an engine in a car (or train,
or airplane) to be a cyclic process. This explains why your car engine
surface gets hot. Work is being done in the cylinders in your engine, but to
make the process cyclic some energy must be lost by heat or the internal
energy would rise.

But car engines are a little bit complicated. Let's do a simpler problem
first.\footnote{%
There are details of car engines that we won't tackle until later. So not
everything in this lecture applies directly to car engines.} Unlike your car
pistons, let's do work slowly and quazi-statically so that the pressure
equalizes every time we move the piston a little bit. The piston moves, and
energy is gained by work, but then lost by heat. We can draw a $P%
%TCIMACRO{\TeXButton{V}{\ooalign{\hfil$V$\hfil\cr\kern0.1em--\hfil\cr}}}%
%BeginExpansion
\ooalign{\hfil$V$\hfil\cr\kern0.1em--\hfil\cr}%
%EndExpansion
$ diagram for this. \FRAME{dtbpFX}{2.0454in}{1.3651in}{0pt}{}{}{Plot}{%
\special{language "Scientific Word";type "MAPLEPLOT";width 2.0454in;height
1.3651in;depth 0pt;display "USEDEF";plot_snapshots TRUE;mustRecompute
FALSE;lastEngine "MuPAD";xmin "2";xmax "10";xviewmin "2";xviewmax
"10";yviewmin "100";yviewmax "900";viewset"XY";rangeset"X";plottype
4;labeloverrides 3;x-label "V(m^3)";y-label "P(Pa)";axesFont "Times New
Roman,12,0000000000,useDefault,normal";numpoints 100;plotstyle
"patch";axesstyle "normal";axestips FALSE;xis \TEXUX{V};var1name
\TEXUX{$V$};function
\TEXUX{$\MATRIX{2,2}{c}\VR{,,c,,,}{,,c,,,}{,,,,,}\HR{,,}\CELL{4}\CELL{%
\allowbreak 500}\CELL{8}\CELL{500}$};linecolor "green";linestyle
1;pointstyle "circle";linethickness 1;lineAttributes "Solid";curveColor
"[flat::RGB:0x00008000]";curveStyle "Line";function
\TEXUX{$\MATRIX{2,2}{c}\VR{,,c,,,}{,,c,,,}{,,,,,}\HR{,,}\CELL{4}\CELL{%
\allowbreak 500}\CELL{8}\CELL{500}$};linecolor "cyan";linestyle 1;pointplot
TRUE;pointstyle "circle";linethickness 1;lineAttributes "Solid";curveColor
"[flat::RGB:0x00008080]";curveStyle "Point";VCamFile
'SYUDSZ1M.xvz';valid_file "T";tempfilename
'SYUDSZ0J.wmf';tempfile-properties "XPR";}}For this situation the pressure
won't change. Let's call such a process where the pressure doesn't change in
the cylinder an \emph{isobaric process. }

\FRAME{dtbpF}{2.226in}{2.1049in}{0pt}{}{}{Figure}{\special{language
"Scientific Word";type "GRAPHIC";maintain-aspect-ratio TRUE;display
"USEDEF";valid_file "T";width 2.226in;height 2.1049in;depth
0pt;original-width 2.9187in;original-height 2.7588in;cropleft "0";croptop
"1";cropright "1";cropbottom "0";tempfilename
'SZ1CK400.wmf';tempfile-properties "XPR";}}

Lets take a close look at the gas within the piston. At equilibrium we know
that the pressure is 
\begin{equation*}
P=\frac{F_{P}}{A}
\end{equation*}%
or the force due to pressure is%
\begin{equation*}
F_{P}=PA
\end{equation*}%
This force is the force of the gas pushing on the piston. We could label it 
\begin{equation*}
F_{pg}=F_{P}
\end{equation*}%
where the subscript $p$ is for piston and the subscript $g$ is for gas. So $%
F_{pg}$ is the force on the piston due to the gas. Notice that the gas
pushes to the right. If we apply an external force, $F_{ex},$ to the piston
slowly so as to allow the system to remain in thermal equilibrium as we go,
then 
\begin{equation*}
\overrightarrow{F}_{ex}=-F_{ex}\hat{x}
\end{equation*}%
because we are pushing in the $-x$-direction. Again, we will move the piston
slowly so we don't violate our quasi-static rule. As we move the piston the
gas inside the piston adjusts in pressure until we are again in equilibrium.
So the magnitudes are nearly the same 
\begin{equation*}
F_{ex}\approx F_{pg}
\end{equation*}%
but the pressure force $F_{pg}$ pushes the opposite direction 
\begin{eqnarray*}
\overrightarrow{F}_{pg} &=&+F_{ex}\hat{x} \\
&=&PA\hat{x}
\end{eqnarray*}%
We should also note that since neither the piston nor the gas are
accelerating, Newton's third law tells us that the force of the gas on the
piston and the force of the piston on the gas are equal in magnitude 
\begin{equation*}
F_{gp}=F_{pg}
\end{equation*}%
but opposite in direction. 
\begin{equation*}
\overrightarrow{F}_{gp}\hat{x}=-F_{pg}\hat{x}
\end{equation*}%
When we move the piston to the left, we have a displacement%
\begin{equation*}
\overrightarrow{dr}=-dx\hat{x}
\end{equation*}%
and when we move the piston to the right 
\begin{equation*}
\overrightarrow{dr}=+dx\hat{x}
\end{equation*}%
The work done \emph{on the gas by the piston would be }%
\begin{equation*}
w_{gp}=\int_{x_{i}}^{x_{f}}\overrightarrow{F}_{ex}\cdot d\overrightarrow{r}
\end{equation*}%
Remembering how to use a dot product we can write this as 
\begin{equation*}
w_{gp}=\int_{x_{i}}^{x_{f}}F_{ex}dx\cos \theta _{Fx}
\end{equation*}%
It is worth remembering that in this formula $F_{ex}$ and $dx$ are
magnitudes. Let's write this as%
\begin{equation*}
w_{gp}=\int_{x_{i}}^{x_{f}}\left\vert F_{ex}\right\vert \left\vert
dx\right\vert \cos \theta _{Fx}
\end{equation*}%
to remind us. Substituting in $\left\vert F_{gp}\right\vert $ for $%
\left\vert F_{ex}\right\vert $ (since they have nearly the same magnitude)
we have 
\begin{eqnarray*}
w_{gp} &=&\int_{x_{i}}^{x_{f}}\left\vert F_{gp}\right\vert \left\vert
dx\right\vert \cos \theta _{Fx} \\
&=&\int_{x_{i}}^{x_{f}}\left\vert PA\right\vert \left\vert dx\right\vert
\cos \theta _{Fx}
\end{eqnarray*}%
Suppose we move the piston to the right as shown in the last figure. The
angle between $F_{ex}$ and $dx$ would be $180\unit{%
%TCIMACRO{\U{b0}}%
%BeginExpansion
{{}^\circ}%
%EndExpansion
}.$ If this seems mysterious, look at the last figure, and remember we are
really working with $F_{ex}.$ We put in $F_{gp}$ because it's magnitude is
the same as that of $F_{ex}.$ So $180\unit{%
%TCIMACRO{\U{b0}}%
%BeginExpansion
{{}^\circ}%
%EndExpansion
}$ is the right angle for our dot product. And $\cos \left( 180\unit{%
%TCIMACRO{\U{b0}}%
%BeginExpansion
{{}^\circ}%
%EndExpansion
}\right) =-1$ so

\begin{equation*}
w_{gp}=-\int_{x_{i}}^{x_{f}}\left\vert PA\right\vert \left\vert dx\right\vert
\end{equation*}%
To get the total work done by the piston on the gas we are integrating over
the entire path the piston travels. Along that path $P$ can only be
positive, and the same is true for $A$. And both $P$ and $A$ are constant,
since we are using our quazi-static process and allowing the pressure to
equalize as we go. So we could write the work done on the gas by the piston
as 
\begin{equation*}
w_{gp}=-PA\int_{x_{i}}^{x_{f}}\left\vert dx\right\vert
\end{equation*}%
but now we need to consider $\left\vert dx\right\vert .$ We know that $dx$
is a small displacement, and displacements can be negative or positive. And
in our case, as the piston goes to the right, $dx$ certainly is positive.
Then we can just drop the absolute value signs. 
\begin{equation*}
w_{gp}=-PA\int_{x_{i}}^{x_{f}}dx
\end{equation*}%
Remember that we are using an isobaric process in compressing the gas That
is, we keep the pressure the same but allow the temperature to change, so
that the volume and temperature change proportionally 
\begin{equation*}
%TCIMACRO{\TeXButton{V}{\ooalign{\hfil$V$\hfil\cr\kern0.1em--\hfil\cr}}}%
%BeginExpansion
\ooalign{\hfil$V$\hfil\cr\kern0.1em--\hfil\cr}%
%EndExpansion
=(\frac{nR}{P})T)
\end{equation*}%
but the pressure is constant. So our integral is just 
\begin{equation*}
w_{gp}=-PA\Delta x
\end{equation*}%
And this tells us something! To expand our gas, we got negative work. That
tells us that our external force is not winning, it is not causing the
motion. We must be getting energy from somewhere else that is making the
piston go out, overcoming the external force.

But now let's say that the external force is winning so we compress our gas. 
\FRAME{dtbpF}{2.0859in}{1.9199in}{0in}{}{}{Figure}{\special{language
"Scientific Word";type "GRAPHIC";maintain-aspect-ratio TRUE;display
"USEDEF";valid_file "T";width 2.0859in;height 1.9199in;depth
0in;original-width 2.047in;original-height 1.8818in;cropleft "0";croptop
"1";cropright "1";cropbottom "0";tempfilename
'THRKT900.wmf';tempfile-properties "XPR";}}

And let's go back to 
\begin{equation*}
w_{gp}=-PA\int_{x_{i}}^{x_{f}}\left\vert dx\right\vert
\end{equation*}%
This time our $dx$ will be negative. This negative plus the one out in front
should give us positive work for the external force. It is winning.

But here we can be a little bit clever. We want to get rid of the absolute
value signs, and if $dx$ really is negative, then the negative hiding inside
of $dx$ really does cancel the negative sign from our dot product. And then 
\begin{equation*}
w_{gp}=-PA\int_{x_{i}}^{x_{f}}dx
\end{equation*}%
gives us positive work \footnote{%
Think of what we did with springs in Principles of Physics I to deal with $%
S=-k\left( x-x_{o}\right) $ in work integrals}. Once again we don't have to
worry about integrating with absolute value signs. So once again our
integral is just 
\begin{eqnarray*}
w_{gp} &=&-PA\int_{x_{i}}^{x_{f}}dx \\
&=&-PA\left( x_{f}-x_{i}\right)
\end{eqnarray*}%
But because $x_{f}<x_{i}$ this integral is not negative! If we compress our
gas, we get a positive amount of work. This is very convenient, we get the
same equation for either expansion or compression of the gas.

We can rewrite this result using what we know about the volume of the gas 
\begin{equation*}
\Delta 
%TCIMACRO{\TeXButton{V}{\ooalign{\hfil$V$\hfil\cr\kern0.1em--\hfil\cr}}}%
%BeginExpansion
\ooalign{\hfil$V$\hfil\cr\kern0.1em--\hfil\cr}%
%EndExpansion
=A\Delta x
\end{equation*}%
and for this case $A$ is a constant, so the work done on the gas by the
piston is 
\begin{eqnarray*}
w_{gp} &=&-P\left( Ax_{f}-Ax_{i}\right) \\
&=&-P\left( 
%TCIMACRO{\TeXButton{V}{\ooalign{\hfil$V$\hfil\cr\kern0.1em--\hfil\cr}}}%
%BeginExpansion
\ooalign{\hfil$V$\hfil\cr\kern0.1em--\hfil\cr}%
%EndExpansion
_{f}-%
%TCIMACRO{\TeXButton{V}{\ooalign{\hfil$V$\hfil\cr\kern0.1em--\hfil\cr}}}%
%BeginExpansion
\ooalign{\hfil$V$\hfil\cr\kern0.1em--\hfil\cr}%
%EndExpansion
_{i}\right)
\end{eqnarray*}%
This gives us an idea. We could write the quantity $Adx$ as $d%
%TCIMACRO{\TeXButton{V}{\ooalign{\hfil$V$\hfil\cr\kern0.1em--\hfil\cr}}}%
%BeginExpansion
\ooalign{\hfil$V$\hfil\cr\kern0.1em--\hfil\cr}%
%EndExpansion
$ and change the limits of integration from $x_{f}$ and $x_{i}$ to $%
%TCIMACRO{\TeXButton{V}{\ooalign{\hfil$V$\hfil\cr\kern0.1em--\hfil\cr}}}%
%BeginExpansion
\ooalign{\hfil$V$\hfil\cr\kern0.1em--\hfil\cr}%
%EndExpansion
_{f}$ and $%
%TCIMACRO{\TeXButton{V}{\ooalign{\hfil$V$\hfil\cr\kern0.1em--\hfil\cr}}}%
%BeginExpansion
\ooalign{\hfil$V$\hfil\cr\kern0.1em--\hfil\cr}%
%EndExpansion
_{i}.$ Then our integral would be 
\begin{eqnarray*}
w_{gp} &=&-\int_{x_{i}}^{x_{f}}PAdx \\
&=&-\int_{%
%TCIMACRO{\TeXButton{V}{\ooalign{\hfil$V$\hfil\cr\kern0.1em--\hfil\cr}}}%
%BeginExpansion
\ooalign{\hfil$V$\hfil\cr\kern0.1em--\hfil\cr}%
%EndExpansion
_{i}}^{%
%TCIMACRO{\TeXButton{V}{\ooalign{\hfil$V$\hfil\cr\kern0.1em--\hfil\cr}}}%
%BeginExpansion
\ooalign{\hfil$V$\hfil\cr\kern0.1em--\hfil\cr}%
%EndExpansion
_{f}}Pd%
%TCIMACRO{\TeXButton{V}{\ooalign{\hfil$V$\hfil\cr\kern0.1em--\hfil\cr}}}%
%BeginExpansion
\ooalign{\hfil$V$\hfil\cr\kern0.1em--\hfil\cr}%
%EndExpansion
\end{eqnarray*}%
This gives the same answer for the work done on the gas for our isobaric case%
\begin{equation*}
w_{gp}=-P\left( 
%TCIMACRO{\TeXButton{V}{\ooalign{\hfil$V$\hfil\cr\kern0.1em--\hfil\cr}}}%
%BeginExpansion
\ooalign{\hfil$V$\hfil\cr\kern0.1em--\hfil\cr}%
%EndExpansion
_{f}-%
%TCIMACRO{\TeXButton{V}{\ooalign{\hfil$V$\hfil\cr\kern0.1em--\hfil\cr}}}%
%BeginExpansion
\ooalign{\hfil$V$\hfil\cr\kern0.1em--\hfil\cr}%
%EndExpansion
_{i}\right)
\end{equation*}

But this is the amount of work that the piston does on the gas. Is the gas
pushing back on the piston? We know it is! But let's review what we know
about work from Principles of Physics I.

Suppose we have two people exerting a force on a box. One is a professional
American football player, The other is a six-year-old. \FRAME{dtbpF}{2.8158in%
}{1.8524in}{0pt}{}{}{Figure}{\special{language "Scientific Word";type
"GRAPHIC";maintain-aspect-ratio TRUE;display "USEDEF";valid_file "T";width
2.8158in;height 1.8524in;depth 0pt;original-width 3.4082in;original-height
2.2329in;cropleft "0";croptop "1";cropright "1";cropbottom "0";tempfilename
'SVAQIE7A.wmf';tempfile-properties "XPR";}}They push in opposite directions.
How much work is being done? To answer this, let's consider a free-body
diagram.\FRAME{dtbpF}{1.8291in}{1.8377in}{0pt}{}{}{Figure}{\special{language
"Scientific Word";type "GRAPHIC";maintain-aspect-ratio TRUE;display
"USEDEF";valid_file "T";width 1.8291in;height 1.8377in;depth
0pt;original-width 1.7919in;original-height 1.8005in;cropleft "0";croptop
"1";cropright "1";cropbottom "0";tempfilename
'SVAQIE7B.wmf';tempfile-properties "XPR";}}where the subscript $B$ is for
the box and the subscript $6$ is for the six-year-old and $L$ is for the
football player (In this case, a (L)inebacker). We can see that there will
be a net force.%
\begin{equation*}
F_{net_{s}}=N_{BL}-N_{B6}-f_{BF}
\end{equation*}%
But which way will the box go? We could guess that the football player will
push the box \emph{and} the six-year-old to the left. So $\Delta s$ will be
negative. Then the work done on the box by the football player will be 
\begin{eqnarray*}
w_{BL} &=&\int_{s_{i}}^{s_{f}}N_{BL}ds \\
&=&N_{BL}\Delta s
\end{eqnarray*}%
but the work done by the child would be 
\begin{eqnarray*}
w_{B6} &=&\int_{s_{i}}^{s_{f}}-N_{B6}ds \\
&=&-N_{BM}\Delta s
\end{eqnarray*}%
The child's work is negative! What can that mean? Well, for starters, the
child's force can't be the force that is making the box move. In fact, the
child's force is another obstacle that the linebacker's force must overcome
to make the box move. This means that the linebacker must do enough work
(push hard enough) to make the box go, \emph{and} to overcome the backward
push of the child.

But that is just our case with the gas and piston The gas pressure force 
\emph{is} pushing on the piston. But in our last example the gas pressure
force on the piston is losing ground (the external force is pushing the
piston to the left). So we would expect the work done by the gas on the
piston to be negative. The force on the piston from the gas is losing. The
force of the piston on the gas is wining.

But $F_{pg}=F_{gp}$ because they are a third law pair. The force on the
piston by the gas also does work. It will have the same magnitude, but will
be in the opposite direction. Since we assumed a quasi-static process, the
two forces must always be nearly the same, 
\begin{equation*}
\left\vert F_{ex}\right\vert =\left\vert F_{pg}\right\vert =\left[ F_{gp}%
\right]
\end{equation*}%
so the amount of work must be nearly the same, but the sign must be negative
for the work on the piston due to the gas's push. We wrote the work done by
the piston on the gas as 
\begin{equation*}
w_{gp}=-\int_{%
%TCIMACRO{\TeXButton{V}{\ooalign{\hfil$V$\hfil\cr\kern0.1em--\hfil\cr}}}%
%BeginExpansion
\ooalign{\hfil$V$\hfil\cr\kern0.1em--\hfil\cr}%
%EndExpansion
_{i}}^{%
%TCIMACRO{\TeXButton{V}{\ooalign{\hfil$V$\hfil\cr\kern0.1em--\hfil\cr}}}%
%BeginExpansion
\ooalign{\hfil$V$\hfil\cr\kern0.1em--\hfil\cr}%
%EndExpansion
_{f}}Pd%
%TCIMACRO{\TeXButton{V}{\ooalign{\hfil$V$\hfil\cr\kern0.1em--\hfil\cr}}}%
%BeginExpansion
\ooalign{\hfil$V$\hfil\cr\kern0.1em--\hfil\cr}%
%EndExpansion
\end{equation*}%
Then the work done by the gas on the piston must be 
\begin{eqnarray*}
w_{pg} &=&-\left( -\int_{%
%TCIMACRO{\TeXButton{V}{\ooalign{\hfil$V$\hfil\cr\kern0.1em--\hfil\cr}}}%
%BeginExpansion
\ooalign{\hfil$V$\hfil\cr\kern0.1em--\hfil\cr}%
%EndExpansion
_{i}}^{%
%TCIMACRO{\TeXButton{V}{\ooalign{\hfil$V$\hfil\cr\kern0.1em--\hfil\cr}}}%
%BeginExpansion
\ooalign{\hfil$V$\hfil\cr\kern0.1em--\hfil\cr}%
%EndExpansion
_{f}}Pd%
%TCIMACRO{\TeXButton{V}{\ooalign{\hfil$V$\hfil\cr\kern0.1em--\hfil\cr}}}%
%BeginExpansion
\ooalign{\hfil$V$\hfil\cr\kern0.1em--\hfil\cr}%
%EndExpansion
\right) \\
&=&\int_{%
%TCIMACRO{\TeXButton{V}{\ooalign{\hfil$V$\hfil\cr\kern0.1em--\hfil\cr}}}%
%BeginExpansion
\ooalign{\hfil$V$\hfil\cr\kern0.1em--\hfil\cr}%
%EndExpansion
_{i}}^{%
%TCIMACRO{\TeXButton{V}{\ooalign{\hfil$V$\hfil\cr\kern0.1em--\hfil\cr}}}%
%BeginExpansion
\ooalign{\hfil$V$\hfil\cr\kern0.1em--\hfil\cr}%
%EndExpansion
_{f}}Pd%
%TCIMACRO{\TeXButton{V}{\ooalign{\hfil$V$\hfil\cr\kern0.1em--\hfil\cr}}}%
%BeginExpansion
\ooalign{\hfil$V$\hfil\cr\kern0.1em--\hfil\cr}%
%EndExpansion
\end{eqnarray*}%
That is, the work done on the gas by the piston is the negative of the work
done on the piston by the gas.

We are principally interested in what happens to the gas, (we studied
objects like the piston in Principles of Physics I) so the first equation is
what we are looking for. This is the internal work done on the gas!

\begin{equation}
w_{gp}=w_{\text{on the gas}}=w_{int}=-\int_{%
%TCIMACRO{\TeXButton{V}{\ooalign{\hfil$V$\hfil\cr\kern0.1em--\hfil\cr}}}%
%BeginExpansion
\ooalign{\hfil$V$\hfil\cr\kern0.1em--\hfil\cr}%
%EndExpansion
_{i}}^{%
%TCIMACRO{\TeXButton{V}{\ooalign{\hfil$V$\hfil\cr\kern0.1em--\hfil\cr}}}%
%BeginExpansion
\ooalign{\hfil$V$\hfil\cr\kern0.1em--\hfil\cr}%
%EndExpansion
_{f}}Pd%
%TCIMACRO{\TeXButton{V}{\ooalign{\hfil$V$\hfil\cr\kern0.1em--\hfil\cr}}}%
%BeginExpansion
\ooalign{\hfil$V$\hfil\cr\kern0.1em--\hfil\cr}%
%EndExpansion
\end{equation}%
And we could also define 
\begin{equation*}
w_{pg}=w_{\text{by the gas}}=\int_{%
%TCIMACRO{\TeXButton{V}{\ooalign{\hfil$V$\hfil\cr\kern0.1em--\hfil\cr}}}%
%BeginExpansion
\ooalign{\hfil$V$\hfil\cr\kern0.1em--\hfil\cr}%
%EndExpansion
_{i}}^{%
%TCIMACRO{\TeXButton{V}{\ooalign{\hfil$V$\hfil\cr\kern0.1em--\hfil\cr}}}%
%BeginExpansion
\ooalign{\hfil$V$\hfil\cr\kern0.1em--\hfil\cr}%
%EndExpansion
_{f}}Pd%
%TCIMACRO{\TeXButton{V}{\ooalign{\hfil$V$\hfil\cr\kern0.1em--\hfil\cr}}}%
%BeginExpansion
\ooalign{\hfil$V$\hfil\cr\kern0.1em--\hfil\cr}%
%EndExpansion
\end{equation*}%
The only difference is the minus sign. We are going to have to be careful!

Still, this looks easy enough, but to integrate we must know how $P$ depends
on $%
%TCIMACRO{\TeXButton{V}{\ooalign{\hfil$V$\hfil\cr\kern0.1em--\hfil\cr}}}%
%BeginExpansion
\ooalign{\hfil$V$\hfil\cr\kern0.1em--\hfil\cr}%
%EndExpansion
$ during the whole process. It was simple for an isobaric process, but it
could be much harder for other processes. We can plot $P$ vs. $%
%TCIMACRO{\TeXButton{V}{\ooalign{\hfil$V$\hfil\cr\kern0.1em--\hfil\cr}}}%
%BeginExpansion
\ooalign{\hfil$V$\hfil\cr\kern0.1em--\hfil\cr}%
%EndExpansion
$ to see how it varies for a process. We would expect the integral to be the
area under a $P%
%TCIMACRO{\TeXButton{V}{\ooalign{\hfil$V$\hfil\cr\kern0.1em--\hfil\cr}}}%
%BeginExpansion
\ooalign{\hfil$V$\hfil\cr\kern0.1em--\hfil\cr}%
%EndExpansion
$ curve. This is very like mechanical work, which was the area under the $F$-%
$x$ curve.\FRAME{dhF}{2.7951in}{2.2883in}{0pt}{}{}{Figure}{\special{language
"Scientific Word";type "GRAPHIC";maintain-aspect-ratio TRUE;display
"USEDEF";valid_file "T";width 2.7951in;height 2.2883in;depth
0pt;original-width 2.7527in;original-height 2.2485in;cropleft "0";croptop
"1";cropright "1";cropbottom "0";tempfilename
'SVAQIE7C.wmf';tempfile-properties "XPR";}}

And we can make a general statement: The work done on a gas in a
quasi-static process that takes the gas from an initial state to a final
state is the negative of the area under the curve on a PV diagram evaluated
between the initial and final states.

Unfortunately, the work done depends on the shape of the curve on the $P%
%TCIMACRO{\TeXButton{V}{\ooalign{\hfil$V$\hfil\cr\kern0.1em--\hfil\cr}}}%
%BeginExpansion
\ooalign{\hfil$V$\hfil\cr\kern0.1em--\hfil\cr}%
%EndExpansion
$ diagram. Several are shown below. Each has a different amount of work done
(notice the different areas under the curves), but the same initial and
final points.

\FRAME{dhF}{3.4999in}{1.0438in}{0pt}{}{}{Figure}{\special{language
"Scientific Word";type "GRAPHIC";maintain-aspect-ratio TRUE;display
"USEDEF";valid_file "T";width 3.4999in;height 1.0438in;depth
0pt;original-width 6.5337in;original-height 1.9285in;cropleft "0";croptop
"1";cropright "1";cropbottom "0";tempfilename
'SVAQIE7D.wmf';tempfile-properties "XPR";}}

Also notice that we have calculated the amount of work done on the gas. For
physicists this is normally what we want.

You might be wondering if that is what we want to build car or jet engines.
Don't we want the gas to expand to push the piston and make the car go? And
of course, you are right in your thinking. Eventually we want to know the
work done by the gas on the piston, $w_{pg}.$ But for now, we will look at $%
w_{gp}.$ And from what we have done we know these two works are very related
for a quasi-static process. 
\begin{equation*}
w_{\text{on the gas}}=-w_{\text{by the gas}}
\end{equation*}

Let's go back to our first law of thermodynamics. 
\begin{equation}
\Delta E_{int}=Q+w_{int}
\end{equation}%
This time, I just have a plus sign on $w_{int.}.$ If we consider work done
on the gas then this work in positive when we add to $\Delta E_{int}.$ The
energy has increased in the gas due to the work done on the gas.%
\begin{equation}
\Delta E_{int}=Q+w_{\text{on the gas}}  \label{First_law_v1}
\end{equation}

We could write this as%
\begin{equation}
\Delta E_{int}=Q-w_{\text{by the gas}}  \label{First_Law_V2}
\end{equation}%
and this would be just as good. Some books use the first (Equation \ref%
{First_law_v1}) and some use the second (Equation \ref{First_Law_V2}). In
our class we will use the first (Equation \ref{First_law_v1}). But as long
as we remember to check for if we have $w_{\text{on the gas}}$ or $w_{\text{%
by the gas}}$ either equation works.

Regardless of how we write our equation, We know how $w_{int}$ changes on a
PV\ diagram. How do $\Delta E_{in}$ and $Q$ change on a PV\ diagram? For
energy transfer by heat, we don't need a pressure or volume change, so a PV
diagram isn't a great way to show $Q.$ But we can consider $\Delta E_{int}$
on a PV\ diagram. We already know that for ideal gasses 
\begin{equation}
E_{int}=\frac{3}{2}k_{B}T
\end{equation}%
So $E_{int}$ changes with $T.$ But there is no $T$ on a PV diagram! For
ideal gasses we know how to solve this problem. Suppose we have $n$ moles
gas in a cylinder. We can solve for $P$ as a function of $V$%
\begin{equation*}
P=\frac{nRT}{%
%TCIMACRO{\TeXButton{V}{\ooalign{\hfil$V$\hfil\cr\kern0.1em--\hfil\cr}}}%
%BeginExpansion
\ooalign{\hfil$V$\hfil\cr\kern0.1em--\hfil\cr}%
%EndExpansion
}
\end{equation*}%
Now we want to let $T$ be constant. Along a line of constant $T$ our $%
E_{int} $ will be constant. 
\begin{equation*}
P=\left( nRT\right) \frac{1}{%
%TCIMACRO{\TeXButton{V}{\ooalign{\hfil$V$\hfil\cr\kern0.1em--\hfil\cr}}}%
%BeginExpansion
\ooalign{\hfil$V$\hfil\cr\kern0.1em--\hfil\cr}%
%EndExpansion
}
\end{equation*}%
Everything in the parenthesis is constant so 
\begin{equation*}
P\propto \frac{1}{%
%TCIMACRO{\TeXButton{V}{\ooalign{\hfil$V$\hfil\cr\kern0.1em--\hfil\cr}}}%
%BeginExpansion
\ooalign{\hfil$V$\hfil\cr\kern0.1em--\hfil\cr}%
%EndExpansion
}
\end{equation*}%
We can plot this for

\begin{equation*}
\begin{tabular}{ll}
$n=1\unit{mol}$ & number of moles \\ 
$R=8.314\frac{\unit{J}}{\unit{mol}\unit{K}}$ & Universal gas constant \\ 
$T=290\unit{K}$ & Temperature in Kelvins%
\end{tabular}%
\end{equation*}

\FRAME{dtbpFX}{3.1669in}{2.111in}{0pt}{}{}{Plot}{\special{language
"Scientific Word";type "MAPLEPLOT";width 3.1669in;height 2.111in;depth
0pt;display "USEDEF";plot_snapshots TRUE;mustRecompute FALSE;lastEngine
"MuPAD";animated TRUE;animationStartTime "0"; animationEndTime "10";
animationFramesPerSecond "10";xmin "0";xmax "40";animationParamMin
"290";animationParamMax "350";xviewmin "0";xviewmax "40";yviewmin
"0";yviewmax "400";viewset"XY";rangeset"X";plottype 4;labeloverrides
3;x-label "V(m^3)";y-label "P(Pa)";axesFont "Times New
Roman,12,0000000000,useDefault,normal";numpoints 100;plotstyle
"patch";axesstyle "normal";axestips FALSE;xis \TEXUX{V};animationParam
\TEXUX{t};var1name \TEXUX{$V$};animationParamList \TEXUX{$t$};function
\TEXUX{$\allowbreak 8.\,\allowbreak 314\frac{1}{V}290$};linecolor
"blue";linestyle 1;pointstyle "point";linethickness 3;lineAttributes
"Solid";var1range "0,40";animationParamRange "290,350";num-x-gridlines
100;curveColor "[flat::RGB:0x000000ff]";curveStyle "Line";VCamFile
'SVAOR55H.xvz';valid_file "T";tempfilename
'SVAORN5O.wmf';tempfile-properties "XPR";}}All along this line $T$ is
constant so $E_{int}$ is constant. We have already seen such a line on a PV
diagram and know it is called an \emph{isotherm}. Now let's plot it again
for a different value of $T$%
\begin{equation*}
\begin{tabular}{ll}
$n=1\unit{mol}$ & number of moles \\ 
$R=8.314\frac{\unit{J}}{\unit{mol}\unit{K}}$ & Universal gas constant \\ 
$T=590\unit{K}$ & Temperature in Kelvins%
\end{tabular}%
\end{equation*}%
\FRAME{dtbpFX}{3.1665in}{2.111in}{0pt}{}{}{Plot}{\special{language
"Scientific Word";type "MAPLEPLOT";width 3.1665in;height 2.111in;depth
0pt;display "USEDEF";plot_snapshots TRUE;mustRecompute FALSE;lastEngine
"MuPAD";animated TRUE;animationStartTime "0"; animationEndTime "10";
animationFramesPerSecond "10";xmin "0";xmax "40";animationParamMin
"290";animationParamMax "350";xviewmin "0";xviewmax "40";yviewmin
"0";yviewmax "400";viewset"XY";rangeset"X";plottype 4;labeloverrides
3;x-label "V(m^3)";y-label "P(Pa)";axesFont "Times New
Roman,12,0000000000,useDefault,normal";numpoints 100;plotstyle
"patch";axesstyle "normal";axestips FALSE;xis \TEXUX{x};animationParam
\TEXUX{t};var1name \TEXUX{$x$};animationParamList \TEXUX{$t$};function
\TEXUX{$\allowbreak 8.\,\allowbreak 314\frac{1}{x}590$};linecolor
"blue";linestyle 1;pointstyle "point";linethickness 3;lineAttributes
"Solid";var1range "0,40";animationParamRange "290,350";num-x-gridlines
100;curveColor "[flat::RGB:0x000000ff]";curveStyle "Line";VCamFile
'SVAOSP5T.xvz';valid_file "T";tempfilename
'SVAORS5P.wmf';tempfile-properties "XPR";}}Notice that the curve looks like
the same shape, but it has been shifted toward the upper right hand corner
of the graph. If we place both the $T=290$ and $T=590$ graphs on the same
plot, we can better see the shift. \FRAME{dtbpFX}{3.1665in}{2.111in}{0pt}{}{%
}{Plot}{\special{language "Scientific Word";type "MAPLEPLOT";width
3.1665in;height 2.111in;depth 0pt;display "USEDEF";plot_snapshots
TRUE;mustRecompute FALSE;lastEngine "MuPAD";animated TRUE;animationStartTime
"0"; animationEndTime "10"; animationFramesPerSecond "10";xmin "0";xmax
"40";animationParamMin "290";animationParamMax "350";xviewmin "0";xviewmax
"40";yviewmin "0";yviewmax "400";viewset"XY";rangeset"X";plottype
4;labeloverrides 3;x-label "V(m^3)";y-label "P(Pa)";axesFont "Times New
Roman,12,0000000000,useDefault,normal";numpoints 100;plotstyle
"patch";axesstyle "normal";axestips FALSE;xis \TEXUX{x};animationParam
\TEXUX{t};var1name \TEXUX{$x$};animationParamList \TEXUX{$t$};function
\TEXUX{$\allowbreak 8.\,\allowbreak 314\frac{1}{x}290$};linecolor
"blue";linestyle 1;pointstyle "point";linethickness 1;lineAttributes
"Solid";var1range "0,40";animationParamRange "290,350";num-x-gridlines
100;curveColor "[flat::RGB:0x000000ff]";curveStyle "Line";function
\TEXUX{$\allowbreak 8.\,\allowbreak 314\frac{1}{x}590$};linecolor
"blue";linestyle 1;pointstyle "point";linethickness 3;lineAttributes
"Solid";var1range "0,40";animationParamRange "290,350";num-x-gridlines
100;curveColor "[flat::RGB:0x000000ff]";curveStyle "Line";VCamFile
'SVAOT55U.xvz';valid_file "T";tempfilename
'SVAORY5Q.wmf';tempfile-properties "XPR";}}The dark curve is our new pot.
The light curve is the old plot for comparison. The curves look a lot alike.
They are essentially the same shape. But their location (how far they are
from the origin) depends on the temperature, with warmer temperatures moving
to the upper right hand part off the graph. So our $P$ vs. $%
%TCIMACRO{\TeXButton{V}{\ooalign{\hfil$V$\hfil\cr\kern0.1em--\hfil\cr}}}%
%BeginExpansion
\ooalign{\hfil$V$\hfil\cr\kern0.1em--\hfil\cr}%
%EndExpansion
$ graph can tell us something about temperature. And because $E_{int}=\frac{3%
}{2}k_{B}T$ these lines and their positions also tell us about the internal
energy. A change in internal energy would be going from one such isothermal
line to another.\FRAME{dtbpF}{3.461in}{2.3116in}{0pt}{}{}{Figure}{\special%
{language "Scientific Word";type "GRAPHIC";maintain-aspect-ratio
TRUE;display "USEDEF";valid_file "T";width 3.461in;height 2.3116in;depth
0pt;original-width 3.4151in;original-height 2.2719in;cropleft "0";croptop
"1";cropright "1";cropbottom "0";tempfilename
'SYU9LY05.wmf';tempfile-properties "XPR";}}Because $\Delta E_{int}$ is path
independent (just depends on $T$), it doesn't matter how we go from $%
E_{1_{int}}$ to $E_{2_{int}}.$ So if we followed this PV path \FRAME{dtbpF}{%
2.8383in}{1.8983in}{0pt}{}{}{Figure}{\special{language "Scientific
Word";type "GRAPHIC";maintain-aspect-ratio TRUE;display "USEDEF";valid_file
"T";width 2.8383in;height 1.8983in;depth 0pt;original-width
2.7951in;original-height 1.8602in;cropleft "0";croptop "1";cropright
"1";cropbottom "0";tempfilename 'SYU9MU06.wmf';tempfile-properties "XPR";}}%
or this PV\ path \FRAME{dtbpF}{2.8677in}{1.9182in}{0pt}{}{}{Figure}{\special%
{language "Scientific Word";type "GRAPHIC";maintain-aspect-ratio
TRUE;display "USEDEF";valid_file "T";width 2.8677in;height 1.9182in;depth
0pt;original-width 2.8253in;original-height 1.8801in;cropleft "0";croptop
"1";cropright "1";cropbottom "0";tempfilename
'SYU9NH07.wmf';tempfile-properties "XPR";}}we would get the same $\Delta
E_{int}.$ A way to think about this is that $\Delta E_{int}$ only depends on
it's starting and end points (like a displacement). So any path from $%
E_{1_{int}}$ to $E_{2_{int}}$ gives the same change in internal energy. 
\FRAME{dtbpF}{2.7181in}{1.8178in}{0pt}{}{}{Figure}{\special{language
"Scientific Word";type "GRAPHIC";maintain-aspect-ratio TRUE;display
"USEDEF";valid_file "T";width 2.7181in;height 1.8178in;depth
0pt;original-width 2.6757in;original-height 1.7798in;cropleft "0";croptop
"1";cropright "1";cropbottom "0";tempfilename
'SYU9QT08.wmf';tempfile-properties "XPR";}}

Looking at an isothermal process can help us understand our other processes.
Suppose we have two different isothermal processes, one at $100\unit{K}$ and
one at $300\unit{K}.$ The isotherms look like this on a $P%
%TCIMACRO{\TeXButton{V}{\ooalign{\hfil$V$\hfil\cr\kern0.1em--\hfil\cr}}}%
%BeginExpansion
\ooalign{\hfil$V$\hfil\cr\kern0.1em--\hfil\cr}%
%EndExpansion
$ diagram: \FRAME{dtbpF}{2.7612in}{1.8316in}{0pt}{}{}{Figure}{\special%
{language "Scientific Word";type "GRAPHIC";maintain-aspect-ratio
TRUE;display "USEDEF";valid_file "T";width 2.7612in;height 1.8316in;depth
0pt;original-width 2.7851in;original-height 1.8378in;cropleft "0";croptop
"1";cropright "1";cropbottom "0";tempfilename
'SYUH880T.wmf';tempfile-properties "XPR";}}Every point on the $300\unit{K}$
line has the same temperature. Every point on the $100\unit{K}$ line has the
same temperature. We can see that points on the $P%
%TCIMACRO{\TeXButton{V}{\ooalign{\hfil$V$\hfil\cr\kern0.1em--\hfil\cr}}}%
%BeginExpansion
\ooalign{\hfil$V$\hfil\cr\kern0.1em--\hfil\cr}%
%EndExpansion
\ $diagram that are closer to the origin must have lower temperatures. We
can see the points farther from the origin have higher temperatures. \FRAME{%
dtbpF}{2.8091in}{1.8636in}{0pt}{}{}{Figure}{\special{language "Scientific
Word";type "GRAPHIC";maintain-aspect-ratio TRUE;display "USEDEF";valid_file
"T";width 2.8091in;height 1.8636in;depth 0pt;original-width
2.8339in;original-height 1.8707in;cropleft "0";croptop "1";cropright
"1";cropbottom "0";tempfilename 'SYUH880S.wmf';tempfile-properties "XPR";}}

Let's go back to energy transfer by heat. We know 
\begin{equation*}
Q=nC_{v}\Delta T
\end{equation*}%
if the volume doesn't change, we now can see a way to determine $Q$ from a
PV\ diagram. \FRAME{dtbpF}{3.1427in}{2.0989in}{0pt}{}{}{Figure}{\special%
{language "Scientific Word";type "GRAPHIC";maintain-aspect-ratio
TRUE;display "USEDEF";valid_file "T";width 3.1427in;height 2.0989in;depth
0pt;original-width 3.589in;original-height 2.3869in;cropleft "0";croptop
"1";cropright "1";cropbottom "0";tempfilename
'SYU9XO0A.wmf';tempfile-properties "XPR";}}But this is just for a process
where volume doesn't change. $Q$ must be different for other paths. Still,
our PV\ diagram does seem to be useful for thinking about thermodynamic
process and for filling in our first law equation 
\begin{equation*}
\Delta E_{int}=Q+w_{int}
\end{equation*}%
We will work with this some more in our next lecture.

\chapter{33 Thermodynamic Processes 2.3.4, 2.3.5}

We found in our last lecture that the first law of thermodynamics gives us
the change in internal energy 
\begin{equation*}
\Delta E_{int}=Q+w_{int}
\end{equation*}%
and that we could illustrate the work, $w_{int},$ for a process on a PV
diagram and we even found a way to show $\Delta E_{int}$ on a PV\ diagram.
We did find a way to show $Q$ for the special case of a process that doesn't
change volume. But that seems so specific. Aren't there infinite different
paths from a specific initial internal energy, $E_{i_{int}}$ to a final
internal energy, $E_{f_{int}}$? Of course, the answer is yes. But many of
these paths from $E_{i_{int}}$ to $E_{f_{int}}$ would require a numerical
solution (this means use a computer to do the calculation). So instead, we
are going to learn a few special processes that are easy to calculate by
hand. We will limit our study mostly to these special processes. In industry
you might need to go beyond these simple processes, but they give us great
insight into how internal energy changes and we can even use them (mostly)
to understand complicated systems like gasoline engines. If you need
something else, you can numerically do the work integral. So let's get
started and define our special processes.

%TCIMACRO{%
%\TeXButton{Fundamental Concepts}{\vspace{0.10in}\hspace{-0.37in}{\Large Fundamental Concepts\vspace{0.2in}}}}%
%BeginExpansion
\vspace{0.10in}\hspace{-0.37in}{\Large Fundamental Concepts\vspace{0.2in}}%
%EndExpansion

\begin{itemize}
\item There is a sign convention for energy transfer. If energy leaves it is
negative, if it comes into the system it is positive.

\item There are two special process molar specific heats for gasses, $C_{V}$
and $C_{P}.$

\item The two molar specific heats are related $C_{P}=C_{V}+R$

\item There are four special processes for which it is easy to find $\Delta
E_{int},$ $Q,$ and $w,$ and three of them are isochoric, isobaric, and
isothermal

\item For an isochoric process $w_{int}=0,$ $Q=nC_{V}\Delta T,$ $\Delta
E_{int}=nC_{V}\Delta T$

\item For an isobaric process $w_{int}=-P\Delta 
%TCIMACRO{\TeXButton{V}{\ooalign{\hfil$V$\hfil\cr\kern0.1em--\hfil\cr}}}%
%BeginExpansion
\ooalign{\hfil$V$\hfil\cr\kern0.1em--\hfil\cr}%
%EndExpansion
,$ $Q=nC_{P}\Delta T,$ $\Delta E_{int}=nC_{P}\Delta T-P\Delta 
%TCIMACRO{\TeXButton{V}{\ooalign{\hfil$V$\hfil\cr\kern0.1em--\hfil\cr}}}%
%BeginExpansion
\ooalign{\hfil$V$\hfil\cr\kern0.1em--\hfil\cr}%
%EndExpansion
$

\item For an isothermal process $w=-nRT\left[ \ln \frac{%
%TCIMACRO{\TeXButton{V}{\ooalign{\hfil$V$\hfil\cr\kern0.1em--\hfil\cr}}}%
%BeginExpansion
\ooalign{\hfil$V$\hfil\cr\kern0.1em--\hfil\cr}%
%EndExpansion
_{f}}{%
%TCIMACRO{\TeXButton{V}{\ooalign{\hfil$V$\hfil\cr\kern0.1em--\hfil\cr}}}%
%BeginExpansion
\ooalign{\hfil$V$\hfil\cr\kern0.1em--\hfil\cr}%
%EndExpansion
_{i}}\right] ,$ $Q=-w,$ $\Delta E_{int}=0$
\end{itemize}

\section{Heat and systems}

Heat is a transfer of energy, so it can't be a property of the system,
itself. It is an interaction of the system with its environment. This means
that heat is not a state variable. It will depend on the process path on a
PV-diagram, like work does. Only the combination 
\begin{equation*}
\Delta E_{int}=Q+w
\end{equation*}%
is path dependent.

We have talked about this before, but it is worth reviewing. It is customary
to refer to energy gained from an environment as positive, and energy lost
to the environment as negative.\FRAME{dtbpF}{2.095in}{2.1154in}{0pt}{}{}{%
Figure}{\special{language "Scientific Word";type
"GRAPHIC";maintain-aspect-ratio TRUE;display "USEDEF";valid_file "T";width
2.095in;height 2.1154in;depth 0pt;original-width 2.1066in;original-height
2.127in;cropleft "0";croptop "1";cropright "1";cropbottom "0";tempfilename
'SYUG5M0Q.wmf';tempfile-properties "XPR";}} This is a sign convention. We
could totally define this the other way. But it is useful for everyone to
agree to one system of signs, and the system used in physics today is that
energy is negative when it leaves a system, and positive when it enters a
system. We will need to remember this for our problems.

It is useful to compare the sign convention for heat with that we built for
work (on the gas).

\begin{equation*}
\begin{tabular}{|c|c|c|}
\hline
& \textbf{Work} & \textbf{Heat} \\ \hline
\begin{tabular}{l}
{\small Interaction} \\ 
{\small mechanism}%
\end{tabular}
& 
\begin{tabular}{l}
{\small Mechanical} \\ 
\end{tabular}
& 
\begin{tabular}{l}
{\small Thermal} \\ 
\end{tabular}
\\ \hline
\begin{tabular}{l}
{\small Process} \\ 
\end{tabular}
& 
\begin{tabular}{l}
{\small Macroscopic forces} \\ 
{\small acting on the system}%
\end{tabular}
& 
\begin{tabular}{l}
{\small Microscopic collisions} \\ 
{\small between gas particles in the system}%
\end{tabular}
\\ \hline
\begin{tabular}{l}
{\small Process} \\ 
{\small Requires}%
\end{tabular}
& 
\begin{tabular}{l}
{\small External force} \\ 
{\small and displacement}%
\end{tabular}
& 
\begin{tabular}{l}
{\small Temperature } \\ 
{\small Difference}%
\end{tabular}
\\ \hline
\begin{tabular}{l}
{\small Positive } \\ 
{\small when}%
\end{tabular}
& 
\begin{tabular}{l}
{\small A gas is compressed: mechanical} \\ 
{\small energy is transferred in to system}%
\end{tabular}
& 
\begin{tabular}{l}
{\small The environment is at a higher } \\ 
{\small temperature than the system}%
\end{tabular}
\\ \hline
\begin{tabular}{l}
{\small Negative} \\ 
{\small when}%
\end{tabular}
& 
\begin{tabular}{l}
{\small A gas expands: mechanical} \\ 
{\small energy is transferred out of the system}%
\end{tabular}
& 
\begin{tabular}{l}
{\small The environment is at a lower } \\ 
{\small temperature than the system}%
\end{tabular}
\\ \hline
\begin{tabular}{l}
{\small Equilibrium} \\ 
{\small condition}%
\end{tabular}
& 
\begin{tabular}{l}
{\small No net force or torque} \\ 
{\small on the system}%
\end{tabular}
& 
\begin{tabular}{l}
{\small System and environment are } \\ 
{\small at the same temperature}%
\end{tabular}
\\ \hline
\end{tabular}%
\end{equation*}

\bigskip It is really important to not use the word \textquotedblleft
heat\textquotedblright\ as we use the word with \textquotedblleft
temperature.\textquotedblright\ Temperature is proportional to the internal
energy, 
\begin{equation*}
E_{int}=\frac{3}{2}k_{B}T
\end{equation*}%
A gas can \textquotedblleft have\textquotedblright\ a temperature our an
amount of internal energy, but it \textquotedblleft
experiences\textquotedblright\ a transfer of energy. That is why it shows up
in the first law as a mechanism for the change in internal energy.%
\begin{equation*}
\Delta E_{int}=Q+w
\end{equation*}%
This is not intuitive because of our left-over language from the caloric
theory. So we must continually be on our guard, because poor language leads
to poor decisions in experiments and engineering designs.

It would help if we all said things like \textquotedblleft to do
heat\textquotedblright\ instead of \textquotedblleft heat\textquotedblright\
when we are warming things. That is what we do with work. We
\textquotedblleft do work\textquotedblright\ so we should \textquotedblleft
do heat,\textquotedblright\ and this may be a good way to think about it,
even though to say \textquotedblleft do heat\textquotedblright\ like this
would make our roommates wonder.

We now know that $E_{int}$ is a state variable. and $\Delta E_{int}$ does
not depend on path. But it is just one of many state variable for a gas. For
ideal gasses, we have 
\begin{equation*}
P,%
%TCIMACRO{\TeXButton{V}{\ooalign{\hfil$V$\hfil\cr\kern0.1em--\hfil\cr}}}%
%BeginExpansion
\ooalign{\hfil$V$\hfil\cr\kern0.1em--\hfil\cr}%
%EndExpansion
,n,T,E_{int}
\end{equation*}%
all as state variables. It takes all of these quantities to describe an
ideal gas.

For example, we know that%
\begin{equation*}
E_{int}=\frac{3}{2}k_{B}T
\end{equation*}%
but this does not tell us $\Delta E_{int}$ for a thermodynamic process. The
quantity $\Delta E_{int}$ only tells us about the change in $E_{int}$. We
cannot tell what path was used to accomplish $\Delta E_{int}.$ Our first law
does not tell us $P,$ $%
%TCIMACRO{\TeXButton{V}{\ooalign{\hfil$V$\hfil\cr\kern0.1em--\hfil\cr}}}%
%BeginExpansion
\ooalign{\hfil$V$\hfil\cr\kern0.1em--\hfil\cr}%
%EndExpansion
,$ or $n$ either. We will have to use the first law in cooperation with the
ideal gas law. So let's do that. Let's define four special processes as
examples. We will also be interested in $Q$ and $w$ for these processes. In
this lecture we will do three processes, and in the next lecture we will
study the fourth.

\subsection{Constant Volume Process}

We ended our last lecture with a constant volume process. We know that a
constant volume process is called an \emph{isochoric process.} Suppose we
have a process that takes us between two states on a $P%
%TCIMACRO{\TeXButton{V}{\ooalign{\hfil$V$\hfil\cr\kern0.1em--\hfil\cr}}}%
%BeginExpansion
\ooalign{\hfil$V$\hfil\cr\kern0.1em--\hfil\cr}%
%EndExpansion
\ $diagram as shown in the next figure. \FRAME{dtbpF}{3.2655in}{2.1854in}{0pt%
}{}{}{Figure}{\special{language "Scientific Word";type
"GRAPHIC";maintain-aspect-ratio TRUE;display "USEDEF";valid_file "T";width
3.2655in;height 2.1854in;depth 0pt;original-width 4.3189in;original-height
2.8807in;cropleft "0";croptop "1";cropright "1";cropbottom "0";tempfilename
'TCKIEW00.wmf';tempfile-properties "XPR";}}We start at one pressure, say $%
P=275\unit{Pa},$ and end at another pressure, say, $P=775\unit{Pa}.$ In this
process, notice that the volume does not change. Would we predict that the
temperature went up or down in taking this path? Well, we went from the
lower to the upper part of the graph, so we expect that the temperature went
up. We could check this by plotting the two isotherms that cross through our
points. \FRAME{dtbpF}{3.301in}{2.2027in}{0pt}{}{}{Figure}{\special{language
"Scientific Word";type "GRAPHIC";maintain-aspect-ratio TRUE;display
"USEDEF";valid_file "T";width 3.301in;height 2.2027in;depth
0pt;original-width 3.525in;original-height 2.3419in;cropleft "0";croptop
"1";cropright "1";cropbottom "0";tempfilename
'TCKIFJ01.wmf';tempfile-properties "XPR";}}Sure enough, our temperature went
up. Starting again with the ideal gas law%
\begin{equation*}
P=\left( \frac{nR}{%
%TCIMACRO{\TeXButton{V}{\ooalign{\hfil$V$\hfil\cr\kern0.1em--\hfil\cr}}}%
%BeginExpansion
\ooalign{\hfil$V$\hfil\cr\kern0.1em--\hfil\cr}%
%EndExpansion
}\right) T
\end{equation*}%
we can see that when $%
%TCIMACRO{\TeXButton{V}{\ooalign{\hfil$V$\hfil\cr\kern0.1em--\hfil\cr}}}%
%BeginExpansion
\ooalign{\hfil$V$\hfil\cr\kern0.1em--\hfil\cr}%
%EndExpansion
$ is constant, $P$ goes like $T,$ so if $P$ went up, so did the temperature.

For this isochoric type process no work is done.\smallskip\ 
\begin{equation*}
w_{\text{on the gas}}=w_{int}=-\int_{%
%TCIMACRO{\TeXButton{V}{\ooalign{\hfil$V$\hfil\cr\kern0.1em--\hfil\cr}}}%
%BeginExpansion
\ooalign{\hfil$V$\hfil\cr\kern0.1em--\hfil\cr}%
%EndExpansion
_{i}}^{%
%TCIMACRO{\TeXButton{V}{\ooalign{\hfil$V$\hfil\cr\kern0.1em--\hfil\cr}}}%
%BeginExpansion
\ooalign{\hfil$V$\hfil\cr\kern0.1em--\hfil\cr}%
%EndExpansion
_{f}}Pd%
%TCIMACRO{\TeXButton{V}{\ooalign{\hfil$V$\hfil\cr\kern0.1em--\hfil\cr}}}%
%BeginExpansion
\ooalign{\hfil$V$\hfil\cr\kern0.1em--\hfil\cr}%
%EndExpansion
\end{equation*}%
but $%
%TCIMACRO{\TeXButton{V}{\ooalign{\hfil$V$\hfil\cr\kern0.1em--\hfil\cr}}}%
%BeginExpansion
\ooalign{\hfil$V$\hfil\cr\kern0.1em--\hfil\cr}%
%EndExpansion
$ doesn't change so $d%
%TCIMACRO{\TeXButton{V}{\ooalign{\hfil$V$\hfil\cr\kern0.1em--\hfil\cr}}}%
%BeginExpansion
\ooalign{\hfil$V$\hfil\cr\kern0.1em--\hfil\cr}%
%EndExpansion
=0.$ and therefore 
\begin{equation*}
w_{int}=0
\end{equation*}%
But we do have energy transfer by heat%
\begin{equation*}
Q=nC_{v}\Delta T
\end{equation*}%
so for our first law 
\begin{eqnarray*}
\Delta E_{int} &=&Q+w \\
&=&nC_{v}\Delta T+0 \\
&=&nC_{v}\Delta T
\end{eqnarray*}%
That seems useful! But how would create such a situation? Consider what
would happen\footnote{%
Only consider this, because it would be irresponsible and dangerous (and
illegal in some states) to do this!} if you placed an aerosol can in a fire.
The aerosol can is rigid, so the volume can't change. But the fire will make
the temperature change, and the pressure will change--that is--until the can
explodes.

We will also need our differential form for the first law for isochoric
processes. 
\begin{equation*}
dE_{V_{int}}=nC_{v}dT
\end{equation*}

\subsection{Constant Pressure Process}

But there is more than one way to add energy to a gas. Suppose we envision
another process where instead of keeping the volume constant we keep the
pressure constant. \FRAME{dtbpF}{2.9723in}{1.9842in}{0in}{}{}{Figure}{%
\special{language "Scientific Word";type "GRAPHIC";maintain-aspect-ratio
TRUE;display "USEDEF";valid_file "T";width 2.9723in;height 1.9842in;depth
0in;original-width 3.0007in;original-height 1.993in;cropleft "0";croptop
"1";cropright "1";cropbottom "0";tempfilename
'SZ3HOQ03.wmf';tempfile-properties "XPR";}}We have already considered this
process. We called in an \emph{isobaric process}. We found that the work
done in an isobaric process is 
\begin{equation*}
w_{\text{on the gas}}=-P\left( 
%TCIMACRO{\TeXButton{V}{\ooalign{\hfil$V$\hfil\cr\kern0.1em--\hfil\cr}}}%
%BeginExpansion
\ooalign{\hfil$V$\hfil\cr\kern0.1em--\hfil\cr}%
%EndExpansion
_{f}-%
%TCIMACRO{\TeXButton{V}{\ooalign{\hfil$V$\hfil\cr\kern0.1em--\hfil\cr}}}%
%BeginExpansion
\ooalign{\hfil$V$\hfil\cr\kern0.1em--\hfil\cr}%
%EndExpansion
_{i}\right)
\end{equation*}%
Let's review that here 
\begin{eqnarray*}
w_{int} &=&-\int_{%
%TCIMACRO{\TeXButton{V}{\ooalign{\hfil$V$\hfil\cr\kern0.1em--\hfil\cr}}}%
%BeginExpansion
\ooalign{\hfil$V$\hfil\cr\kern0.1em--\hfil\cr}%
%EndExpansion
_{i}}^{%
%TCIMACRO{\TeXButton{V}{\ooalign{\hfil$V$\hfil\cr\kern0.1em--\hfil\cr}}}%
%BeginExpansion
\ooalign{\hfil$V$\hfil\cr\kern0.1em--\hfil\cr}%
%EndExpansion
_{f}}Pd%
%TCIMACRO{\TeXButton{V}{\ooalign{\hfil$V$\hfil\cr\kern0.1em--\hfil\cr}} }%
%BeginExpansion
\ooalign{\hfil$V$\hfil\cr\kern0.1em--\hfil\cr}
%EndExpansion
\\
&=&-P\int_{%
%TCIMACRO{\TeXButton{V}{\ooalign{\hfil$V$\hfil\cr\kern0.1em--\hfil\cr}}}%
%BeginExpansion
\ooalign{\hfil$V$\hfil\cr\kern0.1em--\hfil\cr}%
%EndExpansion
_{i}}^{%
%TCIMACRO{\TeXButton{V}{\ooalign{\hfil$V$\hfil\cr\kern0.1em--\hfil\cr}}}%
%BeginExpansion
\ooalign{\hfil$V$\hfil\cr\kern0.1em--\hfil\cr}%
%EndExpansion
_{f}}d%
%TCIMACRO{\TeXButton{V}{\ooalign{\hfil$V$\hfil\cr\kern0.1em--\hfil\cr}} }%
%BeginExpansion
\ooalign{\hfil$V$\hfil\cr\kern0.1em--\hfil\cr}
%EndExpansion
\\
&=&-P\left( 
%TCIMACRO{\TeXButton{V}{\ooalign{\hfil$V$\hfil\cr\kern0.1em--\hfil\cr}}}%
%BeginExpansion
\ooalign{\hfil$V$\hfil\cr\kern0.1em--\hfil\cr}%
%EndExpansion
_{f}-%
%TCIMACRO{\TeXButton{V}{\ooalign{\hfil$V$\hfil\cr\kern0.1em--\hfil\cr}}}%
%BeginExpansion
\ooalign{\hfil$V$\hfil\cr\kern0.1em--\hfil\cr}%
%EndExpansion
_{i}\right) \\
&=&-P\Delta 
%TCIMACRO{\TeXButton{V}{\ooalign{\hfil$V$\hfil\cr\kern0.1em--\hfil\cr}}}%
%BeginExpansion
\ooalign{\hfil$V$\hfil\cr\kern0.1em--\hfil\cr}%
%EndExpansion
\end{eqnarray*}%
We can see that in this process the pressure is not changing. If we go from
left to right along this path, will the temperature rise? We know that going
to the right means higher temperature on a $P%
%TCIMACRO{\TeXButton{V}{\ooalign{\hfil$V$\hfil\cr\kern0.1em--\hfil\cr}}}%
%BeginExpansion
\ooalign{\hfil$V$\hfil\cr\kern0.1em--\hfil\cr}%
%EndExpansion
$ diagram, so yes it will.

\FRAME{dtbpF}{3.0761in}{2.0557in}{0pt}{}{}{Figure}{\special{language
"Scientific Word";type "GRAPHIC";maintain-aspect-ratio TRUE;display
"USEDEF";valid_file "T";width 3.0761in;height 2.0557in;depth
0pt;original-width 3.6375in;original-height 2.4206in;cropleft "0";croptop
"1";cropright "1";cropbottom "0";tempfilename
'SZ3HR706.wmf';tempfile-properties "XPR";}}

For isochoric processes we defined a molar heat capacity at constant volume $%
C_{V}$ that let us find $Q$%
\begin{equation*}
Q=nC_{%
%TCIMACRO{\TeXButton{V}{\ooalign{\hfil$V$\hfil\cr\kern0.1em--\hfil\cr}}}%
%BeginExpansion
\ooalign{\hfil$V$\hfil\cr\kern0.1em--\hfil\cr}%
%EndExpansion
}\Delta T
\end{equation*}%
In our differential first law equation 
\begin{equation*}
dE_{int}=dQ+dw
\end{equation*}%
and we found that for isochoric processes no work was done so 
\begin{equation*}
dE_{int}=dQ=nC_{%
%TCIMACRO{\TeXButton{V}{\ooalign{\hfil$V$\hfil\cr\kern0.1em--\hfil\cr}}}%
%BeginExpansion
\ooalign{\hfil$V$\hfil\cr\kern0.1em--\hfil\cr}%
%EndExpansion
}dT
\end{equation*}%
but for an isobaric process we have 
\begin{equation*}
dE_{int}=dQ-Pd%
%TCIMACRO{\TeXButton{V}{\ooalign{\hfil$V$\hfil\cr\kern0.1em--\hfil\cr}}}%
%BeginExpansion
\ooalign{\hfil$V$\hfil\cr\kern0.1em--\hfil\cr}%
%EndExpansion
\end{equation*}%
It's not the same. What do we do for $dQ$ if the process is isobaric? It
turns out we can play the same trick. Let's define a \emph{molar heat
capacity at constant pressure}, $C_{P}.$ Then 
\begin{equation*}
Q=nC_{P}\Delta T
\end{equation*}%
and 
\begin{equation*}
dQ=nC_{P}dT
\end{equation*}%
We found that for idea gasses 
\begin{equation*}
C_{%
%TCIMACRO{\TeXButton{V}{\ooalign{\hfil$V$\hfil\cr\kern0.1em--\hfil\cr}}}%
%BeginExpansion
\ooalign{\hfil$V$\hfil\cr\kern0.1em--\hfil\cr}%
%EndExpansion
}=\frac{3}{2}R
\end{equation*}%
we need to find a value for $C_{P}.$ Let's use our ideal gas law%
\begin{equation*}
P%
%TCIMACRO{\TeXButton{V}{\ooalign{\hfil$V$\hfil\cr\kern0.1em--\hfil\cr}}}%
%BeginExpansion
\ooalign{\hfil$V$\hfil\cr\kern0.1em--\hfil\cr}%
%EndExpansion
=nRT
\end{equation*}%
and let the volume and Temperature change, but not $P$ nor $n.$%
\begin{equation*}
P\Delta 
%TCIMACRO{\TeXButton{V}{\ooalign{\hfil$V$\hfil\cr\kern0.1em--\hfil\cr}}}%
%BeginExpansion
\ooalign{\hfil$V$\hfil\cr\kern0.1em--\hfil\cr}%
%EndExpansion
=nR\Delta T
\end{equation*}%
For an isobaric process, this is our work! Well, minus our work. We can put
this into our first law equation 
\begin{eqnarray*}
\Delta E_{int} &=&Q-\left( w\right) \\
&=&Q-\left( nR\Delta T\right)
\end{eqnarray*}%
and we can put in $Q=nC_{P}\Delta T$%
\begin{eqnarray*}
\Delta E_{int} &=&\left( nC_{P}\Delta T\right) -nR\Delta T \\
&=&n(C_{P}-R)\Delta T
\end{eqnarray*}%
But $\Delta E_{int}$ is the same for any two processes that go between the
same two temperatures. Suppose we look at a situation where we have the same 
$\Delta E_{int}$ but we make the change two ways, one with an isochoric
process and one with an isobaric process.

\FRAME{dtbpFX}{2.2883in}{1.5264in}{0pt}{}{}{Plot}{\special{language
"Scientific Word";type "MAPLEPLOT";width 2.2883in;height 1.5264in;depth
0pt;display "USEDEF";plot_snapshots TRUE;mustRecompute FALSE;lastEngine
"MuPAD";xmin "2";xmax "10";xviewmin "2";xviewmax "10";yviewmin
"100";yviewmax "1100";viewset"XY";rangeset"X";plottype 4;labeloverrides
3;x-label "V(m^3)";y-label "P(Pa)";axesFont "Times New
Roman,12,0000000000,useDefault,normal";numpoints 100;plotstyle
"patch";axesstyle "normal";axestips FALSE;xis \TEXUX{V};var1name
\TEXUX{$V$};function
\TEXUX{$\MATRIX{2,2}{c}\VR{,,c,,,}{,,c,,,}{,,,,,}\HR{,,}\CELL{4}\CELL{%
\allowbreak 500}\CELL{8}\CELL{500}$};linecolor "green";linestyle
1;pointstyle "circle";linethickness 1;lineAttributes "Solid";curveColor
"[flat::RGB:0x00008000]";curveStyle "Line";function
\TEXUX{$\MATRIX{2,2}{c}\VR{,,c,,,}{,,c,,,}{,,,,,}\HR{,,}\CELL{4}\CELL{%
\allowbreak 500}\CELL{8}\CELL{500}$};linecolor "green";linestyle 1;pointplot
TRUE;pointstyle "circle";linethickness 1;lineAttributes "Solid";curveColor
"[flat::RGB:0x00008000]";curveStyle "Point";function \TEXUX{$\frac{1\left(
8.314\right) \left( 240.\,\allowbreak 56\right) }{V}$};linecolor
"blue";linestyle 2;discont FALSE;pointstyle "circle";linethickness
1;lineAttributes "Dash";var1range "2,10";num-x-gridlines 100;curveColor
"[flat::RGB:0x000000ff]";curveStyle "Line";discont FALSE;function
\TEXUX{$\frac{1\left( 8.314\right) \left( 481.\,\allowbreak 12\right)
}{V}$};linecolor "blue";linestyle 2;pointstyle "circle";linethickness
1;lineAttributes "Dash";var1range "2,10";num-x-gridlines 100;curveColor
"[flat::RGB:0x000000ff]";curveStyle "Line";function
\TEXUX{$\MATRIX{2,2}{c}\VR{,,c,,,}{,,c,,,}{,,,,,}\HR{,,}\CELL{4}\CELL{%
\allowbreak 500}\CELL{4}\CELL{1000.0}$};linecolor "green";linestyle
1;pointstyle "circle";linethickness 1;lineAttributes "Solid";curveColor
"[flat::RGB:0x00008000]";curveStyle "Line";function
\TEXUX{$\MATRIX{2,2}{c}\VR{,,c,,,}{,,c,,,}{,,,,,}\HR{,,}\CELL{4}\CELL{%
\allowbreak 500}\CELL{4}\CELL{1000.0}$};linecolor "green";linestyle
1;pointplot TRUE;pointstyle "circle";linethickness 1;lineAttributes
"Solid";curveColor "[flat::RGB:0x00008000]";curveStyle "Point";VCamFile
'SYUCLQ10.xvz';valid_file "T";tempfilename
'SYUCC90E.wmf';tempfile-properties "XPR";}}We already know that for the
isochoric process 
\begin{equation*}
\Delta E_{%
%TCIMACRO{\TeXButton{V}{\ooalign{\hfil$V$\hfil\cr\kern0.1em--\hfil\cr}}}%
%BeginExpansion
\ooalign{\hfil$V$\hfil\cr\kern0.1em--\hfil\cr}%
%EndExpansion
_{int}}=nC_{v}\Delta T
\end{equation*}%
Then since 
\begin{equation*}
\Delta E_{%
%TCIMACRO{\TeXButton{V}{\ooalign{\hfil$V$\hfil\cr\kern0.1em--\hfil\cr}}}%
%BeginExpansion
\ooalign{\hfil$V$\hfil\cr\kern0.1em--\hfil\cr}%
%EndExpansion
_{int}}=\Delta E_{P_{int}}
\end{equation*}%
then 
\begin{equation*}
nC_{%
%TCIMACRO{\TeXButton{V}{\ooalign{\hfil$V$\hfil\cr\kern0.1em--\hfil\cr}}}%
%BeginExpansion
\ooalign{\hfil$V$\hfil\cr\kern0.1em--\hfil\cr}%
%EndExpansion
}\Delta T=n(C_{P}-R)\Delta T
\end{equation*}%
And we can cancel the $n$'s and the $\Delta T$'s.%
\begin{equation*}
C_{%
%TCIMACRO{\TeXButton{V}{\ooalign{\hfil$V$\hfil\cr\kern0.1em--\hfil\cr}}}%
%BeginExpansion
\ooalign{\hfil$V$\hfil\cr\kern0.1em--\hfil\cr}%
%EndExpansion
}=(C_{P}-R)
\end{equation*}%
which gives

\begin{equation*}
C_{P}=C_{%
%TCIMACRO{\TeXButton{V}{\ooalign{\hfil$V$\hfil\cr\kern0.1em--\hfil\cr}}}%
%BeginExpansion
\ooalign{\hfil$V$\hfil\cr\kern0.1em--\hfil\cr}%
%EndExpansion
}+R
\end{equation*}%
for an ideal gas. So if we know either $C_{P}$ or $C_{%
%TCIMACRO{\TeXButton{V}{\ooalign{\hfil$V$\hfil\cr\kern0.1em--\hfil\cr}}}%
%BeginExpansion
\ooalign{\hfil$V$\hfil\cr\kern0.1em--\hfil\cr}%
%EndExpansion
}$ for a gas, we can always find the other. This works pretty well in
practice for monotonic gasses (and even for diatomic gasses) at standard
temperature and pressure. Not only can we find $Q_{%
%TCIMACRO{\TeXButton{V}{\ooalign{\hfil$V$\hfil\cr\kern0.1em--\hfil\cr}}}%
%BeginExpansion
\ooalign{\hfil$V$\hfil\cr\kern0.1em--\hfil\cr}%
%EndExpansion
}$ and $Q_{P}$ with molar specific heats, but if we know one molar specific
heat, we can easily calculate the other! Since%
\begin{equation*}
C_{%
%TCIMACRO{\TeXButton{V}{\ooalign{\hfil$V$\hfil\cr\kern0.1em--\hfil\cr}}}%
%BeginExpansion
\ooalign{\hfil$V$\hfil\cr\kern0.1em--\hfil\cr}%
%EndExpansion
}=\frac{3}{2}R
\end{equation*}%
then we know 
\begin{equation*}
C_{P}=\frac{5}{2}R
\end{equation*}

So for an isobaric process we have 
\begin{eqnarray*}
\Delta E_{int} &=&Q+w \\
&=&nC_{P}\Delta T-P\Delta 
%TCIMACRO{\TeXButton{V}{\ooalign{\hfil$V$\hfil\cr\kern0.1em--\hfil\cr}}}%
%BeginExpansion
\ooalign{\hfil$V$\hfil\cr\kern0.1em--\hfil\cr}%
%EndExpansion
\end{eqnarray*}

How would we produce such a situation? We saw this in our last lecture, but
let's think about it again. Consider a piston in a cylinder. The cylinder
has a gas inside, and air on the outside\FRAME{dhF}{1.8308in}{1.5212in}{0pt}{%
}{}{Figure}{\special{language "Scientific Word";type
"GRAPHIC";maintain-aspect-ratio TRUE;display "USEDEF";valid_file "T";width
1.8308in;height 1.5212in;depth 0pt;original-width 1.7936in;original-height
1.4849in;cropleft "0";croptop "1";cropright "1";cropbottom "0";tempfilename
'SYUEEH0O.wmf';tempfile-properties "XPR";}}The piston is free to move up or
down. The forces acting on the piston are shown. The weight of the piston
pulls it downward, the force due to pressure of the gas, $F_{p,gas}$ pushes
up, and the air pressure $F_{p,air}$ pushes down. The net force will be 
\begin{equation*}
\Sigma F_{y}=ma_{y}=F_{p,gas}-F_{p,air}-mg
\end{equation*}%
where the mass of the piston is $m.$ We can write this as 
\begin{equation*}
ma_{y}=P_{gas}A-P_{air}A-mg
\end{equation*}%
If the piston is not accelerating , then 
\begin{equation*}
P_{gas}=P_{air}+\frac{mg}{A}
\end{equation*}%
and we find that the pressure of the gas inside is slightly larger than the
outside air pressure. What would happen if we heated the gas? The equation 
\begin{equation*}
P%
%TCIMACRO{\TeXButton{V}{\ooalign{\hfil$V$\hfil\cr\kern0.1em--\hfil\cr}}}%
%BeginExpansion
\ooalign{\hfil$V$\hfil\cr\kern0.1em--\hfil\cr}%
%EndExpansion
=NRT
\end{equation*}%
would suggest that $P$ might change, but since the piston is free to move,
what happens is that the force $F_{p,gas}$ increases when $P_{gas}$
increases, and the piston accelerates upward. It stops when the forces are
again balanced. But that means that $P_{gas}$ will only change for a moment,
and then we will achieve equilibrium at the same old value of $P_{gas}.$ But
something has changed, the volume has increased.

\section{Constant Temperature Process}

We have already done quite a bit of work on a constant temperature process
in our last lecture \FRAME{dtbpFX}{2.4025in}{1.6016in}{0pt}{}{}{Plot}{%
\special{language "Scientific Word";type "MAPLEPLOT";width 2.4025in;height
1.6016in;depth 0pt;display "USEDEF";plot_snapshots TRUE;mustRecompute
FALSE;lastEngine "MuPAD";animated TRUE;animationStartTime "0";
animationEndTime "10"; animationFramesPerSecond "10";xmin "0";xmax
"40";animationParamMin "290";animationParamMax "350";xviewmin "0";xviewmax
"40";yviewmin "0";yviewmax "400";viewset"XY";rangeset"X";plottype
4;labeloverrides 3;x-label "V(m^3)";y-label "P(Pa)";axesFont "Times New
Roman,12,0000000000,useDefault,normal";numpoints 100;plotstyle
"patch";axesstyle "normal";axestips FALSE;xis \TEXUX{V};animationParam
\TEXUX{t};var1name \TEXUX{$V$};animationParamList \TEXUX{$t$};function
\TEXUX{$\allowbreak 8.\,\allowbreak 314\frac{1}{V}290$};linecolor
"blue";linestyle 1;pointstyle "point";linethickness 3;lineAttributes
"Solid";var1range "0,40";animationParamRange "290,350";num-x-gridlines
100;curveColor "[flat::RGB:0x000000ff]";curveStyle "Line";VCamFile
'SYUCQZ13.xvz';valid_file "T";tempfilename
'SYUCQZ0F.wmf';tempfile-properties "XPR";}}We found that $\Delta E_{int}=0$
for an isothermal process. Then 
\begin{equation*}
\Delta E_{int}=Q+w
\end{equation*}%
gives 
\begin{equation*}
Q=-w
\end{equation*}%
But now we need to find $w_{int}$. Let's use our ideal gas law for this
again.

In this case we know $T$ and $n$ are constant, and that suppose we also know 
$P_{i}$ and $%
%TCIMACRO{\TeXButton{V}{\ooalign{\hfil$V$\hfil\cr\kern0.1em--\hfil\cr}}}%
%BeginExpansion
\ooalign{\hfil$V$\hfil\cr\kern0.1em--\hfil\cr}%
%EndExpansion
_{i}$ and $P_{f}$ and $%
%TCIMACRO{\TeXButton{V}{\ooalign{\hfil$V$\hfil\cr\kern0.1em--\hfil\cr}}}%
%BeginExpansion
\ooalign{\hfil$V$\hfil\cr\kern0.1em--\hfil\cr}%
%EndExpansion
_{f}.$ Can we find $\Delta E_{int,}$ $Q$, and $w_{int}$?

We already know how to find $w_{int},$ we integrate. We will need to know
how $P$ varies with $%
%TCIMACRO{\TeXButton{V}{\ooalign{\hfil$V$\hfil\cr\kern0.1em--\hfil\cr}}}%
%BeginExpansion
\ooalign{\hfil$V$\hfil\cr\kern0.1em--\hfil\cr}%
%EndExpansion
$ but once again our ideal gas law helps. 
\begin{equation*}
P=\frac{nRT}{V}
\end{equation*}%
Then, from our work equation%
\begin{eqnarray}
w_{int} &=&-\int_{%
%TCIMACRO{\TeXButton{V}{\ooalign{\hfil$V$\hfil\cr\kern0.1em--\hfil\cr}}}%
%BeginExpansion
\ooalign{\hfil$V$\hfil\cr\kern0.1em--\hfil\cr}%
%EndExpansion
_{i}}^{%
%TCIMACRO{\TeXButton{V}{\ooalign{\hfil$V$\hfil\cr\kern0.1em--\hfil\cr}}}%
%BeginExpansion
\ooalign{\hfil$V$\hfil\cr\kern0.1em--\hfil\cr}%
%EndExpansion
_{f}}Pd%
%TCIMACRO{\TeXButton{V}{\ooalign{\hfil$V$\hfil\cr\kern0.1em--\hfil\cr}} }%
%BeginExpansion
\ooalign{\hfil$V$\hfil\cr\kern0.1em--\hfil\cr}
%EndExpansion
\\
&=&-\int_{%
%TCIMACRO{\TeXButton{V}{\ooalign{\hfil$V$\hfil\cr\kern0.1em--\hfil\cr}}}%
%BeginExpansion
\ooalign{\hfil$V$\hfil\cr\kern0.1em--\hfil\cr}%
%EndExpansion
_{i}}^{%
%TCIMACRO{\TeXButton{V}{\ooalign{\hfil$V$\hfil\cr\kern0.1em--\hfil\cr}}}%
%BeginExpansion
\ooalign{\hfil$V$\hfil\cr\kern0.1em--\hfil\cr}%
%EndExpansion
_{f}}\frac{nRT}{%
%TCIMACRO{\TeXButton{V}{\ooalign{\hfil$V$\hfil\cr\kern0.1em--\hfil\cr}}}%
%BeginExpansion
\ooalign{\hfil$V$\hfil\cr\kern0.1em--\hfil\cr}%
%EndExpansion
}d%
%TCIMACRO{\TeXButton{V}{\ooalign{\hfil$V$\hfil\cr\kern0.1em--\hfil\cr}} }%
%BeginExpansion
\ooalign{\hfil$V$\hfil\cr\kern0.1em--\hfil\cr}
%EndExpansion
\\
&=&-nRT\int_{%
%TCIMACRO{\TeXButton{V}{\ooalign{\hfil$V$\hfil\cr\kern0.1em--\hfil\cr}}}%
%BeginExpansion
\ooalign{\hfil$V$\hfil\cr\kern0.1em--\hfil\cr}%
%EndExpansion
_{i}}^{%
%TCIMACRO{\TeXButton{V}{\ooalign{\hfil$V$\hfil\cr\kern0.1em--\hfil\cr}}}%
%BeginExpansion
\ooalign{\hfil$V$\hfil\cr\kern0.1em--\hfil\cr}%
%EndExpansion
_{f}}\frac{1}{%
%TCIMACRO{\TeXButton{V}{\ooalign{\hfil$V$\hfil\cr\kern0.1em--\hfil\cr}}}%
%BeginExpansion
\ooalign{\hfil$V$\hfil\cr\kern0.1em--\hfil\cr}%
%EndExpansion
}d%
%TCIMACRO{\TeXButton{V}{\ooalign{\hfil$V$\hfil\cr\kern0.1em--\hfil\cr}} }%
%BeginExpansion
\ooalign{\hfil$V$\hfil\cr\kern0.1em--\hfil\cr}
%EndExpansion
\\
&=&-nRT\left. \ln 
%TCIMACRO{\TeXButton{V}{\ooalign{\hfil$V$\hfil\cr\kern0.1em--\hfil\cr}}}%
%BeginExpansion
\ooalign{\hfil$V$\hfil\cr\kern0.1em--\hfil\cr}%
%EndExpansion
\right\vert _{%
%TCIMACRO{\TeXButton{V}{\ooalign{\hfil$V$\hfil\cr\kern0.1em--\hfil\cr}}}%
%BeginExpansion
\ooalign{\hfil$V$\hfil\cr\kern0.1em--\hfil\cr}%
%EndExpansion
_{i}}^{%
%TCIMACRO{\TeXButton{V}{\ooalign{\hfil$V$\hfil\cr\kern0.1em--\hfil\cr}}}%
%BeginExpansion
\ooalign{\hfil$V$\hfil\cr\kern0.1em--\hfil\cr}%
%EndExpansion
_{f}} \\
&=&-nRT\left[ \ln \left( 
%TCIMACRO{\TeXButton{V}{\ooalign{\hfil$V$\hfil\cr\kern0.1em--\hfil\cr}}}%
%BeginExpansion
\ooalign{\hfil$V$\hfil\cr\kern0.1em--\hfil\cr}%
%EndExpansion
_{f}\right) -\ln \left( 
%TCIMACRO{\TeXButton{V}{\ooalign{\hfil$V$\hfil\cr\kern0.1em--\hfil\cr}}}%
%BeginExpansion
\ooalign{\hfil$V$\hfil\cr\kern0.1em--\hfil\cr}%
%EndExpansion
_{-}\right) \right]
\end{eqnarray}%
Using our knowledge of logarithms, we can write this as%
\begin{equation*}
w_{int}=-nRT\left[ \ln \frac{%
%TCIMACRO{\TeXButton{V}{\ooalign{\hfil$V$\hfil\cr\kern0.1em--\hfil\cr}}}%
%BeginExpansion
\ooalign{\hfil$V$\hfil\cr\kern0.1em--\hfil\cr}%
%EndExpansion
_{f}}{%
%TCIMACRO{\TeXButton{V}{\ooalign{\hfil$V$\hfil\cr\kern0.1em--\hfil\cr}}}%
%BeginExpansion
\ooalign{\hfil$V$\hfil\cr\kern0.1em--\hfil\cr}%
%EndExpansion
_{i}}\right]
\end{equation*}%
So then 
\begin{eqnarray*}
Q &=&-w \\
&=&nRT\left[ \ln \frac{%
%TCIMACRO{\TeXButton{V}{\ooalign{\hfil$V$\hfil\cr\kern0.1em--\hfil\cr}}}%
%BeginExpansion
\ooalign{\hfil$V$\hfil\cr\kern0.1em--\hfil\cr}%
%EndExpansion
_{f}}{%
%TCIMACRO{\TeXButton{V}{\ooalign{\hfil$V$\hfil\cr\kern0.1em--\hfil\cr}}}%
%BeginExpansion
\ooalign{\hfil$V$\hfil\cr\kern0.1em--\hfil\cr}%
%EndExpansion
_{i}}\right]
\end{eqnarray*}%
and we did not need our molar specific heat equations for this process.

How do we make such a path happen? Consider our piston again.\FRAME{dtbpF}{%
3.2033in}{2.0384in}{0in}{}{}{Figure}{\special{language "Scientific
Word";type "GRAPHIC";maintain-aspect-ratio TRUE;display "USEDEF";valid_file
"T";width 3.2033in;height 2.0384in;depth 0in;original-width
3.1583in;original-height 2.0003in;cropleft "0";croptop "1";cropright
"1";cropbottom "0";tempfilename 'THTEWZ00.wmf';tempfile-properties "XPR";}}%
This time, let's put it in a vat of ice water. We can compress the gas by
pushing on the piston. But if the temperature is changed by compressing the
gas, it is quickly changed back by the ice water.

We have $\Delta E_{int,}$ $Q$, and $w$ for three of our special processes.
This completes three of our four special processes!

We have one more to go.

\chapter{34 Adiabatic Processes and 2nd Law 2.3.6, 2.4.1}

%TCIMACRO{\TeXButton{\chaptermark{Chapter 36}}{\chaptermark{Chapter 36}}}%
%BeginExpansion
\chaptermark{Chapter 36}%
%EndExpansion

Last lecture we defined three standard processes that we could find
equations for to calculate $Q,$ $w_{int},$ $\Delta E_{int}$ and with the
ideal gas law we could also find $P,$ $%
%TCIMACRO{\TeXButton{V}{\ooalign{\hfil$V$\hfil\cr\kern0.1em--\hfil\cr}}}%
%BeginExpansion
\ooalign{\hfil$V$\hfil\cr\kern0.1em--\hfil\cr}%
%EndExpansion
,$ $n,$ and $T.$ These values describe the state of our system. Those
processes we called isochoric, isobaric and isothermal. Today we will add
one more standard process. And we will ask the question, is there a second
law of thermodynamics?

%TCIMACRO{%
%\TeXButton{Fundamental Concepts}{\vspace{0.10in}\hspace{-0.37in}{\Large Fundamental Concepts\vspace{0.2in}}}}%
%BeginExpansion
\vspace{0.10in}\hspace{-0.37in}{\Large Fundamental Concepts\vspace{0.2in}}%
%EndExpansion

\begin{enumerate}
\item There are four special processes and three of them are isochoric,
isobaric, and isothermal and the last one is adiabatic

\item The ratio of the molar specific heats is $\gamma =C_{P}/C_{V}$

\item The change in internal energy can be expressed as $\Delta
E_{int}=nC_{V}\Delta T$

\item For adiabatic processes we have two new equations $P_{i}%
%TCIMACRO{\TeXButton{V}{\ooalign{\hfil$V$\hfil\cr\kern0.1em--\hfil\cr}}}%
%BeginExpansion
\ooalign{\hfil$V$\hfil\cr\kern0.1em--\hfil\cr}%
%EndExpansion
_{i}^{\gamma }=P_{f}%
%TCIMACRO{\TeXButton{V}{\ooalign{\hfil$V$\hfil\cr\kern0.1em--\hfil\cr}}}%
%BeginExpansion
\ooalign{\hfil$V$\hfil\cr\kern0.1em--\hfil\cr}%
%EndExpansion
_{f}^{\gamma }$ and $T_{f}%
%TCIMACRO{\TeXButton{V}{\ooalign{\hfil$V$\hfil\cr\kern0.1em--\hfil\cr}}}%
%BeginExpansion
\ooalign{\hfil$V$\hfil\cr\kern0.1em--\hfil\cr}%
%EndExpansion
_{f}^{\gamma -1}=T_{i}%
%TCIMACRO{\TeXButton{V}{\ooalign{\hfil$V$\hfil\cr\kern0.1em--\hfil\cr}}}%
%BeginExpansion
\ooalign{\hfil$V$\hfil\cr\kern0.1em--\hfil\cr}%
%EndExpansion
_{i}^{\gamma -1}$
\end{enumerate}

\section{Constant $Q$ Processes}

We found that for isochoric processes there was no work done. And in our
first law equation 
\begin{equation*}
\Delta E_{int}=Q+w_{int}
\end{equation*}%
the work is half the equation. That made $\Delta E_{int}=Q.$ But could we
find a process where $Q$ is zero? Then 
\begin{equation*}
\Delta E_{int}=w_{int}
\end{equation*}%
We call a process where $Q=0$ \emph{adiabatic}.

And adiabatic process does not allow convection or conduction or radiation.
We saw a way to prevent these with our liquid nitrogen dewar. Suppose we
make another cylinder with a piston, but this time we make the cylinder
double walled with vacuum inside and we insolate and make is shiny just like
we did for our dewar. The piston could still do work, but there would be no
energy transfer by heat. \FRAME{dtbpF}{4.7167in}{1.8593in}{0in}{}{}{Figure}{%
\special{language "Scientific Word";type "GRAPHIC";maintain-aspect-ratio
TRUE;display "USEDEF";valid_file "T";width 4.7167in;height 1.8593in;depth
0in;original-width 4.6639in;original-height 1.8222in;cropleft "0";croptop
"1";cropright "1";cropbottom "0";tempfilename
'TCSAAK00.wmf';tempfile-properties "XPR";}}

It might see like we are home free, $Q=0$! We don't need our molar specific
heat equations for an adiabatic process either.

But we don't yet know how to find $w_{int}$ for an adiabatic process, and
therefore we don't know $\Delta E_{int}$. We need to fix this problem. We
could do this by finding an expression for the pressure as a function of
volume along the adiabatic path, then use our work equation. 
\begin{equation*}
w_{int}=-\int_{%
%TCIMACRO{\TeXButton{V}{\ooalign{\hfil$V$\hfil\cr\kern0.1em--\hfil\cr}}}%
%BeginExpansion
\ooalign{\hfil$V$\hfil\cr\kern0.1em--\hfil\cr}%
%EndExpansion
_{i}}^{%
%TCIMACRO{\TeXButton{V}{\ooalign{\hfil$V$\hfil\cr\kern0.1em--\hfil\cr}}}%
%BeginExpansion
\ooalign{\hfil$V$\hfil\cr\kern0.1em--\hfil\cr}%
%EndExpansion
_{f}}Pd%
%TCIMACRO{\TeXButton{V}{\ooalign{\hfil$V$\hfil\cr\kern0.1em--\hfil\cr}}}%
%BeginExpansion
\ooalign{\hfil$V$\hfil\cr\kern0.1em--\hfil\cr}%
%EndExpansion
\end{equation*}

Since%
\begin{eqnarray*}
\Delta E_{int} &=&w_{int} \\
&=&-\int Pd%
%TCIMACRO{\TeXButton{V}{\ooalign{\hfil$V$\hfil\cr\kern0.1em--\hfil\cr}}}%
%BeginExpansion
\ooalign{\hfil$V$\hfil\cr\kern0.1em--\hfil\cr}%
%EndExpansion
\end{eqnarray*}%
This would give us both $w_{int}$ and $\Delta E_{int}$ . But it is not
obvious how to find $P$ as a function of $%
%TCIMACRO{\TeXButton{V}{\ooalign{\hfil$V$\hfil\cr\kern0.1em--\hfil\cr}}}%
%BeginExpansion
\ooalign{\hfil$V$\hfil\cr\kern0.1em--\hfil\cr}%
%EndExpansion
$ for this path. We need a strategy.

Let's plot an adiabatic process on a PV\ diagram to see what it looks like.
Even this is a little hard. An adiabatic process goes from one temperature
to another temperature, and 
\begin{equation*}
\Delta E_{int}\varpropto \Delta T
\end{equation*}%
The adiabatic path looks like this (but it will be a bit before this is
obvious).\FRAME{dtbpFX}{2.2087in}{1.4719in}{0pt}{}{}{Plot}{\special{language
"Scientific Word";type "MAPLEPLOT";width 2.2087in;height 1.4719in;depth
0pt;display "USEDEF";plot_snapshots TRUE;mustRecompute FALSE;lastEngine
"MuPAD";animated TRUE;animationStartTime "0"; animationEndTime "10";
animationFramesPerSecond "10";xmin "0";xmax "100";animationParamMin
"290";animationParamMax "350";xviewmin "0";xviewmax "40";yviewmin
"0";yviewmax "400";viewset"XY";rangeset"X";plottype 4;labeloverrides
3;x-label "V(m^3)";y-label "P(Pa)";axesFont "Times New
Roman,12,0000000000,useDefault,normal";numpoints 100;plotstyle
"patch";axesstyle "normal";axestips FALSE;xis \TEXUX{V};animationParam
\TEXUX{t};var1name \TEXUX{$V$};animationParamList \TEXUX{$t$};function
\TEXUX{$\allowbreak 7500\frac{1}{V^{1.67}}$};linecolor "blue";linestyle
1;pointstyle "point";linethickness 3;lineAttributes "Solid";var1range
"0,100";animationParamRange "290,350";num-x-gridlines 50;curveColor
"[flat::RGB:0x000000ff]";curveStyle "Line";animationStartTime "0";
animationEndTime "10"; animationFramesPerSecond
"10";animationVisibleBeforeStart FALSE;rangeset"XA";function
\TEXUX{$\allowbreak 8.\,\allowbreak 314\frac{1}{V}290$};linecolor
"red";linestyle 2;pointstyle "point";linethickness 1;lineAttributes
"Dash";var1range "0,40";animationParamRange "290,350";num-x-gridlines
100;curveColor "[flat::RGB:0x00ff0000]";curveStyle "Line";function
\TEXUX{$8.\,\allowbreak 314\frac{1}{V}100$};linecolor "red";linestyle
2;pointstyle "point";linethickness 1;lineAttributes "Dash";var1range
"0,40";animationParamRange "290,350";num-x-gridlines 100;curveColor
"[flat::RGB:0x00ff0000]";curveStyle "Line";VCamFile
'SYUDAG1J.xvz';valid_file "T";tempfilename
'SYUDSZ0L.wmf';tempfile-properties "XPR";}}Further remember that $\Delta
E_{int}$ does not depend on path. So we could find $\Delta E_{int}$ for an
adiabatic path, and we could find $\Delta E_{int}$ for a isochoric path, and
as long as the $\Delta T$ was the same for both paths, the $\Delta E_{int}$
would be the same as well! 
\begin{equation*}
\left( \Delta E_{int}\right) _{adiabatic}=\left( \Delta E_{int}\right)
_{isochoric}
\end{equation*}%
Also remember that for an isochoric path there is no work done. So 
\begin{equation*}
Q_{isochoric}=\Delta E_{int}
\end{equation*}%
And we know an equation for $Q_{isochoric}$!%
\begin{equation*}
Q_{isochoric}=nC_{V}\Delta T
\end{equation*}%
So if we know $T_{i}$ and $T_{f}$ for our adiabatic path, then we could find 
$\Delta E_{int}$ by finding the equivalent $Q_{isochoric}$ for the same two
temperatures.%
\begin{equation*}
\Delta E_{int}=nC_{V}\Delta T
\end{equation*}%
This is tremendous! We have an equation for $\Delta E_{int}$, and $%
w_{int}=\Delta E_{int},$ so with $Q=0$ we have a complete set of equations.
Only, we still don't know how to plot this on a PV diagram.

It turns out that we can add two more equations to our adiabatic set that
will really help.

The first is that relationship between pressures and volumes along an
adiabatic path. We now know for adiabatic processes we can say 
\begin{equation*}
\Delta E_{int}=nC_{V}\Delta T
\end{equation*}%
so long as we start with the same $T_{i}$ and end with the same $T_{f}.$ A
very small change in the internal energy would be 
\begin{equation*}
dE=nC_{V}dT
\end{equation*}%
Since for our adiabatic process there is no energy transfer by heat done 
\begin{eqnarray*}
\Delta E_{int} &=&0+w \\
&=&-\int Pd%
%TCIMACRO{\TeXButton{V}{\ooalign{\hfil$V$\hfil\cr\kern0.1em--\hfil\cr}}}%
%BeginExpansion
\ooalign{\hfil$V$\hfil\cr\kern0.1em--\hfil\cr}%
%EndExpansion
\end{eqnarray*}%
A very small amount of internal energy change, $dE_{int}$ would be 
\begin{equation*}
dE_{int}=-Pd%
%TCIMACRO{\TeXButton{V}{\ooalign{\hfil$V$\hfil\cr\kern0.1em--\hfil\cr}}}%
%BeginExpansion
\ooalign{\hfil$V$\hfil\cr\kern0.1em--\hfil\cr}%
%EndExpansion
\end{equation*}%
Let's set these two expressions for $dE_{int}$ equal to each other. 
\begin{equation*}
-Pd%
%TCIMACRO{\TeXButton{V}{\ooalign{\hfil$V$\hfil\cr\kern0.1em--\hfil\cr}}}%
%BeginExpansion
\ooalign{\hfil$V$\hfil\cr\kern0.1em--\hfil\cr}%
%EndExpansion
=nC_{%
%TCIMACRO{\TeXButton{V}{\ooalign{\hfil$V$\hfil\cr\kern0.1em--\hfil\cr}}}%
%BeginExpansion
\ooalign{\hfil$V$\hfil\cr\kern0.1em--\hfil\cr}%
%EndExpansion
}dT
\end{equation*}%
or%
\begin{equation*}
-\frac{Pd%
%TCIMACRO{\TeXButton{V}{\ooalign{\hfil$V$\hfil\cr\kern0.1em--\hfil\cr}}}%
%BeginExpansion
\ooalign{\hfil$V$\hfil\cr\kern0.1em--\hfil\cr}%
%EndExpansion
}{nC_{%
%TCIMACRO{\TeXButton{V}{\ooalign{\hfil$V$\hfil\cr\kern0.1em--\hfil\cr}}}%
%BeginExpansion
\ooalign{\hfil$V$\hfil\cr\kern0.1em--\hfil\cr}%
%EndExpansion
}}=dT
\end{equation*}

Now we know the ideal gas law%
\begin{equation*}
P%
%TCIMACRO{\TeXButton{V}{\ooalign{\hfil$V$\hfil\cr\kern0.1em--\hfil\cr}}}%
%BeginExpansion
\ooalign{\hfil$V$\hfil\cr\kern0.1em--\hfil\cr}%
%EndExpansion
=nRT
\end{equation*}%
if we take a total differential for a sample of $n$ moles of a gas%
\begin{equation*}
Pd%
%TCIMACRO{\TeXButton{V}{\ooalign{\hfil$V$\hfil\cr\kern0.1em--\hfil\cr}}}%
%BeginExpansion
\ooalign{\hfil$V$\hfil\cr\kern0.1em--\hfil\cr}%
%EndExpansion
+%
%TCIMACRO{\TeXButton{V}{\ooalign{\hfil$V$\hfil\cr\kern0.1em--\hfil\cr}}}%
%BeginExpansion
\ooalign{\hfil$V$\hfil\cr\kern0.1em--\hfil\cr}%
%EndExpansion
dP=nRdT
\end{equation*}%
we can substitute in our expression we have for an adiabatic process for $dT$%
\begin{equation*}
Pd%
%TCIMACRO{\TeXButton{V}{\ooalign{\hfil$V$\hfil\cr\kern0.1em--\hfil\cr}}}%
%BeginExpansion
\ooalign{\hfil$V$\hfil\cr\kern0.1em--\hfil\cr}%
%EndExpansion
+%
%TCIMACRO{\TeXButton{V}{\ooalign{\hfil$V$\hfil\cr\kern0.1em--\hfil\cr}}}%
%BeginExpansion
\ooalign{\hfil$V$\hfil\cr\kern0.1em--\hfil\cr}%
%EndExpansion
dP=-nR\frac{Pd%
%TCIMACRO{\TeXButton{V}{\ooalign{\hfil$V$\hfil\cr\kern0.1em--\hfil\cr}}}%
%BeginExpansion
\ooalign{\hfil$V$\hfil\cr\kern0.1em--\hfil\cr}%
%EndExpansion
}{nC_{V}}
\end{equation*}%
or%
\begin{equation*}
Pd%
%TCIMACRO{\TeXButton{V}{\ooalign{\hfil$V$\hfil\cr\kern0.1em--\hfil\cr}}}%
%BeginExpansion
\ooalign{\hfil$V$\hfil\cr\kern0.1em--\hfil\cr}%
%EndExpansion
+%
%TCIMACRO{\TeXButton{V}{\ooalign{\hfil$V$\hfil\cr\kern0.1em--\hfil\cr}}}%
%BeginExpansion
\ooalign{\hfil$V$\hfil\cr\kern0.1em--\hfil\cr}%
%EndExpansion
dP=-R\frac{Pd%
%TCIMACRO{\TeXButton{V}{\ooalign{\hfil$V$\hfil\cr\kern0.1em--\hfil\cr}}}%
%BeginExpansion
\ooalign{\hfil$V$\hfil\cr\kern0.1em--\hfil\cr}%
%EndExpansion
}{C_{%
%TCIMACRO{\TeXButton{V}{\ooalign{\hfil$V$\hfil\cr\kern0.1em--\hfil\cr}}}%
%BeginExpansion
\ooalign{\hfil$V$\hfil\cr\kern0.1em--\hfil\cr}%
%EndExpansion
}}
\end{equation*}

We can use our relationship between $C_{%
%TCIMACRO{\TeXButton{V}{\ooalign{\hfil$V$\hfil\cr\kern0.1em--\hfil\cr}}}%
%BeginExpansion
\ooalign{\hfil$V$\hfil\cr\kern0.1em--\hfil\cr}%
%EndExpansion
}$ and $C_{P}$ that we found earlier for ideal gases%
\begin{equation*}
C_{%
%TCIMACRO{\TeXButton{V}{\ooalign{\hfil$V$\hfil\cr\kern0.1em--\hfil\cr}}}%
%BeginExpansion
\ooalign{\hfil$V$\hfil\cr\kern0.1em--\hfil\cr}%
%EndExpansion
}=C_{P}-R
\end{equation*}%
so%
\begin{equation*}
R=C_{P}-C_{%
%TCIMACRO{\TeXButton{V}{\ooalign{\hfil$V$\hfil\cr\kern0.1em--\hfil\cr}}}%
%BeginExpansion
\ooalign{\hfil$V$\hfil\cr\kern0.1em--\hfil\cr}%
%EndExpansion
}
\end{equation*}%
and let's substitute this in for our $R$ in our total differential equation%
\begin{equation*}
Pd%
%TCIMACRO{\TeXButton{V}{\ooalign{\hfil$V$\hfil\cr\kern0.1em--\hfil\cr}}}%
%BeginExpansion
\ooalign{\hfil$V$\hfil\cr\kern0.1em--\hfil\cr}%
%EndExpansion
+%
%TCIMACRO{\TeXButton{V}{\ooalign{\hfil$V$\hfil\cr\kern0.1em--\hfil\cr}}}%
%BeginExpansion
\ooalign{\hfil$V$\hfil\cr\kern0.1em--\hfil\cr}%
%EndExpansion
dP=-\frac{C_{P}-C_{%
%TCIMACRO{\TeXButton{V}{\ooalign{\hfil$V$\hfil\cr\kern0.1em--\hfil\cr}}}%
%BeginExpansion
\ooalign{\hfil$V$\hfil\cr\kern0.1em--\hfil\cr}%
%EndExpansion
}}{C_{%
%TCIMACRO{\TeXButton{V}{\ooalign{\hfil$V$\hfil\cr\kern0.1em--\hfil\cr}}}%
%BeginExpansion
\ooalign{\hfil$V$\hfil\cr\kern0.1em--\hfil\cr}%
%EndExpansion
}}Pd%
%TCIMACRO{\TeXButton{V}{\ooalign{\hfil$V$\hfil\cr\kern0.1em--\hfil\cr}}}%
%BeginExpansion
\ooalign{\hfil$V$\hfil\cr\kern0.1em--\hfil\cr}%
%EndExpansion
\end{equation*}%
or%
\begin{equation*}
Pd%
%TCIMACRO{\TeXButton{V}{\ooalign{\hfil$V$\hfil\cr\kern0.1em--\hfil\cr}}}%
%BeginExpansion
\ooalign{\hfil$V$\hfil\cr\kern0.1em--\hfil\cr}%
%EndExpansion
+%
%TCIMACRO{\TeXButton{V}{\ooalign{\hfil$V$\hfil\cr\kern0.1em--\hfil\cr}}}%
%BeginExpansion
\ooalign{\hfil$V$\hfil\cr\kern0.1em--\hfil\cr}%
%EndExpansion
dP=\frac{C_{%
%TCIMACRO{\TeXButton{V}{\ooalign{\hfil$V$\hfil\cr\kern0.1em--\hfil\cr}}}%
%BeginExpansion
\ooalign{\hfil$V$\hfil\cr\kern0.1em--\hfil\cr}%
%EndExpansion
}-C_{P}}{C_{%
%TCIMACRO{\TeXButton{V}{\ooalign{\hfil$V$\hfil\cr\kern0.1em--\hfil\cr}}}%
%BeginExpansion
\ooalign{\hfil$V$\hfil\cr\kern0.1em--\hfil\cr}%
%EndExpansion
}}Pd%
%TCIMACRO{\TeXButton{V}{\ooalign{\hfil$V$\hfil\cr\kern0.1em--\hfil\cr}}}%
%BeginExpansion
\ooalign{\hfil$V$\hfil\cr\kern0.1em--\hfil\cr}%
%EndExpansion
\end{equation*}%
This doesn't seem to have helped. To plot our adiabatic process we need $P$
as a function of $%
%TCIMACRO{\TeXButton{V}{\ooalign{\hfil$V$\hfil\cr\kern0.1em--\hfil\cr}}}%
%BeginExpansion
\ooalign{\hfil$V$\hfil\cr\kern0.1em--\hfil\cr}%
%EndExpansion
$ and what we have now is an algebraic mess. But now let's divide through by 
$P%
%TCIMACRO{\TeXButton{V}{\ooalign{\hfil$V$\hfil\cr\kern0.1em--\hfil\cr}}}%
%BeginExpansion
\ooalign{\hfil$V$\hfil\cr\kern0.1em--\hfil\cr}%
%EndExpansion
.$ It's not apparent that this will help, but it does so 
\begin{equation*}
\frac{Pd%
%TCIMACRO{\TeXButton{V}{\ooalign{\hfil$V$\hfil\cr\kern0.1em--\hfil\cr}}}%
%BeginExpansion
\ooalign{\hfil$V$\hfil\cr\kern0.1em--\hfil\cr}%
%EndExpansion
}{P%
%TCIMACRO{\TeXButton{V}{\ooalign{\hfil$V$\hfil\cr\kern0.1em--\hfil\cr}}}%
%BeginExpansion
\ooalign{\hfil$V$\hfil\cr\kern0.1em--\hfil\cr}%
%EndExpansion
}+\frac{%
%TCIMACRO{\TeXButton{V}{\ooalign{\hfil$V$\hfil\cr\kern0.1em--\hfil\cr}}}%
%BeginExpansion
\ooalign{\hfil$V$\hfil\cr\kern0.1em--\hfil\cr}%
%EndExpansion
dP}{P%
%TCIMACRO{\TeXButton{V}{\ooalign{\hfil$V$\hfil\cr\kern0.1em--\hfil\cr}}}%
%BeginExpansion
\ooalign{\hfil$V$\hfil\cr\kern0.1em--\hfil\cr}%
%EndExpansion
}=\frac{C_{%
%TCIMACRO{\TeXButton{V}{\ooalign{\hfil$V$\hfil\cr\kern0.1em--\hfil\cr}}}%
%BeginExpansion
\ooalign{\hfil$V$\hfil\cr\kern0.1em--\hfil\cr}%
%EndExpansion
}-C_{P}}{C_{%
%TCIMACRO{\TeXButton{V}{\ooalign{\hfil$V$\hfil\cr\kern0.1em--\hfil\cr}}}%
%BeginExpansion
\ooalign{\hfil$V$\hfil\cr\kern0.1em--\hfil\cr}%
%EndExpansion
}}\frac{Pd%
%TCIMACRO{\TeXButton{V}{\ooalign{\hfil$V$\hfil\cr\kern0.1em--\hfil\cr}}}%
%BeginExpansion
\ooalign{\hfil$V$\hfil\cr\kern0.1em--\hfil\cr}%
%EndExpansion
}{P%
%TCIMACRO{\TeXButton{V}{\ooalign{\hfil$V$\hfil\cr\kern0.1em--\hfil\cr}}}%
%BeginExpansion
\ooalign{\hfil$V$\hfil\cr\kern0.1em--\hfil\cr}%
%EndExpansion
}
\end{equation*}%
and canceling the extra $%
%TCIMACRO{\TeXButton{V}{\ooalign{\hfil$V$\hfil\cr\kern0.1em--\hfil\cr}}}%
%BeginExpansion
\ooalign{\hfil$V$\hfil\cr\kern0.1em--\hfil\cr}%
%EndExpansion
$'s and $P$'s gives%
\begin{equation*}
\frac{d%
%TCIMACRO{\TeXButton{V}{\ooalign{\hfil$V$\hfil\cr\kern0.1em--\hfil\cr}}}%
%BeginExpansion
\ooalign{\hfil$V$\hfil\cr\kern0.1em--\hfil\cr}%
%EndExpansion
}{%
%TCIMACRO{\TeXButton{V}{\ooalign{\hfil$V$\hfil\cr\kern0.1em--\hfil\cr}}}%
%BeginExpansion
\ooalign{\hfil$V$\hfil\cr\kern0.1em--\hfil\cr}%
%EndExpansion
}+\frac{dP}{P}=\frac{C_{%
%TCIMACRO{\TeXButton{V}{\ooalign{\hfil$V$\hfil\cr\kern0.1em--\hfil\cr}}}%
%BeginExpansion
\ooalign{\hfil$V$\hfil\cr\kern0.1em--\hfil\cr}%
%EndExpansion
}-C_{P}}{C_{%
%TCIMACRO{\TeXButton{V}{\ooalign{\hfil$V$\hfil\cr\kern0.1em--\hfil\cr}}}%
%BeginExpansion
\ooalign{\hfil$V$\hfil\cr\kern0.1em--\hfil\cr}%
%EndExpansion
}}\frac{d%
%TCIMACRO{\TeXButton{V}{\ooalign{\hfil$V$\hfil\cr\kern0.1em--\hfil\cr}}}%
%BeginExpansion
\ooalign{\hfil$V$\hfil\cr\kern0.1em--\hfil\cr}%
%EndExpansion
}{%
%TCIMACRO{\TeXButton{V}{\ooalign{\hfil$V$\hfil\cr\kern0.1em--\hfil\cr}}}%
%BeginExpansion
\ooalign{\hfil$V$\hfil\cr\kern0.1em--\hfil\cr}%
%EndExpansion
}
\end{equation*}%
rearranging terms gives%
\begin{equation*}
\frac{dP}{P}=\left( \frac{C_{%
%TCIMACRO{\TeXButton{V}{\ooalign{\hfil$V$\hfil\cr\kern0.1em--\hfil\cr}}}%
%BeginExpansion
\ooalign{\hfil$V$\hfil\cr\kern0.1em--\hfil\cr}%
%EndExpansion
}-C_{P}}{C_{%
%TCIMACRO{\TeXButton{V}{\ooalign{\hfil$V$\hfil\cr\kern0.1em--\hfil\cr}}}%
%BeginExpansion
\ooalign{\hfil$V$\hfil\cr\kern0.1em--\hfil\cr}%
%EndExpansion
}}-1\right) \frac{d%
%TCIMACRO{\TeXButton{V}{\ooalign{\hfil$V$\hfil\cr\kern0.1em--\hfil\cr}}}%
%BeginExpansion
\ooalign{\hfil$V$\hfil\cr\kern0.1em--\hfil\cr}%
%EndExpansion
}{%
%TCIMACRO{\TeXButton{V}{\ooalign{\hfil$V$\hfil\cr\kern0.1em--\hfil\cr}}}%
%BeginExpansion
\ooalign{\hfil$V$\hfil\cr\kern0.1em--\hfil\cr}%
%EndExpansion
}
\end{equation*}%
Lets get a common denominator for the part in parentheses 
\begin{eqnarray*}
\frac{dP}{P} &=&\left( \frac{C_{%
%TCIMACRO{\TeXButton{V}{\ooalign{\hfil$V$\hfil\cr\kern0.1em--\hfil\cr}}}%
%BeginExpansion
\ooalign{\hfil$V$\hfil\cr\kern0.1em--\hfil\cr}%
%EndExpansion
}-C_{P}}{C_{%
%TCIMACRO{\TeXButton{V}{\ooalign{\hfil$V$\hfil\cr\kern0.1em--\hfil\cr}}}%
%BeginExpansion
\ooalign{\hfil$V$\hfil\cr\kern0.1em--\hfil\cr}%
%EndExpansion
}}-\frac{C_{%
%TCIMACRO{\TeXButton{V}{\ooalign{\hfil$V$\hfil\cr\kern0.1em--\hfil\cr}}}%
%BeginExpansion
\ooalign{\hfil$V$\hfil\cr\kern0.1em--\hfil\cr}%
%EndExpansion
}}{C_{%
%TCIMACRO{\TeXButton{V}{\ooalign{\hfil$V$\hfil\cr\kern0.1em--\hfil\cr}}}%
%BeginExpansion
\ooalign{\hfil$V$\hfil\cr\kern0.1em--\hfil\cr}%
%EndExpansion
}}\right) \frac{d%
%TCIMACRO{\TeXButton{V}{\ooalign{\hfil$V$\hfil\cr\kern0.1em--\hfil\cr}}}%
%BeginExpansion
\ooalign{\hfil$V$\hfil\cr\kern0.1em--\hfil\cr}%
%EndExpansion
}{%
%TCIMACRO{\TeXButton{V}{\ooalign{\hfil$V$\hfil\cr\kern0.1em--\hfil\cr}}}%
%BeginExpansion
\ooalign{\hfil$V$\hfil\cr\kern0.1em--\hfil\cr}%
%EndExpansion
} \\
\frac{dP}{P} &=&\left( -\frac{C_{P}}{C_{%
%TCIMACRO{\TeXButton{V}{\ooalign{\hfil$V$\hfil\cr\kern0.1em--\hfil\cr}}}%
%BeginExpansion
\ooalign{\hfil$V$\hfil\cr\kern0.1em--\hfil\cr}%
%EndExpansion
}}\right) \frac{d%
%TCIMACRO{\TeXButton{V}{\ooalign{\hfil$V$\hfil\cr\kern0.1em--\hfil\cr}}}%
%BeginExpansion
\ooalign{\hfil$V$\hfil\cr\kern0.1em--\hfil\cr}%
%EndExpansion
}{%
%TCIMACRO{\TeXButton{V}{\ooalign{\hfil$V$\hfil\cr\kern0.1em--\hfil\cr}}}%
%BeginExpansion
\ooalign{\hfil$V$\hfil\cr\kern0.1em--\hfil\cr}%
%EndExpansion
}
\end{eqnarray*}%
We know that $\frac{C_{P}}{C_{%
%TCIMACRO{\TeXButton{V}{\ooalign{\hfil$V$\hfil\cr\kern0.1em--\hfil\cr}}}%
%BeginExpansion
\ooalign{\hfil$V$\hfil\cr\kern0.1em--\hfil\cr}%
%EndExpansion
}}$ is just a constant. Let's give it a symbol so we don't have to write out 
$\frac{C_{P}}{C_{%
%TCIMACRO{\TeXButton{V}{\ooalign{\hfil$V$\hfil\cr\kern0.1em--\hfil\cr}}}%
%BeginExpansion
\ooalign{\hfil$V$\hfil\cr\kern0.1em--\hfil\cr}%
%EndExpansion
}}.$ Let's make that symbol a Greek gamma, $\gamma $ 
\begin{equation}
\gamma =\frac{C_{P}}{C_{V}}
\end{equation}%
so we have%
\begin{equation*}
\frac{dP}{P}=-\gamma \frac{d%
%TCIMACRO{\TeXButton{V}{\ooalign{\hfil$V$\hfil\cr\kern0.1em--\hfil\cr}}}%
%BeginExpansion
\ooalign{\hfil$V$\hfil\cr\kern0.1em--\hfil\cr}%
%EndExpansion
}{%
%TCIMACRO{\TeXButton{V}{\ooalign{\hfil$V$\hfil\cr\kern0.1em--\hfil\cr}}}%
%BeginExpansion
\ooalign{\hfil$V$\hfil\cr\kern0.1em--\hfil\cr}%
%EndExpansion
}
\end{equation*}

Which is exciting. We are just down to $P$'s and $%
%TCIMACRO{\TeXButton{V}{\ooalign{\hfil$V$\hfil\cr\kern0.1em--\hfil\cr}}}%
%BeginExpansion
\ooalign{\hfil$V$\hfil\cr\kern0.1em--\hfil\cr}%
%EndExpansion
$'s and a gamma. We may be able to get an expression of $P$ as a function of 
$%
%TCIMACRO{\TeXButton{V}{\ooalign{\hfil$V$\hfil\cr\kern0.1em--\hfil\cr}}}%
%BeginExpansion
\ooalign{\hfil$V$\hfil\cr\kern0.1em--\hfil\cr}%
%EndExpansion
$ yet! We can integrate this last equation%
\begin{eqnarray*}
\int \frac{dP}{P} &=&-\gamma \int \frac{d%
%TCIMACRO{\TeXButton{V}{\ooalign{\hfil$V$\hfil\cr\kern0.1em--\hfil\cr}}}%
%BeginExpansion
\ooalign{\hfil$V$\hfil\cr\kern0.1em--\hfil\cr}%
%EndExpansion
}{%
%TCIMACRO{\TeXButton{V}{\ooalign{\hfil$V$\hfil\cr\kern0.1em--\hfil\cr}}}%
%BeginExpansion
\ooalign{\hfil$V$\hfil\cr\kern0.1em--\hfil\cr}%
%EndExpansion
} \\
\ln \left( P\right) +\zeta _{P} &=&-\gamma \ln \left( 
%TCIMACRO{\TeXButton{V}{\ooalign{\hfil$V$\hfil\cr\kern0.1em--\hfil\cr}}}%
%BeginExpansion
\ooalign{\hfil$V$\hfil\cr\kern0.1em--\hfil\cr}%
%EndExpansion
\right) +\zeta _{V}
\end{eqnarray*}%
where $\zeta _{P}$ and $\zeta _{V}$ are the constants of integration. We can
write this as%
\begin{eqnarray*}
\ln \left( P\right) +\gamma \ln \left( 
%TCIMACRO{\TeXButton{V}{\ooalign{\hfil$V$\hfil\cr\kern0.1em--\hfil\cr}}}%
%BeginExpansion
\ooalign{\hfil$V$\hfil\cr\kern0.1em--\hfil\cr}%
%EndExpansion
\right) &=&-\zeta _{P}+\zeta _{V} \\
&=&\zeta
\end{eqnarray*}%
where $\zeta =\zeta _{V}-\zeta _{P}$ is still just a constant. And now we
need to remember a few things about logs. 
\begin{equation*}
a\ln b=\ln b^{a}
\end{equation*}%
so we can write our equation as 
\begin{equation*}
\ln \left( P\right) +\ln \left( 
%TCIMACRO{\TeXButton{V}{\ooalign{\hfil$V$\hfil\cr\kern0.1em--\hfil\cr}}}%
%BeginExpansion
\ooalign{\hfil$V$\hfil\cr\kern0.1em--\hfil\cr}%
%EndExpansion
^{\gamma }\right) =\zeta
\end{equation*}%
and also remember that for logs 
\begin{equation*}
\ln a+\ln b=\ln \left( ab\right)
\end{equation*}
so our equation becomes just 
\begin{equation*}
\ln \left( P%
%TCIMACRO{\TeXButton{V}{\ooalign{\hfil$V$\hfil\cr\kern0.1em--\hfil\cr}}}%
%BeginExpansion
\ooalign{\hfil$V$\hfil\cr\kern0.1em--\hfil\cr}%
%EndExpansion
^{\gamma }\right) =\zeta
\end{equation*}%
or we can exponentiate both sides 
\begin{equation}
P%
%TCIMACRO{\TeXButton{V}{\ooalign{\hfil$V$\hfil\cr\kern0.1em--\hfil\cr}}}%
%BeginExpansion
\ooalign{\hfil$V$\hfil\cr\kern0.1em--\hfil\cr}%
%EndExpansion
^{\gamma }=e^{\zeta }
\end{equation}%
where $e^{\zeta }$ is still just a constant.

Note that we have done what we set out to do (sort of). We have an
expression of how $P$ varies with $%
%TCIMACRO{\TeXButton{V}{\ooalign{\hfil$V$\hfil\cr\kern0.1em--\hfil\cr}}}%
%BeginExpansion
\ooalign{\hfil$V$\hfil\cr\kern0.1em--\hfil\cr}%
%EndExpansion
.$ 
\begin{equation*}
P=\frac{e^{\zeta }}{%
%TCIMACRO{\TeXButton{V}{\ooalign{\hfil$V$\hfil\cr\kern0.1em--\hfil\cr}}}%
%BeginExpansion
\ooalign{\hfil$V$\hfil\cr\kern0.1em--\hfil\cr}%
%EndExpansion
^{\gamma }}
\end{equation*}%
We could put in $P$ and $%
%TCIMACRO{\TeXButton{V}{\ooalign{\hfil$V$\hfil\cr\kern0.1em--\hfil\cr}}}%
%BeginExpansion
\ooalign{\hfil$V$\hfil\cr\kern0.1em--\hfil\cr}%
%EndExpansion
$ values and solve for the constant and then we could plot our equation to
get an adiabatic curve (which is how I got the adiabatic curve earlier). We
could also put this into our work equation and integrate to find the work
done in an adiabatic process!%
\begin{eqnarray*}
\Delta E_{int} &=&w \\
&=&-\int \frac{e^{\zeta }}{%
%TCIMACRO{\TeXButton{V}{\ooalign{\hfil$V$\hfil\cr\kern0.1em--\hfil\cr}}}%
%BeginExpansion
\ooalign{\hfil$V$\hfil\cr\kern0.1em--\hfil\cr}%
%EndExpansion
^{\gamma }}d%
%TCIMACRO{\TeXButton{V}{\ooalign{\hfil$V$\hfil\cr\kern0.1em--\hfil\cr}}}%
%BeginExpansion
\ooalign{\hfil$V$\hfil\cr\kern0.1em--\hfil\cr}%
%EndExpansion
\end{eqnarray*}

But also note that we have a great conservation equation for an adiabatic
process.

\begin{equation*}
P%
%TCIMACRO{\TeXButton{V}{\ooalign{\hfil$V$\hfil\cr\kern0.1em--\hfil\cr}}}%
%BeginExpansion
\ooalign{\hfil$V$\hfil\cr\kern0.1em--\hfil\cr}%
%EndExpansion
^{\gamma }=\text{constant}
\end{equation*}%
or 
\begin{equation}
P_{i}%
%TCIMACRO{\TeXButton{V}{\ooalign{\hfil$V$\hfil\cr\kern0.1em--\hfil\cr}}}%
%BeginExpansion
\ooalign{\hfil$V$\hfil\cr\kern0.1em--\hfil\cr}%
%EndExpansion
_{i}^{\gamma }=P_{f}%
%TCIMACRO{\TeXButton{V}{\ooalign{\hfil$V$\hfil\cr\kern0.1em--\hfil\cr}}}%
%BeginExpansion
\ooalign{\hfil$V$\hfil\cr\kern0.1em--\hfil\cr}%
%EndExpansion
_{f}^{\gamma }
\end{equation}%
and we know that conservation equations come in very handy in solving
problems! Let's add this last equation to our set of equations for adiabatic
processes.

We can get a second equation for adiabatic processes relating the volume and
temperature. If we know an initial state, 
\begin{equation*}
P_{i}%
%TCIMACRO{\TeXButton{V}{\ooalign{\hfil$V$\hfil\cr\kern0.1em--\hfil\cr}}}%
%BeginExpansion
\ooalign{\hfil$V$\hfil\cr\kern0.1em--\hfil\cr}%
%EndExpansion
_{i}=nRT_{i}
\end{equation*}%
then we can solve the ideal gas law for $P_{i}$ and 
\begin{equation*}
P_{i}=\frac{nRT_{i}}{%
%TCIMACRO{\TeXButton{V}{\ooalign{\hfil$V$\hfil\cr\kern0.1em--\hfil\cr}}}%
%BeginExpansion
\ooalign{\hfil$V$\hfil\cr\kern0.1em--\hfil\cr}%
%EndExpansion
_{i}}
\end{equation*}%
and use this in our adiabatic conservation equation so%
\begin{equation}
P_{i}%
%TCIMACRO{\TeXButton{V}{\ooalign{\hfil$V$\hfil\cr\kern0.1em--\hfil\cr}}}%
%BeginExpansion
\ooalign{\hfil$V$\hfil\cr\kern0.1em--\hfil\cr}%
%EndExpansion
_{i}^{\gamma }=P_{f}%
%TCIMACRO{\TeXButton{V}{\ooalign{\hfil$V$\hfil\cr\kern0.1em--\hfil\cr}}}%
%BeginExpansion
\ooalign{\hfil$V$\hfil\cr\kern0.1em--\hfil\cr}%
%EndExpansion
_{f}^{\gamma }
\end{equation}%
becomes 
\begin{equation*}
\frac{nRT_{f}}{%
%TCIMACRO{\TeXButton{V}{\ooalign{\hfil$V$\hfil\cr\kern0.1em--\hfil\cr}}}%
%BeginExpansion
\ooalign{\hfil$V$\hfil\cr\kern0.1em--\hfil\cr}%
%EndExpansion
_{f}}%
%TCIMACRO{\TeXButton{V}{\ooalign{\hfil$V$\hfil\cr\kern0.1em--\hfil\cr}}}%
%BeginExpansion
\ooalign{\hfil$V$\hfil\cr\kern0.1em--\hfil\cr}%
%EndExpansion
_{f}^{\gamma }=\frac{nRT_{i}}{%
%TCIMACRO{\TeXButton{V}{\ooalign{\hfil$V$\hfil\cr\kern0.1em--\hfil\cr}}}%
%BeginExpansion
\ooalign{\hfil$V$\hfil\cr\kern0.1em--\hfil\cr}%
%EndExpansion
_{i}}%
%TCIMACRO{\TeXButton{V}{\ooalign{\hfil$V$\hfil\cr\kern0.1em--\hfil\cr}}}%
%BeginExpansion
\ooalign{\hfil$V$\hfil\cr\kern0.1em--\hfil\cr}%
%EndExpansion
_{i}^{\gamma }
\end{equation*}%
The $n$'s and $R$'s cancel,%
\begin{equation*}
\frac{T_{f}}{%
%TCIMACRO{\TeXButton{V}{\ooalign{\hfil$V$\hfil\cr\kern0.1em--\hfil\cr}}}%
%BeginExpansion
\ooalign{\hfil$V$\hfil\cr\kern0.1em--\hfil\cr}%
%EndExpansion
_{f}}%
%TCIMACRO{\TeXButton{V}{\ooalign{\hfil$V$\hfil\cr\kern0.1em--\hfil\cr}}}%
%BeginExpansion
\ooalign{\hfil$V$\hfil\cr\kern0.1em--\hfil\cr}%
%EndExpansion
_{f}^{\gamma }=\frac{T_{i}}{%
%TCIMACRO{\TeXButton{V}{\ooalign{\hfil$V$\hfil\cr\kern0.1em--\hfil\cr}}}%
%BeginExpansion
\ooalign{\hfil$V$\hfil\cr\kern0.1em--\hfil\cr}%
%EndExpansion
_{i}}%
%TCIMACRO{\TeXButton{V}{\ooalign{\hfil$V$\hfil\cr\kern0.1em--\hfil\cr}}}%
%BeginExpansion
\ooalign{\hfil$V$\hfil\cr\kern0.1em--\hfil\cr}%
%EndExpansion
_{i}^{\gamma }
\end{equation*}%
And we can combine the $%
%TCIMACRO{\TeXButton{V}{\ooalign{\hfil$V$\hfil\cr\kern0.1em--\hfil\cr}}}%
%BeginExpansion
\ooalign{\hfil$V$\hfil\cr\kern0.1em--\hfil\cr}%
%EndExpansion
$ terms to get 
\begin{equation}
T_{f}%
%TCIMACRO{\TeXButton{V}{\ooalign{\hfil$V$\hfil\cr\kern0.1em--\hfil\cr}}}%
%BeginExpansion
\ooalign{\hfil$V$\hfil\cr\kern0.1em--\hfil\cr}%
%EndExpansion
_{f}^{\gamma -1}=T_{i}%
%TCIMACRO{\TeXButton{V}{\ooalign{\hfil$V$\hfil\cr\kern0.1em--\hfil\cr}}}%
%BeginExpansion
\ooalign{\hfil$V$\hfil\cr\kern0.1em--\hfil\cr}%
%EndExpansion
_{i}^{\gamma -1}
\end{equation}%
This is another conservation equation for adiabatic processes! We will add
this to our equation set for adiabatic processes. These equations are great
shortcuts for problem solving \emph{but remember they are only valid for
adiabatic processes!}

Notice that we never did complete the integral for work%
\begin{equation*}
\Delta E_{int}=w=-\int \frac{e^{\zeta }}{%
%TCIMACRO{\TeXButton{V}{\ooalign{\hfil$V$\hfil\cr\kern0.1em--\hfil\cr}}}%
%BeginExpansion
\ooalign{\hfil$V$\hfil\cr\kern0.1em--\hfil\cr}%
%EndExpansion
^{\gamma }}d%
%TCIMACRO{\TeXButton{V}{\ooalign{\hfil$V$\hfil\cr\kern0.1em--\hfil\cr}}}%
%BeginExpansion
\ooalign{\hfil$V$\hfil\cr\kern0.1em--\hfil\cr}%
%EndExpansion
\end{equation*}%
We totally could at this point, but we already know the answer to be 
\begin{equation*}
\Delta E_{int}=nC_{%
%TCIMACRO{\TeXButton{V}{\ooalign{\hfil$V$\hfil\cr\kern0.1em--\hfil\cr}}}%
%BeginExpansion
\ooalign{\hfil$V$\hfil\cr\kern0.1em--\hfil\cr}%
%EndExpansion
}\Delta T
\end{equation*}%
where the $\Delta T$ is the same as it would be for our actual adiabatic
process. Since we know this and it is easy to calculate, we won't often do
the integral

We still need to find $\gamma $ for our adiabatic processes. We need the
ratio of $C_{P}$ to $C_{%
%TCIMACRO{\TeXButton{V}{\ooalign{\hfil$V$\hfil\cr\kern0.1em--\hfil\cr}}}%
%BeginExpansion
\ooalign{\hfil$V$\hfil\cr\kern0.1em--\hfil\cr}%
%EndExpansion
}$ or $\gamma .$ We now know that for ideal gasses 
\begin{equation*}
C_{%
%TCIMACRO{\TeXButton{V}{\ooalign{\hfil$V$\hfil\cr\kern0.1em--\hfil\cr}}}%
%BeginExpansion
\ooalign{\hfil$V$\hfil\cr\kern0.1em--\hfil\cr}%
%EndExpansion
}=\frac{3}{2}R
\end{equation*}%
\begin{equation*}
C_{P}=\frac{5}{2}R
\end{equation*}%
and we may take the ratio of these values%
\begin{eqnarray*}
\gamma &=&\frac{C_{P}}{C_{%
%TCIMACRO{\TeXButton{V}{\ooalign{\hfil$V$\hfil\cr\kern0.1em--\hfil\cr}}}%
%BeginExpansion
\ooalign{\hfil$V$\hfil\cr\kern0.1em--\hfil\cr}%
%EndExpansion
}} \\
&=&\frac{\frac{5}{2}R}{\frac{3}{2}R} \\
&=&\frac{5}{3}
\end{eqnarray*}

This works well for monotonic gasses, but fails badly for more complex
gasses. So we see our ideal gas formulation starts to break down with more
complex gasses at this point. To go farther, we would need to include the
rotational and vibrational energy of the molecules. Still, we have come a
long ways with our simple ideal gas assumptions!

Note that for solids and liquids $\Delta 
%TCIMACRO{\TeXButton{V}{\ooalign{\hfil$V$\hfil\cr\kern0.1em--\hfil\cr}}}%
%BeginExpansion
\ooalign{\hfil$V$\hfil\cr\kern0.1em--\hfil\cr}%
%EndExpansion
$ is very small so very little work is done. This means that $C_{P}\approx
C_{%
%TCIMACRO{\TeXButton{V}{\ooalign{\hfil$V$\hfil\cr\kern0.1em--\hfil\cr}}}%
%BeginExpansion
\ooalign{\hfil$V$\hfil\cr\kern0.1em--\hfil\cr}%
%EndExpansion
}$ which is why we we could get away with only one value for the molar heat
capacity in the tables for solids and liquids.

\subsection{Path Dependence}

We should do a problem that shows how the path dependence of thermodynamics
works. Notice that how we compress the gas (quickly vs. slowly, etc.) makes
a difference on our PV diagram. For example, quick actions would not let
energy leave by heat. Convection and conduction and radiation take time. So
quick processes would be more like adiabatic processes and make the
temperature rise. Slow actions would allow enough time for energy to leave,
allowing the temperature to stay the same. So when we say there is a
\textquotedblleft path dependence\textquotedblright\ for thermodynamic work,
it is equivalent to saying we have different amounts of work for
\textquotedblleft different paths\textquotedblright\ on a PV-diagram. The
PV-diagram shows the different physical paths as different lines on the
graph. Here is an example.

Suppose we change $1\unit{mol}$ of an ideal gas from one state to another by
two different paths. The two paths are shown in the next figure.\FRAME{dtbpF%
}{3.0894in}{2.0676in}{0pt}{}{}{Figure}{\special{language "Scientific
Word";type "GRAPHIC";maintain-aspect-ratio TRUE;display "USEDEF";valid_file
"T";width 3.0894in;height 2.0676in;depth 0pt;original-width
3.1195in;original-height 2.0782in;cropleft "0";croptop "1";cropright
"1";cropbottom "0";tempfilename 'SUS73WAM.wmf';tempfile-properties "XPR";}}

We recognize the blue path as a combination of an isobaric and an isothermal
process. So for the blue path, the work done will be a combination of%
\begin{eqnarray*}
w_{isobaric} &=&-P\left( 
%TCIMACRO{\TeXButton{V}{\ooalign{\hfil$V$\hfil\cr\kern0.1em--\hfil\cr}}}%
%BeginExpansion
\ooalign{\hfil$V$\hfil\cr\kern0.1em--\hfil\cr}%
%EndExpansion
_{f}-%
%TCIMACRO{\TeXButton{V}{\ooalign{\hfil$V$\hfil\cr\kern0.1em--\hfil\cr}}}%
%BeginExpansion
\ooalign{\hfil$V$\hfil\cr\kern0.1em--\hfil\cr}%
%EndExpansion
_{i}\right) \\
w_{isothermal} &=&nRT\left[ \ln \frac{%
%TCIMACRO{\TeXButton{V}{\ooalign{\hfil$V$\hfil\cr\kern0.1em--\hfil\cr}}}%
%BeginExpansion
\ooalign{\hfil$V$\hfil\cr\kern0.1em--\hfil\cr}%
%EndExpansion
_{i}}{%
%TCIMACRO{\TeXButton{V}{\ooalign{\hfil$V$\hfil\cr\kern0.1em--\hfil\cr}}}%
%BeginExpansion
\ooalign{\hfil$V$\hfil\cr\kern0.1em--\hfil\cr}%
%EndExpansion
_{f}}\right]
\end{eqnarray*}

Suppose that along the isobaric part of the blue path 
\begin{eqnarray*}
P &=&378.29\unit{Pa} \\
%TCIMACRO{\TeXButton{V}{\ooalign{\hfil$V$\hfil\cr\kern0.1em--\hfil\cr}}}%
%BeginExpansion
\ooalign{\hfil$V$\hfil\cr\kern0.1em--\hfil\cr}%
%EndExpansion
_{i} &=&6\unit{m}^{3} \\
%TCIMACRO{\TeXButton{V}{\ooalign{\hfil$V$\hfil\cr\kern0.1em--\hfil\cr}}}%
%BeginExpansion
\ooalign{\hfil$V$\hfil\cr\kern0.1em--\hfil\cr}%
%EndExpansion
_{f} &=&7.\,\allowbreak 538\,4\unit{m}^{3}
\end{eqnarray*}%
and along the isothermal blue path%
\begin{equation*}
T=343\unit{K}
\end{equation*}%
Then the work done on the blue path would be%
\begin{eqnarray*}
w_{blue} &=&-P\left( 
%TCIMACRO{\TeXButton{V}{\ooalign{\hfil$V$\hfil\cr\kern0.1em--\hfil\cr}}}%
%BeginExpansion
\ooalign{\hfil$V$\hfil\cr\kern0.1em--\hfil\cr}%
%EndExpansion
_{f}-%
%TCIMACRO{\TeXButton{V}{\ooalign{\hfil$V$\hfil\cr\kern0.1em--\hfil\cr}}}%
%BeginExpansion
\ooalign{\hfil$V$\hfil\cr\kern0.1em--\hfil\cr}%
%EndExpansion
_{i}\right) +nRT\left[ \ln \frac{%
%TCIMACRO{\TeXButton{V}{\ooalign{\hfil$V$\hfil\cr\kern0.1em--\hfil\cr}}}%
%BeginExpansion
\ooalign{\hfil$V$\hfil\cr\kern0.1em--\hfil\cr}%
%EndExpansion
_{i}}{%
%TCIMACRO{\TeXButton{V}{\ooalign{\hfil$V$\hfil\cr\kern0.1em--\hfil\cr}}}%
%BeginExpansion
\ooalign{\hfil$V$\hfil\cr\kern0.1em--\hfil\cr}%
%EndExpansion
_{f}}\right] \\
&=&-378.29\unit{Pa}\left( 7.\,\allowbreak 538\,4\unit{m}^{3}-6\unit{m}%
^{3}\right) \\
&&+\left( 1\right) \left( 8.314\frac{\unit{J}}{\unit{mol}\unit{K}}\right)
\left( 343\unit{K}\right) \left[ \ln \frac{7.\,\allowbreak 538\,4\unit{m}^{3}%
}{6\unit{m}^{3}}\right] \\
&=&\allowbreak 650.\,\allowbreak 9\unit{J}
\end{eqnarray*}

But the red path is an isovolumetric path. We know that for isovolumetric
paths the work is zero!%
\begin{eqnarray*}
w_{red} &=&-\int_{%
%TCIMACRO{\TeXButton{V}{\ooalign{\hfil$V$\hfil\cr\kern0.1em--\hfil\cr}}}%
%BeginExpansion
\ooalign{\hfil$V$\hfil\cr\kern0.1em--\hfil\cr}%
%EndExpansion
_{1}}^{%
%TCIMACRO{\TeXButton{V}{\ooalign{\hfil$V$\hfil\cr\kern0.1em--\hfil\cr}}}%
%BeginExpansion
\ooalign{\hfil$V$\hfil\cr\kern0.1em--\hfil\cr}%
%EndExpansion
_{2}}Pd%
%TCIMACRO{\TeXButton{V}{\ooalign{\hfil$V$\hfil\cr\kern0.1em--\hfil\cr}} }%
%BeginExpansion
\ooalign{\hfil$V$\hfil\cr\kern0.1em--\hfil\cr}
%EndExpansion
\\
&=&0
\end{eqnarray*}%
and we can see that zero is very different than $651\unit{J}$! Indeed the
work for each PV-diagram path is different. But why?

The red path is what we would get if we put a sealed strong container in a
fire. The fire would provide energy to the gas in the container by heat. But
the strong walls of the container would keep the gas from expanding. The
pressure would rise and the temperature would rise as well. But for the blue
path we envision a movable piston in a cylinder of gas. when the fire adds
energy by heat, the piston rises, keeping the pressure the same. The
temperature does go up. Then we remove the fire, and dunk the cylinder in a
constant temperature bath. And we slowly compress the gas back to its
original volume. By doing this slowly, the work we do adds a little bit of
energy to the gas, but the energy can leave by heat, so the temperature of
the gas stays the same.

Note that these are really different processes! So it is no wonder that the
internal work done is different. And we see all of this with our PV-diagram!

\subsection{Finally finding the speed of sound in air \label%
{Speed_of_sound_derivation}}

This thermodynamics stuff is interesting, but you might wonder how it can be
used to do practical work. We will take on this question in our next lecture
for engines and other thermodynamic machines. But let's use what we know
adiabatic processes and revisit something we know about from earlier in our
course. Let's think about the speed of sound in air.

The disturbance that makes the wave will change the pressure and compress
the gas a bit. So work will be done on the gas. But the disturbance happens
quickly and we wouldn't expect energy transfer by heat. This is pretty close
to an adiabatic process!

We used pipes when we studied sound waves. So, let's consider our air in a
pipe. \FRAME{dtbpF}{3.3244in}{1.1061in}{0pt}{}{}{Figure}{\special{language
"Scientific Word";type "GRAPHIC";maintain-aspect-ratio TRUE;display
"USEDEF";valid_file "T";width 3.3244in;height 1.1061in;depth
0pt;original-width 3.359in;original-height 1.0981in;cropleft "0";croptop
"1";cropright "1";cropbottom "0";tempfilename
'T9AAUO00.wmf';tempfile-properties "XPR";}}The pipe has a cross sectional
area, $A$. And let's look at a small volume of air in the pipe. The darker
region is our small volume of air and let's say it has a mas $\Delta m$
where here the delta means \textquotedblleft a small amount
of\textquotedblright\ and not a difference. The density of the air is given
by 
\begin{equation*}
\rho =\frac{\Delta m}{\Delta 
%TCIMACRO{\TeXButton{V}{\ooalign{\hfil$V$\hfil\cr\kern0.1em--\hfil\cr}}}%
%BeginExpansion
\ooalign{\hfil$V$\hfil\cr\kern0.1em--\hfil\cr}%
%EndExpansion
}
\end{equation*}%
so we can write 
\begin{equation*}
\Delta m=\rho \Delta 
%TCIMACRO{\TeXButton{V}{\ooalign{\hfil$V$\hfil\cr\kern0.1em--\hfil\cr}}}%
%BeginExpansion
\ooalign{\hfil$V$\hfil\cr\kern0.1em--\hfil\cr}%
%EndExpansion
\end{equation*}%
But look at our little mass' volume%
\begin{equation*}
\Delta 
%TCIMACRO{\TeXButton{V}{\ooalign{\hfil$V$\hfil\cr\kern0.1em--\hfil\cr}}}%
%BeginExpansion
\ooalign{\hfil$V$\hfil\cr\kern0.1em--\hfil\cr}%
%EndExpansion
=A\Delta x
\end{equation*}%
We can write $\Delta m$ as 
\begin{equation*}
\Delta m=\rho A\Delta x
\end{equation*}%
Now let's consider the amount of mass that passes a point in our pipe. That
is a mass flow rate and we know that is 
\begin{equation*}
\mathbb{Q}=\frac{dm}{dt}
\end{equation*}%
If we make our small amount of mass very small, we would have 
\begin{equation*}
dm=\rho Adx
\end{equation*}%
and our mass flow rate would be 
\begin{equation*}
\frac{dm}{dt}=\frac{\rho Adx}{dt}=\rho Av
\end{equation*}%
This looks familiar. It is one side of our equation of continuity 
\begin{equation*}
\rho _{in}A_{in}v_{in}=\rho _{out}A_{out}v_{out}
\end{equation*}

Now consider a \textquotedblleft parcel of air.\textquotedblright\ Remember
that this is just a part of the air that we will pay attention to. It is
surrounded by more air and is only different because we are thinking about
it. But consider such a parcel of air and let's make it sort of box shaped
like in the next figure. \FRAME{dtbpF}{3.4042in}{1.4688in}{0pt}{}{}{Figure}{%
\special{language "Scientific Word";type "GRAPHIC";maintain-aspect-ratio
TRUE;display "USEDEF";valid_file "T";width 3.4042in;height 1.4688in;depth
0pt;original-width 3.4406in;original-height 1.4688in;cropleft "0";croptop
"1";cropright "1";cropbottom "0";tempfilename
'T9AAUO01.wmf';tempfile-properties "XPR";}}We could write $\rho _{in}=\rho $
and $v_{in}=v$ and $T_{in}=T$ and we could write $\rho _{out}=\rho +d\rho $, 
$v_{out}=v+dv,$ and $T_{out}=T+dT.$ Let's put these in our continuity
equation.%
\begin{equation*}
\rho Av=\left( \rho +d\rho \right) A\left( v+dv\right)
\end{equation*}%
The area, $A,$ didn't change. So it cancels, and we can do some algebra%
\begin{eqnarray*}
\rho v &=&\left( \rho +d\rho \right) \left( v+dv\right) \\
\rho v &=&\rho v+\rho dv+vd\rho +d\rho dv
\end{eqnarray*}%
We can assume that $d\rho $ will be small if our parcel of air isn't too
big. And likewise $dv$ will be small. So $d\rho dv$ should be very small.
Let's drop this term.%
\begin{equation*}
\rho v=\rho v+\rho dv+vd\rho
\end{equation*}%
and the $\rho v$ terms cancel%
\begin{equation*}
0=\rho dv+vd\rho
\end{equation*}%
so that%
\begin{equation}
\rho dv=-vd\rho  \label{fromcontinuity}
\end{equation}%
This doesn't seem to have helped, but we will use it later. Let's go back to
our parcel of air and this time think of the forces on the input and output
sides. \FRAME{dtbpF}{3.995in}{1.7207in}{0pt}{}{}{Figure}{\special{language
"Scientific Word";type "GRAPHIC";maintain-aspect-ratio TRUE;display
"USEDEF";valid_file "T";width 3.995in;height 1.7207in;depth
0pt;original-width 4.0429in;original-height 1.7252in;cropleft "0";croptop
"1";cropright "1";cropbottom "0";tempfilename
'T9AAUO02.wmf';tempfile-properties "XPR";}}Remember that there is nothing
special about our parcel of air. There is more air all around it. And that
surrounding air will push on our parcel of air. If we look just at the input
and output sides of our parcel, we can see that we will have an input
pressure and an output pressure with $P_{in}=P$ and $P_{out}=P+dP.$ And I\
suppose $dP$ could be positive or negative. We are just saying that the
pressure could be different on the two sides.

We remember that $F=PA$ so we will have a force on the input side%
\begin{equation*}
F_{in}=P_{in}A=PA=Pdydz
\end{equation*}%
and an output force of 
\begin{equation*}
F_{out}=-P_{out}A=-\left( P+dP\right) dydz
\end{equation*}%
on the output side. The net force would be 
\begin{eqnarray*}
F_{net} &=&Pdydz-Pdydz-dPdydz \\
&=&-dPdydz
\end{eqnarray*}%
Then 
\begin{equation*}
ma=-dPdydz
\end{equation*}%
We can divide by the mass to get the acceleration%
\begin{equation*}
a=-\frac{dPdydz}{m}
\end{equation*}%
And now think, 
\begin{equation*}
\rho =\frac{m}{%
%TCIMACRO{\TeXButton{V}{\ooalign{\hfil$V$\hfil\cr\kern0.1em--\hfil\cr}}}%
%BeginExpansion
\ooalign{\hfil$V$\hfil\cr\kern0.1em--\hfil\cr}%
%EndExpansion
}=\frac{m}{dxdydz}
\end{equation*}%
so 
\begin{equation*}
m=\rho dxdydz
\end{equation*}%
so our acceleration is 
\begin{eqnarray*}
a &=&-\frac{dPdydz}{\rho dxdydz} \\
&=&-\frac{dP}{\rho dx}
\end{eqnarray*}%
and $a=dv/dt$%
\begin{equation*}
\frac{dv}{dt}=-\frac{dP}{\rho dx}
\end{equation*}%
or%
\begin{equation*}
dv=-\frac{dPdt}{\rho dx}
\end{equation*}%
And recall that $v=dx/dt$%
\begin{equation*}
dv=-\frac{dP}{\rho v}
\end{equation*}%
or 
\begin{equation*}
\rho vdv=-dP
\end{equation*}

Finely we can use our result from our continuity equation (\ref%
{fromcontinuity}) $\rho dv=-vd\rho $%
\begin{eqnarray*}
v\rho dv &=&-dP \\
v\left( -vd\rho \right) &=&-dP \\
v\left( -vd\rho \right) &=&-dP \\
v^{2} &=&\frac{dP}{d\rho }
\end{eqnarray*}%
and we are close. The wave speed, $v,$ is the square root of the pressure
change over the density change. 
\begin{equation}
v=\sqrt{\frac{dP}{d\rho }}  \label{Sound_speed}
\end{equation}%
This is a little like $\sqrt{B/\rho }.$ So we are close.

Since we are modeling this as an adiabatic process we know 
\begin{equation*}
P%
%TCIMACRO{\TeXButton{V}{\ooalign{\hfil$V$\hfil\cr\kern0.1em--\hfil\cr}}}%
%BeginExpansion
\ooalign{\hfil$V$\hfil\cr\kern0.1em--\hfil\cr}%
%EndExpansion
^{\gamma }=\text{constant}
\end{equation*}%
Let's look at that volume, $%
%TCIMACRO{\TeXButton{V}{\ooalign{\hfil$V$\hfil\cr\kern0.1em--\hfil\cr}}}%
%BeginExpansion
\ooalign{\hfil$V$\hfil\cr\kern0.1em--\hfil\cr}%
%EndExpansion
$. 
\begin{equation*}
%TCIMACRO{\TeXButton{V}{\ooalign{\hfil$V$\hfil\cr\kern0.1em--\hfil\cr}}}%
%BeginExpansion
\ooalign{\hfil$V$\hfil\cr\kern0.1em--\hfil\cr}%
%EndExpansion
=\frac{m}{\rho }
\end{equation*}%
and we can write the mass, $m,$ in terms of the molar mass%
\begin{equation*}
M=\frac{m}{n}
\end{equation*}%
or 
\begin{equation*}
m=nM
\end{equation*}%
So, our volume is 
\begin{equation*}
%TCIMACRO{\TeXButton{V}{\ooalign{\hfil$V$\hfil\cr\kern0.1em--\hfil\cr}}}%
%BeginExpansion
\ooalign{\hfil$V$\hfil\cr\kern0.1em--\hfil\cr}%
%EndExpansion
=\frac{nM}{\rho }
\end{equation*}%
and let's put this into our adiabatic process equation%
\begin{equation*}
P%
%TCIMACRO{\TeXButton{V}{\ooalign{\hfil$V$\hfil\cr\kern0.1em--\hfil\cr}}}%
%BeginExpansion
\ooalign{\hfil$V$\hfil\cr\kern0.1em--\hfil\cr}%
%EndExpansion
^{\gamma }=P\left( \frac{nM}{\rho }\right) ^{\gamma }=\text{constant}
\end{equation*}%
We could write this as 
\begin{equation*}
P\frac{n^{\gamma }M^{\gamma }}{\rho ^{\gamma }}=\text{constant}
\end{equation*}%
and rearrange it to be 
\begin{equation*}
P\frac{1}{\rho ^{\gamma }}=\frac{\text{constant}}{n^{\gamma }M^{\gamma }}
\end{equation*}%
and since $n$ and $M$ are not changing 
\begin{equation*}
P\frac{1}{\rho ^{\gamma }}=\text{new constant}
\end{equation*}

Let's take the logarithm of both sides (yes I\ know it is not obvious that
this will help, but let's do it anyway).%
\begin{eqnarray*}
\ln \left[ P\frac{1}{\rho ^{\gamma }}\right] &=&\ln \left[ \text{new constant%
}\right] \\
\ln P-\gamma \ln \rho &=&\ln \left[ \text{new constant}\right] \\
\ln P-\gamma \ln \rho &=&\text{another new constant}
\end{eqnarray*}%
Now let's take the derivative of both sides%
\begin{equation*}
\frac{d}{d\rho }\left( \ln P-\gamma \ln \rho \right) =\frac{d}{d\rho }\left( 
\text{another new constant}\right)
\end{equation*}%
The right hand side is just zero (which is why we didn't spend time figuring
out the number for the constant!). 
\begin{equation*}
\frac{d}{d\rho }\left( \ln P-\gamma \ln \rho \right) =0
\end{equation*}%
But the left hand side gives 
\begin{equation*}
\frac{1}{P}\frac{dP}{d\rho }-\frac{\gamma }{\rho }\frac{d\rho }{d\rho }=0
\end{equation*}%
and we can rearrange this 
\begin{equation*}
\frac{1}{P}\frac{dP}{d\rho }=\frac{\gamma }{\rho }
\end{equation*}%
\begin{equation}
\frac{dP}{d\rho }=\frac{\gamma }{\rho }P  \label{gas_speed_sound_partial}
\end{equation}%
Now let's pretend that the air is an ideal gas. Then we know that 
\begin{eqnarray*}
PV &=&nRT \\
&=&\frac{m}{M}RT
\end{eqnarray*}%
and we can solve for $P$ and put it into equation (\ref%
{gas_speed_sound_partial}) 
\begin{eqnarray*}
P &=&\frac{m}{VM}RT \\
&=&\frac{\rho }{M}RT
\end{eqnarray*}%
and then equation (\ref{gas_speed_sound_partial} becomes%
\begin{equation}
\frac{dP}{d\rho }=\frac{\gamma }{\rho }\frac{\rho }{M}RT=\frac{\gamma }{M}RT
\end{equation}

and we can put this directly into equation (\ref{Sound_speed})%
\begin{equation*}
v=\sqrt{\frac{dP}{d\rho }}=\sqrt{\frac{\gamma }{M}RT}
\end{equation*}%
We can get our approximation that we used earlier for temperature dependence
for sound speed in air from putting in some numbers. First, let's define $%
T_{o}=273\unit{K}.$ Then 
\begin{eqnarray*}
v &=&\sqrt{\frac{\gamma }{M}RT\frac{T_{o}}{T_{o}}} \\
&=&\sqrt{\frac{\gamma RT_{o}}{M}}\sqrt{\frac{T}{T_{o}}}
\end{eqnarray*}%
It turns out that for air at normal air pressure and at $T=10\unit{%
%TCIMACRO{\U{2103}}%
%BeginExpansion
{}^{\circ}{\rm C}%
%EndExpansion
},$ $\gamma =1.4$ and $M=0.02897\frac{\unit{kg}}{\unit{mol}},$ and of course 
$R=8.31\frac{\unit{J}}{\unit{mol}\unit{K}}$ then our speed is

\begin{eqnarray*}
v &=&\sqrt{\frac{\left( 1.4\right) \left( 8.31\frac{\unit{J}}{\unit{mol}%
\unit{K}}\right) \left( 273\unit{K}\right) }{0.02897\frac{\unit{kg}}{\unit{%
mol}}}}\sqrt{\frac{T}{T_{o}}} \\
&=&\allowbreak 331.\,\allowbreak 1\frac{\unit{m}}{\unit{s}}\sqrt{\frac{T}{%
T_{o}}} \\
&=&\allowbreak 331.\,\allowbreak 1\frac{\unit{m}}{\unit{s}}\sqrt{\frac{T}{273%
\unit{K}}}
\end{eqnarray*}%
where $T$ is in kelvins. We can change this to having $T$ in Celsius if we 
\begin{eqnarray*}
v &=&\allowbreak 331.\,\allowbreak 1\frac{\unit{m}}{\unit{s}}\sqrt{\frac{%
T_{o}+T_{c}}{T_{o}}} \\
&=&\allowbreak 331.\,\allowbreak 1\frac{\unit{m}}{\unit{s}}\sqrt{1+\frac{%
T_{c}}{T_{o}}}
\end{eqnarray*}%
which is what we used before. But now we can go back to 
\begin{equation*}
v=\sqrt{\frac{\gamma }{M}RT}
\end{equation*}%
and use this for any gas at (mostly) any temperature (so long as our gas is
close to being and ideal gas). We could even extend this to diatomic gasses
by using our extended $C_{%
%TCIMACRO{\TeXButton{V}{\ooalign{\hfil$V$\hfil\cr\kern0.1em--\hfil\cr}}}%
%BeginExpansion
\ooalign{\hfil$V$\hfil\cr\kern0.1em--\hfil\cr}%
%EndExpansion
}$ and $C_{P}$ values to get an extended $\gamma $ that includes vibration
and rotation of the diatomic molecules.

\chapter{35 Producing Useful Work 2.4.2, 2.4.3}

%TCIMACRO{\TeXButton{\chaptermark{Chapter 38}}{\chaptermark{Chapter 38}}}%
%BeginExpansion
\chaptermark{Chapter 38}%
%EndExpansion

So far we have talked about doing work on gas in cylinders. That may not
seem all that useful. But if we think that car engines have gas in cylinders
with pistons, then it may seem like this could be important.

%TCIMACRO{%
%\TeXButton{Fundamental Concepts}{\vspace{0.10in}\hspace{-0.37in}{\Large Fundamental Concepts\vspace{0.2in}}}}%
%BeginExpansion
\vspace{0.10in}\hspace{-0.37in}{\Large Fundamental Concepts\vspace{0.2in}}%
%EndExpansion

\begin{itemize}
\item Irreversability

\item It is impossible to make a cyclical machine that transfers energy
spontaneously from a colder object to a hotter object

\item $W_{\text{useful }}=-W_{\text{on the gas}}$

\item A heat reservoir is a source of thermal energy that can be tapped
without causing any measurable change in temperature.

\item Energy Transfer Diagrams

\item The Second Law of Thermodynamics

\item Heat Engines

\item Refrigerators

\item Efficiency of a heat engine, $\eta =\frac{W_{eng}}{\left\vert
Q_{h}\right\vert }$

\item Coefficient of Performance (COP) for refrigerators
\end{itemize}

\section{Irreversibility}

We should ask, are our special processes reversible? We found that for a
free expansion the process was not reversible.

\FRAME{dtbpF}{3.7896in}{2.2035in}{0pt}{}{}{Figure}{\special{language
"Scientific Word";type "GRAPHIC";maintain-aspect-ratio TRUE;display
"USEDEF";valid_file "T";width 3.7896in;height 2.2035in;depth
0pt;original-width 3.7421in;original-height 2.1629in;cropleft "0";croptop
"1";cropright "1";cropbottom "0";tempfilename
'SYUH1H0R.wmf';tempfile-properties "XPR";}}But how about for our special
processes: It would make our math easier if they were reversible.

For an isochoric process we know $w_{int}=0,$ $Q=nC_{V}\Delta T,$ $\Delta
E_{int}=nC_{V}\Delta T.$ It seems like we can change the pressure of a gas,
say by transferring energy by heat, and let the temperature rise, but then
remove the extra energy by heat (cool the cylinder of gas) back to the same
pressure all without changing the volume. It seems like this would be
reversible.

For an isobaric process we know $w_{int}=-P\Delta 
%TCIMACRO{\TeXButton{V}{\ooalign{\hfil$V$\hfil\cr\kern0.1em--\hfil\cr}}}%
%BeginExpansion
\ooalign{\hfil$V$\hfil\cr\kern0.1em--\hfil\cr}%
%EndExpansion
,$ $Q=nC_{P}\Delta T,$ $\Delta E_{int}=nC_{P}\Delta T-P\Delta 
%TCIMACRO{\TeXButton{V}{\ooalign{\hfil$V$\hfil\cr\kern0.1em--\hfil\cr}}}%
%BeginExpansion
\ooalign{\hfil$V$\hfil\cr\kern0.1em--\hfil\cr}%
%EndExpansion
.$ It seems like we could make a cylinder with a free floating position that
could go up or down. We could transfer energy by heat and let the volume
increase. But then transfer energy back out of the cylinder to restore it to
the same pressure. This seems reversible. For an isothermal process $w=-nRT%
\left[ \ln \frac{%
%TCIMACRO{\TeXButton{V}{\ooalign{\hfil$V$\hfil\cr\kern0.1em--\hfil\cr}}}%
%BeginExpansion
\ooalign{\hfil$V$\hfil\cr\kern0.1em--\hfil\cr}%
%EndExpansion
_{f}}{%
%TCIMACRO{\TeXButton{V}{\ooalign{\hfil$V$\hfil\cr\kern0.1em--\hfil\cr}}}%
%BeginExpansion
\ooalign{\hfil$V$\hfil\cr\kern0.1em--\hfil\cr}%
%EndExpansion
_{i}}\right] ,$ $Q=-w,$ $\Delta E_{int}=0.$ Could we do work on a cylinder
of gas that is in an ice water bath and transfer energy out by heat such
that the temperature stays the same, then pull the piston back out allowing
energy to transfer back in by heat, still keeping the temperature the same?
This also seems possible.

Finally, for an adiabatic process we now know that $w=\Delta E_{int},$ $Q=0$
and $\Delta E_{int}=nC_{v}\Delta t$ for an equivalent process. So we
insolate our cylinder and do some work. That increases the internal energy
(and the temperature) but we pull the piston back out and the internal
energy reverts back. These all seem like reversible processes.

But notice that in each case, though the cylinder contents came back to the
same state, the surroundings would not be the same. It each case energy was
transferred either work or by heat. And the surroundings had to supply that
work or energy transfer by heat. And we expect that in energy transfer by
heat that not only was our cylinder affected by a stove or hotplate, the air
in the room would also have a larger temperature after our energy transfer.
In the case of work, our diligent coresearcher that kept pushing and pulling
pistons is now tired. The surroundings are different. Causing a change in
the surroundings is something we need to think about, and it gives us our
first version of the \emph{Second Law of Thermodynamics.}

\section{Useful work.}

Suppose that we want to produce useful work from a thermodynamic device. For
building such a device we have a nomenclature issue.\FRAME{dhF}{4.8775in}{%
1.8273in}{0pt}{}{}{Figure}{\special{language "Scientific Word";type
"GRAPHIC";maintain-aspect-ratio TRUE;display "USEDEF";valid_file "T";width
4.8775in;height 1.8273in;depth 0pt;original-width 4.8239in;original-height
1.7902in;cropleft "0";croptop "1";cropright "1";cropbottom "0";tempfilename
'SUS73WC1.wmf';tempfile-properties "XPR";}}Consider our piston. We can do
work on the gas in our piston. And let's be specific, we want to do work on
the gas and that is what we have been calling $w_{int}.$ 
\begin{equation*}
w_{\text{on the gas}}=w_{int}
\end{equation*}%
If we add energy by heat, and the piston is free to move, then we do an
amount of work on the gas 
\begin{equation*}
w=-P\left( 
%TCIMACRO{\TeXButton{V}{\ooalign{\hfil$V$\hfil\cr\kern0.1em--\hfil\cr}}}%
%BeginExpansion
\ooalign{\hfil$V$\hfil\cr\kern0.1em--\hfil\cr}%
%EndExpansion
_{f}-%
%TCIMACRO{\TeXButton{V}{\ooalign{\hfil$V$\hfil\cr\kern0.1em--\hfil\cr}}}%
%BeginExpansion
\ooalign{\hfil$V$\hfil\cr\kern0.1em--\hfil\cr}%
%EndExpansion
_{i}\right)
\end{equation*}%
This is a negative amount of work, since $%
%TCIMACRO{\TeXButton{V}{\ooalign{\hfil$V$\hfil\cr\kern0.1em--\hfil\cr}}}%
%BeginExpansion
\ooalign{\hfil$V$\hfil\cr\kern0.1em--\hfil\cr}%
%EndExpansion
_{f}$ is lager than $%
%TCIMACRO{\TeXButton{V}{\ooalign{\hfil$V$\hfil\cr\kern0.1em--\hfil\cr}}}%
%BeginExpansion
\ooalign{\hfil$V$\hfil\cr\kern0.1em--\hfil\cr}%
%EndExpansion
_{i}$ and there is that minus sign out in front. We defined work done \emph{%
on the gas} as positive if the gas was being compressed, and the gas is not
being compressed. So the work done on the gas is negative. But we could also
define the work done \emph{by the gas}. That is the negative of the work
done on the gas. 
\begin{equation*}
w_{\text{on the gas}}=-w_{\text{by the gas}}
\end{equation*}

But think of a steam engine. We want the force \emph{by the gas} on the
piston, because this pushing of the gas on the piston is what we can use to
push the wheel of a steam engine, or a turbine, etc. That is negative work
on the gas. And it is inconvenient to have the useful work be negative.
Let's define the useful work done by a thermodynamic device as the negative
of work done on the gas.%
\begin{equation*}
w_{\text{useful }}=w_{\text{by the gas}}=-w_{\text{on the gas}}
\end{equation*}

But, $w_{\text{useful }}$is kind of a awkward name. Usually a machine that
does useful work is called an engine, so let's call our useful work $w_{eng}$
for \textquotedblleft work done by an engine.\textquotedblright

\begin{equation*}
w_{eng}=w_{\text{by the gas}}=-w_{int}
\end{equation*}%
Then, the useful work that our cylinder and piston do is minus the work we
have been calculating. We need to be careful with minus signs now.

\section{The Second Law of Thermodynamics}

We did a lot with our first law of thermodynamics (using our ideal gas law).
But there is a second law of thermodynamics. It is a little harder to state
in a nice equation. So we are going to build up to it a little at a time.
Our first version is due to a researcher called Clausius, so it is called
the Clausius' statement of the Second Law of Thermodynamics.

\begin{center}
\rule{300pt}{1pt}

It is impossible to make a cyclical machine that transfers energy
spontaneously from a colder object to a hotter object

\rule{300pt}{1pt}
\end{center}

This seems obvious. Refrigerators have to be plugged in and so do air
conditioners. They require an input of energy to run. But this is just what
we have been discussing. For a refrigerator or an air conditioner to work
the surroundings have to change. Energy in the form of electricity has to be
supplied. And that not only changes the surroundings it also changes our
budget. In our next lectures, we will form more mathematical statements of
the second law.

\section{Reservoirs}

Think of a water reservoir. In the western US we know what these are. It is
a large man-made body of water that is designed to provide irrigation water
for farms and homes.

\FRAME{dtbpF}{3.7697in}{2.834in}{0in}{}{}{Figure}{\special{language
"Scientific Word";type "GRAPHIC";maintain-aspect-ratio TRUE;display
"USEDEF";valid_file "T";width 3.7697in;height 2.834in;depth
0in;original-width 3.7222in;original-height 2.7916in;cropleft "0";croptop
"1";cropright "1";cropbottom "0";tempfilename
'TCSBPA01.wmf';tempfile-properties "XPR";}}

Think of what happens when you water your garden with water from the
reservoir. Some water does leave from the reservoir, but compared to the
whole reservoir of water, one garden watering makes very little difference.
Often there is more water coming into the reservoir, so really there is
negligible water loss from your individual garden.\footnote{%
There is some danger in thinking of reservoirs this way. People at one time
thought of the great lakes and oceans as reservoirs that were just too big
for us to have any effect on them. But there are a lot of us, and if we all
dump our trash in these reservoirs we can have a terrible effect on them. It
has taken years to clean up the great lakes so they are safe to swim in and
safe to fish. The effect from one person is not enough to change the state
of such a reservoir, but the effects from millions (or billions) is enough
to change it's state. We should follow God's counsel and be good stewards of
the Earth.}

A heat reservoir is like this. A hot reservoir is a source of thermal energy
that can be tapped without causing any measurable change in temperature to
the reservoir, itself. We can also have a cold reservoir that has a large
lack of thermal energy. A small amount of thermal energy transferred to this
cold reservoir won't measurably change its temperature.

The next figure is called an \emph{Energy Transfer Diagram}.

\FRAME{dhF}{1.1917in}{1.6873in}{0pt}{}{}{Figure}{\special{language
"Scientific Word";type "GRAPHIC";maintain-aspect-ratio TRUE;display
"USEDEF";valid_file "T";width 1.1917in;height 1.6873in;depth
0pt;original-width 1.8351in;original-height 2.6091in;cropleft "0";croptop
"1";cropright "1";cropbottom "0";tempfilename
'SUS73WC2.wmf';tempfile-properties "XPR";}}It shows energy moving by heat
from a hot reservoir to a cold reservoir.

Using our understanding of the second law of thermodynamics, we can see that
an amount of energy, $Q_{h}$, can flow by heat from the hot reservoir and an
amount of energy$,$ $Q_{c},$ can flow by heat to the cold reservoir. We
expect that if there is no other mechanism for dissipating the energy, then $%
Q_{c}=Q_{h}.$ We will define both $Q_{h}$ and $Q_{c}$ as positive
quantities, like vector magnitudes If they are negative, we will explicitly
write a negative sign in front of the variable $Q_{h}$ or $Q_{c}$.

Look at this energy transfer diagram.\FRAME{dhF}{2.1525in}{2.1482in}{0in}{}{%
}{Figure}{\special{language "Scientific Word";type
"GRAPHIC";maintain-aspect-ratio TRUE;display "USEDEF";valid_file "T";width
2.1525in;height 2.1482in;depth 0in;original-width 2.1136in;original-height
2.1093in;cropleft "0";croptop "1";cropright "1";cropbottom "0";tempfilename
'SUS73WC3.wmf';tempfile-properties "XPR";}}

This diagram says that energy flows from the cold reservoir to the hot
reservoir. From our experience, this doesn't happen. And our bran new Second
Law of Thermodynamics says this won't happen! We only spontaneously
experience energy coming from hot things and going to cold things. It would
take work to make this energy transfer happen, but we don't have any work on
the diagram. So this energy transfer diagram shows us a process that does
not happen.

Here is another energy transfer diagram\FRAME{dhF}{1.5316in}{1.9493in}{0pt}{%
}{}{Figure}{\special{language "Scientific Word";type
"GRAPHIC";maintain-aspect-ratio TRUE;display "USEDEF";valid_file "T";width
1.5316in;height 1.9493in;depth 0pt;original-width 2.6974in;original-height
3.442in;cropleft "0";croptop "1";cropright "1";cropbottom "0";tempfilename
'SUS73WC4.wmf';tempfile-properties "XPR";}}This diagram says that we do have
work, but that the work is directly converted into heat. This is a common
experience. Rub your hands together. Your hands warm up. The work has
created more internal energy, But that energy quickly dissipates into the
room The room temperature is not measurably affected by this increase in
internal energy of the air. The air in the room is a cold reservoir.

Here is another diagram.\FRAME{dhF}{1.9666in}{1.9631in}{0pt}{}{}{Figure}{%
\special{language "Scientific Word";type "GRAPHIC";maintain-aspect-ratio
TRUE;display "USEDEF";valid_file "T";width 1.9666in;height 1.9631in;depth
0pt;original-width 2.4613in;original-height 2.4569in;cropleft "0";croptop
"1";cropright "1";cropbottom "0";tempfilename
'SUS73WC5.wmf';tempfile-properties "XPR";}}This diagram says that energy
from the hot reservoir is converted directly into useful work. We expect
that we could make a transfer of energy produce work. Think of our piston in
the gas cylinder. But this diagram show a $100\%$ conversion of our energy
transfer by heat becoming useful work. This is not something we expect to
happen. Think about it, in a combustion engine we do get energy by heat by
burning fuel, but we don't just get useful work (the car moves). We also
heat up the outside of the engine. If you touch your engine after it has run
for a while you can burn your hand. Not all of the energy from the burning
fuel became useful work. We know the reason for this already. We need to
remove energy by heat to make our engine have a cyclic process so it doesn't
melt down. Some of the energy from the burning fuel must leave to the cold
reservoir as waste heat. But our diagram has no waste heat. This isn't
possible.

We could define a generic device that takes energy by heat and converts it
into useful work. We call such a device a \emph{heat engine}. The heat
engine described by the last diagram would be a perfect engine, converting
all the energy transferred from the hot reservoir into useful work. And from
our experience, this doesn't seem possible.

Our first law of thermodynamics told us about energy conservation, but it
did not tell us what thermodynamic processes are possible. We need another
law of thermodynamics that tells us what process we can expect to work. We
have found two that don't work. So which processes can work?

Our Clausius' statement of the second law of thermodynamics told us one of
these engine designs didn't work. And this is the job of the second law of
thermodynamics, to tell us what to expect will actually happen. In terms of
heat engines, we could add a new statement of our second law:

\begin{center}
\rule{300pt}{1pt}

Second Law of Thermodynamics: Perfect engines are not possible

\rule{300pt}{1pt}
\end{center}

\section{Heat Engines}

Using our energy transfer diagram, let's introduce a more practical engine.
Here is the diagram. \FRAME{dhF}{1.5454in}{1.8697in}{0pt}{}{}{Figure}{%
\special{language "Scientific Word";type "GRAPHIC";maintain-aspect-ratio
TRUE;display "USEDEF";valid_file "T";width 1.5454in;height 1.8697in;depth
0pt;original-width 2.655in;original-height 3.2197in;cropleft "0";croptop
"1";cropright "1";cropbottom "0";tempfilename
'SUS73WC6.wmf';tempfile-properties "XPR";}}We have added the engine
explicitly on the diagram. The engine takes energy from the heat reservoir
and turns it into useful work. But there is some energy that is waisted and
this energy ends up in the cold reservoir. Using our first law of
thermodynamics, the amount of energy that becomes work must be 
\begin{equation*}
w_{\text{engine }}=Q_{h}-Q_{c}
\end{equation*}%
We call $Q_{h}-Q_{c}$ the net heat%
\begin{equation*}
Q_{net}=Q_{h}-Q_{c}
\end{equation*}%
Notice that $Q_{h}$ is positive \emph{for the engine} and $Q_{c}$ is
negative \emph{for the engine.}

A useful heat engine is a device that takes in energy by heat and operates
in a cycle, If it did not periodically return to its original state, then it
could not keep going. The basics of the cycle are as follows

\begin{enumerate}
\item The working substance absorbs energy by heat from the high temperature
reservoir

\item Work is done by the engine

\item Energy is expelled by heat to a low temperature reservoir
\end{enumerate}

\subsection{Example: Steam engine}

An old fashioned train engine is a great example of a heat engine. The
picture on the left is such an old steam locomotive. On the right is a
schematic of how the steam engine works.

\FRAME{dtbpF}{4.804in}{1.7936in}{0pt}{}{}{Figure}{\special{language
"Scientific Word";type "GRAPHIC";maintain-aspect-ratio TRUE;display
"USEDEF";valid_file "T";width 4.804in;height 1.7936in;depth
0pt;original-width 6.5051in;original-height 2.4102in;cropleft "0";croptop
"1";cropright "1";cropbottom "0";tempfilename
'TCSC5X02.wmf';tempfile-properties "XPR";}}

In a steam engine

\begin{enumerate}
\item Water in the boiler absorbs energy by heat from burning fuel (hot
reservoir) and evaporates to steam

\item The steam does work by expanding against a piston

\item The steam cools and condenses releasing energy by heat to the outside
air (cold. reservoir) , and the liquid water returns to the boiler.
\end{enumerate}

Since the engine goes through a cyclical process, 
\begin{equation*}
\Delta E_{int}=0
\end{equation*}%
for the water in the tank. Its initial and final internal energies are the
same, so 
\begin{equation}
Q_{net}=w_{eng}
\end{equation}

The work done by the engine equals the \emph{net }energy absorbed by the
engine. The work is equal to the area enclosed by the curve of the PV
diagram.\FRAME{dhF}{2.1785in}{1.5601in}{0pt}{}{}{Figure}{\special{language
"Scientific Word";type "GRAPHIC";maintain-aspect-ratio TRUE;display
"USEDEF";valid_file "T";width 2.1785in;height 1.5601in;depth
0pt;original-width 4.0041in;original-height 2.8591in;cropleft "0";croptop
"1";cropright "1";cropbottom "0";tempfilename
'SUS73WC8.wmf';tempfile-properties "XPR";}}

Thermal efficiency is defined as the ratio of the net work done by the
engine during one cycle to the energy input at the higher temperature

\begin{eqnarray}
\eta &=&\frac{w_{eng}}{\left\vert Q_{h}\right\vert } \\
&=&\frac{Q_{net}}{Q_{h}} \\
&=&\frac{\left\vert Q_{h}\right\vert -\left\vert Q_{c}\right\vert }{Q_{h}} \\
&=&1-\frac{\left\vert Q_{c}\right\vert }{\left\vert Q_{h}\right\vert }
\end{eqnarray}%
We can think of this as 
\begin{equation}
\eta =\frac{\text{what you gain in work}}{\text{what you gave in energy by
heat}}
\end{equation}%
Sometimes efficiency is given the symbol $e,$ but this is easy to mistake
for the base of the natural logarithm, so let's use $\eta .$

For a car internal combustion engine 
\begin{equation}
\eta \approx 20\%
\end{equation}%
is a very good number (most student and faculty cars are far less efficient
than this).

Suppose that we have a $100\%$ efficient engine. We know that this is not
possible, but consider what this would mean. 
\begin{eqnarray}
\eta &=&1-\frac{\left\vert Q_{c}\right\vert }{\left\vert Q_{h}\right\vert }
\\
&=&1
\end{eqnarray}%
this can only happen when $\left\vert Q_{c}\right\vert =0$ which means that
no energy is transferred out of the engine by heat -- which is not really
possible. The fact that real engines have efficiencies much lower than $1$
leads to another statement of the second law of thermodynamics.

\FRAME{dhF}{2.9352in}{2.93in}{0pt}{}{}{Figure}{\special{language "Scientific
Word";type "GRAPHIC";maintain-aspect-ratio TRUE;display "USEDEF";valid_file
"T";width 2.9352in;height 2.93in;depth 0pt;original-width
2.8911in;original-height 2.8859in;cropleft "0";croptop "1";cropright
"1";cropbottom "0";tempfilename 'SUS73WC9.wmf';tempfile-properties "XPR";}}

\begin{center}
\rule{300pt}{1pt}

Second Law of Thermodynamics: It is impossible to construct a heat engine
that, operating in a cycle, inputs energy by heat from a hot reservoir and
converts the energy entirely into useful work.

\rule{300pt}{1pt}
\end{center}

This is called the Kelvin-Plank form of the second law, and it means that 
\begin{equation}
w_{eng}<\left\vert Q_{h}\right\vert
\end{equation}%
so the device represented by our drawing above is impossible to achieve.

\section{Heat Pumps and Refrigerators}

Let's look at a special form of a heat engine, a head pump.

\FRAME{dhF}{2.041in}{2.6507in}{0pt}{}{}{Figure}{\special{language
"Scientific Word";type "GRAPHIC";maintain-aspect-ratio TRUE;display
"USEDEF";valid_file "T";width 2.041in;height 2.6507in;depth
0pt;original-width 2.002in;original-height 2.6091in;cropleft "0";croptop
"1";cropright "1";cropbottom "0";tempfilename
'SUS73WCA.wmf';tempfile-properties "XPR";}}A heat pump is an engine run in
reverse. We wish to pump energy by heat from a cold reservoir to the hot
reservoir. To do this we must do work \emph{on }the heat pump. This is not a
natural process, so by adding in work we are adding energy to make it happen.

A window air conditioner unit is a good example of a heat pump. It removes
energy by heat from the inside of your apartment and transfers that energy
to the hot outside of your apparent. To make this happen, you need
mechanical device that does work.

A refrigerator is another example of a heat pump. It transfers energy by
heat from the cold interior to the warmer environment of your apartment
kitchen. It would be great if we could transfer heat from our cold reservoir
(room of a house or freezer compartment of a fridge) to a hot reservoir
without doing work, but this is not possible. This fact was discovered by
Clausius and lead to the Clausius' statement of the second law that we
already know, but we can update it so it is in terms of engines:

\begin{center}
\rule{300pt}{1pt}

Second Law of Thermodynamics: It is impossible to construct a cyclical
machine whose sole effect is to transfer energy continuously by heat from
one object to another object at a higher temperature without the input of
energy by work.

\rule{300pt}{1pt}
\end{center}

Or more simply, as we have already said: Energy does not transfer
spontaneously by heat from a cold object to a hot object.

\subsection{Coefficient of Performance}

How well a heat pump works can be described by the \emph{coefficient of
performance} (COP).%
\begin{eqnarray}
COP_{h} &=&\frac{\text{energy transferred at high temperature}}{\text{work
done on heat pump}} \\
&=&\frac{\left\vert Q_{h}\right\vert }{w}
\end{eqnarray}

$COP$ is similar to efficiency. Think of our air conditioner. $Q_{h}$ is
typically higher than $w.$ That is, we usually transfer more energy by heat
to the outside world than we provide in work. The outside part of the air
conditioner does indeed get hotter than the outside air temperature. That is
why energy will transfer by heat. Values of $COP$ are generally greater than 
$1,$ though it is possible for them to be less than $1.$ Of course we would
like the $COP$ to be as high as possible

\FRAME{dtbpFU}{2.469in}{2.5581in}{0pt}{\Qcb{{\protect\small Refrigerator
Design. Red = Gas at high pressure and very high temperature. Pink = Liquid
at high pressure and high temperature. Blue = Liquid at low pressure and
very low temperature. Light Blue = Gas at low pressure and low temperature.
(Drawing courtesy Jahobr)}}}{}{Figure}{\special{language "Scientific
Word";type "GRAPHIC";maintain-aspect-ratio TRUE;display "USEDEF";valid_file
"T";width 2.469in;height 2.5581in;depth 0pt;original-width
2.4871in;original-height 2.5776in;cropleft "0";croptop "1";cropright
"1";cropbottom "0";tempfilename 'SZEJYT03.wmf';tempfile-properties "XPR";}}

For refrigeration, the important thing is how much energy we transfer out of
the cool area, $Q_{c}$, so for cooling the $COP$ is 
\begin{equation}
COP_{c}=\frac{\left\vert Q_{c}\right\vert }{w}
\end{equation}%
A good refrigerator should have a high COP, Typical values are 5 or 6.

Now we know a about heat engines, we should look at how a more realistic
heat engine works. In our next lecture we will take on an engine like you
might find on a lawnmower or a go-cart.

\chapter{36 Carnot, the Limit of Efficiency for Heat Engines 2.4.4, 2.4.5}

Last lecture we studied a somewhat realistic heat engine. We found we got
something like a $56.5\%$ efficiency. But is this really realistic? Could
you build such an engine? We need a way to tell what the maximum efficiency
for a heat engine would be. Then, if our design exceeds this maximum
efficiency, we know we have a mistake in our design. Our goal in a heat
engine problem is always to find for each process the values of $\Delta
E_{int},$ $Q,$ $w_{int},$ and for the whole engine $w_{eng},$ $Q_{h}$, $%
Q_{c} $. We also want to find the ideal gas state variables $P,$ $%
%TCIMACRO{\TeXButton{V}{\ooalign{\hfil$V$\hfil\cr\kern0.1em--\hfil\cr}}}%
%BeginExpansion
\ooalign{\hfil$V$\hfil\cr\kern0.1em--\hfil\cr}%
%EndExpansion
,$ $n,$ and $T$ at the transition points between processes. But we need the
maximum efficiency to know if our result is realistic. Let's use all of
these techniques to find our maximum efficiency.

%TCIMACRO{%
%\TeXButton{Fundamental Concepts}{\vspace{0.10in}\hspace{-0.37in}{\Large Fundamental Concepts\vspace{0.2in}}}}%
%BeginExpansion
\vspace{0.10in}\hspace{-0.37in}{\Large Fundamental Concepts\vspace{0.2in}}%
%EndExpansion

\begin{itemize}
\item The Carnot cycle represents the maximum theoretical limit of
efficiency for a heat engine.

\item The efficiency of the Carnot cycle depends only on the two extreme
temperatures $\eta _{C}=1-\frac{T_{c}}{T_{h}}.$

\item An engine design that is more efficient that a Carnot engine must have
a fundamental flaw in it's design.
\end{itemize}

%TCIMACRO{%
%\TeXButton{Question 123.19.1}{\marginpar {
%\hspace{-0.5in}
%\begin{minipage}[t]{1in}
%\small{Question 123.19.1}
%\end{minipage}
%}}}%
%BeginExpansion
\marginpar {
\hspace{-0.5in}
\begin{minipage}[t]{1in}
\small{Question 123.19.1}
\end{minipage}
}%
%EndExpansion

\section{Carnot's Theorem}

We now know that we can design actual engines that power cars and useful
machines by combining several of our special thermodynamic process into
larger, more complex processes that form a cycle. You might guess that there
are more than just isochoric, isobaric, isothermal, and adiabatic processes
that exist. A good question is what processes should we combine to form our
engine?

Past researchers have asked this question over and over and many new engines
have been developed. Of course we want our engines to be as efficient as
possible, and the second law of thermodynamics says the engines can't be
100\% efficient. But is there a \textquotedblleft most
efficient\textquotedblright\ engine, one that no other engine can beat in
efficiency? If there is, we could use this to judge new engine designs. Any
engine with a better efficiency would be a mistake that won't work. And
designs that approach this maximum efficiency would be a possible
improvement in technology. There is such a maximally efficient engine.

Carnot was an Engineer. He thought of an engine that would operate on an
ideal reversible cycle. He hit upon the most efficient cycle possible.
Sadly, it is not possible to create his engine in practice. Carnot's engine
assumed no friction and ideal thermodynamic processes. So technically it is
more efficient than is actually possible.

So why is it useful to know about Carnot's engine design? The Carnot cycle
provides us with a useful upper limit for all practical engines that is easy
to calculate. To tell if a design is practical, we can compare it to the
Carnot cycle. If our design looks like it will out perform the Carnot
engine, we know we have done something wrong. Or conversely, if our system
design requires an engine to be better than a Carnot engine, we need to go
back to the drawing board.

Knowing just this much we can state Carnot's theorem.

\begin{center}
\rule{300pt}{1pt}

No real heat engine operating between two energy reservoirs can be more

efficient than a Carnot engine operating between the same two reservoirs

\rule{300pt}{1pt}

All real engines are less efficient than a Carnot engine because they do not

operate through a reversible cycle. The efficiency of a real engine is

reduced by friction, energy losses through conduction, etc.

\rule{300pt}{1pt}
\end{center}

\section{The Carnot Cycle}

Carnot found a set of special processes that when combined made a
theoretical heat engine. And he proved mathematically that this particular
combination is theoretically the most efficient possible. Let's take a look
at this engine design.

The Carnot cycle is shown in the figure. It consists of the following parts:

\begin{enumerate}
\item Process $A\rightarrow B$ is an isothermal expansion

\item Process $B\rightarrow C$ is an adiabatic expansion

\item Process $C\rightarrow D$ is an isothermal contraction

\item Process $D\rightarrow A$ is an adiabatic compression
\end{enumerate}

\FRAME{dhF}{3.774in}{2.623in}{0pt}{}{}{Figure}{\special{language "Scientific
Word";type "GRAPHIC";maintain-aspect-ratio TRUE;display "USEDEF";valid_file
"T";width 3.774in;height 2.623in;depth 0pt;original-width
3.7256in;original-height 2.5815in;cropleft "0";croptop "1";cropright
"1";cropbottom "0";tempfilename 'SUS73WCJ.wmf';tempfile-properties "XPR";}}

Let's take an example. Let's choose a Carnot engine that operates between
the same two extreme temperatures as our Otto cycle example from our last
lecture. 
\begin{eqnarray*}
T_{c} &=&293\unit{K} \\
T_{h} &=&1023\unit{K}
\end{eqnarray*}%
Let's suppose the carnot engine is made from a piston with a monotonic idea
gas. Let's choose the minimum and maximum volumes to be 
\begin{eqnarray*}
%TCIMACRO{\TeXButton{V}{\ooalign{\hfil$V$\hfil\cr\kern0.1em--\hfil\cr}}}%
%BeginExpansion
\ooalign{\hfil$V$\hfil\cr\kern0.1em--\hfil\cr}%
%EndExpansion
_{A} &=&0.001\unit{m}^{3} \\
%TCIMACRO{\TeXButton{V}{\ooalign{\hfil$V$\hfil\cr\kern0.1em--\hfil\cr}}}%
%BeginExpansion
\ooalign{\hfil$V$\hfil\cr\kern0.1em--\hfil\cr}%
%EndExpansion
_{C} &=&0.015\unit{m}^{3}
\end{eqnarray*}%
and let's have 
\begin{equation*}
n=1\unit{mol}
\end{equation*}%
of gas just to make the math easy. For these numbers, our graph looks like
this:\FRAME{dtbpF}{4.0269in}{2.7612in}{0pt}{}{}{Figure}{\special{language
"Scientific Word";type "GRAPHIC";maintain-aspect-ratio TRUE;display
"USEDEF";valid_file "T";width 4.0269in;height 2.7612in;depth
0pt;original-width 4.0748in;original-height 2.7851in;cropleft "0";croptop
"1";cropright "1";cropbottom "0";tempfilename
'SUS73WCK.wmf';tempfile-properties "XPR";}}Our job is to find $P,$ $%
%TCIMACRO{\TeXButton{V}{\ooalign{\hfil$V$\hfil\cr\kern0.1em--\hfil\cr}}}%
%BeginExpansion
\ooalign{\hfil$V$\hfil\cr\kern0.1em--\hfil\cr}%
%EndExpansion
,$ $T,$ $w$, $Q,$ and $\Delta E_{int}$ for each part of the cycle and $Q_{h}$%
, $Q_{c},$ and $w_{eng}$ for the whole cycle.

Let's take each process separately like we did for the Otto cycle.

\section{Process $A\rightarrow B$ isothermal expansion}

The gas is placed in contact with the high temperature reservoir, $T_{h}.$
The gas absorbs heat $|Q_{h}|.$ The gas does work $w_{AB}$ in raising the
piston\FRAME{dtbpF}{4.9644in}{3.3288in}{0pt}{}{}{Figure}{\special{language
"Scientific Word";type "GRAPHIC";maintain-aspect-ratio TRUE;display
"USEDEF";valid_file "T";width 4.9644in;height 3.3288in;depth
0pt;original-width 5.0301in;original-height 3.3634in;cropleft "0";croptop
"1";cropright "1";cropbottom "0";tempfilename
'SUS73WCL.wmf';tempfile-properties "XPR";}}

Remember for isothermal processes we can write%
\begin{equation*}
\frac{P_{i}%
%TCIMACRO{\TeXButton{V}{\ooalign{\hfil$V$\hfil\cr\kern0.1em--\hfil\cr}}}%
%BeginExpansion
\ooalign{\hfil$V$\hfil\cr\kern0.1em--\hfil\cr}%
%EndExpansion
_{i}}{P_{f}%
%TCIMACRO{\TeXButton{V}{\ooalign{\hfil$V$\hfil\cr\kern0.1em--\hfil\cr}}}%
%BeginExpansion
\ooalign{\hfil$V$\hfil\cr\kern0.1em--\hfil\cr}%
%EndExpansion
_{f}}=\frac{nRT}{nRT}
\end{equation*}%
because the temperature does not change, so%
\begin{equation*}
P_{i}%
%TCIMACRO{\TeXButton{V}{\ooalign{\hfil$V$\hfil\cr\kern0.1em--\hfil\cr}}}%
%BeginExpansion
\ooalign{\hfil$V$\hfil\cr\kern0.1em--\hfil\cr}%
%EndExpansion
_{i}=P_{f}%
%TCIMACRO{\TeXButton{V}{\ooalign{\hfil$V$\hfil\cr\kern0.1em--\hfil\cr}}}%
%BeginExpansion
\ooalign{\hfil$V$\hfil\cr\kern0.1em--\hfil\cr}%
%EndExpansion
_{f}
\end{equation*}%
Along Path $AB$ this would be%
\begin{equation*}
P_{A}%
%TCIMACRO{\TeXButton{V}{\ooalign{\hfil$V$\hfil\cr\kern0.1em--\hfil\cr}}}%
%BeginExpansion
\ooalign{\hfil$V$\hfil\cr\kern0.1em--\hfil\cr}%
%EndExpansion
_{A}=P_{B}%
%TCIMACRO{\TeXButton{V}{\ooalign{\hfil$V$\hfil\cr\kern0.1em--\hfil\cr}}}%
%BeginExpansion
\ooalign{\hfil$V$\hfil\cr\kern0.1em--\hfil\cr}%
%EndExpansion
_{B}
\end{equation*}%
And we know for an isothermal process that 
\begin{equation*}
w_{AB}=nRT_{A}\ln \left( \frac{%
%TCIMACRO{\TeXButton{V}{\ooalign{\hfil$V$\hfil\cr\kern0.1em--\hfil\cr}}}%
%BeginExpansion
\ooalign{\hfil$V$\hfil\cr\kern0.1em--\hfil\cr}%
%EndExpansion
_{A}}{%
%TCIMACRO{\TeXButton{V}{\ooalign{\hfil$V$\hfil\cr\kern0.1em--\hfil\cr}}}%
%BeginExpansion
\ooalign{\hfil$V$\hfil\cr\kern0.1em--\hfil\cr}%
%EndExpansion
_{B}}\right)
\end{equation*}%
and that%
\begin{equation*}
\Delta E_{int}=0
\end{equation*}%
so%
\begin{equation*}
\left\vert Q_{AB}\right\vert =\left\vert -w_{AB}\right\vert
\end{equation*}%
and we can identify this as $\left\vert Q_{h}\right\vert $ and $T_{A}=T_{h}$

We can complete state $A$ using the ideal gas law%
\begin{eqnarray*}
P_{A} &=&\frac{nRT_{h}}{%
%TCIMACRO{\TeXButton{V}{\ooalign{\hfil$V$\hfil\cr\kern0.1em--\hfil\cr}}}%
%BeginExpansion
\ooalign{\hfil$V$\hfil\cr\kern0.1em--\hfil\cr}%
%EndExpansion
_{A}}=\frac{\left( 1\unit{mol}\right) \left( 8.314\frac{\unit{J}}{\unit{mol}%
\unit{K}}\right) \left( 1023\unit{K}\right) }{0.001\unit{m}^{3}} \\
&=&8.\,\allowbreak 505\,2\times 10^{6}\unit{Pa}
\end{eqnarray*}

To complete state $B$ we need the volume and the pressure. To find the
pressure at point $B$ we can consider that point $B$ is on the from both
path $AB$ and on the line from path $BC.$ At point $B$ the lines cross so
both the adiabat and the isotherm have the same pressure and volume. So we
know from the idea gas law 
\begin{equation*}
P_{B}=\frac{nRT_{h}}{%
%TCIMACRO{\TeXButton{V}{\ooalign{\hfil$V$\hfil\cr\kern0.1em--\hfil\cr}}}%
%BeginExpansion
\ooalign{\hfil$V$\hfil\cr\kern0.1em--\hfil\cr}%
%EndExpansion
_{B}}
\end{equation*}%
and from our equations for adiabatic processes 
\begin{equation*}
P_{B}=\frac{K_{BC}}{%
%TCIMACRO{\TeXButton{V}{\ooalign{\hfil$V$\hfil\cr\kern0.1em--\hfil\cr}}}%
%BeginExpansion
\ooalign{\hfil$V$\hfil\cr\kern0.1em--\hfil\cr}%
%EndExpansion
_{B}^{\gamma }}
\end{equation*}%
where the constant $K_{BC}=P_{B}%
%TCIMACRO{\TeXButton{V}{\ooalign{\hfil$V$\hfil\cr\kern0.1em--\hfil\cr}}}%
%BeginExpansion
\ooalign{\hfil$V$\hfil\cr\kern0.1em--\hfil\cr}%
%EndExpansion
_{B}^{\gamma }$ from our adiabatic equations. But we don't know $K_{BC}$, $%
P_{B},$ or $%
%TCIMACRO{\TeXButton{V}{\ooalign{\hfil$V$\hfil\cr\kern0.1em--\hfil\cr}}}%
%BeginExpansion
\ooalign{\hfil$V$\hfil\cr\kern0.1em--\hfil\cr}%
%EndExpansion
_{B}$

We do know that at $C$ 
\begin{equation*}
P_{C}=\frac{nRT_{c}}{%
%TCIMACRO{\TeXButton{V}{\ooalign{\hfil$V$\hfil\cr\kern0.1em--\hfil\cr}}}%
%BeginExpansion
\ooalign{\hfil$V$\hfil\cr\kern0.1em--\hfil\cr}%
%EndExpansion
_{C}}
\end{equation*}%
and that%
\begin{equation*}
P_{C}=\frac{K_{BC}}{%
%TCIMACRO{\TeXButton{V}{\ooalign{\hfil$V$\hfil\cr\kern0.1em--\hfil\cr}}}%
%BeginExpansion
\ooalign{\hfil$V$\hfil\cr\kern0.1em--\hfil\cr}%
%EndExpansion
_{c}^{\gamma }}
\end{equation*}%
and at point $C$ and we do know $%
%TCIMACRO{\TeXButton{V}{\ooalign{\hfil$V$\hfil\cr\kern0.1em--\hfil\cr}}}%
%BeginExpansion
\ooalign{\hfil$V$\hfil\cr\kern0.1em--\hfil\cr}%
%EndExpansion
_{C,}$ so we can find $K_{BC}.$ We set the two expressions for $P_{C}$ equal%
\begin{equation*}
\frac{nRT_{c}}{%
%TCIMACRO{\TeXButton{V}{\ooalign{\hfil$V$\hfil\cr\kern0.1em--\hfil\cr}}}%
%BeginExpansion
\ooalign{\hfil$V$\hfil\cr\kern0.1em--\hfil\cr}%
%EndExpansion
_{C}}=\frac{K_{BC}}{%
%TCIMACRO{\TeXButton{V}{\ooalign{\hfil$V$\hfil\cr\kern0.1em--\hfil\cr}}}%
%BeginExpansion
\ooalign{\hfil$V$\hfil\cr\kern0.1em--\hfil\cr}%
%EndExpansion
_{C}^{\gamma }}
\end{equation*}%
\begin{eqnarray*}
K_{BC} &=&\frac{nRT_{c}}{%
%TCIMACRO{\TeXButton{V}{\ooalign{\hfil$V$\hfil\cr\kern0.1em--\hfil\cr}}}%
%BeginExpansion
\ooalign{\hfil$V$\hfil\cr\kern0.1em--\hfil\cr}%
%EndExpansion
_{C}}%
%TCIMACRO{\TeXButton{V}{\ooalign{\hfil$V$\hfil\cr\kern0.1em--\hfil\cr}}}%
%BeginExpansion
\ooalign{\hfil$V$\hfil\cr\kern0.1em--\hfil\cr}%
%EndExpansion
_{C}^{\gamma } \\
&=&nRT_{c}%
%TCIMACRO{\TeXButton{V}{\ooalign{\hfil$V$\hfil\cr\kern0.1em--\hfil\cr}}}%
%BeginExpansion
\ooalign{\hfil$V$\hfil\cr\kern0.1em--\hfil\cr}%
%EndExpansion
_{C}^{\gamma -1}
\end{eqnarray*}%
and we have all these numbers.%
\begin{eqnarray*}
K_{BC} &=&nRT_{c}%
%TCIMACRO{\TeXButton{V}{\ooalign{\hfil$V$\hfil\cr\kern0.1em--\hfil\cr}}}%
%BeginExpansion
\ooalign{\hfil$V$\hfil\cr\kern0.1em--\hfil\cr}%
%EndExpansion
_{C}^{\gamma -1}=\left( 1\unit{mol}\right) \left( 8.314\frac{\unit{J}}{\unit{%
mol}\unit{K}}\right) \left( 293\unit{K}\right) \left( 0.015\unit{m}%
^{3}\right) ^{\frac{2}{3}} \\
&=&148.\,\allowbreak 16\unit{J}\unit{m}^{2}
\end{eqnarray*}
Then, moving back to point $B$ we can set the two equations for $P_{B}$
equal 
\begin{eqnarray*}
P_{B} &=&\frac{nRT_{h}}{%
%TCIMACRO{\TeXButton{V}{\ooalign{\hfil$V$\hfil\cr\kern0.1em--\hfil\cr}}}%
%BeginExpansion
\ooalign{\hfil$V$\hfil\cr\kern0.1em--\hfil\cr}%
%EndExpansion
_{B}} \\
P_{B} &=&\frac{K_{BC}}{%
%TCIMACRO{\TeXButton{V}{\ooalign{\hfil$V$\hfil\cr\kern0.1em--\hfil\cr}}}%
%BeginExpansion
\ooalign{\hfil$V$\hfil\cr\kern0.1em--\hfil\cr}%
%EndExpansion
_{B}^{\gamma }}
\end{eqnarray*}%
to get 
\begin{equation*}
\frac{nRT_{h}}{%
%TCIMACRO{\TeXButton{V}{\ooalign{\hfil$V$\hfil\cr\kern0.1em--\hfil\cr}}}%
%BeginExpansion
\ooalign{\hfil$V$\hfil\cr\kern0.1em--\hfil\cr}%
%EndExpansion
_{B}}=\frac{K_{BC}}{%
%TCIMACRO{\TeXButton{V}{\ooalign{\hfil$V$\hfil\cr\kern0.1em--\hfil\cr}}}%
%BeginExpansion
\ooalign{\hfil$V$\hfil\cr\kern0.1em--\hfil\cr}%
%EndExpansion
_{B}^{\gamma }}
\end{equation*}%
then 
\begin{equation*}
\frac{%
%TCIMACRO{\TeXButton{V}{\ooalign{\hfil$V$\hfil\cr\kern0.1em--\hfil\cr}}}%
%BeginExpansion
\ooalign{\hfil$V$\hfil\cr\kern0.1em--\hfil\cr}%
%EndExpansion
_{B}^{\gamma }}{%
%TCIMACRO{\TeXButton{V}{\ooalign{\hfil$V$\hfil\cr\kern0.1em--\hfil\cr}}}%
%BeginExpansion
\ooalign{\hfil$V$\hfil\cr\kern0.1em--\hfil\cr}%
%EndExpansion
_{B}}=\frac{K_{BC}}{nRT_{h}}
\end{equation*}%
\smallskip \strut \smallskip 
\begin{equation*}
%TCIMACRO{\TeXButton{V}{\ooalign{\hfil$V$\hfil\cr\kern0.1em--\hfil\cr}}}%
%BeginExpansion
\ooalign{\hfil$V$\hfil\cr\kern0.1em--\hfil\cr}%
%EndExpansion
_{B}^{\gamma -1}=\frac{K_{BC}}{nRT_{h}}
\end{equation*}%
Which gives us an awkward 
\begin{equation*}
%TCIMACRO{\TeXButton{V}{\ooalign{\hfil$V$\hfil\cr\kern0.1em--\hfil\cr}}}%
%BeginExpansion
\ooalign{\hfil$V$\hfil\cr\kern0.1em--\hfil\cr}%
%EndExpansion
_{B}=\sqrt[\gamma -1]{\frac{K_{BC}}{nRT_{h}}}
\end{equation*}%
but it works.%
\begin{eqnarray*}
%TCIMACRO{\TeXButton{V}{\ooalign{\hfil$V$\hfil\cr\kern0.1em--\hfil\cr}}}%
%BeginExpansion
\ooalign{\hfil$V$\hfil\cr\kern0.1em--\hfil\cr}%
%EndExpansion
_{B} &=&\sqrt[\frac{2}{3}]{\frac{148.\,\allowbreak 16\unit{J}\unit{m}^{2}}{%
\left( 1\unit{mol}\right) \left( 8.314\frac{\unit{J}}{\unit{mol}\unit{K}}%
\right) \left( 1023\unit{K}\right) }} \\
&=&2.\,\allowbreak 299\,2\times 10^{-3}\unit{m}^{3}
\end{eqnarray*}%
And we are back to the ideal gas law to find the pressure%
\begin{eqnarray*}
P_{B} &=&\frac{nRT_{B}}{%
%TCIMACRO{\TeXButton{V}{\ooalign{\hfil$V$\hfil\cr\kern0.1em--\hfil\cr}}}%
%BeginExpansion
\ooalign{\hfil$V$\hfil\cr\kern0.1em--\hfil\cr}%
%EndExpansion
_{B}} \\
&=&\frac{\left( 1\unit{mol}\right) \left( 8.314\frac{\unit{J}}{\unit{mol}%
\unit{K}}\right) \left( 1023\unit{K}\right) }{2.\,\allowbreak 299\,2\times
10^{-3}\unit{m}^{3}} \\
&=&\allowbreak 3.\,\allowbreak 699\,2\times 10^{6}\unit{Pa}
\end{eqnarray*}%
Armed with all of this, we can compute the work (and heat) for process $AB.$%
\begin{equation*}
w_{AB}=nRT_{h}\ln \left( \frac{%
%TCIMACRO{\TeXButton{V}{\ooalign{\hfil$V$\hfil\cr\kern0.1em--\hfil\cr}}}%
%BeginExpansion
\ooalign{\hfil$V$\hfil\cr\kern0.1em--\hfil\cr}%
%EndExpansion
_{A}}{%
%TCIMACRO{\TeXButton{V}{\ooalign{\hfil$V$\hfil\cr\kern0.1em--\hfil\cr}}}%
%BeginExpansion
\ooalign{\hfil$V$\hfil\cr\kern0.1em--\hfil\cr}%
%EndExpansion
_{B}}\right)
\end{equation*}%
\begin{eqnarray*}
w_{AB} &=&\left( 1\unit{mol}\right) \left( 8.314\frac{\unit{J}}{\unit{mol}%
\unit{K}}\right) \left( \allowbreak 1023\unit{K}\right) \ln \left( \frac{%
0.001\unit{m}^{3}}{2.\,\allowbreak 299\,2\times 10^{-3}\unit{m}^{3}}\right)
\\
&=&-7081.\,\allowbreak 1\unit{J}
\end{eqnarray*}%
\begin{equation*}
Q_{AB}=-w_{AB}=7081.\,\allowbreak 1\unit{J}
\end{equation*}%
and this heat is positive so it must be part of $Q_{h}$

\begin{equation*}
\begin{tabular}{|c|c|c|c|c|c|c|}
\hline
\textbf{State} & $\mathbf{T}\left( \unit{K}\right) $ & $\mathbf{P}\left( 
\unit{Pa}\right) $ & $%
%TCIMACRO{\TeXButton{V}{\ooalign{\hfil$V$\hfil\cr\kern0.1em--\hfil\cr}}}%
%BeginExpansion
\ooalign{\hfil$V$\hfil\cr\kern0.1em--\hfil\cr}%
%EndExpansion
\left( \unit{m}^{3}\right) $ & $\mathbf{n}\left( \unit{mol}\right) $ & 
\textbf{--} & \textbf{--} \\ \hline
$\mathbf{A}$ & $1023$ & $8.\,\allowbreak 505\,2\times 10^{6}$ & $0.001$ & $1$
& \textbf{--} & \textbf{--} \\ \hline
$\mathbf{B}$ & $1023$ & $3.\,\allowbreak 699\,2\times 10^{6}$ & $%
2.\,\allowbreak 299\,2\times 10^{-3}$ & $1$ & \textbf{--} & \textbf{--} \\ 
\hline
\textbf{Process} & $\mathbf{Q}\left( \unit{J}\right) $ & $\mathbf{w}%
_{int}\left( \unit{J}\right) $ & $\mathbf{\Delta E}_{int}\left( \unit{J}%
\right) $ & $\mathbf{w}_{eng}\left( \unit{J}\right) $ & $\mathbf{Q}_{h}$ & $%
\mathbf{Q}_{c}$ \\ \hline
$\mathbf{AB}$ & $7081.\,\allowbreak 1$ & $-7081.\,\allowbreak 1$ & $0$ & $%
7081.\,\allowbreak 1$ & $7081.\,\allowbreak 1$ & $0$ \\ \hline
\end{tabular}%
\end{equation*}

\section{Process $B\rightarrow C$ adiabatic expansion}

The base of the cylinder is replaced by a thermally nonconducting wall. No
energy enters or leaves the system by heat. 
\begin{equation*}
Q_{BC}=0
\end{equation*}

The temperature falls from $T_{h}$ to $T_{c}.$ The there is work $w_{BC}$
done on the gas.\FRAME{dtbpF}{5.3086in}{3.822in}{0pt}{}{}{Figure}{\special%
{language "Scientific Word";type "GRAPHIC";maintain-aspect-ratio
TRUE;display "USEDEF";valid_file "T";width 5.3086in;height 3.822in;depth
0pt;original-width 5.3813in;original-height 3.8663in;cropleft "0";croptop
"1";cropright "1";cropbottom "0";tempfilename
'SUS73WCM.wmf';tempfile-properties "XPR";}}

Remember that for an adiabatic process 
\begin{equation*}
P_{i}%
%TCIMACRO{\TeXButton{V}{\ooalign{\hfil$V$\hfil\cr\kern0.1em--\hfil\cr}}}%
%BeginExpansion
\ooalign{\hfil$V$\hfil\cr\kern0.1em--\hfil\cr}%
%EndExpansion
_{i}^{\gamma }=P_{f}%
%TCIMACRO{\TeXButton{V}{\ooalign{\hfil$V$\hfil\cr\kern0.1em--\hfil\cr}}}%
%BeginExpansion
\ooalign{\hfil$V$\hfil\cr\kern0.1em--\hfil\cr}%
%EndExpansion
_{f}^{\gamma }
\end{equation*}%
and we usually find a more convenient method for calculating $\Delta
E_{int}. $ We choose an isochoric path, even though it isn't right. But $%
\Delta E_{int}$ is the same for any path between the same two temperatures,
so

\begin{equation*}
\Delta E_{int}=nC_{V}\Delta T
\end{equation*}%
We can find the pressure at $C$ using the ideal gas law%
\begin{eqnarray*}
P_{C} &=&\frac{nRT_{c}}{%
%TCIMACRO{\TeXButton{V}{\ooalign{\hfil$V$\hfil\cr\kern0.1em--\hfil\cr}}}%
%BeginExpansion
\ooalign{\hfil$V$\hfil\cr\kern0.1em--\hfil\cr}%
%EndExpansion
_{C}}=\frac{\left( 1\unit{mol}\right) \left( 8.314\frac{\unit{J}}{\unit{mol}%
\unit{K}}\right) \left( \allowbreak 293\unit{K}\right) }{0.015\unit{m}^{3}}
\\
&=&1.\,\allowbreak 624\times 10^{5}\unit{Pa}
\end{eqnarray*}%
and%
\begin{eqnarray*}
\Delta E_{int} &=&\left( 1\unit{mol}\right) \frac{3}{2}\left( 8.314\frac{%
\unit{J}}{\unit{mol}\unit{K}}\right) \left( 293\unit{K}-1023\unit{K}\right)
\\
&=&-9103.\,\allowbreak 8\unit{J}
\end{eqnarray*}%
We can summarize all this.%
\begin{equation*}
\begin{tabular}{|c|c|c|c|c|c|c|}
\hline
\textbf{State} & $\mathbf{T}\left( \unit{K}\right) $ & $\mathbf{P}\left( 
\unit{Pa}\right) $ & $%
%TCIMACRO{\TeXButton{V}{\ooalign{\hfil$V$\hfil\cr\kern0.1em--\hfil\cr}}}%
%BeginExpansion
\ooalign{\hfil$V$\hfil\cr\kern0.1em--\hfil\cr}%
%EndExpansion
\left( \unit{m}^{3}\right) $ & $\mathbf{n}\left( \unit{mol}\right) $ & 
\textbf{--} & \textbf{--} \\ \hline
$\mathbf{B}$ & $1023$ & $3.\,\allowbreak 699\,2\times 10^{6}$ & $%
2.\,\allowbreak 299\,2\times 10^{-3}$ & $1$ & \textbf{--} & \textbf{--} \\ 
\hline
$\mathbf{C}$ & $273$ & $1.\,\allowbreak 624\times 10^{5}$ & $0.015\unit{m}%
^{3}$ & $1$ & \textbf{--} & \textbf{--} \\ \hline
\textbf{Process} & $\mathbf{Q}\left( \unit{J}\right) $ & $\mathbf{w}%
_{int}\left( \unit{J}\right) $ & $\mathbf{\Delta E}_{int}\left( \unit{J}%
\right) $ & $\mathbf{w}_{eng}\left( \unit{J}\right) $ & $\mathbf{Q}_{h}$ & $%
\mathbf{Q}_{c}$ \\ \hline
$\mathbf{BC}$ & $0$ & $-9103.\,\allowbreak 8$ & $-9103.\,\allowbreak 8$ & $%
9103.\,\allowbreak 8$ & $0$ & $0$ \\ \hline
\end{tabular}%
\end{equation*}

\section{Process $C\rightarrow D$ isothermal compression}

The gas is placed in contact with the cold temperature reservoir. The gas
expels energy $\left\vert Q_{c}\right\vert $ and work $W_{CD}$ is done on
the gas.\FRAME{dtbpF}{5.6518in}{3.9905in}{0pt}{}{}{Figure}{\special{language
"Scientific Word";type "GRAPHIC";maintain-aspect-ratio TRUE;display
"USEDEF";valid_file "T";width 5.6518in;height 3.9905in;depth
0pt;original-width 5.7317in;original-height 4.0375in;cropleft "0";croptop
"1";cropright "1";cropbottom "0";tempfilename
'SUS73WCN.wmf';tempfile-properties "XPR";}}The process is isothermal, so $%
T_{C}=T_{D}$ and $nRT_{C}=nRT_{D}$ so we have 
\begin{equation}
P_{C}%
%TCIMACRO{\TeXButton{V}{\ooalign{\hfil$V$\hfil\cr\kern0.1em--\hfil\cr}}}%
%BeginExpansion
\ooalign{\hfil$V$\hfil\cr\kern0.1em--\hfil\cr}%
%EndExpansion
_{C}=P_{D}%
%TCIMACRO{\TeXButton{V}{\ooalign{\hfil$V$\hfil\cr\kern0.1em--\hfil\cr}}}%
%BeginExpansion
\ooalign{\hfil$V$\hfil\cr\kern0.1em--\hfil\cr}%
%EndExpansion
_{D}
\end{equation}%
and 
\begin{equation}
w_{CD}=nRT_{c}\ln \left( \frac{%
%TCIMACRO{\TeXButton{V}{\ooalign{\hfil$V$\hfil\cr\kern0.1em--\hfil\cr}}}%
%BeginExpansion
\ooalign{\hfil$V$\hfil\cr\kern0.1em--\hfil\cr}%
%EndExpansion
_{C}}{%
%TCIMACRO{\TeXButton{V}{\ooalign{\hfil$V$\hfil\cr\kern0.1em--\hfil\cr}}}%
%BeginExpansion
\ooalign{\hfil$V$\hfil\cr\kern0.1em--\hfil\cr}%
%EndExpansion
_{D}}\right)
\end{equation}%
again%
\begin{equation}
\Delta E_{int}=0
\end{equation}%
so%
\begin{equation}
\left\vert Q_{CD}\right\vert =\left\vert -w_{CD}\right\vert
\end{equation}%
and we can identify this as $\left\vert Q_{c}\right\vert $

We again need to find a volume, $%
%TCIMACRO{\TeXButton{V}{\ooalign{\hfil$V$\hfil\cr\kern0.1em--\hfil\cr}}}%
%BeginExpansion
\ooalign{\hfil$V$\hfil\cr\kern0.1em--\hfil\cr}%
%EndExpansion
_{D}$ and we can use the same method as before, finding $%
%TCIMACRO{\TeXButton{V}{\ooalign{\hfil$V$\hfil\cr\kern0.1em--\hfil\cr}}}%
%BeginExpansion
\ooalign{\hfil$V$\hfil\cr\kern0.1em--\hfil\cr}%
%EndExpansion
_{B}.$ At position $D$ 
\begin{equation*}
P_{D}=\frac{nRT_{c}}{%
%TCIMACRO{\TeXButton{V}{\ooalign{\hfil$V$\hfil\cr\kern0.1em--\hfil\cr}}}%
%BeginExpansion
\ooalign{\hfil$V$\hfil\cr\kern0.1em--\hfil\cr}%
%EndExpansion
_{D}}
\end{equation*}%
and 
\begin{equation*}
P_{D}=\frac{K_{DA}}{%
%TCIMACRO{\TeXButton{V}{\ooalign{\hfil$V$\hfil\cr\kern0.1em--\hfil\cr}}}%
%BeginExpansion
\ooalign{\hfil$V$\hfil\cr\kern0.1em--\hfil\cr}%
%EndExpansion
_{D}^{\gamma }}
\end{equation*}%
but we don't know $K_{DA}$, $P_{D},$ or $%
%TCIMACRO{\TeXButton{V}{\ooalign{\hfil$V$\hfil\cr\kern0.1em--\hfil\cr}}}%
%BeginExpansion
\ooalign{\hfil$V$\hfil\cr\kern0.1em--\hfil\cr}%
%EndExpansion
_{D}$

But we also know that at $A$ 
\begin{eqnarray*}
P_{A} &=&\frac{nRT_{h}}{%
%TCIMACRO{\TeXButton{V}{\ooalign{\hfil$V$\hfil\cr\kern0.1em--\hfil\cr}}}%
%BeginExpansion
\ooalign{\hfil$V$\hfil\cr\kern0.1em--\hfil\cr}%
%EndExpansion
_{A}} \\
P_{A} &=&\frac{K_{DA}}{%
%TCIMACRO{\TeXButton{V}{\ooalign{\hfil$V$\hfil\cr\kern0.1em--\hfil\cr}}}%
%BeginExpansion
\ooalign{\hfil$V$\hfil\cr\kern0.1em--\hfil\cr}%
%EndExpansion
_{A}^{\gamma }}
\end{eqnarray*}%
and at point $A$ and we do know $%
%TCIMACRO{\TeXButton{V}{\ooalign{\hfil$V$\hfil\cr\kern0.1em--\hfil\cr}}}%
%BeginExpansion
\ooalign{\hfil$V$\hfil\cr\kern0.1em--\hfil\cr}%
%EndExpansion
_{A,}$ so we can find $K_{DA}$%
\begin{equation*}
\frac{nRT_{h}}{%
%TCIMACRO{\TeXButton{V}{\ooalign{\hfil$V$\hfil\cr\kern0.1em--\hfil\cr}}}%
%BeginExpansion
\ooalign{\hfil$V$\hfil\cr\kern0.1em--\hfil\cr}%
%EndExpansion
_{A}}=\frac{K_{DA}}{%
%TCIMACRO{\TeXButton{V}{\ooalign{\hfil$V$\hfil\cr\kern0.1em--\hfil\cr}}}%
%BeginExpansion
\ooalign{\hfil$V$\hfil\cr\kern0.1em--\hfil\cr}%
%EndExpansion
_{A}^{\gamma }}
\end{equation*}%
so%
\begin{equation*}
K_{DA}=nRT_{h}%
%TCIMACRO{\TeXButton{V}{\ooalign{\hfil$V$\hfil\cr\kern0.1em--\hfil\cr}}}%
%BeginExpansion
\ooalign{\hfil$V$\hfil\cr\kern0.1em--\hfil\cr}%
%EndExpansion
_{A}^{\gamma -1}
\end{equation*}%
or%
\begin{eqnarray*}
K_{DA} &=&\left( 1\unit{mol}\right) \left( 8.314\frac{\unit{J}}{\unit{mol}%
\unit{K}}\right) \left( 1023\unit{K}\right) \left( 0.001\unit{m}^{3}\right)
^{\frac{2}{3}} \\
&=&85.\,\allowbreak 052\unit{J}\unit{m}^{2}
\end{eqnarray*}%
Then moving back to point $D$ we have 
\begin{equation*}
\frac{nRT_{c}}{%
%TCIMACRO{\TeXButton{V}{\ooalign{\hfil$V$\hfil\cr\kern0.1em--\hfil\cr}}}%
%BeginExpansion
\ooalign{\hfil$V$\hfil\cr\kern0.1em--\hfil\cr}%
%EndExpansion
_{D}}=\frac{K_{DA}}{%
%TCIMACRO{\TeXButton{V}{\ooalign{\hfil$V$\hfil\cr\kern0.1em--\hfil\cr}}}%
%BeginExpansion
\ooalign{\hfil$V$\hfil\cr\kern0.1em--\hfil\cr}%
%EndExpansion
_{D}^{\gamma }}
\end{equation*}%
and 
\begin{equation*}
%TCIMACRO{\TeXButton{V}{\ooalign{\hfil$V$\hfil\cr\kern0.1em--\hfil\cr}}}%
%BeginExpansion
\ooalign{\hfil$V$\hfil\cr\kern0.1em--\hfil\cr}%
%EndExpansion
_{D}=\sqrt[\gamma -1]{\frac{K_{DA}}{nRT_{c}}}
\end{equation*}%
Numerically we get%
\begin{eqnarray*}
%TCIMACRO{\TeXButton{V}{\ooalign{\hfil$V$\hfil\cr\kern0.1em--\hfil\cr}}}%
%BeginExpansion
\ooalign{\hfil$V$\hfil\cr\kern0.1em--\hfil\cr}%
%EndExpansion
_{D} &=&\sqrt[\frac{2}{3}]{\frac{85.\,\allowbreak 052\unit{J}\unit{m}^{2}}{%
\left( 1\unit{mol}\right) \left( 8.314\frac{\unit{J}}{\unit{mol}\unit{K}}%
\right) \left( 293\unit{K}\right) }} \\
&=&6.\,\allowbreak 524\,1\times 10^{-3}\unit{m}^{3}
\end{eqnarray*}%
And we are back to the ideal gas law to find the pressure%
\begin{eqnarray*}
P_{D} &=&\frac{\left( 1\unit{mol}\right) \left( 8.314\frac{\unit{J}}{\unit{%
mol}\unit{K}}\right) \left( \allowbreak 293\unit{K}\right) }{6.\,\allowbreak
524\,1\times 10^{-3}\unit{m}^{3}} \\
&=&3.\,\allowbreak 733\,9\times 10^{5}
\end{eqnarray*}%
And so we can find the work as%
\begin{eqnarray*}
w_{CD} &=&\left( 1\unit{mol}\right) \left( 8.314\frac{\unit{J}}{\unit{mol}%
\unit{K}}\right) \left( \allowbreak 273\unit{K}\right) \ln \left( \frac{0.015%
\unit{m}^{3}}{6.\,\allowbreak 524\,1\times 10^{-3}\unit{m}^{3}}\right) \\
&=&1889.\,\allowbreak 7\unit{J}
\end{eqnarray*}%
And our summary is as follows

\begin{equation*}
\begin{tabular}{|c|c|c|c|c|c|c|}
\hline
\textbf{State} & $\mathbf{T}\left( \unit{K}\right) $ & $\mathbf{P}\left( 
\unit{Pa}\right) $ & $%
%TCIMACRO{\TeXButton{V}{\ooalign{\hfil$V$\hfil\cr\kern0.1em--\hfil\cr}}}%
%BeginExpansion
\ooalign{\hfil$V$\hfil\cr\kern0.1em--\hfil\cr}%
%EndExpansion
\left( \unit{m}^{3}\right) $ & $\mathbf{n}\left( \unit{mol}\right) $ & 
\textbf{--} & \textbf{--} \\ \hline
$\mathbf{C}$ & $293$ & $1.\,\allowbreak 624\times 10^{5}$ & $0.015\unit{m}%
^{3}$ & $1$ & \textbf{--} & \textbf{--} \\ \hline
$\mathbf{D}$ & $293$ & $3.\,\allowbreak 733\,9\times 10^{5}$ & $%
6.\,\allowbreak 524\,1\times 10^{-3}$ & $1$ & \textbf{--} & \textbf{--} \\ 
\hline
\textbf{Process} & $\mathbf{Q}\left( \unit{J}\right) $ & $\mathbf{w}%
_{int}\left( \unit{J}\right) $ & $\mathbf{\Delta E}_{int}\left( \unit{J}%
\right) $ & $\mathbf{w}_{eng}\left( \unit{J}\right) $ & $\mathbf{Q}_{h}$ & $%
\mathbf{Q}_{c}$ \\ \hline
$\mathbf{CD}$ & $-1889.\,\allowbreak 7$ & $1889.\,\allowbreak 7$ & $0$ & $0$
& $0$ & $1889.\,\allowbreak 7$ \\ \hline
\end{tabular}%
\end{equation*}

\section{Process $D\rightarrow A$ adiabatic compression}

The gas cylinder is again placed against a thermally nonconducting wall so
no heat is exchanged with the surroundings. The temperature of the gas
increases from $T_{c}$ to $T_{h}.$The work done on the gas is $w_{DA}$\FRAME{%
dtbpF}{5.2731in}{4.034in}{0pt}{}{}{Figure}{\special{language "Scientific
Word";type "GRAPHIC";maintain-aspect-ratio TRUE;display "USEDEF";valid_file
"T";width 5.2731in;height 4.034in;depth 0pt;original-width
5.3458in;original-height 4.0819in;cropleft "0";croptop "1";cropright
"1";cropbottom "0";tempfilename 'SUS73WCO.wmf';tempfile-properties "XPR";}}%
We know $Q=0$ and that $\Delta E_{int}=w_{int}.$ Again we can use%
\begin{equation*}
\Delta E_{int}=nC_{V}\Delta T
\end{equation*}%
\begin{eqnarray*}
\Delta E_{int} &=&\left( 1\unit{mol}\right) \frac{3}{2}\left( 8.314\frac{%
\unit{J}}{\unit{mol}\unit{K}}\right) \left( 1023\unit{K}-293\unit{K}\right)
\\
&=&9103.\,\allowbreak 8
\end{eqnarray*}%
so then 
\begin{equation*}
\begin{tabular}{|c|c|c|c|c|c|c|}
\hline
\textbf{State} & $\mathbf{T}\left( \unit{K}\right) $ & $\mathbf{P}\left( 
\unit{Pa}\right) $ & $%
%TCIMACRO{\TeXButton{V}{\ooalign{\hfil$V$\hfil\cr\kern0.1em--\hfil\cr}}}%
%BeginExpansion
\ooalign{\hfil$V$\hfil\cr\kern0.1em--\hfil\cr}%
%EndExpansion
\left( \unit{m}^{3}\right) $ & $\mathbf{n}\left( \unit{mol}\right) $ & 
\textbf{--} & \textbf{--} \\ \hline
$\mathbf{D}$ & $293$ & $3.\,\allowbreak 733\,9\times 10^{5}$ & $%
6.\,\allowbreak 524\,1\times 10^{-3}$ & $1$ & \textbf{--} & \textbf{--} \\ 
\hline
$\mathbf{A}$ & $1023$ & $8.\,\allowbreak 505\,2\times 10^{6}$ & $0.001$ & $1$
& \textbf{--} & \textbf{--} \\ \hline
\textbf{Process} & $\mathbf{Q}\left( \unit{J}\right) $ & $\mathbf{W}%
_{int}\left( \unit{J}\right) $ & $\mathbf{\Delta E}_{int}\left( \unit{J}%
\right) $ & $\mathbf{W}_{eng}\left( \unit{J}\right) $ & $\mathbf{Q}_{h}$ & $%
\mathbf{Q}_{c}$ \\ \hline
$\mathbf{DA}$ & $0$ & $9103.\,\allowbreak 8$ & $9103.\,\allowbreak 8$ & $%
-9103.\,\allowbreak 8$ & $0$ & $0$ \\ \hline
\end{tabular}%
\end{equation*}%
All together we have%
\begin{equation*}
\begin{tabular}{|c|c|c|c|c|}
\hline
\textbf{State} & $\mathbf{T}\left( \unit{K}\right) $ & $\mathbf{P}\left( 
\unit{Pa}\right) $ & $%
%TCIMACRO{\TeXButton{V}{\ooalign{\hfil$V$\hfil\cr\kern0.1em--\hfil\cr}}}%
%BeginExpansion
\ooalign{\hfil$V$\hfil\cr\kern0.1em--\hfil\cr}%
%EndExpansion
\left( \unit{m}^{3}\right) $ & $\mathbf{n}\left( \unit{mol}\right) $ \\ 
\hline
$\mathbf{A}$ & $1023$ & $8.\,\allowbreak 505\,2\times 10^{6}$ & $0.001$ & $1$
\\ \hline
$\mathbf{B}$ & $1023$ & $3.\,\allowbreak 699\,2\times 10^{6}$ & $%
2.\,\allowbreak 299\,2\times 10^{-3}$ & $1$ \\ \hline
$\mathbf{C}$ & $293$ & $1.\,\allowbreak 624\times 10^{5}$ & $0.015$ & $1$ \\ 
\hline
$\mathbf{D}$ & $293$ & $3.\,\allowbreak 733\,9\times 10^{5}$ & $%
6.\,\allowbreak 524\,1\times 10^{-3}$ & $1$ \\ \hline
\end{tabular}%
\end{equation*}%
And our internal energy, heat, and work are%
\begin{equation*}
\begin{tabular}{|c|c|c|c|c|c|c|}
\hline
\textbf{Process} & $\mathbf{Q}\left( \unit{J}\right) $ & $\mathbf{W}%
_{int}\left( \unit{J}\right) $ & $\mathbf{\Delta E}_{int}\left( \unit{J}%
\right) $ & $\mathbf{W}_{eng}\left( \unit{J}\right) $ & $\mathbf{Q}_{h}$ & $%
\mathbf{Q}_{c}$ \\ \hline
$\mathbf{AB}$ & $7081.\,\allowbreak 1$ & $-7081.\,\allowbreak 13$ & $0$ & $%
7081.\,\allowbreak 1$ & $7081.\,\allowbreak 1$ & $0$ \\ \hline
$\mathbf{BC}$ & $0$ & $-9103.\,\allowbreak 8$ & $-9103.\,\allowbreak 8$ & $%
9103.\,\allowbreak 8$ & $0$ & $0$ \\ \hline
$\mathbf{CD}$ & $-1889.\,\allowbreak 7$ & $1889.\,\allowbreak 7$ & $0$ & $%
-1889.\,\allowbreak 7$ & $0$ & $1889.\,\allowbreak 7$ \\ \hline
$\mathbf{DA}$ & $0$ & $9103.\,\allowbreak 8$ & $9103.\,\allowbreak 8$ & $%
-9103.\,\allowbreak 8$ & $0$ & $0$ \\ \hline
Total & $5191.\,\allowbreak 4$ & $-5191.\,\allowbreak 4$ & $0$ & $%
5192.\,\allowbreak 2$ & $7081.\,\allowbreak 1$ & $1889.\,\allowbreak 7$ \\ 
\hline
\end{tabular}%
\end{equation*}

We can see that process $A\rightarrow B$ provided all of $Q_{h}$ and process 
$C\rightarrow D$ provided all of $Q_{C}.$Let's label this in our PV\
diagram. \FRAME{dhF}{3.2292in}{2.2468in}{0pt}{}{}{Figure}{\special{language
"Scientific Word";type "GRAPHIC";maintain-aspect-ratio TRUE;display
"USEDEF";valid_file "T";width 3.2292in;height 2.2468in;depth
0pt;original-width 3.1842in;original-height 2.207in;cropleft "0";croptop
"1";cropright "1";cropbottom "0";tempfilename
'SUS73WCP.wmf';tempfile-properties "XPR";}}

\section{Carnot Efficiency}

We can also find the efficiency of the Carnot cycle. The net work is $%
w_{ABCDA}$ and the efficiency is 
\begin{equation*}
\eta =\frac{w_{eng}}{\left\vert Q_{h}\right\vert }=1-\frac{\left\vert
Q_{c}\right\vert }{\left\vert Q_{h}\right\vert }
\end{equation*}

No heat is transferred during the adiabatic processes, so the only heat
transfer is during the isothermal processes, we have%
\begin{eqnarray*}
\eta  &=&1-\frac{\left\vert nRT_{c}\ln \left( \frac{%
%TCIMACRO{\TeXButton{V}{\ooalign{\hfil$V$\hfil\cr\kern0.1em--\hfil\cr}}}%
%BeginExpansion
\ooalign{\hfil$V$\hfil\cr\kern0.1em--\hfil\cr}%
%EndExpansion
_{C}}{%
%TCIMACRO{\TeXButton{V}{\ooalign{\hfil$V$\hfil\cr\kern0.1em--\hfil\cr}}}%
%BeginExpansion
\ooalign{\hfil$V$\hfil\cr\kern0.1em--\hfil\cr}%
%EndExpansion
_{D}}\right) \right\vert }{\left\vert nRT_{h}\ln \left( \frac{%
%TCIMACRO{\TeXButton{V}{\ooalign{\hfil$V$\hfil\cr\kern0.1em--\hfil\cr}}}%
%BeginExpansion
\ooalign{\hfil$V$\hfil\cr\kern0.1em--\hfil\cr}%
%EndExpansion
_{A}}{%
%TCIMACRO{\TeXButton{V}{\ooalign{\hfil$V$\hfil\cr\kern0.1em--\hfil\cr}}}%
%BeginExpansion
\ooalign{\hfil$V$\hfil\cr\kern0.1em--\hfil\cr}%
%EndExpansion
_{B}}\right) \right\vert } \\
&=&1-\frac{\left\vert T_{c}\ln \left( \frac{%
%TCIMACRO{\TeXButton{V}{\ooalign{\hfil$V$\hfil\cr\kern0.1em--\hfil\cr}}}%
%BeginExpansion
\ooalign{\hfil$V$\hfil\cr\kern0.1em--\hfil\cr}%
%EndExpansion
_{C}}{%
%TCIMACRO{\TeXButton{V}{\ooalign{\hfil$V$\hfil\cr\kern0.1em--\hfil\cr}}}%
%BeginExpansion
\ooalign{\hfil$V$\hfil\cr\kern0.1em--\hfil\cr}%
%EndExpansion
_{D}}\right) \right\vert }{\left\vert T_{h}\ln \left( \frac{%
%TCIMACRO{\TeXButton{V}{\ooalign{\hfil$V$\hfil\cr\kern0.1em--\hfil\cr}}}%
%BeginExpansion
\ooalign{\hfil$V$\hfil\cr\kern0.1em--\hfil\cr}%
%EndExpansion
_{A}}{%
%TCIMACRO{\TeXButton{V}{\ooalign{\hfil$V$\hfil\cr\kern0.1em--\hfil\cr}}}%
%BeginExpansion
\ooalign{\hfil$V$\hfil\cr\kern0.1em--\hfil\cr}%
%EndExpansion
_{B}}\right) \right\vert }
\end{eqnarray*}%
Using the adiabatic temperature-volume relationships we can find expressions
for the hot and cold temperatures%
\begin{equation*}
T_{h}%
%TCIMACRO{\TeXButton{V}{\ooalign{\hfil$V$\hfil\cr\kern0.1em--\hfil\cr}}}%
%BeginExpansion
\ooalign{\hfil$V$\hfil\cr\kern0.1em--\hfil\cr}%
%EndExpansion
_{A}^{\gamma -1}=T_{c}%
%TCIMACRO{\TeXButton{V}{\ooalign{\hfil$V$\hfil\cr\kern0.1em--\hfil\cr}}}%
%BeginExpansion
\ooalign{\hfil$V$\hfil\cr\kern0.1em--\hfil\cr}%
%EndExpansion
_{D}^{\gamma -1}
\end{equation*}%
and%
\begin{equation*}
T_{h}%
%TCIMACRO{\TeXButton{V}{\ooalign{\hfil$V$\hfil\cr\kern0.1em--\hfil\cr}}}%
%BeginExpansion
\ooalign{\hfil$V$\hfil\cr\kern0.1em--\hfil\cr}%
%EndExpansion
_{B}^{\gamma -1}=T_{c}%
%TCIMACRO{\TeXButton{V}{\ooalign{\hfil$V$\hfil\cr\kern0.1em--\hfil\cr}}}%
%BeginExpansion
\ooalign{\hfil$V$\hfil\cr\kern0.1em--\hfil\cr}%
%EndExpansion
_{C}^{\gamma -1}
\end{equation*}%
We could divide these two equations to get 
\begin{equation*}
\frac{T_{h}%
%TCIMACRO{\TeXButton{V}{\ooalign{\hfil$V$\hfil\cr\kern0.1em--\hfil\cr}}}%
%BeginExpansion
\ooalign{\hfil$V$\hfil\cr\kern0.1em--\hfil\cr}%
%EndExpansion
_{A}^{\gamma -1}}{T_{h}%
%TCIMACRO{\TeXButton{V}{\ooalign{\hfil$V$\hfil\cr\kern0.1em--\hfil\cr}}}%
%BeginExpansion
\ooalign{\hfil$V$\hfil\cr\kern0.1em--\hfil\cr}%
%EndExpansion
_{B}^{\gamma -1}}=\frac{T_{c}%
%TCIMACRO{\TeXButton{V}{\ooalign{\hfil$V$\hfil\cr\kern0.1em--\hfil\cr}}}%
%BeginExpansion
\ooalign{\hfil$V$\hfil\cr\kern0.1em--\hfil\cr}%
%EndExpansion
_{D}^{\gamma -1}}{T_{c}%
%TCIMACRO{\TeXButton{V}{\ooalign{\hfil$V$\hfil\cr\kern0.1em--\hfil\cr}}}%
%BeginExpansion
\ooalign{\hfil$V$\hfil\cr\kern0.1em--\hfil\cr}%
%EndExpansion
_{C}^{\gamma -1}}
\end{equation*}%
or%
\begin{equation*}
\frac{%
%TCIMACRO{\TeXButton{V}{\ooalign{\hfil$V$\hfil\cr\kern0.1em--\hfil\cr}}}%
%BeginExpansion
\ooalign{\hfil$V$\hfil\cr\kern0.1em--\hfil\cr}%
%EndExpansion
_{A}^{\gamma -1}}{%
%TCIMACRO{\TeXButton{V}{\ooalign{\hfil$V$\hfil\cr\kern0.1em--\hfil\cr}}}%
%BeginExpansion
\ooalign{\hfil$V$\hfil\cr\kern0.1em--\hfil\cr}%
%EndExpansion
_{B}^{\gamma -1}}=\frac{%
%TCIMACRO{\TeXButton{V}{\ooalign{\hfil$V$\hfil\cr\kern0.1em--\hfil\cr}}}%
%BeginExpansion
\ooalign{\hfil$V$\hfil\cr\kern0.1em--\hfil\cr}%
%EndExpansion
_{D}^{\gamma -1}}{%
%TCIMACRO{\TeXButton{V}{\ooalign{\hfil$V$\hfil\cr\kern0.1em--\hfil\cr}}}%
%BeginExpansion
\ooalign{\hfil$V$\hfil\cr\kern0.1em--\hfil\cr}%
%EndExpansion
_{C}^{\gamma -1}}
\end{equation*}%
which gives us%
\begin{equation*}
\sqrt[\gamma -1]{\frac{%
%TCIMACRO{\TeXButton{V}{\ooalign{\hfil$V$\hfil\cr\kern0.1em--\hfil\cr}}}%
%BeginExpansion
\ooalign{\hfil$V$\hfil\cr\kern0.1em--\hfil\cr}%
%EndExpansion
_{A}^{\gamma -1}}{%
%TCIMACRO{\TeXButton{V}{\ooalign{\hfil$V$\hfil\cr\kern0.1em--\hfil\cr}}}%
%BeginExpansion
\ooalign{\hfil$V$\hfil\cr\kern0.1em--\hfil\cr}%
%EndExpansion
_{B}^{\gamma -1}}}=\sqrt[\gamma -1]{\frac{%
%TCIMACRO{\TeXButton{V}{\ooalign{\hfil$V$\hfil\cr\kern0.1em--\hfil\cr}}}%
%BeginExpansion
\ooalign{\hfil$V$\hfil\cr\kern0.1em--\hfil\cr}%
%EndExpansion
_{D}^{\gamma -1}}{%
%TCIMACRO{\TeXButton{V}{\ooalign{\hfil$V$\hfil\cr\kern0.1em--\hfil\cr}}}%
%BeginExpansion
\ooalign{\hfil$V$\hfil\cr\kern0.1em--\hfil\cr}%
%EndExpansion
_{C}^{\gamma -1}}}
\end{equation*}%
But if we take both sides to the $\gamma -1$ power we have just%
\begin{equation*}
\frac{%
%TCIMACRO{\TeXButton{V}{\ooalign{\hfil$V$\hfil\cr\kern0.1em--\hfil\cr}}}%
%BeginExpansion
\ooalign{\hfil$V$\hfil\cr\kern0.1em--\hfil\cr}%
%EndExpansion
_{A}}{%
%TCIMACRO{\TeXButton{V}{\ooalign{\hfil$V$\hfil\cr\kern0.1em--\hfil\cr}}}%
%BeginExpansion
\ooalign{\hfil$V$\hfil\cr\kern0.1em--\hfil\cr}%
%EndExpansion
_{B}}=\frac{%
%TCIMACRO{\TeXButton{V}{\ooalign{\hfil$V$\hfil\cr\kern0.1em--\hfil\cr}}}%
%BeginExpansion
\ooalign{\hfil$V$\hfil\cr\kern0.1em--\hfil\cr}%
%EndExpansion
_{D}}{%
%TCIMACRO{\TeXButton{V}{\ooalign{\hfil$V$\hfil\cr\kern0.1em--\hfil\cr}}}%
%BeginExpansion
\ooalign{\hfil$V$\hfil\cr\kern0.1em--\hfil\cr}%
%EndExpansion
_{C}}
\end{equation*}%
Using this we find that our efficiency equation is 
\begin{eqnarray*}
\eta  &=&1-\frac{\left\vert T_{c}\ln \left( \frac{%
%TCIMACRO{\TeXButton{V}{\ooalign{\hfil$V$\hfil\cr\kern0.1em--\hfil\cr}}}%
%BeginExpansion
\ooalign{\hfil$V$\hfil\cr\kern0.1em--\hfil\cr}%
%EndExpansion
_{C}}{%
%TCIMACRO{\TeXButton{V}{\ooalign{\hfil$V$\hfil\cr\kern0.1em--\hfil\cr}}}%
%BeginExpansion
\ooalign{\hfil$V$\hfil\cr\kern0.1em--\hfil\cr}%
%EndExpansion
_{D}}\right) \right\vert }{\left\vert T_{h}\ln \left( \frac{%
%TCIMACRO{\TeXButton{V}{\ooalign{\hfil$V$\hfil\cr\kern0.1em--\hfil\cr}}}%
%BeginExpansion
\ooalign{\hfil$V$\hfil\cr\kern0.1em--\hfil\cr}%
%EndExpansion
_{D}}{%
%TCIMACRO{\TeXButton{V}{\ooalign{\hfil$V$\hfil\cr\kern0.1em--\hfil\cr}}}%
%BeginExpansion
\ooalign{\hfil$V$\hfil\cr\kern0.1em--\hfil\cr}%
%EndExpansion
_{C}}\right) \right\vert } \\
&=&1-\frac{\left\vert T_{c}\ln \left( \frac{%
%TCIMACRO{\TeXButton{V}{\ooalign{\hfil$V$\hfil\cr\kern0.1em--\hfil\cr}}}%
%BeginExpansion
\ooalign{\hfil$V$\hfil\cr\kern0.1em--\hfil\cr}%
%EndExpansion
_{C}}{%
%TCIMACRO{\TeXButton{V}{\ooalign{\hfil$V$\hfil\cr\kern0.1em--\hfil\cr}}}%
%BeginExpansion
\ooalign{\hfil$V$\hfil\cr\kern0.1em--\hfil\cr}%
%EndExpansion
_{D}}\right) \right\vert }{\left\vert -T_{h}\ln \left( \frac{%
%TCIMACRO{\TeXButton{V}{\ooalign{\hfil$V$\hfil\cr\kern0.1em--\hfil\cr}}}%
%BeginExpansion
\ooalign{\hfil$V$\hfil\cr\kern0.1em--\hfil\cr}%
%EndExpansion
_{C}}{%
%TCIMACRO{\TeXButton{V}{\ooalign{\hfil$V$\hfil\cr\kern0.1em--\hfil\cr}}}%
%BeginExpansion
\ooalign{\hfil$V$\hfil\cr\kern0.1em--\hfil\cr}%
%EndExpansion
_{D}}\right) \right\vert }
\end{eqnarray*}%
or, removing the minus sign because of the absolute values%
\begin{equation*}
\eta =1-\frac{T_{c}}{T_{h}}
\end{equation*}%
where the temperatures, of course, are in Kelvin. This means that for this
engine%
\begin{equation*}
\frac{\left\vert Q_{c}\right\vert }{\left\vert Q_{h}\right\vert }=\frac{T_{c}%
}{T_{h}}
\end{equation*}%
and 
\begin{equation}
\eta _{C}=1-\frac{T_{c}}{T_{h}}
\end{equation}%
This is the general form for the efficiency of a Carnot cycle. This would be
the most efficient any heat engine could be that operates between the
temperatures $T_{c}$ and $T_{h}.$

This achieved our goal. We have a way to find the maximum efficiency for a
heat engine. And along the way we made several statements of the second law
of thermodynamics. 

\chapter{37 An example: The Otto Cycle}

We now know about heat engine. But so far we have been very vague about what
happens in the yellow engine circle in our energy transfer diagrams. Let's
look at a (mostly) real example of a heat engine. This model of a heat
engine is due to a researcher named Otto, so it is called the \emph{Otto
Cycle}.

%TCIMACRO{%
%\TeXButton{Fundamental Concepts}{\vspace{0.10in}\hspace{-0.37in}{\Large Fundamental Concepts\vspace{0.2in}}}}%
%BeginExpansion
\vspace{0.10in}\hspace{-0.37in}{\Large Fundamental Concepts\vspace{0.2in}}%
%EndExpansion

\begin{itemize}
\item To describe a heat engine cycle you need to find $P,$ $%
%TCIMACRO{\TeXButton{V}{\ooalign{\hfil$V$\hfil\cr\kern0.1em--\hfil\cr}}}%
%BeginExpansion
\ooalign{\hfil$V$\hfil\cr\kern0.1em--\hfil\cr}%
%EndExpansion
,$ $n,$ $T,$ $\Delta E_{int},$ $Q,$ $W,$ for each part of the cycle and $%
W_{eng},$ $Q_{h}$, $Q_{c}$ for the entire cycle.

\item The Otto Cycle
\end{itemize}

\section{Calculating the Otto Cycle}

A cyclic process is made from more than one thermodynamic process. The Otto
cycle technically has sic processes. To design an engine we need to describe
all of these processes and make sure they work together to produce a net
useful work.

\FRAME{dtbpF}{4.8888in}{4.0075in}{0pt}{}{}{Figure}{\special{language
"Scientific Word";type "GRAPHIC";maintain-aspect-ratio TRUE;display
"USEDEF";valid_file "T";width 4.8888in;height 4.0075in;depth
0pt;original-width 4.8352in;original-height 3.9583in;cropleft "0";croptop
"1";cropright "1";cropbottom "0";tempfilename
'TJXN1F0M.wmf';tempfile-properties "XPR";}}

The $P%
%TCIMACRO{\TeXButton{V}{\ooalign{\hfil$V$\hfil\cr\kern0.1em--\hfil\cr}}}%
%BeginExpansion
\ooalign{\hfil$V$\hfil\cr\kern0.1em--\hfil\cr}%
%EndExpansion
\ $diagram above shows the Otto cycle. This cycle is approximately what
happens in a gasoline engine, like those in cars or, more like what we find
in lawn mowers. It consists of four main processes and two pseudo processes.
All of these processes correspond to a piston moving in a cylinder (our
favorite thermodynamic system!). Here is a picture of what is happening in
the engine. \FRAME{dtbpF}{5.5729in}{1.3589in}{0pt}{}{}{Figure}{\special%
{language "Scientific Word";type "GRAPHIC";maintain-aspect-ratio
TRUE;display "USEDEF";valid_file "T";width 5.5729in;height 1.3589in;depth
0pt;original-width 7.2928in;original-height 1.7553in;cropleft "0";croptop
"1";cropright "1";cropbottom "0";tempfilename
'TJXN1F0N.wmf';tempfile-properties "XPR";}}

To understand the Otto cycle let's do the job of finding $\Delta E_{int},$ $%
Q,$ and $w_{int}$ for each part of the cycle and then describe the state of
the gas at the beginning and end of each cycle by finding $P,$ $%
%TCIMACRO{\TeXButton{V}{\ooalign{\hfil$V$\hfil\cr\kern0.1em--\hfil\cr}}}%
%BeginExpansion
\ooalign{\hfil$V$\hfil\cr\kern0.1em--\hfil\cr}%
%EndExpansion
,$ $n,$ $T,$ and $E_{int}$. For our engine, suppose the compression ratio is 
\begin{equation*}
\frac{%
%TCIMACRO{\TeXButton{V}{\ooalign{\hfil$V$\hfil\cr\kern0.1em--\hfil\cr}}}%
%BeginExpansion
\ooalign{\hfil$V$\hfil\cr\kern0.1em--\hfil\cr}%
%EndExpansion
_{A}}{%
%TCIMACRO{\TeXButton{V}{\ooalign{\hfil$V$\hfil\cr\kern0.1em--\hfil\cr}}}%
%BeginExpansion
\ooalign{\hfil$V$\hfil\cr\kern0.1em--\hfil\cr}%
%EndExpansion
_{B}}=8.00
\end{equation*}%
Let's take 
\begin{eqnarray*}
%TCIMACRO{\TeXButton{V}{\ooalign{\hfil$V$\hfil\cr\kern0.1em--\hfil\cr}}}%
%BeginExpansion
\ooalign{\hfil$V$\hfil\cr\kern0.1em--\hfil\cr}%
%EndExpansion
_{A} &=&500\unit{cm}^{3}=\allowbreak 0.000\,5\unit{m}^{3} \\
P_{A} &=&100\unit{kPa} \\
T_{A} &=&293\unit{K}
\end{eqnarray*}%
Then we can find $%
%TCIMACRO{\TeXButton{V}{\ooalign{\hfil$V$\hfil\cr\kern0.1em--\hfil\cr}}}%
%BeginExpansion
\ooalign{\hfil$V$\hfil\cr\kern0.1em--\hfil\cr}%
%EndExpansion
_{B}$ 
\begin{equation*}
%TCIMACRO{\TeXButton{V}{\ooalign{\hfil$V$\hfil\cr\kern0.1em--\hfil\cr}}}%
%BeginExpansion
\ooalign{\hfil$V$\hfil\cr\kern0.1em--\hfil\cr}%
%EndExpansion
_{B}=\frac{%
%TCIMACRO{\TeXButton{V}{\ooalign{\hfil$V$\hfil\cr\kern0.1em--\hfil\cr}}}%
%BeginExpansion
\ooalign{\hfil$V$\hfil\cr\kern0.1em--\hfil\cr}%
%EndExpansion
_{A}}{8}=\allowbreak 62.\,\allowbreak 5\unit{cm}^{3}=\allowbreak
6.\,\allowbreak 25\times 10^{-5}\unit{m}^{3}
\end{equation*}%
and we will need to measure the temperature at $C.$ Suppose we measure 
\begin{equation*}
T_{C}=1023\unit{K}
\end{equation*}%
and the intake temperature to be 
\begin{equation*}
T_{I}=293\unit{K}
\end{equation*}
Further suppose our engine uses a diatomic ideal gas with 
\begin{eqnarray*}
C_{%
%TCIMACRO{\TeXButton{V}{\ooalign{\hfil$V$\hfil\cr\kern0.1em--\hfil\cr}}}%
%BeginExpansion
\ooalign{\hfil$V$\hfil\cr\kern0.1em--\hfil\cr}%
%EndExpansion
} &=&\frac{5}{2}R \\
R &=&8.314\frac{\unit{J}}{\unit{mol}\unit{K}}
\end{eqnarray*}%
so 
\begin{equation*}
\gamma =1.4
\end{equation*}%
Let's follow the process through to see what the state variables are at each
point on the PV\ diagram, and what $Q_{h}$ and $Q_{c}$ are for each process,
and what the work, $W,$ is and what the change in internal energy will be, $%
\Delta E_{int}.$ Finally, we will find the efficiency of the engine.\FRAME{%
dtbpF}{5.0246in}{2.93in}{0pt}{}{}{Figure}{\special{language "Scientific
Word";type "GRAPHIC";maintain-aspect-ratio TRUE;display "USEDEF";valid_file
"T";width 5.0246in;height 2.93in;depth 0pt;original-width
2.9836in;original-height 1.7279in;cropleft "0";croptop "1";cropright
"1";cropbottom "0";tempfilename 'OttoProcessAB_2.wmf';tempfile-properties
"XNPR";}}

\textbf{Process }$A\rightarrow B$\textbf{\ is an adiabatic compression.} The
piston moves upward, compressing the fuel-air mixture. This happens rapidly,
so there is no time for energy to leave by heat, $Q=0.$ That makes it close
enough to being adiabatic that we can say this process is approximately
adiabatic. In this process $\Delta E_{int}=w_{int}.$

We can use one of our adiabatic equations 
\begin{equation*}
P_{i}%
%TCIMACRO{\TeXButton{V}{\ooalign{\hfil$V$\hfil\cr\kern0.1em--\hfil\cr}}}%
%BeginExpansion
\ooalign{\hfil$V$\hfil\cr\kern0.1em--\hfil\cr}%
%EndExpansion
_{i}^{\gamma }=P_{f}%
%TCIMACRO{\TeXButton{V}{\ooalign{\hfil$V$\hfil\cr\kern0.1em--\hfil\cr}}}%
%BeginExpansion
\ooalign{\hfil$V$\hfil\cr\kern0.1em--\hfil\cr}%
%EndExpansion
_{f}^{\gamma }
\end{equation*}%
to find the final pressure for this process at $B$%
\begin{equation*}
P_{B}=P_{A}\frac{%
%TCIMACRO{\TeXButton{V}{\ooalign{\hfil$V$\hfil\cr\kern0.1em--\hfil\cr}}}%
%BeginExpansion
\ooalign{\hfil$V$\hfil\cr\kern0.1em--\hfil\cr}%
%EndExpansion
_{A}^{\gamma }}{%
%TCIMACRO{\TeXButton{V}{\ooalign{\hfil$V$\hfil\cr\kern0.1em--\hfil\cr}}}%
%BeginExpansion
\ooalign{\hfil$V$\hfil\cr\kern0.1em--\hfil\cr}%
%EndExpansion
_{B}^{\gamma }}
\end{equation*}%
But of course this works for any point on the line from $A$ to $B$ so 
\begin{equation*}
P=P_{A}\frac{%
%TCIMACRO{\TeXButton{V}{\ooalign{\hfil$V$\hfil\cr\kern0.1em--\hfil\cr}}}%
%BeginExpansion
\ooalign{\hfil$V$\hfil\cr\kern0.1em--\hfil\cr}%
%EndExpansion
_{A}^{\gamma }}{%
%TCIMACRO{\TeXButton{V}{\ooalign{\hfil$V$\hfil\cr\kern0.1em--\hfil\cr}}}%
%BeginExpansion
\ooalign{\hfil$V$\hfil\cr\kern0.1em--\hfil\cr}%
%EndExpansion
^{\gamma }}
\end{equation*}

We can use this form for $P$ to find the work for process $A\rightarrow B$ 
\begin{eqnarray*}
w_{AB} &=&-\int_{%
%TCIMACRO{\TeXButton{V}{\ooalign{\hfil$V$\hfil\cr\kern0.1em--\hfil\cr}}}%
%BeginExpansion
\ooalign{\hfil$V$\hfil\cr\kern0.1em--\hfil\cr}%
%EndExpansion
_{A}}^{%
%TCIMACRO{\TeXButton{V}{\ooalign{\hfil$V$\hfil\cr\kern0.1em--\hfil\cr}}}%
%BeginExpansion
\ooalign{\hfil$V$\hfil\cr\kern0.1em--\hfil\cr}%
%EndExpansion
_{B}}Pd%
%TCIMACRO{\TeXButton{V}{\ooalign{\hfil$V$\hfil\cr\kern0.1em--\hfil\cr}} }%
%BeginExpansion
\ooalign{\hfil$V$\hfil\cr\kern0.1em--\hfil\cr}
%EndExpansion
\\
&=&-\int_{%
%TCIMACRO{\TeXButton{V}{\ooalign{\hfil$V$\hfil\cr\kern0.1em--\hfil\cr}}}%
%BeginExpansion
\ooalign{\hfil$V$\hfil\cr\kern0.1em--\hfil\cr}%
%EndExpansion
_{A}}^{%
%TCIMACRO{\TeXButton{V}{\ooalign{\hfil$V$\hfil\cr\kern0.1em--\hfil\cr}}}%
%BeginExpansion
\ooalign{\hfil$V$\hfil\cr\kern0.1em--\hfil\cr}%
%EndExpansion
_{B}}P_{A}\frac{%
%TCIMACRO{\TeXButton{V}{\ooalign{\hfil$V$\hfil\cr\kern0.1em--\hfil\cr}}}%
%BeginExpansion
\ooalign{\hfil$V$\hfil\cr\kern0.1em--\hfil\cr}%
%EndExpansion
_{A}^{\gamma }}{%
%TCIMACRO{\TeXButton{V}{\ooalign{\hfil$V$\hfil\cr\kern0.1em--\hfil\cr}}}%
%BeginExpansion
\ooalign{\hfil$V$\hfil\cr\kern0.1em--\hfil\cr}%
%EndExpansion
^{\gamma }}d%
%TCIMACRO{\TeXButton{V}{\ooalign{\hfil$V$\hfil\cr\kern0.1em--\hfil\cr}} }%
%BeginExpansion
\ooalign{\hfil$V$\hfil\cr\kern0.1em--\hfil\cr}
%EndExpansion
\\
&=&-P_{A}%
%TCIMACRO{\TeXButton{V}{\ooalign{\hfil$V$\hfil\cr\kern0.1em--\hfil\cr}}}%
%BeginExpansion
\ooalign{\hfil$V$\hfil\cr\kern0.1em--\hfil\cr}%
%EndExpansion
_{A}^{\gamma }\int_{%
%TCIMACRO{\TeXButton{V}{\ooalign{\hfil$V$\hfil\cr\kern0.1em--\hfil\cr}}}%
%BeginExpansion
\ooalign{\hfil$V$\hfil\cr\kern0.1em--\hfil\cr}%
%EndExpansion
_{A}}^{%
%TCIMACRO{\TeXButton{V}{\ooalign{\hfil$V$\hfil\cr\kern0.1em--\hfil\cr}}}%
%BeginExpansion
\ooalign{\hfil$V$\hfil\cr\kern0.1em--\hfil\cr}%
%EndExpansion
_{B}}\frac{1}{%
%TCIMACRO{\TeXButton{V}{\ooalign{\hfil$V$\hfil\cr\kern0.1em--\hfil\cr}}}%
%BeginExpansion
\ooalign{\hfil$V$\hfil\cr\kern0.1em--\hfil\cr}%
%EndExpansion
^{\gamma }}d%
%TCIMACRO{\TeXButton{V}{\ooalign{\hfil$V$\hfil\cr\kern0.1em--\hfil\cr}} }%
%BeginExpansion
\ooalign{\hfil$V$\hfil\cr\kern0.1em--\hfil\cr}
%EndExpansion
\\
&=&-P_{A}^{\gamma }%
%TCIMACRO{\TeXButton{V}{\ooalign{\hfil$V$\hfil\cr\kern0.1em--\hfil\cr}}}%
%BeginExpansion
\ooalign{\hfil$V$\hfil\cr\kern0.1em--\hfil\cr}%
%EndExpansion
_{A}^{\gamma }\left( \frac{%
%TCIMACRO{\TeXButton{V}{\ooalign{\hfil$V$\hfil\cr\kern0.1em--\hfil\cr}}}%
%BeginExpansion
\ooalign{\hfil$V$\hfil\cr\kern0.1em--\hfil\cr}%
%EndExpansion
^{1-\gamma }}{1-\gamma }\right\vert _{%
%TCIMACRO{\TeXButton{V}{\ooalign{\hfil$V$\hfil\cr\kern0.1em--\hfil\cr}}}%
%BeginExpansion
\ooalign{\hfil$V$\hfil\cr\kern0.1em--\hfil\cr}%
%EndExpansion
_{A}}^{%
%TCIMACRO{\TeXButton{V}{\ooalign{\hfil$V$\hfil\cr\kern0.1em--\hfil\cr}}}%
%BeginExpansion
\ooalign{\hfil$V$\hfil\cr\kern0.1em--\hfil\cr}%
%EndExpansion
_{B}} \\
&=&-\frac{P_{A}%
%TCIMACRO{\TeXButton{V}{\ooalign{\hfil$V$\hfil\cr\kern0.1em--\hfil\cr}}}%
%BeginExpansion
\ooalign{\hfil$V$\hfil\cr\kern0.1em--\hfil\cr}%
%EndExpansion
_{A}^{\gamma }}{1-\gamma }\left( 
%TCIMACRO{\TeXButton{V}{\ooalign{\hfil$V$\hfil\cr\kern0.1em--\hfil\cr}}}%
%BeginExpansion
\ooalign{\hfil$V$\hfil\cr\kern0.1em--\hfil\cr}%
%EndExpansion
_{B}^{1-\gamma }-%
%TCIMACRO{\TeXButton{V}{\ooalign{\hfil$V$\hfil\cr\kern0.1em--\hfil\cr}}}%
%BeginExpansion
\ooalign{\hfil$V$\hfil\cr\kern0.1em--\hfil\cr}%
%EndExpansion
_{A}^{1-\gamma }\right)
\end{eqnarray*}%
We can rearrange this into a more convenient formula%
\begin{eqnarray*}
w_{AB} &=&\frac{P_{A}%
%TCIMACRO{\TeXButton{V}{\ooalign{\hfil$V$\hfil\cr\kern0.1em--\hfil\cr}}}%
%BeginExpansion
\ooalign{\hfil$V$\hfil\cr\kern0.1em--\hfil\cr}%
%EndExpansion
_{A}^{\gamma }}{\gamma -1}\left( \frac{1}{%
%TCIMACRO{\TeXButton{V}{\ooalign{\hfil$V$\hfil\cr\kern0.1em--\hfil\cr}}}%
%BeginExpansion
\ooalign{\hfil$V$\hfil\cr\kern0.1em--\hfil\cr}%
%EndExpansion
_{B}^{\gamma -1}}-\frac{1}{%
%TCIMACRO{\TeXButton{V}{\ooalign{\hfil$V$\hfil\cr\kern0.1em--\hfil\cr}}}%
%BeginExpansion
\ooalign{\hfil$V$\hfil\cr\kern0.1em--\hfil\cr}%
%EndExpansion
_{A}^{\gamma -1}}\right) \\
&=&\frac{1}{\gamma -1}\left( \frac{P_{A}%
%TCIMACRO{\TeXButton{V}{\ooalign{\hfil$V$\hfil\cr\kern0.1em--\hfil\cr}}}%
%BeginExpansion
\ooalign{\hfil$V$\hfil\cr\kern0.1em--\hfil\cr}%
%EndExpansion
_{A}^{\gamma }}{%
%TCIMACRO{\TeXButton{V}{\ooalign{\hfil$V$\hfil\cr\kern0.1em--\hfil\cr}}}%
%BeginExpansion
\ooalign{\hfil$V$\hfil\cr\kern0.1em--\hfil\cr}%
%EndExpansion
_{B}^{\gamma -1}}-\frac{P_{A}%
%TCIMACRO{\TeXButton{V}{\ooalign{\hfil$V$\hfil\cr\kern0.1em--\hfil\cr}}}%
%BeginExpansion
\ooalign{\hfil$V$\hfil\cr\kern0.1em--\hfil\cr}%
%EndExpansion
_{A}^{\gamma }}{%
%TCIMACRO{\TeXButton{V}{\ooalign{\hfil$V$\hfil\cr\kern0.1em--\hfil\cr}}}%
%BeginExpansion
\ooalign{\hfil$V$\hfil\cr\kern0.1em--\hfil\cr}%
%EndExpansion
_{A}^{\gamma -1}}\right) \\
&=&\frac{1}{\gamma -1}\left( \frac{P_{A}%
%TCIMACRO{\TeXButton{V}{\ooalign{\hfil$V$\hfil\cr\kern0.1em--\hfil\cr}}}%
%BeginExpansion
\ooalign{\hfil$V$\hfil\cr\kern0.1em--\hfil\cr}%
%EndExpansion
_{A}^{\gamma }}{%
%TCIMACRO{\TeXButton{V}{\ooalign{\hfil$V$\hfil\cr\kern0.1em--\hfil\cr}}}%
%BeginExpansion
\ooalign{\hfil$V$\hfil\cr\kern0.1em--\hfil\cr}%
%EndExpansion
_{B}^{\gamma }%
%TCIMACRO{\TeXButton{V}{\ooalign{\hfil$V$\hfil\cr\kern0.1em--\hfil\cr}}}%
%BeginExpansion
\ooalign{\hfil$V$\hfil\cr\kern0.1em--\hfil\cr}%
%EndExpansion
_{B}^{-1}}-\frac{P_{A}%
%TCIMACRO{\TeXButton{V}{\ooalign{\hfil$V$\hfil\cr\kern0.1em--\hfil\cr}}}%
%BeginExpansion
\ooalign{\hfil$V$\hfil\cr\kern0.1em--\hfil\cr}%
%EndExpansion
_{A}^{\gamma }}{%
%TCIMACRO{\TeXButton{V}{\ooalign{\hfil$V$\hfil\cr\kern0.1em--\hfil\cr}}}%
%BeginExpansion
\ooalign{\hfil$V$\hfil\cr\kern0.1em--\hfil\cr}%
%EndExpansion
_{A}^{\gamma -1}}\right)
\end{eqnarray*}%
And once again using 
\begin{equation*}
P_{B}%
%TCIMACRO{\TeXButton{V}{\ooalign{\hfil$V$\hfil\cr\kern0.1em--\hfil\cr}}}%
%BeginExpansion
\ooalign{\hfil$V$\hfil\cr\kern0.1em--\hfil\cr}%
%EndExpansion
_{B}^{\gamma }=P_{A}%
%TCIMACRO{\TeXButton{V}{\ooalign{\hfil$V$\hfil\cr\kern0.1em--\hfil\cr}}}%
%BeginExpansion
\ooalign{\hfil$V$\hfil\cr\kern0.1em--\hfil\cr}%
%EndExpansion
_{A}^{\gamma }
\end{equation*}%
or 
\begin{equation*}
P_{A}=P_{B}\frac{%
%TCIMACRO{\TeXButton{V}{\ooalign{\hfil$V$\hfil\cr\kern0.1em--\hfil\cr}}}%
%BeginExpansion
\ooalign{\hfil$V$\hfil\cr\kern0.1em--\hfil\cr}%
%EndExpansion
_{B}^{\gamma }}{%
%TCIMACRO{\TeXButton{V}{\ooalign{\hfil$V$\hfil\cr\kern0.1em--\hfil\cr}}}%
%BeginExpansion
\ooalign{\hfil$V$\hfil\cr\kern0.1em--\hfil\cr}%
%EndExpansion
_{A}^{\gamma }}
\end{equation*}%
We get 
\begin{equation*}
w_{AB}=\frac{1}{\gamma -1}\left( \frac{P_{B}\frac{%
%TCIMACRO{\TeXButton{V}{\ooalign{\hfil$V$\hfil\cr\kern0.1em--\hfil\cr}}}%
%BeginExpansion
\ooalign{\hfil$V$\hfil\cr\kern0.1em--\hfil\cr}%
%EndExpansion
_{B}^{\gamma }}{%
%TCIMACRO{\TeXButton{V}{\ooalign{\hfil$V$\hfil\cr\kern0.1em--\hfil\cr}}}%
%BeginExpansion
\ooalign{\hfil$V$\hfil\cr\kern0.1em--\hfil\cr}%
%EndExpansion
_{A}^{\gamma }}%
%TCIMACRO{\TeXButton{V}{\ooalign{\hfil$V$\hfil\cr\kern0.1em--\hfil\cr}}}%
%BeginExpansion
\ooalign{\hfil$V$\hfil\cr\kern0.1em--\hfil\cr}%
%EndExpansion
_{A}^{\gamma }}{%
%TCIMACRO{\TeXButton{V}{\ooalign{\hfil$V$\hfil\cr\kern0.1em--\hfil\cr}}}%
%BeginExpansion
\ooalign{\hfil$V$\hfil\cr\kern0.1em--\hfil\cr}%
%EndExpansion
_{B}^{\gamma }%
%TCIMACRO{\TeXButton{V}{\ooalign{\hfil$V$\hfil\cr\kern0.1em--\hfil\cr}}}%
%BeginExpansion
\ooalign{\hfil$V$\hfil\cr\kern0.1em--\hfil\cr}%
%EndExpansion
_{B}^{-1}}-\frac{P_{A}%
%TCIMACRO{\TeXButton{V}{\ooalign{\hfil$V$\hfil\cr\kern0.1em--\hfil\cr}}}%
%BeginExpansion
\ooalign{\hfil$V$\hfil\cr\kern0.1em--\hfil\cr}%
%EndExpansion
_{A}^{\gamma }}{%
%TCIMACRO{\TeXButton{V}{\ooalign{\hfil$V$\hfil\cr\kern0.1em--\hfil\cr}}}%
%BeginExpansion
\ooalign{\hfil$V$\hfil\cr\kern0.1em--\hfil\cr}%
%EndExpansion
_{A}^{\gamma -1}}\right)
\end{equation*}%
A lot of things cancel. We get%
\begin{eqnarray*}
w_{AB} &=&\frac{1}{\gamma -1}\left( \frac{P_{B}}{%
%TCIMACRO{\TeXButton{V}{\ooalign{\hfil$V$\hfil\cr\kern0.1em--\hfil\cr}}}%
%BeginExpansion
\ooalign{\hfil$V$\hfil\cr\kern0.1em--\hfil\cr}%
%EndExpansion
_{B}^{-1}}-\frac{P_{A}%
%TCIMACRO{\TeXButton{V}{\ooalign{\hfil$V$\hfil\cr\kern0.1em--\hfil\cr}}}%
%BeginExpansion
\ooalign{\hfil$V$\hfil\cr\kern0.1em--\hfil\cr}%
%EndExpansion
_{A}^{\gamma }}{%
%TCIMACRO{\TeXButton{V}{\ooalign{\hfil$V$\hfil\cr\kern0.1em--\hfil\cr}}}%
%BeginExpansion
\ooalign{\hfil$V$\hfil\cr\kern0.1em--\hfil\cr}%
%EndExpansion
_{A}^{\gamma -1}}\right) \\
&=&\frac{1}{\gamma -1}\left( P_{B}%
%TCIMACRO{\TeXButton{V}{\ooalign{\hfil$V$\hfil\cr\kern0.1em--\hfil\cr}}}%
%BeginExpansion
\ooalign{\hfil$V$\hfil\cr\kern0.1em--\hfil\cr}%
%EndExpansion
_{B}-P_{A}%
%TCIMACRO{\TeXButton{V}{\ooalign{\hfil$V$\hfil\cr\kern0.1em--\hfil\cr}}}%
%BeginExpansion
\ooalign{\hfil$V$\hfil\cr\kern0.1em--\hfil\cr}%
%EndExpansion
_{A}\right)
\end{eqnarray*}%
which we can generalize into yet another basic equation for adiabatic
processes!%
\begin{equation*}
w_{if}=\left( \frac{1}{\gamma -1}\right) \left( P_{f}%
%TCIMACRO{\TeXButton{V}{\ooalign{\hfil$V$\hfil\cr\kern0.1em--\hfil\cr}}}%
%BeginExpansion
\ooalign{\hfil$V$\hfil\cr\kern0.1em--\hfil\cr}%
%EndExpansion
_{f}-P_{i}%
%TCIMACRO{\TeXButton{V}{\ooalign{\hfil$V$\hfil\cr\kern0.1em--\hfil\cr}}}%
%BeginExpansion
\ooalign{\hfil$V$\hfil\cr\kern0.1em--\hfil\cr}%
%EndExpansion
_{i}\right)
\end{equation*}%
Then for process $A\rightarrow B$ we have%
\begin{equation*}
w_{AB}=\left( \frac{1}{\gamma -1}\right) \left( P_{B}%
%TCIMACRO{\TeXButton{V}{\ooalign{\hfil$V$\hfil\cr\kern0.1em--\hfil\cr}}}%
%BeginExpansion
\ooalign{\hfil$V$\hfil\cr\kern0.1em--\hfil\cr}%
%EndExpansion
_{B}-P_{A}%
%TCIMACRO{\TeXButton{V}{\ooalign{\hfil$V$\hfil\cr\kern0.1em--\hfil\cr}}}%
%BeginExpansion
\ooalign{\hfil$V$\hfil\cr\kern0.1em--\hfil\cr}%
%EndExpansion
_{A}\right)
\end{equation*}%
and we can use our equation for $P_{B}$ that we got before, $P_{B}=P_{A}%
\frac{%
%TCIMACRO{\TeXButton{V}{\ooalign{\hfil$V$\hfil\cr\kern0.1em--\hfil\cr}}}%
%BeginExpansion
\ooalign{\hfil$V$\hfil\cr\kern0.1em--\hfil\cr}%
%EndExpansion
_{A}^{\gamma }}{%
%TCIMACRO{\TeXButton{V}{\ooalign{\hfil$V$\hfil\cr\kern0.1em--\hfil\cr}}}%
%BeginExpansion
\ooalign{\hfil$V$\hfil\cr\kern0.1em--\hfil\cr}%
%EndExpansion
_{B}^{\gamma }}$ so that our $w_{AB}$ equation becomes 
\begin{equation*}
w_{AB}=\left( \frac{1}{\gamma -1}\right) \left( P_{A}\frac{%
%TCIMACRO{\TeXButton{V}{\ooalign{\hfil$V$\hfil\cr\kern0.1em--\hfil\cr}}}%
%BeginExpansion
\ooalign{\hfil$V$\hfil\cr\kern0.1em--\hfil\cr}%
%EndExpansion
_{A}^{\gamma }}{%
%TCIMACRO{\TeXButton{V}{\ooalign{\hfil$V$\hfil\cr\kern0.1em--\hfil\cr}}}%
%BeginExpansion
\ooalign{\hfil$V$\hfil\cr\kern0.1em--\hfil\cr}%
%EndExpansion
_{B}^{\gamma }}%
%TCIMACRO{\TeXButton{V}{\ooalign{\hfil$V$\hfil\cr\kern0.1em--\hfil\cr}}}%
%BeginExpansion
\ooalign{\hfil$V$\hfil\cr\kern0.1em--\hfil\cr}%
%EndExpansion
_{B}-P_{A}%
%TCIMACRO{\TeXButton{V}{\ooalign{\hfil$V$\hfil\cr\kern0.1em--\hfil\cr}}}%
%BeginExpansion
\ooalign{\hfil$V$\hfil\cr\kern0.1em--\hfil\cr}%
%EndExpansion
_{A}\right)
\end{equation*}%
And from our compression ratio equation 
\begin{equation*}
%TCIMACRO{\TeXButton{V}{\ooalign{\hfil$V$\hfil\cr\kern0.1em--\hfil\cr}}}%
%BeginExpansion
\ooalign{\hfil$V$\hfil\cr\kern0.1em--\hfil\cr}%
%EndExpansion
_{B}=\frac{%
%TCIMACRO{\TeXButton{V}{\ooalign{\hfil$V$\hfil\cr\kern0.1em--\hfil\cr}}}%
%BeginExpansion
\ooalign{\hfil$V$\hfil\cr\kern0.1em--\hfil\cr}%
%EndExpansion
_{A}}{8.00}
\end{equation*}%
so $w_{AB}$ becomes 
\begin{eqnarray*}
w_{AB} &=&\left( \frac{1}{\gamma -1}\right) \left( P_{A}\frac{%
%TCIMACRO{\TeXButton{V}{\ooalign{\hfil$V$\hfil\cr\kern0.1em--\hfil\cr}}}%
%BeginExpansion
\ooalign{\hfil$V$\hfil\cr\kern0.1em--\hfil\cr}%
%EndExpansion
_{A}^{\gamma }}{\left( \frac{%
%TCIMACRO{\TeXButton{V}{\ooalign{\hfil$V$\hfil\cr\kern0.1em--\hfil\cr}}}%
%BeginExpansion
\ooalign{\hfil$V$\hfil\cr\kern0.1em--\hfil\cr}%
%EndExpansion
_{A}}{8}\right) ^{\gamma }}\frac{%
%TCIMACRO{\TeXButton{V}{\ooalign{\hfil$V$\hfil\cr\kern0.1em--\hfil\cr}}}%
%BeginExpansion
\ooalign{\hfil$V$\hfil\cr\kern0.1em--\hfil\cr}%
%EndExpansion
_{A}}{8.00}-P_{A}%
%TCIMACRO{\TeXButton{V}{\ooalign{\hfil$V$\hfil\cr\kern0.1em--\hfil\cr}}}%
%BeginExpansion
\ooalign{\hfil$V$\hfil\cr\kern0.1em--\hfil\cr}%
%EndExpansion
_{A}\right) \\
&=&\left( \frac{1}{\gamma -1}\right) \left( 8^{\gamma }P_{A}\left( \frac{%
%TCIMACRO{\TeXButton{V}{\ooalign{\hfil$V$\hfil\cr\kern0.1em--\hfil\cr}}}%
%BeginExpansion
\ooalign{\hfil$V$\hfil\cr\kern0.1em--\hfil\cr}%
%EndExpansion
_{A}}{8}\right) -P_{A}%
%TCIMACRO{\TeXButton{V}{\ooalign{\hfil$V$\hfil\cr\kern0.1em--\hfil\cr}}}%
%BeginExpansion
\ooalign{\hfil$V$\hfil\cr\kern0.1em--\hfil\cr}%
%EndExpansion
_{A}\right) \\
&=&\left( \frac{%
%TCIMACRO{\TeXButton{V}{\ooalign{\hfil$V$\hfil\cr\kern0.1em--\hfil\cr}}}%
%BeginExpansion
\ooalign{\hfil$V$\hfil\cr\kern0.1em--\hfil\cr}%
%EndExpansion
_{A}}{\gamma -1}\right) \left( 8^{\gamma -1}P_{A}-P_{A}\right) \\
&=&\left( \frac{%
%TCIMACRO{\TeXButton{V}{\ooalign{\hfil$V$\hfil\cr\kern0.1em--\hfil\cr}}}%
%BeginExpansion
\ooalign{\hfil$V$\hfil\cr\kern0.1em--\hfil\cr}%
%EndExpansion
_{A}P_{A}}{\gamma -1}\right) \left( 8^{\gamma -1}-1\right)
\end{eqnarray*}%
So the work is 
\begin{equation*}
w_{AB}=\left( \frac{%
%TCIMACRO{\TeXButton{V}{\ooalign{\hfil$V$\hfil\cr\kern0.1em--\hfil\cr}}}%
%BeginExpansion
\ooalign{\hfil$V$\hfil\cr\kern0.1em--\hfil\cr}%
%EndExpansion
_{A}P_{A}}{\gamma -1}\right) \left( 8^{\gamma -1}-1\right)
\end{equation*}%
and we know all of these numbers so 
\begin{eqnarray*}
w_{AB} &=&\left( \frac{\left( \allowbreak 0.000\,5\unit{m}^{3}\right) \left(
100\unit{kPa}\right) }{1.4-1}\right) \left( 8^{1.4-1}-1\right) \\
&=&\allowbreak 162.\,\allowbreak 17\unit{J}
\end{eqnarray*}

But this is the work done on the gas, the useful work is the negative of this%
\begin{equation*}
w_{eng_{AB}}=-\allowbreak 162.\,\allowbreak 17\unit{J}
\end{equation*}%
Which tells us that it took energy to make the process $A\rightarrow B$
happen. You might need a pull cord to start this engine or have an electric
starter motor. The change in internal energy is 
\begin{equation*}
\Delta E_{AB}=\allowbreak 162.\,\allowbreak 17\unit{J}
\end{equation*}%
and we can find 
\begin{equation*}
E_{int}=\frac{5}{2}nRT
\end{equation*}%
but we need $n.$ We can get it from what we know at $A$ using the ideal gas
law%
\begin{eqnarray*}
n &=&\frac{P_{A}%
%TCIMACRO{\TeXButton{V}{\ooalign{\hfil$V$\hfil\cr\kern0.1em--\hfil\cr}}}%
%BeginExpansion
\ooalign{\hfil$V$\hfil\cr\kern0.1em--\hfil\cr}%
%EndExpansion
_{A}}{RT_{A}} \\
&=&\frac{\left( 100\unit{kPa}\right) \left( \allowbreak 0.000\,5\unit{m}%
^{3}\right) }{\left( 8.314\frac{\unit{J}}{\unit{mol}\unit{K}}\right) \left(
293\unit{K}\right) } \\
&=&\allowbreak 2.\,\allowbreak 052\,5\times 10^{-2}\unit{mol}
\end{eqnarray*}%
Then%
\begin{eqnarray*}
E_{A} &=&\frac{5}{2}nRT_{A} \\
&=&\frac{5}{2}\left( \allowbreak 2.\,\allowbreak 052\,5\times 10^{-2}\unit{%
mol}\right) \left( 8.314\frac{\unit{J}}{\unit{mol}\unit{K}}\right) \left( 293%
\unit{K}\right) \\
&=&\allowbreak 125.\,\allowbreak 00\unit{J}
\end{eqnarray*}

We used 
\begin{equation*}
P_{B}=P_{A}\frac{%
%TCIMACRO{\TeXButton{V}{\ooalign{\hfil$V$\hfil\cr\kern0.1em--\hfil\cr}}}%
%BeginExpansion
\ooalign{\hfil$V$\hfil\cr\kern0.1em--\hfil\cr}%
%EndExpansion
_{A}^{\gamma }}{%
%TCIMACRO{\TeXButton{V}{\ooalign{\hfil$V$\hfil\cr\kern0.1em--\hfil\cr}}}%
%BeginExpansion
\ooalign{\hfil$V$\hfil\cr\kern0.1em--\hfil\cr}%
%EndExpansion
_{B}^{\gamma }}
\end{equation*}%
but we should find a numeric answer for $P_{B}.$ Let's use our compression
ratio equation for $%
%TCIMACRO{\TeXButton{V}{\ooalign{\hfil$V$\hfil\cr\kern0.1em--\hfil\cr}}}%
%BeginExpansion
\ooalign{\hfil$V$\hfil\cr\kern0.1em--\hfil\cr}%
%EndExpansion
_{B}$ again 
\begin{eqnarray*}
P_{B} &=&P_{A}\frac{%
%TCIMACRO{\TeXButton{V}{\ooalign{\hfil$V$\hfil\cr\kern0.1em--\hfil\cr}}}%
%BeginExpansion
\ooalign{\hfil$V$\hfil\cr\kern0.1em--\hfil\cr}%
%EndExpansion
_{A}^{\gamma }}{\left( \frac{%
%TCIMACRO{\TeXButton{V}{\ooalign{\hfil$V$\hfil\cr\kern0.1em--\hfil\cr}}}%
%BeginExpansion
\ooalign{\hfil$V$\hfil\cr\kern0.1em--\hfil\cr}%
%EndExpansion
_{A}}{8}\right) ^{\gamma }} \\
&=&P_{A}\frac{%
%TCIMACRO{\TeXButton{V}{\ooalign{\hfil$V$\hfil\cr\kern0.1em--\hfil\cr}}}%
%BeginExpansion
\ooalign{\hfil$V$\hfil\cr\kern0.1em--\hfil\cr}%
%EndExpansion
_{A}^{\gamma }}{\frac{%
%TCIMACRO{\TeXButton{V}{\ooalign{\hfil$V$\hfil\cr\kern0.1em--\hfil\cr}}}%
%BeginExpansion
\ooalign{\hfil$V$\hfil\cr\kern0.1em--\hfil\cr}%
%EndExpansion
_{A}^{\gamma }}{8^{\gamma }}} \\
&=&8^{\gamma }P_{A} \\
&=&8^{1.4}\left( 100\unit{kPa}\right) \\
&=&\allowbreak 1837.\,\allowbreak 9\unit{kPa}
\end{eqnarray*}%
And finally, from the ideal gas law 
\begin{eqnarray*}
T_{B} &=&\frac{P_{B}%
%TCIMACRO{\TeXButton{V}{\ooalign{\hfil$V$\hfil\cr\kern0.1em--\hfil\cr}}}%
%BeginExpansion
\ooalign{\hfil$V$\hfil\cr\kern0.1em--\hfil\cr}%
%EndExpansion
_{B}}{nR} \\
&=&\frac{\left( 1837.\,\allowbreak 9\unit{kPa}\right) \left( 6.\,25\times
10^{-5}\unit{m}^{3}\right) }{\left( 2.\,052\,5\times 10^{-2}\unit{mol}%
\right) \left( 8.314\frac{\unit{J}}{\unit{mol}\unit{K}}\right) } \\
&=&673.\,\allowbreak 15\unit{K}
\end{eqnarray*}%
At last we can find $E_{B}$%
\begin{eqnarray*}
E_{B} &=&\frac{5}{2}nRT_{B} \\
&=&\frac{5}{2}\left( \allowbreak 2.\,\allowbreak 052\,5\times 10^{-2}\unit{%
mol}\right) \left( 8.314\frac{\unit{J}}{\unit{mol}\unit{K}}\right) \left(
673.\,\allowbreak 15\unit{K}\right) \\
&=&\allowbreak 287.\,\allowbreak 17\unit{J}
\end{eqnarray*}%
what we have learned is summarized in this table

\begin{equation*}
\begin{tabular}{|c|c|c|c|c|c|c|}
\hline
\textbf{State} & $\mathbf{T}\left( \unit{K}\right) $ & $\mathbf{P}\left( 
\unit{kPa}\right) $ & $%
%TCIMACRO{\TeXButton{V}{\ooalign{\hfil$V$\hfil\cr\kern0.1em--\hfil\cr}}}%
%BeginExpansion
\ooalign{\hfil$V$\hfil\cr\kern0.1em--\hfil\cr}%
%EndExpansion
\left( \unit{m}^{3}\right) $ & $\mathbf{n}\left( \unit{mol}\right) $ & $%
\mathbf{E}_{int}\left( \unit{J}\right) $ & \textbf{--} \\ \hline
$\mathbf{A}$ & $293$ & $100$ & $0.000\,5$ & $2.\,052\,5\times 10^{-2}$ & $%
\allowbreak 125.\,\allowbreak 00$ & \textbf{--} \\ \hline
$\mathbf{B}$ & $673.\,\allowbreak 15$ & $\allowbreak 1837.\,\allowbreak 9$ & 
$6.\,25\times 10^{-5}$ & $2.\,052\,5\times 10^{-2}$ & $287.\,\allowbreak 17$
& \textbf{--} \\ \hline
\textbf{Process} & $\mathbf{Q}\left( \unit{J}\right) $ & $\mathbf{w}%
_{int}\left( \unit{J}\right) $ & $\mathbf{\Delta E}_{int}\left( \unit{J}%
\right) $ & $\mathbf{w}_{eng}\left( \unit{J}\right) $ & $\mathbf{Q}_{h}$ & $%
\mathbf{Q}_{c}$ \\ \hline
$\mathbf{AB}$ & $0$ & $\allowbreak 162.\,\allowbreak 17\unit{J}$ & $%
\allowbreak 162.\,\allowbreak 17\unit{J}$ & $-\allowbreak 162.\,\allowbreak
17\unit{J}$ & $\mathbf{0}$ & $\mathbf{0}$ \\ \hline
\end{tabular}%
\end{equation*}

That was just one process!

\textbf{Process }$B\rightarrow C$\textbf{\ is an isochoric process.} \FRAME{%
dtbpF}{3.3304in}{2.9568in}{0pt}{}{}{Figure}{\special{language "Scientific
Word";type "GRAPHIC";maintain-aspect-ratio TRUE;display "USEDEF";valid_file
"T";width 3.3304in;height 2.9568in;depth 0pt;original-width
3.2846in;original-height 2.9136in;cropleft "0";croptop "1";cropright
"1";cropbottom "0";tempfilename 'TJXN1F0O.wmf';tempfile-properties "XPR";}}%
This is where the spark plug ignites the air-fuel mixture created by the
fuel injectors. A great deal of energy is released in a hurry. The
temperature jumps up and the pressure does too, but the volume is the same
because this process happens in a flash--literally!

So this must be an isochoric process. We expect $w_{BC}=0,$ so $\Delta
E_{BC}=Q_{BC}.$ The pressure increases, so we expect $Q_{BC}$ to be
positive. That means this energy is entering our engine. This must be part
of $Q_{h}$! We know the temperature at $C$ and we know the volume $%
%TCIMACRO{\TeXButton{V}{\ooalign{\hfil$V$\hfil\cr\kern0.1em--\hfil\cr}}}%
%BeginExpansion
\ooalign{\hfil$V$\hfil\cr\kern0.1em--\hfil\cr}%
%EndExpansion
_{C}=%
%TCIMACRO{\TeXButton{V}{\ooalign{\hfil$V$\hfil\cr\kern0.1em--\hfil\cr}}}%
%BeginExpansion
\ooalign{\hfil$V$\hfil\cr\kern0.1em--\hfil\cr}%
%EndExpansion
_{B}$ so we can find the final pressure%
\begin{eqnarray*}
P_{C} &=&\frac{nRT_{C}}{%
%TCIMACRO{\TeXButton{V}{\ooalign{\hfil$V$\hfil\cr\kern0.1em--\hfil\cr}}}%
%BeginExpansion
\ooalign{\hfil$V$\hfil\cr\kern0.1em--\hfil\cr}%
%EndExpansion
_{C}} \\
&=&\frac{\left( 2.\,052\,5\times 10^{-2}\right) \left( 8.314\frac{\unit{J}}{%
\unit{mol}\unit{K}}\right) \left( 1023\right) }{6.\,25\times 10^{-5}\unit{m}%
^{3}} \\
&=&2.\,\allowbreak 793\,1\times 10^{6}\unit{Pa} \\
&=&2793.\,\allowbreak 1\unit{kPa}
\end{eqnarray*}%
We can find the internal energy at $C$ now%
\begin{eqnarray*}
E_{C} &=&\frac{5}{2}nRT_{C} \\
&=&\frac{5}{2}\left( \allowbreak 2.\,\allowbreak 052\,5\times 10^{-2}\unit{%
mol}\right) \left( 8.314\frac{\unit{J}}{\unit{mol}\unit{K}}\right) \left(
1023\unit{K}\right) \\
&=&\allowbreak 436.\,\allowbreak 42\unit{J}
\end{eqnarray*}%
We can use 
\begin{eqnarray*}
Q &=&nC_{V}\Delta T \\
&=&n\frac{5}{2}R\left( T_{C}-T_{B}\right)
\end{eqnarray*}%
For the process $BC$ to find the energy transferred by heat%
\begin{eqnarray*}
Q_{BC} &=&\left( 2.\,052\,5\times 10^{-2}\unit{mol}\right) \frac{5}{2}\left(
8.314\frac{\unit{J}}{\unit{mol}\unit{K}}\right) \left( 1023\unit{K}%
-673.\,\allowbreak 15\unit{K}\right) \\
&=&149.\,\allowbreak 25\unit{J}
\end{eqnarray*}%
We already mentioned that $Q_{BC}=\Delta E_{int_{BC}}$ and $w_{BC}=0,$ so we
can complete our table of state variables and process variables for this
process. 
\begin{equation*}
\begin{tabular}{|c|c|c|c|c|c|c|}
\hline
\textbf{State} & $\mathbf{T}\left( \unit{K}\right) $ & $\mathbf{P}\left( 
\unit{kPa}\right) $ & $%
%TCIMACRO{\TeXButton{V}{\ooalign{\hfil$V$\hfil\cr\kern0.1em--\hfil\cr}}}%
%BeginExpansion
\ooalign{\hfil$V$\hfil\cr\kern0.1em--\hfil\cr}%
%EndExpansion
\left( \unit{m}^{3}\right) $ & $\mathbf{n}\left( \unit{mol}\right) $ & $%
\mathbf{E}_{int}\left( \unit{J}\right) $ & \textbf{--} \\ \hline
$\mathbf{B}$ & $673.\,\allowbreak 15$ & $\allowbreak 1837.\,\allowbreak 9$ & 
$6.\,25\times 10^{-5}$ & $2.\,052\,5\times 10^{-2}$ & $287.\,\allowbreak 17$
& \textbf{--} \\ \hline
$\mathbf{C}$ & $1023$ & $2793.\,\allowbreak 1$ & $6.\,25\times 10^{-5}$ & $%
2.\,052\,5\times 10^{-2}$ & $436.\,\allowbreak 42$ & \textbf{--} \\ \hline
\textbf{Process} & $\mathbf{Q}\left( \unit{J}\right) $ & $\mathbf{w}%
_{int}\left( \unit{J}\right) $ & $\mathbf{\Delta E}_{int}\left( \unit{J}%
\right) $ & $\mathbf{w}_{eng}\left( \unit{J}\right) $ & $\mathbf{Q}_{h}$ & $%
Q_{c}$ \\ \hline
$\mathbf{BC}$ & $149.\,\allowbreak 25$ & $0$ & $149.\,\allowbreak 25$ & $%
\mathbf{0}$ & $149.\,\allowbreak 25$ & $0$ \\ \hline
\end{tabular}%
\end{equation*}

\FRAME{dtbpF}{4.4227in}{3.1384in}{0pt}{}{}{Figure}{\special{language
"Scientific Word";type "GRAPHIC";maintain-aspect-ratio TRUE;display
"USEDEF";valid_file "T";width 4.4227in;height 3.1384in;depth
0pt;original-width 2.1862in;original-height 1.5437in;cropleft "0";croptop
"1";cropright "1";cropbottom "0";tempfilename
'TJXN1F0P.wmf';tempfile-properties "XPR";}}

\textbf{Process }$C\rightarrow D$\textbf{\ is another adiabatic process},
this time an expansion. This is sometimes called the \textquotedblleft power
stroke\textquotedblright\ because the hot, high pressure gas pushes the
piston down. This push is what makes the car or lawn mower (snow blower) go.
We know that $Q_{CD}=0$ so $\Delta E_{CD}=w_{CD}$

To find the ideal gas state variables at $D$ let's start with 
\begin{equation*}
P_{i}%
%TCIMACRO{\TeXButton{V}{\ooalign{\hfil$V$\hfil\cr\kern0.1em--\hfil\cr}}}%
%BeginExpansion
\ooalign{\hfil$V$\hfil\cr\kern0.1em--\hfil\cr}%
%EndExpansion
_{i}^{\gamma }=P_{f}%
%TCIMACRO{\TeXButton{V}{\ooalign{\hfil$V$\hfil\cr\kern0.1em--\hfil\cr}}}%
%BeginExpansion
\ooalign{\hfil$V$\hfil\cr\kern0.1em--\hfil\cr}%
%EndExpansion
_{f}^{\gamma }
\end{equation*}%
again, then 
\begin{equation*}
P_{D}=P_{C}\frac{%
%TCIMACRO{\TeXButton{V}{\ooalign{\hfil$V$\hfil\cr\kern0.1em--\hfil\cr}}}%
%BeginExpansion
\ooalign{\hfil$V$\hfil\cr\kern0.1em--\hfil\cr}%
%EndExpansion
_{C}^{\gamma }}{%
%TCIMACRO{\TeXButton{V}{\ooalign{\hfil$V$\hfil\cr\kern0.1em--\hfil\cr}}}%
%BeginExpansion
\ooalign{\hfil$V$\hfil\cr\kern0.1em--\hfil\cr}%
%EndExpansion
_{D}^{\gamma }}
\end{equation*}%
but from our PV\ diagram we see $%
%TCIMACRO{\TeXButton{V}{\ooalign{\hfil$V$\hfil\cr\kern0.1em--\hfil\cr}}}%
%BeginExpansion
\ooalign{\hfil$V$\hfil\cr\kern0.1em--\hfil\cr}%
%EndExpansion
_{C}=%
%TCIMACRO{\TeXButton{V}{\ooalign{\hfil$V$\hfil\cr\kern0.1em--\hfil\cr}}}%
%BeginExpansion
\ooalign{\hfil$V$\hfil\cr\kern0.1em--\hfil\cr}%
%EndExpansion
_{B}$ and $%
%TCIMACRO{\TeXButton{V}{\ooalign{\hfil$V$\hfil\cr\kern0.1em--\hfil\cr}}}%
%BeginExpansion
\ooalign{\hfil$V$\hfil\cr\kern0.1em--\hfil\cr}%
%EndExpansion
_{D}=%
%TCIMACRO{\TeXButton{V}{\ooalign{\hfil$V$\hfil\cr\kern0.1em--\hfil\cr}}}%
%BeginExpansion
\ooalign{\hfil$V$\hfil\cr\kern0.1em--\hfil\cr}%
%EndExpansion
_{A}$. This makes sense because the cylinder and piston system maximum and
minimum volumes don't change. So we can write%
\begin{equation*}
P_{D}=P_{C}\frac{%
%TCIMACRO{\TeXButton{V}{\ooalign{\hfil$V$\hfil\cr\kern0.1em--\hfil\cr}}}%
%BeginExpansion
\ooalign{\hfil$V$\hfil\cr\kern0.1em--\hfil\cr}%
%EndExpansion
_{B}^{\gamma }}{%
%TCIMACRO{\TeXButton{V}{\ooalign{\hfil$V$\hfil\cr\kern0.1em--\hfil\cr}}}%
%BeginExpansion
\ooalign{\hfil$V$\hfil\cr\kern0.1em--\hfil\cr}%
%EndExpansion
_{A}^{\gamma }}=8^{-\gamma }P_{C}
\end{equation*}%
and knowing this we can find the temperature at $D$ using the ideal gas law%
\begin{equation*}
T_{D}=\frac{P_{D}%
%TCIMACRO{\TeXButton{V}{\ooalign{\hfil$V$\hfil\cr\kern0.1em--\hfil\cr}}}%
%BeginExpansion
\ooalign{\hfil$V$\hfil\cr\kern0.1em--\hfil\cr}%
%EndExpansion
_{D}}{nR}
\end{equation*}%
Let's numerically calculate $P_{D}$ and $T_{D}$ at this point.%
\begin{eqnarray*}
P_{D} &=&8^{-1.4}\left( 2793.\,\allowbreak 1\right) \\
&=&151.\,\allowbreak 97\unit{kPa}
\end{eqnarray*}%
and 
\begin{eqnarray*}
T_{D} &=&\frac{\left( 151.\,\allowbreak 97\unit{kPa}\right) \left( 0.000\,5%
\unit{m}^{3}\right) }{\left( 2.\,052\,5\times 10^{-2}\unit{mol}\right)
\left( 8.314\frac{\unit{J}}{\unit{mol}\unit{K}}\right) } \\
&=&445.\,\allowbreak 28\unit{K}
\end{eqnarray*}%
And while we are at it we can complete our state $D$ by calculating $E_{D}$%
\begin{eqnarray*}
E_{D} &=&\frac{5}{2}nRT_{D} \\
&=&\frac{5}{2}\left( \allowbreak 2.\,\allowbreak 052\,5\times 10^{-2}\unit{%
mol}\right) \left( 8.314\frac{\unit{J}}{\unit{mol}\unit{K}}\right) \left(
445.\,\allowbreak 28\unit{K}\right) \\
&=&189.\,\allowbreak 96\unit{J}
\end{eqnarray*}%
Now for the work done, we again use our new adiabatic equation 
\begin{eqnarray*}
w_{CD} &=&\left( \frac{1}{\gamma -1}\right) \left( P_{D}%
%TCIMACRO{\TeXButton{V}{\ooalign{\hfil$V$\hfil\cr\kern0.1em--\hfil\cr}}}%
%BeginExpansion
\ooalign{\hfil$V$\hfil\cr\kern0.1em--\hfil\cr}%
%EndExpansion
_{D}-P_{C}%
%TCIMACRO{\TeXButton{V}{\ooalign{\hfil$V$\hfil\cr\kern0.1em--\hfil\cr}}}%
%BeginExpansion
\ooalign{\hfil$V$\hfil\cr\kern0.1em--\hfil\cr}%
%EndExpansion
_{C}\right) \\
&=&\left( \frac{1}{\gamma -1}\right) \left( 8^{-\gamma }P_{C}%
%TCIMACRO{\TeXButton{V}{\ooalign{\hfil$V$\hfil\cr\kern0.1em--\hfil\cr}}}%
%BeginExpansion
\ooalign{\hfil$V$\hfil\cr\kern0.1em--\hfil\cr}%
%EndExpansion
_{A}-P_{C}%
%TCIMACRO{\TeXButton{V}{\ooalign{\hfil$V$\hfil\cr\kern0.1em--\hfil\cr}}}%
%BeginExpansion
\ooalign{\hfil$V$\hfil\cr\kern0.1em--\hfil\cr}%
%EndExpansion
_{B}\right) \\
&=&\left( \frac{P_{C}}{\gamma -1}\right) \left( 8^{-\gamma }%
%TCIMACRO{\TeXButton{V}{\ooalign{\hfil$V$\hfil\cr\kern0.1em--\hfil\cr}}}%
%BeginExpansion
\ooalign{\hfil$V$\hfil\cr\kern0.1em--\hfil\cr}%
%EndExpansion
_{A}-%
%TCIMACRO{\TeXButton{V}{\ooalign{\hfil$V$\hfil\cr\kern0.1em--\hfil\cr}}}%
%BeginExpansion
\ooalign{\hfil$V$\hfil\cr\kern0.1em--\hfil\cr}%
%EndExpansion
_{B}\right)
\end{eqnarray*}%
all of which we know, so%
\begin{eqnarray*}
W_{CD} &=&\left( \frac{2793.1\unit{kPa}}{1.4-1}\right) \left( 8^{-1.4}\left(
0.000\,5\unit{m}^{3}\right) -6.\,25\times 10^{-5}\unit{m}^{3}\right) \\
&=&-246.\,\allowbreak 46\unit{J}
\end{eqnarray*}%
which completes our set for process $CD$ 
\begin{equation*}
\begin{tabular}{|c|c|c|c|c|c|c|}
\hline
\textbf{State} & $\mathbf{T}\left( \unit{K}\right) $ & $\mathbf{P}\left( 
\unit{kPa}\right) $ & $%
%TCIMACRO{\TeXButton{V}{\ooalign{\hfil$V$\hfil\cr\kern0.1em--\hfil\cr}}}%
%BeginExpansion
\ooalign{\hfil$V$\hfil\cr\kern0.1em--\hfil\cr}%
%EndExpansion
\left( \unit{m}^{3}\right) $ & $\mathbf{n}\left( \unit{mol}\right) $ & $%
\mathbf{E}_{int}\left( \unit{J}\right) $ & \textbf{--} \\ \hline
$\mathbf{C}$ & $1023$ & $2793.\,\allowbreak 1$ & $6.\,25\times 10^{-5}$ & $%
2.\,052\,5\times 10^{-2}$ & $436.\,\allowbreak 42$ & \textbf{--} \\ \hline
$\mathbf{D}$ & $445.\,\allowbreak 28$ & $151.\,\allowbreak 97$ & $0.000\,5$
& $2.\,052\,5\times 10^{-2}$ & $189.\,\allowbreak 96\unit{J}$ & \textbf{--}
\\ \hline
\textbf{Process} & $\mathbf{Q}\left( \unit{J}\right) $ & $\mathbf{w}%
_{int}\left( \unit{J}\right) $ & $\mathbf{\Delta E}_{int}\left( \unit{J}%
\right) $ & $\mathbf{w}_{eng}\left( \unit{J}\right) $ & $\mathbf{Q}_{h}$ & $%
\mathbf{Q}_{c}$ \\ \hline
$\mathbf{CD}$ & $0$ & $-246.\,\allowbreak 46$ & $-246.\,\allowbreak 46$ & $%
246.\,\allowbreak 46$ & $\mathbf{0}$ & $\mathbf{0}$ \\ \hline
\end{tabular}%
\end{equation*}

\textbf{Process }$D\rightarrow A$\textbf{\ is an isochoric process.} We are
nearing the end of the cycle, In process $D\rightarrow A$ a valve opens
letting the pressure drop quickly. \FRAME{dtbpF}{3.0891in}{2.5927in}{0pt}{}{%
}{Figure}{\special{language "Scientific Word";type
"GRAPHIC";maintain-aspect-ratio TRUE;display "USEDEF";valid_file "T";width
3.0891in;height 2.5927in;depth 0pt;original-width 3.6123in;original-height
3.0268in;cropleft "0";croptop "1";cropright "1";cropbottom "0";tempfilename
'TJXPL40R.wmf';tempfile-properties "XPR";}}Since this happens quickly, the
piston does not have time to change position, so the volume does not change.
Process $DA$ is (almost) isochoric. We already know $P_{A},$ $%
%TCIMACRO{\TeXButton{V}{\ooalign{\hfil$V$\hfil\cr\kern0.1em--\hfil\cr}}}%
%BeginExpansion
\ooalign{\hfil$V$\hfil\cr\kern0.1em--\hfil\cr}%
%EndExpansion
_{A},$ $n,$ and $T_{A}$ because that is where we started and this is a
cyclic process. But we need to find the process work, heat and change in
internal energy. Since the volume does not change, $W_{DA}=0$ and $\Delta
E_{DA}=Q_{DA}$ which we can find using 
\begin{eqnarray*}
Q_{DA} &=&nC_{V}\Delta T \\
&=&n\frac{5}{2}R\left( T_{A}-T_{D}\right)
\end{eqnarray*}%
which gives%
\begin{eqnarray*}
Q_{DA} &=&\left( 2.\,052\,5\times 10^{-2}\unit{mol}\right) \frac{5}{2}\left(
8.314\frac{\unit{J}}{\unit{mol}\unit{K}}\right) \left( 293\unit{K}%
-445.\,\allowbreak 28\unit{K}\right) \\
&=&-64.\,\allowbreak 964\unit{J}
\end{eqnarray*}%
Which tells us that we have lost energy from the engine by heat. So $Q_{DA}$
must contribute to $Q_{c}.$ 
\begin{equation*}
\begin{tabular}{|c|c|c|c|c|c|c|}
\hline
\textbf{State} & $\mathbf{T}\left( \unit{K}\right) $ & $\mathbf{P}\left( 
\unit{kPa}\right) $ & $%
%TCIMACRO{\TeXButton{V}{\ooalign{\hfil$V$\hfil\cr\kern0.1em--\hfil\cr}}}%
%BeginExpansion
\ooalign{\hfil$V$\hfil\cr\kern0.1em--\hfil\cr}%
%EndExpansion
\left( \unit{m}^{3}\right) $ & $\mathbf{n}\left( \unit{mol}\right) $ & $%
\mathbf{E}_{int}\left( \unit{J}\right) $ & \textbf{--} \\ \hline
$\mathbf{D}$ & $445.\,\allowbreak 28$ & $151.\,\allowbreak 97$ & $0.000\,5$
& $2.\,052\,5\times 10^{-2}$ & $189.\,\allowbreak 96\unit{J}$ & \textbf{--}
\\ \hline
$\mathbf{A}$ & $293$ & $100$ & $0.000\,5$ & $2.\,052\,5\times 10^{-2}$ & $%
\allowbreak 125.\,\allowbreak 00$ & \textbf{--} \\ \hline
\textbf{Process} & $\mathbf{Q}\left( \unit{J}\right) $ & $\mathbf{w}%
_{int}\left( \unit{J}\right) $ & $\mathbf{\Delta E}_{int}\left( \unit{J}%
\right) $ & $\mathbf{w}_{eng}\left( \unit{J}\right) $ & $\mathbf{Q}_{h}$ & $%
\mathbf{Q}_{c}$ \\ \hline
$\mathbf{DA}$ & $-64.\,\allowbreak 964$ & $0$ & $-64.\,\allowbreak 964$ & $0$
& $0$ & $64.\,\allowbreak 964$ \\ \hline
\end{tabular}%
\end{equation*}

\textbf{Process }$A\rightarrow I$ \textbf{\ is not a normal thermodynamic
process (and neither is process }$I\rightarrow A$\textbf{)}\FRAME{dhF}{%
3.2569in}{2.8323in}{0pt}{}{}{Figure}{\special{language "Scientific
Word";type "GRAPHIC";maintain-aspect-ratio TRUE;display "USEDEF";valid_file
"T";width 3.2569in;height 2.8323in;depth 0pt;original-width
3.2119in;original-height 2.789in;cropleft "0";croptop "1";cropright
"1";cropbottom "0";tempfilename 'TJXPL40S.wmf';tempfile-properties "XPR";}}

But we are not done. The opening of the valve let the spent fuel begin to
leave as exhaust. But we can help remove the exhaust by pushing the piston
up. This is done at constant pressure, because the valve is open to the
exhaust system air pressure. So this is almost an isobaric process. But it's
not quite isobaric because we will lose gas by venting it. The number of
moles, $n,$ will change. We can find the new number of moles using the ideal
gas law%
\begin{equation*}
P_{I}%
%TCIMACRO{\TeXButton{V}{\ooalign{\hfil$V$\hfil\cr\kern0.1em--\hfil\cr}}}%
%BeginExpansion
\ooalign{\hfil$V$\hfil\cr\kern0.1em--\hfil\cr}%
%EndExpansion
_{I}=nRT
\end{equation*}%
since $P_{I}=P_{A}$ and $%
%TCIMACRO{\TeXButton{V}{\ooalign{\hfil$V$\hfil\cr\kern0.1em--\hfil\cr}}}%
%BeginExpansion
\ooalign{\hfil$V$\hfil\cr\kern0.1em--\hfil\cr}%
%EndExpansion
_{I}=%
%TCIMACRO{\TeXButton{V}{\ooalign{\hfil$V$\hfil\cr\kern0.1em--\hfil\cr}}}%
%BeginExpansion
\ooalign{\hfil$V$\hfil\cr\kern0.1em--\hfil\cr}%
%EndExpansion
_{B}$ because the piston minimum volume is the same each time the piston
goes up. Then we an use our compression ratio for $%
%TCIMACRO{\TeXButton{V}{\ooalign{\hfil$V$\hfil\cr\kern0.1em--\hfil\cr}}}%
%BeginExpansion
\ooalign{\hfil$V$\hfil\cr\kern0.1em--\hfil\cr}%
%EndExpansion
_{A}$ and $%
%TCIMACRO{\TeXButton{V}{\ooalign{\hfil$V$\hfil\cr\kern0.1em--\hfil\cr}}}%
%BeginExpansion
\ooalign{\hfil$V$\hfil\cr\kern0.1em--\hfil\cr}%
%EndExpansion
_{I}$%
\begin{equation*}
\frac{%
%TCIMACRO{\TeXButton{V}{\ooalign{\hfil$V$\hfil\cr\kern0.1em--\hfil\cr}}}%
%BeginExpansion
\ooalign{\hfil$V$\hfil\cr\kern0.1em--\hfil\cr}%
%EndExpansion
_{A}}{%
%TCIMACRO{\TeXButton{V}{\ooalign{\hfil$V$\hfil\cr\kern0.1em--\hfil\cr}}}%
%BeginExpansion
\ooalign{\hfil$V$\hfil\cr\kern0.1em--\hfil\cr}%
%EndExpansion
_{I}}=8.00
\end{equation*}%
or, solving for $%
%TCIMACRO{\TeXButton{V}{\ooalign{\hfil$V$\hfil\cr\kern0.1em--\hfil\cr}}}%
%BeginExpansion
\ooalign{\hfil$V$\hfil\cr\kern0.1em--\hfil\cr}%
%EndExpansion
_{I}$ 
\begin{equation*}
%TCIMACRO{\TeXButton{V}{\ooalign{\hfil$V$\hfil\cr\kern0.1em--\hfil\cr}}}%
%BeginExpansion
\ooalign{\hfil$V$\hfil\cr\kern0.1em--\hfil\cr}%
%EndExpansion
_{I}=8/%
%TCIMACRO{\TeXButton{V}{\ooalign{\hfil$V$\hfil\cr\kern0.1em--\hfil\cr}}}%
%BeginExpansion
\ooalign{\hfil$V$\hfil\cr\kern0.1em--\hfil\cr}%
%EndExpansion
_{A}
\end{equation*}%
which we could calculate, but it is just the same as $%
%TCIMACRO{\TeXButton{V}{\ooalign{\hfil$V$\hfil\cr\kern0.1em--\hfil\cr}}}%
%BeginExpansion
\ooalign{\hfil$V$\hfil\cr\kern0.1em--\hfil\cr}%
%EndExpansion
_{B}$

We want $n_{I},$ so let's use the ideal gas law for points $A$ and $I.$ We
get 
\begin{equation*}
P_{A}%
%TCIMACRO{\TeXButton{V}{\ooalign{\hfil$V$\hfil\cr\kern0.1em--\hfil\cr}}}%
%BeginExpansion
\ooalign{\hfil$V$\hfil\cr\kern0.1em--\hfil\cr}%
%EndExpansion
_{A}=n_{A}RT_{A}
\end{equation*}%
and 
\begin{equation*}
P_{I}%
%TCIMACRO{\TeXButton{V}{\ooalign{\hfil$V$\hfil\cr\kern0.1em--\hfil\cr}}}%
%BeginExpansion
\ooalign{\hfil$V$\hfil\cr\kern0.1em--\hfil\cr}%
%EndExpansion
_{I}=n_{I}RT_{I}
\end{equation*}%
We can divide these two equations%
\begin{equation*}
\frac{P_{A}%
%TCIMACRO{\TeXButton{V}{\ooalign{\hfil$V$\hfil\cr\kern0.1em--\hfil\cr}}}%
%BeginExpansion
\ooalign{\hfil$V$\hfil\cr\kern0.1em--\hfil\cr}%
%EndExpansion
_{A}}{P_{I}%
%TCIMACRO{\TeXButton{V}{\ooalign{\hfil$V$\hfil\cr\kern0.1em--\hfil\cr}}}%
%BeginExpansion
\ooalign{\hfil$V$\hfil\cr\kern0.1em--\hfil\cr}%
%EndExpansion
_{I}}=\frac{n_{A}RT_{A}}{n_{I}RT_{I}}
\end{equation*}%
and note that we have our compression ratio hidden in the left hand side,
and recall that $P_{I}=P_{A}$ so the left hand side is just $8.$ And on the
right hand side $T_{A}=T_{I}$ because we didn't change the temperature of
the gas, we just moved it. 
\begin{equation*}
8=\frac{n_{A}}{n_{I}}
\end{equation*}%
We can solve this for the number of moles at piston $I$%
\begin{equation*}
n_{I}=\frac{n_{A}}{8}
\end{equation*}%
And we have all the numbers, so 
\begin{eqnarray*}
n_{I} &=&\frac{\left( 2.\,052\,5\times 10^{-2}\unit{mol}\right) }{8} \\
&=&\allowbreak 2.\,\allowbreak 565\,6\times 10^{-3}\unit{mol}
\end{eqnarray*}%
And then $E_{I}$ will be 
\begin{eqnarray*}
E_{D} &=&\frac{5}{2}n_{I}RT_{I} \\
&=&\frac{5}{2}\left( 2.\,\allowbreak 565\,6\times 10^{-3}\unit{mol}\right)
\left( 8.314\frac{\unit{J}}{\unit{mol}\unit{K}}\right) \left( 293\unit{K}%
\right)  \\
&=&15.\,\allowbreak 625\unit{J}
\end{eqnarray*}%
The work done for a true isobaric process is 
\begin{equation*}
w_{AI}=P_{A}\left( 
%TCIMACRO{\TeXButton{V}{\ooalign{\hfil$V$\hfil\cr\kern0.1em--\hfil\cr}}}%
%BeginExpansion
\ooalign{\hfil$V$\hfil\cr\kern0.1em--\hfil\cr}%
%EndExpansion
_{I}-%
%TCIMACRO{\TeXButton{V}{\ooalign{\hfil$V$\hfil\cr\kern0.1em--\hfil\cr}}}%
%BeginExpansion
\ooalign{\hfil$V$\hfil\cr\kern0.1em--\hfil\cr}%
%EndExpansion
_{A}\right) 
\end{equation*}%
but this assumed that $n$ did not change. Really we are just moving the gas,
and it doesn't take much of a force to move the little bit of gas. So to a
good approximation we can say 
\begin{eqnarray*}
w_{int} &=&F\Delta x \\
&\approx &0
\end{eqnarray*}%
and we can ignore $w_{AI}.$

For finding $Q_{AI}$ we would like to use 
\begin{equation*}
Q=nC_{P}\Delta T
\end{equation*}%
but we quickly run into trouble again because $n$ changed! We are using
convection to remove heat instead of conduction. We can get around this by
finding $\Delta E_{int}$ directly. 
\begin{eqnarray*}
\Delta E_{AI} &=&E_{I}-E_{A} \\
&=&15.\,\allowbreak 625\unit{J}-125.\,\allowbreak 00\unit{J} \\
&=&-109.\,\allowbreak 38\unit{J}
\end{eqnarray*}%
From this and the first law%
\begin{equation*}
\Delta E_{int}=Q+w_{int}
\end{equation*}%
We can find $Q$%
\begin{eqnarray*}
Q &\approx &\Delta E_{int} \\
&\approx &-109.\,\allowbreak 38\unit{J}
\end{eqnarray*}

Energy is leaving with the gas, so we seem to pick up a contribution to $%
Q_{c}.$ But really this is not part of our thermodynamic process, because we
mechanically removed the energy by removing the gas. Let's denote this by
putting our energy in parenthesis to remind us that this is a more
mechanical than thermal process.%
\begin{equation*}
\begin{tabular}{|c|c|c|c|c|c|c|}
\hline
\textbf{State} & $\mathbf{T}\left( \unit{K}\right) $ & $\mathbf{P}\left( 
\unit{kPa}\right) $ & $%
%TCIMACRO{\TeXButton{V}{\ooalign{\hfil$V$\hfil\cr\kern0.1em--\hfil\cr}}}%
%BeginExpansion
\ooalign{\hfil$V$\hfil\cr\kern0.1em--\hfil\cr}%
%EndExpansion
\left( \unit{m}^{3}\right) $ & $\mathbf{n}\left( \unit{mol}\right) $ & $%
\mathbf{E}_{int}\left( \unit{J}\right) $ & \textbf{--} \\ \hline
$\mathbf{A}$ & $293$ & $100$ & $0.000\,5$ & $2.\,052\,5\times 10^{-2}$ & $%
\allowbreak 125.\,\allowbreak 00$ & \textbf{--} \\ \hline
$\mathbf{I}$ & $293$ & $100$ & $6.\,25\times 10^{-5}$ & $2.\,\allowbreak
565\,6\times 10^{-3}$ & $15.\,\allowbreak 625\unit{J}$ & \textbf{--} \\ 
\hline
\textbf{Process} & $\mathbf{Q}\left( \unit{J}\right) $ & $\mathbf{w}%
_{int}\left( \unit{J}\right) $ & $\mathbf{\Delta E}_{int}\left( \unit{J}%
\right) $ & $\mathbf{w}_{eng}\left( \unit{J}\right) $ & $\mathbf{Q}_{h}$ & $%
\mathbf{Q}_{c}$ \\ \hline
$\mathbf{AI}$ & $\left( -109.\,\allowbreak 38\unit{J}\right) $ & $\sim 0$ & $%
\left( -109.\,\allowbreak 38\right) $ & $\sim 0$ & $0$ & $\left(
109.\,\allowbreak 38\right) $ \\ \hline
\end{tabular}%
\end{equation*}%
$\allowbreak $\FRAME{dhF}{3.0606in}{3.2932in}{0pt}{}{}{Figure}{\special%
{language "Scientific Word";type "GRAPHIC";maintain-aspect-ratio
TRUE;display "USEDEF";valid_file "T";width 3.0606in;height 3.2932in;depth
0pt;original-width 3.0173in;original-height 3.2474in;cropleft "0";croptop
"1";cropright "1";cropbottom "0";tempfilename
'TJXPL40T.wmf';tempfile-properties "XPR";}}Now the exhaust valve closes, and
the intake valve opens. The new fuel mixture is brought in. We essentially
reverse process $AI$ in process $IA.$ 
\begin{equation*}
\begin{tabular}{|c|c|c|c|c|c|c|}
\hline
\textbf{State} & $\mathbf{T}\left( \unit{K}\right) $ & $\mathbf{P}\left( 
\unit{kPa}\right) $ & $%
%TCIMACRO{\TeXButton{V}{\ooalign{\hfil$V$\hfil\cr\kern0.1em--\hfil\cr}}}%
%BeginExpansion
\ooalign{\hfil$V$\hfil\cr\kern0.1em--\hfil\cr}%
%EndExpansion
\left( \unit{m}^{3}\right) $ & $\mathbf{n}\left( \unit{mol}\right) $ & $%
\mathbf{E}_{int}\left( \unit{J}\right) $ & \textbf{--} \\ \hline
$\mathbf{I}$ & $293$ & $100$ & $6.\,25\times 10^{-5}$ & $2.\,\allowbreak
565\,7\times 10^{-3}$ & $15.\,\allowbreak 625\unit{J}$ & \textbf{--} \\ 
\hline
$\mathbf{A}$ & $293$ & $100$ & $0.000\,5$ & $2.\,052\,5\times 10^{-2}$ & $%
\allowbreak 125.\,\allowbreak 00$ & \textbf{--} \\ \hline
\textbf{Process} & $\mathbf{Q}\left( \unit{J}\right) $ & $\mathbf{w}%
_{int}\left( \unit{J}\right) $ & $\mathbf{\Delta E}_{int}\left( \unit{J}%
\right) $ & $\mathbf{w}_{eng}\left( \unit{J}\right) $ & $\mathbf{Q}_{h}$ & $%
\mathbf{Q}_{c}$ \\ \hline
$\mathbf{IA}$ & $\left( 109.\,\allowbreak 38\unit{J}\right) $ & $\sim 0$ & $%
\left( 109.\,\allowbreak 38\unit{J}\right) $ & $\sim 0$ & $\left(
-109.\,\allowbreak 38\unit{J}\right) $ & $0$ \\ \hline
\end{tabular}%
\end{equation*}%
But there is a difference. The new gas has useful fuel. and that means we
are ready to complete the cycle again by compressing the gas, igniting it,
and receiving the power stroke and exhaust. Note that Processes $AI$ and $IA$
have a net zero effect thermodynamically, but are very important to making
your car or lawn mower go!

Putting it all together gives our PV\ diagram for the whole cycle. Note that
since $n$ changed for the $AI-IA$ processes, it is not really possible to
include it on the PV\ diagram. That is why it has a dashed line. We know the
volume and pressure along the path, but the predicted temperature from the
graph may not be right because $n$ changed too.\FRAME{dtbpF}{3.4895in}{%
2.4751in}{0in}{}{}{Figure}{\special{language "Scientific Word";type
"GRAPHIC";maintain-aspect-ratio TRUE;display "USEDEF";valid_file "T";width
3.4895in;height 2.4751in;depth 0in;original-width 3.4437in;original-height
2.4336in;cropleft "0";croptop "1";cropright "1";cropbottom "0";tempfilename
'TJXPL40U.wmf';tempfile-properties "XPR";}}

Let's summarize what we found in large tables. First the ideal gas state
variables and $E_{int}$

\begin{equation*}
\begin{tabular}{|c|c|c|c|c|c|}
\hline
\textbf{State} & $\mathbf{T}\left( \unit{K}\right) $ & $\mathbf{P}\left( 
\unit{kPa}\right) $ & $%
%TCIMACRO{\TeXButton{V}{\ooalign{\hfil$V$\hfil\cr\kern0.1em--\hfil\cr}}}%
%BeginExpansion
\ooalign{\hfil$V$\hfil\cr\kern0.1em--\hfil\cr}%
%EndExpansion
\left( \unit{m}^{3}\right) $ & $\mathbf{n}\left( \unit{mol}\right) $ & $%
\mathbf{E}_{int}\left( \unit{J}\right) $ \\ \hline
$\mathbf{A}$ & $293$ & $100$ & $0.000\,5$ & $2.\,052\,5\times 10^{-2}$ & $%
\allowbreak 125.\,\allowbreak 00$ \\ \hline
$\mathbf{B}$ & $673.\,\allowbreak 15$ & $\allowbreak 1837.\,\allowbreak 9$ & 
$6.\,25\times 10^{-5}$ & $2.\,052\,5\times 10^{-2}$ & $287.\,\allowbreak 17$
\\ \hline
$\mathbf{C}$ & $1023$ & $2793.\,\allowbreak 1$ & $6.\,25\times 10^{-5}$ & $%
2.\,052\,5\times 10^{-2}$ & $436.\,\allowbreak 42$ \\ \hline
$\mathbf{D}$ & $445.\,\allowbreak 28$ & $151.\,\allowbreak 97$ & $0.000\,5$
& $2.\,052\,5\times 10^{-2}$ & $189.\,\allowbreak 96\unit{J}$ \\ \hline
$\mathbf{I}$ & $293$ & $100$ & $6.\,25\times 10^{-5}$ & $2.\,\allowbreak
565\,6\times 10^{-3}$ & $15.\,\allowbreak 625$ \\ \hline
\end{tabular}%
\end{equation*}

and now the values of $Q,$ $w,$ and $E_{int}$ for the processes. We have
included $w_{eng},$ $Q_{h}$ and $Q_{c}$ as well 
\begin{equation*}
\begin{tabular}{|c|c|c|c|c|c|c|}
\hline
\textbf{Process} & $\mathbf{Q}\left( \unit{J}\right) $ & $\mathbf{w}%
_{int}\left( \unit{J}\right) $ & $\mathbf{\Delta E}_{int}\left( \unit{J}%
\right) $ & $\mathbf{w}_{eng}\left( \unit{J}\right) $ & $\mathbf{Q}_{h}$ & $%
\mathbf{Q}_{c}$ \\ \hline
$\mathbf{AB}$ & $0$ & $\allowbreak 162.\,\allowbreak 17$ & $\allowbreak
162.\,\allowbreak 17$ & $-\allowbreak 162.\,\allowbreak 17$ & $\mathbf{0}$ & 
$\mathbf{0}$ \\ \hline
$\mathbf{BC}$ & $149.\,\allowbreak 25$ & $0$ & $149.\,\allowbreak 25$ & $%
\mathbf{0}$ & $149.\,\allowbreak 25$ & $0$ \\ \hline
$\mathbf{CD}$ & $0$ & $-246.\,\allowbreak 46$ & $-246.\,\allowbreak 46$ & $%
246.\,\allowbreak 46$ & $\mathbf{0}$ & $\mathbf{0}$ \\ \hline
$\mathbf{DA}$ & $-64.\,\allowbreak 964$ & $0$ & $-64.\,\allowbreak 964$ & $0$
& $0$ & $64.\,\allowbreak 964$ \\ \hline
$\mathbf{AI}$ & $\left( -109.\,\allowbreak 38\right) $ & $\left( \sim
0\right) $ & $\left( -109.\,\allowbreak 38\right) $ & $\left( \sim 0\right) $
& $\left( 0\right) $ & $\left( 109.\,\allowbreak 38\right) $ \\ \hline
$\mathbf{IA}$ & $\left( 109.\,\allowbreak 38\right) $ & $\left( \sim
0\right) $ & $\left( 109.\,\allowbreak 38\right) $ & $\left( \sim 0\right) $
& $\left( -109.\,\allowbreak 38\right) $ & $\left( 0\right) $ \\ \hline
Whole Cycle & $84.\,\allowbreak 286$ & $-84.\,\allowbreak 29$ & $0$ & $%
\allowbreak 84.\,\allowbreak 29$ & $149.\,\allowbreak 25$ & $%
64.\,\allowbreak 964$ \\ \hline
\end{tabular}%
\end{equation*}%
Note that the last row of the table is a total for the entire cycle. We see
that we have useful work of $84.3\unit{J}$ and that $Q_{h}=258.\,\allowbreak
63\unit{J}$ and $Q_{c}=174.\,\allowbreak 34.$ We can use this to identify
the processes that make $Q_{h}$ and $Q_{c}$ happen on our graph. The useful
work is the difference between the work done in the two adiabatic processes. 
\FRAME{dtbpF}{4.6172in}{3.6927in}{0pt}{}{}{Figure}{\special{language
"Scientific Word";type "GRAPHIC";maintain-aspect-ratio TRUE;display
"USEDEF";valid_file "T";width 4.6172in;height 3.6927in;depth
0pt;original-width 4.5653in;original-height 3.6452in;cropleft "0";croptop
"1";cropright "1";cropbottom "0";tempfilename
'TJXPL40V.wmf';tempfile-properties "XPR";}}We can see that it takes some
work to make process $AB$ happen. And recall that that is why we need a
starter motor or a pull cord to get our engine started. Once the engine gets
going the momentum from the piston and crank shaft keep the engine running.
But still, process $AB$ requires energy input by work, so it is like a
resistance for the cycle.

\section{Otto Efficiency}

What is the efficiency of this cycle? We don't have a really good way with
our simple view of thermodynamics to accurately account for the change in
chemical potential energy involved with state $I.$ The energy transfer in
processes $AI$ and $IA$ are, really mechanical energy processes. So for this
class it might be more fair to say that processes $AI$ and $IA$ have no net
effect (except for refueling) and exclude their effect from $Q_{c}$ and $%
Q_{h}$ then ignoring the $I$ processes%
\begin{eqnarray*}
\eta &=&\frac{w_{eng}}{Q_{h}} \\
&=&\frac{84.\,\allowbreak 29}{149.\,\allowbreak 25} \\
&=&\allowbreak 0.564\,76
\end{eqnarray*}

For a gasoline engine, it is often quoted that the efficiency is%
\begin{equation*}
\eta _{otto}=1-\frac{1}{\left( \frac{%
%TCIMACRO{\TeXButton{V}{\ooalign{\hfil$V$\hfil\cr\kern0.1em--\hfil\cr}}}%
%BeginExpansion
\ooalign{\hfil$V$\hfil\cr\kern0.1em--\hfil\cr}%
%EndExpansion
_{A}}{%
%TCIMACRO{\TeXButton{V}{\ooalign{\hfil$V$\hfil\cr\kern0.1em--\hfil\cr}}}%
%BeginExpansion
\ooalign{\hfil$V$\hfil\cr\kern0.1em--\hfil\cr}%
%EndExpansion
_{B}}\right) ^{\gamma -1}}
\end{equation*}%
Let's test this. 
\begin{eqnarray*}
\eta _{otto} &=&1-\frac{1}{\left( \frac{0.000\,5}{6.\,25\times 10^{-5}}%
\right) ^{1.4-1}} \\
&=&0.564\,72\allowbreak
\end{eqnarray*}%
which is nearly the same. But you might guess that this is very optimistic.
Efficiencies for actual engines come in at values around $20\%.$

This was a lot of work, but then engine design is a lot of work! And that is
what they pay engineers to do.

But we should ask, it this engine possible? Does it violate our maximum
efficiency criterion from a Carnot cycle operating between the same two
temperatures?

If you were paying close attention you would have noticed that we chose the
temperatures for out Otto cycle engine to be the same as the temperatures
for the Carnot cycle from last lecture. The the Carnot efficiency would be

\begin{eqnarray*}
\eta _{C} &=&1-\frac{293\unit{K}}{1023\unit{K}} \\
&=&0.713\,59
\end{eqnarray*}%
which is larger than our 
\begin{equation*}
\eta _{otto}=0.564\,72\allowbreak 
\end{equation*}%
And because 
\begin{equation*}
\eta _{otto}<\eta _{C}
\end{equation*}%
our Otto cycle isn't ruled out by our Carnot criterion. 

So the Otto cycle example might be achievable. Of course, we can tell that
our Otto cycle is still an idealization. We recognize that we neglected all
friction, and assumed perfect adiabatic processes. But we can see that there
is some hope of making our Otto cycle-based engines closer to the $56.4\%$
efficiency that we calculated. But no Otto cycle that operates between the
same $T_{c}=293\unit{K}$ and $T_{h}=1023\unit{K}$ can ever be better than $%
71\%$ efficient because that is the maximum efficiency for a heat engine
operating between these two temperatures.

This example demonstrated how we can use the Carnot cycle to test a somewhat
real design, and as we developed the Carnot criterion we did make some
statements about the second law of thermodynamics.  But if you were left
with the feeling that we really didn't determine what the second law meant
by looking at the physics of what is happening to our gas, you were right.
Let's take that on in our last lecture.

\chapter{38 Microscopic Energy Transfer Considerations 2.4.6, 2.4.7}

So far, our statements of the second law of thermodynamics have been very
practical. They tell us what will and what won't happen. Like we won't be
able to make an engine that is more efficient than a Carnot engine. But
surely the reason for what happens and what doesn't has to do with the
molecules and forces that act on them. We found that we could describe
things like pressure and temperature in terms of what happened to the
molecules. It is time to do tat for the second law of thermodynamics. This
won't be a complete description of the second law, but will be a good start
for your upper division thermodynamics classes.

%TCIMACRO{\TeXButton{\chaptermark{Chapter 40}}{\chaptermark{Chapter 40}}}%
%BeginExpansion
\chaptermark{Chapter 40}%
%EndExpansion

%TCIMACRO{%
%\TeXButton{Fundamental Concepts}{\vspace{0.10in}\hspace{-0.37in}{\Large Fundamental Concepts\vspace{0.2in}}}}%
%BeginExpansion
\vspace{0.10in}\hspace{-0.37in}{\Large Fundamental Concepts\vspace{0.2in}}%
%EndExpansion

\begin{enumerate}
\item Microscopic idea of Energy Transfer

\item Microscopic idea of disorder

\item Entropy
\end{enumerate}

\section{Microscopic View of Energy Transfer}

Suppose we have a chamber divided into two volumes. The volumes are
separated by a membrane. If they have different temperatures, we expect that
energy will transfer across the membrane from one volume to the other, but
how?

Think of a basket ball court in a gym. In our gyms, we have large nets that
divide courts to keep balls from passing from one court to another. Now
think of millions of basket balls in each side of the gym net. The basket
balls all are moving. Some of the basket balls will strike the net. If a
ball on the other side strikes the net at a particular place and time , and
a ball from the other side of the net strikes at the same place and time,
the collision between the balls will transfer energy. This would be a
conservation of momentum and energy problem. The slower ball would gain
energy, the faster ball would loose energy. \FRAME{dhF}{2.9914in}{2.2329in}{%
0pt}{}{}{Figure}{\special{language "Scientific Word";type
"GRAPHIC";maintain-aspect-ratio TRUE;display "USEDEF";valid_file "T";width
2.9914in;height 2.2329in;depth 0pt;original-width 2.9473in;original-height
2.1923in;cropleft "0";croptop "1";cropright "1";cropbottom "0";tempfilename
'SW9PJL1B.wmf';tempfile-properties "XPR";}}

This is analogous to our gasses in chambers. The molecules, like the
basketballs, transfer energy by collision.

The internal energy for each side of the chamber depends on both $N,$ then
number of molecules, and $T,$ the temperature.

\begin{eqnarray*}
E_{1i} &=&\frac{3}{2}N_{1}k_{B}T_{1i} \\
E_{2i} &=&\frac{3}{2}N_{2}k_{B}T_{2i}
\end{eqnarray*}

The total energy is 
\begin{equation*}
E_{tot}=E_{1i}+E_{2i}
\end{equation*}%
if our container does not change volume (no work done) and is insolated ($%
Q=0 $) then $E_{tot}$ will not change. But eventually the two sides come to
the same temperature%
\begin{equation*}
T_{f}=T_{1f}=T_{2f}
\end{equation*}%
This happens when the average energy of the molecules are the same. On
average, neither side gains energy due to the collisions on the membrane
(think of the basketballs bouncing off each other through the net). We can
rewrite our energy equation as the average energy per molecule for the final
condition as 
\begin{equation*}
\bar{E}_{tot}=\bar{E}_{1f}=\bar{E}_{2f}
\end{equation*}%
because the average energies will all be the same. We could say that the
total energy of side $1$ is just about $N_{1}$ times the average energy per
molecule on side $1.$ 
\begin{equation*}
E_{1f}=N_{1}\frac{3}{2}k_{B}T_{1f}=N_{1}\bar{E}_{1}
\end{equation*}%
We could do the same for side $2.$%
\begin{equation*}
E_{2f}=N_{2}\frac{3}{2}k_{B}T_{2f}=N_{2}\bar{E}_{2}
\end{equation*}%
We can identify the average energy per molecule for each side. 
\begin{eqnarray*}
\bar{E}_{1} &=&\frac{3}{2}k_{B}T_{1f} \\
\bar{E}_{2} &=&\frac{3}{2}k_{B}T_{2f}
\end{eqnarray*}%
The we could write these average energies per molecule as 
\begin{eqnarray*}
\bar{E}_{1} &=&\frac{E_{1f}}{N_{1}} \\
\bar{E}_{2} &=&\frac{E_{2f}}{N_{2}}
\end{eqnarray*}%
and thinking about it this would mean 
\begin{equation*}
\bar{E}_{tot}=\frac{E_{tot}}{N_{1}+N_{2}}
\end{equation*}%
The total average energy per molecule%
\begin{equation*}
\bar{E}_{tot}=\bar{E}_{1f}=\bar{E}_{2f}
\end{equation*}%
becomes 
\begin{equation*}
\frac{E_{tot}}{N_{1}+N_{2}}=\frac{E_{1f}}{N_{1}}=\frac{E_{2f}}{N_{2}}
\end{equation*}%
Which gives a set of expressions for the final energy of each side%
\begin{eqnarray*}
E_{1f} &=&\frac{N_{1}E_{tot}}{N_{1}+N_{2}} \\
E_{2f} &=&\frac{N_{2}E_{tot}}{N_{1}+N_{2}}
\end{eqnarray*}%
and we can verify that 
\begin{eqnarray*}
E_{tot} &=&E_{1f}+E_{2f} \\
&=&\frac{N_{1}E_{tot}}{N_{1}+N_{2}}+\frac{N_{2}E_{tot}}{N_{1}+N_{2}} \\
&=&E_{tot}
\end{eqnarray*}%
Let's take what we have learned and mathematically express this as
macroscopic formula and show it at a microscopic level as a problem example.

Let's show that if $Q_{tot}=0,$ and no work is done, then%
\begin{equation*}
Q_{1}=-Q_{2}
\end{equation*}%
This is a calorimetry problem! And we know how to do calorimetry problems.

Start with the first law, 
\begin{eqnarray*}
Q_{1} &=&\Delta E_{1}=E_{1f}-E_{1i} \\
Q_{2} &=&\Delta E_{2}=E_{2f}-E_{2i}
\end{eqnarray*}%
Since no work is done. Let's take $Q_{1}$ first 
\begin{equation*}
Q_{1}=E_{1f}-E_{1i}
\end{equation*}%
And we can use our equation for $E_{1f}$%
\begin{equation*}
E_{1f}=\frac{N_{1}E_{tot}}{N_{1}+N_{2}}
\end{equation*}%
to write $Q_{1}$ as 
\begin{equation*}
Q_{1}=\frac{N_{1}\left( E_{1i}+E_{2i}\right) }{N_{1}+N_{2}}-E_{1i}
\end{equation*}%
This is fairly simple, but we can use some algebra to make it a little more
complicated%
\begin{eqnarray*}
Q_{1} &=&\frac{N_{1}\left( E_{1i}+E_{2i}\right) }{N_{1}+N_{2}}-E_{1i}\frac{%
N_{1}+N_{2}}{N_{1}+N_{2}} \\
&=&\frac{N_{1}E_{1i}+N_{1}E_{2i}-E_{1i}N_{1}-E_{1i}N_{2}}{N_{1}+N_{2}} \\
&=&\frac{N_{1}E_{2i}-E_{1i}N_{2}}{N_{1}+N_{2}}
\end{eqnarray*}%
At this point, let's play a mathematical trick. Note that 
\begin{equation*}
0=\frac{N_{2}E_{2i}}{N_{1}+N_{2}}-\frac{N_{2}E_{2i}}{N_{1}+N_{2}}
\end{equation*}%
This seems like a funny way to write $0,$ but it works. Now let's take 
\begin{eqnarray*}
Q_{1} &=&\frac{N_{1}E_{2i}-E_{1i}N_{2}}{N_{1}+N_{2}} \\
&=&\frac{N_{1}E_{2i}}{N_{1}+N_{2}}-\frac{E_{1i}N_{2}}{N_{1}+N_{2}}+0
\end{eqnarray*}%
Adding zero does not change the value of $Q_{1}$ so%
\begin{eqnarray*}
Q_{1} &=&\frac{N_{1}E_{2i}}{N_{1}+N_{2}}-\frac{E_{1i}N_{2}}{N_{1}+N_{2}}+%
\frac{N_{2}E_{2i}}{N_{1}+N_{2}}-\frac{N_{2}E_{2i}}{N_{1}+N_{2}} \\
&=&\frac{N_{1}E_{2i}}{N_{1}+N_{2}}+\frac{N_{2}E_{2i}}{N_{1}+N_{2}}-\frac{%
E_{1i}N_{2}}{N_{1}+N_{2}}-\frac{N_{2}E_{2i}}{N_{1}+N_{2}} \\
&=&\frac{N_{1}E_{2i}+N_{2}E_{2i}}{N_{1}+N_{2}}-\frac{E_{1i}N_{2}}{N_{1}+N_{2}%
}-\frac{N_{2}E_{2i}}{N_{1}+N_{2}}
\end{eqnarray*}%
The first term is just $E_{2i},$ and using $E_{tot}=E_{1i}+E_{2i}$ again we
see 
\begin{eqnarray*}
Q_{1} &=&E_{2i}-\left( \frac{\left( E_{1i}N_{2}+N_{2}E_{2i}\right) }{%
N_{1}+N_{2}}\right) \\
&=&E_{2i}-\left( \frac{N_{2}\left( E_{1i}+E_{2i}\right) }{N_{1}+N_{2}}\right)
\\
&=&E_{2i}-\left( \frac{N_{2}\left( E_{tot}\right) }{N_{1}+N_{2}}\right) \\
&=&E_{2i}-E_{2f} \\
&=&-Q_{2}
\end{eqnarray*}%
which is what we expect. We have shown our calorimetry equation using the
ideas of microscopic energy transport for ideal gasses. This was
interesting, but we still don't have a microscopic statement of the second
law of thermodynamics. Let's consider motion of the molecules again.

\section{The arrow of time.}

Why does time go only one way? The answer to this question may surprise you. 
\emph{We don't know.} There is nothing in our laws of motion that tells us
that time should only go one way. Take motion under constant acceleration.%
\begin{equation*}
y_{f}=y_{i}+v_{o}\Delta t+\frac{1}{2}a\Delta t^{2}
\end{equation*}%
and let's start with $t_{i}=0$ so 
\begin{equation*}
y_{f}=y_{i}+v_{o}t_{f}+\frac{1}{2}at_{f}^{2}
\end{equation*}%
this equation describes the motion of an object under constant acceleration,
like that of gravity near the Earth's surface. But nothing in the equation
tells us that time must take on ever grater positive values.

The truth is that there is only one physical law that tells us that time
goes only one way, and that law is an empirical law that says essentially
\textquotedblleft time only goes one way.\textquotedblright\ We experience
time going one way, so we developed a law that says so. This is not a
terribly convincing argument. But this is the state of physics.

Some processes (really all processes) can't actually run forward or
backward. Idealized processes can be run in reverse, but not real process.
We call processes that can only work one direction \emph{irreversible
processes}. If we shot a movie of an irreversible process and ran the movie
backwards, it would be funny, because we would know that the reverse process
can't happen.

Let's take an example. Suppose our process is dropping a marker on the
floor. Because the marker bounces are not perfectly elastic, the potential
energy of the marker is converted into heat energy. If we added the lost
energy back to the marker on the floor, it won't bounce up to a height for
us to catch. This process of dropping the marker is irreversible.

To make this idea of irreversible processes and the arrow of time sound a
little better, let's study thermodynamic situations and see that at least
the arrow of time makes sense. To do this, let's take a statistical approach
(because we have learned that in thermodynamics we tend to use averages or
rms values).

\section{Disorder}

Suppose we have a system that has particles and places for those particles
to be placed. We can start with a small space (see figure below). If I have
one particle moving randomly, and divide my space into two compartments,
then I have a probability of $\frac{1}{2}$ that the particle will be in the
left compartment.\FRAME{dhF}{3.8199in}{0.4532in}{0pt}{}{}{Figure}{\special%
{language "Scientific Word";type "GRAPHIC";maintain-aspect-ratio
TRUE;display "USEDEF";valid_file "T";width 3.8199in;height 0.4532in;depth
0pt;original-width 7.5766in;original-height 0.8743in;cropleft "0";croptop
"1";cropright "1";cropbottom "0";tempfilename
'SW9PJL1C.wmf';tempfile-properties "XPR";}}Now if I have two particles, the
probability that two particles will be in the left compartment is $\left( 
\frac{1}{2}\right) ^{2}$ or $\frac{1}{4}.$ \FRAME{dhF}{2.8928in}{0.8233in}{%
0pt}{}{}{Figure}{\special{language "Scientific Word";type
"GRAPHIC";maintain-aspect-ratio TRUE;display "USEDEF";valid_file "T";width
2.8928in;height 0.8233in;depth 0pt;original-width 2.8496in;original-height
0.7904in;cropleft "0";croptop "1";cropright "1";cropbottom "0";tempfilename
'SW9PJL1D.wmf';tempfile-properties "XPR";}}We can think of this by looking
at the ways we can put particles in the container. There are four ways to
arrange the two particles in the two compartments. One of the four ways has
both particles in the left compartment.

I can keep adding particles. For three particles the probability of being in
the left side is $\left( \frac{1}{2}\right) ^{3}$ or $\frac{1}{8}.$ Here are
the eight ways we can put the particles in the two compartments. And only
one of those ways has all three particles on the left side. \FRAME{dhF}{%
2.7769in}{1.6042in}{0pt}{}{}{Figure}{\special{language "Scientific
Word";type "GRAPHIC";maintain-aspect-ratio TRUE;display "USEDEF";valid_file
"T";width 2.7769in;height 1.6042in;depth 0pt;original-width
4.1295in;original-height 2.373in;cropleft "0";croptop "1";cropright
"1";cropbottom "0";tempfilename 'SW9PJL1E.wmf';tempfile-properties "XPR";}}%
For $50$ particles the probability of all $50$ being on the left side would
be 
\begin{equation*}
\left( \frac{1}{2}\right) ^{50}=\allowbreak \frac{1}{1125\,899\,\allowbreak
906\,842\,624}
\end{equation*}

We can extend this thinking to divisions of particles. Suppose I have $100$
particles. We know that they will have a distribution of speeds
(Maxwell-Boltzman distribution). What is the probability that the $50$
fastest particles will be in the left side and the $50$ slowest will be on
the right side? This is the product of the last result with itself 
\begin{eqnarray*}
\left( \frac{1}{2}\right) ^{50}\left( \frac{1}{2}\right) ^{50} &=& \\
&=&\allowbreak \frac{1}{1267\,\allowbreak 650\,600\,228\,\allowbreak
229\,401\,496\,\allowbreak 703\,205\,376} \\
&=&\allowbreak 7.\,\allowbreak 888\,6\times 10^{-31}
\end{eqnarray*}%
This is not very probable. And we are only working with $100$ particles and
two sides of a box!

The point to this exercise is to realize there are very may \emph{states} or
ways to place particles in a divided box. Each configuration is equally
probable. Suppose we want a particular outcome, like the fastest $50$
molecules on one side and the slowest $50$ molecules on the other. There are
many many more ways to $\emph{not}$ have this configuration that there are
to have this configuration.

Each configuration of molecules and locations is called a \emph{microstate}.
There are a huge number of possible microstates in our example.

The collection of microstates is called a \emph{macrostate}. There is more
than one configuration of particles in our fast-slow division that places
all the fast molecules on one side and all the slow molecules on another
side. But the collection of all microstates that make up this division would
be a macrostate. For example, we could have an even distribution of
particles (within their fast/slow groups) on each side, or we could have all
the fastest particles in each group on the top, and the slowest on the
bottom within each major fast or slow side of the box. These two microstates
sill preserve the division between the $50$ fastest on the left and the
slowest $50$ on the right. They just rearrange the molecules within their
groups.

\section{Extension to large spaces and many particles}

Now let's take a room full of molecules. We count the number of microstates
by considering the number of places we can put a molecule. Each molecule has
a volume $%
%TCIMACRO{\TeXButton{V}{\ooalign{\hfil$V$\hfil\cr\kern0.1em--\hfil\cr}}}%
%BeginExpansion
\ooalign{\hfil$V$\hfil\cr\kern0.1em--\hfil\cr}%
%EndExpansion
_{m}.$ So there are about 
\begin{equation}
\emph{w=\,}\frac{%
%TCIMACRO{\TeXButton{V}{\ooalign{\hfil$V$\hfil\cr\kern0.1em--\hfil\cr}}}%
%BeginExpansion
\ooalign{\hfil$V$\hfil\cr\kern0.1em--\hfil\cr}%
%EndExpansion
}{%
%TCIMACRO{\TeXButton{V}{\ooalign{\hfil$V$\hfil\cr\kern0.1em--\hfil\cr}}}%
%BeginExpansion
\ooalign{\hfil$V$\hfil\cr\kern0.1em--\hfil\cr}%
%EndExpansion
_{m}}
\end{equation}%
places to put the molecule that are distinct.\footnote{%
It's really more than this if you tight pack the molecules, like in a
crystal structure. But this approximation is close enough for our study of
the second law of thermodynamics.} The $w$ is not work, nor is it weight
force. We are reusing a symbol (again). It is the number of places that we
can put a molecule.

Now suppose we have $N$ molecules. If we ignore the probability that we
might have two molecules in one space (which is different than the case
above!) we have 
\begin{equation}
\left( \emph{w}\right) ^{N}=\left( \frac{%
%TCIMACRO{\TeXButton{V}{\ooalign{\hfil$V$\hfil\cr\kern0.1em--\hfil\cr}}}%
%BeginExpansion
\ooalign{\hfil$V$\hfil\cr\kern0.1em--\hfil\cr}%
%EndExpansion
}{%
%TCIMACRO{\TeXButton{V}{\ooalign{\hfil$V$\hfil\cr\kern0.1em--\hfil\cr}}}%
%BeginExpansion
\ooalign{\hfil$V$\hfil\cr\kern0.1em--\hfil\cr}%
%EndExpansion
_{m}}\right) ^{N}
\end{equation}%
ways to place these molecules among our $\emph{w}$ states. I'll call this 
\begin{equation}
\emph{W}=\left( \frac{%
%TCIMACRO{\TeXButton{V}{\ooalign{\hfil$V$\hfil\cr\kern0.1em--\hfil\cr}}}%
%BeginExpansion
\ooalign{\hfil$V$\hfil\cr\kern0.1em--\hfil\cr}%
%EndExpansion
_{i}}{%
%TCIMACRO{\TeXButton{V}{\ooalign{\hfil$V$\hfil\cr\kern0.1em--\hfil\cr}}}%
%BeginExpansion
\ooalign{\hfil$V$\hfil\cr\kern0.1em--\hfil\cr}%
%EndExpansion
_{m}}\right) ^{N}
\end{equation}%
where $%
%TCIMACRO{\TeXButton{V}{\ooalign{\hfil$V$\hfil\cr\kern0.1em--\hfil\cr}}}%
%BeginExpansion
\ooalign{\hfil$V$\hfil\cr\kern0.1em--\hfil\cr}%
%EndExpansion
_{i}$ is a particular volume (we will let it change to $%
%TCIMACRO{\TeXButton{V}{\ooalign{\hfil$V$\hfil\cr\kern0.1em--\hfil\cr}}}%
%BeginExpansion
\ooalign{\hfil$V$\hfil\cr\kern0.1em--\hfil\cr}%
%EndExpansion
_{f}$ in just a minute).

\section{A Measure of Disorder}

If all the molecules were on the left side of a room, that would be an
unusually ordered situation. Normally the molecules are spread out through
the whole volume of a room, so it would be very unusual to have such a high
degree of order as to have all the molecules only on the left. Having the
molecules spread out across the whole room but with the fastest molecules on
the left and the slowest on the right is also an unusually ordered
situation, but it seems less ordered than having all the molecules on just
one side.

It would be convenient to find a way to measure the amount of disorder. That
is, we want to know how messed up the pattern of molecules in the room is.
Then we could see how a thermodynamic process affects the amount of
disorder. The traditional way to do this is to define a quantity%
\begin{equation}
S\equiv k_{B}\ln \left( W\right)
\end{equation}%
This has the property that $S$ increases as the number of states, $W,$
increases. The logarithm is convenient because $W$ is often a very big
number for real systems and the $\ln \left( W\right) $ is a much smaller
number to deal with.

Let's take an example, our old friend, the adiabatic free expansion.

\FRAME{dtbpF}{2.335in}{1.7487in}{0pt}{}{}{Figure}{\special{language
"Scientific Word";type "GRAPHIC";maintain-aspect-ratio TRUE;display
"USEDEF";valid_file "T";width 2.335in;height 1.7487in;depth
0pt;original-width 3.3589in;original-height 2.5088in;cropleft "0";croptop
"1";cropright "1";cropbottom "0";tempfilename
'SW9PJL1F.wmf';tempfile-properties "XPR";}}

For this case we start with all the atoms on one side of a membrane (our
super ordered situation). No energy transfers out by heat. No work is done,
we just puncture the membrane in the center and let the gas flow from one
side to the other. The gas molecules already have kinetic energy, so we
don't have to force them through the whole. The ones that happen to be
headed toward the whole just go through the hole.

This will mix up the molecules, changing the amount of order. How does the
disorder change in this process? Consider the number of states.%
\begin{equation*}
\emph{W}_{i}=\left( \frac{%
%TCIMACRO{\TeXButton{V}{\ooalign{\hfil$V$\hfil\cr\kern0.1em--\hfil\cr}}}%
%BeginExpansion
\ooalign{\hfil$V$\hfil\cr\kern0.1em--\hfil\cr}%
%EndExpansion
_{i}}{%
%TCIMACRO{\TeXButton{V}{\ooalign{\hfil$V$\hfil\cr\kern0.1em--\hfil\cr}}}%
%BeginExpansion
\ooalign{\hfil$V$\hfil\cr\kern0.1em--\hfil\cr}%
%EndExpansion
_{m}}\right) ^{N}
\end{equation*}%
\begin{equation*}
\emph{W}_{f}=\left( \frac{%
%TCIMACRO{\TeXButton{V}{\ooalign{\hfil$V$\hfil\cr\kern0.1em--\hfil\cr}}}%
%BeginExpansion
\ooalign{\hfil$V$\hfil\cr\kern0.1em--\hfil\cr}%
%EndExpansion
_{f}}{%
%TCIMACRO{\TeXButton{V}{\ooalign{\hfil$V$\hfil\cr\kern0.1em--\hfil\cr}}}%
%BeginExpansion
\ooalign{\hfil$V$\hfil\cr\kern0.1em--\hfil\cr}%
%EndExpansion
_{m}}\right) ^{N}
\end{equation*}

We can see that $%
%TCIMACRO{\TeXButton{V}{\ooalign{\hfil$V$\hfil\cr\kern0.1em--\hfil\cr}}}%
%BeginExpansion
\ooalign{\hfil$V$\hfil\cr\kern0.1em--\hfil\cr}%
%EndExpansion
_{f}$ is bigger than $%
%TCIMACRO{\TeXButton{V}{\ooalign{\hfil$V$\hfil\cr\kern0.1em--\hfil\cr}}}%
%BeginExpansion
\ooalign{\hfil$V$\hfil\cr\kern0.1em--\hfil\cr}%
%EndExpansion
_{i}.$ We can write our disorder measure as 
\begin{equation*}
S_{i}\equiv k_{B}\ln \left( \emph{W}_{i}\right)
\end{equation*}%
in the beginning and 
\begin{equation*}
S_{f}\equiv k_{B}\ln \left( \emph{W}_{f}\right)
\end{equation*}%
after we puncture the membrane. Then the change is order is 
\begin{eqnarray*}
\Delta S &=&S_{f}-S_{i} \\
&=&k_{B}\ln \left( \emph{W}_{f}\right) -k_{B}\ln \left( \emph{W}_{i}\right)
\\
&=&k_{B}\ln \left( \frac{\emph{W}_{f}}{\emph{W}_{i}}\right)
\end{eqnarray*}%
If we look at the ratio of the number of states%
\begin{eqnarray*}
\frac{\emph{W}_{f}}{\emph{W}_{i}} &=&\frac{\left( \frac{%
%TCIMACRO{\TeXButton{V}{\ooalign{\hfil$V$\hfil\cr\kern0.1em--\hfil\cr}}}%
%BeginExpansion
\ooalign{\hfil$V$\hfil\cr\kern0.1em--\hfil\cr}%
%EndExpansion
_{f}}{%
%TCIMACRO{\TeXButton{V}{\ooalign{\hfil$V$\hfil\cr\kern0.1em--\hfil\cr}}}%
%BeginExpansion
\ooalign{\hfil$V$\hfil\cr\kern0.1em--\hfil\cr}%
%EndExpansion
_{m}}\right) ^{N}}{\left( \frac{%
%TCIMACRO{\TeXButton{V}{\ooalign{\hfil$V$\hfil\cr\kern0.1em--\hfil\cr}}}%
%BeginExpansion
\ooalign{\hfil$V$\hfil\cr\kern0.1em--\hfil\cr}%
%EndExpansion
_{i}}{%
%TCIMACRO{\TeXButton{V}{\ooalign{\hfil$V$\hfil\cr\kern0.1em--\hfil\cr}}}%
%BeginExpansion
\ooalign{\hfil$V$\hfil\cr\kern0.1em--\hfil\cr}%
%EndExpansion
_{m}}\right) ^{N}} \\
&=&\left( \frac{%
%TCIMACRO{\TeXButton{V}{\ooalign{\hfil$V$\hfil\cr\kern0.1em--\hfil\cr}}}%
%BeginExpansion
\ooalign{\hfil$V$\hfil\cr\kern0.1em--\hfil\cr}%
%EndExpansion
_{f}}{%
%TCIMACRO{\TeXButton{V}{\ooalign{\hfil$V$\hfil\cr\kern0.1em--\hfil\cr}}}%
%BeginExpansion
\ooalign{\hfil$V$\hfil\cr\kern0.1em--\hfil\cr}%
%EndExpansion
_{i}}\right) ^{N}
\end{eqnarray*}%
so 
\begin{eqnarray}
\Delta S &=&k_{B}\ln \left( \left( \frac{%
%TCIMACRO{\TeXButton{V}{\ooalign{\hfil$V$\hfil\cr\kern0.1em--\hfil\cr}}}%
%BeginExpansion
\ooalign{\hfil$V$\hfil\cr\kern0.1em--\hfil\cr}%
%EndExpansion
_{f}}{%
%TCIMACRO{\TeXButton{V}{\ooalign{\hfil$V$\hfil\cr\kern0.1em--\hfil\cr}}}%
%BeginExpansion
\ooalign{\hfil$V$\hfil\cr\kern0.1em--\hfil\cr}%
%EndExpansion
_{i}}\right) ^{N}\right)  \notag \\
&=&Nk_{B}\ln \left( \left( \frac{%
%TCIMACRO{\TeXButton{V}{\ooalign{\hfil$V$\hfil\cr\kern0.1em--\hfil\cr}}}%
%BeginExpansion
\ooalign{\hfil$V$\hfil\cr\kern0.1em--\hfil\cr}%
%EndExpansion
_{f}}{%
%TCIMACRO{\TeXButton{V}{\ooalign{\hfil$V$\hfil\cr\kern0.1em--\hfil\cr}}}%
%BeginExpansion
\ooalign{\hfil$V$\hfil\cr\kern0.1em--\hfil\cr}%
%EndExpansion
_{i}}\right) \right)
\end{eqnarray}%
We should ask if the disorder increased or decreased. Since $%
%TCIMACRO{\TeXButton{V}{\ooalign{\hfil$V$\hfil\cr\kern0.1em--\hfil\cr}}}%
%BeginExpansion
\ooalign{\hfil$V$\hfil\cr\kern0.1em--\hfil\cr}%
%EndExpansion
_{f}$ is larger for our case than $%
%TCIMACRO{\TeXButton{V}{\ooalign{\hfil$V$\hfil\cr\kern0.1em--\hfil\cr}}}%
%BeginExpansion
\ooalign{\hfil$V$\hfil\cr\kern0.1em--\hfil\cr}%
%EndExpansion
_{i},$ then $\Delta S$ will be a positive number. This means the disorder
has increased in this case. That is just what we expected.

It is a good question to ask, will $\Delta S$ ever be negative? In other
words, will the gas ever line up on only one side every again? If the answer
is no, then this is an irreversible process. We already know that that this
specific process is irreversible. We would have to do work to pump the gas
back into one side. The gas won't spontaneously reorder itself.

Of course, we could set up a pump and pump our molecules back to one side of
the container. That would lower the disorder, making $\Delta S$ negative.
But now that we see that disorder is a function of molecules and their
energy, think of the molecules around the pump. The pump has waste heat $%
Q_{c}.$ That waste heat will add energy to the air molecules. That energy
will be kinetic energy and now we have stirred up the air around the pump.
We have increased the disorder for the pump environment. If we add $\Delta
S_{container}+\Delta S_{environment}$ we will find that the total change in
disorder for all systems is always positive or zero. 
\begin{equation*}
\Delta S\geq 0
\end{equation*}

\section{Definition of entropy}

We have learned to use one state variable, $\Delta E_{int}.$ The trick is to
realize that $\Delta E_{int}$ only depends on $\Delta T$ and use any easy
path to calculate it's value. That is why we can use 
\begin{equation*}
\Delta E_{int}=nC_{V}\Delta T
\end{equation*}%
no matter what process we have. State variables, like $\Delta E_{int}$ give
us powerful ways to calculate what will happen in a system.

There is another state variable in thermodynamics, \emph{change in entropy}, 
$\Delta S.$ We will again find that we can calculate $\Delta S$ by any path
because it depends only on the starting and ending points. And from our
previous discussion we recognize that this new state variable is just a
change in our disorder parameter! 
\begin{equation*}
\Delta S=k_{B}\ln \left( \frac{\emph{W}_{f}}{\emph{W}_{i}}\right)
\end{equation*}%
Note that this is a difference between two disorder parameters. We will call
our disorder parameter the \emph{entropy} of the system. 
\begin{equation}
S\equiv k_{B}\ln \left( W\right)
\end{equation}%
Entropy was originally defined at the macroscopic level. Statistical
Mechanics (the microscopic study of thermodynamics) gives us our modern
definition we have followed in this discussion. We can state this definition
as

\begin{center}
\rule{300pt}{1pt}

Entropy is a measure of the amount of disorder in a system and its
environment.

\rule{300pt}{1pt}
\end{center}

\subsection{Entropy and the Second Law}

In our example, because, $%
%TCIMACRO{\TeXButton{V}{\ooalign{\hfil$V$\hfil\cr\kern0.1em--\hfil\cr}}}%
%BeginExpansion
\ooalign{\hfil$V$\hfil\cr\kern0.1em--\hfil\cr}%
%EndExpansion
_{f}>%
%TCIMACRO{\TeXButton{V}{\ooalign{\hfil$V$\hfil\cr\kern0.1em--\hfil\cr}}}%
%BeginExpansion
\ooalign{\hfil$V$\hfil\cr\kern0.1em--\hfil\cr}%
%EndExpansion
_{i}$ we found that $\Delta S$ will be positive. This implies that the
disorder has increased. We can see that this is true, because there are now
molecules distributed among many more states (position, occupation
combinations) than before.

We will find this to be true in general, for irreversible processes (all
natural process that we know about). The total entropy of an isolated system
that undergoes a change cannot decrease.

What about non-isolated systems? We must then consider the entropy of the
system \emph{and} its surroundings. We again find that entropy of a system 
\emph{and} its surroundings cannot decrease. But the entropy of a system can
decrease if the entropy of the surroundings increases more so the total
effect is an increase in entropy (Think of the pump in our last example).

These statements constitute the most commonly seen form of the second law of
thermodynamics.

\begin{center}
\rule{300pt}{1pt}

Second Law of Thermodynamics: The entropy of a system and its surroundings
increases over time. 
\begin{equation*}
\Delta S\geq 0
\end{equation*}

\rule{300pt}{1pt}
\end{center}

\subsection{Macroscopic view}

We don't really want to count the microstates of a process and look for
changes in their number to do calculations. It would get tedious very fast!
And so far our definition of $\Delta S$ was done only for a free expansion.
We expect that in some way disorder would increase if we transferred in
energy by heat, or did some kind of work on the system. We need a
macroscopic version of our definition of entropy.

Let's look at the definition and then show that it works. The macroscopic
change in entropy is given by%
\begin{equation}
\Delta S=\frac{\Delta Q_{r}}{T}
\end{equation}%
where $Q_{r}$ is the energy transferred by heat through a reversible path
and $T$ is the absolute temperature (in $\unit{K}$) of the system. The
assumption of a single temperature is because we have assumed $\Delta Q_{r}$
is small. But why do we use a reversible path?

Suppose we have an irreversible process. It might be hard to compute $Q.$
But remember that, like $\Delta E_{int},$ $\Delta S$ does not depend on
path! It depends only on the initial and final states. So we may calculate $%
\Delta S$ using any convenient path, and a reversible path is more
convenient.

Note that macroscopically we have not defined entropy but the change in
entropy, $\Delta S$. This is the meaningful quantity (like $\Delta E_{int})$
for thermodynamic calculations.

Let's do an example: What is the change in entropy when a block of ice
melts? And before we do the math, let's think about it. The block of ice has
all the water molecules in one section of the room. That is highly ordered.
When the ice melts, the molecules can spread across the floor of the room.
That would be less ordered. The disorder increases, so the entropy
increases. Now the math.

Before we said that $\Delta Q_{r}$ was small enough that we could say that $%
T $ doesn't change while the energy $\Delta Q_{r}$ is transferred. We should
probably call such a small transfer of energy $dQ_{r}.$ Then a small change
in entropy would be 
\begin{equation*}
dS=\frac{dQ_{r}}{T}
\end{equation*}%
and we could add up all these small changes of entropy while the ice melts
to get the total entropy change for all the ice melting. 
\begin{equation*}
\Delta S=\int \frac{dQ_{r}}{T}
\end{equation*}%
Note that while the ice is melting, we will have ice water at a constant $%
T_{melt}=273\unit{K}.$ So our temperature really isn't changing! Then our
integral is 
\begin{equation*}
\Delta S=\frac{1}{T_{melt}}\int dQ_{r}
\end{equation*}%
and the integral is not too hard%
\begin{equation*}
\Delta S=\frac{Q_{r}}{T_{melt}}
\end{equation*}%
and for melting we know what $Q_{r}$ could be. We can freeze water as well
as melt ice. So we could model this as a reversible process. The amount of
energy transfer needed to melt the ice is just 
\begin{equation*}
Q=mL
\end{equation*}%
where $m$ is the mass of the ice and $L$ is the latent heat of fusion for
water. So 
\begin{equation*}
\Delta S=\frac{mL}{T_{melt}}
\end{equation*}

This wasn't hard at all!

Let's do another example to tie our macroscopic $\Delta S$ to our
microscopic $\Delta S.$ Let's find the change in entropy for a gas brought
through an isothermal expansion. Again we have%
\begin{equation*}
\Delta S=\int \frac{dQ_{r}}{T}
\end{equation*}%
and since the process is isothermal, $T$ is constant%
\begin{equation*}
\Delta S=\frac{1}{T}\int_{0}^{Q_{r}}dQ_{r}
\end{equation*}%
and the integral would seem easy. It is just $Q_{r},$ but it is a particular 
$Q_{r}$ for an isothermal process. For an isothermal process we know%
\begin{equation*}
\Delta E_{int}=0
\end{equation*}%
and 
\begin{equation*}
\Delta E_{int}=Q+w
\end{equation*}%
so 
\begin{equation*}
Q=-w
\end{equation*}%
and further we know that for an isothermal process 
\begin{eqnarray*}
w &=&-\int Pd%
%TCIMACRO{\TeXButton{V}{\ooalign{\hfil$V$\hfil\cr\kern0.1em--\hfil\cr}} }%
%BeginExpansion
\ooalign{\hfil$V$\hfil\cr\kern0.1em--\hfil\cr}
%EndExpansion
\\
&=&-nRT\ln \left( \frac{%
%TCIMACRO{\TeXButton{V}{\ooalign{\hfil$V$\hfil\cr\kern0.1em--\hfil\cr}}}%
%BeginExpansion
\ooalign{\hfil$V$\hfil\cr\kern0.1em--\hfil\cr}%
%EndExpansion
_{f}}{%
%TCIMACRO{\TeXButton{V}{\ooalign{\hfil$V$\hfil\cr\kern0.1em--\hfil\cr}}}%
%BeginExpansion
\ooalign{\hfil$V$\hfil\cr\kern0.1em--\hfil\cr}%
%EndExpansion
_{i}}\right)
\end{eqnarray*}%
so the $Q_{r}$ for our process could be just 
\begin{equation*}
Q=nRT\ln \left( \frac{%
%TCIMACRO{\TeXButton{V}{\ooalign{\hfil$V$\hfil\cr\kern0.1em--\hfil\cr}}}%
%BeginExpansion
\ooalign{\hfil$V$\hfil\cr\kern0.1em--\hfil\cr}%
%EndExpansion
_{f}}{%
%TCIMACRO{\TeXButton{V}{\ooalign{\hfil$V$\hfil\cr\kern0.1em--\hfil\cr}}}%
%BeginExpansion
\ooalign{\hfil$V$\hfil\cr\kern0.1em--\hfil\cr}%
%EndExpansion
_{i}}\right)
\end{equation*}%
Then 
\begin{eqnarray*}
\Delta S &=&\frac{1}{T}nRT\ln \left( \frac{%
%TCIMACRO{\TeXButton{V}{\ooalign{\hfil$V$\hfil\cr\kern0.1em--\hfil\cr}}}%
%BeginExpansion
\ooalign{\hfil$V$\hfil\cr\kern0.1em--\hfil\cr}%
%EndExpansion
_{f}}{%
%TCIMACRO{\TeXButton{V}{\ooalign{\hfil$V$\hfil\cr\kern0.1em--\hfil\cr}}}%
%BeginExpansion
\ooalign{\hfil$V$\hfil\cr\kern0.1em--\hfil\cr}%
%EndExpansion
_{i}}\right) \\
&=&nR\ln \left( \frac{%
%TCIMACRO{\TeXButton{V}{\ooalign{\hfil$V$\hfil\cr\kern0.1em--\hfil\cr}}}%
%BeginExpansion
\ooalign{\hfil$V$\hfil\cr\kern0.1em--\hfil\cr}%
%EndExpansion
_{f}}{%
%TCIMACRO{\TeXButton{V}{\ooalign{\hfil$V$\hfil\cr\kern0.1em--\hfil\cr}}}%
%BeginExpansion
\ooalign{\hfil$V$\hfil\cr\kern0.1em--\hfil\cr}%
%EndExpansion
_{i}}\right)
\end{eqnarray*}%
and remember that $nR=Nk_{B}$ so 
\begin{equation*}
\Delta S=Nk_{B}\ln \left( \frac{%
%TCIMACRO{\TeXButton{V}{\ooalign{\hfil$V$\hfil\cr\kern0.1em--\hfil\cr}}}%
%BeginExpansion
\ooalign{\hfil$V$\hfil\cr\kern0.1em--\hfil\cr}%
%EndExpansion
_{f}}{%
%TCIMACRO{\TeXButton{V}{\ooalign{\hfil$V$\hfil\cr\kern0.1em--\hfil\cr}}}%
%BeginExpansion
\ooalign{\hfil$V$\hfil\cr\kern0.1em--\hfil\cr}%
%EndExpansion
_{i}}\right)
\end{equation*}%
which is just what we found from our microscopic view of entropy change! Now
we realize that when we let the gas flow from one side of our divide box to
the other, we did this isothermally. And that is true, the gas did not
change the average kinetic energy per molecule as it flowed from one side of
the box to the other. Since our system was isolated, $Q=0$ and $w=0$ for the
whole system, so $\Delta E_{int}=0$ and the system did not change
temperature. But surely something changed! And that something was $\Delta S!$

\subsection{Heat death of the Universe}

We should ask the question, is the universe an isolated system? We don't
know the answer to this for sure\footnote{%
After all, Spiderman\ and his colleagues seem to be able to come and go from
different \textquotedblleft universes\textquotedblright\ all the time.}, but
we believe this is true that the universe is isolated.\footnote{%
Dispite what the television shows seem to say.} That means that the change
in entropy of the universe will always be positive. That means the universe
is \textquotedblleft running down.\textquotedblright\ The ultimate end of
this increase in energy would mean that all mater in the universe will be
spread out and cold. It will be disordered. This is called the \emph{heat
death of the universe}.

Note that this is a little different than the eschatological view that we
have in the Church of Jesus Christ of Latter Day Saints.

Throughout my career, many colleagues and acquaintances have said it is
foolish to believe in a church and a resurrection because the second law of
thermodynamics tells us that the universe will die in a cold blur. How do we
respond?

I can only give my opinion. But I think it is worth discussing, as an
example of how we think about both science and our faith in God, so here is
a disclaimer for what follows.

\begin{center}
\rule{300pt}{1pt}

WARNING: This is not an official statement of doctrine. This is an answer to
a question that many scientists must face from colleagues who will question
their faith using the principals of thermodynamics. Official statements by
the General Authorities supersede any opinions expressed here.

\rule{300pt}{1pt}
\end{center}

\subsection{The First Law and LDS\ Thought\protect\footnote{%
Comparing physical theory to LDS\ thought has it's dangers, I am aware. But
as scientists, we can't escape thinking about this. Just as a caution, John
A. Widsoe wrote a book called \emph{Joseph Smith as Scientist} in which
Widsoe tried very hard to show that LDS\ though is in harmony with Universal
Ether Theory. But we now know Universal Ether Theory is not correct, so
being in harmony with it is meaningless. We should be cautious!}}

I suppose we should ask, do we believe in the first law of thermodynamics?
If this law does not work for us, we have no reason to expect the second law
to work. The first law is basically conservation of energy. Here are some
quotes to consider.

\begin{quote}
And there stood one among them that was like unto God, and he said unto
those who were with him: We will go down, for there is space there, and we
will take of these materials, and we will make an earth whereon these may
dwell; (Abraham 3:24 )

You ask the learned doctors why the say the world was made out of nothing;
and they will answer, `Doesn't the Bible say He created the world?' And they
infer, from the word create, that it must have been made out of nothing.
Now, the word create came from the word baurau which does not mean create
out of nothing; it means to organize; the same as a man would organize
materials and build a ship. (Joseph Smith, King Follett Discourse)
\end{quote}

The idea of the first law of thermodynamics seems in harmony with LDS
thought.

So let's take on the Second Law.

\subsection{The Second Law and LDS\ Thought}

Again here are some quotes

\begin{quote}
Wherefore, the first judgment which came upon man must needs have remained
to an endless duration. And if so, this flesh must have laid down to rot and
to crumble to its mother earth, to rise no more. (2 Nephi 9:7)

And our spirits must have become like unto [Satan], and we become devils,
angels to a devil, to be shut out from the presence of our God, and to
remain with the father of lies, in misery, like unto himself. (2 Nephi 9:9)
\end{quote}

So our bodies fall apart, with the molecules spreading to more and more
states. And even our spirit seems to go occupy some kind of low energy state
away from God. This seems disordered. We seem to accept the concept of the
second law as well! --This will be a surprise to our non-member friends.

But then how can we expect eternal life and exaltation? Does something
counter the second law (a Third or Fourth Law of Thermodynamics)?

\subsubsection{An Additional Law of Thermodynamics?}

Again let's go to the scriptures. Jacob describes another force not
recognized by current scientific theory that counteracts this tendency to
decay

\begin{quote}
Wherefore, it must needs be an infinite atonement---save it should be an
infinite atonement this corruption could not put on incorruption. Wherefore,
the first judgment which came upon man must needs have remained to an
endless duration. And if so, this flesh must have laid down to rot and to
crumble to its mother earth, to rise no more. (2 Nephi 9:7-8)

O how great the goodness of our God, who prepareth a way for our escape from
the grasp of this awful monster; yea, that monster, death and hell, which I
call the death of the body, and also the death of the spirit. (2 Nephi 9:10)
\end{quote}

Atonement means to bring together. Like putting back all the atoms into on
side of our container after they have expanded, or like putting a body back
together after it decays, or putting a spirit back together with God. This
is an apt description of a process or law that would counter or balance the
second law\footnote{%
Nibley, The meaning of the temple, reprinted in the Collected Works of Hugh
Nibley: volume 12, Deseret Book, Salt Lake City, UT, 1992}. Since the
atonement can only be accessed through Christ, it is not likely to be
discovered in a laboratory, but it is a very real, physical, process. (see
for example Luke 24:39)

This is not official doctrine from the Church, but is a way to think about
the second law until we get direct revelation on the subject. We really
should not be worried about a conflict between thermodynamics and our
doctrine. I don't believe one exists, as you can see. It is totally OK\ for
us (our bodies and spirits) to have a lower entropy while the surrounding
universe has a larger entropy. The second law is not violated.

\section{Last Thoughts}

We are done with our introductory study of thermodynamics. Many of you will
have junior level thermodynamics classes, so you are likely to see more of
this. But we can stop for now. For physics majors the next physics class is
likely to be Principles of Physics III (PH220), electricity and magnetism. I
would prefer to call this class \textquotedblleft introduction to field
theory\textquotedblright\ because we really study electric and magnetic
fields for the first time. We will take up the topic of wave motion at the
end of Principles of Physics III to mathematically describe light as
electromagnetic waves again. And then we will describe everything using
waves in Principles of Physics IV - called \textquotedblleft Modern
Physics,\textquotedblright\ as our first introduction to quantum mechanics.

I hope you have enjoyed learning about waves, optics, and thermodynamics. I
hope you will see the world a little differently (and think about the
process of seeing the world!). For physics majors, there are still so many
cool things to learn. If you are leaving physics at this point, you can
continue to learn about the world and the universe by reading and watching
physics related material. If you are innovating new techniques or products,
you might even want to hire the services of a physicist to help with the
underlying principles for your new venture!

\appendix

\label{AppendixMarker}

\chapter{Experiment 1}

\section{Introduction}

We have learned that 
\begin{equation*}
\omega ^{2}=\frac{k}{m}
\end{equation*}%
and we know that 
\begin{equation*}
\omega =2\pi f
\end{equation*}%
and 
\begin{equation*}
T=\frac{1}{f}
\end{equation*}%
so we would predict that 
\begin{eqnarray}
T &=&\frac{1}{\frac{\omega }{2\pi }}=\frac{2\pi }{\omega }  \label{period} \\
&=&2\pi \sqrt{\frac{m}{k}}  \notag
\end{eqnarray}%
We would like to verify that this is true experimentally.

Unfortunately, springs don't come marked with their spring constants. You
will need to measure your spring constant. A way to do this is to hang your
spring motionless. When the mass is not moving Newton's second law tells us 
\begin{equation*}
\sum_{y}F=ma_{y}=F_{S}-F_{g}
\end{equation*}%
With no change in velocity, the acceleration must be zero 
\begin{equation*}
0=F_{S}-F_{g}
\end{equation*}%
so%
\begin{eqnarray*}
F_{S} &=&F_{g} \\
k\Delta y &=&mg
\end{eqnarray*}%
or%
\begin{equation}
k=\frac{mg}{\Delta y}  \label{konstant}
\end{equation}

We believe that the equation of motion for our mass-spring system would be 
\begin{eqnarray}
y\left( t\right) &=&y_{\max }\cos \left( \omega t+\phi _{o}\right)
\label{motion} \\
&=&y_{\max }\cos \left( \frac{2\pi }{T}t+\phi _{o}\right)  \notag
\end{eqnarray}

\section{Equipment}

We have springs, masses, and motion detectors. Your motion detector should
be placed below the mass as was demonstrated in class. Be careful to not
drop the mass on the motion detector.

The motion detectors connect to our laptop computers. Use the LoggerPro
program to collect data. Your lab instructor will help you set up the
program and show you how to collect data.

\subsection{Assignment}

\begin{enumerate}
\item Find your spring constant

\item Using equation (\ref{period}) predict the period of oscillation of
your mass and spring.

\item Suspend your mass from the spring and capture several cycles of the
motion using the motion detector.

\item Verify that equation (\ref{motion}) is a good model for our physical
situation by fitting a sinusoidal curve to the data (ask instructor for
help).

\item Using your curve fit equation, compare your predicted period to your
measured period. Do this by matching the fit equation to our theoretical
equation of motion 
\begin{eqnarray*}
Y &=&A\ast \cos (Bx+C)+D \\
y\left( t\right) &=&y_{\max }\cos \left( \frac{2\pi }{T}t+\phi _{o}\right)
\end{eqnarray*}%
we can see that 
\begin{equation*}
B=\frac{2\pi }{T}
\end{equation*}%
and solve for $T.$

\item Do your predicted and measured periods match?

\item If they don't match, consider your inputs in equation (\ref{period})
and try again.
\end{enumerate}

\chapter{Experiment 2}

\section{Introduction}

\FRAME{dhF}{4.7668in}{2.0513in}{0pt}{}{}{Figure}{\special{language
"Scientific Word";type "GRAPHIC";maintain-aspect-ratio TRUE;display
"USEDEF";valid_file "T";width 4.7668in;height 2.0513in;depth
0pt;original-width 4.7124in;original-height 2.0124in;cropleft "0";croptop
"1";cropright "1";cropbottom "0";tempfilename
'SUS73WCQ.wmf';tempfile-properties "XPR";}}We studied calorimetry earlier in
the semester. The idea behind calorimetry is just conservation of energy. A
calorimeter is supposed to be an insulated system, so energy can't leave by
heat.

Your job today is to measure the latent heat of fusion of ice/water.

Some helpful reminders are the specific heat equation%
\begin{equation*}
Q=m_{i}c_{i}\Delta T
\end{equation*}%
and the latent heat equation.%
\begin{eqnarray*}
Q &=&\pm L_{f}M\qquad \text{melt/freeze} \\
Q &=&\pm L_{v}M\qquad \text{boil/condense}
\end{eqnarray*}%
From our table we used in class, we know what the value should be.

\begin{equation*}
\begin{tabular}{|l|l|l|l|l|}
\hline
\textbf{Substance} & $T_{m}\left( \unit{%
%TCIMACRO{\U{2103}}%
%BeginExpansion
{}^{\circ}{\rm C}%
%EndExpansion
}\right) $ & $L_{f}\left( \frac{\unit{J}}{\unit{kg}}\right) $ & $T_{m}\left( 
\unit{%
%TCIMACRO{\U{2103}}%
%BeginExpansion
{}^{\circ}{\rm C}%
%EndExpansion
}\right) $ & $L_{v}\left( \frac{\unit{J}}{\unit{kg}}\right) $ \\ \hline
Nitrogen $\left( N_{2}\right) $ & $-210$ & $0.26\times 10^{5}$ & $-196$ & $%
1.99\times 10^{5}$ \\ \hline
Ethyl alcohol & $-114$ & $1.09\times 10^{5}$ & $78$ & $8.79\times 10^{5}$ \\ 
\hline
Mercury & $-39$ & $0.11\times 10^{5}$ & $357$ & $2.96\times 10^{5}$ \\ \hline
Water & $0$ & $3.33\times 10^{5}$ & $100$ & $22.6\times 10^{5}$ \\ \hline
Lead & $328$ & $0.25\times 10^{5}$ & $1750$ & $8.58\times 10^{5}$ \\ \hline
\end{tabular}%
\end{equation*}%
Find the latent heat of fusion for water/ice and calculate a percent error
to see how well you did%
\begin{equation*}
\%Error=\frac{measured-acutal}{actual}
\end{equation*}

\section{Equipment}

I will have ice water available in a tub. I will also have calorimeter cups,
thermometers, and triple beam balances for finding mass. Water will be
available through our store room. Let me know if you need something else.

\section{Assignment}

Measure the latent heat of fusion of water.

\begin{enumerate}
\item Place some ice in water to form a slush (this will give you ice at $0%
\unit{%
%TCIMACRO{\U{2103}}%
%BeginExpansion
{}^{\circ}{\rm C}%
%EndExpansion
}$).

\item Place some water into the calorimeter cup, and measure its temperature
and mass.

\item Place a piece of ice into the water in the calorimeter. You will need
to find a way to measure its mass.

\item Measure the temperature change of the water after the ice has melted.

\item Calculate the latent heat, and compare to the accepted value.

\item Turn in your answer, $\%Error,$ and all your group names on a single
sheet of paper (the lab instructions sheet is fine)
\end{enumerate}

\chapter{Experiment 3}

\section{Introduction}

A converging lens is one which is thicker at the center than at the edge,
and which converges incident parallel rays to a real focus on the opposite
side of the lens from the object. A diverging lens is thinner at the center
than at the edge and diverges the light from a virtual focus on the same
side of the lens as the object. The optic axis of a lens is a line drawn
through the center of the lens perpendicular to the face of the lens. The
principal focal point is a point on the principal axis through which
incident rays parallel to the principal axis pass, or appear to pass, after
refraction by the lens. The focal length, $f$, of the lens is the distance
from the optical center of the lens to the principal focal point. The
reciprocal of the focal length in meters is called the power of the lens and
is expressed in diopters (d).

\begin{equation*}
P=\frac{1}{f\left( \unit{m}\right) }
\end{equation*}

If $s$ is the distance of the object from the optical center of the lens, $%
s^{\prime }$ the distance of the image from the lens, and $f$ the focal
length, the thin lens equation is%
\begin{equation}
\frac{1}{s}+\frac{1}{s^{\prime }}=\frac{1}{f}  \label{lens}
\end{equation}

The magnification is given by

\begin{equation}
m=-\frac{s^{\prime }}{s}  \label{magnification}
\end{equation}

For the case where more than one lens is present, Equation (\ref%
{magnification}) is applied successively to each lens, treating the image
formed by one lens as the object for the next lens, keeping in mind that the
distances in Equation (\ref{magnification}) are measured with respect to
each individual lens, and not a single fixed point.

\section{Equipment}

We have $2\unit{m}$ sticks, lens and holders, paper screens, and light
sources.

\section{Assignment}

Verify the lens equation (\ref{lens}), and the image magnification equation (%
\ref{magnification}), for the lens with the longer focal length. I suggest
you use at least three object distances greater than two focal lengths. I
suggest you measure the image distance and the image height. I suggest you
also measure at least three object distances between one and two focal
lengths. Measure the image distance and the image height.

Plot your magnification results. Interpret the results, does the lens
equation work?

Compute the magnification in each case and verify that equation (\ref%
{magnification}) holds.

The lens equation is supposed to be good for virtual images as well as real
images. We need to verify the equation for this case as well. I suggest you
do this by measuring at least one object distance closer to the lens than
one focal length, locate and measure the image distance. Measure the image
height. Verify that equations (\ref{lens}) and (\ref{magnification}) are
valid for this case. Include this case on your graph.

\QTP{Appendix}
Problem Types

Simple Harmonic Motion (SHM)

\begin{eqnarray*}
x\left( t\right) &=&x_{\max }\cos \left( \omega t+\phi _{o}\right) \\
v\left( t\right) &=&-\omega x_{\max }\sin \left( \omega t+\phi _{o}\right) \\
a\left( t\right) &=&-\omega ^{2}x_{\max }\cos \left( \omega t+\phi
_{o}\right)
\end{eqnarray*}%
\begin{equation*}
\omega =2\pi f
\end{equation*}%
\begin{equation*}
f=\frac{1}{T}
\end{equation*}%
\begin{equation*}
v_{\max }=\omega x_{\max }
\end{equation*}%
\begin{equation*}
a_{\max }=\omega ^{2}x_{\max }
\end{equation*}

Energy in SHM%
\begin{eqnarray*}
K &=&\frac{1}{2}mv^{2} \\
&=&\frac{1}{2}kx_{\max }^{2}\sin ^{2}\left( \omega t+\phi \right)
\end{eqnarray*}

\begin{eqnarray*}
U_{s} &=&\frac{1}{2}kx^{2} \\
&=&\frac{1}{2}kx_{\max }^{2}\cos ^{2}\left( \omega t+\phi \right)
\end{eqnarray*}

Angular Frequency and oscillation

For a mass sprig system%
\begin{equation*}
\omega =\sqrt{\frac{k}{m}}
\end{equation*}

for a pendulum%
\begin{equation*}
\omega =\sqrt{\frac{g}{L}}
\end{equation*}%
for a physical pendulum 
\begin{equation*}
\omega =\sqrt{\frac{mgd}{\mathbb{I}}}
\end{equation*}

Damped oscillation

\begin{equation*}
x\left( t\right) =x_{\max }e^{-\frac{b}{2m}t}\cos \left( \omega t+\phi
_{o}\right)
\end{equation*}%
\begin{equation*}
\omega =\sqrt{\frac{k}{m}-\left( \frac{b}{2m}\right) ^{2}}
\end{equation*}%
\begin{equation*}
x\left( t\right) =x_{\max }e^{-\frac{b}{2m}t}\cos \left( \left( \sqrt{\frac{k%
}{m}-\left( \frac{b}{2m}\right) ^{2}}\right) t+\phi _{o}\right)
\end{equation*}

Forced oscillation

\begin{equation*}
\mathbf{F}_{d}=-b\mathbf{v}
\end{equation*}%
\begin{equation*}
F\left( t\right) =F_{o}\sin \left( \omega _{f}t\right)
\end{equation*}%
\begin{equation*}
\Sigma F=F_{o}\sin \left( \omega _{f}t\right) -kx-bv_{x}=ma
\end{equation*}

\begin{equation*}
x\left( t\right) =A\cos \left( \omega _{f}t+\phi \right)
\end{equation*}%
\begin{equation*}
A=\frac{\frac{F_{o}}{m}}{\sqrt{\left( \omega _{f}^{2}-\omega _{o}^{2}\right)
^{2}+\left( \frac{b\omega _{f}}{m}\right) ^{2}}}
\end{equation*}%
\begin{equation*}
\omega _{o}=\sqrt{\frac{k}{m}}
\end{equation*}

Waves

\begin{equation*}
y\left( t\right) =A\cos \left( ax-\omega t+\phi \right)
\end{equation*}%
\begin{equation*}
k=\frac{2\pi }{\lambda }
\end{equation*}%
\begin{equation*}
v=\frac{\lambda }{T}
\end{equation*}%
\begin{equation*}
v=\lambda f
\end{equation*}

Waves on strings

\begin{equation*}
v=\sqrt{\frac{T_{s}}{\mu }}
\end{equation*}%
\begin{equation*}
\mu =\frac{m}{L}
\end{equation*}

Sound waves

\begin{eqnarray*}
v_{sound} &=&v_{o}\sqrt{1+\frac{T_{c}}{T_{o}}} \\
v_{o} &=&331\frac{\unit{m}}{\unit{s}} \\
T_{o} &=&273
\end{eqnarray*}

\begin{equation*}
P\equiv \frac{F}{A}
\end{equation*}%
\begin{equation*}
s\left( x,t\right) =s_{\max }\cos \left( kx-\omega t+\phi _{o}\right)
\end{equation*}

Spherical waves

\begin{equation*}
y=A\left( r\right) \sin \left( kr-\omega t+\phi _{o}\right)
\end{equation*}

Light waves

\begin{equation*}
n=\frac{c}{v}
\end{equation*}%
\begin{equation*}
c=299792458\frac{\unit{m}}{\unit{s}}
\end{equation*}%
Power and Intensity%
\begin{equation*}
\mathcal{P}=\frac{\Delta E}{\Delta t}
\end{equation*}%
\begin{equation*}
\mathcal{I}\equiv \frac{\mathcal{P}}{A}
\end{equation*}

\begin{equation*}
A=4\pi r^{2}
\end{equation*}

Sound Intensity Level

\begin{equation*}
\beta =10\log _{10}\left( \frac{I}{I_{o}}\right)
\end{equation*}%
\begin{equation*}
I_{o}=1\times 10^{-12}\frac{\unit{W}}{\unit{m}^{2}}
\end{equation*}

Doppler Effect

\begin{equation*}
f^{\prime }=\frac{v+v_{d}}{v}f\text{ \qquad observer moving toward the source%
}
\end{equation*}%
\begin{equation*}
f^{\prime }=\frac{v-v_{d}}{v}f\text{ \qquad observer moving away from the
source}
\end{equation*}%
\begin{equation*}
f^{\prime }=\frac{v}{v-v_{s}}f\qquad \text{Source moving toward observer}
\end{equation*}%
\begin{equation*}
f^{\prime }=\frac{v}{v+v_{s}}f\qquad \text{Source moving away from observer}
\end{equation*}%
\begin{equation*}
f^{\prime }=\frac{v\pm v_{d}}{v\mp v_{s}}f
\end{equation*}

Superposition

\begin{equation*}
\sin a+\sin b=2\cos \left( \frac{a-b}{2}\right) \sin \left( \frac{a+b}{2}%
\right)
\end{equation*}%
\begin{eqnarray*}
y_{r} &=&y_{\max }\sin \left( kx-\omega t\right) +A\sin \left( kx-\omega
t+\phi _{o}\right) \\
&=&2y_{\max }\cos \left( \frac{\phi _{o}}{2}\right) \sin \left( kx-\omega t+%
\frac{\phi _{o}}{2}\right)
\end{eqnarray*}%
Standing waves%
\begin{equation*}
\sin \left( a\pm b\right) =\sin \left( a\right) \cos \left( b\right) \pm
\cos \left( a\right) \sin \left( b\right)
\end{equation*}%
\begin{eqnarray*}
y_{r} &=&y_{\max }\sin \left( kx-\omega t\right) +y_{\max }\sin \left(
kx+\omega t\right) \\
&=&\left( 2y_{\max }\sin \left( kx\right) \right) \cos \left( \omega t\right)
\end{eqnarray*}%
\begin{equation*}
x=n\frac{\lambda }{4}
\end{equation*}%
\begin{eqnarray*}
f_{n} &=&\frac{v}{\lambda _{n}}= \\
&=&\frac{n}{2L}v \\
&=&\frac{n}{2L}\sqrt{\frac{T}{\mu }}
\end{eqnarray*}

Standing waves in pipes

Open on both ends

\begin{eqnarray*}
\lambda _{n} &=&\frac{2}{n}L \\
f_{n} &=&n\frac{v}{2L}\qquad n=1,2,3,4\ldots
\end{eqnarray*}%
Closed on one end%
\begin{eqnarray*}
\lambda _{n} &=&\frac{4}{n}L \\
f_{n} &=&n\frac{v}{4L}\qquad n=1,3,5\ldots
\end{eqnarray*}

\backmatter

\chapter{Afterword}

The back matter often includes one or more of an index, an afterword,
acknowledgements, a bibliography, a colophon, or any other similar item. In
the back matter, chapters do not produce a chapter number, but they are
entered in the table of contents. If you are not using anything in the back
matter, you can delete the back matter 
%TCIMACRO{\TeXButton{TeX}{\TeX{}} }%
%BeginExpansion
\TeX{}
%EndExpansion
field and everything that follows it.

\end{document}
